% Generated by the blueprint tool (blueprint-pdf/1) from a blueprint snapshot; edit the
% Markdown sources, not this file.  Compile with `latexmk -pdf` or by running
% `pdflatex` twice (the second run settles the contents and the page references).
% Images, if any, are read from assets/ next to this file.
\documentclass[11pt]{article}
\usepackage[T1]{fontenc}
\usepackage[utf8]{inputenc}
\usepackage{lmodern}
\usepackage{textcomp}
\usepackage[a4paper,margin=2.5cm,headheight=14pt]{geometry}
\usepackage{amsmath,amssymb,amsthm}
\usepackage{amscd}  % KaTeX's CD diagrams; part of amsmath, so always there
\usepackage{mathtools}
\usepackage{bm}
\usepackage{graphicx}
\usepackage{xcolor}
\usepackage{array,booktabs,tabularx}
\usepackage{fancyvrb}
\usepackage[normalem]{ulem}
\usepackage{microtype}
% Optional, and only with their fonts: a small TeX installation has the
% packages without the fonts, which fails at the first formula.
\IfFileExists{mathrsfs.sty}{\IfFileExists{rsfs10.tfm}{\usepackage{mathrsfs}}{}}{}
\IfFileExists{stmaryrd.sty}{\IfFileExists{stmary10.tfm}{\usepackage{stmaryrd}}{}}{}
\IfFileExists{cancel.sty}{\usepackage{cancel}}{}
\usepackage[hyphens]{url}
\usepackage[pdfusetitle,bookmarksnumbered=false]{hyperref}
\usepackage{bookmark}
\definecolor{bplink}{HTML}{1F4E9C}
\definecolor{bpmuted}{HTML}{5F676E}
\hypersetup{colorlinks=true,linkcolor=bplink,urlcolor=bplink,citecolor=bplink}
\setlength{\parindent}{0pt}
\setlength{\parskip}{0.55em plus 0.2em minus 0.1em}
% Lists as tight as the page's.
\makeatletter
\def\@listi{\leftmargin\leftmargini\parsep 0pt\topsep 0.3em plus 0.1em\itemsep 0.1em plus 0.1em}
\let\@listI\@listi
\def\@listii{\leftmargin\leftmarginii\labelwidth\leftmarginii\advance\labelwidth-\labelsep\parsep 0pt\topsep 0.1em\itemsep 0.1em}
% An ordered list that does not start at 1.
\newcommand{\bpstartat}[1]{\setcounter{\@enumctr}{\numexpr#1-1\relax}}
\makeatother
\setcounter{tocdepth}{3}
\emergencystretch=3em
\allowdisplaybreaks

% A numbered statement, set as a paper sets one: "Theorem 2.1 (Title)."
% #1 is the lead ("Theorem~2.1"), #2 what follows it in roman.
\newenvironment{bpstatement}[2]{\par\addvspace{0.6\baselineskip}\noindent\textbf{#1}#2\textbf{.}\enspace\ignorespaces}
  {\par\addvspace{0.4\baselineskip}}
\newenvironment{bpstatementit}[2]{\begin{bpstatement}{#1}{#2}\itshape\ignorespaces}{\end{bpstatement}}
% A step: a relation with prose of its own, read under the object it starts from.
\newenvironment{bpstep}[1]{\par\addvspace{0.3\baselineskip}\begin{list}{}{\leftmargin=1.6em\rightmargin=0pt\topsep=0pt}\item\relax
  {\small\color{bpmuted}#1}\par\nopagebreak}{\end{list}\par}
% Notes about where an object is written out.
\newenvironment{bpnote}{\par\small\color{bpmuted}\itshape}{\par}
\newcommand{\bpomitted}[1]{\par{\small\color{bpmuted}\itshape #1\par}}
\newcommand{\bpbodyheading}[1]{\par\addvspace{0.4\baselineskip}\noindent\textbf{#1}\par\nopagebreak}
\newcommand{\bprule}{\par\medskip\noindent\hfil\rule{0.3\linewidth}{0.4pt}\hfil\par\medskip}
\newcommand{\bpboundary}[1]{\par\nopagebreak{\small\color{bpmuted}#1\par}\nopagebreak}
% A reference to something this export leaves out, and one that resolves nowhere.
\newcommand{\bpout}[1]{#1\textsuperscript{\hyperlink{bp-outside}{\ensuremath{\ast}}}}
\newcommand{\bpbroken}[1]{\textcolor{red!70!black}{[#1?]}}
\newcommand{\bpslug}[1]{\texttt{#1}}
% A bibliography entry: label, anchor, entry; the label hangs in the margin.
\newcommand{\bpbibitem}[3]{\par\setbox0\hbox{[#1]\enspace}\ifdim\wd0<5em\setbox0\hbox to 5em{[#1]\hfil}\fi
  \hangindent=\wd0 \hangafter=1 \noindent\box0\phantomsection\label{#2}#3\par\smallskip}
% A character the export knows no way to set.
\newcommand{\bpglyph}[1]{\fbox{\scriptsize #1}}
% Images, at their own size unless they are wider than the line.
\newsavebox{\bpimagebox}
\newcommand{\bpimage}[1]{\sbox{\bpimagebox}{\includegraphics{#1}}%
  \ifdim\wd\bpimagebox>\linewidth\resizebox{\linewidth}{!}{\usebox{\bpimagebox}}\else\usebox{\bpimagebox}\fi}
\newcommand{\bpmissingimage}[1]{\fbox{\parbox{0.8\linewidth}{\centering\small\color{bpmuted}Image not included: #1}}}
\DefineVerbatimEnvironment{bpverbatim}{Verbatim}{fontsize=\small,xleftmargin=1em}
\renewcommand{\qedsymbol}{\ensuremath{\blacksquare}}
\hypersetup{pdftitle={Weil Conjectures — Part Ⅰ The six operations},pdfsubject={Part Ⅰ The six operations; all levels},pdfcreator={blueprint}}
% KaTeX commands the website knows
\providecommand{\R}{\mathbb{R}}\providecommand{\N}{\mathbb{N}}\providecommand{\Z}{\mathbb{Z}}
\providecommand{\Q}{\mathbb{Q}}\providecommand{\C}{\mathbb{C}}
\providecommand{\Reals}{\mathbb{R}}\providecommand{\reals}{\mathbb{R}}\providecommand{\Complex}{\mathbb{C}}
\providecommand{\cnums}{\mathbb{C}}\providecommand{\natnums}{\mathbb{N}}\providecommand{\Bbb}[1]{\mathbb{#1}}
\providecommand{\bold}[1]{\mathbf{#1}}\providecommand{\frak}[1]{\mathfrak{#1}}\providecommand{\mathbbm}[1]{\mathbb{#1}}
\providecommand{\lang}{\langle}\providecommand{\rang}{\rangle}\providecommand{\bra}[1]{\langle #1\rvert}
\providecommand{\ket}[1]{\lvert #1\rangle}\providecommand{\braket}[1]{\langle #1\rangle}
\providecommand{\set}[1]{\{#1\}}\providecommand{\Set}[1]{\left\{#1\right\}}
\providecommand{\infin}{\infty}\providecommand{\isin}{\in}\providecommand{\plusmn}{\pm}\providecommand{\sub}{\subset}
\providecommand{\sube}{\subseteq}\providecommand{\supe}{\supseteq}\providecommand{\exist}{\exists}
\providecommand{\bull}{\bullet}\providecommand{\larr}{\leftarrow}\providecommand{\rarr}{\rightarrow}
\providecommand{\lrarr}{\leftrightarrow}\providecommand{\Larr}{\Leftarrow}\providecommand{\Rarr}{\Rightarrow}
\providecommand{\Lrarr}{\Leftrightarrow}\providecommand{\uarr}{\uparrow}\providecommand{\darr}{\downarrow}
\providecommand{\harr}{\leftrightarrow}\providecommand{\lArr}{\Leftarrow}\providecommand{\rArr}{\Rightarrow}
\providecommand{\hArr}{\Leftrightarrow}\providecommand{\alef}{\aleph}\providecommand{\alefsym}{\aleph}
\providecommand{\real}{\Re}\providecommand{\image}{\Im}\providecommand{\weierp}{\wp}
\providecommand{\Alpha}{\mathrm{A}}\providecommand{\Beta}{\mathrm{B}}\providecommand{\Epsilon}{\mathrm{E}}
\providecommand{\Zeta}{\mathrm{Z}}\providecommand{\Eta}{\mathrm{H}}\providecommand{\Iota}{\mathrm{I}}
\providecommand{\Kappa}{\mathrm{K}}\providecommand{\Mu}{\mathrm{M}}\providecommand{\Nu}{\mathrm{N}}
\providecommand{\Omicron}{\mathrm{O}}\providecommand{\Rho}{\mathrm{P}}\providecommand{\Tau}{\mathrm{T}}
\providecommand{\Chi}{\mathrm{X}}\providecommand{\omicron}{o}\providecommand{\gt}{>}\providecommand{\lt}{<}
\providecommand{\And}{\&}\providecommand{\hdashline}{\hline}\providecommand{\operatornamewithlimits}[1]{\operatorname*{#1}}
\providecommand{\llbracket}{[\![}\providecommand{\rrbracket}{]\!]}\providecommand{\cancel}[1]{#1}
\providecommand{\mathscr}{\mathcal}
% the characters beyond ASCII this document uses
\DeclareUnicodeCharacter{2160}{\textup{I}}
\DeclareUnicodeCharacter{2161}{\textup{II}}
\DeclareUnicodeCharacter{2162}{\textup{III}}
% the project's macros ([katex.macros] in blueprint.toml)
\providecommand{\A}{}\renewcommand{\A}{\mathbf A}
\providecommand{\AEn}{}\renewcommand{\AEn}{A_{E,n}}
\providecommand{\Aaf}{}\renewcommand{\Aaf}{\mathbf A}
\providecommand{\Ab}{}\renewcommand{\Ab}{\mathsf{Ab}}
\providecommand{\AffSch}{}\renewcommand{\AffSch}{\operatorname{AffSch}}
\providecommand{\Aone}{}\renewcommand{\Aone}{\mathbf A^1}
\providecommand{\Bl}{}\renewcommand{\Bl}{\operatorname{Bl}}
\providecommand{\Br}{}\renewcommand{\Br}{\mathrm{Br}}
\providecommand{\C}{}\renewcommand{\C}{\mathbb C}
\providecommand{\Cech}{}\renewcommand{\Cech}{\check{H}}
\providecommand{\Coext}{}\renewcommand{\Coext}{\operatorname{Coext}}
\providecommand{\Cone}{}\renewcommand{\Cone}{\operatorname{Cone}}
\providecommand{\Cons}{}\renewcommand{\Cons}{\operatorname{Cons}}
\providecommand{\Cref}{}\renewcommand{\Cref}[1]{\text{#1}}
\providecommand{\D}{}\renewcommand{\D}{D}
\providecommand{\DLAFT}{}\renewcommand{\DLAFT}{\operatorname{FT}}
\providecommand{\DLAFr}{}\renewcommand{\DLAFr}{\operatorname{Fr}}
\providecommand{\DLAQ}{}\renewcommand{\DLAQ}{{\overline{\mathbf Q}_\ell}}
\providecommand{\DLRQ}{}\renewcommand{\DLRQ}{{\overline{\mathbf Q}_\ell}}
\providecommand{\Dbc}{}\renewcommand{\Dbc}{D^b_c}
\providecommand{\Dc}{}\renewcommand{\Dc}{D_{\mathrm{cons}}}
\providecommand{\Dctf}{}\renewcommand{\Dctf}{D^b_{ctf}}
\providecommand{\Dic}{}\renewcommand{\Dic}{D_{\IndCons}}
\providecommand{\Dplus}{}\renewcommand{\Dplus}{D^{+}}
\providecommand{\End}{}\renewcommand{\End}{\operatorname{End}}
\providecommand{\Et}{}\renewcommand{\Et}{\mathrm{\acute{e}t}}
\providecommand{\Ex}{}\renewcommand{\Ex}{\mathrm{Ex}}
\providecommand{\Ext}{}\renewcommand{\Ext}{\operatorname{Ext}}
\providecommand{\F}{}\renewcommand{\F}{\mathbb F}
\providecommand{\FT}{}\renewcommand{\FT}{\operatorname{FT}}
\providecommand{\Fbar}{}\renewcommand{\Fbar}{{\overline{\mathbf F}_q}}
\providecommand{\Fell}{}\renewcommand{\Fell}{\mathbf F_\ell}
\providecommand{\Fib}{}\renewcommand{\Fib}{\operatorname{Fib}}
\providecommand{\Fix}{}\renewcommand{\Fix}{\operatorname{Fix}}
\providecommand{\Fp}{}\renewcommand{\Fp}{{\mathbf F_{p}}}
\providecommand{\Fq}{}\renewcommand{\Fq}{{\mathbf F_q}}
\providecommand{\Fqn}{}\renewcommand{\Fqn}{{\mathbf F_{q^{n}}}}
\providecommand{\Frob}{}\renewcommand{\Frob}{\operatorname{Frob}}
\providecommand{\GL}{}\renewcommand{\GL}{\operatorname{GL}}
\providecommand{\Gal}{}\renewcommand{\Gal}{\operatorname{Gal}}
\providecommand{\Hens}{}\renewcommand{\Hens}{\operatorname{Hens}}
\providecommand{\Hilb}{}\renewcommand{\Hilb}{\operatorname{Hilb}}
\providecommand{\Hom}{}\renewcommand{\Hom}{\operatorname{Hom}}
\providecommand{\IndCons}{}\renewcommand{\IndCons}{\operatorname{IndCons}}
\providecommand{\Kar}{}\renewcommand{\Kar}{\operatorname{Kar}}
\providecommand{\LT}{}\renewcommand{\LT}{\operatorname{LT}}
\providecommand{\Lie}{}\renewcommand{\Lie}{\operatorname{Lie}}
\providecommand{\Loc}{}\renewcommand{\Loc}{\operatorname{Loc}}
\providecommand{\Mod}{}\renewcommand{\Mod}{\mathsf{Mod}}
\providecommand{\N}{}\renewcommand{\N}{\mathbb{N}}
\providecommand{\Needspace}{}\renewcommand{\Needspace}{\par}
\providecommand{\O}{}\renewcommand{\O}{{\mathcal{O}}}
\providecommand{\Ob}{}\renewcommand{\Ob}{\mathrm{Ob}}
\providecommand{\Oe}{}\renewcommand{\Oe}{\mathcal O_E}
\providecommand{\PP}{}\renewcommand{\PP}{\mathbf{P}}
\providecommand{\Perf}{}\renewcommand{\Perf}{\operatorname{Perf}}
\providecommand{\Pic}{}\renewcommand{\Pic}{\mathrm{Pic}}
\providecommand{\Pm}{}\renewcommand{\Pm}{\mathbf P^{m}}
\providecommand{\Pone}{}\renewcommand{\Pone}{\mathbf P^{1}}
\providecommand{\Pqn}{}\renewcommand{\Pqn}{\mathbf P^{n}}
\providecommand{\ProFin}{}\renewcommand{\ProFin}{\operatorname{ProFin}}
\providecommand{\Proet}{}\renewcommand{\Proet}{\mathrm{pro\acute{e}t}}
\providecommand{\Q}{}\renewcommand{\Q}{\mathbb Q}
\providecommand{\Qb}{}\renewcommand{\Qb}{\overline{\mathbf Q}_{\ell}}
\providecommand{\Qbar}{}\renewcommand{\Qbar}{\overline{\mathbf Q}}
\providecommand{\QbarX}{}\renewcommand{\QbarX}{\overline{\mathbf Q}_{\ell,X}}
\providecommand{\Qbarl}{}\renewcommand{\Qbarl}{{\overline{\mathbf Q}_{\ell}}}
\providecommand{\Qell}{}\renewcommand{\Qell}{{\mathbf{Q}_{\ell}}}
\providecommand{\Ql}{}\renewcommand{\Ql}{\mathbf Q_{\ell}}
\providecommand{\Qlb}{}\renewcommand{\Qlb}{{\overline{\mathbf Q}_\ell}}
\providecommand{\Qlbar}{}\renewcommand{\Qlbar}{\overline{\mathbf Q}_\ell}
\providecommand{\R}{}\renewcommand{\R}{\mathbb{R}}
\providecommand{\RG}{}\renewcommand{\RG}{\operatorname{\mathrm{R}\Gamma}}
\providecommand{\RGamma}{}\renewcommand{\RGamma}{R\Gamma}
\providecommand{\RGammac}{}\renewcommand{\RGammac}{R\Gamma_c}
\providecommand{\RHom}{}\renewcommand{\RHom}{R\mathcal{H}om}
\providecommand{\RHomX}{}\renewcommand{\RHomX}{R\mathcal H om_X}
\providecommand{\Rat}{}\renewcommand{\Rat}{\operatorname{IsRat}}
\providecommand{\Res}{}\renewcommand{\Res}{\operatorname{Res}}
\providecommand{\Rlim}{}\renewcommand{\Rlim}{R\!\varprojlim}
\providecommand{\SIm}{}\renewcommand{\SIm}{\operatorname{SIm}}
\providecommand{\Sch}{}\renewcommand{\Sch}{\mathsf{Sch}}
\providecommand{\Sh}{}\renewcommand{\Sh}{\operatorname{Sh}}
\providecommand{\Span}{}\renewcommand{\Span}{\operatorname{span}}
\providecommand{\Spec}{}\renewcommand{\Spec}{\operatorname{Spec}}
\providecommand{\Tr}{}\renewcommand{\Tr}{\operatorname{Tr}}
\providecommand{\Trsm}{}\renewcommand{\Trsm}{\operatorname{Tr}^{\mathrm{sm}}}
\providecommand{\Wt}{}\renewcommand{\Wt}{\operatorname{Wt}}
\providecommand{\Z}{}\renewcommand{\Z}{\mathbb{Z}}
\providecommand{\Zar}{}\renewcommand{\Zar}{\mathrm{zar}}
\providecommand{\Zell}{}\renewcommand{\Zell}{{\mathbf{Z}_{\ell}}}
\providecommand{\Zl}{}\renewcommand{\Zl}{{\mathbf Z_\ell}}
\providecommand{\affet}{}\renewcommand{\affet}[1]{{#1}^{\mathrm{aff}}_{\mathrm{ét}}}
\providecommand{\affproet}{}\renewcommand{\affproet}[1]{{#1}^{\mathrm{aff}}_{\mathrm{proét}}}
\providecommand{\boxt}{}\renewcommand{\boxt}{\boxtimes_A^{\mathrm L}}
\providecommand{\cF}{}\renewcommand{\cF}{\mathcal{F}}
\providecommand{\cHom}{}\renewcommand{\cHom}{R\mathcal{H}\!om}
\providecommand{\cI}{}\renewcommand{\cI}{\mathcal{I}}
\providecommand{\cL}{}\renewcommand{\cL}{\mathcal L}
\providecommand{\cM}{}\renewcommand{\cM}{\mathcal M}
\providecommand{\cN}{}\renewcommand{\cN}{\mathcal N}
\providecommand{\cO}{}\renewcommand{\cO}{\mathcal O}
\providecommand{\cT}{}\renewcommand{\cT}{\mathcal T}
\providecommand{\checked}{}\renewcommand{\checked}{\color{red}\aftergroup\restorecolor}
\providecommand{\coeff}{}\renewcommand{\coeff}{\operatorname{coeff}}
\providecommand{\coker}{}\renewcommand{\coker}{\operatorname{coker}}
\providecommand{\colim}{}\renewcommand{\colim}{\operatorname{colim}}
\providecommand{\constructedby}{}\renewcommand{\constructedby}[1]{\par\smallskip\noindent\textcolor{purple!60!black}{\textsf{Constructed by:} }
  \footnotesize\sloppy\nolinkurl{#1}\par}
\providecommand{\deepresult}{}\renewcommand{\deepresult}[1]{\par\smallskip\noindent\textcolor{orange!45!black}{
  \textsf{Blueprint role: named deep theorem node; see the cluster proof architecture.}}
  \par\noindent\textcolor{orange!45!black}{\textsf{Mathematical source:} }#1\par}
\providecommand{\dt}{}\renewcommand{\dt}{\mathbin{\otimes^{\mathbf L}}}
\providecommand{\dw}{}\renewcommand{\dw}{\operatorname{dw}}
\providecommand{\et}{}\renewcommand{\et}[1]{{#1}_{\mathrm{ét}}}
\providecommand{\ev}{}\renewcommand{\ev}{\operatorname{ev}}
\providecommand{\fib}{}\renewcommand{\fib}{\operatorname{fib}}
\providecommand{\formalizationnote}{}\renewcommand{\formalizationnote}[1]{\begin{quote}\small\raggedright\textbf{Formalization note.} #1\end{quote}}
\providecommand{\fpqc}{}\renewcommand{\fpqc}{X_{\mathrm{fpqc}}}
\providecommand{\gammacont}{}\renewcommand{\gammacont}{\Gamma_{\mathrm{cont}}}
\providecommand{\gm}{}\renewcommand{\gm}{\mathbf G_m}
\providecommand{\hocolim}{}\renewcommand{\hocolim}{\operatorname{hocolim}}
\providecommand{\holim}{}\renewcommand{\holim}{\operatorname{holim}}
\providecommand{\hrB}{}\renewcommand{\hrB}{\@ifnextchar]{\hrEnd}{\hrStep}}
\providecommand{\hrEnd}{}\renewcommand{\hrEnd}{\@firstoftwo{}}
\providecommand{\hrStep}{}\renewcommand{\hrStep}{\@firstoftwo{\hrB}}
\providecommand{\id}{}\renewcommand{\id}{\operatorname{id}}
\providecommand{\kunFl}{}\renewcommand{\kunFl}{\mathbf F_\ell}
\providecommand{\kunZn}{}\renewcommand{\kunZn}{\mathbf Z/n\mathbf Z}
\providecommand{\kuncd}{}\renewcommand{\kuncd}{\operatorname{cd}}
\providecommand{\kundim}{}\renewcommand{\kundim}{\dim}
\providecommand{\kunotimes}{}\renewcommand{\kunotimes}[1]{\otimes^{\mathbf L}_{#1}}
\providecommand{\kunsupp}{}\renewcommand{\kunsupp}{\operatorname{Supp}}
\providecommand{\lotimes}{}\renewcommand{\lotimes}{\otimes^{\mathrm L}}
\providecommand{\m}{}\renewcommand{\m}{{\mathfrak{m}}}
\providecommand{\n}{}\renewcommand{\n}{{\mathfrak{n}}}
\providecommand{\nagataGap}{}\renewcommand{\nagataGap}[1]{\par\noindent\textbf{Formalization gap.} #1\par}
\providecommand{\openobligation}{}\renewcommand{\openobligation}[1]{\par\smallskip\noindent\textcolor{red!65!black}{\textsf{Open obligation:} }
  \footnotesize\sloppy\nolinkurl{#1}\par}
\providecommand{\p}{}\renewcommand{\p}{{\mathfrak{p}}}
\providecommand{\pr}{}\renewcommand{\pr}{\operatorname{pr}}
\providecommand{\preshv}{}\renewcommand{\preshv}[1]{\mathrm{PreShv}(#1)}
\providecommand{\pro}{}\renewcommand{\pro}[1]{\mathrm{Pro}(#1)}
\providecommand{\proet}{}\renewcommand{\proet}{X_{\mathrm{proét}}}
\providecommand{\q}{}\renewcommand{\q}{{\mathfrak{q}}}
\providecommand{\rad}{}\renewcommand{\rad}{\operatorname{rad}}
\providecommand{\rank}{}\renewcommand{\rank}{\operatorname{rank}}
\providecommand{\red}{}\renewcommand{\red}{\color{red}}
\providecommand{\restorecolor}{}\renewcommand{\restorecolor}{\color{black}}
\providecommand{\retiredobligation}{}\renewcommand{\retiredobligation}[1]{\par\smallskip\noindent\textcolor{red!75!black}{\textsf{Retired invalid route:} }
  \footnotesize\sloppy\nolinkurl{#1}\par}
\providecommand{\rk}{}\renewcommand{\rk}{\operatorname{rk}}
\providecommand{\sep}{}\renewcommand{\sep}{\textnormal{sep}}
\providecommand{\shv}{}\renewcommand{\shv}[1]{\mathrm{Shv}(#1)}
\providecommand{\spec}{}\renewcommand{\spec}[1]{\mathrm{Spec}(#1)}
\providecommand{\stacks}{}\renewcommand{\stacks}[1]{\href{https://stacks.math.columbia.edu/tag/#1}{\textnormal{Stacks, Tag~#1}}}
\providecommand{\stacksproject}{}\renewcommand{\stacksproject}[1]{{\cite[\href{https://stacks.math.columbia.edu/tag/#1}{Tag #1}]{stacks-project}}}
\providecommand{\statusnote}{}\renewcommand{\statusnote}[1]{\begin{quote}\small\raggedright\textbf{Status.} #1\end{quote}}
\providecommand{\tildeideal}{}\renewcommand{\tildeideal}[2]{\widetilde{(#1, #2)}}
\providecommand{\tildering}{}\renewcommand{\tildering}[3]{{#1}_{\widetilde{(#2, #3)}}}
\providecommand{\tr}{}\renewcommand{\tr}{\operatorname{tr}}
\providecommand{\wt}{}\renewcommand{\wt}{\operatorname{wt}}

\begin{document}
\pagestyle{plain}
\begin{center}
{\LARGE\bfseries Weil Conjectures\par}
\vspace{0.7em}
{\large Part Ⅰ: The six operations\par}
\vspace{0.5em}
{\small\color{bpmuted}Snapshot 2026-09-26, revision 8e4cb88e2a\par}
\end{center}
\vspace{0.5em}
{\small\tableofcontents}
\bigskip\hrule\bigskip

% object: part-The-six-operations (blueprint/part-The-six-operations/_section.md)
\part*{\texorpdfstring{Part Ⅰ\quad The six operations}{Part Ⅰ The six operations}}
\phantomsection\label{bp1}%
\addcontentsline{toc}{part}{\texorpdfstring{Part Ⅰ\quad The six operations}{Part Ⅰ The six operations}}
% text: prose-P1-01
\phantomsection\label{bp2}%
On the étale site the constant sheaf \(\mathbf Z_\ell\) on \(\operatorname{Spec}\mathbf F_q\) has \(H^1=0\), although \(\varprojlim_nH^1(\operatorname{Spec}\mathbf F_q,\mathbf Z/\ell^n)=\mathbf Z_\ell\) (Section \bpout{1.2}). Moreover, Frobenius eigenvalues live in \(\overline{\mathbf Q}_\ell\), but SGA 4 constructs \(Rf_!\) and \(f^!\) and proves the finiteness theorems only for \(\mathbf Z/\ell^n\) [\hyperref[bp702]{SGA4}], [\hyperref[bp703]{SGA4½}]. So the theory we need is not the one SGA 4 provides. This part has to overcome two difficulties. The first concerns the coefficients: we must pass to the limit over \(n\) without losing cohomology. Classically one does this with systems \((F_n)\) and makes every derived construction by hand [\hyperref[bp704]{Eke90}]. We work instead on the pro-étale site of Bhatt and Scholze (Chapter \hyperref[bp3]{2}). There \(\mathbf Z_\ell\) is a sheaf of rings with the expected cohomology, and one definition serves all the coefficient rings

\[
A,\ \mathbf Z_\ell,\ \mathcal O_E,\ \mathbf Q_\ell,\ E,\ \overline{\mathbf Q}_\ell,
\]

compatibly with every change of coefficients between them (Section \hyperref[bp608]{7.2}).

What the rest of the notes needs from this part is the table of Section \bpout{1.3}: for every morphism \(f\) of schemes of finite type over \(k\), the operations \(f^*\), \(Rf_*\), \(Rf_!\), \(f^!\), \(\otimes^{\mathbf L}\) and \(\cHom\), with their adjunctions, base change and projection formulas, and the finiteness statement that each carries constructible complexes to constructible complexes (Theorem \hyperref[bp413]{5.2} and Proposition \hyperref[bp605]{7.1}). The later parts use these operations as their basic tools. For example, the trace formula of Part II sums trace functions over the fibres of \(f\) by applying \(Rf_!\), and Poincaré duality comes from the computation of \(f^!\Lambda\) for the structure morphism \(f\) of a smooth variety. Similarly, the Fourier transform of Part III is built from \(f^*\), \(\otimes^{\mathbf L}\) and \(Rf_!\).

The second difficulty is to show that \(f^!\) preserves constructible complexes, and it would be there even with finite coefficients. It arises because \(f^!\) is given only by its universal property, and already for a closed immersion \(i\) with open complement \(j\) its value is not obvious: we know only that \(i^!K\) sits in a triangle \(i^!K\to i^*K\to i^*Rj_*j^*K\). To see that it is constructible we need to know that \(Rj_*\) preserves constructibility, and this is Deligne's generic base change theorem. Accordingly, we construct \(f^!\) in Chapter \hyperref[bp407]{5} and prove that it preserves constructible complexes in Chapter \hyperref[bp513]{6}, as sketched in Section \bpout{1.4}.

The two difficulties are handled separately, so that each geometric argument is made only once. The geometric inputs, namely proper base change (Chapter \hyperref[bp285]{4}) and the computation of \(f^!\) on charts (Chapter \hyperref[bp513]{6}), are proved with torsion coefficients on the étale site, where SGA 4 applies. They are then transferred to \(\ell\)-adic coefficients without being proved again. Two facts make this transfer possible (Chapters \hyperref[bp193]{3} and \hyperref[bp601]{Chapter 7}). First, a constructible complex over \(\mathcal O_E\) is determined by its reductions modulo \(\pi_E^n\) (Proposition \hyperref[bp226]{3.11}), so its properties can be read off from torsion complexes. Second, whether a natural map between two such constructions is an isomorphism can be checked with finite coefficients (Proposition \hyperref[bp634]{7.12}). For example, smooth base change for \(\mathbf Z_\ell\), \(E\) and \(\overline{\mathbf Q}_\ell\) follows from smooth base change for finite rings in this way (Section \hyperref[bp677]{8.8}).

% object: chap-The-pro-etale-site (blueprint/chap-The-pro-etale-site/_section.md), content hidden
\section*{\texorpdfstring{2\quad The pro-étale site}{2 The pro-étale site}}
\phantomsection\label{bp3}%
\addcontentsline{toc}{section}{\texorpdfstring{2\quad The pro-étale site}{2 The pro-étale site}}
% text: prose-02-01
\phantomsection\label{bp4}%
On the étale site the constant sheaf \(\mathbf Z_\ell\) forgets that \(\mathbf Z_\ell\) is profinite. The pro-étale site \(X_{\mathrm{pro}\acute{e}t}\) of Bhatt and Scholze [\hyperref[bp705]{BS15}] repairs this. Its objects are the weakly étale \(X\)-schemes, and its coverings are fpqc coverings (Definition \hyperref[bp14]{2.4}).

We need three properties of this site.

\begin{itemize}
\item (a) \emph{Limits.} For a tower \(\cdots\to F_2\to F_1\) of surjections of sheaves, cohomology commutes with the limit: \(R\Gamma(X,\varprojlim_nF_n)\simeq R\varprojlim_nR\Gamma(X,F_n)\).
\item (b) \emph{Topology.} The sheaf \(\mathbf Z_\ell\) remembers the topology of \(\mathbf Z_\ell\).
\item (c) \emph{Comparison.} Sheaves coming from the étale site keep their étale cohomology, so that the results of SGA 4 still apply to them.
\end{itemize}

Property (b) can be seen on a point. Over a separably closed field the topos of \(X_{\mathrm{pro}\acute{e}t}\) is, up to a set-theoretic cutoff, that of condensed sets (Proposition \hyperref[bp153]{2.50}, proved in Subsection \hyperref[bp151]{2.12.1} below), and there the sections of \(\mathbf Z_\ell\) over a profinite set \(S\) are all continuous maps \(S\to\mathbf Z_\ell\), not only the locally constant ones.

\begin{center}\small
\begin{tabularx}{\linewidth}{>{\raggedright\arraybackslash}X>{\raggedright\arraybackslash}X>{\raggedright\arraybackslash}X}
\toprule
 & on \(X_{\acute{e}t}\) & on \(X_{\mathrm{pro}\acute{e}t}\) \\
\midrule
objects & étale \(X\)-schemes & weakly étale \(X\)-schemes \\
local objects & strictly henselian local rings & w-contractible rings \\
topos of a geometric point & sets & condensed sets \\
the sheaf \(\mathbf Z_\ell\) on a point & \(\mathbf Z_\ell\) with discrete topology & continuous maps to \(\mathbf Z_\ell\) \\
\bottomrule
\end{tabularx}
\end{center}

Property (a) fails on \(X_{\acute{e}t}\), for the following reason. A surjection of sheaves is surjective on sections only after refining the cover. A countable product of surjections needs one refinement for each factor, and these infinitely many refinements need not have a common refinement among étale covers. So a countable product of surjections need not be surjective, \(R^1\varprojlim_nF_n\) need not vanish, and \(\varprojlim_nF_n\) differs from \(R\varprojlim_nF_n\). For \(F_n=\mathbf Z/\ell^n\) on \(\operatorname{Spec}\mathbf F_q\) this is the lost \(H^1\). On \(X_{\mathrm{pro}\acute{e}t}\) there are enough objects on which no refinement is needed: the spectra of w-contractible rings, roughly families of strictly henselian local rings parametrized by an extremally disconnected profinite set. Every cover of such a spectrum splits, so every surjection of sheaves is surjective on its sections there, and limits of sheaves can be computed there as limits of sets and modules. However, building these objects takes two theorems from outside algebraic geometry: Gleason's theorem on extremally disconnected spaces and Olivier's theorem on weakly étale maps. Once we have them, property (a) follows in a few lines (Proposition \hyperref[bp172]{2.55}). In particular, over \(\mathbf F_q\) it gives back \(H^1=\mathbf Z_\ell\) (Example \hyperref[bp178]{2.57}).

For property (c) we compare the two sites through the pullback functor \(\nu_X^*\) from sheaves on \(X_{\acute{e}t}\) to sheaves on \(X_{\mathrm{pro}\acute{e}t}\). Its image consists of the sheaves that we call \emph{classical}: a sheaf on \(X_{\mathrm{pro}\acute{e}t}\) is classical if its sections over a limit \(U=\varprojlim_iU_i\) of affine étale \(X\)-schemes are the colimit of its sections over the \(U_i\) (Definition \hyperref[bp91]{2.30}). In other words, a classical sheaf is determined by its values on étale \(X\)-schemes, as an étale sheaf is. The comparison has four parts.

\begin{itemize}
\item \(\nu_X^*\) is an equivalence from étale sheaves onto classical sheaves (Theorem \hyperref[bp95]{2.32}).
\item It preserves the cohomology of bounded below complexes (Theorem \hyperref[bp103]{2.35}).
\item It commutes with direct images along quasi-compact quasi-separated morphisms (Theorem \hyperref[bp113]{2.38}).
\item For a finite coefficient ring, a constructible sheaf on \(X_{\mathrm{pro}\acute{e}t}\) is classical, and \(\nu_X^*\) is an equivalence from \(D^b_c(X_{\acute{e}t},\Lambda)\) onto \(\Dbc(X,\Lambda)\) (Theorem \hyperref[bp129]{2.43}). So with finite coefficients the two sites have the same constructible complexes.
\end{itemize}

Later, \(f^!\) is built by Brown representability on the unbounded derived category, so we also need the comparison for unbounded complexes. With finite coefficients it holds as well, because affine schemes étale over a scheme of finite type over \(k\) have uniformly bounded cohomological dimension. Each cohomology group of an unbounded complex therefore depends only on a bounded below truncation, where the comparison is already known (Proposition \hyperref[bp146]{2.49}). The sheaf \(\mathbf Z_\ell\), however, is not classical (Example \hyperref[bp97]{2.33}). This is what the computation of \(H^1(\operatorname{Spec}\mathbf F_q,\mathbf Z_\ell)\) in Section \bpout{1.2} needs: \(H^1\) comes out as \(\mathbf Z_\ell\) only when sections see the topology of \(\mathbf Z_\ell\). So the comparison applies to the \(\mathbf Z/\ell^n\) but not to \(\mathbf Z_\ell\) itself, and we reach \(\mathbf Z_\ell\) from them by derived completion (Section \hyperref[bp179]{2.14}).

Throughout this chapter, schemes are arbitrary unless a result states a finiteness hypothesis, because objects of \(X_{\mathrm{pro}\acute{e}t}\) are in general neither Noetherian nor of finite type, even when \(X\) is a variety. The standing assumption that schemes are of finite type over \(k\) resumes in Chapter \hyperref[bp193]{3}.

% object: sec-Weakly-etale-morphisms (blueprint/chap-The-pro-etale-site/sec-Weakly-etale-morphisms/_section.md), content hidden
\subsection*{\texorpdfstring{2.1\quad Weakly étale morphisms}{2.1 Weakly étale morphisms}}
\phantomsection\label{bp5}%
\addcontentsline{toc}{subsection}{\texorpdfstring{2.1\quad Weakly étale morphisms}{2.1 Weakly étale morphisms}}
% text: prose-02-02
\phantomsection\label{bp6}%
The objects we want are inverse limits of étale \(X\)-schemes. Being such a limit is a condition on a presentation, which is awkward to work with, so we take as objects the weakly étale \(X\)-schemes of the following definition. This class is defined by flatness conditions. For this reason it is stable under base change and composition, and membership in it is local for the Zariski topology (Lemma \hyperref[bp12]{2.3}). We lose little by the change: after a faithfully flat refinement every weakly étale scheme becomes a limit of étale schemes (Theorem \hyperref[bp42]{2.14}).

Typical examples are the schemes that the étale site can only reach through colimits. The strict henselization \(\operatorname{Spec}\mathcal O^{\mathrm{sh}}_{X,\bar x}\) is the limit of the affine étale neighbourhoods of \(\bar x\), and on \(X_{\acute{e}t}\) it appears only through the stalk at \(\bar x\). On \(X_{\mathrm{pro}\acute{e}t}\) it is an object. Similarly \(\operatorname{Spec}\overline{\mathbf F}_q=\varprojlim_m\operatorname{Spec}\mathbf F_{q^m}\) is an object over \(\operatorname{Spec}\mathbf F_q\), and for a profinite set \(S=\varprojlim S_i\) the scheme \(X\times S=\varprojlim_i\coprod_{S_i}X\) is an object over \(X\). Through the last family the site sees profinite sets, and hence the topology of \(\mathbf Z_\ell\).

% object: definition-ind-etale-pro-etale-weakly (blueprint/chap-The-pro-etale-site/sec-Weakly-etale-morphisms/definition-ind-etale-pro-etale-weakly/_section.md)
\phantomsection\label{bp7}%
\begin{bpstatement}{Definition~2.1}{ (Ind-étale, pro-étale and weakly étale morphisms)}
A ring map \(R\to B\) is \emph{ind-étale} if \(B=\varinjlim_\alpha B_\alpha\) for a filtered system of étale \(R\)-algebras \(B_\alpha\). A \emph{pro-étale affine over \(X\)} is a scheme \(U=\varprojlim_iU_i\) that is the limit of a cofiltered system of affine schemes \(U_i\) étale over \(X\). If \(X=\operatorname{Spec}(R)\) is affine, the pro-étale affines over \(X\) are the schemes

\[
\operatorname{Spec}(B)=\varprojlim_\alpha
\operatorname{Spec}(B_\alpha)\longrightarrow X
\]

with \(R\to B\) ind-étale.

A scheme morphism \(u:U\to X\) is \emph{weakly étale} if \(u\) and its diagonal \(\Delta_u:U\to U\times_XU\) are flat.
\end{bpstatement}

% text: prose-02-03
\phantomsection\label{bp8}%
The three classes of ring maps are related as follows.

% object: proposition-etale-implications (blueprint/chap-The-pro-etale-site/sec-Weakly-etale-morphisms/proposition-etale-implications/_section.md)
\phantomsection\label{bp9}%
\begin{bpstatementit}{Proposition~2.2}{ (Étale, ind-étale and weakly étale ring maps)}
Let \(R\to B\) be a ring map.

\begin{enumerate}
\item If \(R\to B\) is étale, it is ind-étale; if it is ind-étale, it is weakly étale.
\item If \(R\to B\) is weakly étale and of finite presentation, it is étale.
\end{enumerate}
\end{bpstatementit}

% proof: proof/proposition-etale-implications
\begin{proof}[{Proof}]\phantomsection\label{bp10}%
The first implication in (1) is immediate from the definition, since an étale algebra is the colimit of the constant system on itself. For the second, a filtered colimit of étale algebras is flat, and its multiplication map is the filtered colimit of the multiplication maps \(B_\alpha\otimes_RB_\alpha\to B_\alpha\), which are flat; see [Stacks, \href{https://stacks.math.columbia.edu/tag/092A}{092A}].

For (2) we show that \(R\to B\) is unramified, which together with flatness and finite presentation means that it is étale. The multiplication map \(B\otimes_RB\to B\) is surjective, and it is flat by assumption. It is of finite presentation because \(B\) is of finite presentation over \(R\). Hence the diagonal of \(\operatorname{Spec}(B)\to\operatorname{Spec}(R)\) is a closed immersion that is flat and of finite presentation. Its ideal \(I\) is then finitely generated with \(I=I^2\), hence generated by an idempotent, so the diagonal is also an open immersion. A morphism locally of finite type whose diagonal is an open immersion is unramified, and this proves (2).
\end{proof}

% text: prose-02-04
\phantomsection\label{bp11}%
For instance, \(\mathbf F_q\to\overline{\mathbf F}_q\) is ind-étale, as the colimit of the finite fields \(\mathbf F_{q^m}\), but it is not étale, because \(\overline{\mathbf F}_q\) is not of finite type over \(\mathbf F_q\).

In (2) the flatness of the diagonal takes the place of unramifiedness; flatness of \(R\to B\) alone would not suffice. The converse of the second implication in (1) fails in general: a weakly étale map need not be a filtered colimit of étale maps. It does become ind-étale after a faithfully flat refinement (Theorem \hyperref[bp42]{2.14} below, and [Stacks, \href{https://stacks.math.columbia.edu/tag/097Z}{097Z}]). Geometrically, a weakly étale scheme over \(X\) is covered, in the pro-étale topology, by a limit of étale \(X\)-schemes.

% object: lemma-stability (blueprint/chap-The-pro-etale-site/sec-Weakly-etale-morphisms/lemma-stability/_section.md)
\phantomsection\label{bp12}%
\begin{bpstatementit}{Lemma~2.3}{ (Stability)}
Weakly étale morphisms are stable under base change and composition, and the property is local on source and target for the Zariski topology. Every \(X\)-morphism between two weakly étale \(X\)-schemes is weakly étale. Ind-étale ring maps are stable under base change and composition. On affine charts, a morphism \(\operatorname{Spec}(B)\to\operatorname{Spec}(R)\) is weakly étale if and only if both

\[
R\longrightarrow B,
\qquad B\otimes_RB\longrightarrow B
\]

are flat.
\end{bpstatementit}

% proof: proof/lemma-stability
\begin{proof}[{Proof}]\phantomsection\label{bp13}%
Since weakly étale means that a morphism and its diagonal are flat, each assertion about weakly étale morphisms comes down to writing the relevant diagonal in terms of diagonals and base changes that we already know to be flat.

Base change preserves flatness of a morphism and its diagonal. For a composite \(V\to U\to X\), the morphism \(V\to X\) is flat as a composite of flat morphisms. Its diagonal \(\Delta_{V/X}\) is \(\Delta_{V/U}\) followed by \(V\times_UV\to V\times_XV\), and the second map is a base change of \(\Delta_{U/X}\). Both maps are flat, and hence so is \(\Delta_{V/X}\).

Let \(h:U\to V\) be an \(X\)-morphism between weakly étale \(X\)-schemes. We factor \(h\) as the graph \(U\to U\times_XV\) followed by the projection \(U\times_XV\to V\). The graph is a base change of the diagonal of \(V/X\), which is flat, and the projection is a base change of \(U\to X\), which is also flat. Their composite \(h\) is therefore flat. Similarly, the diagonal \(\Delta_h:U\to
U\times_VU\) is the base change of \(\Delta_{U/X}\) along the monomorphism \(U\times_VU\to U\times_XU\), so it is flat as well. Finally, Zariski locality follows from the fact that flatness is Zariski local.

On affine charts, the diagonal of \(\operatorname{Spec}(B)\to\operatorname{Spec}(R)\) is the closed immersion \(\operatorname{Spec}(B)\to\operatorname{Spec}(B\otimes_RB)\) given by the multiplication map, so the morphism is weakly étale exactly when \(R\to B\) and \(B\otimes_RB\to B\) are flat. By Zariski locality on source and target, this describes weakly étale morphisms on any affine charts.

For ind-étale ring maps, base change is clear because tensor products commute with filtered colimits. For composition, let \(R\to B\) and \(B\to C\) be ind-étale, with \(B=\varinjlim_\alpha B_\alpha\). Each étale \(B\)-algebra \(C'\) occurring in a presentation of \(C\) is of finite presentation over \(B\), hence is the base change of an étale algebra \(C'_\alpha\) over some \(B_\alpha\). Since \(B_\alpha\) is étale over \(R\), so is \(C'_\alpha\), and \(C'\) is the filtered colimit of the étale \(R\)-algebras \(C'_\alpha\otimes_{B_\alpha}B_\beta\) for \(\beta\geq\alpha\). Thus \(C\) is a filtered colimit of ind-étale \(R\)-algebras, and such a colimit is again ind-étale [Stacks, \href{https://stacks.math.columbia.edu/tag/097L}{097L}].
\end{proof}

% object: definition-small-pro-etale-site (blueprint/chap-The-pro-etale-site/sec-Weakly-etale-morphisms/definition-small-pro-etale-site/_section.md)
\phantomsection\label{bp14}%
\begin{bpstatement}{Definition~2.4}{ (The small pro-étale site)}
The objects of \(X_{\mathrm{pro}\acute{e}t}\) are weakly étale \(X\)-schemes, with all \(X\)-morphisms as arrows. A family \(\{U_i\to U\}_{i\in I}\) is a covering if it is an fpqc covering. Explicitly, for every affine open \(W\subset U\) there must be finitely many indices \(i_1,\ldots,i_m\) and affine opens \(V_j\subset U_{i_j}\times_UW\) with

\[
W=\bigcup_{j=1}^m\operatorname{image}(V_j\to W).
\]

The indices may depend on \(W\), and \(I\) may be infinite. Equivalently one may use finitely many quasi-compact opens above each \(W\). See [Stacks, \href{https://stacks.math.columbia.edu/tag/0989}{0989}, \href{https://stacks.math.columbia.edu/tag/098A}{098A}].
\end{bpstatement}

% text: prose-02-05
\phantomsection\label{bp15}%
The finite refinement condition is part of the topology, and a jointly surjective family that does not satisfy it is not a covering. We record two consequences of the definition.

% object: proposition-pro-etale-coverings (blueprint/chap-The-pro-etale-site/sec-Weakly-etale-morphisms/proposition-pro-etale-coverings/_section.md)
\phantomsection\label{bp16}%
\begin{bpstatementit}{Proposition~2.5}{ (Coverings in the pro-étale site)}
Let \(U\) be an object of \(X_{\mathrm{pro}\acute{e}t}\).

\begin{enumerate}
\item If \(U\) is affine, every covering of \(U\) can be refined by a finite affine covering, whose disjoint union is affine and faithfully flat over \(U\).
\item Every étale covering and every Zariski covering of \(U\) is a pro-étale covering.
\end{enumerate}
\end{bpstatementit}

% proof: proof/proposition-pro-etale-coverings
\begin{proof}[{Proof}]\phantomsection\label{bp17}%
Part (1) is the covering condition of Definition \hyperref[bp14]{2.4} with \(W=U\). For (2), recall that flat morphisms locally of finite presentation are open. By quasi-compactness we then find the required finite refinements; see [Stacks, \href{https://stacks.math.columbia.edu/tag/098B}{098B}].
\end{proof}

% text: prose-02-06
\phantomsection\label{bp18}%
By (2) the pro-étale topology is finer than the étale topology. By (1), over an affine object a covering may always be replaced by a single faithfully flat map. For example, the single map \(\operatorname{Spec}\overline{\mathbf F}_q\to\operatorname{Spec}\mathbf F_q\) is a pro-étale covering, although it is not an étale one. These coverings define a topos with the expected functoriality in \(X\).

% object: proposition-topos-functoriality (blueprint/chap-The-pro-etale-site/sec-Weakly-etale-morphisms/proposition-topos-functoriality/_section.md)
\phantomsection\label{bp19}%
\begin{bpstatementit}{Proposition~2.6}{ (The topos and its functoriality)}
The above coverings define a subcanonical site. Its category of sheaves \(\Sh(X_{\mathrm{pro}\acute{e}t})\) is a Grothendieck topos. Every scheme morphism \(f:X\to Y\) induces a geometric morphism

\[
f_{\mathrm{pro}\acute{e}t}:\Sh(X_{\mathrm{pro}\acute{e}t})\longrightarrow\Sh(Y_{\mathrm{pro}\acute{e}t}),
\qquad f^{-1}\dashv f_*.
\]

Here \(f^{-1}:\Sh(Y_{\mathrm{pro}\acute{e}t})\to\Sh(X_{\mathrm{pro}\acute{e}t})\) preserves finite limits and all colimits, and is exact on abelian sheaves. For a second morphism \(g:Y\to Z\) there is a canonical isomorphism \((g\circ f)^{-1}\simeq f^{-1}\circ g^{-1}\).
\end{bpstatementit}

% proof: proof/proposition-topos-functoriality
\begin{proof}[{Proof}]\phantomsection\label{bp20}%
There are two things to check: that the coverings form a subcanonical topology, and that base change of objects is a morphism of sites.

Identity covers, composites of covers, and base changes of covers are again covers. For the underlying morphisms this is Lemma \hyperref[bp12]{2.3}, and for the finite refinement condition it is part of the fpqc covering axioms. Representable presheaves satisfy descent for these covers by fpqc descent for morphisms of schemes, so they are sheaves, and the topology is subcanonical.

For \(f:X\to Y\), the functor \(V\mapsto V\times_YX\) sends weakly étale \(Y\)-schemes to weakly étale \(X\)-schemes by Lemma \hyperref[bp12]{2.3}. It preserves fibre products and the final object, hence finite limits, and it carries covering families to covering families because fpqc coverings are stable under base change. Such a functor defines a geometric morphism, whose inverse image \(f^{-1}\) is left exact and, as a left adjoint, preserves all colimits. Exactness on abelian sheaves follows, and compatibility with composition holds because \((V\times_ZY)\times_YX\simeq V\times_ZX\) canonically. This is the small-site version of the constructions in [Stacks, \href{https://stacks.math.columbia.edu/tag/0988}{0988}].
\end{proof}


% object: sec-w-local-and-w-contractible-rings (blueprint/chap-The-pro-etale-site/sec-w-local-and-w-contractible-rings/_section.md), content hidden
\subsection*{\texorpdfstring{2.2\quad w-local and w-contractible rings}{2.2 w-local and w-contractible rings}}
\phantomsection\label{bp21}%
\addcontentsline{toc}{subsection}{\texorpdfstring{2.2\quad w-local and w-contractible rings}{2.2 w-local and w-contractible rings}}
% text: prose-02-07
\phantomsection\label{bp22}%
To compute limits of sheaves by taking sections, we want objects over which every cover splits, so that sheafification changes nothing there. On the étale site the spectrum of a strictly henselian local ring, such as the étale local ring of a variety at a geometric point, has this property, but such spectra do not cover \(X\). Instead, the rings we need are families of strictly henselian local rings over a profinite set of components. Moreover, we require this set to be extremally disconnected, which means that the closure of every open subset is open, because under this condition a cover splits. Namely, over each component of \(\operatorname{Spec}(R)\) we choose a component of the cover lying above it, and extremal disconnectedness lets us make this choice continuously. The model case is the ring of locally constant functions from a profinite set to a separably closed field (Example \hyperref[bp28]{2.9}), which is a family of copies of one field indexed by the profinite set. Below we define these rings, the w-contractible rings, by a hierarchy of local conditions, and we show that every ring has a faithfully flat ind-étale cover by one of them (Theorem \hyperref[bp39]{2.13}).

Recall that spectral spaces are quasi-compact and sober, with a basis of quasi-compact opens closed under finite intersections, and that the space of connected components of a spectral space is profinite.

% object: definition-w-local-w-strictly-local (blueprint/chap-The-pro-etale-site/sec-w-local-and-w-contractible-rings/definition-w-local-w-strictly-local/_section.md)
\phantomsection\label{bp23}%
\begin{bpstatement}{Definition~2.7}{ (w-local, w-strictly local and w-contractible rings)}
A topological space is \emph{extremally disconnected} if the closure of every open subset is open. A spectral space \(T\) is \emph{w-local} if its set \(T^c\) of closed points is closed and every point specializes to exactly one closed point. A ring \(R\) is w-local if \(\operatorname{Spec}(R)\) is w-local. It is \emph{w-strictly local} if it is w-local and \(R_{\mathfrak m}\) is strictly henselian for every maximal ideal \(\mathfrak m\). It is \emph{w-contractible} if it is w-strictly local and the profinite space \(\pi_0(\operatorname{Spec}(R))\) is extremally disconnected.
\end{bpstatement}

% text: prose-02-08
\phantomsection\label{bp24}%
Closedness of \(T^c\) does not follow from uniqueness of closed specializations, which is why the definition requires it separately. The basic properties of these notions are as follows.

% object: proposition-w-local-components (blueprint/chap-The-pro-etale-site/sec-w-local-and-w-contractible-rings/proposition-w-local-components/_section.md)
\phantomsection\label{bp25}%
\begin{bpstatementit}{Proposition~2.8}{ (Connected components of w-local spaces and rings)}
\begin{enumerate}
\item If \(T\) is a w-local spectral space, the natural map \(T^c\to\pi_0(T)\) is a homeomorphism.
\item If \(R\) is w-local, the component of \(\operatorname{Spec}(R)\) containing a closed point \(\mathfrak m\) is the image of \(\operatorname{Spec}(R_{\mathfrak m})\). In particular, if \(R\) is w-strictly local, every component is the spectrum of a strictly henselian local ring.
\item A strictly henselian local ring is w-contractible.
\end{enumerate}
\end{bpstatementit}

% proof: proof/proposition-w-local-components
\begin{proof}[{Proof}]\phantomsection\label{bp26}%
For (1), each component consists of the generalizations of its unique closed point, and \(T^c\) is compact. Part (2) follows, and (3) holds because the spectrum of a strictly henselian local ring has one closed point and one component. See [Stacks, \href{https://stacks.math.columbia.edu/tag/096A}{096A}, \href{https://stacks.math.columbia.edu/tag/0968}{0968}, \href{https://stacks.math.columbia.edu/tag/0969}{0969}].
\end{proof}

% text: prose-02-09
\phantomsection\label{bp27}%
So a w-strictly local ring is a family of strictly henselian local rings indexed by the profinite set of its components, and a w-contractible ring is such a family indexed by an extremally disconnected profinite set.

% object: example-locally-constant-functions-w-local (blueprint/chap-The-pro-etale-site/sec-w-local-and-w-contractible-rings/example-locally-constant-functions-w-local/_section.md)
\phantomsection\label{bp28}%
\begin{bpstatement}{Example~2.9}{ (Locally constant functions)}
Let \(k'\) be a separably closed field, let \(S=\varprojlim_iS_i\) be a profinite set written as a limit of finite sets, and let \(R=LC(S,k')\) be the ring of locally constant functions \(S\to k'\). Then \(R=\varinjlim_ik'^{\,S_i}\), and \(\operatorname{Spec}(k'^{\,S_i})\) is the finite discrete set \(S_i\), so \(\operatorname{Spec}(R)=\varprojlim_iS_i=S\): the point \(s\in S\) is the maximal ideal of functions vanishing at \(s\). Every point is closed and is its own component, so \(R\) is w-local with \(\pi_0(\operatorname{Spec}(R))=S\). The local ring at \(s\) is the colimit of the local rings \(k'\) of the \(k'^{\,S_i}\) at the image of \(s\), that is, \(k'\), which is strictly henselian. Thus \(R\) is w-strictly local: it is the family of copies of \(\operatorname{Spec}k'\) indexed by \(S\). It is w-contractible exactly when \(S\) is extremally disconnected. In this example the definitions are statements about \(S\) alone. Over \(\operatorname{Spec}k'\) every pro-étale affine is of this form (Proposition \hyperref[bp153]{2.50} below).
\end{bpstatement}

% text: prose-02-10
\phantomsection\label{bp29}%
The constructions below change the set of components while keeping the local rings, and for this we use the following class of maps. A ring map \(R\to A\) is a \emph{local isomorphism} if \(\operatorname{Spec}(A)\) is covered by principal opens \(D(g)\), each mapping isomorphically onto an open subscheme of \(\operatorname{Spec}(R)\) [Stacks, \href{https://stacks.math.columbia.edu/tag/096E}{096E}], and it is \emph{ind-Zariski} if it is a filtered colimit of local isomorphisms [Stacks, \href{https://stacks.math.columbia.edu/tag/096M}{096M}]. For instance, \(R\to R\times R\) is a local isomorphism. Taking the colimit over the finite quotients \(S_i\) of a profinite set \(S\), we see that \(R\to\varinjlim_iR^{S_i}=LC(S,R)\) is ind-Zariski. Its spectrum is \(\operatorname{Spec}(R)\times S\). A localization \(R\to R_{\mathfrak p}\) at a prime is ind-Zariski as well, as the colimit of the \(R_f\) with \(f\notin\mathfrak p\). The property we use is that an ind-Zariski map does not change local rings: each local ring of \(A\) is a local ring of \(R\).

We use extremally disconnected spaces through their projectivity. In Proposition \hyperref[bp36]{2.12} we use it to split a surjection of component spaces, and in Theorem \hyperref[bp39]{2.13} we cover a profinite component space by an extremally disconnected one (Lemma \hyperref[bp33]{2.11}).

% object: proposition-gleason-s-theorem (blueprint/chap-The-pro-etale-site/sec-w-local-and-w-contractible-rings/proposition-gleason-s-theorem/_section.md)
\phantomsection\label{bp30}%
\begin{bpstatementit}{Proposition~2.10}{ (Gleason's theorem)}
A compact Hausdorff space \(S\) is extremally disconnected if and only if it is projective with respect to continuous surjections of compact Hausdorff spaces: for every continuous surjection \(T_1\to T_2\) of compact Hausdorff spaces, every continuous map \(S\to T_2\) lifts to a continuous map \(S\to T_1\).
\end{bpstatementit}

% proof: proof/proposition-gleason-s-theorem
\begin{proof}[{Proof}]\phantomsection\label{bp31}%
This is the projectivity theorem for compact Hausdorff spaces; see [Stacks, \href{https://stacks.math.columbia.edu/tag/08YN}{08YN}]. Here we sketch the implication from extremal disconnectedness to projectivity, which is the direction used later. To lift a map \(S\to T_2\) along a surjection \(T_1\to T_2\), one pulls the surjection back to \(S\) and chooses a minimal closed subspace of the fibre product that still surjects onto \(S\). Then extremal disconnectedness forces this irreducible surjection to be a homeomorphism, and its inverse gives the lift.
\end{proof}

% text: prose-02-11
\phantomsection\label{bp32}%
To use the proposition we need enough extremally disconnected spaces. They come from Stone–Čech compactifications of discrete sets.

% object: lemma-extremally-disconnected-covers (blueprint/chap-The-pro-etale-site/sec-w-local-and-w-contractible-rings/lemma-extremally-disconnected-covers/_section.md)
\phantomsection\label{bp33}%
\begin{bpstatementit}{Lemma~2.11}{ (Extremally disconnected covers)}
Every compact Hausdorff space \(T\), in particular every profinite space, admits a continuous surjection \(S\to T\) with \(S\) compact Hausdorff and extremally disconnected. Every extremally disconnected compact Hausdorff space is profinite.
\end{bpstatementit}

% proof: proof/lemma-extremally-disconnected-covers
\begin{proof}[{Proof}]\phantomsection\label{bp34}%
Write \(T^{\delta}\) for the underlying set of \(T\) with the discrete topology and take \(S=\beta(T^{\delta})\), its Stone–Čech compactification. The Stone–Čech compactification of a discrete space is extremally disconnected, and the identity \(T^{\delta}\to T\) extends uniquely to a continuous surjection \(\beta(T^{\delta})\to T\); see [Stacks, \href{https://stacks.math.columbia.edu/tag/090D}{090D}]. An extremally disconnected compact Hausdorff space is totally disconnected, since the closure of an open set is clopen, and hence it is profinite.
\end{proof}

% text: prose-02-12
\phantomsection\label{bp35}%
A finite discrete space is extremally disconnected, while the Cantor set \(\{0,1\}^{\mathbf N}\) is not. To see this, let \(V_n\) be the clopen set of sequences beginning with \(0^n1\), and let \(U\) be the union of the \(V_n\) for odd \(n\). The closure of \(U\) is \(U\cup\{000\cdots\}\), and it is not open, because every neighbourhood of \(000\cdots\) contains some \(V_n\) with \(n\) even, which is disjoint from the closure. More generally, infinite extremally disconnected compact spaces have no nontrivial convergent sequences, and therefore they are never metrizable. In practice they arise as Stone–Čech compactifications.

Over a w-contractible ring an ind-étale cover splits at each closed point by strict henselianity, and Gleason's theorem lets us choose these local splittings continuously in the point.

% object: proposition-splitting-ind-etale-covers (blueprint/chap-The-pro-etale-site/sec-w-local-and-w-contractible-rings/proposition-splitting-ind-etale-covers/_section.md)
\phantomsection\label{bp36}%
\begin{bpstatementit}{Proposition~2.12}{ (Splitting ind-étale covers)}
If \(R\) is w-contractible, every faithfully flat ind-étale ring map \(R\to B\) has an \(R\)-algebra retraction \(B\to R\).
\end{bpstatementit}

% proof: proof/proposition-splitting-ind-etale-covers
\begin{proof}[{Proof (Proof outline)}]\phantomsection\label{bp37}%
We first replace \(B\) by an ind-Zariski \(B\)-algebra, so that its spectrum is w-local and its closed points lie exactly over the closed points of \(\operatorname{Spec}(R)\); such a refinement does not change the local rings. Let \(\mathfrak q\subset B\) lie over a closed point \(\mathfrak m\) of \(\operatorname{Spec}(R)\). Since \(R_{\mathfrak m}\) is strictly henselian, the ind-étale local homomorphism \(R_{\mathfrak m}\to B_{\mathfrak q}\) is an isomorphism. Hence every component of \(\operatorname{Spec}(B)\) maps isomorphically onto the component of \(\operatorname{Spec}(R)\) below it, and a retraction amounts to a continuous choice of one component above each component of \(\operatorname{Spec}(R)\). Such a choice exists because, by Proposition \hyperref[bp30]{2.10}, the surjection \(\pi_0(\operatorname{Spec}(B))\to\pi_0(\operatorname{Spec}(R))\) of profinite component spaces has a section. Given such a section, we restrict \(B\) to the corresponding closed set of components by a filtered colimit of localizations at idempotents. The result is an \(R\)-algebra isomorphic to \(R\), and this gives the retraction. For details see [Stacks, \href{https://stacks.math.columbia.edu/tag/097D}{097D}, \href{https://stacks.math.columbia.edu/tag/09AZ}{09AZ}, \href{https://stacks.math.columbia.edu/tag/0982}{0982}].
\end{proof}

% text: prose-02-13
\phantomsection\label{bp38}%
The proposition would be of little use if w-contractible rings were rare. The next theorem shows that every ring has a faithfully flat ind-étale cover by one, so that, after a single refinement, every affine object of \(X_{\mathrm{pro}\acute{e}t}\) is replaced by one over which covers split.

% object: theorem-existence-contractible-covers (blueprint/chap-The-pro-etale-site/sec-w-local-and-w-contractible-rings/theorem-existence-contractible-covers/_section.md)
\phantomsection\label{bp39}%
\begin{bpstatementit}{Theorem~2.13}{ (Existence of contractible covers)}
Every ring \(R\) admits a faithfully flat ind-étale algebra \(D\) that is w-contractible.
\end{bpstatementit}

% proof: proof/theorem-existence-contractible-covers
\begin{proof}[{Proof}]\phantomsection\label{bp40}%
The construction is that of [\hyperref[bp705]{BS15}, Theorem 2.4.9] and [Stacks, \href{https://stacks.math.columbia.edu/tag/0983}{0983}]. We first make the local rings strictly henselian, and then replace the profinite space of components by an extremally disconnected one surjecting onto it, which does not change the local rings.

In the first stage we construct a faithfully flat ind-étale \(R\)-algebra \(C\) that is w-strictly local. To do so we first replace \(R\) by an ind-Zariski algebra whose spectrum is w-local, and then take a filtered colimit of étale neighbourhoods of the closed points, which makes the local rings at the closed points strictly henselian. This is the construction of simultaneous strict localizations in [Stacks, \href{https://stacks.math.columbia.edu/tag/097X}{097X}]. Each step is a filtered colimit of étale algebras, so \(C\) is ind-étale over \(R\) by Lemma \hyperref[bp12]{2.3}.

For the second stage, choose a surjection \(S\to\pi_0(\operatorname{Spec}(C))\) with \(S\) extremally disconnected using Lemma \hyperref[bp33]{2.11}, and realize this map by an ind-Zariski \(C\)-algebra \(D\) such that

\[
|\operatorname{Spec}(D)|
\simeq |\operatorname{Spec}(C)|
\times_{\pi_0(\operatorname{Spec}(C))}S
\]

as topological spaces. Concretely, write \(S\) as a limit of finite quotients, use finite products of \(C\), restrict to the relevant closed sets of components by idempotent localizations, and take a filtered colimit; see [Stacks, \href{https://stacks.math.columbia.edu/tag/097D}{097D}]. This construction identifies local rings, is surjective when \(S\) is, and preserves w-locality. Hence \(C\to D\) is faithfully flat and the local rings at closed points of \(D\) remain strictly henselian. The space of components of \(\operatorname{Spec}(D)\) is \(S\), so \(D\) is w-contractible, and the composite \(R\to C\to D\) is ind-étale and faithfully flat.
\end{proof}

% text: prose-02-14
\phantomsection\label{bp41}%
To compute on \(X_{\mathrm{pro}\acute{e}t}\) we want the limits of étale affines to form a basis (Proposition \hyperref[bp46]{2.15}), so that étale sheaves can be evaluated on every object by a colimit. Objects of \(X_{\mathrm{pro}\acute{e}t}\) are only weakly étale, while the covers above are ind-étale, and the next theorem bridges the two.

% object: theorem-pro-etale-refinements-weakly-etale (blueprint/chap-The-pro-etale-site/sec-w-local-and-w-contractible-rings/theorem-pro-etale-refinements-weakly-etale/_section.md)
\phantomsection\label{bp42}%
\begin{bpstatementit}{Theorem~2.14}{ (Pro-étale refinements of weakly étale maps)}
If \(R\to B\) is weakly étale, there is a faithfully flat ind-étale map \(B\to B'\) such that \(R\to B'\) is ind-étale.
\end{bpstatementit}

% proof: proof/theorem-pro-etale-refinements-weakly-etale
\begin{proof}[{Proof}]\phantomsection\label{bp43}%
Weakly étale and ind-étale maps differ only in how local rings are glued. Namely, after strict localization of the base a weakly étale map identifies local rings, that is, for every prime \(\mathfrak q\) of \(B\) lying over \(\mathfrak p\subset R\) the induced map of local rings \(R_{\mathfrak p}\to B_{\mathfrak q}\) is an isomorphism. Such a map becomes ind-Zariski after refinement.

Applying the construction of simultaneous strict localizations to the pair \(R\to B\) we obtain compatible ind-étale refinements \(R\to C\) and \(B\to D\), with both spectra w-local and with closed points of \(D\) mapping to closed points of \(C\). The map \(C\to D\) is still weakly étale by Lemma \hyperref[bp12]{2.3}, and the local rings of \(C\) at closed points are strictly henselian. Olivier's rigidity theorem says that a weakly étale local map out of a strictly henselian local ring is an isomorphism; see [Stacks, \href{https://stacks.math.columbia.edu/tag/092Z}{092Z}]. Applying it at the closed points of \(D\), which lie over closed points of \(C\), we get isomorphisms of local rings there. Every point of the w-local \(\operatorname{Spec}(D)\) is a generalization of a closed point, and its local ring is a localization of the local ring at that closed point. Hence \(C\to D\) identifies all local rings. For a ring map that identifies local rings there is a faithfully flat ind-Zariski refinement of the target over which the map becomes ind-Zariski [Stacks, \href{https://stacks.math.columbia.edu/tag/097G}{097G}]. Applied to \(C\to D\), this gives a faithfully flat ind-Zariski map \(D\to D'\) such that \(C\to D'\) is ind-Zariski, and we take \(B'=D'\).

The algebra \(B'\) has the two required properties. The map \(B\to B'\) is the composite \(B\to D\to D'\) of two faithfully flat ind-étale maps, and \(R\to B'\) is the composite \(R\to C\to D'\) of an ind-étale and an ind-Zariski map. Both composites are ind-étale by Lemma \hyperref[bp12]{2.3}. For details see [Stacks, \href{https://stacks.math.columbia.edu/tag/097Z}{097Z}].
\end{proof}


% object: sec-Affine-bases (blueprint/chap-The-pro-etale-site/sec-Affine-bases/_section.md), content hidden
\subsection*{\texorpdfstring{2.3\quad Affine bases}{2.3 Affine bases}}
\phantomsection\label{bp44}%
\addcontentsline{toc}{subsection}{\texorpdfstring{2.3\quad Affine bases}{2.3 Affine bases}}
% text: prose-02-15
\phantomsection\label{bp45}%
To compute with sheaves on \(X_{\mathrm{pro}\acute{e}t}\) it suffices to evaluate on two kinds of affines. Every object of \(X_{\mathrm{pro}\acute{e}t}\) is covered by pro-étale affines (Theorem \hyperref[bp42]{2.14}), and these are in turn covered by w-contractible affines (Theorem \hyperref[bp39]{2.13}), over which taking sections is exact.

The two kinds of affines are used for different purposes. On a pro-étale affine \(U=\varprojlim_iU_i\), a sheaf pulled back from \(X_{\acute{e}t}\) can be evaluated by a colimit, and for this reason the comparison theorems of Section \hyperref[bp89]{2.7} are proved on such affines. On a w-contractible affine, by contrast, every epimorphism of sheaves is surjective on sections, so the homological results of the next three sections are proved there.

In Lemma \hyperref[bp49]{2.16} we turn the splitting of ring maps in Proposition \hyperref[bp36]{2.12} into a statement about sheaves. The idea is this. A finite cover of a w-contractible affine has a section. Such a section amounts to a decomposition of the affine into finitely many clopen pieces, each mapping into one member of the cover. Local lifts along an epimorphism can then be glued over these pieces without any further refinement.

% object: proposition-pro-etale-affines-form-basis (blueprint/chap-The-pro-etale-site/sec-Affine-bases/proposition-pro-etale-affines-form-basis/_section.md)
\phantomsection\label{bp46}%
\begin{bpstatementit}{Proposition~2.15}{ (Pro-étale affines form a basis)}
Every object of \(X_{\mathrm{pro}\acute{e}t}\) has a pro-étale covering by pro-étale affines over \(X\). These form a full subsite whose category of sheaves, of sets or of abelian groups, is equivalent to \(\Sh(X_{\mathrm{pro}\acute{e}t})\) by restriction.
\end{bpstatementit}

% proof: proof/proposition-pro-etale-affines-form-basis
\begin{proof}[{Proof}]\phantomsection\label{bp47}%
The first assertion is the refinement theorem applied chart by chart, and the second is then formal.

Given an object \(U\) of \(X_{\mathrm{pro}\acute{e}t}\), cover \(X\) by affine opens and cover their inverse images in \(U\) by affine opens. These Zariski coverings are pro-étale coverings, so it suffices to cover each chart. On such a chart the map is a weakly étale ring map \(R\to B\). By Theorem \hyperref[bp42]{2.14} there is a single faithfully flat cover \(\operatorname{Spec}(B')\to\operatorname{Spec}(B)\) such that \(\operatorname{Spec}(B')\) is a limit of affine schemes étale over \(\operatorname{Spec}(R)\), hence over \(X\). The equivalence of categories of sheaves follows from the comparison lemma for sites; see [\hyperref[bp705]{BS15}, Lemma 4.2.4].
\end{proof}

% text: prose-02-16
\phantomsection\label{bp48}%
Over \(X=\operatorname{Spec}\mathbf F_q\), for example, the weakly étale \(X\)-scheme \(\operatorname{Spec}\overline{\mathbf F}_q\) is already a pro-étale affine, and \(X\) itself is covered by it. Applying Theorem \hyperref[bp39]{2.13} to the charts of such a basis, we pass from pro-étale affines to w-contractible ones. The next two results show that the spectra of w-contractible rings have the splitting property for sheaves and that they cover every object.

% object: lemma-w-contractible-affines-projective (blueprint/chap-The-pro-etale-site/sec-Affine-bases/lemma-w-contractible-affines-projective/_section.md)
\phantomsection\label{bp49}%
\begin{bpstatementit}{Lemma~2.16}{ (Sections of epimorphisms onto w-contractible affines)}
Let \(U=\operatorname{Spec}(D)\) be an object of \(X_{\mathrm{pro}\acute{e}t}\) with \(D\) w-contractible. Then every epimorphism of sheaves \(F\to h_U\) has a section. Equivalently, \(h_U\) is projective in \(\Sh(X_{\mathrm{pro}\acute{e}t})\). More generally, for an object \(W\) of \(X_{\mathrm{pro}\acute{e}t}\) the following are equivalent: every epimorphism of sheaves of sets \(F\to h_W\) admits a section; \(h_W\) is projective in \(\Sh(X_{\mathrm{pro}\acute{e}t})\); every epimorphism \(F\to G\) induces a surjection \(F(W)\to G(W)\).
\end{bpstatementit}

% proof: proof/lemma-w-contractible-affines-projective
\begin{proof}[{Proof}]\phantomsection\label{bp50}%
For the equivalence of the three conditions on an object \(W\): by the Yoneda lemma, \(h_W\) is projective exactly when every epimorphism \(F\to G\) induces a surjection \(F(W)=\Hom(h_W,F)\to\Hom(h_W,G)=G(W)\). If every epimorphism onto \(h_W\) splits, pull an epimorphism \(F\to G\) back along a section \(h_W\to G\); the resulting epimorphism onto \(h_W\) splits, so the section lifts to \(F\). Conversely, if every epimorphism is surjective on sections over \(W\), then \(\mathrm{id}_{h_W}\) lifts along any epimorphism onto \(h_W\).

To see that covers of \(U\) split, note that by Theorem \hyperref[bp42]{2.14} a faithfully flat weakly étale \(D\)-algebra can be refined to a faithfully flat ind-étale \(D\)-algebra. This refinement splits by Proposition \hyperref[bp36]{2.12}, and hence so does the original map. Now, given a cover of the affine \(U\), we refine it by a finite affine cover \(\{V_j\to U\}\). Its disjoint union \(\coprod_jV_j\) is a single faithfully flat weakly étale affine cover, which therefore has a section \(s:U\to\coprod_jV_j\). The sets \(s^{-1}(V_j)\) then form a finite clopen partition of \(U\), and on \(s^{-1}(V_j)\) the inclusion into \(U\) factors through \(V_j\).

If \(F\to h_U\) is an epimorphism, then by the definition of an epimorphism of sheaves \(\mathrm{id}_U\) lifts after passing to a covering of \(U\). We refine this covering by the clopen partition just constructed, so that on each piece \(s^{-1}(V_j)\) we have a lift. Since the pieces are disjoint and cover \(U\), the sheaf condition glues these lifts with no compatibility to check, and the result is a global lift of \(\mathrm{id}_U\), that is, a section. The same argument, after pulling back an arbitrary epimorphism along a map from \(h_U\), proves projectivity. See [Stacks, \href{https://stacks.math.columbia.edu/tag/098F}{098F}, \href{https://stacks.math.columbia.edu/tag/0990}{0990}].
\end{proof}

% object: theorem-enough-affine-weakly-contractible-objects (blueprint/chap-The-pro-etale-site/sec-Affine-bases/theorem-enough-affine-weakly-contractible-objects/_section.md)
\phantomsection\label{bp51}%
\begin{bpstatementit}{Theorem~2.17}{ (Enough affine weakly contractible objects)}
Every object of \(X_{\mathrm{pro}\acute{e}t}\) has a covering by affines \(U=\operatorname{Spec}(D)\) with \(D\) w-contractible.
\end{bpstatementit}

% proof: proof/theorem-enough-affine-weakly-contractible-objects
\begin{proof}[{Proof}]\phantomsection\label{bp52}%
Cover the given object by pro-étale affines \(\operatorname{Spec}(B)\) (Proposition \hyperref[bp46]{2.15}), and for each of them choose, by Theorem \hyperref[bp39]{2.13}, a faithfully flat ind-étale \(B\)-algebra \(D\) that is w-contractible. Then \(\operatorname{Spec}(D)\to\operatorname{Spec}(B)\) is weakly étale by Proposition \hyperref[bp9]{2.2}, so \(\operatorname{Spec}(D)\) is an object of \(X_{\mathrm{pro}\acute{e}t}\) by Lemma \hyperref[bp12]{2.3}, and as a faithfully flat map of affines it is a covering. Composing these coverings gives the result.
\end{proof}


% object: sec-Weakly-contractible-objects (blueprint/chap-The-pro-etale-site/sec-Weakly-contractible-objects/_section.md), content hidden
\subsection*{\texorpdfstring{2.4\quad Weakly contractible objects}{2.4 Weakly contractible objects}}
\phantomsection\label{bp53}%
\addcontentsline{toc}{subsection}{\texorpdfstring{2.4\quad Weakly contractible objects}{2.4 Weakly contractible objects}}
% text: prose-02-17
\phantomsection\label{bp54}%
On a topological space we check exactness at points, and on the étale site at geometric points. On \(X_{\mathrm{pro}\acute{e}t}\), instead, we check exactness, and compute limits and colimits, by evaluating on weakly contractible objects. Over these objects we can compute in the ordinary derived categories of sheaves of modules by taking sections. Until Section \hyperref[bp89]{2.7}, \(\mathcal R\) denotes a sheaf of rings on \(X_{\mathrm{pro}\acute{e}t}\), for instance the pullback of a condensed ring (Section \hyperref[bp149]{2.12}).

The main result is Corollary \hyperref[bp59]{2.20}. On a weakly contractible affine \(W\), taking sections is exact and commutes with all limits and colimits of module sheaves. Moreover, these functors, for all such \(W\), are jointly conservative.

Stalk functors behave worse. Indeed, a stalk functor is exact and commutes with colimits, but not with infinite products. Moreover, products of sheaves on \(X_{\acute{e}t}\) need not be exact either, because a section of a product of surjections lifts in each factor only after its own refinement of the cover, and infinitely many refinements cannot in general be made at once. Evaluation at \(W\) avoids both problems: it commutes with products because limits of sheaves are computed on sections, and it is exact because covers of \(W\) split.

In the simplest case, \(X\) a geometric point, the site is that of profinite sets (Proposition \hyperref[bp153]{2.50}, proved in Subsection \hyperref[bp151]{2.12.1} below). The weakly contractible affines are then the extremally disconnected ones, among them \(\beta(T^\delta)\) for a set \(T\). Indeed, a surjection of profinite sets onto such an \(S\) has a section by Gleason's theorem.

% object: definition-weak-contractibility (blueprint/chap-The-pro-etale-site/sec-Weakly-contractible-objects/definition-weak-contractibility/_section.md)
\phantomsection\label{bp55}%
\begin{bpstatement}{Definition~2.18}{ (Weak contractibility)}
An object \(W\) of \(X_{\mathrm{pro}\acute{e}t}\) is \emph{weakly contractible} if every epimorphism of sheaves of sets \(F\to h_W\) admits a section, where \(h_W\) is the representable sheaf associated to \(W\). By Lemma \hyperref[bp49]{2.16}, equivalently \(h_W\) is projective in \(\Sh(X_{\mathrm{pro}\acute{e}t})\), or every epimorphism \(F\to G\) induces a surjection \(F(W)\to G(W)\).
\end{bpstatement}

% text: prose-02-18
\phantomsection\label{bp56}%
This is the categorical definition of [\hyperref[bp705]{BS15}, §3.2].

% object: example-weak-contractibility-point (blueprint/chap-The-pro-etale-site/sec-Weakly-contractible-objects/example-weak-contractibility-point/_section.md)
\phantomsection\label{bp57}%
\begin{bpstatement}{Example~2.19}{ (Weak contractibility over a point)}
Let \(X\) be a geometric point, so that the affine objects of \(X_{\mathrm{pro}\acute{e}t}\) are the profinite sets (Proposition \hyperref[bp153]{2.50} below). A continuous surjection \(S\to T\) of profinite sets is a covering, so \(h_S\to h_T\) is an epimorphism, and by the Yoneda lemma a section of it is a continuous map \(T\to S\) that is a section of the surjection. If \(T\) is finite, such a section exists: choose a preimage of each point; it is continuous because \(T\) is discrete. If \(T=\{1,2,3,\ldots\}\cup\{\infty\}\) is a convergent sequence, it need not exist. Let \(S=T_0\amalg T_1\), where \(T_0\) is the convergent sequence \(\{2,4,6,\ldots\}\cup\{\infty_0\}\) and \(T_1\) is \(\{1,3,5,\ldots\}\cup\{\infty_1\}\), and map \(S\to T\) by the identity on integers and \(\infty_0,\infty_1\mapsto\infty\). A section must send each integer \(n\) to its only preimage \(n\), and then the even integers converge to \(\infty_0\) while the odd ones converge to \(\infty_1\), so no value at \(\infty\) makes the section continuous. Hence \(T\) is not weakly contractible. In agreement with Gleason's theorem, \(T\) is not extremally disconnected: the set of odd integers is open, and its closure, which contains \(\infty\), is not open, because every neighbourhood of \(\infty\) contains even integers.
\end{bpstatement}

% text: prose-02-19
\phantomsection\label{bp58}%
Note that the splitting condition concerns a covering family as a whole. If \(W\) is disconnected, a lift of \(\mathrm{id}_{h_W}\) along a covering family may use different members of the family on different clopen pieces of \(W\).

By Lemma \hyperref[bp49]{2.16}, the spectrum of a w-contractible ring is weakly contractible. By Theorem \hyperref[bp51]{2.17}, such affines cover every object of \(X_{\mathrm{pro}\acute{e}t}\). For abelian sheaves we use these two facts in the following form.

% object: corollary-evaluation-weakly-contractible-objects (blueprint/chap-The-pro-etale-site/sec-Weakly-contractible-objects/corollary-evaluation-weakly-contractible-objects/_section.md)
\phantomsection\label{bp59}%
\begin{bpstatementit}{Corollary~2.20}{ (Evaluation on weakly contractible objects)}
Let \(W\) be a weakly contractible object of \(X_{\mathrm{pro}\acute{e}t}\), for instance \(W=\operatorname{Spec}(D)\) with \(D\) w-contractible.

\begin{enumerate}
\item The functor \(F\mapsto F(W)\) from abelian sheaves on \(X_{\mathrm{pro}\acute{e}t}\) to abelian groups is exact. Hence \(H^q(W,F)=0\) for \(q>0\), and \(R\Gamma(W,K)\simeq K(W)\) for every complex \(K\), with the right side computed term by term.
\item If \(W\) is affine, evaluation at \(W\) preserves all limits and all colimits of \(\mathcal R\)-modules.
\item Exactness of complexes and isomorphisms of \(\mathcal R\)-module sheaves can be tested by evaluation on all affine weakly contractible objects of \(X_{\mathrm{pro}\acute{e}t}\).
\end{enumerate}
\end{bpstatementit}

% proof: proof/corollary-evaluation-weakly-contractible-objects
\begin{proof}[{Proof}]\phantomsection\label{bp60}%
All three parts follow from the fact that every epimorphism onto a sheaf is surjective on sections over \(W\), together with the fact that such \(W\) cover every object. Exactness and conservativity follow directly, and for colimits we need in addition that \(W\) is quasi-compact.

(1) Taking sections is left exact for any object of a site, since it preserves kernels. By weak contractibility every epimorphism of sheaves is surjective on sections over \(W\), so the sections functor is also right exact. An exact functor carries acyclic complexes to acyclic complexes, and so it agrees with its derived functor, also for unbounded complexes. In particular \(H^q(W,F)=0\) for \(q>0\).

(2) Limits of sheaves are computed on sections, so only colimits need an argument. We first consider a direct sum \(\bigoplus_iM_i\). The direct sum sheaf is the sheafification of the presheaf direct sum, so a section over \(W\) is, locally on \(W\), a finite sum of sections of the \(M_i\). Since \(W\) is affine, hence quasi-compact as an object of the site, a finite covering suffices, and a single finite set of indices \(i\) serves for all members of this covering. By the sheaf condition, the section is then a section of the corresponding finite subsum, and finite sums commute with evaluation. It follows that evaluation preserves arbitrary direct sums [Stacks, \href{https://stacks.math.columbia.edu/tag/0935}{0935}]. An arbitrary colimit of modules is the cokernel of a map between direct sums, and evaluation preserves cokernels by (1). Hence it preserves all colimits.

(3) Suppose that a sequence \(F'\to F\to F''\) with composite zero becomes exact on all affine weakly contractible \(W\). A local section of the kernel of \(F\to F''\) is, after restriction to a covering by such \(W\), the image of a section of \(F'\), because these objects cover every object. Hence it lies in the image sheaf, and the sequence is exact. Applied to kernels and cokernels of a map, this criterion shows that a map which is an isomorphism on all such \(W\) is an isomorphism, and applied in each degree it gives the statement for complexes; see [Stacks, \href{https://stacks.math.columbia.edu/tag/0945}{0945}].
\end{proof}

% text: prose-02-20
\phantomsection\label{bp61}%
The quasi-compactness used here is quasi-compactness as an object of the site: coverings admit finite refinements. Affine objects of \(X_{\mathrm{pro}\acute{e}t}\) have this property. The affine weakly contractible objects are quasi-compact, quasi-separated and projective, and they cover every object.

Cohomology on other objects need not vanish. For instance, a fibre product of two weakly contractible affines need not be weakly contractible, and for this reason the Čech complex of a cover by them does not in general compute cohomology.


% object: sec-Repleteness-and-limits (blueprint/chap-The-pro-etale-site/sec-Repleteness-and-limits/_section.md), content hidden
\subsection*{\texorpdfstring{2.5\quad Repleteness and limits}{2.5 Repleteness and limits}}
\phantomsection\label{bp62}%
\addcontentsline{toc}{subsection}{\texorpdfstring{2.5\quad Repleteness and limits}{2.5 Repleteness and limits}}
% text: prose-02-21
\phantomsection\label{bp63}%
We need inverse limits to pass from torsion to adic coefficients, for example to form \(\varprojlim_n\mathbf Z/\ell^n\), and the first use of weakly contractible objects is to show that such limits are well behaved. In a general topos a limit of epimorphisms need not be an epimorphism, because the local lifts at each stage may require ever finer covers. On \(X_{\mathrm{pro}\acute{e}t}\), however, this does not happen, since surjectivity and exactness can be checked on sections over weakly contractible objects, and there no refinement is needed. As a result, the limit of a tower of epimorphisms maps epimorphically onto each term (Proposition \hyperref[bp65]{2.22}). For the same reason products are exact (Proposition \hyperref[bp68]{2.23}).

Exactness of products means that derived limits can be computed as for abelian groups. Namely, the derived limit of a tower is the fibre of \(1-T\) on the product. Its cohomology is therefore given by the Milnor sequence of Proposition \hyperref[bp71]{2.24}, and \(R^p\varprojlim\) vanishes for \(p\geq2\). For a tower of epimorphisms, such as \((\mathbf Z/\ell^n)_X\), moreover \(R^1\varprojlim\) vanishes, and the derived limit is the ordinary limit. In the next section we also deduce left completeness from the Milnor sequence. With these facts we compute the cohomology of \(\mathbf Z_\ell\) as continuous cohomology (Proposition \hyperref[bp172]{2.55}).

% object: definition-repleteness (blueprint/chap-The-pro-etale-site/sec-Repleteness-and-limits/definition-repleteness/_section.md)
\phantomsection\label{bp64}%
\begin{bpstatement}{Definition~2.21}{ (Repleteness)}
A topos \(\mathcal T\) is \emph{replete} if, for every inverse sequence \(\cdots\to F_2\to F_1\to F_0\) of sheaves of sets with epimorphic transition maps, the projection \(\varprojlim_nF_n\to F_m\) is an epimorphism for every \(m\).
\end{bpstatement}

% object: proposition-repleteness-pro-etale-topos (blueprint/chap-The-pro-etale-site/sec-Repleteness-and-limits/proposition-repleteness-pro-etale-topos/_section.md)
\phantomsection\label{bp65}%
\begin{bpstatementit}{Proposition~2.22}{ (Repleteness of the pro-étale topos)}
For every scheme \(X\), the topos \(\Sh(X_{\mathrm{pro}\acute{e}t})\) is replete.
\end{bpstatementit}

% proof: proof/proposition-repleteness-pro-etale-topos
\begin{proof}[{Proof}]\phantomsection\label{bp66}%
Evaluate the sequence on a weakly contractible \(W\). Since every epimorphism of sheaves is surjective on sections over \(W\), the transition maps \(F_{n+1}(W)\to F_n(W)\) are surjective. Given a section at stage \(m\), we choose lifts successively at stages \(m+1,m+2,\ldots\) (by countable choice), while the earlier stages are determined by the transition maps. This gives a compatible section of \(\varprojlim_nF_n(W)\). Since limits are computed on sections, the projection is surjective on \(W\).

An arbitrary \(U\) is covered by affine weakly contractible objects, so every section of \(F_m\) over \(U\) lifts to \(\varprojlim_nF_n\) on each member of such a cover. Hence the projection is an epimorphism of sheaves; see [\hyperref[bp705]{BS15}, §§3.1–3.2].
\end{proof}

% text: prose-02-22
\phantomsection\label{bp67}%
For sheaves of modules the same lifting argument, applied coordinatewise rather than along a tower, shows that products are exact. This lets us compute derived limits by termwise products of complexes.

% object: proposition-exact-products (blueprint/chap-The-pro-etale-site/sec-Repleteness-and-limits/proposition-exact-products/_section.md)
\phantomsection\label{bp68}%
\begin{bpstatementit}{Proposition~2.23}{ (Exact products)}
Arbitrary set-indexed products are exact in \(\operatorname{Mod}(X_{\mathrm{pro}\acute{e}t},\mathcal R)\). Products in \(D(X_{\mathrm{pro}\acute{e}t},\mathcal R)\) are computed by taking termwise products of representative complexes, and their cohomology sheaves are the products of the cohomology sheaves.
\end{bpstatementit}

% proof: proof/proposition-exact-products
\begin{proof}[{Proof}]\phantomsection\label{bp69}%
By Corollary corollary-evaluation-weakly-contractible-objects, exactness of sheaves can be checked on affine weakly contractible \(W\), and by part (1) evaluation at such \(W\) is exact. Since limits are computed on sections, a product of short exact sequences therefore evaluates on \(W\) to a product of short exact sequences of \(\mathcal R(W)\)-modules, which is exact because by the axiom of choice we may lift a tuple coordinate by coordinate. This proves the first assertion.

The derived product is the termwise product of K-injective representatives (complexes \(I\) such that maps from acyclic complexes to \(I\) are null-homotopic, the unbounded analogue of bounded-below injective complexes). Since products are exact, a product of quasi-isomorphisms is a quasi-isomorphism: its cone is the product of the cones, which is acyclic. We may therefore replace the K-injective representatives by any quasi-isomorphic ones. Hence termwise products compute the product in \(D(X_{\mathrm{pro}\acute{e}t},\mathcal R)\), and they commute with cohomology because products are exact; see [Stacks, \href{https://stacks.math.columbia.edu/tag/0947}{0947}, \href{https://stacks.math.columbia.edu/tag/08TB}{08TB}].
\end{proof}

% text: prose-02-23
\phantomsection\label{bp70}%
Derived limits of countable towers have an explicit model. For a countable tower of modules, write \(t_n:M_{n+1}\to M_n\) and let \(T:\prod_nM_n\to\prod_nM_n\) be \(T((m_n)_n)=(t_n(m_{n+1}))_n\). For a tower \((K_n)\) in \(D(X_{\mathrm{pro}\acute{e}t},\mathcal R)\), represented by a tower of complexes \(K_{n+1}\to K_n\), the same formula defines \(T\) termwise, and the derived inverse limit is the homotopy fibre of \(1-T\)

\[
R\varprojlim_n K_n
=
\fib\Bigl(
\prod_n K_n\xrightarrow{\,1-T\,}\prod_n K_n
\Bigr)=\Cone(1-T)[-1].
\tag{2.1}
\]

The cone is formed on the level of complexes, so that the construction is functorial in the tower. By Proposition \hyperref[bp68]{2.23} the termwise products are the derived products, and so, by [Stacks, \href{https://stacks.math.columbia.edu/tag/0940}{0940}], this cone agrees with the derived inverse limit.

% object: proposition-milnor-sequence (blueprint/chap-The-pro-etale-site/sec-Repleteness-and-limits/proposition-milnor-sequence/_section.md)
\phantomsection\label{bp71}%
\begin{bpstatementit}{Proposition~2.24}{ (The Milnor sequence)}
For a countable tower of \(\mathcal R\)-modules, \(R^p\varprojlim=0\) for \(p>1\), and

\[
R^1\varprojlim_n M_n
=\coker\left(\prod_n M_n\xrightarrow{\,1-T\,}\prod_n M_n\right).
\]

For a tower \((K_n)\) of complexes there is a short exact sequence of sheaves

\[
0\longrightarrow R^1\varprojlim_n H^{q-1}(K_n)
\longrightarrow H^q(R\varprojlim_nK_n)
\longrightarrow \varprojlim_n H^q(K_n)
\longrightarrow 0.
\tag{2.2}
\]

If the transitions of a tower of modules are epimorphisms, its derived limit is its ordinary limit in degree zero.
\end{bpstatementit}

% proof: proof/proposition-milnor-sequence
\begin{proof}[{Proof}]\phantomsection\label{bp72}%
Everything follows from the triangle (2.1). Indeed, for a tower of modules, \(R\varprojlim_nM_n\) is the fibre of a map between two sheaves in degree \(0\), so it has cohomology only in degrees \(0\) and \(1\), namely the kernel and the cokernel of \(1-T\). Here the kernel is \(\varprojlim_nM_n\), and the cokernel is \(R^1\varprojlim_nM_n\), which proves the first two assertions. Similarly, for a tower of complexes, the long exact cohomology sequence of (2.1) reads

\[
\prod_nH^{q-1}(K_n)\xrightarrow{1-T}\prod_nH^{q-1}(K_n)
\to H^q(R\varprojlim_nK_n)\to\prod_nH^q(K_n)\xrightarrow{1-T}\prod_nH^q(K_n),
\]

where we used that by Proposition \hyperref[bp68]{2.23} the cohomology of each product is the product of the cohomology sheaves. The cokernel on the left is \(R^1\varprojlim_nH^{q-1}(K_n)\) and the kernel on the right is \(\varprojlim_nH^q(K_n)\), which gives (2.2).

Suppose now that the transition maps of a tower of modules are epimorphisms. We show that \(1-T\) is an epimorphism; then \(R^1\varprojlim_nM_n=0\) and the derived limit is \(\varprojlim_nM_n\). By Corollary \hyperref[bp59]{2.20} it suffices to check surjectivity on sections over every weakly contractible \(W\), where the maps \(t_n\) are surjective. Given \((a_n)\), choose \(m_0\) arbitrarily and then, recursively, \(m_{n+1}\) with \(t_n(m_{n+1})=m_n-a_n\). Then \((1-T)(m_n)_n=(a_n)_n\). See [Stacks, \href{https://stacks.math.columbia.edu/tag/0947}{0947}].
\end{proof}

% object: example-nonzero-rlim-one (blueprint/chap-The-pro-etale-site/sec-Repleteness-and-limits/example-nonzero-rlim-one/_section.md)
\phantomsection\label{bp73}%
\begin{bpstatement}{Example~2.25}{ (A nonzero $R^1\varprojlim$)}
Without surjective transition maps the derived limit can be nonzero in degree one. To see this, work over a geometric point and evaluate on the final object, which is weakly contractible, so that we may compute with abelian groups. Let \(M_n=\mathbf Z\) for all \(n\), with transition maps multiplication by \(\ell\). Then the maps \(M_n\to\mathbf Z\), \(m\mapsto\ell^nm\), form a map to the constant tower \(\mathbf Z\) whose cokernel is the tower \((\mathbf Z/\ell^n)\). The constant tower has limit \(\mathbf Z\) and no \(R^1\varprojlim\), and \(\varprojlim_nM_n=0\), because an element of \(M_0\) in the image of every \(M_n\) is divisible by every power of \(\ell\). The long exact sequence of \(R\varprojlim\) therefore reads

\[
0\longrightarrow\mathbf Z\longrightarrow\mathbf Z_\ell\longrightarrow R^1\varprojlim_nM_n\longrightarrow0,
\]

so the sections of \(R^1\varprojlim_nM_n\) over the point form \(\mathbf Z_\ell/\mathbf Z\), an uncountable group. If the \(M_n\) are the groups \(H^{q-1}\) of a tower of complexes, this group appears in \(H^q\) of the derived limit by (2.2).
\end{bpstatement}

% text: prose-02-24
\phantomsection\label{bp74}%
More generally, the \(R^1\varprojlim\) term vanishes if, on every affine weakly contractible \(W\), the tower of section modules satisfies the Mittag–Leffler condition. To see this, note that evaluation on \(W\) preserves cokernels (Corollary \hyperref[bp59]{2.20}). So the cokernel of \(1-T\) evaluated on \(W\) is the cokernel of \(1-T\) on the tower \((M_n(W))\), which vanishes by the Mittag–Leffler lemma [Stacks, \href{https://stacks.math.columbia.edu/tag/07KV}{07KV}]. As evaluation on these \(W\) detects vanishing, the \(R^1\varprojlim\) term vanishes.


% object: sec-Left-completeness-and-hypercovers (blueprint/chap-The-pro-etale-site/sec-Left-completeness-and-hypercovers/_section.md), content hidden
\subsection*{\texorpdfstring{2.6\quad Left completeness and hypercovers}{2.6 Left completeness and hypercovers}}
\phantomsection\label{bp75}%
\addcontentsline{toc}{subsection}{\texorpdfstring{2.6\quad Left completeness and hypercovers}{2.6 Left completeness and hypercovers}}
% text: prose-02-25
\phantomsection\label{bp76}%
An unbounded complex on \(X_{\mathrm{pro}\acute{e}t}\) is determined by its bounded-below truncations, and its cohomology is computed from theirs by a derived limit. On the étale site of a scheme of infinite cohomological dimension this can fail, whereas on \(X_{\mathrm{pro}\acute{e}t}\) it holds because products are exact, so that the Milnor sequence applies.

Left completeness (Proposition \hyperref[bp77]{2.26}) lets us treat an unbounded complex as the derived limit of bounded below ones, to which injective resolutions and the comparison theorems apply. In this way we prove the unbounded comparison of Proposition \hyperref[bp146]{2.49}.

For bounded below complexes we also give an explicit model of cohomology. On a quasi-compact quasi-separated object \(U\), Proposition \hyperref[bp83]{2.28} provides a hypercover all of whose terms \(W_r\) are affine weakly contractible. Since these terms have no higher cohomology, the hypercover spectral sequence collapses, and \(R\Gamma(U,K)\) is the total complex of the sections \(K^b(W_r)\). This model has a consequence we use later. Namely, in each total degree only finitely many terms occur, and evaluation on each \(W_r\) commutes with filtered colimits. Hence cohomology commutes with filtered colimits of uniformly bounded below complexes (Corollary \hyperref[bp86]{2.29}). We use this to show that cohomology and direct images commute with inverting \(\ell\) and with extension of scalars to \(\overline{\mathbf Q}_\ell\), the colimit over the finite extensions \(E'\) of \(E\).

% object: proposition-left-completeness (blueprint/chap-The-pro-etale-site/sec-Left-completeness-and-hypercovers/proposition-left-completeness/_section.md)
\phantomsection\label{bp77}%
\begin{bpstatementit}{Proposition~2.26}{ (Left completeness)}
The standard \(t\)-structure on \(D(X_{\mathrm{pro}\acute{e}t},\mathcal R)\) is left complete: for every complex \(K\) the natural map

\[
K\longrightarrow R\varprojlim_{n\geq0}\tau_{\geq-n}K
\]

is an isomorphism.
\end{bpstatementit}

% proof: proof/proposition-left-completeness
\begin{proof}[{Proof}]\phantomsection\label{bp78}%
We compute the cohomology sheaves of the derived limit with the Milnor sequence and find that they are those of \(K\).

Represent \(K\) by a complex and form its canonical truncation tower, with maps \(K\to\tau_{\geq-n}K\) and \(\tau_{\geq-(n+1)}K\to\tau_{\geq-n}K\). For fixed \(q\), the sheaves \(H^q(\tau_{\geq-n}K)\) and \(H^{q-1}(\tau_{\geq-n}K)\) are eventually constant, with identity transition maps, equal to the corresponding cohomology sheaves of \(K\), so the \(R^1\) term in (2.2) vanishes and the limit term is \(H^q(K)\). As the natural map induces this identification in every degree, it is an isomorphism; see Stacks, [0947](https://stacks.math.columbia.edu/tag/0947).
\end{proof}

% text: prose-02-26
\phantomsection\label{bp79}%
To compute the cohomology of an unbounded complex from its truncations, we need derived functors on the whole unbounded derived category that commute with derived inverse limits.

% object: proposition-unbounded-functors-products (blueprint/chap-The-pro-etale-site/sec-Left-completeness-and-hypercovers/proposition-unbounded-functors-products/_section.md)
\phantomsection\label{bp80}%
\begin{bpstatementit}{Proposition~2.27}{ (Unbounded derived functors preserve products and limits)}
Let \(U\) be an object of \(X_{\mathrm{pro}\acute{e}t}\) and \(f:X\to Y\) a morphism of schemes. The functor \(R\Gamma(U,-)\), and the functor \(Rf_*\) from \(D(X_{\mathrm{pro}\acute{e}t},\mathcal R)\) to the derived category of sheaves of \(f_*\mathcal R\)-modules on \(Y_{\mathrm{pro}\acute{e}t}\), are defined on the unbounded derived categories. They preserve products and derived inverse limits.
\end{bpstatementit}

% proof: proof/proposition-unbounded-functors-products
\begin{proof}[{Proof}]\phantomsection\label{bp81}%
K-injective resolutions of arbitrary complexes exist in the Grothendieck category of \(\mathcal R\)-modules [Stacks, \href{https://stacks.math.columbia.edu/tag/079P}{079P}], which defines both functors. Both are right adjoints on derived categories: \(R\Gamma(U,-)\) is right adjoint to the functor sending a complex of abelian groups \(M\) to \(\mathcal R[h_U]\otimes^{\mathbf L}_{\mathbf Z}M\), where \(\mathcal R[h_U]\) is the free \(\mathcal R\)-module on \(h_U\), and \(Rf_*\) is right adjoint to the derived pullback of modules. Right adjoints preserve products. Since (2.1) is built from products and a cone, an exact functor preserving products also preserves derived inverse limits [Stacks, \href{https://stacks.math.columbia.edu/tag/0940}{0940}].
\end{proof}

% text: prose-02-27
\phantomsection\label{bp82}%
For bounded-below complexes, weakly contractible objects provide explicit cochain complexes computing cohomology.

% object: proposition-weakly-contractible-hypercovers (blueprint/chap-The-pro-etale-site/sec-Left-completeness-and-hypercovers/proposition-weakly-contractible-hypercovers/_section.md)
\phantomsection\label{bp83}%
\begin{bpstatementit}{Proposition~2.28}{ (Weakly contractible hypercovers)}
Let \(U\) be a quasi-compact, quasi-separated object of \(X_{\mathrm{pro}\acute{e}t}\). There is a hypercover \(W_\bullet\to U\) whose term in every simplicial degree is one affine weakly contractible object. If \(K\in D^+(X_{\mathrm{pro}\acute{e}t},\mathcal R)\) is represented by a complex \(K^\bullet\) vanishing below degree \(a\), then

\[
R\Gamma(U,K)
\simeq \operatorname{Tot}\bigl(K^b(W_r)\bigr)_{r\geq0,\ b\geq a}.
\tag{2.3}
\]

The cosimplicial differential is the alternating sum of pullbacks along face maps, combined with \((-1)^r\) times the differential of \(K^\bullet\). In total degree \(q\) this complex has the finite sum \(\bigoplus_{0\leq r\leq q-a}K^{q-r}(W_r)\).
\end{bpstatementit}

% proof: proof/proposition-weakly-contractible-hypercovers
\begin{proof}[{Proof}]\phantomsection\label{bp84}%
We build a hypercover whose terms have no higher cohomology, so that the hypercover spectral sequence collapses to the complex of sections. By quasi-compactness each term can be taken to be a single affine.

We start with a weakly contractible cover of \(U\), and at each stage we cover the matching object by weakly contractible affines, as for fibre products above. The matching objects are finite limits of quasi-compact, quasi-separated schemes, so they are again quasi-compact and quasi-separated and are covered by finitely many affine weakly contractible objects. Taking finite disjoint unions, we get a single affine term in each degree; see [Stacks, \href{https://stacks.math.columbia.edu/tag/09A1}{09A1}].

For a sheaf \(M\), the hypercover spectral sequence has \(E_1^{r,s}=H^s(W_r,M)\), and the rows with \(s>0\) vanish by Corollary corollary-evaluation-weakly-contractible-objects. The remaining row is the cosimplicial complex of sections. For bounded-below \(K^\bullet\), choose a bounded-below injective resolution \(K^\bullet\to I^\bullet\). Since evaluation on \(W_r\) is exact, \(K^\bullet(W_r)\to I^\bullet(W_r)\) is a quasi-isomorphism. After a shift, both double complexes lie in the first quadrant, so their total complexes are quasi-isomorphic, and hypercover descent for \(I^\bullet\) identifies the second with \(R\Gamma(U,K)\). See [Stacks, \href{https://stacks.math.columbia.edu/tag/09A2}{09A2}].
\end{proof}

% text: prose-02-28
\phantomsection\label{bp85}%
For a single sheaf this is the cochain complex \(M(W_0)\to M(W_1)\to M(W_2)\to\cdots\). It may be nonzero in every degree, so the hypercover gives no bound on cohomological dimension. For unbounded \(K\), by left completeness (Proposition \hyperref[bp77]{2.26}) and since \(R\Gamma(U,-)\) commutes with derived inverse limits, we get \(R\Gamma(U,K)\simeq
R\varprojlim_n R\Gamma(U,\tau_{\geq-n}K)\), and the hypercover complex computes each term.

% object: corollary-continuity-cohomology (blueprint/chap-The-pro-etale-site/sec-Left-completeness-and-hypercovers/corollary-continuity-cohomology/_section.md)
\phantomsection\label{bp86}%
\begin{bpstatementit}{Corollary~2.29}{ (Continuity of cohomology)}
Suppose \(U\) is quasi-compact and quasi-separated and \((K_i^\bullet)\) is a filtered diagram of complexes of \(\mathcal R\)-modules vanishing below one fixed degree \(a\). Then for every \(q\),

\[
\varinjlim_i H^q(U,K_i^\bullet)
\simeq H^q(U,\varinjlim_iK_i^\bullet).
\]

If \(f:X\to Y\) is quasi-compact and quasi-separated, the analogous map \(\varinjlim_iR^qf_*K_i^\bullet\to
R^qf_*(\varinjlim_iK_i^\bullet)\) is an isomorphism.
\end{bpstatementit}

% proof: proof/corollary-continuity-cohomology
\begin{proof}[{Proof}]\phantomsection\label{bp87}%
We compute both sides with one hypercover by weakly contractible affines, chosen independently of \(i\). Because of the common lower bound \(a\), each total degree \(q\) of (2.3) is a finite sum, of the terms \(K_i^{q-r}(W_r)\) with \(0\leq r\leq q-a\), and the range of \(r\) is the same for all \(i\).

By Corollary \hyperref[bp59]{2.20}, sections on each \(W_r\) commute with filtered colimits, hence so does the finite sum in each total degree of (2.3), and the hypercover complex of \(\varinjlim_iK_i^\bullet\) is the filtered colimit of the hypercover complexes of the \(K_i^\bullet\). The first assertion follows because filtered colimits of modules are exact and so commute with cohomology.

The second assertion reduces to the first. Recall that \(R^qf_*K\) is the sheaf associated to the presheaf \(V\mapsto H^q(X\times_YV,K)\), and that the affine objects \(V\) of \(Y_{\mathrm{pro}\acute{e}t}\) form a basis. For such \(V\) the object \(X\times_YV\) is quasi-compact and quasi-separated, so the first assertion applies to it, and sheafification commutes with colimits.
\end{proof}

% text: prose-02-29
\phantomsection\label{bp88}%
In particular, \(H^q(U,-)\) commutes with direct sums of sheaves of modules. Quasi-compactness of \(U\) is needed already for \(H^0\). Let \(X\) be a geometric point and \(U=\coprod_{n\in\mathbf N}X\), and consider the direct sum \(\bigoplus_{i\in\mathbf N}\mathbf Z\) of copies of the constant sheaf. Its sections over \(U\) are \(\prod_n\bigoplus_i\mathbf Z\), which contains the element whose \(n\)th component is the \(n\)th basis vector, whereas \(\bigoplus_i H^0(U,\mathbf Z)=\bigoplus_i\prod_n\mathbf Z\) consists of the elements with only finitely many nonzero indices \(i\).

For direct sums of unbounded complexes the common lower bound is replaced by uniform bounds on cohomological dimension (Proposition \hyperref[bp146]{2.49}).


% object: sec-Classical-sheaves (blueprint/chap-The-pro-etale-site/sec-Classical-sheaves/_section.md), content hidden
\subsection*{\texorpdfstring{2.7\quad Classical sheaves}{2.7 Classical sheaves}}
\phantomsection\label{bp89}%
\addcontentsline{toc}{subsection}{\texorpdfstring{2.7\quad Classical sheaves}{2.7 Classical sheaves}}
% text: prose-02-30
\phantomsection\label{bp90}%
From here until Section \hyperref[bp170]{2.13}, \(\Lambda\) denotes a discrete commutative ring and \(\Lambda_X\) its constant sheaf on \(X_{\mathrm{pro}\acute{e}t}\). To use the torsion arguments of étale cohomology, we must know which pro-étale sheaves come from étale sheaves. The main result answers this: étale sheaves embed fully faithfully, and the essential image consists of the sheaves whose sections on a limit \(\varprojlim U_i\) are the colimit of the sections on the \(U_i\). In particular, the reductions \(\mathbf Z/\ell^n\) lie in this image, while the \(\ell\)-adic coefficient sheaves themselves, such as \(\mathbf Z_{\ell,X}\), do not (Example \hyperref[bp97]{2.33}).

The colimit condition is forced on us by the way étale sheaves extend to limits. If \(F\) is an étale sheaf and \(\bar x\) a geometric point, the only reasonable value of the pullback of \(F\) on the strict henselization \(\operatorname{Spec}\mathcal O^{\mathrm{sh}}_{X,\bar x}=\varprojlim(U,\bar u)\) is \(\varinjlim F(U)=F_{\bar x}\). In this way stalks of étale sheaves become sections of pro-étale sheaves. Proposition \hyperref[bp92]{2.31} defines the pullback \(\nu_X^*\) by this rule, and Theorem \hyperref[bp95]{2.32} characterizes its image by it. The part that is not formal is cohomology, because pro-étale covers are not étale covers. In the next section we show that they create no new cohomology classes for classical sheaves.

% object: definition-classical-pro-etale-sheaves (blueprint/chap-The-pro-etale-site/sec-Classical-sheaves/definition-classical-pro-etale-sheaves/_section.md)
\phantomsection\label{bp91}%
\begin{bpstatement}{Definition~2.30}{ (Classical pro-étale sheaves)}
A sheaf \(F\) on \(X_{\mathrm{pro}\acute{e}t}\) is \emph{classical} if, for every pro-étale affine presentation \(U=\varprojlim_i U_i\) with \(U_i\) affine and étale over \(X\), its restriction maps induce a bijection

\[
\varinjlim_i F(U_i)\xrightarrow{\ \sim\ }F(U).
\tag{2.4}
\]

A complex is classical if all its cohomology sheaves are classical.
\end{bpstatement}

% object: proposition-comparison-morphism-inverse-image (blueprint/chap-The-pro-etale-site/sec-Classical-sheaves/proposition-comparison-morphism-inverse-image/_section.md)
\phantomsection\label{bp92}%
\begin{bpstatementit}{Proposition~2.31}{ (The comparison morphism and its inverse image)}
Inclusion of étale objects defines a morphism of topoi

\[
\nu_X:X_{\mathrm{pro}\acute{e}t}\longrightarrow X_{\acute{e}t}.
\]

Its direct image \(\nu_{X,*}:\Sh(X_{\mathrm{pro}\acute{e}t})\to\Sh(X_{\acute{e}t})\) is restriction to the étale subsite. Its inverse image \(\nu_X^*:\Sh(X_{\acute{e}t})\to\Sh(X_{\mathrm{pro}\acute{e}t})\) is exact and has the description

\[
(\nu_X^*F)(h:V\to X)
=\Gamma(V_{\acute{e}t},h_{\acute{e}t}^*F),
\tag{2.5}
\]

where \(h_{\acute{e}t}^*:\Sh(X_{\acute{e}t})\to\Sh(V_{\acute{e}t})\) is the étale inverse image along the weakly étale map \(h\). Both functors carry sheaves of \(\Lambda\)-modules to sheaves of \(\Lambda\)-modules, so they restrict to functors between \(\operatorname{Mod}(X_{\acute{e}t},\Lambda)\) and \(\operatorname{Mod}(X_{\mathrm{pro}\acute{e}t},\Lambda)\). In particular, for \(U=\varprojlim_iU_i\) as above, we have

\[
(\nu_X^*F)(U)=\varinjlim_i F(U_i).
\tag{2.6}
\]
\end{bpstatementit}

% proof: proof/proposition-comparison-morphism-inverse-image
\begin{proof}[{Proof}]\phantomsection\label{bp93}%
The morphism \(\nu_X\) is formal. For the formula, we write down the right-hand side of (2.5), check that it is a sheaf, and check that it has the universal property of the left adjoint of \(\nu_{X,*}\).

Since an étale morphism is weakly étale and an étale covering is a pro-étale covering, the inclusion is a continuous functor preserving finite limits, and it defines \(\nu_X\). For (2.5), we define the right-hand side for every weakly étale \(V\). Since pro-étale coverings are fpqc coverings, fpqc descent for étale sheaves shows that this rule is already a sheaf, so no sheafification is needed. If we pull back the restriction of a pro-étale sheaf \(G\) to \(X_{\acute{e}t}\), we get the counit \(\nu_X^*\nu_{X,*}G\to G\). On an étale object \(V\) the rule returns \(F(V)\), and this is the unit \(F\to\nu_{X,*}\nu_X^*F\). The two adjunction identities follow by restriction and pullback, so the rule is \(\nu_X^*\) [Stacks, \href{https://stacks.math.columbia.edu/tag/0GLZ}{0GLZ}].

For an affine limit \(U\), formula (2.6) follows from the fact that sections of an étale sheaf commute with affine inverse limits [Stacks, \href{https://stacks.math.columbia.edu/tag/099S}{099S}].
\end{proof}

% text: prose-02-31
\phantomsection\label{bp94}%
For \(X=\operatorname{Spec}\mathbf F_q\) and \(U=\operatorname{Spec}\overline{\mathbf F}_q\), formula (2.6) says that \((\nu_X^*F)(U)\) is the colimit of the \(F(\operatorname{Spec}\mathbf F_{q^m})\), that is, the underlying set of the Galois module corresponding to \(F\). Formula (2.6) also shows at once that every sheaf \(\nu_X^*F\) is classical. The next theorem says that there are no other classical sheaves.

% object: theorem-classical-sheaves-essential-image (blueprint/chap-The-pro-etale-site/sec-Classical-sheaves/theorem-classical-sheaves-essential-image/_section.md)
\phantomsection\label{bp95}%
\begin{bpstatementit}{Theorem~2.32}{ (Classical sheaves as the essential image)}
The functor \(\nu_X^*\) is fully faithful. Its essential image is exactly the classical sheaves of Definition \hyperref[bp91]{2.30}. For modules it is an exact equivalence between \(\operatorname{Mod}(X_{\acute{e}t},\Lambda)\) and the full subcategory of classical \(\Lambda_X\)-module sheaves on \(X_{\mathrm{pro}\acute{e}t}\).
\end{bpstatementit}

% proof: proof/theorem-classical-sheaves-essential-image
\begin{proof}[{Proof}]\phantomsection\label{bp96}%
A left adjoint is fully faithful exactly when its unit is an isomorphism, and the essential image is then the set of objects on which the counit is an isomorphism. We check both on pro-étale affines, where (2.6) computes \(\nu_X^*\).

Evaluating (2.5) on an étale object \(V\) gives \(F(V)\), so \(F\xrightarrow{\sim}\nu_{X,*}\nu_X^*F\), which proves full faithfulness. Moreover, every \(\nu_X^*F\) is classical by (2.6). Conversely, let \(G\) be classical. On a pro-étale affine the counit \(\nu_X^*\nu_{X,*}G\to G\) is

\[
(\nu_X^*\nu_{X,*}G)(U)
=\varinjlim_iG(U_i)\longrightarrow G(U),
\]

which is an isomorphism by (2.4). Since pro-étale affines form a covering basis (Proposition \hyperref[bp46]{2.15}), the counit is an isomorphism of sheaves.
\end{proof}

% object: example-coefficient-sheaf-outside-classical-part (blueprint/chap-The-pro-etale-site/sec-Classical-sheaves/example-coefficient-sheaf-outside-classical-part/_section.md)
\phantomsection\label{bp97}%
\begin{bpstatement}{Example~2.33}{ (A coefficient sheaf outside the classical part)}
Let \(X\) be a geometric point and let \(S=\mathbf Z_\ell\) be regarded as a profinite object of \(X_{\mathrm{pro}\acute{e}t}\). The completed coefficient sheaf \(\mathbf Z_{\ell,X}\) has sections \(\operatorname{Cont}(S,\mathbf Z_\ell)\) on \(S\), among them the identity function. The pullback \(\nu_X^*\) from \(X_{\acute{e}t}\) of the constant sheaf with value the \emph{discrete} abelian group \(\mathbf Z_\ell\) has only locally constant functions as sections, and these have finite image on the compact space \(S\). Hence the identity function does not come from any finite stage of the profinite presentation of \(S\), and \(\mathbf Z_{\ell,X}\) is not classical. Its reductions \(\mathbf Z/\ell^n\), on the other hand, are classical. So étale arguments apply at each finite level, while the coefficient sheaf itself is defined once on the pro-étale site.
\end{bpstatement}


% object: sec-Comparison-with-the-etale-site (blueprint/chap-The-pro-etale-site/sec-Comparison-with-the-etale-site/_section.md), content hidden
\subsection*{\texorpdfstring{2.8\quad Comparison with the étale site}{2.8 Comparison with the étale site}}
\phantomsection\label{bp98}%
\addcontentsline{toc}{subsection}{\texorpdfstring{2.8\quad Comparison with the étale site}{2.8 Comparison with the étale site}}
% text: prose-02-32
\phantomsection\label{bp99}%
For a bounded-below complex \(K\) on \(X_{\acute{e}t}\), the pro-étale cohomology of \(\nu_X^*K\) is the étale cohomology of \(K\) (Theorem \hyperref[bp103]{2.35}). Bounds on cohomological dimension are needed only for unbounded complexes (Section \hyperref[bp135]{2.11}). With this theorem, statements of SGA 4 about the cohomology of torsion sheaves, such as \(\mathbf Z/\ell^n\) on a variety, can be quoted on the pro-étale site.

The theorem rests on two facts. A pro-étale cover of an affine can be refined to a limit of étale covers, and sections of a classical sheaf on a limit are the colimit of sections. It follows that the Čech complex of a pro-étale cover is a filtered colimit of étale Čech complexes. Because of this, an injective étale sheaf stays acyclic after pullback (Lemma \hyperref[bp100]{2.34}), and resolving by injectives gives the theorem. For non-classical coefficients the comparison fails, as the computation of \(H^1(\operatorname{Spec}\mathbf F_q,\mathbf Z_\ell)\) in Section \bpout{1.2} shows.

It follows that the cohomology of a pro-étale affine is the colimit of the cohomology of its étale approximations (Corollary \hyperref[bp106]{2.36}). Moreover, classical complexes in \(D^+\) are exactly the pullbacks of étale ones (Corollary \hyperref[bp109]{2.37}), and this is what we need for proper base change on the pro-étale site.

% object: lemma-acyclicity-pulled-back-injectives (blueprint/chap-The-pro-etale-site/sec-Comparison-with-the-etale-site/lemma-acyclicity-pulled-back-injectives/_section.md)
\phantomsection\label{bp100}%
\begin{bpstatementit}{Lemma~2.34}{ (Acyclicity of pulled-back injectives)}
If \(I\) is an injective sheaf of \(\Lambda\)-modules on \(X_{\acute{e}t}\), then \(R^q\nu_{X,*}\nu_X^*I=0\) for \(q>0\).
\end{bpstatementit}

% proof: proof/lemma-acyclicity-pulled-back-injectives
\begin{proof}[{Proof}]\phantomsection\label{bp101}%
The sheaf \(R^q\nu_{X,*}\nu_X^*I\) is associated to \(V\mapsto H^q(V_{\mathrm{pro}\acute{e}t},\nu_X^*I)\), so it suffices to show that positive cohomology vanishes on affine étale objects \(V\). The restriction of \(I\) to \(V_{\acute{e}t}\) is injective, since restriction has an exact left adjoint, namely extension by zero. Hence we may assume \(X\) affine. By Theorem \hyperref[bp42]{2.14}, every pro-étale covering of an affine object is refined by a finite covering by affines that are ind-étale over it. By the Čech acyclicity criterion for a basis, it therefore suffices to show that the Čech complexes of \(\nu_X^*I\) for such coverings have no positive cohomology.

As in [Stacks, \href{https://stacks.math.columbia.edu/tag/099U}{099U}], we claim that the Čech complex of such a covering is a filtered colimit of Čech complexes of étale coverings. Write each covering map as an inverse limit of affine étale maps over the object being covered. By (2.6), the Čech complex of the covering is the filtered colimit of the Čech complexes of these étale approximations. Each finitely presented approximation descends to a common étale stage over the base, becomes a covering after increasing that stage, and its nerve descends by taking fibre products. Each augmented Čech complex at a finite stage is obtained by applying \(\Hom(-,I)\) to the exact augmented complex of free sheaves of \(\Lambda\)-modules on the nerve, so it has no positive cohomology because \(I\) is injective. Filtered colimits are exact, so the same holds for the pro-étale covering, and the criterion implies the required vanishing.
\end{proof}

% text: prose-02-33
\phantomsection\label{bp102}%
Since injective sheaves compute derived functors, the lemma suffices to compare derived direct images, and we obtain the comparison of cohomology by resolving a bounded-below complex by injectives.

% object: theorem-bounded-below-cohomology-comparison (blueprint/chap-The-pro-etale-site/sec-Comparison-with-the-etale-site/theorem-bounded-below-cohomology-comparison/_section.md)
\phantomsection\label{bp103}%
\begin{bpstatementit}{Theorem~2.35}{ (Bounded-below cohomology comparison)}
For \(K\in D^+(X_{\acute{e}t},\Lambda)\) the unit of the adjunction between \(\nu_X^*:D^+(X_{\acute{e}t},\Lambda)\to D^+(X_{\mathrm{pro}\acute{e}t},\Lambda)\) and \(R\nu_{X,*}\) is an isomorphism

\[
K\xrightarrow{\ \sim\ }R\nu_{X,*}\nu_X^*K.
\tag{2.7}
\]

Consequently \(\nu_X^*:D^+(X_{\acute{e}t},\Lambda)\to D^+(X_{\mathrm{pro}\acute{e}t},\Lambda)\) is fully faithful, and there is a canonical isomorphism

\[
R\Gamma(X_{\mathrm{pro}\acute{e}t},\nu_X^*K)
\simeq R\Gamma(X_{\acute{e}t},K).
\]
\end{bpstatementit}

% proof: proof/theorem-bounded-below-cohomology-comparison
\begin{proof}[{Proof}]\phantomsection\label{bp104}%
We compute \(R\nu_{X,*}\nu_X^*K\) with an injective resolution of \(K\) on the étale site. By the lemma its pullback is still acyclic for \(\nu_{X,*}\), and by full faithfulness on sheaves the result is the resolution itself.

Choose a bounded-below injective resolution \(K\to I^\bullet\). Since \(\nu_X^*\) is exact, \(\nu_X^*I^\bullet\) represents \(\nu_X^*K\), and by Lemma \hyperref[bp100]{2.34} its terms are \(\nu_{X,*}\)-acyclic, so \(R\nu_{X,*}\nu_X^*K\) is computed by the complex \(\nu_{X,*}\nu_X^*I^\bullet\). By Theorem \hyperref[bp95]{2.32} the unit \(I^q\to\nu_{X,*}\nu_X^*I^q\) is an isomorphism for every \(q\), so this complex is \(I^\bullet\), which proves (2.7). Since \(\Gamma(X_{\mathrm{pro}\acute{e}t},-)=\Gamma(X_{\acute{e}t},-)\circ\nu_{X,*}\), the statement on global sections follows by composing the derived functors; see [Stacks, \href{https://stacks.math.columbia.edu/tag/099V}{099V}, \href{https://stacks.math.columbia.edu/tag/099W}{099W}].
\end{proof}

% text: prose-02-34
\phantomsection\label{bp105}%
Applied to a pro-étale affine, the theorem computes its cohomology from the étale schemes of which it is the limit.

% object: corollary-cohomology-pro-etale-affine (blueprint/chap-The-pro-etale-site/sec-Comparison-with-the-etale-site/corollary-cohomology-pro-etale-affine/_section.md)
\phantomsection\label{bp106}%
\begin{bpstatementit}{Corollary~2.36}{ (Cohomology of a pro-étale affine)}
Let \(K\in D^+(X_{\acute{e}t},\Lambda)\), and let \(U=\varprojlim_iU_i\) be a pro-étale affine with \(U_i\) affine and étale over \(X\). Then there is an isomorphism

\[
R\Gamma(U_{\mathrm{pro}\acute{e}t},(\nu_X^*K)|_U)
\simeq \varinjlim_i R\Gamma((U_i)_{\acute{e}t},K|_{U_i}).
\tag{2.8}
\]
\end{bpstatementit}

% proof: proof/corollary-cohomology-pro-etale-affine
\begin{proof}[{Proof}]\phantomsection\label{bp107}%
By the theorem applied to \(U\), the left side is \(R\Gamma(U_{\acute{e}t},K|_U)\), so the assertion is about étale cohomology alone. For a sheaf, continuity of étale cohomology under affine limits [Stacks, \href{https://stacks.math.columbia.edu/tag/03Q6}{03Q6}] identifies this with the colimit in each degree. For a bounded-below complex, only finitely many terms of the hypercohomology spectral sequence contribute to each total degree, and filtered colimits are exact.
\end{proof}

% text: prose-02-35
\phantomsection\label{bp108}%
To apply statements proved on the étale site, such as proper base change, on \(X_{\mathrm{pro}\acute{e}t}\), we also need the essential image of \(\nu_X^*\) on \(D^+\). It consists of the classical complexes. Write \(D^+_{\mathrm{class}}(X_{\mathrm{pro}\acute{e}t},\Lambda)\) for the full subcategory of \(D^+(X_{\mathrm{pro}\acute{e}t},\Lambda)\) of classical complexes.

% object: corollary-bounded-below-classical-complexes (blueprint/chap-The-pro-etale-site/sec-Comparison-with-the-etale-site/corollary-bounded-below-classical-complexes/_section.md)
\phantomsection\label{bp109}%
\begin{bpstatementit}{Corollary~2.37}{ (Bounded-below classical complexes)}
Classical module sheaves form a weak Serre subcategory, closed under kernels, cokernels and extensions. There is an equivalence

\[
D^+_{\mathrm{class}}(X_{\mathrm{pro}\acute{e}t},\Lambda)
\simeq D^+(X_{\acute{e}t},\Lambda),
\]

with inverse functors \(R\nu_{X,*}\) and \(\nu_X^*\).
\end{bpstatementit}

% proof: proof/corollary-bounded-below-classical-complexes
\begin{proof}[{Proof}]\phantomsection\label{bp110}%
We need the first assertion so that “classical complex” defines a triangulated subcategory, and for extensions we use the vanishing of \(R^1\nu_{X,*}\) from the theorem. The equivalence then follows by comparing cohomology sheaves.

For kernels and cokernels the first assertion follows from exactness and full faithfulness of \(\nu_X^*\). Indeed, a map between classical sheaves is \(\nu_X^*\) of a map of étale sheaves, and \(\nu_X^*\) carries its kernel and cokernel to the kernel and cokernel. Given an extension \(0\to\nu_X^*F\to G\to\nu_X^*F'\to0\), we apply \(\nu_{X,*}\). Since \(R^1\nu_{X,*}\nu_X^*F=0\) by Theorem \hyperref[bp103]{2.35}, the result \(0\to F\to\nu_{X,*}G\to F'\to0\) is exact, and the counit maps its pullback to the original extension. This map is the identity on the outer terms, so by the five lemma \(G\simeq\nu_X^*\nu_{X,*}G\) is classical.

The functor \(\nu_X^*\) lands in classical complexes and is fully faithful on \(D^+\) by Theorem \hyperref[bp103]{2.35}, so the equivalence amounts to the counit \(\nu_X^*R\nu_{X,*}K\to K\) being an isomorphism for \(K\) in \(D^+_{\mathrm{class}}\). By (2.7), every classical sheaf \(\nu_X^*F\) is \(\nu_{X,*}\)-acyclic, so the hypercohomology spectral sequence for \(R\nu_{X,*}K\) degenerates and \(H^q(R\nu_{X,*}K)=\nu_{X,*}H^q(K)\). On \(H^q\) the counit is therefore the counit \(\nu_X^*\nu_{X,*}H^q(K)\to H^q(K)\) for a classical sheaf, which is an isomorphism by Theorem \hyperref[bp95]{2.32}; see [Stacks, \href{https://stacks.math.columbia.edu/tag/09B1}{09B1}].
\end{proof}


% object: sec-Direct-images (blueprint/chap-The-pro-etale-site/sec-Direct-images/_section.md), content hidden
\subsection*{\texorpdfstring{2.9\quad Direct images}{2.9 Direct images}}
\phantomsection\label{bp111}%
\addcontentsline{toc}{subsection}{\texorpdfstring{2.9\quad Direct images}{2.9 Direct images}}
% text: prose-02-36
\phantomsection\label{bp112}%
Proper base change is proved on the étale site (Chapter \hyperref[bp285]{4}), while \(Rf_!\) is defined by direct images on the pro-étale site. To transfer proper base change, and to compute these direct images, we compare direct images on the two sites.

Let \(g:X\to Y\) be a morphism of schemes. Base change of objects of the sites defines a commutative square of topoi

\[
\begin{CD}
X_{\mathrm{pro}\acute{e}t}@>{\nu_X}>>X_{\acute{e}t}\\
@V{g_{\mathrm{pro}\acute{e}t}}VV @VV{g_{\acute{e}t}}V\\
Y_{\mathrm{pro}\acute{e}t}@>{\nu_Y}>>Y_{\acute{e}t}.
\end{CD}
\]

It induces an isomorphism \(g_{\mathrm{pro}\acute{e}t}^*\nu_Y^*\simeq\nu_X^*g_{\acute{e}t}^*\) of functors \(\operatorname{Mod}(Y_{\acute{e}t},\Lambda)\to\operatorname{Mod}(X_{\mathrm{pro}\acute{e}t},\Lambda)\), compatible with composition. Since all four functors are exact, it is also an isomorphism of functors \(D(Y_{\acute{e}t},\Lambda)\to D(X_{\mathrm{pro}\acute{e}t},\Lambda)\). For direct images the comparison needs a hypothesis on \(g\). The comparison map \(\nu_Y^*Rg_{\acute{e}t,*}\to Rg_{\mathrm{pro}\acute{e}t,*}\nu_X^*\) of functors \(D^+(X_{\acute{e}t},\Lambda)\to D^+(Y_{\mathrm{pro}\acute{e}t},\Lambda)\) is defined as the adjoint of the composite of \(g_{\mathrm{pro}\acute{e}t}^*\nu_Y^*Rg_{\acute{e}t,*}\simeq\nu_X^*g_{\acute{e}t}^*Rg_{\acute{e}t,*}\) with \(\nu_X^*\) applied to the counit of the adjunction \(g_{\acute{e}t}^*\dashv Rg_{\acute{e}t,*}\).

Evaluated on a pro-étale affine \(V=\varprojlim_iV_i\) over \(Y\), the comparison map is the map from \(\varinjlim_iH^q(X\times_YV_i,F)\) to \(H^q(X\times_YV,F)\). So the theorem below is the continuity of étale cohomology under affine limits, and like it, it needs \(g\) to be quasi-compact. Without this the map fails to be surjective already for \(q=0\). For example, let \(Y\) be affine and connected, \(X=\coprod_{n\in\mathbf N}Y\) and \(F\) the constant sheaf \(\mathbf Z/2\). Then \(g_*F\) is the product of countably many copies of \(\mathbf Z/2\) on \(Y\). On \(V=Y\times S\) with \(S\) an infinite profinite set, the pro-étale side has as sections all sequences of locally constant functions \(S\to\mathbf Z/2\), while the classical side only has the sequences that factor through a common finite quotient of \(S\).

% object: theorem-comparison-direct-images (blueprint/chap-The-pro-etale-site/sec-Direct-images/theorem-comparison-direct-images/_section.md)
\phantomsection\label{bp113}%
\begin{bpstatementit}{Theorem~2.38}{ (Comparison of direct images)}
If \(g\) is quasi-compact and quasi-separated, the canonical natural transformation

\[
\nu_Y^*Rg_{\acute{e}t,*}K\longrightarrow
Rg_{\mathrm{pro}\acute{e}t,*}\nu_X^*K
\tag{2.9}
\]

is an isomorphism for \(K\in D^+(X_{\acute{e}t},\Lambda)\).
\end{bpstatementit}

% proof: proof/theorem-comparison-direct-images
\begin{proof}[{Proof}]\phantomsection\label{bp114}%
We work over an affine open of \(Y\) and evaluate on the covering basis \(V=\varprojlim_iV_i\), with \(V_i\) affine and étale over \(Y\). By (2.6) and Theorem \hyperref[bp103]{2.35}, for a sheaf \(F\) the map in degree \(q\) is the continuity map

\[
\varinjlim_i H^q_{\acute{e}t}(X\times_YV_i,F)
\longrightarrow
H^q_{\acute{e}t}(X\times_YV,F).
\]

Since \(g\) is quasi-compact and quasi-separated, so are its base changes, and the map is an isomorphism by continuity of étale cohomology under affine limits [Stacks, \href{https://stacks.math.columbia.edu/tag/03Q6}{03Q6}]. Using bounded-below resolutions, or the convergent hypercohomology spectral sequences, we pass from sheaves to \(K\); see [Stacks, \href{https://stacks.math.columbia.edu/tag/09A6}{09A6}].
\end{proof}

% text: prose-02-37
\phantomsection\label{bp115}%
Since every bounded-below classical complex is of the form \(\nu_X^*K\) (Corollary \hyperref[bp109]{2.37}), the theorem shows that, for \(g\) quasi-compact and quasi-separated, \(Rg_{\mathrm{pro}\acute{e}t,*}\) carries bounded-below classical complexes to classical complexes.

In later arguments we transport maps built from units, counits and composition from the étale to the pro-étale site, and for this we record how the comparison map behaves under these operations. Write \(\beta_g:g_{\mathrm{pro}\acute{e}t}^*\nu_Y^*\simeq\nu_X^*g_{\acute{e}t}^*\) for the isomorphism above and \(c_g\) for the map (2.9).

% object: lemma-comparison-direct-images-compatibilities (blueprint/chap-The-pro-etale-site/sec-Direct-images/lemma-comparison-direct-images-compatibilities/_section.md)
\phantomsection\label{bp116}%
\begin{bpstatementit}{Lemma~2.39}{ (Compatibilities of the comparison of direct images)}
Let \(g:X\to Y\) be a morphism of schemes.

\begin{enumerate}
\item For \(L\in D^+(Y_{\acute{e}t},\Lambda)\), the composite \(\nu_Y^*L\to\nu_Y^*Rg_{\acute{e}t,*}g_{\acute{e}t}^*L\xrightarrow{c_g}
  Rg_{\mathrm{pro}\acute{e}t,*}\nu_X^*g_{\acute{e}t}^*L\) of \(\nu_Y^*\) of the unit with \(c_g\) equals the unit \(\nu_Y^*L\to Rg_{\mathrm{pro}\acute{e}t,*}g_{\mathrm{pro}\acute{e}t}^*\nu_Y^*L\) followed by \(Rg_{\mathrm{pro}\acute{e}t,*}(\beta_g)\).
\item For \(K\in D^+(X_{\acute{e}t},\Lambda)\), the composite \(g_{\mathrm{pro}\acute{e}t}^*\nu_Y^*Rg_{\acute{e}t,*}K\xrightarrow{g^*c_g}
  g_{\mathrm{pro}\acute{e}t}^*Rg_{\mathrm{pro}\acute{e}t,*}\nu_X^*K\to\nu_X^*K\) with the counit equals \(\nu_X^*\) of the counit \(g_{\acute{e}t}^*Rg_{\acute{e}t,*}K\to K\) precomposed with \(\beta_g\).
\item For a second morphism \(h:Y\to Z\), under the identifications \(R(h\circ g)_*\simeq Rh_*Rg_*\) on both sites, \(c_{h\circ g}\) is the composite \[
   \nu_Z^*Rh_{\acute{e}t,*}Rg_{\acute{e}t,*}K\xrightarrow{c_h}
   Rh_{\mathrm{pro}\acute{e}t,*}\nu_Y^*Rg_{\acute{e}t,*}K\xrightarrow{Rh_{\mathrm{pro}\acute{e}t,*}(c_g)}
   Rh_{\mathrm{pro}\acute{e}t,*}Rg_{\mathrm{pro}\acute{e}t,*}\nu_X^*K .
   \]
\end{enumerate}
\end{bpstatementit}

% proof: proof/lemma-comparison-direct-images-compatibilities
\begin{proof}[{Proof}]\phantomsection\label{bp117}%
Part (2) is the definition of \(c_g\) as an adjoint. For (1), by (2) and the naturality of \(\beta_g\), the adjoint of the composite in (1) is \(\beta_g\) followed by \(\nu_X^*\) of the composite \(g^*_{\acute{e}t}L\to g^*_{\acute{e}t}Rg_{\acute{e}t,*}g^*_{\acute{e}t}L\to g^*_{\acute{e}t}L\), which is the identity by the triangle identity. Hence the adjoint is \(\beta_g\), whose adjoint is the map on the right of (1). For (3), the isomorphisms \(\beta\) are compatible with composition, \(\beta_{h\circ g}\) being \(\beta_g\) composed with \(g_{\mathrm{pro}\acute{e}t}^*(\beta_h)\), and the adjoint of a composite of maps defined as in (2) is the composite of their adjoints.
\end{proof}

% text: prose-02-38
\phantomsection\label{bp118}%
In the comparison for constructible sheaves below we argue by induction on strata, using the functors of localization triangles. The next proposition identifies these functors on the two sites.

% object: proposition-comparison-immersions (blueprint/chap-The-pro-etale-site/sec-Direct-images/proposition-comparison-immersions/_section.md)
\phantomsection\label{bp119}%
\begin{bpstatementit}{Proposition~2.40}{ (Comparison for immersions)}
\begin{enumerate}
\item For an open immersion \(j:U\hookrightarrow X\), there is a canonical isomorphism \[
   \nu_X^*j_{\acute{e}t,!}\simeq j_{\mathrm{pro}\acute{e}t,!}\nu_U^*\quad\text{of functors }\operatorname{Mod}(U_{\acute{e}t},\Lambda)\to\operatorname{Mod}(X_{\mathrm{pro}\acute{e}t},\Lambda),
   \tag{2.10}
   \] and hence of functors \(D(U_{\acute{e}t},\Lambda)\to D(X_{\mathrm{pro}\acute{e}t},\Lambda)\).
\item For a closed immersion \(i:Z\hookrightarrow X\), the map (2.9) is an isomorphism \(\nu_X^*Ri_{\acute{e}t,*}\simeq Ri_{\mathrm{pro}\acute{e}t,*}\nu_Z^*\) of functors \(D^+(Z_{\acute{e}t},\Lambda)\to D^+(X_{\mathrm{pro}\acute{e}t},\Lambda)\), and \(i_{\mathrm{pro}\acute{e}t}^*\nu_X^*\simeq\nu_Z^*i_{\acute{e}t}^*\). For \(K\in D^+(X_{\acute{e}t},\Lambda)\), these isomorphisms and (2.10) carry \(\nu_X^*\) of the localization triangle \(j_{\acute{e}t,!}j_{\acute{e}t}^*K\to K\to Ri_{\acute{e}t,*}i_{\acute{e}t}^*K\to\) to a distinguished triangle \(j_{\mathrm{pro}\acute{e}t,!}j_{\mathrm{pro}\acute{e}t}^*\nu_X^*K\to\nu_X^*K\to Ri_{\mathrm{pro}\acute{e}t,*}i_{\mathrm{pro}\acute{e}t}^*\nu_X^*K\to\) whose first two maps are the counit and the unit of the adjunctions \(j_!\dashv j^*\) and \(i^*\dashv Ri_*\) on \(X_{\mathrm{pro}\acute{e}t}\).
\item Let \(f:X\to Y\) be a separated morphism of finite type with a factorization \(f=\bar f\circ j\), where \(j\) is an open immersion and \(\bar f\) is proper. Then there is a canonical isomorphism \(\nu_Y^*R\bar f_{\acute{e}t,*}j_{\acute{e}t,!}\simeq R\bar f_{\mathrm{pro}\acute{e}t,*}j_{\mathrm{pro}\acute{e}t,!}\nu_X^*\) of functors \(D^+(X_{\acute{e}t},\Lambda)\to D^+(Y_{\mathrm{pro}\acute{e}t},\Lambda)\).
\end{enumerate}
\end{bpstatementit}

% proof: proof/proposition-comparison-immersions
\begin{proof}[{Proof}]\phantomsection\label{bp120}%
For (1), both functors are left adjoint to restriction followed by \(\nu_{U,*}\), and this also identifies the adjunction maps [Stacks, \href{https://stacks.math.columbia.edu/tag/09AD}{09AD}]. As the functors are exact, the identity passes to derived categories. In (2), the first identity is Theorem \hyperref[bp113]{2.38}, and the second is the comparison of inverse images. Since \(\nu_X^*\) is exact, it carries the localization triangle to a distinguished triangle. Its first map is carried to the counit on \(X_{\mathrm{pro}\acute{e}t}\) because the isomorphism of (1) identifies the adjunction maps, and its second map to the unit by Lemma lemma-comparison-direct-images-compatibilities, applied to \(g=i\) and \(L=K\). Part (3) follows from (1) and Theorem \hyperref[bp113]{2.38} for the proper, hence quasi-compact and quasi-separated, map \(\bar f\).
\end{proof}


% object: sec-Constructible-sheaves (blueprint/chap-The-pro-etale-site/sec-Constructible-sheaves/_section.md), content hidden
\subsection*{\texorpdfstring{2.10\quad Constructible sheaves}{2.10 Constructible sheaves}}
\phantomsection\label{bp121}%
\addcontentsline{toc}{subsection}{\texorpdfstring{2.10\quad Constructible sheaves}{2.10 Constructible sheaves}}
% text: prose-02-39
\phantomsection\label{bp122}%
The six operations act on constructible complexes, and their finiteness theorems are proved on the étale site. To use these theorems we show that constructible pro-étale sheaves are classical, so that they correspond to constructible étale sheaves, and that the comparison preserves finite Tor amplitude.

A sheaf of modules on \(X_{\mathrm{pro}\acute{e}t}\) is locally constant if it is trivialized by a pro-étale covering, and it is of finite presentation if its constant models are finitely presented \(\Lambda\)-modules. Building on this, for Noetherian \(\Lambda\), a sheaf is constructible if, locally on \(X\), there is a finite partition into constructible locally closed strata such that it is locally constant of finite type on each stratum. We write \(\Cons(X,\Lambda)\) for the category of constructible sheaves and \(\Dbc(X,\Lambda)\) for the bounded complexes with constructible cohomology sheaves. For a discrete ring these notions agree with Definition \hyperref[bp199]{3.1}, which is stated for all coefficient rings in Section \hyperref[bp197]{3.2} below. These conditions refer only to pro-étale covers, so the main step is to descend torsors from pro-étale to étale covers.

A priori there could be more pro-étale local systems than étale ones, since the covers are larger. For finite groups, however, there are not. Indeed, a torsor under a finite group is described by a cocycle on one pro-étale affine cover. That cover is a limit of étale covers and the group is classical, so the cocycle is defined at a finite stage. For \(\mathbf Z_\ell\), by contrast, there are more local systems, and this is the reason for using the pro-étale site. For example, over \(\operatorname{Spec}\mathbf F_q\), a rank-one \(\mathbf Z_\ell\)-local system on which Frobenius acts by a unit \(u\in\mathbf Z_\ell^\times\) of infinite order is trivialized by a pro-étale cover but by no étale one. Our comparison therefore concerns only finite coefficients, and we study \(\ell\)-adic local systems through their reductions.

What we use later is the case of a finite ring \(A\). For such \(A\), Theorem \hyperref[bp129]{2.43} and Corollary \hyperref[bp132]{2.44} identify the category \(D^b_{\mathrm{ctf}}(X,A)\) on the pro-étale site with the classical one.

% object: lemma-torsors-under-classical-groups (blueprint/chap-The-pro-etale-site/sec-Constructible-sheaves/lemma-torsors-under-classical-groups/_section.md)
\phantomsection\label{bp123}%
\begin{bpstatementit}{Lemma~2.41}{ (Torsors under classical groups)}
Every pro-étale torsor under a classical sheaf of groups is classical, and hence has a trivializing étale covering.
\end{bpstatementit}

% proof: proof/lemma-torsors-under-classical-groups
\begin{proof}[{Proof}]\phantomsection\label{bp124}%
Write the group as \(\nu_X^*G_0\). We may work on an affine open of \(X\). Using Theorem \hyperref[bp42]{2.14} and finite affine refinements, we refine a trivializing covering to a single faithfully flat ind-étale affine \(U\to X\), and we write \(U=\varprojlim_iU_i\) with \(U_i\) affine étale over \(X\). A trivialization over \(U\) determines a transition element in \(\nu_X^*G_0(U\times_XU)\). By continuity for affine limits it descends to \(G_0(U_i\times_XU_i)\) for some \(i\), and the cocycle identity on the triple product holds after enlarging \(i\). Since \(U\to X\) is surjective, each \(U_i\to X\) is an étale covering. Gluing the trivial torsor along this cocycle, we get an étale torsor whose pullback is isomorphic to the given torsor. Hence the given torsor is classical by Theorem \hyperref[bp95]{2.32}, and it is trivialized by the étale covering \(U_i\to X\). See [Stacks, \href{https://stacks.math.columbia.edu/tag/099X}{099X}].
\end{proof}

% text: prose-02-40
\phantomsection\label{bp125}%
The hypothesis that the group is classical is essential. For instance, the group \(\mathbf Z_\ell^\times\) is not classical, and the rank-one \(\mathbf Z_\ell\)-local system on \(\operatorname{Spec}\mathbf F_q\) described above is a \(\mathbf Z_\ell^\times\)-torsor with no étale trivialization. On the other hand, for locally constant sheaves with finitely presented constant models over a discrete ring, the relevant group is the automorphism group of the model, which is discrete, and the lemma applies.

% object: proposition-locally-constant-sheaves-finite-presentation (blueprint/chap-The-pro-etale-site/sec-Constructible-sheaves/proposition-locally-constant-sheaves-finite-presentation/_section.md)
\phantomsection\label{bp126}%
\begin{bpstatementit}{Proposition~2.42}{ (Locally constant sheaves of finite presentation)}
Every locally constant pro-étale \(\Lambda_X\)-module sheaf of finite presentation is classical. In particular, it has a trivializing étale covering, and \(\nu_X^*\) identifies the étale local systems of finite presentation with these pro-étale ones.
\end{bpstatementit}

% proof: proof/proposition-locally-constant-sheaves-finite-presentation
\begin{proof}[{Proof}]\phantomsection\label{bp127}%
A locally constant sheaf with model \(M\) is the same as a torsor under the automorphism sheaf of the constant sheaf \(\underline M\), so the proposition follows from the lemma once we know that this automorphism sheaf is classical.

Decompose \(X\) into the open pieces on which the constant model has a fixed isomorphism type, and fix a finitely presented module \(M\). A presentation \(\Lambda^m\to\Lambda^n\to M\to0\) identifies the internal endomorphism sheaf of \(\underline M\) with the kernel of \(\underline M^{\,n}\to\underline M^{\,m}\). Since the constant sheaf functor is exact, this kernel is the constant sheaf on \(\operatorname{End}_\Lambda(M)\), and so the automorphism sheaf of \(\underline M\) is the constant discrete group \(\operatorname{Aut}_\Lambda(M)\). Hence \(\mathit{Isom}(\underline M,F)\) is a torsor under this classical group, and by Lemma \hyperref[bp123]{2.41} it has sections étale locally. Since \(F\) is obtained from this torsor by twisting \(\underline M\), it is the pullback of the corresponding twist of the constant étale sheaf, and so it is classical. This is the descent argument of [Stacks, \href{https://stacks.math.columbia.edu/tag/099Y}{099Y}].
\end{proof}

% text: prose-02-41
\phantomsection\label{bp128}%
To pass from a single stratum to arbitrary constructible sheaves, note that for a finite ring \(A\) the proposition applies to every finite local system of \(A\)-modules. A constructible sheaf is built from such local systems on the strata of a stratification, by extension by zero and extensions. Both operations preserve classical sheaves. For extension by zero this follows from the comparisons of Section \hyperref[bp111]{2.9}. For extensions it is Corollary \hyperref[bp109]{2.37}, which uses that \(R^1\nu_{X,*}\) vanishes on classical sheaves.

% object: theorem-comparison-constructible-sheaves (blueprint/chap-The-pro-etale-site/sec-Constructible-sheaves/theorem-comparison-constructible-sheaves/_section.md)
\phantomsection\label{bp129}%
\begin{bpstatementit}{Theorem~2.43}{ (Comparison for constructible sheaves)}
Let \(X\) be quasi-compact and quasi-separated. For a discrete Noetherian \(\Lambda\), every constructible pro-étale \(\Lambda_X\)-module sheaf is classical. The functor \(\nu_X^*\) gives an equivalence from constructible étale sheaves of \(\Lambda\)-modules to \(\Cons(X,\Lambda)\), and an equivalence \(D^b_c(X_{\acute{e}t},\Lambda)\to\Dbc(X,\Lambda)\) of bounded constructible derived categories.
\end{bpstatementit}

% proof: proof/theorem-comparison-constructible-sheaves
\begin{proof}[{Proof}]\phantomsection\label{bp130}%
The preceding proposition treats the case of a single stratum. For a finite stratification we argue locally, by induction on the number of strata. Once the sheaves are known to be classical, the equivalences follow from full faithfulness of \(\nu_X^*\) and its version for \(D^+\).

Given a point, we shrink to an affine open on which the stratum containing it is closed, with inclusion \(i\) and open complement \(j\). The exact sequence

\[
0\longrightarrow j_!j^*F\longrightarrow F
\longrightarrow i_*i^*F\longrightarrow0
\]

reduces the assertion to fewer strata, because on pro-étale sheaves of modules the functors \(j_!\) and \(i_*\) preserve classical sheaves, by (2.10) and Theorem \hyperref[bp113]{2.38}, and classical sheaves are closed under extensions (Corollary \hyperref[bp109]{2.37}). Since \(\nu_X^*\) is fully faithful, sheaves that are classical on the members of a Zariski covering are classical, so the local statement suffices. The derived statement follows from Corollary \hyperref[bp109]{2.37} and bounded truncation; see [Stacks, \href{https://stacks.math.columbia.edu/tag/09AK}{09AK}, \href{https://stacks.math.columbia.edu/tag/09AL}{09AL}].
\end{proof}

% text: prose-02-42
\phantomsection\label{bp131}%
The six operations act on complexes of finite Tor amplitude, so we also need to know that the equivalence just proved respects this condition. A priori the pro-étale condition is stronger, since it requires a Tor bound against all pro-étale modules, most of which are not classical.

% object: corollary-finite-tor-amplitude-under-comparison (blueprint/chap-The-pro-etale-site/sec-Constructible-sheaves/corollary-finite-tor-amplitude-under-comparison/_section.md)
\phantomsection\label{bp132}%
\begin{bpstatementit}{Corollary~2.44}{ (Finite Tor amplitude under comparison)}
Let \(X\) be quasi-compact and quasi-separated and let \(\Lambda\) be a discrete Noetherian ring. The equivalence \(\nu_X^*:D^b_c(X_{\acute{e}t},\Lambda)\to\Dbc(X,\Lambda)\) of Theorem \hyperref[bp129]{2.43} restricts to an equivalence

\[
D^b_{\mathrm{ctf}}(X_{\acute{e}t},\Lambda)\simeq D^b_{\mathrm{ctf}}(X,\Lambda).
\]
\end{bpstatementit}

% proof: proof/corollary-finite-tor-amplitude-under-comparison
\begin{proof}[{Proof}]\phantomsection\label{bp133}%
For one direction we represent an étale complex by a bounded flat complex and pull it back, and for the other we use the fact that \(\nu_X^*\) detects vanishing.

The functor \(\nu_X^*\) is exact (Proposition \hyperref[bp92]{2.31}), and since \(\nu_X^*\Lambda=\Lambda_X\) it is symmetric monoidal on sheaves of modules. The pullback of a flat sheaf of modules is flat [Stacks, \href{https://stacks.math.columbia.edu/tag/05VD}{05VD}]. An étale complex of Tor amplitude in \([a,b]\) is represented by a complex of flat sheaves in degrees \([a,b]\), so its pullback is a complex of flat pro-étale sheaves in the same degrees and has Tor amplitude in \([a,b]\) against every pro-étale module, not only against pulled-back ones. Conversely, let \(K\) be an étale complex. Its Tor amplitude can be read off from \(K\otimes^{\mathbf L}_\Lambda M\) for arbitrary étale sheaves of modules \(M\). The inverse image of this complex is \(\nu_X^*K\otimes^{\mathbf L}_{\Lambda_X}\nu_X^*M\), and since \(\nu_X^*\) is exact and fully faithful, it detects vanishing of cohomology. Hence a Tor bound on \(\nu_X^*K\) implies the same bound on \(K\).
\end{proof}

% text: prose-02-43
\phantomsection\label{bp134}%
In particular, the corollary applies to every finite commutative \(\ell\)-primary ring \(A\).


% object: sec-Unbounded-complexes (blueprint/chap-The-pro-etale-site/sec-Unbounded-complexes/_section.md), content hidden
\subsection*{\texorpdfstring{2.11\quad Unbounded complexes}{2.11 Unbounded complexes}}
\phantomsection\label{bp135}%
\addcontentsline{toc}{subsection}{\texorpdfstring{2.11\quad Unbounded complexes}{2.11 Unbounded complexes}}
% text: prose-02-44
\phantomsection\label{bp136}%
Proper base change needs the comparison with the étale site only for bounded-below complexes. Brown representability, however, takes place in the unbounded derived category, so we extend the comparison to unbounded complexes. The difficulty is that a morphism from a sheaf into an unbounded complex \(K\) involves in principle all the cohomology sheaves of \(K\) at once, and so does a cohomology sheaf of \(Rg_*K\). Both depend on only finitely many of them if Ext groups and cohomology of sheaves vanish above a fixed degree, uniformly in the sheaf. The smallest such degree is the cohomological dimension defined below.

Without such a bound, \(\Hom\) out of the constant sheaf need not commute with direct sums. For example, take \(\ell=2\) and \(S=\operatorname{Spec}\mathbf R\). Here étale cohomology is the cohomology of \(\operatorname{Gal}(\mathbf C/\mathbf R)=\mathbf Z/2\), and \(H^q(\mathbf Z/2,\mathbf F_2)=\mathbf F_2\) for every \(q\geq0\). Thus \(\operatorname{cd}_2(S)=\infty\), and \(\Hom(\mathbf F_2,-)\) does not commute with the direct sum \(\bigoplus_{n\geq0}\mathbf F_2[n]\) (see the opening of Chapter \hyperref[bp407]{5}).

Over a field of finite cohomological dimension, the étale cohomology of schemes of finite type is uniformly bounded. In that case every computation in a fixed degree can be carried out on a bounded-below truncation. We use the bounds in this way in Proposition \hyperref[bp146]{2.49} below and, later, in the proofs of compactness and of the preservation of direct sums by \(Rp_*\) (Lemmas \hyperref[bp426]{5.6} and \hyperref[bp463]{Lemma 5.17}).

% object: definition-cohomological-dimension (blueprint/chap-The-pro-etale-site/sec-Unbounded-complexes/definition-cohomological-dimension/_section.md)
\phantomsection\label{bp137}%
\begin{bpstatement}{Definition~2.45}{ (Cohomological dimension)}
For a qcqs scheme \(S\) and a qcqs morphism \(g:X\to Y\), the \emph{\(\ell\)-cohomological dimensions} are

\[
\begin{aligned}
\operatorname{cd}_\ell(S)&=\sup\{q: H^q(S_{\mathrm{pro}\acute{e}t},F)\neq0
\text{ for some classical $\ell$-power torsion sheaf $F$}\},\\
\operatorname{cd}_\ell(g)&=\sup\{q: R^qg_*F\neq0
\text{ for some classical $\ell$-power torsion sheaf $F$ on }X_{\mathrm{pro}\acute{e}t}\}.
\end{aligned}
\]

For a field \(k'\) we write \(\operatorname{cd}_\ell(k')=\operatorname{cd}_\ell(\operatorname{Spec}k')\).
\end{bpstatement}

% text: prose-02-45
\phantomsection\label{bp138}%
By Theorems \hyperref[bp103]{2.35} and \hyperref[bp113]{Theorem 2.38}, these are the usual \(\ell\)-cohomological dimensions computed on the étale sites. So the classical bounds apply. If \(k'\) is a field with \(\operatorname{cd}_\ell(k')<\infty\), then \(\operatorname{cd}_\ell(S)\) and \(\operatorname{cd}_\ell(g)\) are finite for all schemes \(S\) and morphisms \(g\) of finite type over \(k'\) [Stacks, \href{https://stacks.math.columbia.edu/tag/0F10}{0F10}]. For our field \(k\) the hypothesis holds by the following lemma.

% object: lemma-cohomological-dimension-base-field (blueprint/chap-The-pro-etale-site/sec-Unbounded-complexes/lemma-cohomological-dimension-base-field/_section.md)
\phantomsection\label{bp139}%
\begin{bpstatementit}{Lemma~2.46}{ (Cohomological dimension of the base field)}
If \(k'\) is an algebraically closed field, then \(\operatorname{cd}_\ell(k')=0\). If \(k'\) is a finite field, then \(\operatorname{cd}_\ell(k')\le1\).
\end{bpstatementit}

% proof: proof/lemma-cohomological-dimension-base-field
\begin{proof}[{Proof}]\phantomsection\label{bp140}%
Étale cohomology of \(\operatorname{Spec}k'\) is Galois cohomology [Stacks, \href{https://stacks.math.columbia.edu/tag/03QQ}{03QQ}]. For algebraically closed \(k'\) the Galois group is trivial. For finite \(k'\) it is \(\widehat{\mathbf Z}\), which has cohomological dimension one (see Serre, \emph{Galois cohomology}, II.§3).
\end{proof}

% text: prose-02-46
\phantomsection\label{bp141}%
Over \(\overline{\mathbf F}_q\), for example, a variety of dimension \(n\) has \(\ell\)-cohomological dimension at most \(2n\), and the curve \(\mathbf P^1\) attains this bound, since \(H^2(\mathbf P^1,\mathbf Z/\ell)\neq0\). Over \(\mathbf F_q\), however, one more degree is possible. Indeed, for \(\ell\neq p\) the Hochschild–Serre spectral sequence gives \(H^3(\mathbf P^1_{\mathbf F_q},\mu_\ell)\simeq
H^1(\operatorname{Spec}\mathbf F_q,H^2(\mathbf P^1_{\overline{\mathbf F}_q},\mu_\ell))
=H^1(\operatorname{Spec}\mathbf F_q,\mathbf Z/\ell)=\mathbf Z/\ell\).

To extend the comparison theorems to unbounded complexes with these bounds, we need two facts. First, the bounds hold uniformly over the affine objects on which we take sections (Lemma \hyperref[bp142]{2.47}). Second, on the pro-étale side a cohomology group of an unbounded complex in a fixed degree only sees a bounded-below truncation (Lemma \hyperref[bp144]{2.48}).

% object: lemma-uniform-bounds-pro-etale-affines (blueprint/chap-The-pro-etale-site/sec-Unbounded-complexes/lemma-uniform-bounds-pro-etale-affines/_section.md)
\phantomsection\label{bp142}%
\begin{bpstatementit}{Lemma~2.47}{ (Uniform bounds)}
Let \(X\) and \(Y\) be of finite type over \(k\), let \(A\) be a finite commutative \(\ell\)-primary ring, let \(g:X\to Y\) be of finite type, and put \(d=\dim(X)+\operatorname{cd}_\ell(k)\). There is an integer \(c\), depending only on \(X\), such that for every classical sheaf \(F\) of \(A\)-modules on \(X_{\mathrm{pro}\acute{e}t}\), every pro-étale affine \(U\) over \(X\) and every pro-étale affine \(V\) over \(Y\),

\[
H^q(U,F)=0\ \text{ for }q>d,
\qquad
H^q(X\times_YV,F)=0\ \text{ for }q>c.
\]

Moreover \(\operatorname{cd}_\ell(g)<\infty\).
\end{bpstatementit}

% proof: proof/lemma-uniform-bounds-pro-etale-affines
\begin{proof}[{Proof}]\phantomsection\label{bp143}%
We write both cohomologies as colimits of étale cohomology of schemes of finite type over \(k\) and bound these uniformly. Write \(U=\varprojlim_iU_i\) and \(V=\varprojlim_iV_i\) with \(U_i\) affine étale over \(X\) and \(V_i\) affine étale over \(Y\), and write \(F=\nu_X^*F_{\acute{e}t}\) (Theorem \hyperref[bp95]{2.32}). By Theorem \hyperref[bp103]{2.35}, applied to the schemes \(U\) and \(X\times_YV\), the cohomology of \(F\) on them is the étale cohomology of \(F_{\acute{e}t}\). As \(U=\varprojlim_iU_i\) and \(X\times_YV=\varprojlim_iX\times_YV_i\) are limits of qcqs schemes with affine transition maps, [Stacks, \href{https://stacks.math.columbia.edu/tag/03Q6}{03Q6}] identifies this with the colimit of the étale cohomology on the \(U_i\), respectively on the \(X\times_YV_i\). Each \(U_i\) is affine of finite type over \(k\) with \(\dim(U_i)\leq\dim(X)\), so by [Stacks, \href{https://stacks.math.columbia.edu/tag/0F0V}{0F0V}] we have \(\operatorname{cd}_\ell(U_i)\leq d\) for sheaves of \(A\)-modules. To bound cohomology on \(X\times_YV_i\) independently of \(i\), we choose a finite affine open cover of \(X\), with \(r\) members. Since \(X\) and \(Y\) are separated over \(k\), its intersections and their pullbacks to the affine \(V_i\) are affine; each pullback is étale over the corresponding intersection and has dimension at most \(\dim(X)\). The alternating Čech complex has length \(r-1\), so the Čech spectral sequence gives the bound \(c=d+r-1\). The last assertion is [Stacks, \href{https://stacks.math.columbia.edu/tag/0F10}{0F10}], which applies because \(g\) is of finite type over a field with \(\operatorname{cd}_\ell(k)<\infty\) (Lemma \hyperref[bp139]{2.46}).
\end{proof}

% object: lemma-truncation-pro-etale-sections (blueprint/chap-The-pro-etale-site/sec-Unbounded-complexes/lemma-truncation-pro-etale-sections/_section.md)
\phantomsection\label{bp144}%
\begin{bpstatementit}{Lemma~2.48}{ (Truncation on the pro-étale site)}
Let \(A\) be a finite commutative \(\ell\)-primary ring, let \(W\) be an object of \(X_{\mathrm{pro}\acute{e}t}\), and let \(e\) be an integer such that \(H^q(W,F)=0\) for \(q>e\) and every classical sheaf \(F\) of \(A\)-modules on \(X_{\mathrm{pro}\acute{e}t}\). Let \(L\in D(X_{\mathrm{pro}\acute{e}t},A)\) have classical cohomology sheaves, with \(H^q(L)=0\) for \(q>b\). Then \(H^jR\Gamma(W,L)=0\) for \(j>b+e+1\).
\end{bpstatementit}

% proof: proof/lemma-truncation-pro-etale-sections
\begin{proof}[{Proof}]\phantomsection\label{bp145}%
By left completeness (Proposition \hyperref[bp77]{2.26}), \(L\simeq R\varprojlim_n\tau_{\geq-n}L\), and derived sections preserve derived inverse limits, so \(R\Gamma(W,L)\simeq R\varprojlim_nR\Gamma(W,\tau_{\geq-n}L)\). Each \(\tau_{\geq-n}L\) is bounded with classical cohomology sheaves in degrees at most \(b\), so by the hypothesis and the hypercohomology spectral sequence, \(R\Gamma(W,\tau_{\geq-n}L)\) has no cohomology above \(b+e\). The Milnor sequence for this tower of complexes of abelian groups, proved as in Proposition \hyperref[bp71]{2.24}, expresses \(H^j\) of the derived limit through \(\varprojlim^1\) of \(H^{j-1}\) and \(\varprojlim\) of \(H^j\) of the terms. Both vanish for \(j>b+e+1\).
\end{proof}

% object: proposition-unbounded-comparison-finite-coefficients (blueprint/chap-The-pro-etale-site/sec-Unbounded-complexes/proposition-unbounded-comparison-finite-coefficients/_section.md)
\phantomsection\label{bp146}%
\begin{bpstatementit}{Proposition~2.49}{ (Unbounded comparison for finite coefficients)}
Let \(X\) and \(Y\) be of finite type over the base field \(k\), let \(A\) be a finite commutative \(\ell\)-primary ring, and let \(g:X\to Y\) be of finite type. Then \(\nu_X^*:D(X_{\acute{e}t},A)\to D(X_{\mathrm{pro}\acute{e}t},A)\) is fully faithful and preserves arbitrary direct sums. The map (2.9), now a natural transformation of functors \(D(X_{\acute{e}t},A)\to D(Y_{\mathrm{pro}\acute{e}t},A)\), is an isomorphism for every \(K\in D(X_{\acute{e}t},A)\).
\end{bpstatementit}

% proof: proof/proposition-unbounded-comparison-finite-coefficients
\begin{proof}[{Proof}]\phantomsection\label{bp147}%
By the two lemmas, both sides in a fixed degree depend only on a bounded-below truncation of \(K\), where Theorems \hyperref[bp103]{2.35} and \hyperref[bp113]{Theorem 2.38} apply. Let \(d\) and \(c\) be as in Lemma \hyperref[bp142]{2.47}, fix \(K\in
D(X_{\acute{e}t},A)\) and a degree \(m\), and consider the triangle \(\tau_{<N}K\to K\to\tau_{\geq N}K\) for an integer \(N\) chosen below.

Full faithfulness of \(\nu_X^*\) means that the unit \(K\to R\nu_{X,*}\nu_X^*K\) is an isomorphism, and it suffices to check this on the maps \(H^mR\Gamma(U,K)\to H^mR\Gamma(U,\nu_X^*K)\) for affine étale \(U\) over \(X\). We first reduce to bounded-below complexes. On the étale side, \(\Gamma(U,-)\) has cohomological dimension at most \(d\) on sheaves of \(A\)-modules, and by [Stacks, \href{https://stacks.math.columbia.edu/tag/07K7}{07K7}] the group \(H^mR\Gamma(U,K)\) depends only on \(\tau_{\geq N}K\) as soon as \(N\leq m-d\). On the pro-étale side, since \(\nu_X^*\) is exact it commutes with truncations, and Lemma \hyperref[bp144]{2.48}, applied with \(e=d\) to \(L=\nu_X^*\tau_{<N}K\) and \(b=N-1\), shows that \(H^m\) and \(H^{m+1}\) of \(R\Gamma(U,L)\) vanish as soon as \(N<m-d\). Hence for such \(N\) both sides are unchanged when \(K\) is replaced by \(\tau_{\geq N}K\), and for this bounded-below complex the map is an isomorphism by Theorem \hyperref[bp103]{2.35}.

For direct images, we evaluate (2.9) on a pro-étale affine \(V\) over \(Y\); these objects form a basis of \(Y_{\mathrm{pro}\acute{e}t}\). As before, we bound the contribution of \(\tau_{<N}K\) on both sides. On the étale side, the complex \(Rg_{\acute{e}t,*}\tau_{<N}K\) has cohomology sheaves in degrees at most \(N-1+\operatorname{cd}_\ell(g)\), so \(\nu_Y^*Rg_{\acute{e}t,*}\tau_{<N}K\) has classical cohomology in the same range, and Lemma \hyperref[bp144]{2.48} on \(Y\), with \(e=\dim(Y)+\operatorname{cd}_\ell(k)\) (Lemma \hyperref[bp142]{2.47} for \(Y\)), shows that its derived sections over \(V\) vanish above degree \(N+\operatorname{cd}_\ell(g)+\dim(Y)+\operatorname{cd}_\ell(k)\). On the other side, \(R\Gamma(V,Rg_{\mathrm{pro}\acute{e}t,*}\nu_X^*\tau_{<N}K)=R\Gamma(X\times_YV,\nu_X^*\tau_{<N}K)\), and the same lemma with \(W=X\times_YV\) and \(e=c\) gives vanishing above degree \(N+c\). Therefore, choosing \(N\) so small that \(H^m\) and \(H^{m+1}\) vanish on both sides, both sides of (2.9), evaluated on \(V\) in degree \(m\), are unchanged when \(K\) is replaced by \(\tau_{\geq N}K\), and for this bounded-below complex the map is an isomorphism by Theorem \hyperref[bp113]{2.38}.

The functor \(\nu_X^*\) preserves direct sums because it is an exact left adjoint.
\end{proof}

% text: prose-02-47
\phantomsection\label{bp148}%
Since by Theorem \hyperref[bp95]{2.32} the classical sheaves of \(A\)-modules are the essential image of \(\nu_X^*\) on sheaves, and \(\nu_X^*\) commutes with direct sums of sheaves, classical sheaves of \(A\)-modules are closed under arbitrary direct sums in \(\operatorname{Mod}(X_{\mathrm{pro}\acute{e}t},A)\).


% object: sec-Coefficient-rings (blueprint/chap-The-pro-etale-site/sec-Coefficient-rings/_section.md), content hidden
\subsection*{\texorpdfstring{2.12\quad Coefficient rings}{2.12 Coefficient rings}}
\phantomsection\label{bp149}%
\addcontentsline{toc}{subsection}{\texorpdfstring{2.12\quad Coefficient rings}{2.12 Coefficient rings}}
% text: prose-02-48
\phantomsection\label{bp150}%
We want pullback along morphisms and change of coefficients to be compatible without any bookkeeping. For this we define each coefficient ring once, on a point, and pull it back to every scheme along a single morphism of topoi. This gives \(A_X\), \(\mathcal O_{E,X}\), \(E_X\) and \(\QbarX\). Here \(A\) is a finite commutative ring killed by a power of \(\ell\), with the discrete topology, and \(A_X\) is the usual constant sheaf. For the point we take the site of profinite sets described below. By Proposition \hyperref[bp153]{2.50} it is the pro-étale site of a separably closed field. Unlike the étale site of the field, it keeps the topology of the coefficient rings.

% object: subsec-Condensed-rings (blueprint/chap-The-pro-etale-site/sec-Coefficient-rings/subsec-Condensed-rings/_section.md), content hidden
\subsubsection*{\texorpdfstring{2.12.1\quad Condensed rings}{2.12.1 Condensed rings}}
\phantomsection\label{bp151}%
\addcontentsline{toc}{subsubsection}{\texorpdfstring{2.12.1\quad Condensed rings}{2.12.1 Condensed rings}}
% text: prose-02-49
\phantomsection\label{bp152}%
The constant sheaf with value the abstract ring \(\mathbf Z_\ell\) forgets the topology of \(\mathbf Z_\ell\), and on \(\operatorname{Spec}\mathbf F_q\) it has \(H^1=0\) (Section \bpout{1.2}). We therefore need a notion of coefficient ring on a geometric point that retains the topology. The pro-étale site of a separably closed field is the site of profinite sets (Proposition \hyperref[bp153]{2.50} below). So a coefficient ring must at least be a sheaf of rings on profinite sets, that is, a condensed ring. The continuous maps from profinite sets into \(\mathbf Z_\ell\) form such a sheaf. With it, \(H^1\) of \(\mathbf F_q\) becomes \(\mathbf Z_\ell\) (Example \hyperref[bp178]{2.57}).

Let \(*_{\mathrm{pro}\acute{e}t}\) be the site of profinite sets, with finite jointly surjective families as covers. A \emph{condensed ring} is a sheaf of rings on \(*_{\mathrm{pro}\acute{e}t}\). A topological ring \(V\) defines a condensed ring \(\underline V\), whose sections over \(S\) are the continuous maps \(S\to V\). This is a sheaf because a surjection of profinite sets is a quotient map.

% object: proposition-pro-etale-site-separably-closed (blueprint/chap-The-pro-etale-site/sec-Coefficient-rings/subsec-Condensed-rings/proposition-pro-etale-site-separably-closed/_section.md)
\phantomsection\label{bp153}%
\begin{bpstatementit}{Proposition~2.50}{ (The pro-étale site of a separably closed field)}
Let \(\Omega\) be a separably closed field. The category of pro-étale affines over \(\operatorname{Spec}\Omega\) is equivalent to the category of profinite sets, with finite jointly surjective families as covers. Consequently its topos is, up to a set-theoretic cutoff, the topos of condensed sets. A discrete ring \(A\) defines the classical constant sheaf, while a profinite or locally profinite coefficient ring defines its sheaf of continuous functions on profinite sets.
\end{bpstatementit}

% proof: proof/proposition-pro-etale-site-separably-closed
\begin{proof}[{Proof}]\phantomsection\label{bp154}%
Since a finitely presented étale \(\Omega\)-algebra is a finite product of copies of \(\Omega\), filtered colimits of such algebras correspond, under \(\operatorname{Spec}\), to inverse limits of finite sets. Under this correspondence, faithfully flat covering maps become surjections of profinite sets, and by quasi-compactness every cover can be refined by a finite one, so the sites are equivalent. The topoi then agree because the pro-étale affines form a basis and sheaves are determined by their values on them. Finally, the last assertion is the definition of the condensed ring \(\underline R\) of a topological ring \(R\) given above, whose sections over a profinite set \(S\) are the continuous maps \(S\to R\); for discrete \(R\) these are the locally constant functions, which are the sections of the classical constant sheaf.
\end{proof}

% text: prose-02-50
\phantomsection\label{bp155}%
The étale topos of \(\Omega\) is the topos of sets, whereas by the proposition its pro-étale topos is the topos of condensed sets. The final object of \(*_{\mathrm{pro}\acute{e}t}\) is weakly contractible, so taking global sections is exact. However, a condensed sheaf is not determined by its global sections. For classical sheaves geometric stalks still suffice, whereas a general sheaf of modules on the pro-étale site has to be evaluated on all weakly contractible affines (Corollary \hyperref[bp59]{2.20}). For example, the completed coefficient sheaf \(\mathbf Z_{\ell,X}\) of Example \hyperref[bp97]{2.33} is of this general kind: its sections on a profinite set \(S\) are all continuous functions \(S\to\mathbf Z_\ell\).

Fix an algebraic closure \(\overline{\mathbf Q}_\ell\) of \(\mathbf Q_\ell\). Throughout, \(E\) denotes a finite extension of \(\mathbf Q_\ell\) inside \(\overline{\mathbf Q}_\ell\), with ring of integers \(\mathcal O_E\) and uniformizer \(\pi_E\). For \(\mathcal O_E\) and \(E\) the \(\ell\)-adic topologies define condensed rings \(\underline{\mathcal O_E}\) and \(\underline E\). For \(\overline{\mathbf Q}_\ell\), however, the topology causes a problem. To see why, recall that classically a \(\overline{\mathbf Q}_\ell\)-sheaf is an \(E\)-sheaf for some finite \(E\) with scalars extended, and all arguments with lattices and reduction modulo \(\pi_E\) take place over that \(E\). They cannot be made over \(\overline{\mathbf Q}_\ell\) itself, because \(\overline{\mathbf Q}_\ell\) is not complete and its valuation ring is not a discrete valuation ring. But the sheaf of all continuous maps to \(\overline{\mathbf Q}_\ell\) has sections that are defined over no finite \(E\) (see below), so it would lose this reduction. To keep it available, we take instead the filtered colimit over the finite coefficient fields. Using the valuation topology on \(\overline{\mathbf Q}_\ell\), we put

\[
\underline{\overline{\mathbf Q}_\ell}(S)=
\{a:S\to\overline{\mathbf Q}_\ell\text{ continuous}:a(S)\subset E\text{ for some finite }E/\mathbf Q_\ell\}
=\varinjlim_E\underline E(S).
\]

This is already a sheaf, because a cover of a profinite set can be refined by a finite one, and finitely many finite extensions of \(\mathbf Q_\ell\) lie in a common one. Since filtered colimits of sheaves on \(*_{\mathrm{pro}\acute{e}t}\) are computed sectionwise, \(\underline{\overline{\mathbf Q}_\ell}=\varinjlim_E\underline E\) as condensed rings.

This condensed ring is smaller than \(\underline V\) for \(V=\overline{\mathbf Q}_\ell\) with the valuation topology. Indeed, algebraic numbers of unbounded degree can tend to \(0\). Such a sequence is a continuous map from a convergent sequence to \(\overline{\mathbf Q}_\ell\) with image in no finite extension. With our choice, Lemma \hyperref[bp265]{3.23} shows that every constructible \(\overline{\mathbf Q}_\ell\)-complex and every morphism between such complexes descends to a finite \(E\).


% object: subsec-Coefficient-sheaves (blueprint/chap-The-pro-etale-site/sec-Coefficient-rings/subsec-Coefficient-sheaves/_section.md), content hidden
\subsubsection*{\texorpdfstring{2.12.2\quad Coefficient sheaves}{2.12.2 Coefficient sheaves}}
\phantomsection\label{bp156}%
\addcontentsline{toc}{subsubsection}{\texorpdfstring{2.12.2\quad Coefficient sheaves}{2.12.2 Coefficient sheaves}}
% text: prose-02-51
\phantomsection\label{bp157}%
A profinite set \(S\) defines a pro-étale scheme \(\operatorname{Spec}LC(S,\mathbf Z)\) over \(\mathbf Z\), where \(LC(S,\mathbf Z)\) is the ring of locally constant functions. The functor \(S\mapsto\operatorname{Spec}LC(S,\mathbf Z)\) is continuous and preserves finite limits, as one checks on finite quotients of \(S\), so it induces a morphism of topoi \(\Sh((\operatorname{Spec}\mathbf Z)_{\mathrm{pro}\acute{e}t})\to
\Sh(*_{\mathrm{pro}\acute{e}t})\) [Stacks, \href{https://stacks.math.columbia.edu/tag/00X6}{00X6}]. Composing with the structure morphism of \(X\), we obtain

\[
p_X:\Sh(X_{\mathrm{pro}\acute{e}t})\longrightarrow\Sh(*_{\mathrm{pro}\acute{e}t}),
\]

and for a condensed ring \(\Lambda\) we write \(\Lambda_X=p_X^*\Lambda\). For a morphism \(f:X\to Y\) we have \(p_Y\circ f=p_X\), and hence a canonical isomorphism \(f^{-1}\Lambda_Y\simeq\Lambda_X\), compatible with identities and composition. Likewise, a morphism \(\Lambda\to\Lambda'\) of condensed rings pulls back to \(\Lambda_X\to\Lambda'_X\), which defines extension of scalars \(\Lambda'_X\otimes_{\Lambda_X}-\) from \(\Lambda_X\)-modules to \(\Lambda'_X\)-modules on \(X_{\mathrm{pro}\acute{e}t}\). Thus both pullback along morphisms and change of coefficients are obtained from \(p_X^*\).

The next lemma identifies these rings with the completed coefficient rings of classical \(\ell\)-adic cohomology.

% object: lemma-comparison-completed-coefficients (blueprint/chap-The-pro-etale-site/sec-Coefficient-rings/subsec-Coefficient-sheaves/lemma-comparison-completed-coefficients/_section.md)
\phantomsection\label{bp158}%
\begin{bpstatementit}{Lemma~2.51}{ (Comparison with completed coefficients)}
Write \(\mathcal O_{E,X}\), \(E_X\) and \(\QbarX\) for the pullbacks to \(X_{\mathrm{pro}\acute{e}t}\) of \(\underline{\mathcal O_E}\), \(\underline E\) and \(\underline{\overline{\mathbf Q}_\ell}\). There are canonical isomorphisms of sheaves of rings

\[
\mathcal O_{E,X}\simeq\varprojlim_n(\mathcal O_E/\pi_E^n)_X,
\qquad E_X\simeq\mathcal O_{E,X}[1/\ell],
\qquad \QbarX\simeq\varinjlim_E E_X,
\]

where \((\mathcal O_E/\pi_E^n)_X\) is the constant sheaf. They are compatible with inverse image along every morphism of schemes.
\end{bpstatementit}

% proof: proof/lemma-comparison-completed-coefficients
\begin{proof}[{Proof}]\phantomsection\label{bp159}%
The functor \(p_X^*\) commutes with colimits, so the second and third isomorphisms reduce to identities on the point. The first is an inverse limit, and since \(p_X^*\) need not commute with inverse limits, we use instead that \(\underline{\mathcal O_E}\) is representable.

On the point, \(\underline{\mathcal O_E}\) is represented by the profinite set \(\mathcal O_E=\varprojlim_n\mathcal O_E/\pi_E^n\), so its pullback to \(X\) is represented by the pro-étale \(X\)-scheme \(\varprojlim_n\coprod_{\mathcal O_E/\pi_E^n}X\). By the Yoneda lemma the sheaf this scheme represents is \(\varprojlim_n(\mathcal O_E/\pi_E^n)_X\), and the ring structures match.

A continuous map from a profinite set to \(E\) has compact image, and hence lands in \(\ell^{-m}\mathcal O_E\) for some \(m\). Thus \(\underline E=\underline{\mathcal O_E}[1/\ell]\) on the point. Since inverting \(\ell\) is a filtered colimit and \(p_X^*\) commutes with colimits, we obtain the second isomorphism. In the same way the third follows from \(\underline{\overline{\mathbf Q}_\ell}=\varinjlim_E\underline E\). All three isomorphisms are obtained by applying \(p_X^*\), and \(p_Y\circ f=p_X\) for every morphism \(f:X\to Y\), so they are compatible with pullback.
\end{proof}

% text: prose-02-52
\phantomsection\label{bp160}%
The first isomorphism describes \(\mathcal O_{E,X}\) as a limit of its finite quotients. To compute derived reductions modulo \(\pi_E^n\) we also need these quotients as cokernels of multiplication by \(\pi_E^n\), which the next corollary provides.

% object: corollary-finite-quotients-their-transition-maps (blueprint/chap-The-pro-etale-site/sec-Coefficient-rings/subsec-Coefficient-sheaves/corollary-finite-quotients-their-transition-maps/_section.md)
\phantomsection\label{bp161}%
\begin{bpstatementit}{Corollary~2.52}{ (Finite quotients and their transition maps)}
For every \(n\geq1\) there is a canonical exact sequence on \(X_{\mathrm{pro}\acute{e}t}\)

\[
0\longrightarrow\mathcal O_{E,X}
\xrightarrow{\pi_E^n}\mathcal O_{E,X}
\longrightarrow(\mathcal O_E/\pi_E^n)_X\longrightarrow0.
\]

Thus the quotient is represented by the two-term complex \([\mathcal O_{E,X}\xrightarrow{\pi_E^n}\mathcal O_{E,X}]\) in degrees \(-1,0\). Under this representation, reduction from level \(n+1\) to level \(n\) is the chain map \((\pi_E,\mathrm{id})\). These resolutions and maps commute with inverse image along every morphism of schemes.
\end{bpstatementit}

% proof: proof/corollary-finite-quotients-their-transition-maps
\begin{proof}[{Proof}]\phantomsection\label{bp162}%
It suffices to prove exactness on the point, since \(p_X^*\) is exact and inverse image along \(f:X\to Y\) satisfies \(p_Y\circ f=p_X\). Now, on a profinite set \(S\), multiplication by \(\pi_E^n\) is injective on continuous functions \(S\to\mathcal O_E\), and its image is the kernel of reduction. Moreover, reduction is surjective because a continuous function \(S\to\mathcal O_E/\pi_E^n\) takes finitely many values, and choosing a representative of each gives a continuous lift. Finally, the chain map \((\pi_E,\mathrm{id})\) commutes with the differentials because \(\pi_E^{n}\pi_E=\pi_E^{n+1}\), and its degree-zero component is the identity, so it induces the quotient projection.
\end{proof}

% text: prose-02-53
\phantomsection\label{bp163}%
Using the two-term resolution we see that, for every complex \(K\) of \(\mathcal O_{E,X}\)-modules, the derived reduction \(K\otimes^{\mathbf L}(\mathcal O_E/\pi_E^n)_X\) is the cone of \(\pi_E^n:K\to K\) (Lemma \hyperref[bp190]{2.62}).

We call the rings \(A,\mathbf Z_\ell,\mathcal O_E,\mathbf Q_\ell,E,\overline{\mathbf Q}_\ell\) the \emph{admissible} coefficient rings. Here \(E=\mathbf Q_\ell\) is allowed, so that \(\mathbf Z_\ell\) and \(\mathbf Q_\ell\) are the case \(E=\mathbf Q_\ell\) of \(\mathcal O_E\) and \(E\). For \(E=\mathbf Q_\ell\) we write \(\mathbf Z_{\ell,X}=\mathcal O_{\mathbf Q_\ell,X}\) and \(\mathbf Q_{\ell,X}=(\mathbf Q_\ell)_X\). By the comparison lemma above,

\[
\mathbf Z_{\ell,X}\simeq\varprojlim_n(\mathbf Z/\ell^n)_X,
\qquad
\mathbf Q_{\ell,X}\simeq\mathbf Z_{\ell,X}[1/\ell].
\]

For a condensed ring \(\Lambda\) we write \(\operatorname{Mod}(X_{\mathrm{pro}\acute{e}t},\Lambda)\) for the category of \(\Lambda_X\)-modules. Its derived category is \(D(X_{\mathrm{pro}\acute{e}t},\Lambda)\).

Duality and the trace map need, besides the constant sheaf, its twists by roots of unity. The reason is already visible on \(\mathbf P^1\) over an algebraically closed field \(k\), for \(n\) invertible in \(k\). The Kummer sequence \(0\to\mu_n\to\mathbf G_m\xrightarrow{n}\mathbf G_m\to0\) gives an isomorphism \(H^2(\mathbf P^1,\mu_n)\simeq\operatorname{Pic}(\mathbf P^1)/n=\mathbf Z/n\), given by the degree of a line bundle. This isomorphism involves no choice, and it sends the class of a point to \(1\). An identification of \(H^2(\mathbf P^1,\mathbf Z/n)\) with \(\mathbf Z/n\), on the other hand, requires an isomorphism \(\mu_n\simeq\mathbf Z/n\), that is, a choice of a primitive \(n\)th root of unity. Frobenius does not respect such a choice.

To see this, let \(k\) be an algebraic closure of \(\mathbf F_q\). The arithmetic Frobenius acts on \(\mu_{\ell^n}(k)\) by \(\zeta\mapsto\zeta^q\). The geometric Frobenius is the inverse of the arithmetic Frobenius in \(\operatorname{Gal}(k/\mathbf F_q)\). Hence it acts on \(\Lambda(1)\) by \(q^{-1}\), and on \(H^2(\mathbf P^1,\Lambda)\simeq\Lambda(-1)\) by \(q\). This agrees with \(\#\mathbf P^1(\mathbf F_q)=1+q\).

Over \(\mathbf C\) the twist is \(2\pi i\mathbf Z\). The exponential sequence identifies \(H^2(\mathbf P^1(\mathbf C),2\pi i\mathbf Z)\) with \(\mathbf Z\). An identification of \(H^2(\mathbf P^1(\mathbf C),\mathbf Z)\) with \(\mathbf Z\), by contrast, is a choice of orientation, and the complex orientation depends on the choice of \(i\).

With twists, the trace map of \(\mathbf P^1\) (Proposition \hyperref[bp534]{6.7}) and smooth purity \(f^!\simeq f^*(d)[2d]\) (Theorem \bpout{9.3}) are canonical and commute with Frobenius (Proposition \bpout{10.20}). This is where the factor \(q^d\) in Poincaré duality (c) comes from.

% object: definition-tate-twists (blueprint/chap-The-pro-etale-site/sec-Coefficient-rings/subsec-Coefficient-sheaves/definition-tate-twists/_section.md)
\phantomsection\label{bp164}%
\begin{bpstatement}{Definition~2.53}{ (Tate twists)}
Let \(X\) be a scheme on which \(\ell\) is invertible, and let \(\mu_{\ell^n,X}\) be the sheaf of \(\ell^n\)th roots of unity on \(X_{\mathrm{pro}\acute{e}t}\). Then it is locally isomorphic to \((\mathbf Z/\ell^n\mathbf Z)_X\). For a finite ring \(A\) with \(\ell^rA=0\) put \(A_X(1)=\mu_{\ell^r,X}\otimes_{\mathbf Z/\ell^r\mathbf Z}A_X\), which by the Kummer maps is independent of \(r\). Put \(\mathcal O_{E,X}(1)=\varprojlim_n(\mathcal O_E/\ell^n)_X(1)\), and obtain \(E_X(1)\) and \(\QbarX(1)\) by inverting \(\ell\) and passing to the colimit over \(E\); equivalently, \(\QbarX(1)=\QbarX\otimes_{\mathbf Z_{\ell,X}}\mathbf Z_{\ell,X}(1)\) with \(\mathbf Z_{\ell,X}(1)=\varprojlim_n\mu_{\ell^n,X}\). Each \(\Lambda_X(1)\) is locally free of rank one. For \(d\ge0\), \(\Lambda_X(d)\) is the \(d\)-fold tensor power, \(\Lambda_X(-d)\) its dual, and \(K(d)=K\otimes^{\mathbf L}_{\Lambda_X}\Lambda_X(d)\).
\end{bpstatement}


% object: subsec-Inverse-image (blueprint/chap-The-pro-etale-site/sec-Coefficient-rings/subsec-Inverse-image/_section.md), content hidden
\subsubsection*{\texorpdfstring{2.12.3\quad Inverse image}{2.12.3 Inverse image}}
\phantomsection\label{bp165}%
\addcontentsline{toc}{subsubsection}{\texorpdfstring{2.12.3\quad Inverse image}{2.12.3 Inverse image}}
% text: prose-02-54
\phantomsection\label{bp166}%
Every base change argument below compares a module over \(\Lambda_Y\) with its pullback to \(X\), and for this we want the pullback of a \(\Lambda_Y\)-module to be its inverse image as a sheaf, with no tensor product. For coherent sheaves this fails, since \(f^*M=\mathcal O_X\otimes_{f^{-1}\mathcal O_Y}f^{-1}M\) and \(f^{-1}\mathcal O_Y\ne\mathcal O_X\) in general. For our coefficient sheaves the two notions agree, because all coefficient sheaves are pulled back from the point. The same holds for their Tate twists, where \(\ell\) is invertible.

The second part of the next lemma says that extension of scalars to a finite extension \(E'/E\) is a finite direct sum. We use this, for instance, in the proof of Proposition \hyperref[bp405]{4.40}.

% object: lemma-inverse-image-coefficient-sheaves (blueprint/chap-The-pro-etale-site/sec-Coefficient-rings/subsec-Inverse-image/lemma-inverse-image-coefficient-sheaves/_section.md)
\phantomsection\label{bp167}%
\begin{bpstatementit}{Lemma~2.54}{ (Inverse image of the coefficient sheaves)}
Let \(f:X\to Y\) be a morphism of schemes.

\begin{enumerate}
\item For every admissible \(\Lambda\), the canonical map \[
   f^{-1}\Lambda_Y\longrightarrow\Lambda_X
   \] is an isomorphism of sheaves of rings. If \(\ell\) is invertible on \(Y\), the same holds for its Tate twists, including the twists of a finite ring \(A\).
\item If \(E\subset E'\) is a finite extension, then \(\mathcal O_{E',X}\) is a finite free \(\mathcal O_{E,X}\)-module and \(E'_X\) is a finite free \(E_X\)-module.
\end{enumerate}
\end{bpstatementit}

% proof: proof/lemma-inverse-image-coefficient-sheaves
\begin{proof}[{Proof}]\phantomsection\label{bp168}%
(1) For every condensed ring \(\Lambda\), the identity \(p_Y\circ f=p_X\) shows that

\[
f^{-1}\Lambda_Y=f^{-1}p_Y^*\Lambda\simeq p_X^*\Lambda=\Lambda_X,
\]

and this proves the assertion for the coefficient rings themselves.

For the twists, \(\ell\) is invertible on \(Y\), hence on \(X\), as Definition \hyperref[bp164]{2.53} requires. The torsion twists \(\mu_{\ell^n}\) are then represented by finite étale group schemes, and these commute with base change. Their inverse limit is represented by the inverse limit of these schemes, so that its inverse image is represented by the base change [Stacks, \href{https://stacks.math.columbia.edu/tag/09A5}{09A5}]. This proves the assertion for \(\mathcal O_E(1)\). For a finite \(A\) killed by \(\ell^n\), its twist is obtained from the twist at finite level by extending scalars from \(\mathbf Z/\ell^n\mathbf Z\), so inverse image commutes with it as well. The other twists are obtained by inverting \(\ell\), extending scalars, and taking tensor powers and duals.

(2) An \(\mathcal O_E\)-basis of \(\mathcal O_{E'}\) is an isomorphism of topological \(\mathcal O_E\)-modules \(\mathcal O_{E'}\simeq\mathcal O_E^{[E':E]}\), hence of condensed modules \(\underline{\mathcal O_{E'}}\simeq\underline{\mathcal O_E}^{[E':E]}\). Pull it back along \(p_X\), and invert \(\ell\) for the fields.
\end{proof}

% text: prose-02-55
\phantomsection\label{bp169}%
By (1), the pullback of modules over these sheaves, defined in Section \hyperref[bp195]{3.1} below, is the exact functor \(f^{-1}\).



% object: sec-Continuous-cohomology (blueprint/chap-The-pro-etale-site/sec-Continuous-cohomology/_section.md), content hidden
\subsection*{\texorpdfstring{2.13\quad Continuous cohomology}{2.13 Continuous cohomology}}
\phantomsection\label{bp170}%
\addcontentsline{toc}{subsection}{\texorpdfstring{2.13\quad Continuous cohomology}{2.13 Continuous cohomology}}
% text: prose-02-56
\phantomsection\label{bp171}%
The comparison theorems above concern discrete coefficients, whereas the integral and rational \(\ell\)-adic coefficient rings of Section \hyperref[bp149]{2.12} are defined directly on the pro-étale site, because the discrete constant sheaves with the same values have the wrong cohomology (Example \hyperref[bp178]{2.57}). The main result of this section, Proposition \hyperref[bp172]{2.55}, says that pro-étale cohomology with coefficients \(\mathbf Z_\ell\) is the classical continuous cohomology, that is, the derived limit of cohomology with coefficients \(\mathbf Z/\ell^n\).

The sheaf \(\mathbf Z_{\ell,X}\) is the inverse limit of the classical sheaves \((\mathbf Z/\ell^n)_X\), with surjective transition maps (Lemma \hyperref[bp158]{2.51}). For such a limit, the cohomology is the derived limit of the cohomology of the terms. When the terms are classical, this is continuous étale cohomology.

% object: proposition-continuous-cohomology-pro-etale-cohomology (blueprint/chap-The-pro-etale-site/sec-Continuous-cohomology/proposition-continuous-cohomology-pro-etale-cohomology/_section.md)
\phantomsection\label{bp172}%
\begin{bpstatementit}{Proposition~2.55}{ (Continuous cohomology as pro-étale cohomology)}
Let \(X\) be a scheme and \((F_n)\) an inverse system of abelian sheaves on \(X_{\mathrm{pro}\acute{e}t}\) with epimorphic transition maps. Put \(F=\varprojlim_nF_n\) in the pro-étale topos. Then, in the derived category of abelian groups,

\[
R\Gamma(X_{\mathrm{pro}\acute{e}t},F)
\simeq R\varprojlim_n R\Gamma(X_{\mathrm{pro}\acute{e}t},F_n).
\]
\end{bpstatementit}

% proof: proof/proposition-continuous-cohomology-pro-etale-cohomology
\begin{proof}[{Proof}]\phantomsection\label{bp173}%
Since the transition maps are epimorphisms, \(F\simeq R\varprojlim_nF_n\) by the last assertion of Proposition \hyperref[bp71]{2.24}. Moreover, \(R\Gamma(X_{\mathrm{pro}\acute{e}t},-)\) commutes with derived inverse limits since it is a right adjoint.
\end{proof}

% text: prose-02-57
\phantomsection\label{bp174}%
When the \(F_n\) come from the étale site, as the sheaves \((\mathbf Z/\ell^n)_X\) do, each term on the right is étale cohomology by the comparison theorem, and the proposition identifies pro-étale cohomology of the limit with continuous étale cohomology \(R\varprojlim_nR\Gamma(X_{\acute{e}t},F_n)\).

% object: corollary-continuous-etale-cohomology (blueprint/chap-The-pro-etale-site/sec-Continuous-cohomology/corollary-continuous-etale-cohomology/_section.md)
\phantomsection\label{bp175}%
\begin{bpstatementit}{Corollary~2.56}{ (Continuous étale cohomology)}
In Proposition \hyperref[bp172]{2.55}, suppose that \(F_n=\nu_X^*F_{n,\acute{e}t}\) for an inverse system \((F_{n,\acute{e}t})\) of abelian sheaves on \(X_{\acute{e}t}\) with epimorphic transition maps. Then \(R\Gamma(X_{\mathrm{pro}\acute{e}t},F)\) is the continuous étale cohomology of \((F_{n,\acute{e}t})\): the right derived functor of \((G_n)\mapsto\Gamma(X_{\acute{e}t},\varprojlim_nG_n)\) on the abelian category of inverse systems, evaluated at \((F_{n,\acute{e}t})\).
\end{bpstatementit}

% proof: proof/corollary-continuous-etale-cohomology
\begin{proof}[{Proof}]\phantomsection\label{bp176}%
By Theorem \hyperref[bp103]{2.35}, each \(R\Gamma(X_{\mathrm{pro}\acute{e}t},F_n)\) is \(R\Gamma(X_{\acute{e}t},F_{n,\acute{e}t})\). The derived functor of \((G_n)\mapsto\Gamma(X_{\acute{e}t},\varprojlim_nG_n)\) is the composite of the derived functors of sections and of inverse limit, so its value at \((F_{n,\acute{e}t})\) is \(R\varprojlim_nR\Gamma(X_{\acute{e}t},F_{n,\acute{e}t})\) [\hyperref[bp705]{BS15}, Proposition 5.6.2].
\end{proof}

% text: prose-02-58
\phantomsection\label{bp177}%
Applying the Milnor sequence (2.2) to the tower of complexes \(R\Gamma(X_{\mathrm{pro}\acute{e}t},F_n)\) and using the proposition, we obtain for every \(q\) a natural exact sequence

\[
0\longrightarrow R^1\varprojlim_nH^{q-1}(X_{\mathrm{pro}\acute{e}t},F_n)
\longrightarrow H^q(X_{\mathrm{pro}\acute{e}t},F)
\longrightarrow\varprojlim_nH^q(X_{\mathrm{pro}\acute{e}t},F_n)
\longrightarrow0.
\tag{2.11}
\]

The first term can be nonzero, because the transition maps between the cohomology groups need not be surjective even if those between the sheaves \(F_n\) are. It vanishes if the groups \(H^{q-1}(X_{\mathrm{pro}\acute{e}t},F_n)\) satisfy the Mittag–Leffler condition. This holds for instance if they are finite, since then descending chains of images stabilize at each level.

% object: example-first-cohomology-finite-field (blueprint/chap-The-pro-etale-site/sec-Continuous-cohomology/example-first-cohomology-finite-field/_section.md)
\phantomsection\label{bp178}%
\begin{bpstatement}{Example~2.57}{ (The first cohomology of a finite field)}
For \(X=\operatorname{Spec}\mathbf F_q\) and \(\ell\nmid q\),

\[
H^1(X_{\mathrm{pro}\acute{e}t},\mathbf Z_{\ell,X})\simeq\mathbf Z_\ell,
\qquad H^1(X_{\mathrm{pro}\acute{e}t},\mathbf Q_{\ell,X})\simeq\mathbf Q_\ell.
\]

At each finite level the first cohomology is \(\operatorname{Hom}_{\mathrm{cont}}(\widehat{\mathbf Z},
\mathbf Z/\ell^n)=\mathbf Z/\ell^n\), and the groups in degree zero are finite, so that their \(\varprojlim^1\) vanishes. Hence the sequence (2.11) for \(F_n=(\mathbf Z/\ell^n)_X\) proves the integral statement. Since \(X\) is qcqs, its pro-étale cohomology commutes with filtered colimits [Stacks, \href{https://stacks.math.columbia.edu/tag/09A3}{09A3}], in particular with inverting \(\ell\), and the rational statement follows. For the constant étale sheaf with value the discrete group \(\mathbf Q_\ell\), on the other hand, the first cohomology is \(\operatorname{Hom}_{\mathrm{cont}}(\widehat{\mathbf Z},\mathbf Q_\ell)=0\), because a continuous homomorphism from a profinite group to a discrete torsion-free group has finite, hence zero, image.
\end{bpstatement}


% object: sec-Extension-by-zero-and-completion (blueprint/chap-The-pro-etale-site/sec-Extension-by-zero-and-completion/_section.md), content hidden
\subsection*{\texorpdfstring{2.14\quad Extension by zero and completion}{2.14 Extension by zero and completion}}
\phantomsection\label{bp179}%
\addcontentsline{toc}{subsection}{\texorpdfstring{2.14\quad Extension by zero and completion}{2.14 Extension by zero and completion}}
% text: prose-02-59
\phantomsection\label{bp180}%
The integral coefficient sheaves are inverse limits. To restrict integral complexes to the strata of a stratification, and to extend them by zero, we therefore need the functors \(j_!\), \(j^*\), \(i^*\) and \(i_*\) to commute with limits. On the pro-étale site this holds, even for inverse image along a closed immersion, and from this we deduce the localization triangles and construct the local cohomology functor. With these tools we can argue with closed supports before the exceptional inverse image and proper base change are available.

That \(i^*\) commutes with limits is special to the pro-étale site. On \(X_{\acute{e}t}\) the stalk of \(i^*F\) at a point of \(Z\) is the stalk of \(F\). This is a colimit, and colimits do not commute with infinite products. On \(X_{\mathrm{pro}\acute{e}t}\), by contrast, the value of \(i^*M\) on a weakly contractible affine over \(Z\) is the value of \(M\) on one affine object of \(X_{\mathrm{pro}\acute{e}t}\), henselian along \(Z\) (Proposition \hyperref[bp181]{2.58}), and evaluation on one object preserves limits. For \(\mathcal O_{E,X}=\varprojlim(\mathcal O_E/\pi_E^n)_X\) this means that \(i^*\) applied to the limit is the limit of the restrictions \((\mathcal O_E/\pi_E^n)_Z\). The same holds for any integral complex given as a derived limit of its reductions.

A complex over \(\mathcal O_E\) is determined by its reductions modulo \(\pi_E^n\) when it is derived complete (Definition \hyperref[bp188]{2.61} below). Combined with the computation of the cohomology of these reductions on the étale site, this is how we reduce integral statements to finite coefficients throughout the notes.

% object: proposition-extension-zero-restriction-closed-subscheme (blueprint/chap-The-pro-etale-site/sec-Extension-by-zero-and-completion/proposition-extension-zero-restriction-closed-subscheme/_section.md)
\phantomsection\label{bp181}%
\begin{bpstatementit}{Proposition~2.58}{ (Extension by zero and restriction to a closed subscheme)}
Let \(X\) be quasi-compact and quasi-separated, \(j:U\hookrightarrow X\) a quasi-compact open immersion, and \(i:Z\hookrightarrow X\) its closed complement. For a sheaf of rings \(\mathcal R\) on \(X_{\mathrm{pro}\acute{e}t}\), we use its inverse images as the coefficient rings on \(U\) and \(Z\). Thus \(i^*\) goes from \(\mathcal R\)-module sheaves on \(X_{\mathrm{pro}\acute{e}t}\) to \(i^*\mathcal R\)-module sheaves on \(Z_{\mathrm{pro}\acute{e}t}\), \(i_*\) goes back, and \(j^*\), \(j_!\) relate \(\mathcal R\)-module sheaves on \(X_{\mathrm{pro}\acute{e}t}\) and \(j^*\mathcal R\)-module sheaves on \(U_{\mathrm{pro}\acute{e}t}\). Then \(i^*\) is exact and preserves all small limits and colimits; \(i_*\) is exact and fully faithful. The functors \(j^*\) and \(j_!\) are exact, and every \(\mathcal R\)-module sheaf \(M\) has the exact sequence

\[
0\longrightarrow j_!j^*M\longrightarrow M
\longrightarrow i_*i^*M\longrightarrow0.
\]

Moreover \(j_!\) preserves products and derived inverse limits of module complexes. Finally, for a morphism \(g:X'\to X\) with base changes \(j':U'\hookrightarrow X'\) and \(g':U'\to U\), the canonical map \(j'_!g'^*\to g^*j_!\) is an isomorphism.
\end{bpstatementit}

% proof: proof/proposition-extension-zero-restriction-closed-subscheme
\begin{proof}[{Proof}]\phantomsection\label{bp182}%
The special feature of the pro-étale site is that \(i^*\) preserves limits; we prove this first, deduce exactness of \(j_!\) from the localization sequence, obtain the statement about products and derived limits from a triangle expressing \(j_!\) through \(Rj_*\), \(i^*\) and \(i_*\), and check base change by adjunction.

\emph{Step 1: the functors \(i^*\) and \(i_*\).} The functor \(i^*\) preserves limits, because evaluating \(i^*M\) on a weakly contractible affine over \(Z\) is the same as evaluating \(M\) on an affine object of \(X_{\mathrm{pro}\acute{e}t}\), and evaluation commutes with limits. The assertion is local on \(X\). Let \(V\) be a weakly contractible object of the affine ind-étale basis of \(Z\). It determines an affine object \(v(V)\) in the corresponding basis of \(X\), henselian along its inverse image of \(Z\): the functor \(v\) is the fully faithful left adjoint to base change of these affine bases, and \(v(V)\times_XZ\simeq V\) [Stacks, \href{https://stacks.math.columbia.edu/tag/09A9}{09A9}]. The pullback of a sheaf, evaluated on \(V\), is its evaluation on \(v(V)\), because every covering of a weakly contractible object splits, so that sheafification does not change values there. Since evaluation commutes with limits and these objects detect isomorphisms, \(i^*\) preserves limits. Being a left adjoint, it also preserves colimits. Since it preserves both limits and colimits, it is exact. Lifting covers along \(v\) shows that \(i_*\) is exact and that \(i^*i_*=\mathrm{id}\) and \(j^*i_*=0\). Finally, the localization sequence is exact on the pro-étale site [Stacks, \href{https://stacks.math.columbia.edu/tag/09AH}{09AH}], that is, the adjunction map \(M\to i_*i^*M\) is surjective with kernel \(j_!j^*M\).

\emph{Step 2: exactness of \(j^*\) and \(j_!\).} The functor \(j^*\) is restriction to \(U_{\mathrm{pro}\acute{e}t}\), hence exact. The functor \(j_!\) is exact because its composites \(j^*j_!\) and \(i^*j_!\) are exact and the pair \((j^*,i^*)\) detects exactness. Here the first composite is \(j^*j_!=\mathrm{id}\), and the second is \(i^*j_!=0\). Indeed, in the localization sequence of \(j_!N\) the first map \(j_!j^*j_!N\to j_!N\) is an isomorphism since \(j^*j_!N=N\). Hence \(i_*i^*j_!N=0\), and \(i^*j_!N=0\) because \(i_*\) is fully faithful. The pair \((j^*,i^*)\) detects exactness, because both functors are exact, hence commute with taking homology of a complex of sheaves, while a sheaf \(M\) with \(j^*M=0\) and \(i^*M=0\) is zero by the localization sequence.

\emph{Step 3: products and derived inverse limits.} For \(K\in D(U_{\mathrm{pro}\acute{e}t},j^*\mathcal R)\), applying the localization sequence to a complex representing \(Rj_*K\) and using \(j^*Rj_*K\simeq K\), we obtain the triangle

\[
j_!K\longrightarrow Rj_*K\longrightarrow i_*i^*Rj_*K
\longrightarrow j_!K[1].
\]

The last two terms preserve products and derived inverse limits, because \(Rj_*\) and \(i_*\) are right adjoints and \(i^*\) is exact and preserves limits by Step 1; hence so does their fibre \(j_!\). As \(j_!\) is exact and products are exact, the corresponding statements about limits of sheaves of modules follow.

\emph{Step 4: base change.} The map \(j'_!g'^*\to g^*j_!\) is an isomorphism because its right adjoint is. Indeed, \(j'_!g'^*\) is left adjoint to \(g'_*j'^*\) and \(g^*j_!\) is left adjoint to \(j^*g_*\), and the right adjoint of the map is the base change map \(j^*g_*\to g'_*j'^*\), which is an isomorphism since both sides send a sheaf on \(X'_{\mathrm{pro}\acute{e}t}\) to its restriction to the objects lying over \(U'\). The same holds on derived categories, since the coefficient rings are inverse images, so that \(g^*\), \(g'^*\) and \(j'_!\) are exact. The argument with the functor \(v\) is [\hyperref[bp705]{BS15}, §6.1]; see also [Stacks, \href{https://stacks.math.columbia.edu/tag/09AA}{09AA}, \href{https://stacks.math.columbia.edu/tag/09AH}{09AH}, \href{https://stacks.math.columbia.edu/tag/09BM}{09BM}] and [Stacks, \href{https://stacks.math.columbia.edu/tag/09BP}{09BP}].
\end{proof}

% text: prose-02-60
\phantomsection\label{bp183}%
For arguments with closed supports we need a right adjoint of \(i_*\) before \(i^!\) is constructed by Brown representability. We build it from \(i^*\), \(j^*\) and \(Rj_*\) alone. We keep the notation of Proposition \hyperref[bp181]{2.58}.

% object: definition-local-cohomology (blueprint/chap-The-pro-etale-site/sec-Extension-by-zero-and-completion/definition-local-cohomology/_section.md)
\phantomsection\label{bp184}%
\begin{bpstatement}{Definition~2.59}{ (Local cohomology)}
The \emph{local cohomology functor} \(R\underline\Gamma_Z:D(X_{\mathrm{pro}\acute{e}t},\mathcal R)\to D(Z_{\mathrm{pro}\acute{e}t},i^*\mathcal R)\) is the right derived functor of the functor of sections supported on \(Z\),

\[
F\longmapsto i^*\ker(F\longrightarrow j_*j^*F),
\]

where \(F\to j_*j^*F\) is the unit of \(j^*\dashv j_*\).
\end{bpstatement}

% object: lemma-local-cohomology-right-adjoint (blueprint/chap-The-pro-etale-site/sec-Extension-by-zero-and-completion/lemma-local-cohomology-right-adjoint/_section.md)
\phantomsection\label{bp185}%
\begin{bpstatementit}{Lemma~2.60}{ (Local cohomology is right adjoint to $i_*$)}
The functor \(R\underline\Gamma_Z\) is right adjoint to the exact fully faithful functor \(i_*:D(Z_{\mathrm{pro}\acute{e}t},i^*\mathcal R)\to D(X_{\mathrm{pro}\acute{e}t},\mathcal R)\). It is computed by

\[
R\underline\Gamma_ZK\simeq i^*\fib(K\longrightarrow Rj_*j^*K),
\]

where \(K\to Rj_*j^*K\) is the unit of \(j^*\dashv Rj_*\).
\end{bpstatementit}

% proof: proof/lemma-local-cohomology-right-adjoint
\begin{proof}[{Proof}]\phantomsection\label{bp186}%
We first compute \(R\underline\Gamma_Z\). Represent \(K\) by a K-injective complex \(J\) with injective terms, so that \(Rj_*j^*K\) is represented by \(j_*j^*J\). The unit \(J\to j_*j^*J\) is surjective in each degree: for an injective \(\mathcal R\)-module \(I\) and an object \(V\), a section of \(j_*j^*I\) over \(V\) is a map \(j_!j^*\mathcal R_V\to I\), and it extends along the monomorphism \(j_!j^*\mathcal R_V\to\mathcal R_V\). Hence the kernel of \(J\to j_*j^*J\), taken termwise, represents the fibre of the unit, and applying the exact functor \(i^*\) gives \(R\underline\Gamma_ZK\simeq i^*\fib(K\to Rj_*j^*K)\).

Let \(F=\fib(K\to Rj_*j^*K)\). Since \(j^*i_*M=0\), we have \(\operatorname{Hom}(i_*M,Rj_*j^*K)=\operatorname{Hom}(j^*i_*M,j^*K)=0\) in all degrees, so maps \(i_*M\to K\) are the same as maps \(i_*M\to F\). The restriction \(j^*F\) vanishes, because \(j^*\) of the unit is an isomorphism, so the localization sequence gives \(F\simeq i_*i^*F\), and \(\operatorname{Hom}(i_*M,i_*i^*F)=\operatorname{Hom}(M,i^*F)\) since \(i_*\) is fully faithful.
\end{proof}

% text: prose-02-61
\phantomsection\label{bp187}%
Once \(i^!\) exists, it agrees with \(R\underline\Gamma_Z\) on the localizing category on which \(i^!\) is defined, since both are right adjoints of \(i_*\) there.

We treat integral coefficients through their reductions modulo powers of \(\pi_E\). This is possible for complexes that are recovered from their reductions, the derived complete ones. Derived completeness is the analogue for complexes of the condition \(M=\varprojlim_nM/\pi_E^nM\) on a module. For example, \(\mathcal O_{E,X}\) itself is derived complete. The sheaf \(E_X\) is not, since all its reductions vanish.

% object: definition-derived-adic-completeness (blueprint/chap-The-pro-etale-site/sec-Extension-by-zero-and-completion/definition-derived-adic-completeness/_section.md)
\phantomsection\label{bp188}%
\begin{bpstatement}{Definition~2.61}{ (Derived adic completeness)}
Let \(E/\mathbf Q_\ell\) be finite, with uniformizer \(\pi_E\), and let \(\mathcal O_{E,X}=p_X^*\underline{\mathcal O_E}\) and \((\mathcal O_E/\pi_E^n)_X\) be the integral and finite coefficient sheaves of Section \hyperref[bp149]{2.12}. A complex \(K\in D(X_{\mathrm{pro}\acute{e}t},\mathcal O_E)\) is \emph{derived complete} if its canonical map

\[
K\longrightarrow R\varprojlim_n
(K\otimes^{\mathbf L}_{\mathcal O_{E,X}}(\mathcal O_E/\pi_E^{n})_X)
\]

is an isomorphism.
\end{bpstatement}

% text: prose-02-62
\phantomsection\label{bp189}%
We write \(K_n=K\otimes^{\mathbf L}_{\mathcal O_{E,X}}(\mathcal O_E/\pi_E^{n})_X\) for the reductions of \(K\). The canonical transition maps \(K_{n+1}\to K_n\) induce isomorphisms \(K_{n+1}\otimes^{\mathbf L}_{(\mathcal O_E/\pi_E^{n+1})_X}(\mathcal O_E/\pi_E^{n})_X\simeq K_n\), by associativity of the derived tensor product. We call \((K_n)\), with these isomorphisms, the \emph{tower of reductions} of \(K\).

% object: lemma-reduction-derived-nakayama (blueprint/chap-The-pro-etale-site/sec-Extension-by-zero-and-completion/lemma-reduction-derived-nakayama/_section.md)
\phantomsection\label{bp190}%
\begin{bpstatementit}{Lemma~2.62}{ (Reduction and derived Nakayama)}
Let \(K\in D(X_{\mathrm{pro}\acute{e}t},\mathcal O_E)\). Then \(K_n\) is the cone of \(\pi_E^n:K\to K\). If \(K\) is derived complete and \(K_1=0\), then \(K=0\). Consequently a morphism between derived complete complexes is an isomorphism whenever the induced map on \(K_1\) is an isomorphism.
\end{bpstatementit}

% proof: proof/lemma-reduction-derived-nakayama
\begin{proof}[{Proof}]\phantomsection\label{bp191}%
The first assertion follows from the two-term resolution of \((\mathcal O_E/\pi_E^n)_X\) in Corollary \hyperref[bp161]{2.52}. Tensoring \(K\) with the short exact sequence \(0\to\pi_E^n\mathcal O_E/\pi_E^{n+1}\to\mathcal O_E/\pi_E^{n+1}\to
\mathcal O_E/\pi_E^n\to0\), whose first term is isomorphic to \(\mathcal O_E/\pi_E\), gives a triangle \(K_1\to K_{n+1}\to K_n\). By induction \(K_1=0\) implies \(K_n=0\) for all \(n\). Then the derived limit of the \(K_n\) vanishes, and \(K=0\) by completeness.

For the last assertion, let \(u:K\to L\) be a morphism of derived complete complexes and \(C\) its cone. Derived completeness is preserved by cones, since both sides of the completion map are exact functors of \(K\), so \(C\) is derived complete. Reduction modulo \(\pi_E\) is exact, so \(C_1\) is the cone of \(u_1:K_1\to L_1\), which vanishes when \(u_1\) is an isomorphism. Hence \(C=0\), and \(u\) is an isomorphism.
\end{proof}

% text: prose-02-63
\phantomsection\label{bp192}%
Over a point, for \(K=\mathcal O_E\) placed in degree \(0\), the reductions are \(K_n=\mathcal O_E/\pi_E^n\), and the lemma is the classical Nakayama lemma for the complete local ring \(\mathcal O_E\). The hypothesis of completeness cannot be dropped, since \(E\neq0\) but \(E/\pi_E=0\).

Let \(K\) be derived complete. Since \(R\Gamma(X_{\mathrm{pro}\acute{e}t},-)\) is a right adjoint, it commutes with the derived limit, and therefore

\[
R\Gamma(X_{\mathrm{pro}\acute{e}t},K)
\simeq R\varprojlim_nR\Gamma(X_{\mathrm{pro}\acute{e}t},K_n).
\]

When the \(K_n\) are classical bounded-below complexes, each term on the right can be computed on the étale site by Theorem \hyperref[bp103]{2.35}. In this way we compute the cohomology of \(K\) from its tower of reductions \((K_n)\).

In Section \hyperref[bp197]{3.2} we define constructible complexes over \(\mathcal O_E\), and we show that for Noetherian \(X\) they are exactly the derived complete complexes \(K\) whose reduction \(K_1\) is constructible (Proposition \hyperref[bp226]{3.11}). Derived complete complexes can equally be studied over the constant sheaf with value the discrete ring \(\mathcal O_E\) (Lemma lemma-integral-models).

The rational coefficient sheaves are obtained by inverting \(\ell\) and passing to the colimit over \(E\): \(E_X\simeq\mathcal O_{E,X}[1/\ell]\) and \(\QbarX\simeq\varinjlim_E E_X\) (Lemma \hyperref[bp158]{2.51}). Chapter \hyperref[bp193]{3} also compares constructible complexes over \(E_X\) and \(\QbarX\) with the classical adic categories. It does so by lattices and by descent to a finite extension.



% object: chap-Constructible-complexes (blueprint/chap-Constructible-complexes/_section.md), content hidden
\section*{\texorpdfstring{3\quad Constructible complexes}{3 Constructible complexes}}
\phantomsection\label{bp193}%
\addcontentsline{toc}{section}{\texorpdfstring{3\quad Constructible complexes}{3 Constructible complexes}}
% text: prose-03-01
\phantomsection\label{bp194}%
A lisse \(E\)-sheaf is the analogue of a local system: pro-étale locally it is isomorphic to \(E_X^r\) for some \(r\) (Definition \hyperref[bp201]{3.2} below). Inside such a sheaf \(L\), a lattice is a locally free \(\mathcal O_E\)-subsheaf \(L_0\) with \(L_0[1/\ell]=L\), and the first point of this chapter is that a lattice need not exist. For an example, glue \(0\) and \(\infty\) on \(\mathbf P^1\) to obtain a nodal curve \(X\). Then the infinite chain \(Y\) of projective lines, with \(\infty\) on each line glued to \(0\) on the next, maps to \(X\) by an étale covering. Moreover, the group \(\mathbf Z\) acts on \(Y\) by translation, simply transitively on the fibres, so \(Y\to X\) is an étale \(\mathbf Z\)-torsor. Using this torsor and letting the generator of \(\mathbf Z\) act on \(E\) by multiplication by \(\ell\), we obtain a lisse \(E\)-sheaf of rank one. It has no lattice, because \(\ell\) is not a unit of \(\mathcal O_E\), so that multiplication by \(\ell\) does not carry any lattice onto itself. The sheaf comes from the character \(n\mapsto\ell^n\) of \(\mathbf Z\), and since this character does not extend continuously to \(\widehat{\mathbf Z}\), the sheaf has no classical counterpart as a system \((F_n)\) (see the discussion before Lemma \hyperref[bp269]{3.24}; the example is in [\hyperref[bp705]{BS15}]).

This matters because every geometric input to the next chapters (proper base change, the cohomology of curves, Deligne's generic base change) is a theorem about torsion sheaves, whereas a rational sheaf has no reduction modulo \(\ell\): since \(E_X\) is \(\ell\)-divisible, we have \(E_X\otimes^{\mathbf L}\mathcal O_E/\pi_E=0\). To use the torsion theory we must therefore find integral objects inside rational ones, and then pass to their torsion reductions. The example above shows that this cannot be done globally.

Classically each coefficient ring comes with its own theory:

\begin{center}\small
\begin{tabularx}{\linewidth}{>{\raggedright\arraybackslash}X>{\raggedright\arraybackslash}X}
\toprule
coefficients & classical constructible objects \\
\midrule
finite ring \(A\) & constructible étale sheaves of \(A\)-modules \\
\(\mathbf Z_\ell\), \(\mathcal O_E\) & systems \((F_n)\) of constructible   \(\mathcal O_E/\pi_E^n\)-sheaves \\
 & with \(F_{n+1}\otimes\mathcal O_E/\pi_E^n\simeq F_n\) \\
\(\mathbf Q_\ell\), \(E\) & the same systems up to isogeny, that is, with \(\ell\) inverted \\
\(\overline{\mathbf Q}_\ell\) & a colimit of the categories for all finite \(E/\mathbf Q_\ell\) \\
\bottomrule
\end{tabularx}
\end{center}

On the pro-étale site, by contrast, one local condition covers all of them. Namely, a \(\Lambda_X\)-module is constructible if, on the strata of a finite stratification of \(X\), it is pro-étale locally of the form \(\Lambda_X\otimes_{\Lambda(*)}P\) for a finitely presented module \(P\) over \(\Lambda(*)\), the ring of sections over the point (Definition \hyperref[bp199]{3.1}). For finite \(A\) these are the classical constructible sheaves, while for \(E\) and \(\overline{\mathbf Q}_\ell\) they are lisse on the strata. For \(\mathcal O_E\), on the other hand, the modules \(P\) may have torsion, for example \(P=\mathcal O_E/\pi_E\).

The six operations will act on \(D^b_{\mathrm{ctf}}(X,\Lambda)\), bounded complexes with constructible cohomology and finite Tor amplitude. For \(\Lambda\) not a finite ring the Tor condition is automatic, that is, \(D^b_{\mathrm{ctf}}(X,\Lambda)=\Dbc(X,\Lambda)\) (Proposition \hyperref[bp232]{3.12}). For a finite ring, however, it matters: over \(\mathbf Z/\ell^2\) the constructible sheaf \(\mathbf F_\ell\) has infinite Tor dimension (Section \hyperref[bp197]{3.2} below).

Over \(\mathcal O_E\) the passage to torsion is a criterion: for Noetherian \(X\), a complex is constructible exactly when it is derived complete and its reduction modulo \(\pi_E\) is constructible (Proposition \hyperref[bp226]{3.11}). The criterion is not formal, because the local models of the reductions modulo \(\pi_E^n\) must be made compatible in \(n\) (Proposition \hyperref[bp223]{3.10}, Section \hyperref[bp205]{3.3}). This description is due to Bhatt and Scholze [\hyperref[bp705]{BS15}].

Over \(E\) and \(\overline{\mathbf Q}_\ell\), as the nodal curve shows, lattices exist only after stratifying, and this is enough. After descending from \(\overline{\mathbf Q}_\ell\) to a finite extension \(E\) of \(\mathbf Q_\ell\) and refining a stratification, a lisse sheaf has a lattice on each stratum (Lemma \hyperref[bp269]{3.24}), and the reductions of the lattice modulo \(\pi_E^n\) are classical torsion sheaves (Lemma \hyperref[bp274]{3.25}).

On a connected normal scheme no stratification is needed: there a lattice exists globally, and lisse \(E\)-sheaves are continuous representations of \(\pi_1(X,\bar x)\) (Proposition \bpout{13.10}, proved in Part III).

Dévissage along strata also shows that \(\Hom\) out of a constructible complex commutes with filtered colimits of complexes bounded below by a common degree (Lemma \hyperref[bp260]{3.22}); this continuity is used, together with descent to a finite \(E\), in Chapter \hyperref[bp407]{5}.

Throughout this chapter schemes are again separated and of finite type over \(k\), as in our conventions.

% object: sec-Derived-functors (blueprint/chap-Constructible-complexes/sec-Derived-functors/_section.md), content hidden
\subsection*{\texorpdfstring{3.1\quad Derived functors}{3.1 Derived functors}}
\phantomsection\label{bp195}%
\addcontentsline{toc}{subsection}{\texorpdfstring{3.1\quad Derived functors}{3.1 Derived functors}}
% text: prose-03-02
\phantomsection\label{bp196}%
Four of the six operations, \(f^*\), \(Rf_*\), \(\otimes^{\mathbf L}\) and internal Hom, exist for every condensed ring \(\Lambda\) on the whole unbounded derived categories by the usual resolutions. The remaining two need more work: \(Rf_!\) is built from proper direct image and extension by zero (Section \hyperref[bp459]{5.4}), and its right adjoint comes from Brown representability.

Sheaves of \(\Lambda_X\)-modules can be viewed as module objects over the ring object \(\Lambda_X\) in \(\Sh(X_{\mathrm{pro}\acute{e}t},\operatorname{Ab})\), with the tensor product over \(\Lambda_X\) given by the coequalizer of the two actions \(M\otimes\Lambda_X\otimes N\rightrightarrows M\otimes N\). The category of \(\Lambda_X\)-modules is Grothendieck abelian. Hence complexes have K-injective resolutions, which define \(Rf_*\) and \(\cHom\) [Stacks, \href{https://stacks.math.columbia.edu/tag/07A5}{07A5}, \href{https://stacks.math.columbia.edu/tag/08J7}{08J7}]. They also have K-flat resolutions, which define \(\otimes^{\mathbf L}\), and \(K\otimes^{\mathbf L}-\) is left adjoint to \(\cHom(K,-)\) [Stacks, \href{https://stacks.math.columbia.edu/tag/06YL}{06YL}].

The pullback of a \(\Lambda_Y\)-module \(M\) along \(f:X\to Y\) is \(f^*M=\Lambda_X\otimes_{f^{-1}\Lambda_Y}f^{-1}M\). Since \(f^{-1}\Lambda_Y\simeq\Lambda_X\) (Lemma \hyperref[bp167]{2.54}), this is just \(f^{-1}M\). Therefore \(f^*:\operatorname{Mod}(Y_{\mathrm{pro}\acute{e}t},\Lambda)\to
\operatorname{Mod}(X_{\mathrm{pro}\acute{e}t},\Lambda)\) is exact, and on derived categories it is left adjoint to \(Rf_*:D(X_{\mathrm{pro}\acute{e}t},\Lambda)\to
D(Y_{\mathrm{pro}\acute{e}t},\Lambda)\). For an open immersion \(j:U\to X\), extension by zero \(j_!:\operatorname{Mod}(U_{\mathrm{pro}\acute{e}t},\Lambda)\to\operatorname{Mod}(X_{\mathrm{pro}\acute{e}t},\Lambda)\) is exact and left adjoint to \(j^*\).

Constructible sheaves will be built from lisse sheaves on the strata of a stratification, each extended by zero to all of \(X\) (Lemma \hyperref[bp248]{3.18} and Corollary \hyperref[bp253]{3.20}). We fix this extension by zero now. Let \(a:S\hookrightarrow X\) be a locally closed immersion, for instance the inclusion of a stratum. Then \(S\) is open in its closure, so \(a\) factors as \(a=i\circ j\), with \(j:S\hookrightarrow\overline S\) the open immersion into the closure and \(i:\overline S\hookrightarrow X\) the closed immersion. We put

\[
a_!=i_*j_!:\operatorname{Mod}(S_{\mathrm{pro}\acute{e}t},\Lambda)\longrightarrow\operatorname{Mod}(X_{\mathrm{pro}\acute{e}t},\Lambda).
\]

For a \(\Lambda_S\)-module \(L\), the sheaf \(a_!L\) restricts to \(L\) on \(S\), and its restrictions to \(\overline S\setminus S\) and to \(X\setminus\overline S\) are zero. Indeed \(a^*a_!=j^*i^*i_*j_!=\mathrm{id}\), the sheaf \(j_!L\) vanishes on \(\overline S\setminus S\), and \(i_*\) of a sheaf vanishes on \(X\setminus\overline S\) (Proposition \hyperref[bp181]{2.58}, which applies because the open immersions involved are quasi-compact, \(X\) being Noetherian). The same proposition shows that \(j_!\) and \(i_*\) are exact, so \(a_!\) is exact. It therefore acts on complexes term by term and defines \(a_!:D(S_{\mathrm{pro}\acute{e}t},\Lambda)\to D(X_{\mathrm{pro}\acute{e}t},\Lambda)\) with no derived functor needed. Once \(Rf_!\) is constructed (Section \hyperref[bp459]{5.4}), this functor is \(Ra_!\).


% object: sec-Constructibility (blueprint/chap-Constructible-complexes/sec-Constructibility/_section.md), content hidden
\subsection*{\texorpdfstring{3.2\quad Constructibility}{3.2 Constructibility}}
\phantomsection\label{bp197}%
\addcontentsline{toc}{subsection}{\texorpdfstring{3.2\quad Constructibility}{3.2 Constructibility}}
% text: prose-03-03
\phantomsection\label{bp198}%
With finite coefficients a constructible sheaf is, on the strata of a finite stratification, locally constant with finite stalks. The six operations, however, will act on constructible complexes of finite Tor amplitude for every coefficient ring, so we need a definition of constructibility that covers all coefficient rings at once. We give it by a local condition (Definition \hyperref[bp199]{3.1}). This condition has two parts, a stratification and a local model that replaces local constancy, and the complexes we use must in addition have finite Tor amplitude. As we explain below, each of the three conditions answers a specific need. For \(\mathcal O_E\) we then show that constructibility can be tested on the reduction modulo \(\pi_E\) (Proposition \hyperref[bp226]{3.11}). For a morphism \(g:Z\to X\) and a \(\Lambda_X\)-module \(M\) we write \(M|_Z=g^*M\), a \(\Lambda_Z\)-module.

A stratification cannot be avoided, already for sheaves coming from geometry: for \(j:\mathbf G_m\hookrightarrow\mathbf A^1\), the sheaf \(j_!\Lambda\) is zero at the origin and \(\Lambda\) elsewhere, so it is not locally constant on \(\mathbf A^1\).

Local constancy itself has to be rephrased. For instance, for \(\Lambda=\mathbf Z_\ell\), “locally constant” with value a \(\mathbf Z_\ell\)-module would mean locally a constant sheaf with value the discrete group, but on \(\operatorname{Spec}\mathbf F_q\) such a sheaf has \(H^1=0\) (Section \bpout{1.2}), whereas the answer we want is \(\mathbf Z_\ell\). We therefore ask instead that \(M\) be locally of the form \(\Lambda_U\otimes_{\Lambda(*)}P\), with \(P\) a finitely presented module over the ring \(\Lambda(*)\) of sections over the point, as in the opening of this chapter. This keeps the \(\ell\)-adic topology. For finite \(\Lambda\) it is local constancy, and for \(E\) it says that \(M\) is locally free of finite rank. In this formulation, finite presentation of \(P\) replaces finiteness of the stalks.

Finally, the six operations should commute with change of coefficients, and this needs finite Tor amplitude. Over \(\mathbf Z/\ell^2\) the sheaf \(\mathbf F_\ell\) is constructible, but \(\mathbf F_\ell\otimes^{\mathbf L}_{\mathbf Z/\ell^2}\mathbf F_\ell\) has cohomology \(\mathbf F_\ell\) in every degree \(\leq0\). With the Tor condition imposed, by contrast, the category \(D^b_{\mathrm{ctf}}\) is stable under change of coefficients and under \(\otimes^{\mathbf L}\), and for \(\mathcal O_E\), \(E\) and \(\overline{\mathbf Q}_\ell\) the Tor condition is automatic (Proposition \hyperref[bp232]{3.12}).

% object: definition-constructible-sheaves (blueprint/chap-Constructible-complexes/sec-Constructibility/definition-constructible-sheaves/_section.md)
\phantomsection\label{bp199}%
\begin{bpstatement}{Definition~3.1}{ (Constructible sheaves)}
Let \(\Lambda\) be a commutative condensed ring.

\begin{enumerate}
\item Consider the following condition on a \(\Lambda_X\)-module \(M\): there is a pro-étale cover \(\{U_i\to X\}\) and finitely presented modules \(P_i\) over the ring \(\Lambda(*)\) of sections over the point such that \[
   M|_{U_i}\simeq\Lambda_{U_i}\otimes_{\Lambda(*)}P_i.
   \] Here \(\Lambda(*)\) and \(P_i\) are constant sheaves on \(U_i\), acting through \(\Lambda(*)\to\Lambda_{U_i}\). Equivalently, on each \(U_i\) the module \(M\) has a presentation by a finite matrix with entries in \(\Lambda(*)\).
\item For qcqs \(X\), a \(\Lambda_X\)-module \(M\) is \emph{constructible} if there is a finite stratification of \(X\) by constructible locally closed subschemes such that on each stratum \(S\) the restriction \(M|_S\) satisfies the condition of (1). On an arbitrary scheme this condition is imposed locally on affine opens.
\end{enumerate}

We write \(\Cons(X,\Lambda)\subset\operatorname{Mod}(X_{\mathrm{pro}\acute{e}t},\Lambda)\) for the full subcategory of constructible sheaves, and \(\Dbc(X,\Lambda)\subset
D(X_{\mathrm{pro}\acute{e}t},\Lambda)\) for the full subcategory of bounded complexes with constructible cohomology sheaves. A complex \(K\) has \emph{Tor amplitude in \([a,b]\)} if \(K\otimes^{\mathbf L}_{\Lambda_X}N\) has cohomology only in \([a,b]\) for every sheaf \(N\) of \(\Lambda_X\)-modules. We write \(D^b_{\mathrm{ctf}}(X,\Lambda)\) for the full subcategory of \(\Dbc(X,\Lambda)\) of complexes of finite Tor amplitude.
\end{bpstatement}

% text: prose-03-04
\phantomsection\label{bp200}%
The modules that satisfy the condition of (1) on all of \(X\) are the analogue of local systems, and they are the building blocks on the strata.

% object: definition-lisse-sheaves (blueprint/chap-Constructible-complexes/sec-Constructibility/definition-lisse-sheaves/_section.md)
\phantomsection\label{bp201}%
\begin{bpstatement}{Definition~3.2}{ (Lisse sheaves)}
Let \(\Lambda\) be a commutative condensed ring. A \(\Lambda_X\)-module \(M\) is \emph{lisse} if it satisfies the condition of Definition definition-constructible-sheaves on \(X\) itself, without a stratification. A complex \(K\in D(X_{\mathrm{pro}\acute{e}t},\Lambda)\) is \emph{lisse} if its cohomology sheaves are lisse.
\end{bpstatement}

% text: prose-03-05
\phantomsection\label{bp202}%
For the individual coefficient rings the definitions take a familiar form.

% object: remark-constructible-sheaves-coefficient-rings (blueprint/chap-Constructible-complexes/sec-Constructibility/remark-constructible-sheaves-coefficient-rings/_section.md)
\phantomsection\label{bp203}%
\begin{bpstatement}{Remark~3.3}{ (Constructible sheaves for the individual coefficient rings)}
\begin{enumerate}
\item For Noetherian \(X\), every locally closed subset is constructible, so in Definition definition-constructible-sheaves any finite stratification by locally closed subschemes may be used.
\item For finite \(A\) and qcqs \(X\), \(D^b_{\mathrm{ctf}}(X,A)\) is the pullback of the usual étale category \(D^b_{\mathrm{ctf}}(X_{\acute{e}t},A)\) by Corollary \hyperref[bp132]{2.44}.
\item If \(\Lambda=E\) or \(\overline{\mathbf Q}_\ell\), the ring \(\Lambda(*)\) is a field, so each \(P_i\) is a finite-dimensional vector space, and a lisse module is locally free of finite rank. In other words, a constructible \(E\)- or \(\overline{\mathbf Q}_\ell\)-sheaf is one that is lisse on the strata of a finite stratification [\hyperref[bp705]{BS15}, Definition 6.8.6].
\item If \(\Lambda=\mathcal O_E\), locally on a stratum a constructible \(\mathcal O_{E,X}\)-module is \(P\otimes_{\mathcal O_E}\mathcal O_{E,X}\) for a finitely generated \(\mathcal O_E\)-module \(P\). For example, \(\mathcal O_{E,X}/\pi_E\) is constructible.
\item For the admissible coefficients, the constructible sheaves form an abelian category. For finite \(A\) this is the classical theory of sheaves with finite stalks, and the other cases are Lemma \hyperref[bp248]{3.18}, Corollary \hyperref[bp253]{3.20} and Lemma \hyperref[bp471]{5.19}, which we prove in Subsection \hyperref[bp469]{5.4.2}.
\item For finite \(A\), a lisse \(A_X\)-module is the same as a locally constant sheaf with finite stalks. For \(\Lambda=\mathcal O_E\) a lisse module may have torsion, such as \(\mathcal O_{E,X}/\pi_E\); the lisse \(\mathcal O_{E,X}\)-modules without torsion are the locally free ones of finite rank, since a finitely generated torsion free \(\mathcal O_E\)-module is free.
\end{enumerate}
\end{bpstatement}

% text: prose-03-06
\phantomsection\label{bp204}%
Part (5) is special to the admissible coefficients. For a general condensed ring, kernels of maps between constructible sheaves can fail to be constructible, and then \(\Dbc(X,\Lambda)\) is not triangulated.


% object: sec-Integral-constructible-complexes (blueprint/chap-Constructible-complexes/sec-Integral-constructible-complexes/_section.md), content hidden
\subsection*{\texorpdfstring{3.3\quad Integral constructible complexes}{3.3 Integral constructible complexes}}
\phantomsection\label{bp205}%
\addcontentsline{toc}{subsection}{\texorpdfstring{3.3\quad Integral constructible complexes}{3.3 Integral constructible complexes}}
% text: prose-03-07
\phantomsection\label{bp206}%
Constructibility over \(\mathcal O_E\) is defined by a local condition, but in practice we know a complex \(K\in D(X_{\mathrm{pro}\acute{e}t},\mathcal O_E)\) through its reductions \(K_n=K\otimes^{\mathbf L}_{\mathcal O_{E,X}}(\mathcal O_E/\pi_E^n)_X\) of Definition \hyperref[bp188]{2.61}. The criterion below says that \(K\) is constructible exactly when it is derived complete and \(K_1\) is constructible. Write \(\mathbf F_E=\mathcal O_E/\pi_E\) for the residue field, so that \(K_1\in D(X_{\mathrm{pro}\acute{e}t},\mathbf F_E)\).

To prove the criterion we must produce, on each stratum, one complex of finite free \(\mathcal O_E\)-modules that is locally a model for all the \(K_n\) at once. Each \(K_n\) separately is easy to control: its bounds and its stratification are inherited from \(K_1\) by dévissage along the powers of \(\pi_E\), and on a stratum it is locally isomorphic to a constant perfect complex. The difficulty is that these local isomorphisms are chosen independently at each level, so they have to be made compatible before we pass to the limit. On a weakly contractible cover we do this by a Mittag–Leffler argument, which needs a uniform rate. Such a rate exists because a bounded complex of finite free modules over the discrete valuation ring \(\mathcal O_E\) has finitely generated cohomology, and the torsion of this cohomology is killed by a fixed power of \(\pi_E\). Throughout, we use the étale site only at the finite levels, through Theorem \hyperref[bp129]{2.43}.

% object: example-integral-criterion-point (blueprint/chap-Constructible-complexes/sec-Integral-constructible-complexes/example-integral-criterion-point/_section.md)
\phantomsection\label{bp207}%
\begin{bpstatement}{Example~3.4}{ (The criterion over a point)}
On the spectrum of an algebraically closed field, take for \(K\) the sheaf attached to an \(\mathcal O_E\)-module. The computation is then one with modules, and it shows where the uniform rate comes from. Resolving \(\mathcal O_E/\pi_E^e\) by \([\mathcal O_E\xrightarrow{\pi_E^e}\mathcal O_E]\) in degrees \(-1,0\) and reducing modulo \(\pi_E\), we find

\[
\mathcal O_E\otimes^{\mathbf L}_{\mathcal O_E}\mathbf F_E=\mathbf F_E,
\qquad
\mathcal O_E/\pi_E^e\otimes^{\mathbf L}_{\mathcal O_E}\mathbf F_E
=\bigl[\mathbf F_E\xrightarrow{\,0\,}\mathbf F_E\bigr]
=\mathbf F_E\oplus\mathbf F_E[1]\quad(e\geq1),
\]

since \(\pi_E^e\) becomes zero modulo \(\pi_E\). For \(P=\mathcal O_E^r\oplus\bigoplus_{i=1}^s\mathcal O_E/\pi_E^{e_i}\) the reduction \(K_1\) has \(H^0=\mathbf F_E^{r+s}\) and \(H^{-1}=\mathbf F_E^s\). Thus \(H^{-1}(K_1)\) counts the cyclic torsion summands. For \(K=\bigoplus_{n\geq1}\mathcal O_E/\pi_E^n\) we get \(K_1=\bigoplus_{n\geq1}(\mathbf F_E\oplus\mathbf F_E[1])\). This is not constructible, because its cohomology is an infinite-dimensional vector space in degrees \(0\) and \(-1\). In other words, a constructible \(K_1\) allows only finitely many torsion summands. For \(P\) the single power \(\pi_E^{\max e_i}\) then kills all the torsion, and this is the uniform rate.

The completeness hypothesis cannot be dropped either. The module \(E\) is \(\pi_E\)-divisible, so \(E\otimes^{\mathbf L}\mathcal O_E/\pi_E^n=0\) for all \(n\). Hence \(K_1=0\) is constructible, yet \(E\) is not a finitely generated \(\mathcal O_E\)-module. The criterion does not apply because \(E\) is not derived complete: its derived completion \(R\varprojlim_n(E\otimes^{\mathbf L}\mathcal O_E/\pi_E^n)\) is zero.
\end{bpstatement}

% object: lemma-bounds-finite-reductions (blueprint/chap-Constructible-complexes/sec-Integral-constructible-complexes/lemma-bounds-finite-reductions/_section.md)
\phantomsection\label{bp208}%
\begin{bpstatementit}{Lemma~3.5}{ (Bounds for the finite reductions)}
Let \(X\) be Noetherian and let \(K\in D(X_{\mathrm{pro}\acute{e}t},\mathcal O_E)\) be derived complete with \(K_1\in\Dbc(X,\mathbf F_E)\). Choose \([a,b]\) containing the cohomology of \(K_1\). Then every \(K_n\) lies in \(\Dbc(X,\mathcal O_E/\pi_E^n)\) and has Tor amplitude in \([a,b]\), and \(K\) has cohomology only in \([a,b]\). If the cohomology sheaves of \(K_1\) are locally constant on a stratum, then so are those of every \(K_n\).
\end{bpstatementit}

% proof: proof/lemma-bounds-finite-reductions
\begin{proof}[{Proof}]\phantomsection\label{bp209}%
Everything about \(K_n\) is obtained from \(K_1\) by induction on \(n\), using that \(K_n\) is an extension of \(K_{n-1}\) by \(K_1\). The Tor bound needs an extra argument with minimal resolutions, and the bound on \(K\) itself then follows from the Milnor sequence.

The filtration of \(\mathcal O_E/\pi_E^n\) by the powers of \(\pi_E\) has graded pieces \(\mathbf F_E\). Tensoring its short exact sequences with \(K\) over \(\mathcal O_{E,X}\), we get triangles \(K_1\to K_{n}\to K_{n-1}\), up to the identification \(\pi_E^{n-1}\mathcal O_E/\pi_E^n\simeq\mathbf F_E\). Moreover, complexes with constructible cohomology in \([a,b]\) are closed under extensions in triangles, by the long exact cohomology sequence and the fact that constructible sheaves form a weak Serre subcategory. Induction on \(n\) now shows that \(K_n\) is bounded in \([a,b]\) with constructible cohomology. Since locally constant finite sheaves of modules also form a weak Serre subcategory, the same induction works on a fixed stratum, and a single stratification works for all \(n\). By the comparison theorem for constructible sheaves (Theorem \hyperref[bp129]{2.43}), these finite sheaves and complexes are classical, that is, pulled back from the étale site.

For the Tor bound we use \(K_n\otimes^{\mathbf L}_{\mathcal O_E/\pi_E^n}\mathbf F_E\simeq K_1\). At an étale geometric point \(\bar x\), take a bounded above resolution of the stalk of \(K_n\) by finite free modules over the Artinian local ring \(\mathcal O_E/\pi_E^n\). Cancelling a two-term identity summand whenever a differential matrix has a unit entry, we obtain a minimal resolution, whose differentials have entries in \((\pi_E)\), so they vanish after tensoring with \(\mathbf F_E\). Since stalks commute with derived tensor products, the minimal resolution tensored with \(\mathbf F_E\) computes the stalk of \(K_1\) at \(\bar x\), with zero differential. Its rank in degree \(q\) is therefore \(\dim_{\mathbf F_E}H^q(K_1)_{\bar x}\). This vanishes outside \([a,b]\), and we conclude that the stalk has a bounded finite free resolution in \([a,b]\). Such a bound on stalks gives the same Tor bound on the étale site, and Corollary \hyperref[bp132]{2.44} carries it to the pro-étale site, for Tor against arbitrary sheaves of modules there.

By completeness \(K=R\varprojlim K_n\), so the Milnor sequence (2.2) gives the lower bound and the vanishing above \(b+1\). For \(H^{b+1}(K)\) we use the top degree, where \(H^b(K_n)=H^b(K_{n+1})\otimes_{\mathcal O_E/\pi_E^{n+1}}\mathcal O_E/\pi_E^n\), so the transition maps are epimorphisms. Their \(R^1\varprojlim\) vanishes by the last assertion of Proposition \hyperref[bp71]{2.24}, and this means that \(H^{b+1}(K)=0\).
\end{proof}

% text: prose-03-08
\phantomsection\label{bp210}%
To make the local trivializations of the \(K_n\) compatible, we need two facts about complexes over \(\mathcal O_E\): a compatible system of perfect complexes over the rings \(\mathcal O_E/\pi_E^n\) comes from a single bounded finite free complex, and the images of the cohomology at finite levels stabilize at a uniform rate.

% object: lemma-perfect-complexes-over-oe (blueprint/chap-Constructible-complexes/sec-Integral-constructible-complexes/lemma-perfect-complexes-over-oe/_section.md)
\phantomsection\label{bp211}%
\begin{bpstatementit}{Lemma~3.6}{ (Perfect complexes over $\mathcal O_E$)}
For \(n\ge1\) let \(M_n\in D(\mathcal O_E/\pi_E^n)\) be perfect, with maps \(M_{n+1}\to M_n\) inducing isomorphisms \(M_{n+1}\otimes^{\mathbf L}_{\mathcal O_E/\pi_E^{n+1}}\mathcal O_E/\pi_E^n\simeq M_n\). Then there is a bounded complex \(P^\bullet\) of finite free \(\mathcal O_E\)-modules with isomorphisms \(P^\bullet/\pi_E^n\simeq M_n\) in \(D(\mathcal O_E/\pi_E^n)\), compatible with the transition maps.
\end{bpstatementit}

% proof: proof/lemma-perfect-complexes-over-oe
\begin{proof}[{Proof}]\phantomsection\label{bp212}%
We use that a minimal complex over the local ring \(\mathcal O_E/\pi_E^n\) is determined up to isomorphism of complexes, not only up to homotopy, so that the transition maps can be made into equalities of matrices and passed to the limit.

Represent each \(M_n\) by a minimal finite free complex, obtained by cancelling unit entries in the differential matrices. Its reduction modulo \(\pi_E\) has zero differential, so its rank in degree \(q\) is \(\dim_{\mathbf F_E}H^q(M_1)\), which is independent of \(n\) and nonzero only in finitely many degrees. The given isomorphism \(M_{n+1}/\pi_E^n\simeq M_n\) between minimal complexes is represented by a chain homotopy equivalence \(f\), with homotopy inverse \(g\). Since both complexes have zero differential modulo \(\pi_E\), the relations \(gf-1=dh+hd\) and \(fg-1=dh'+h'd\) show that \(f\) is an isomorphism modulo \(\pi_E\) in each degree. By Nakayama's lemma its matrices are then invertible over \(\mathcal O_E/\pi_E^n\). Lifting them to invertible matrices over \(\mathcal O_E/\pi_E^{n+1}\) and changing bases, we may assume that the differential at level \(n+1\) reduces exactly to the differential at level \(n\), with the given transition in the derived category. Induction on \(n\) gives compatible matrices, and their limits define \(P^\bullet\). Its differential squares to zero because it does so modulo every \(\pi_E^n\), and the identifications with the \(M_n\) are compatible by construction.
\end{proof}

% text: prose-03-09
\phantomsection\label{bp213}%
The second fact is used in the proof of Lemma \hyperref[bp221]{3.9} for the internal Hom complex \((P^\bullet)^\vee\otimes P^\bullet\), whose \(H^0\) at level \(n\) consists of the maps \(P^\bullet/\pi_E^n\to P^\bullet/\pi_E^n\) in the derived category. There it says that a map at level \(n\) which lifts far enough lifts to all levels, and how far is enough depends only on the torsion of the cohomology.

% object: lemma-images-cohomology-finite-levels (blueprint/chap-Constructible-complexes/sec-Integral-constructible-complexes/lemma-images-cohomology-finite-levels/_section.md)
\phantomsection\label{bp214}%
\begin{bpstatementit}{Lemma~3.7}{ (Images of cohomology at finite levels)}
Let \(Q^\bullet\) be a bounded complex of finite free \(\mathcal O_E\)-modules, and let \(\pi_E^e\) kill the torsion of \(H^{q+1}(Q^\bullet)\). Then for \(m\ge n+e\)

\[
\operatorname{im}\bigl(H^q(Q^\bullet/\pi_E^m)\to
H^q(Q^\bullet/\pi_E^n)\bigr)
=\operatorname{im}\bigl(H^q(Q^\bullet)\to
H^q(Q^\bullet/\pi_E^n)\bigr).
\]
\end{bpstatementit}

% text: prose-03-10
\phantomsection\label{bp215}%
Such an \(e\) exists, since \(\mathcal O_E\) is a Noetherian discrete valuation ring and \(H^{q+1}(Q^\bullet)\) is finitely generated.

% proof: proof/lemma-images-cohomology-finite-levels
\begin{proof}[{Proof of Lemma 3.7}]\phantomsection\label{bp216}%
The two-term resolution of \(\mathcal O_E/\pi_E^n\) gives an exact sequence

\[
0\to H^q(Q^\bullet)/\pi_E^n\to H^q(Q^\bullet/\pi_E^n)
\to H^{q+1}(Q^\bullet)[\pi_E^n]\to0.
\]

For the transition from level \(m\) to level \(n\), the induced map on the right-hand terms is multiplication by \(\pi_E^{m-n}\), because the corresponding map of resolutions is the identity in degree zero and multiplication by \(\pi_E^{m-n}\) in degree \(-1\). If \(m-n\ge e\), this map is zero, since \(H^{q+1}(Q^\bullet)[\pi_E^m]\) consists of torsion elements and \(\pi_E^e\) kills them. The image of \(H^q(Q^\bullet/\pi_E^m)\) therefore lies in the left-hand term \(H^q(Q^\bullet)/\pi_E^n\) at level \(n\), which is the image of \(H^q(Q^\bullet)\). The reverse inclusion holds because the map from \(H^q(Q^\bullet)\) factors through level \(m\).
\end{proof}

% text: prose-03-11
\phantomsection\label{bp217}%
The bound \(m\geq n+e\) cannot be improved. To see this, take \(Q^\bullet=[\mathcal O_E\xrightarrow{\pi_E^e}\mathcal O_E]\) in degrees \(0,1\), so that \(H^0(Q^\bullet)=0\) and \(H^1(Q^\bullet)=\mathcal O_E/\pi_E^e\). For \(m\geq e\) the group \(H^0(Q^\bullet/\pi_E^m)\) is \(\pi_E^{m-e}\mathcal O_E/\pi_E^m\), and its image at level \(n\) is generated by \(\pi_E^{m-e}\) modulo \(\pi_E^n\). This image is zero exactly when \(m\geq n+e\), in agreement with the image of \(H^0(Q^\bullet)=0\).

At a single finite level the local structure is a statement about the étale site: a perfect complex with locally constant cohomology is locally constant. On a connected scheme its local model is moreover the same at every point.

% object: lemma-local-constancy-finite-level (blueprint/chap-Constructible-complexes/sec-Integral-constructible-complexes/lemma-local-constancy-finite-level/_section.md)
\phantomsection\label{bp218}%
\begin{bpstatementit}{Lemma~3.8}{ (Locally constant perfect complexes)}
Let \(S\) be a connected Noetherian scheme, let \(A\) be a finite local ring, and let \(L\in D^b(S_{\acute{e}t},A)\) have finite Tor amplitude and cohomology sheaves that are locally constant with finite stalks. Let \(\bar s\) be a geometric point of \(S\) and let \(M\) be a bounded complex of finite free \(A\)-modules with \(M\simeq L_{\bar s}\) in \(D(A)\). Then \(L\) is étale locally on \(S\) isomorphic to the constant complex \(\underline M\).
\end{bpstatementit}

% proof: proof/lemma-local-constancy-finite-level
\begin{proof}[{Proof}]\phantomsection\label{bp219}%
We first show that a model of a stalk spreads out to a neighbourhood, and then use connectedness to see that the model is the same everywhere.

Every stalk \(L_{\bar t}\) has finite Tor amplitude and finite cohomology, hence is perfect over the local ring \(A\), that is, represented by a bounded complex \(M_{\bar t}\) of finite free \(A\)-modules. An isomorphism \(M_{\bar t}\simeq L_{\bar t}\) is then a section at \(\bar t\) of the zeroth cohomology sheaf of \(R\cHom(\underline{M_{\bar t}},L)\); here the stalks of this sheaf are computed stalkwise, since \(\underline{M_{\bar t}}\) is a bounded complex of finite free modules. We extend the section to an étale neighbourhood of \(\bar t\) and lift it there to a morphism \(\underline{M_{\bar t}}\to L\). Its cone has bounded locally constant finite cohomology and zero stalk at \(\bar t\). Hence its finitely many cohomology sheaves vanish after shrinking the neighbourhood, and \(L\) is isomorphic to \(\underline{M_{\bar t}}\) there. In particular, the stalks of \(L\) at all geometric points over the image of this neighbourhood are isomorphic to \(M_{\bar t}\) in \(D(A)\).

For each isomorphism class \(c\) in \(D(A)\), let \(S_c\subset S\) be the set of points at whose geometric points the stalk of \(L\) lies in \(c\). By the previous paragraph and openness of étale morphisms, each \(S_c\) is open. Moreover, the sets \(S_c\) are disjoint and cover \(S\). Since \(S\) is connected, \(S=S_c\) for the class \(c\) of \(M\). Thus every point has an étale neighbourhood on which \(L\simeq\underline M\); compare [Stacks, \href{https://stacks.math.columbia.edu/tag/09BI}{09BI}].
\end{proof}

% text: prose-03-12
\phantomsection\label{bp220}%
On a weakly contractible object the local isomorphisms of the finite levels become global, and the only remaining problem is to choose them compatibly. For this we use the uniform rate of Lemma \hyperref[bp214]{3.7}.

% object: lemma-integral-model-weakly-contractible (blueprint/chap-Constructible-complexes/sec-Integral-constructible-complexes/lemma-integral-model-weakly-contractible/_section.md)
\phantomsection\label{bp221}%
\begin{bpstatementit}{Lemma~3.9}{ (Integral models on weakly contractible schemes)}
Let \(W\) be an affine weakly contractible scheme, let \(K\in D(W_{\mathrm{pro}\acute{e}t},\mathcal O_E)\) be derived complete, and let \(P^\bullet\) be a bounded complex of finite free \(\mathcal O_E\)-modules. Suppose that for every \(n\geq1\) there is an isomorphism \((P^\bullet/\pi_E^n)_W\simeq K_n\) in \(D(W_{\mathrm{pro}\acute{e}t},\mathcal O_E/\pi_E^n)\). Then \(K\simeq\mathcal O_{E,W}\otimes_{\mathcal O_E}P^\bullet\).
\end{bpstatementit}

% proof: proof/lemma-integral-model-weakly-contractible
\begin{proof}[{Proof}]\phantomsection\label{bp222}%
The hypothesis gives an isomorphism \((P^\bullet/\pi_E^n)_W\simeq K_n\) at each level \(n\), but these isomorphisms need not be compatible with the reduction maps, so they do not assemble to a map into \(K=R\varprojlim_nK_n\). What we need is a compatible family of maps, that is, an element of the inverse limit of the Hom groups between the reductions. Such an element lifts to a map \(\mathcal O_{E,W}\otimes_{\mathcal O_E}P^\bullet\to K\), because \(K\) is derived complete. To find one that is an isomorphism at level \(1\), we show that these Hom groups satisfy the Mittag–Leffler condition with a rate \(e\) independent of \(n\). Derived Nakayama then shows that the lift is an isomorphism.

Consider the tower of complexes

\[
\cHom_{(\mathcal O_E/\pi_E^n)_W}\bigl((P^\bullet/\pi_E^n)_W,K_n\bigr)
\in D(W_{\mathrm{pro}\acute{e}t},\mathcal O_E/\pi_E^n),
\]

with transition maps induced by the reduction maps \(K_{n+1}\to K_n\), and the tower of groups \(\Hom_{D(W_{\mathrm{pro}\acute{e}t},\mathcal O_E/\pi_E^n)}\bigl((P^\bullet/\pi_E^n)_W,K_n\bigr)\) obtained by applying \(H^0R\Gamma(W,-)\). Since \(P^\bullet\) is finite free, the first tower is the tower of reductions of \((\mathcal O_{E,W}\otimes_{\mathcal O_E}(P^\bullet)^\vee)\otimes^{\mathbf L}_{\mathcal O_{E,W}}K\).

These Hom groups can be infinite, so the Mittag–Leffler condition has to come from a uniform bound. Let \(Q^\bullet=(P^\bullet)^\vee\otimes_{\mathcal O_E}P^\bullet\) and choose \(e\) such that \(\pi_E^e\) kills the torsion of \(H^1(Q^\bullet)\). Let \(N\ge m\ge n+e\) and choose one isomorphism \((P^\bullet/\pi_E^N)_W\simeq K_N\). Its reductions to levels \(m\) and \(n\) identify the three internal Hom complexes at levels \(N\to m\to n\) with the constant complexes \(Q^\bullet/\pi_E^N\to Q^\bullet/\pi_E^m\to Q^\bullet/\pi_E^n\). Their zeroth cohomology sheaves are therefore the constant sheaves on the groups \(H^0(Q^\bullet/\pi_E^N)\to H^0(Q^\bullet/\pi_E^m)\to H^0(Q^\bullet/\pi_E^n)\). By Lemma \hyperref[bp214]{3.7} the images from levels \(N\) and \(m\) at level \(n\) agree, as subsheaves. Since \(W\) is weakly contractible, taking sections over \(W\) is exact, and so it commutes with taking cohomology and images. Hence the images in the Hom group at level \(n\) from levels \(N\) and \(m\) coincide as well. Since \(e\) does not depend on \(n\), \(m\) or \(N\), we obtain the Mittag–Leffler condition: for every \(n\) the images at level \(n\) are constant from level \(n+e\) on.

To lift, note that each Hom group contains an isomorphism, and its reduction to level \(1\) is again an isomorphism. The stable image at level \(1\) is the image from level \(1+e\), so it contains an invertible class. By the Mittag–Leffler property this class lifts to a compatible family \((\varphi_n)_n\) in the inverse limit of the Hom groups. The Milnor sequence gives a surjection

\[
\begin{aligned}
&H^0R\varprojlim_n R\Gamma\bigl(W,\cHom_{(\mathcal O_E/\pi_E^n)_W}
((P^\bullet/\pi_E^n)_W,K_n)\bigr)\\
&\qquad\longrightarrow
\varprojlim_n\Hom_{D(W_{\mathrm{pro}\acute{e}t},\mathcal O_E/\pi_E^n)}
\bigl((P^\bullet/\pi_E^n)_W,K_n\bigr),
\end{aligned}
\]

whose source is \(\Hom_{D(W_{\mathrm{pro}\acute{e}t},\mathcal O_E)}(\mathcal O_{E,W}\otimes_{\mathcal O_E}P^\bullet,K)\). This holds because derived sections commute with inverse limits, \(K\) is complete, and tensoring with the finite free dual of \(P^\bullet\) commutes with inverse limits. The family \((\varphi_n)\) therefore lifts to a morphism \(\mathcal O_{E,W}\otimes_{\mathcal O_E}P^\bullet\to K\) whose first reduction is an isomorphism. Since both objects are complete, derived Nakayama (Lemma \hyperref[bp190]{2.62}) shows that this morphism is an isomorphism. Here completeness of the source follows termwise from \(\mathcal O_{E,W}=\varprojlim(\mathcal O_E/\pi_E^n)_W\). Indeed, the transition maps are epimorphisms, so this limit agrees with the derived limit by Proposition \hyperref[bp71]{2.24}.
\end{proof}

% object: proposition-local-structure-integral-constructible-complexes (blueprint/chap-Constructible-complexes/sec-Integral-constructible-complexes/proposition-local-structure-integral-constructible-complexes/_section.md)
\phantomsection\label{bp223}%
\begin{bpstatementit}{Proposition~3.10}{ (Local structure of integral constructible complexes)}
Let \(X\) be Noetherian and let \(K\in D(X_{\mathrm{pro}\acute{e}t},\mathcal O_E)\) be derived complete with \(K_1\in\Dbc(X,\mathbf F_E)\). Then there is a finite stratification of \(X\) such that on each stratum \(S\), pro-étale locally,

\[
K|_S\simeq \mathcal O_{E,S}\otimes_{\mathcal O_E}P^\bullet
\]

for a bounded complex \(P^\bullet\) of finite free \(\mathcal O_E\)-modules.
\end{bpstatementit}

% proof: proof/proposition-local-structure-integral-constructible-complexes
\begin{proof}[{Proof}]\phantomsection\label{bp224}%
We first trivialize each \(K_n\) locally, separately for each \(n\), and then make the trivializations compatible. To begin, choose a finite stratification on which all cohomology sheaves of \(K_1\) are locally constant, and refine it by connected components; there are finitely many of these, since the strata are Noetherian.

For every stratum \(S\) the restriction \(K|_S\) is again derived complete with constructible first reduction. Indeed, restriction to an open subscheme preserves limits, and so does restriction to a closed subscheme on the pro-étale site (Proposition \hyperref[bp181]{2.58}). Moreover, both restrictions are exact and commute with the two-term cones defining the reductions. From now on we fix a connected stratum \(S\).

\emph{Local trivializations at each finite level.} By Lemma \hyperref[bp208]{3.5}, each \(K_n|_S\) is classical and of finite Tor amplitude, with locally constant finite cohomology, so we may work on \(S_{\acute{e}t}\). Fix a geometric point \(\bar s\). Then the stalk complex at \(\bar s\) has a minimal finite free model \(M_n\) in the common interval \([a,b]\). By Lemma \hyperref[bp218]{3.8} the complex \(K_n|_S\) is étale locally isomorphic to the constant complex \((M_n)_S\).

On stalks at \(\bar s\), the reduction maps \(K_{n+1}\to K_n\) induce maps \(M_{n+1}\to M_n\) inducing isomorphisms \(M_{n+1}\otimes^{\mathbf L}_{\mathcal O_E/\pi_E^{n+1}}\mathcal O_E/\pi_E^n\simeq M_n\), because the étale stalk functor is exact and symmetric monoidal. By Lemma \hyperref[bp211]{3.6} there is a bounded finite free complex \(P^\bullet\) with \(P^\bullet/\pi_E^n\simeq M_n\). Back on \(S_{\mathrm{pro}\acute{e}t}\), we get for each \(n\) local isomorphisms \((P^\bullet/\pi_E^n)_S\simeq K_n|_S\) in \(D(S_{\mathrm{pro}\acute{e}t},\mathcal O_E/\pi_E^n)\), chosen separately for each \(n\).

\emph{The integral isomorphism.} Choose an affine weakly contractible cover \(W\to S\). For each \(n\), the local isomorphism becomes an isomorphism \((P^\bullet/\pi_E^n)_W\simeq K_n|_W\) on all of \(W\), since by weak contractibility a covering by trivializing charts splits. Moreover, the restriction \(K|_W\) is derived complete, and its reductions are the \(K_n|_W\). By Lemma \hyperref[bp221]{3.9} we conclude that \(K|_W\simeq\mathcal O_{E,W}\otimes_{\mathcal O_E}P^\bullet\).
\end{proof}

% text: prose-03-13
\phantomsection\label{bp225}%
The local structure is the hard half of the following criterion.

% object: proposition-integral-constructibility-criterion (blueprint/chap-Constructible-complexes/sec-Integral-constructible-complexes/proposition-integral-constructibility-criterion/_section.md)
\phantomsection\label{bp226}%
\begin{bpstatementit}{Proposition~3.11}{ (The integral constructibility criterion)}
Let \(X\) be Noetherian. The following conditions on \(K\in D(X_{\mathrm{pro}\acute{e}t},\mathcal O_E)\) are equivalent:

\begin{enumerate}
\item \(K\in\Dbc(X,\mathcal O_E)\);
\item \(K\) is derived complete and \(K_1\in\Dbc(X,\mathbf F_E)\);
\item there is a finite stratification of \(X\) such that on each stratum \(S\), pro-étale locally, \(K|_S\simeq\mathcal O_{E,S}\otimes_{\mathcal O_E}P^\bullet\) for a bounded complex \(P^\bullet\) of finite free \(\mathcal O_E\)-modules;
\item \(K\in D^b_{\mathrm{ctf}}(X,\mathcal O_E)\).
\end{enumerate}
\end{bpstatementit}

% text: prose-03-14
\phantomsection\label{bp227}%
In the classical adic setting, this description of the integral constructible category is [\hyperref[bp705]{BS15}, Proposition 6.8.11] and [Stacks, \href{https://stacks.math.columbia.edu/tag/09C7}{09C7}].

% proof: proof/proposition-integral-constructibility-criterion
\begin{proof}[{Proof of Proposition 3.11}]\phantomsection\label{bp228}%
The substantial implication is (2)\(\Rightarrow\)(3), which is Proposition \hyperref[bp223]{3.10}. The others hold because the complexes \(\mathcal O_{E,S}\otimes_{\mathcal O_E}P^\bullet\) are complete, have finite Tor amplitude, and have constructible reductions. These properties survive extension by zero from strata. We prove (1)\(\Rightarrow\)(2)\(\Rightarrow\)(3)\(\Rightarrow\)(1) and (3)\(\Leftrightarrow\)(4).

\emph{(1) implies (2).} We first treat a single constructible sheaf, locally on a stratum, then extend by zero from the strata, and finally pass to complexes.

A finitely generated \(\mathcal O_E\)-module is a finite sum of copies of \(\mathcal O_E\) and of modules \(\mathcal O_E/\pi_E^r\), so its sheaf has a two-term finite free resolution. The coefficient sheaf \(\mathcal O_{E,U_i}\) on a member \(U_i\) of the cover in Definition definition-constructible-sheaves is derived complete. Indeed, it is the limit of its finite quotients by Lemma \hyperref[bp158]{2.51}, and since the transition maps between these quotients are surjective, this limit agrees with the derived limit by Proposition \hyperref[bp71]{2.24}. As completeness is preserved by finite sums and cones, the modules \(\mathcal O_{E,U_i}\otimes_{\mathcal O_E}P_i\) of Definition definition-constructible-sheaves are complete, and their reductions are bounded locally constant finite complexes. We may check these assertions on a pro-étale cover, because restriction to an object is a localization of the topos, and it preserves limits and the cones defining the reductions.

For a constructible sheaf, choose a finite stratification as in the definition. Then for a stratum \(a:S\hookrightarrow X\), factored as \(a=i\circ j\) as in Section \hyperref[bp195]{3.1}, the open immersion \(j\) is quasi-compact because \(X\) is Noetherian. Extension by zero \(a_!=i_*j_!\) preserves completeness, because it is exact, commutes with the cones defining the reductions, and preserves limits. Here, for \(j_!\) the last point is the assertion about limits in Proposition \hyperref[bp181]{2.58}, see [Stacks, \href{https://stacks.math.columbia.edu/tag/09BP}{09BP}], while for \(i_*\) it holds because \(i_*\) is exact and a right adjoint. Using the finite filtration by open and closed pieces, we see that the whole sheaf is complete with constructible finite reductions. Applying this to the finitely many cohomology sheaves of \(K\) and using the bounded Postnikov triangles, we obtain (2).

\emph{(2) implies (3).} This is Proposition \hyperref[bp223]{3.10}.

\emph{(3) implies (1).} Each of the finitely many strata is quasi-compact, so the local isomorphisms of (3) can be taken on a finite covering by quasi-compact pro-étale objects, and their bounds give a common cohomological interval. Since the map from the constant sheaf with value the discrete ring \(\mathcal O_E\) to \(\mathcal O_{E,S}\) is flat [Stacks, \href{https://stacks.math.columbia.edu/tag/099L}{099L}], the cohomology of \(\mathcal O_{E,S}\otimes_{\mathcal O_E}P^\bullet\) is \(\mathcal O_{E,S}\otimes_{\mathcal O_E}H^q(P^\bullet)\). Here \(H^q(P^\bullet)\) is finitely generated over \(\mathcal O_E\), so these sheaves satisfy the condition of Definition definition-constructible-sheaves. This proves (1).

\emph{(3) is equivalent to (4).} The complexes \(\mathcal O_{E,S}\otimes_{\mathcal O_E}P^\bullet\) have finite Tor amplitude, and since the stratification is finite there is a common interval. Since the restrictions to the strata are exact and jointly conservative, they detect vanishing after tensoring with any sheaf of modules, and (3) implies (4). Conversely, (4) implies (1) by definition of \(D^b_{\mathrm{ctf}}\), and hence (3) by the implications already proved.
\end{proof}

% text: prose-03-15
\phantomsection\label{bp229}%
Lemma \hyperref[bp274]{3.25} compares \(D^b_{\mathrm{ctf}}(X,\mathcal O_E)\) with the classical category of constructible \(\mathcal O_E\)-complexes [Stacks, \href{https://stacks.math.columbia.edu/tag/09C1}{09C1}], which is defined with modules over the constant sheaf with value the discrete ring \(\mathcal O_E\).


% object: sec-Finite-Tor-amplitude (blueprint/chap-Constructible-complexes/sec-Finite-Tor-amplitude/_section.md), content hidden
\subsection*{\texorpdfstring{3.4\quad Finite Tor amplitude}{3.4 Finite Tor amplitude}}
\phantomsection\label{bp230}%
\addcontentsline{toc}{subsection}{\texorpdfstring{3.4\quad Finite Tor amplitude}{3.4 Finite Tor amplitude}}
% text: prose-03-16
\phantomsection\label{bp231}%
For the integral and rational rings the finite Tor amplitude condition is automatic.

% object: proposition-ctf-equals-bc (blueprint/chap-Constructible-complexes/sec-Finite-Tor-amplitude/proposition-ctf-equals-bc/_section.md)
\phantomsection\label{bp232}%
\begin{bpstatementit}{Proposition~3.12}{ (Constructible complexes have finite Tor amplitude for $\ell$-adic coefficients)}
For \(\Lambda=\mathbf Z_\ell,\mathcal O_E,\mathbf Q_\ell,E,\overline{\mathbf Q}_\ell\), we have

\[
D^b_{\mathrm{ctf}}(X,\Lambda)=\Dbc(X,\Lambda).
\]
\end{bpstatementit}

% proof: proof/proposition-ctf-equals-bc
\begin{proof}[{Proof}]\phantomsection\label{bp233}%
For \(\mathcal O_E\), and hence for \(\mathbf Z_\ell\), this is the equivalence of (1) and (4) in the criterion above. For the fields we show that every constructible sheaf is flat, which bounds the Tor amplitude of a bounded complex by its cohomological amplitude.

For a field \(\Lambda(*)\), a constructible sheaf \(M\) is, on each stratum and locally on a cover \(\{U_i\}\), of the form \(\Lambda_{U_i}\otimes_{\Lambda(*)}P_i\) with \(P_i\) a finite-dimensional vector space, so it is locally a finite free \(\Lambda_{U_i}\)-module. Restriction to the strata is exact and the restrictions are jointly conservative, as in the proof of the criterion. Hence every constructible sheaf has Tor amplitude in \([0,0]\). Finally, a bounded complex whose finitely many cohomology sheaves have finite Tor amplitude has finite Tor amplitude, by the triangles of its canonical filtration.
\end{proof}

% text: prose-03-17
\phantomsection\label{bp234}%
For a finite ring that is not regular the condition is needed, as the \(\mathbf Z/\ell^2\)-module \(\mathbf F_\ell\) of Section \hyperref[bp197]{3.2} shows. Moreover, finite Tor amplitude is a condition on the whole complex, not on its cohomology sheaves. The category \(D^b_{\mathrm{ctf}}\) is thick, but in general it is not stable under standard truncations. For example, for \(A=\mathbf Z/\ell^2\) the complex \([A\xrightarrow{\ell}A]\) is perfect, whereas its nonzero cohomology modules are \(\mathbf F_\ell\), which are not perfect. For this reason, the compactness proofs for finite \(A\) work with bounded flat constructible models, and standard truncations are taken in \(\Dbc\). The models exist by the next lemma.

% object: lemma-flat-models-generators (blueprint/chap-Constructible-complexes/sec-Finite-Tor-amplitude/lemma-flat-models-generators/_section.md)
\phantomsection\label{bp235}%
\begin{bpstatementit}{Lemma~3.13}{ (Flat models and generators)}
Every \(C\in D^b_{\mathrm{ctf}}(X,A)\) on \(X_{\mathrm{pro}\acute{e}t}\) is represented by a bounded complex of flat constructible \(A\)-sheaves. A flat constructible sheaf has a finite stratification on whose strata it is locally of the form \(\underline M\) with \(M\) a finite projective \(A\)-module. Consequently, \(D^b_{\mathrm{ctf}}(X,A)\) is generated, using shifts, finite extensions and direct summands, by sheaves \(a_!L\), where \(a:S\hookrightarrow X\) is locally closed and \(L\) is locally of the form \(\underline M\) with \(M\) a finite projective \(A\)-module.
\end{bpstatementit}

% proof: proof/lemma-flat-models-generators
\begin{proof}[{Proof}]\phantomsection\label{bp236}%
The flat model exists on the étale site by a classical result, and we pull it back. The generators then come from two filtrations: the stupid filtration of the model, and the filtration of each term by strata.

By Corollary \hyperref[bp132]{2.44}, \(C=\nu_X^*C_{\acute{e}t}\) with \(C_{\acute{e}t}\in D^b_{\mathrm{ctf}}(X_{\acute{e}t},A)\), where \(\nu_X^*:D(X_{\acute{e}t},A)\to D(X_{\mathrm{pro}\acute{e}t},A)\) is the comparison inverse image. By [Stacks, \href{https://stacks.math.columbia.edu/tag/0F4M}{0F4M}] the complex \(C_{\acute{e}t}\) has a bounded flat model, and since \(\nu_X^*\) is exact on sheaves, its pullback is a bounded flat constructible model on \(X_{\mathrm{pro}\acute{e}t}\). We choose a common stratification on which its finitely many terms are locally constant, and their stalks are then finite projective, because a finite flat module over the Noetherian ring \(A\) is finite projective. If \(F\) is a flat constructible sheaf, then the exact sequences for the successive open and closed strata express \(F\) as a finite extension of the sheaves \(a_!L\) attached to its strata. Similarly, the stupid filtration of the bounded flat complex expresses \(C\) as a finite extension of shifts of its \emph{terms}, and shifts and finite extensions preserve finite Tor amplitude. The same argument proves the corresponding statement on \(X_{\acute{e}t}\), which we use in the proof of Proposition \hyperref[bp451]{5.14} before passing to the pro-étale site.
\end{proof}


% object: sec-Field-coefficients (blueprint/chap-Constructible-complexes/sec-Field-coefficients/_section.md), content hidden
\subsection*{\texorpdfstring{3.5\quad Field coefficients}{3.5 Field coefficients}}
\phantomsection\label{bp237}%
\addcontentsline{toc}{subsection}{\texorpdfstring{3.5\quad Field coefficients}{3.5 Field coefficients}}
% text: prose-03-18
\phantomsection\label{bp238}%
Until the end of this section, \(\Lambda=E\) or \(\overline{\mathbf Q}_\ell\), so that \(D^b_{\mathrm{ctf}}=\Dbc\) and we can argue with constructible cohomology sheaves and standard truncations. Recall (after Definition \hyperref[bp201]{3.2}) that a \(\Lambda_X\)-module is constructible exactly when it is lisse, that is locally free of finite rank for the pro-étale topology, on each stratum of a finite stratification of \(X\). By dévissage we reduce statements about \(\Dbc(X,\Lambda)\) to the simplest constructible sheaves, namely lisse sheaves extended by zero from a stratum, the sheaves \(a_!L\) of Section \hyperref[bp195]{3.1}. This dévissage first splits a complex into its cohomology sheaves by truncation, and then filters each constructible sheaf by strata.

For this dévissage we need constructible sheaves to form an abelian category (Lemma \hyperref[bp248]{3.18}). In the proof we reduce a morphism of lisse \(E\)-sheaves to a morphism of integral lattices, and for \(\overline{\mathbf Q}_\ell\) we first descend the morphism to a finite \(E\). To prepare for this, in the next two lemmas we construct the lattices locally and carry out the descent.

% object: definition-lattice (blueprint/chap-Constructible-complexes/sec-Field-coefficients/definition-lattice/_section.md)
\phantomsection\label{bp239}%
\begin{bpstatement}{Definition~3.14}{ (Lattices)}
Let \(E/\mathbf Q_\ell\) be finite and let \(L\) be a lisse \(E_X\)-module. A \emph{lattice} in \(L\) is a lisse \(\mathcal O_{E,X}\)-submodule \(M\subset L\) with \(M[1/\ell]=L\).
\end{bpstatement}

% object: lemma-lattices-lisse-e-sheaves (blueprint/chap-Constructible-complexes/sec-Field-coefficients/lemma-lattices-lisse-e-sheaves/_section.md)
\phantomsection\label{bp240}%
\begin{bpstatementit}{Lemma~3.15}{ (Lattices in lisse $E$-sheaves)}
Let \(E/\mathbf Q_\ell\) be finite and let \(L\) be a lisse \(E_X\)-module. The sheaf on \(X_{\mathrm{pro}\acute{e}t}\) of lattices in \(L\) is represented by a separated surjective étale \(X\)-scheme, with a universal lattice. In particular, \(L\) admits a lattice étale locally on \(X\).
\end{bpstatementit}

% proof: proof/lemma-lattices-lisse-e-sheaves
\begin{proof}[{Proof}]\phantomsection\label{bp241}%
Where \(L\) is free of rank \(r\), a lattice is the same as a lattice in \(E^r\), so the sheaf of lattices is locally the constant sheaf on the discrete set \(\GL_r(E)/\GL_r(\mathcal O_E)\). The representing scheme is obtained by descending the corresponding disjoint unions of copies of the trivializing objects.

The rank of \(L\) is a locally constant function on \(|X|\). Indeed, on a trivializing covering the subsets of constant rank are open, closed and saturated, and they descend by fpqc descent of open subsets. So \(X\) is the disjoint union of the open and closed subsets on which the rank is constant. Since the sheaf of lattices restricts to the sheaf of lattices on each of them, we may assume that \(L\) has constant rank \(r\).

We now construct the representing scheme. On a pro-étale affine \(V\) where \(L\) is free of rank \(r\), the sheaf of lattices is the constant sheaf on the set

\[
\GL_r(E)/\GL_r(\mathcal O_E)
\]

of lattices in \(E^r\), which is discrete because \(\GL_r(\mathcal O_E)\) is an open subgroup of \(\GL_r(E)\). To glue these constant sheaves, recall from Lemma \hyperref[bp158]{2.51} that on an affine object, sections of \(E_X\) are continuous functions to \(E\), with bounded denominators by quasi-compactness. A change of basis on \(V\times_XV\) is therefore a continuous function to \(\GL_r(E)\). Since the stabilizer of a lattice is open, the image of each fixed lattice under this function is locally constant on \(V\times_XV\). Now consider the disjoint union of copies of \(V\times_XV\) indexed by the lattices in \(E^r\). On each component, the preimages of the target components are then open and closed and form a disjoint open cover, so the induced map from this disjoint union to itself is a morphism of schemes. In this way the cocycle of the transition matrices defines a descent datum of schemes for these disjoint unions of copies of \(V\).

For affine \(V=\operatorname{Spec}B\), the disjoint union of copies of \(V\) indexed by a set \(I\) is ind-quasi-affine, since its finite subunions are affine open subschemes and together they cover it. Since descent along fpqc coverings is effective for ind-quasi-affine morphisms [Stacks, \href{https://stacks.math.columbia.edu/tag/0APK}{0APK}], the descent datum above is effective. Taking a finite affine trivializing covering of each affine open of \(X\) and gluing, we therefore obtain a scheme representing the sheaf of lattices. This scheme is separated and étale over \(X\) by [Stacks, \href{https://stacks.math.columbia.edu/tag/02KU}{02KU}, \href{https://stacks.math.columbia.edu/tag/02VN}{02VN}], and it is surjective because every free module has a lattice. By descent of sheaves of modules, this scheme carries a universal lattice.
\end{proof}

% text: prose-03-19
\phantomsection\label{bp242}%
For example, in rank one the lattices in \(E\) are the \(\pi_E^m\mathcal O_E\) for \(m\in\mathbf Z\), so the discrete set of lattices is \(E^\times/\mathcal O_E^\times\simeq\mathbf Z\). Hence for \(L=E_X\) the representing scheme is \(\coprod_{m\in\mathbf Z}X\).

% object: example-lattices-nodal-curve (blueprint/chap-Constructible-complexes/sec-Field-coefficients/example-lattices-nodal-curve/_section.md)
\phantomsection\label{bp243}%
\begin{bpstatement}{Example~3.16}{ (The scheme of lattices on the nodal curve)}
Let \(E=\mathbf Q_\ell\) and let \(L\) be the sheaf on the nodal curve \(X\) from the start of this chapter, obtained from \(E_Y\) on the chain of lines \(Y\to X\) by letting the generator of \(\mathbf Z\) act by \(\ell\). Over \(Y\) the lattices in \(L|_Y=E_Y\) are the constant lattices \(\ell^m\mathbf Z_\ell\), so over \(Y\) the scheme of lattices is \(\coprod_{m\in\mathbf Z}Y\). The generator of \(\mathbf Z\) moves the copy indexed by \(m\) to the copy indexed by \(m+1\) while translating \(Y\). Hence the scheme of lattices over \(X\), which is the quotient by \(\mathbf Z\), is a single copy of \(Y\). It is surjective and étale over \(X\), as the lemma says, but it has no section. Indeed, a section would trivialize the \(\mathbf Z\)-torsor \(Y\to X\), whereas \(Y\) is connected and \(\coprod_{\mathbf Z}X\) is not. This is the absence of a global lattice seen geometrically, and the universal lattice lives on \(Y\), where \(L\) becomes constant.
\end{bpstatement}

% text: prose-03-20
\phantomsection\label{bp244}%
With lattices we pass from \(E\) to \(\mathcal O_E\). For \(\overline{\mathbf Q}_\ell\) we first need to come down to a finite extension \(E\), and on a quasi-compact scheme this only involves finitely many matrix entries.

% object: lemma-descent-lisse-qbar-sheaves-finite (blueprint/chap-Constructible-complexes/sec-Field-coefficients/lemma-descent-lisse-qbar-sheaves-finite/_section.md)
\phantomsection\label{bp245}%
\begin{bpstatementit}{Lemma~3.17}{ (Descent of lisse $\Qbar$-sheaves)}
A finite collection of lisse \(\QbarX\)-modules and morphisms between them is obtained by extension of scalars from lisse \(E_X\)-modules and their morphisms, for one finite extension \(E/\mathbf Q_\ell\).
\end{bpstatementit}

% proof: proof/lemma-descent-lisse-qbar-sheaves-finite
\begin{proof}[{Proof}]\phantomsection\label{bp246}%
The data are described by finitely many matrices with entries in \(\QbarX\) over finitely many quasi-compact objects, and each entry already lies in some \(E_X\). We make this precise.

Since \(X\) is quasi-compact, we may choose a finite covering by quasi-compact pro-étale affines on which all the modules are free; such refinements exist by [Stacks, \href{https://stacks.math.columbia.edu/tag/0988}{0988}]. Their pairwise intersections \(U_i\times_XU_j\) are quasi-compact since \(X\) is quasi-separated, and quasi-separated since they are subschemes of the affine schemes \(U_i\times U_j\). On a quasi-compact quasi-separated object, taking sections commutes with filtered colimits, because coverings admit finite quasi-compact refinements and the sheaf condition involves only finite limits [Stacks, \href{https://stacks.math.columbia.edu/tag/0739}{0739}]. By Lemma \hyperref[bp158]{2.51} we have \(\QbarX=\varinjlim_EE_X\). Hence the finitely many transition matrices, their inverses, and the matrices of the given morphisms have entries in \(E_X\) for a single finite \(E\). Moreover, their cocycle and compatibility identities already hold over \(E_X\), since \(E_X\to\QbarX\) is injective. Descent of sheaves of modules now gives the required lisse modules and morphisms.
\end{proof}

% text: prose-03-21
\phantomsection\label{bp247}%
With these two lemmas we can show that constructible sheaves with field coefficients form an abelian category, and obtain the filtration by extensions by zero of lisse sheaves that the dévissage uses.

% object: lemma-constructible-sheaves (blueprint/chap-Constructible-complexes/sec-Field-coefficients/lemma-constructible-sheaves/_section.md)
\phantomsection\label{bp248}%
\begin{bpstatementit}{Lemma~3.18}{ (Abelian category of constructible sheaves)}
Let \(\Lambda=E\) or \(\overline{\mathbf Q}_\ell\).

\begin{enumerate}
\item \(\Cons(X,\Lambda)\) is closed under kernels, cokernels and extensions in \(\operatorname{Mod}(X_{\mathrm{pro}\acute{e}t},\Lambda)\). In particular it is an abelian category, and its inclusion into \(\operatorname{Mod}(X_{\mathrm{pro}\acute{e}t},\Lambda)\) is exact.
\item \(\Cons(X,\Lambda)\) is essentially small.
\item Every \(F\in\Cons(X,\Lambda)\) has a finite filtration \(0=F_0\subset F_1\subset\cdots\subset F_m=F\) whose subquotients \(F_k/F_{k-1}\) are of the form \(a_!L\), with \(a:S\hookrightarrow X\) a locally closed immersion and \(L\) a lisse \(\Lambda_S\)-module.
\end{enumerate}
\end{bpstatementit}

% text: prose-03-22
\phantomsection\label{bp249}%
If \(X\) is a curve, part (3) amounts to an exact sequence \(0\to j_!L\to F\to i_*i^*F\to0\), where \(j\) is the inclusion of a dense open subset on which \(F\) is lisse, \(L=j^*F\), and \(i\) is the inclusion of the finitely many closed points of the complement.

% proof: proof/lemma-constructible-sheaves
\begin{proof}[{Proof of Lemma 3.18}]\phantomsection\label{bp250}%
Part (1) is the substantial one. All three closure properties are local on a common stratification, so it suffices to treat lisse sheaves on a single stratum, and there the only difficulty is the kernel of a morphism. Parts (2) and (3), by contrast, are formal.

(1) Given finitely many constructible sheaves, we refine their stratifications to a common one. Then all three closure properties can be checked on a single stratum \(S\).

\emph{Extensions.} Let \(0\to F'\to F\to F''\to0\) be exact with \(F'\) and \(F''\) lisse on \(S\). After passing to a cover, \(F'\) and \(F''\) are free. Since \(F\to F''\) is surjective, after a further cover a basis of \(F''\) lifts to sections of \(F\). Hence the sequence splits locally, and \(F\) is locally free of finite rank.

\emph{Kernels and cokernels.} Locally, a morphism of lisse sheaves is a matrix of continuous functions to \(E\), and a priori its kernel need not be locally free. Once we choose lattices, however, the morphism becomes a constant matrix on a suitable pro-étale cover, because its reductions are maps of finite locally constant sheaves. The kernel and cokernel of a constant matrix over \(E\) are constant of finite rank, and since being lisse is a local property for the pro-étale topology, this suffices.

First let \(\Lambda=E\) and let \(f:L\to L'\) be a morphism of lisse sheaves on \(S\). By Lemma \hyperref[bp240]{3.15} we may replace \(S\) by a connected quasi-compact affine étale neighbourhood \(V\) on which \(L\) and \(L'\) contain lattices \(M\) and \(N\). On a finite quasi-compact pro-étale covering trivializing both lattices, the matrix entries of \(f\) have a common denominator, so \(h=\pi_E^mf\) maps \(M\) into \(N\) for some \(m\ge0\).

The reductions \(M_n=M/\pi_E^nM\) and \(N_n=N/\pi_E^nN\) are finite locally constant étale sheaves, and \(M=\varprojlim M_n\), \(N=\varprojlim N_n\) by Lemma \hyperref[bp158]{2.51}, checked where \(M\) and \(N\) are free. Fix a geometric point \(\bar v\) of \(V\). Then the completed fibres \(M_{\bar v}=\varprojlim(M_n)_{\bar v}\) and \(N_{\bar v}=\varprojlim(N_n)_{\bar v}\) are finite free \(\mathcal O_E\)-modules with continuous \(\pi_1(V,\bar v)\)-actions, and the reductions of \(h\) give an equivariant map \(h_{\bar v}\) between them. To make \(h\) constant, note that each \(M_n\) and \(N_n\) becomes constant on a connected finite étale cover of the connected scheme \(V\), and we may choose these covers to form a tower. Their inverse limit \(V'\to V\) is a connected pro-étale covering, so maps between constant sheaves on \(V'\) are constant. On \(V'\) the lattices \(M\) and \(N\) are identified with the constant sheaves attached to \(M_{\bar v}\) and \(N_{\bar v}\), and \(h\) with the constant matrix \(h_{\bar v}\). Hence on \(V'\) the kernel and the cokernel of \(f=\pi_E^{-m}h\) are the constant sheaves attached to the kernel and the cokernel of \(h_{\bar v}\otimes_{\mathcal O_E}E\). In particular they are lisse.

For \(\Lambda=\overline{\mathbf Q}_\ell\), we descend \(L\), \(L'\) and \(f\) to a finite extension \(E/\mathbf Q_\ell\) by Lemma \hyperref[bp245]{3.17}. Since the functor \(\QbarX\otimes_{E_X}-\) is exact, it takes the kernel and cokernel over \(E\) to those of \(f\). Finally, restriction to strata is exact, so the kernel and cokernel of a morphism of constructible sheaves on \(X\) are constructible.

(2) We choose the sites in a fixed universe. Up to isomorphism, a constructible sheaf is then determined by a finite stratification and, on each stratum, descent data of finite rank on a covering, and these data form a set.

(3) By Noetherian induction we refine the stratification of \(F\) so that there is an open stratum \(j:U\hookrightarrow X\) with closed complement \(i:Z\hookrightarrow X\) that is again a union of strata. The exact sequence \(0\to j_!j^*F\to F\to i_*i^*F\to0\) separates the stratum \(U\) from the rest, and here \(j^*F\) is lisse. Applying the same step to \(i^*F\) on \(Z\) and iterating, we obtain a finite filtration of \(F\) whose subquotients are \(a_!(F|_S)\) for the strata \(a:S\hookrightarrow X\), and each \(F|_S\) is lisse.
\end{proof}

% object: example-running-example (blueprint/chap-Constructible-complexes/sec-Field-coefficients/example-running-example/_section.md)
\phantomsection\label{bp251}%
\begin{bpstatement}{Example~3.19}{ (A running example)}
We use the same few schemes throughout:

\[
\{0\}\xrightarrow{\ i\ }\mathbf A^1_k\xleftarrow{\ j\ }\mathbf G_{m,k},
\qquad \mathbf A^1_k\xrightarrow{\ j_\infty\ }\mathbf P^1_k .
\]

The sheaf \(j_!\Lambda_{\mathbf G_m}\) is constructible for the stratification \(\mathbf A^1=\mathbf G_m\sqcup\{0\}\). Its stalk is \(\Lambda\) at a geometric point of \(\mathbf G_m\) and \(0\) at the origin, so it is not lisse on \(\mathbf A^1\), and the sequence of part (3) is the localization sequence

\[
0\to j_!\Lambda_{\mathbf G_m}\to\Lambda_{\mathbf A^1}\to i_*\Lambda\to0 .
\]

We return to this example for compact supports (after Lemma \hyperref[bp526]{6.4}), for \(i^!\) (Example \hyperref[bp591]{6.25}) and for the trace (Section \bpout{9.6}).
\end{bpstatement}

% text: prose-03-23
\phantomsection\label{bp252}%
\textbf{Exercises.} (1) Compute the stalks of \(a_!L\) for a locally closed immersion \(a\) and a lisse \(L\). (2) For \(K=\Lambda_X[1]\), in which degree is the nonzero cohomology sheaf? (3) Show that \(j_!\Lambda_{\mathbf G_m}\) has no nonzero global sections on \(\mathbf A^1_k\).

For complexes we get the following. Since constructible sheaves form an abelian subcategory, complexes with constructible cohomology are stable under cones, and the filtration of part (3) together with the truncation triangles generates them from the sheaves \(a_!L\).

% object: corollary-coefficient-heart-b (blueprint/chap-Constructible-complexes/sec-Field-coefficients/corollary-coefficient-heart-b/_section.md)
\phantomsection\label{bp253}%
\begin{bpstatementit}{Corollary~3.20}{ (Generators of $D^b_c$ with field coefficients)}
Let \(\Lambda=E\) or \(\overline{\mathbf Q}_\ell\). Then \(\Dbc(X,\Lambda)\) is a triangulated subcategory of \(D(X_{\mathrm{pro}\acute{e}t},\Lambda)\), and every object of \(\Dbc(X,\Lambda)\) is obtained by finitely many cones from objects \(a_!L[n]\), where \(a:S\hookrightarrow X\) is a locally closed immersion, \(L\) is a lisse \(\Lambda_S\)-module and \(n\in\mathbf Z\).
\end{bpstatementit}

% proof: proof/corollary-coefficient-heart-b
\begin{proof}[{Proof}]\phantomsection\label{bp254}%
The first assertion is closure under cones, which we check on cohomology sheaves. Let \(K\to K'\to K''\to K[1]\) be a distinguished triangle with \(K,K'\) in \(\Dbc(X,\Lambda)\). By the long exact cohomology sequence, each \(\mathcal H^i(K'')\) is an extension of \(\ker(\mathcal H^{i+1}K\to\mathcal H^{i+1}K')\) by \(\coker(\mathcal H^iK\to\mathcal H^iK')\). Both are constructible by Lemma lemma-constructible-sheaves, and hence so is \(\mathcal H^i(K'')\). Closure under shifts is clear, since shifts only renumber the cohomology sheaves.

For the second assertion, we argue by induction on the number of nonzero cohomology sheaves of \(K\). If \(i\) is the largest degree with \(\mathcal H^i(K)\ne0\), the truncation triangle \(\tau_{<i}K\to K\to\mathcal H^i(K)[-i]\to\) exhibits \(K\) as the cone of a map \(\mathcal H^i(K)[-i-1]\to\tau_{<i}K\), and \(\tau_{<i}K\) has fewer nonzero cohomology sheaves. This reduces the claim to a single shifted constructible sheaf, and for it the filtration from Lemma lemma-constructible-sheaves yields finitely many short exact sequences, that is, distinguished triangles, with terms \(a_!L\).
\end{proof}


% object: sec-Continuity-of-Ext (blueprint/chap-Constructible-complexes/sec-Continuity-of-Ext/_section.md), content hidden
\subsection*{\texorpdfstring{3.6\quad Continuity of Ext}{3.6 Continuity of Ext}}
\phantomsection\label{bp255}%
\addcontentsline{toc}{subsection}{\texorpdfstring{3.6\quad Continuity of Ext}{3.6 Continuity of Ext}}
% text: prose-03-24
\phantomsection\label{bp256}%
A morphism of \(\overline{\mathbf Q}_\ell\)-complexes should be defined over some finite extension \(E\) of \(\mathbf Q_\ell\). Since \(\QbarX=\varinjlim_EE_X\), this amounts to asking that Hom out of a constructible complex commute with a filtered colimit, where a filtered colimit of complexes means the termwise colimit of a filtered diagram of complexes of sheaves. We prove this continuity for complexes with a common lower bound. In this form it is the pro-étale analogue of the classical fact that cohomology of a quasi-compact quasi-separated scheme commutes with filtered colimits. Besides the descent to a finite coefficient field below, continuity is used in two places. The first is the uniform bound for Ext groups, which is proved first for integral lattices, so we need continuity for \(\mathcal O_E\) as well. The second is the compactness argument for \(E\) and \(\overline{\mathbf Q}_\ell\) (Section \hyperref[bp415]{5.2}).

For cohomology and higher direct images the statement is Corollary \hyperref[bp86]{2.29}, which holds for every sheaf of rings. To reach Hom, recall that a constructible sheaf has a finite filtration by sheaves \(a_!L\) with \(L\) lisse on a stratum \(S\) (Lemma lemma-constructible-sheaves). Using sections with support in a closed subscheme, we turn Hom out of \(a_!L\) into cohomology on \(S\) of a complex built from the target by \(Rj_*\) and restrictions, and the comparison then goes through a hypercohomology spectral sequence. Because of the common lower bound, only finitely many of its terms contribute to a given degree, and each of these terms commutes with the colimit.

% object: lemma-sections-support-closed-subscheme (blueprint/chap-Constructible-complexes/sec-Continuity-of-Ext/lemma-sections-support-closed-subscheme/_section.md)
\phantomsection\label{bp257}%
\begin{bpstatementit}{Lemma~3.21}{ (Sections with closed support)}
Let \(\Lambda\) be admissible, let \(i:Z\hookrightarrow X\) be a closed immersion with open complement \(j:V\hookrightarrow X\), and let \(R\underline\Gamma_Z:D(X_{\mathrm{pro}\acute{e}t},\Lambda)\to D(Z_{\mathrm{pro}\acute{e}t},\Lambda)\) be the local cohomology functor (Definition \hyperref[bp184]{2.59}),

\[
R\underline\Gamma_Z K=i^*\fib(K\longrightarrow Rj_*j^*K).
\]

For a sheaf \(M\) of \(\Lambda_X\)-modules on \(X_{\mathrm{pro}\acute{e}t}\), let \(\underline\Gamma_Z(M)=i^*\ker(M\to j_*j^*M)\), the sheaf on \(Z\) of sections of \(M\) with support in \(Z\).

\begin{enumerate}
\item \(R\underline\Gamma_Z\) is the right derived functor of \(\underline\Gamma_Z\). Consequently the adjunction of Lemma \hyperref[bp185]{2.60} holds in the form \(R\!\operatorname{Hom}_X(i_*N,K)=R\!\operatorname{Hom}_Z(N,R\underline\Gamma_ZK)\).
\item Let \(a=i\circ u:S\hookrightarrow X\) with \(u:S\hookrightarrow Z\) an open immersion, and let \(L\) be a finite locally projective \(\Lambda_S\)-module with dual \(L^\vee\). Then naturally in \(K\in D(X_{\mathrm{pro}\acute{e}t},\Lambda)\) \[
   R\!\operatorname{Hom}_X(a_!L,K)
   =R\Gamma\bigl(S,L^\vee\otimes u^*R\underline\Gamma_ZK\bigr).
   \tag{3.1}
   \]
\end{enumerate}
\end{bpstatementit}

% proof: proof/lemma-sections-support-closed-subscheme
\begin{proof}[{Proof}]\phantomsection\label{bp258}%
Lemma \hyperref[bp185]{2.60} gives the adjunction between \(i_*\) and \(R\underline\Gamma_Z\) on derived categories. What is new here is that \(R\underline\Gamma_Z\) is the right derived functor of \(\underline\Gamma_Z\), and this upgrades the adjunction to \(R\!\operatorname{Hom}\) complexes and yields formula (3.1). For this we show two things: \(\underline\Gamma_Z\) is right adjoint to the exact functor \(i_*\) on sheaves, and \(M\to j_*j^*M\) is surjective for injective \(M\). The second means that on injective sheaves the kernel of this map is also its fibre.

We use that \(i^*\) and \(i_*\) are exact, that \(i^*i_*=\mathrm{id}\), and that every sheaf \(M\) sits in the localization sequence \(0\to j_!j^*M\to M\to i_*i^*M\to0\) [Stacks, \href{https://stacks.math.columbia.edu/tag/09AA}{09AA}, \href{https://stacks.math.columbia.edu/tag/09AH}{09AH}, \href{https://stacks.math.columbia.edu/tag/09BM}{09BM}]. The last fact shows that a sheaf whose restriction to \(V\) vanishes is of the form \(i_*\) of a sheaf on \(Z\). In particular \(\ker(M\to j_*j^*M)=i_*\underline\Gamma_Z(M)\).

\emph{Adjunction on sheaves.} The functor \(\underline\Gamma_Z\) is right adjoint to \(i_*:\operatorname{Mod}(Z_{\mathrm{pro}\acute{e}t},\Lambda)\to\operatorname{Mod}(X_{\mathrm{pro}\acute{e}t},\Lambda)\). Indeed, a map \(i_*N\to M\) composes to zero with \(M\to j_*j^*M\), because \(j^*i_*N=0\), so it factors through \(\ker(M\to j_*j^*M)=i_*\underline\Gamma_Z(M)\). Since \(i_*\) is fully faithful, such factorizations are the maps \(N\to\underline\Gamma_Z(M)\). Moreover, as \(i_*\) is exact, its right adjoint \(\underline\Gamma_Z\) preserves K-injective complexes.

\emph{Surjectivity on injectives.} For an injective sheaf \(I\) the unit \(I\to j_*j^*I\) of the adjunction \(j^*\dashv j_*\) on \(\operatorname{Mod}(X_{\mathrm{pro}\acute{e}t},\Lambda)\) is surjective, because it has a right inverse. To construct it, extend the inclusion \(j_!j^*I\to I\) along \(j_!j^*I\hookrightarrow j_*j^*I\) to a map \(s:j_*j^*I\to I\), which is possible because \(I\) is injective. Then the composite \(j_*j^*I\xrightarrow{s}I\to j_*j^*I\) restricts to the identity on \(V\). Since \(j_*\) is fully faithful, an endomorphism of \(j_*j^*I\) is determined by its restriction to \(V\), and so this composite is the identity.

\emph{Comparison with the fibre.} The right derived functor of \(\underline\Gamma_Z\) agrees with \(i^*\fib(K\to Rj_*j^*K)\). To see this, represent \(K\) by a K-injective complex \(I^\bullet\) with injective terms [Stacks, \href{https://stacks.math.columbia.edu/tag/079P}{079P}]. Then \(\underline\Gamma_Z(I^\bullet)\) computes the right derived functor of \(\underline\Gamma_Z\) at \(K\). On the other side, the complex \(j^*I^\bullet\) is K-injective, because the left adjoint \(j_!\) of \(j^*\) is exact. Hence \(j_*j^*I^\bullet\) represents \(Rj_*j^*K\). By the surjectivity just proved, the sequences

\[
0\to i_*\underline\Gamma_Z(I^n)\to I^n\to j_*j^*I^n\to0
\]

are exact termwise, and they define a distinguished triangle \(i_*\underline\Gamma_Z(I^\bullet)\to K\to Rj_*j^*K\to i_*\underline\Gamma_Z(I^\bullet)[1]\). Applying \(i^*\) and using \(i^*i_*=\mathrm{id}\) identifies \(\underline\Gamma_Z(I^\bullet)\) with \(i^*\fib(K\to Rj_*j^*K)\).

\emph{Part (1).} With \(I^\bullet\) as above, the complex \(\underline\Gamma_Z(I^\bullet)\) is K-injective by the adjunction on sheaves, and it represents \(R\underline\Gamma_ZK\) by the comparison with the fibre. The adjunction on complexes therefore gives

\[
R\!\operatorname{Hom}_X(i_*N,K)=\Hom^\bullet(i_*N,I^\bullet)
=\Hom^\bullet(N,\underline\Gamma_Z(I^\bullet)),
\]

whose right-hand side is \(R\!\operatorname{Hom}_Z(N,R\underline\Gamma_ZK)\).

\emph{Formula (3.1).} Since \(a_!=i_*u_!\), part (1), the adjunction \(u_!\dashv u^*\) and the identity \(\cHom(L,-)=L^\vee\otimes-\) for the finite locally projective sheaf \(L\) give \(R\!\operatorname{Hom}_X(i_*u_!L,K)=R\!\operatorname{Hom}_Z(u_!L,R\underline\Gamma_ZK)
=R\Gamma(S,L^\vee\otimes u^*R\underline\Gamma_ZK)\).
\end{proof}

% text: prose-03-25
\phantomsection\label{bp259}%
Formula (3.1) expresses Hom out of \(a_!L\) through cohomology on the stratum \(S\) and the functors \(i^*\), \(u^*\) and \(Rj_*\). Each of these commutes with filtered colimits of uniformly bounded below complexes, and from this we deduce the continuity of Ext.

% object: lemma-continuity-ext-out-constructible-complex (blueprint/chap-Constructible-complexes/sec-Continuity-of-Ext/lemma-continuity-ext-out-constructible-complex/_section.md)
\phantomsection\label{bp260}%
\begin{bpstatementit}{Lemma~3.22}{ (Continuity of Ext)}
Let \((K_\alpha^\bullet)\) be a filtered diagram of complexes of \(\Lambda_X\)-modules that all vanish below a fixed degree, and let \(M\) be one of the following:

\begin{enumerate}
\item an object of \(\Dbc(X,\Lambda)\), where \(\Lambda=E\) or \(\overline{\mathbf Q}_\ell\);
\item a finite extension of sheaves \(a_!L_0\), with \(a:S\hookrightarrow X\) locally closed and \(L_0\) a finite locally free \(\mathcal O_{E,S}\)-module, where \(\Lambda=\mathcal O_E\).
\end{enumerate}

Then for every integer \(n\) the natural map

\[
\varinjlim_\alpha\Hom(M,K_\alpha^\bullet[n])
\longrightarrow
\Hom\bigl(M,(\varinjlim_\alpha K_\alpha^\bullet)[n]\bigr)
\]

is an isomorphism.
\end{bpstatementit}

% text: prose-03-26
\phantomsection\label{bp261}%
Applied to sheaves, which are complexes in degree \(0\), the lemma says that for \(M\in\Cons(X,\Lambda)\), \(\Lambda=E\) or \(\overline{\mathbf Q}_\ell\), the functors \(\Ext^n(M,-)\) commute with filtered colimits of \(\Lambda_X\)-modules.

% proof: proof/lemma-continuity-ext-out-constructible-complex
\begin{proof}[{Proof of Lemma 3.22}]\phantomsection\label{bp262}%
We reduce to a single generator \(a_!L\) and then use (3.1) to turn Hom into cohomology on the stratum \(S\), where Corollary \hyperref[bp86]{2.29} applies. The common lower bound is used at the end, to compare the abutments of two spectral sequences.

In both cases \(M\) is obtained by finitely many extensions and shifts from sheaves \(a_!L\), with \(a:S\hookrightarrow X\) a locally closed immersion and \(L\) a finite locally free \(\Lambda_S\)-module: in case (1) by Corollary \hyperref[bp253]{3.20}, and in case (2) by assumption. Both sides are cohomological functors of \(M\). By the five lemma, the class of objects \(M\) for which the claim holds is therefore closed under extensions and shifts, so it suffices to treat \(M=a_!L\), with \(a=i\circ
u:S\hookrightarrow Z\hookrightarrow X\) as in Lemma \hyperref[bp257]{3.21}. The only property of \(L\) used below is that \(L^\vee\otimes-\) is exact.

By (3.1), we must show that the cohomology of \(R\Gamma(S,L^\vee\otimes u^*R\underline\Gamma_ZK)\) commutes with the colimit, where \(K=\varinjlim_\alpha K_\alpha^\bullet\). To this end we apply the exact functor \(L^\vee\otimes u^*i^*(-):D(X_{\mathrm{pro}\acute{e}t},\Lambda)\to D(S_{\mathrm{pro}\acute{e}t},\Lambda)\) to the triangle \(i_*R\underline\Gamma_ZK\to K\to Rj_*j^*K\to\) and to the corresponding triangles for the \(K_\alpha^\bullet\). By the long exact sequence and the five lemma, it suffices to treat the two terms

\[
R\Gamma\bigl(S,L^\vee\otimes u^*i^*K\bigr)
\quad\text{and}\quad
R\Gamma\bigl(S,L^\vee\otimes u^*i^*Rj_*j^*K\bigr).
\]

For the first term, Corollary \hyperref[bp86]{2.29} on \(S\) applies, because the functor \(L^\vee\otimes u^*i^*(-)\) is exact, commutes with colimits and preserves the lower bound. For the second term, consider the hypercohomology spectral sequence

\[
E_2^{p,r}=H^p\bigl(S,L^\vee\otimes u^*i^*R^rj_*(j^*K)\bigr)
\Longrightarrow
H^{p+r}\bigl(S,L^\vee\otimes u^*i^*Rj_*j^*K\bigr)
\]

which is functorial in \(K\). The sheaves \(R^rj_*(j^*K_\alpha)\) commute with the colimit by the second part of Corollary \hyperref[bp86]{2.29}, and \(H^p(S,-)\) does by its first part. Hence, after passing to the colimit, the map of spectral sequences is an isomorphism on \(E_2\). The abutments agree as well, since by the common lower bound only finitely many \(E_2^{p,r}\) contribute to a given total degree.
\end{proof}


% object: sec-Descent-to-a-finite-extension (blueprint/chap-Constructible-complexes/sec-Descent-to-a-finite-extension/_section.md), content hidden
\subsection*{\texorpdfstring{3.7\quad Descent to a finite extension}{3.7 Descent to a finite extension}}
\phantomsection\label{bp263}%
\addcontentsline{toc}{subsection}{\texorpdfstring{3.7\quad Descent to a finite extension}{3.7 Descent to a finite extension}}
% text: prose-03-27
\phantomsection\label{bp264}%
The lattice arguments below need the discrete valuation ring \(\mathcal O_E\) with its uniformizer \(\pi_E\), but the valuation ring of \(\overline{\mathbf Q}_\ell\) is not of this kind. Since \(\overline{\mathbf Q}_\ell\) is the filtered union of the finite extensions \(E\) of \(\mathbf Q_\ell\), we show instead that every object and every morphism of \(\Dbc(X,\overline{\mathbf Q}_\ell)\) is defined over some finite \(E\).

For \(C\in D(X_{\mathrm{pro}\acute{e}t},E)\) and a finite extension \(E'/E\), the extensions of scalars are the complexes \(C_{E'}=E'_X\otimes_{E_X}C\) and \(C_{\overline{\mathbf Q}_\ell}=\QbarX\otimes_{E_X}C\). These tensor products need not be derived, since \(E'_X\) and \(\QbarX\) are flat over \(E_X\).

The classical counterpart is the fact that a continuous representation of a profinite group on \(\overline{\mathbf Q}_\ell^r\), with the \(\ell\)-adic topology, takes values in \(\GL_r(E)\) for some finite \(E/\mathbf Q_\ell\). Similarly, a constructible complex, or a morphism between two of them, is described by finitely many matrices with entries in \(\QbarX\), and on a quasi-compact scheme finitely many such entries already lie in a single \(E_X\). We use the next lemma whenever a statement for \(\overline{\mathbf Q}_\ell\) is reduced to one for a finite \(E\), as in Lemma lemma-integral-models and in the uniform Ext bound (Lemma \hyperref[bp426]{5.6}).

% object: lemma-descent-finite-coefficient-field (blueprint/chap-Constructible-complexes/sec-Descent-to-a-finite-extension/lemma-descent-finite-coefficient-field/_section.md)
\phantomsection\label{bp265}%
\begin{bpstatementit}{Lemma~3.23}{ (Descent to a finite coefficient field)}
\begin{enumerate}
\item For \(C,D\in\Dbc(X,E)\) and every integer \(r\), extension of scalars induces an isomorphism \[
   \varinjlim_{E'\supset E}
   \Hom_{D(X_{\mathrm{pro}\acute{e}t},E')}(C_{E'},D_{E'}[r])
   \;\xrightarrow{\ \sim\ }\;
   \Hom_{D(X_{\mathrm{pro}\acute{e}t},\overline{\mathbf Q}_\ell)}(C_{\overline{\mathbf Q}_\ell},D_{\overline{\mathbf Q}_\ell}[r]),
   \] the colimit running over finite extensions \(E'\) of \(E\) in \(\overline{\mathbf Q}_\ell\).
\item Every object of \(\Dbc(X,\overline{\mathbf Q}_\ell)\) is isomorphic to \(K_{\overline{\mathbf Q}_\ell}\) for some \(K\in\Dbc(X,E)\) and some finite \(E/\mathbf Q_\ell\).
\end{enumerate}
\end{bpstatementit}

% proof: proof/lemma-descent-finite-coefficient-field
\begin{proof}[{Proof}]\phantomsection\label{bp266}%
Morphisms descend because \(\QbarX\) is the filtered colimit of the \(E'_X\) and Hom out of a constructible complex commutes with this colimit (Lemma \hyperref[bp260]{3.22}). Objects are built from lisse sheaves by finitely many extensions, each classified by a morphism, and so they descend as well.

(1) Extension of scalars \(E'_X\otimes_{E_X}-\) is left adjoint to restriction of scalars \(D(X_{\mathrm{pro}\acute{e}t},E')\to D(X_{\mathrm{pro}\acute{e}t},E)\). Hence \(\Hom_{D(X_{\mathrm{pro}\acute{e}t},E')}(C_{E'},D_{E'}[r])=\Hom_{D(X_{\mathrm{pro}\acute{e}t},E)}(C,D_{E'}[r])\), and similarly for \(\overline{\mathbf Q}_\ell\). Represent \(D\) by a bounded complex \(D^\bullet\). Then the complexes \(D^\bullet\otimes_{E_X}E'_X\) form a filtered system of complexes of \(E_X\)-modules, all vanishing below the lower bound of \(D^\bullet\), and their colimit is \(D^\bullet\otimes_{E_X}\QbarX\), which represents \(D_{\overline{\mathbf Q}_\ell}\). Applying Lemma \hyperref[bp260]{3.22} to \(C\in\Dbc(X,E)\), we get

\[
\varinjlim_{E'}\Hom_{D(X_{\mathrm{pro}\acute{e}t},E)}(C,D_{E'}[r])
\xrightarrow{\ \sim\ }\Hom_{D(X_{\mathrm{pro}\acute{e}t},E)}(C,D_{\overline{\mathbf Q}_\ell}[r]),
\]

which is the claim after the adjunction.

(2) First let \(K\) be a constructible \(\QbarX\)-module, with filtration \(N_1\subset\cdots\subset N_m=K\) as in Lemma lemma-constructible-sheaves. We descend \(N_k\) by induction on \(k\). The subquotient \(N_k/N_{k-1}=a_!L\) descends, because \(L\) does by Lemma \hyperref[bp245]{3.17} and \(a_!\) commutes with extension of scalars. Suppose \(N_{k-1}\) and \(N_k/N_{k-1}\) descend to \(E\). Then the extension \(0\to N_{k-1}\to N_k\to N_k/N_{k-1}\to0\) gives a distinguished triangle, classified by a morphism \(N_k/N_{k-1}\to N_{k-1}[1]\). By (1) this morphism descends after a finite extension of \(E\). The fibre of the descended morphism, that is, its cone shifted by \([-1]\), is a model of \(N_k\).

For a bounded complex \(K\), we argue in the same way by induction on the number of nonzero cohomology sheaves. Let \(i\) be the largest degree with \(\mathcal H^i(K)\neq0\), and consider the truncation triangle \(\tau_{<i}K\to K\to\mathcal H^i(K)[-i]\to\). Once \(\tau_{<i}K\) and \(\mathcal H^i(K)\) have models over a finite \(E\), the connecting morphism \(\mathcal H^i(K)[-i-1]\to\tau_{<i}K\) descends after a finite extension by (1), and the cone of the descended morphism is a model of \(K\). Each step replaces \(E\) by a finite extension, so after finitely many steps \(K\) is defined over a single finite \(E\).
\end{proof}


% object: sec-Lattices-after-stratification (blueprint/chap-Constructible-complexes/sec-Lattices-after-stratification/_section.md), content hidden
\subsection*{\texorpdfstring{3.8\quad Lattices after stratification}{3.8 Lattices after stratification}}
\phantomsection\label{bp267}%
\addcontentsline{toc}{subsection}{\texorpdfstring{3.8\quad Lattices after stratification}{3.8 Lattices after stratification}}
% text: prose-03-28
\phantomsection\label{bp268}%
We pass from statements with coefficients \(\mathcal O_E/\pi_E^n\) to statements with coefficients \(E\) through integral lattices. A lisse \(E_X\)-module may fail to have a lattice on all of \(X\), but it has one on each stratum of a finite stratification (Lemma \hyperref[bp269]{3.24}), and this is enough for dévissage arguments. In particular, from the stratified statement we obtain the integral generators of Lemma lemma-integral-models. Lemma \hyperref[bp269]{3.24} is also the first step in the proofs of the uniform Ext bound and of the degree bound for proper direct image in Chapter \hyperref[bp407]{5}. On a connected normal scheme, by contrast, there is a global lattice, and lisse \(E\)-sheaves are representations of \(\pi_1(X,\bar x)\); we prove this where it is used, in Proposition \bpout{13.10}.

On the nodal curve \(X\) of Example \hyperref[bp243]{3.16}, the sheaf \(L\) has no global lattice, and one stratification step repairs this. On the open stratum \(X\smallsetminus\{\text{node}\}\simeq\mathbf G_m\) the torsor is trivial, so \(L\) is constant there and has a lattice. The node, with its reduced structure, is the spectrum of a field, which is normal, and hence the restriction of \(L\) to the node has a lattice as well.

% object: lemma-lattices-after-finite-stratification (blueprint/chap-Constructible-complexes/sec-Lattices-after-stratification/lemma-lattices-after-finite-stratification/_section.md)
\phantomsection\label{bp269}%
\begin{bpstatementit}{Lemma~3.24}{ (Lattices after a finite stratification)}
Let \(L\) be a lisse \(E_X\)-module. There is a finite stratification of \(X\) by locally closed subschemes \(S\) such that each \(L|_S\) contains a lisse \(\mathcal O_{E,S}\)-submodule \(L_0\) with \(L_0[1/\ell]=L|_S\).
\end{bpstatementit}

% text: prose-03-29
\phantomsection\label{bp270}%
In the proof we first find a lattice after a surjective étale base change and then push it down by the trace. For the first step we use that the lattices in \(E^r\) form a discrete set on which \(\GL_r(E)\) acts continuously, which makes the sheaf of lattices in \(L\) an étale \(X\)-scheme. For the second step the base change must be finite étale, which we arrange by restricting to the strata of a finite stratification.

% proof: proof/lemma-lattices-after-finite-stratification
\begin{proof}[{Proof of Lemma 3.24}]\phantomsection\label{bp271}%
Let \(Y\to X\) be the étale \(X\)-scheme representing lattices in \(L\) (Lemma \hyperref[bp240]{3.15}), with its universal lattice.

\emph{Pushing down by the trace.} Choose finitely many quasi-compact affine opens of this scheme whose images cover \(X\), and let \(p:V\to X\) be their disjoint union. Then \(p\) is quasi-compact, separated, surjective and étale, and the universal lattice gives a lattice \(M_0\subset p^*L\). By [Stacks, \href{https://stacks.math.columbia.edu/tag/07RY}{07RY}, \href{https://stacks.math.columbia.edu/tag/07RZ}{07RZ}], \(X\) has a finite stratification by locally closed subschemes over each of which \(p\) is finite étale and surjective. On such a stratum \(S\), let \(\operatorname{Tr}:p_*p^*L\to L\) be the trace of the finite étale map \(p\) over \(S\), and define \(L_0\subset L|_S\) as the image of the composite

\[
p_*M_0\longrightarrow p_*p^*L\xrightarrow{\operatorname{Tr}}L.
\]

To see that \(L_0\) is a lattice, we split the finite étale cover locally on \(S\) into \(d\) copies of \(S\). Then \(p_*M_0\) is a direct sum of \(d\) lattices \(M_1,\dots,M_d\) in \(L\), and the trace is their sum, so \(L_0=M_1+\cdots+M_d\). If we further trivialize \(L\) pro-étale locally, each \(M_k\) becomes a locally constant function to the discrete set of lattices in \(E^r\), and after a common refinement each \(M_k\) is the constant sheaf attached to a lattice \(N_k\subset E^r\). Then \(L_0\) is the constant sheaf attached to \(N_1+\cdots+N_d\), which is a finite free \(\mathcal O_E\)-lattice spanning \(E^r\). Hence \(L_0\) is lisse and \(L_0[1/\ell]=L|_S\).
\end{proof}


% object: sec-Integral-models (blueprint/chap-Constructible-complexes/sec-Integral-models/_section.md), content hidden
\subsection*{\texorpdfstring{3.9\quad Integral models}{3.9 Integral models}}
\phantomsection\label{bp272}%
\addcontentsline{toc}{subsection}{\texorpdfstring{3.9\quad Integral models}{3.9 Integral models}}
% text: prose-03-30
\phantomsection\label{bp273}%
To deduce statements over \(E\) and \(\overline{\mathbf Q}_\ell\) from statements over \(\mathcal O_E\), we need integral objects that generate \(\Dbc(X,E)\) and \(\Dbc(X,\overline{\mathbf Q}_\ell)\). We also want to know that our integral constructible complexes are those of the Stacks Project, so that results proved there apply, and we prove both in the next lemma. Recall from Proposition \hyperref[bp226]{3.11} in Section \hyperref[bp205]{3.3} that an object lies in \(\Dbc(X,\mathcal O_E)=D^b_{\mathrm{ctf}}(X,\mathcal O_E)\) precisely when it is derived complete and its reduction modulo \(\pi_E\) is a bounded constructible complex pulled back from \(X_{\acute{e}t}\). The Stacks Project works instead with the category \(\Dc(X,\mathcal O_E)\) of [Stacks, \href{https://stacks.math.columbia.edu/tag/09C1}{09C1}], built from modules over the constant sheaf \(\mathcal O_E\) rather than over \(\mathcal O_{E,X}\).

% object: lemma-integral-models (blueprint/chap-Constructible-complexes/sec-Integral-models/lemma-integral-models/_section.md)
\phantomsection\label{bp274}%
\begin{bpstatementit}{Lemma~3.25}{ (Integral models)}
\begin{enumerate}
\item Restriction of scalars along the map from the constant sheaf \(\mathcal O_E\) to \(\mathcal O_{E,X}\) is an equivalence between the derived complete objects of \(D(X_{\mathrm{pro}\acute{e}t},\mathcal O_E)\) and the derived complete complexes of modules over the constant sheaf \(\mathcal O_E\).
\item An object \(K\) of \(D(X_{\mathrm{pro}\acute{e}t},\mathcal O_E)\) lies in \(\Dbc(X,\mathcal O_E)\) if and only if its restriction of scalars lies in \(\Dc(X,\mathcal O_E)\).
\item \(\Dbc(-,\mathcal O_E)\) is preserved by \(f^*:D(Y_{\mathrm{pro}\acute{e}t},\mathcal O_E)\to D(X_{\mathrm{pro}\acute{e}t},\mathcal O_E)\) for every morphism \(f:X\to Y\), by \(\otimes^{\mathbf L}_{\mathcal O_{E,X}}\) on \(D(X_{\mathrm{pro}\acute{e}t},\mathcal O_E)\), and by \(a_!:D(S_{\mathrm{pro}\acute{e}t},\mathcal O_E)\to D(X_{\mathrm{pro}\acute{e}t},\mathcal O_E)\) for every locally closed immersion \(a:S\hookrightarrow X\).
\item Every object of \(\Dbc(X,E)\) is obtained by finitely many cones and shifts from objects \((a_!L_0)[1/\ell]\), where \(a:S\hookrightarrow X\) is a locally closed immersion and \(L_0\) is a lisse \(\mathcal O_{E,S}\)-module that is locally free of finite rank. Every object of \(\Dbc(X,\overline{\mathbf Q}_\ell)\) is obtained in the same way from the objects \(\QbarX\otimes_{\mathcal O_{E,X}}a_!L_0\), for some finite \(E/\mathbf Q_\ell\).
\end{enumerate}
\end{bpstatementit}

% text: prose-03-31
\phantomsection\label{bp275}%
In (4) each \(L_0\) satisfies the condition of Definition definition-constructible-sheaves, so it lies in \(\Cons(S,\mathcal O_E)\). By (3) the objects \(a_!L_0\) then lie in \(\Dbc(X,\mathcal O_E)\). Their reductions modulo \(\pi_E^n\) are bounded constructible complexes with coefficients in \(\mathcal O_E/\pi_E^n\), which are pulled back from \(X_{\acute{e}t}\) (Corollary \hyperref[bp132]{2.44}).

% proof: proof/lemma-integral-models
\begin{proof}[{Proof of Lemma 3.25}]\phantomsection\label{bp276}%
In (1) and (2) we compare two coefficient sheaves related by a flat map that is an isomorphism modulo every \(\pi_E^n\), so they have the same derived complete objects. For (3) we use the two characterizations of integral constructibility, while for (4) we combine the dévissage of Corollary \hyperref[bp253]{3.20} with the lattices of Lemma \hyperref[bp269]{3.24}.

(1) and (2). The map from the constant sheaf \(\mathcal O_E\) to \(\mathcal O_{E,X}\) is flat and induces isomorphisms modulo every \(\pi_E^n\) [Stacks, \href{https://stacks.math.columbia.edu/tag/099L}{099L}]. Hence by [Stacks, \href{https://stacks.math.columbia.edu/tag/099I}{099I}] restriction of scalars is an equivalence on derived complete objects, which is (1). By [Stacks, \href{https://stacks.math.columbia.edu/tag/09C1}{09C1}], an object of \(\Dc(X,\mathcal O_E)\) is a derived complete complex whose reduction modulo \(\pi_E\) is constructible of finite Tor dimension over \(\mathbf F_E\). Since over the field \(\mathbf F_E\) every bounded complex has finite Tor dimension, this description agrees with the description of \(\Dbc(X,\mathcal O_E)\) in Proposition \hyperref[bp226]{3.11}, and together with (1) this gives (2).

(3) By Proposition \hyperref[bp226]{3.11}, an object of \(D(X_{\mathrm{pro}\acute{e}t},\mathcal O_E)\) lies in \(\Dbc(X,\mathcal O_E)\) if and only if it satisfies condition (2) there: it is derived complete with constructible reduction modulo \(\pi_E\). Equivalently, it satisfies condition (3): there is a finite stratification such that on each stratum \(S\), pro-étale locally, \(K|_S\simeq\mathcal O_{E,S}\otimes_{\mathcal O_E}P^\bullet\) for a bounded complex \(P^\bullet\) of finite free \(\mathcal O_E\)-modules; here the implication from (2) to (3) is Proposition \hyperref[bp223]{3.10}. For each of the three operations we therefore check whichever of these two conditions is convenient.

\emph{Inverse image.} By Lemma \hyperref[bp167]{2.54}, \(f^*\) sends finite free \(\mathcal O_{E,Y}\)-modules to finite free \(\mathcal O_{E,X}\)-modules, so it preserves the local description of condition (3).

\emph{Extension by zero.} The functor \(a_!=i_*j_!\) preserves derived completeness. Indeed, \(j_!\) for a quasi-compact open immersion \(j\) preserves limits by [Stacks, \href{https://stacks.math.columbia.edu/tag/09BP}{09BP}], and \(i_*\) is exact and preserves limits by [Stacks, \href{https://stacks.math.columbia.edu/tag/09AA}{09AA}, \href{https://stacks.math.columbia.edu/tag/09BM}{09BM}]. Moreover, both functors commute with reduction modulo \(\pi_E\) and preserve constructibility of the reductions, so \(a_!\) preserves condition (2).

\emph{Tensor product.} Let \(K,M\in\Dbc(X,\mathcal O_E)\), and choose a stratification for \(K\) as above. Derived completeness and constructibility of the reduction modulo \(\pi_E\) are both preserved by cones, so \(\Dbc(X,\mathcal O_E)\) is closed under cones. Moreover, the localization triangles along the strata exhibit \(K\) as a finite iterated extension of the objects \(a_!a^*K\), for the strata \(a:S\hookrightarrow X\), and the functor \(-\otimes^{\mathbf L}M\) preserves these triangles. It therefore suffices to show that \((a_!a^*K)\otimes^{\mathbf L}M\) lies in \(\Dbc(X,\mathcal O_E)\).

For this we use the projection formula for a locally closed immersion: for \(L\in D(S_{\mathrm{pro}\acute{e}t},\mathcal O_E)\) and \(M\in D(X_{\mathrm{pro}\acute{e}t},\mathcal O_E)\) the canonical map

\[
(a_!L)\otimes^{\mathbf L}M\longrightarrow a_!(L\otimes^{\mathbf L}a^*M)
\]

is an isomorphism. Indeed, write \(a=i\circ j\) with \(j\) open and \(i\) closed, so that \(a_!=i_*j_!\). On \(S\) the map restricts to the identity of \(L\otimes^{\mathbf L}a^*M\), while on every complementary stratum both sides restrict to zero. Since restriction to the strata is exact and jointly conservative by [Stacks, \href{https://stacks.math.columbia.edu/tag/09AD}{09AD}], the map is an isomorphism.

Given the parts on inverse image and extension by zero, we are reduced to showing that \(a^*K\otimes^{\mathbf L}a^*M\) lies in \(\Dbc(S,\mathcal O_E)\). Tensoring with \(a^*K\) preserves derived completeness, because \(a^*K\) is pro-étale locally a bounded complex of finite free modules. Moreover, the reduction modulo \(\pi_E\) of the tensor product is the tensor product of the reductions, and this is constructible.

(4) For \(\overline{\mathbf Q}_\ell\), we first descend to a finite \(E\) by Lemma \hyperref[bp265]{3.23}. By Corollary \hyperref[bp253]{3.20}, every object of \(\Dbc(X,E)\) is obtained by finitely many cones and shifts from sheaves \(a_!L\) with \(L\) lisse. We refine the stratum \(S\) by a finite stratification as in Lemma \hyperref[bp269]{3.24}. This replaces \(a_!L\) by finitely many extensions of such sheaves, and on each new stratum \(L\) contains a lattice \(L_0\) with \(L_0[1/\ell]=L\), so that \(a_!L=(a_!L_0)[1/\ell]\). It remains to see that \(L_0\) is integrally constructible. The reduction of \(L_0\) modulo \(\pi_E^n\) is a locally constant sheaf of finitely presented \(\mathcal O_E/\pi_E^n\)-modules, so it comes from the étale site by [Stacks, \href{https://stacks.math.columbia.edu/tag/099Y}{099Y}]. Thus \(L_0\in\Dbc(S,\mathcal O_E)\), and \(a_!L_0\in\Dbc(X,\mathcal O_E)\) by (3).
\end{proof}

% text: prose-03-32
\phantomsection\label{bp277}%
On \(\mathbf A^1\), part (4) is visible in the running example. There the constant sheaf \(E_{\mathbf A^1}\) is the cone of a map \(i_*E[-1]\to j_!E_{\mathbf G_m}\), and up to shift both terms are of the form \((a_!L_0)[1/\ell]\) with \(L_0\) the constant lattice \(\mathcal O_E\) on the stratum.


% object: sec-Inverting-ell (blueprint/chap-Constructible-complexes/sec-Inverting-ell/_section.md), content hidden
\subsection*{\texorpdfstring{3.10\quad Inverting $\ell$}{3.10 Inverting ell}}
\phantomsection\label{bp278}%
\addcontentsline{toc}{subsection}{\texorpdfstring{3.10\quad Inverting $\ell$}{3.10 Inverting ell}}
% text: prose-03-33
\phantomsection\label{bp279}%
Passing from \(\mathcal O_E\) to \(E\) means inverting \(\ell\), and on derived categories this is a Verdier localization.

% object: proposition-inverting-ell-full-derived-categories (blueprint/chap-Constructible-complexes/sec-Inverting-ell/proposition-inverting-ell-full-derived-categories/_section.md)
\phantomsection\label{bp280}%
\begin{bpstatementit}{Proposition~3.26}{ (Inverting $\ell$ on derived categories)}
Extension of scalars \(E_X\otimes_{\mathcal O_{E,X}}-:
D(X_{\mathrm{pro}\acute{e}t},\mathcal O_E)\to D(X_{\mathrm{pro}\acute{e}t},E)\) identifies \(D(X_{\mathrm{pro}\acute{e}t},E)\) with the Verdier quotient of \(D(X_{\mathrm{pro}\acute{e}t},\mathcal O_E)\) by the complexes \(K\) with \(H^q(K)[1/\ell]=0\) for all \(q\).
\end{bpstatementit}

% proof: proof/proposition-inverting-ell-full-derived-categories
\begin{proof}[{Proof}]\phantomsection\label{bp281}%
A functor with a fully faithful right adjoint is the Verdier localization at its kernel, so we show that restriction of scalars is fully faithful and compute the kernel.

Inverting \(\ell\) is a filtered colimit along multiplication by \(\ell\), hence exact. Therefore \(E_X=\mathcal O_{E,X}[1/\ell]\) is flat over \(\mathcal O_{E,X}\), and \(H^q(E_X\otimes K)=H^q(K)[1/\ell]\). Restriction of scalars is right adjoint to extension of scalars, and the counit is an isomorphism, because \(\ell\) already acts invertibly on an \(E_X\)-module. Restriction is therefore fully faithful, and extension of scalars is the Verdier localization at its kernel. By the formula for \(H^q(E_X\otimes K)\), this kernel consists of the complexes \(K\) with \(H^q(K)[1/\ell]=0\) for all \(q\).
\end{proof}

% text: prose-03-34
\phantomsection\label{bp282}%
The kernel contains more than the complexes killed by a power of \(\ell\), for example \(\bigoplus_{n\ge1}\mathcal O_{E,X}/\ell^n\). It does so even among constructible sheaves: on \(\bigsqcup_{n\ge1}\operatorname{Spec}k\), the sheaf equal to \(\mathcal O_E/\ell^n\) on the \(n\)-th point is killed by inverting \(\ell\) but by no power of \(\ell\). When the underlying space of \(X\) is Noetherian, however, this cannot happen in \(\Dbc(X,\mathcal O_E)\), and there inverting \(\ell\) amounts to tensoring the Hom groups with \(E\).

% object: proposition-inverting-ell (blueprint/chap-Constructible-complexes/sec-Inverting-ell/proposition-inverting-ell/_section.md)
\phantomsection\label{bp283}%
\begin{bpstatementit}{Proposition~3.27}{ (Inverting $\ell$)}
Assume that the underlying topological space of \(X\) is Noetherian (for instance, that \(X\) is of finite type over a field). A complex in \(\Dbc(X,\mathcal O_E)\) that is killed by inverting \(\ell\) is killed by a fixed power of \(\ell\), and extension of scalars induces an equivalence

\[
\Dbc(X,\mathcal O_E)[\ell^{-1}]\xrightarrow{\ \sim\ }\Dbc(X,E),
\]

where the source category has the objects of \(\Dbc(X,\mathcal O_E)\) and Hom groups tensored with \(E\).
\end{bpstatementit}

% proof: proof/proposition-inverting-ell
\begin{proof}[{Proof}]\phantomsection\label{bp284}%
Both assertions are [\hyperref[bp705]{BS15}, Proposition 6.8.14(2)]. We prove the first, which is where the Noetherian hypothesis enters. Let \(K\in\Dbc(X,\mathcal O_E)\). By Definition \hyperref[bp199]{3.1}, there is a finite stratification such that the restriction of each of the finitely many cohomology sheaves of \(K\) to each stratum \(S\) is, on quasi-compact pro-étale charts \(U\), of the form \(\mathcal O_{E,U}\otimes_{\mathcal O_E}P\) with \(P\) finitely generated. Suppose that \(K[1/\ell]=0\). Extension of scalars to \(E_X\) is exact, so \(E_U\otimes_{\mathcal O_E}P=0\), and each such \(P\) is torsion. The underlying space of \(X\) is Noetherian, so each stratum is quasi-compact and finitely many charts suffice. With finitely many \(P\), one power of \(\ell\) kills the restrictions of all cohomology sheaves of \(K\) to all strata. By the localization sequences of Proposition \hyperref[bp181]{2.58}, a larger power kills every cohomology sheaf of \(K\). Since \(K\) is bounded, a further power kills \(K\) itself, by induction along its truncation triangles.

Full faithfulness says that for \(K,L\in\Dbc(X,\mathcal O_E)\) every morphism \(E_X\otimes K\to E_X\otimes L\) is \(\ell^{-c}\) times the image of a morphism \(K\to L\) for some \(c\geq0\), and that a morphism \(K\to L\) becomes zero after extension of scalars only if some power of \(\ell\) kills it. For essential surjectivity, recall from Lemma lemma-integral-models that \(\Dbc(X,E)\) is generated by objects coming from \(\Dbc(X,\mathcal O_E)\). By the description of the Hom groups, a connecting morphism between them lifts to an integral one after multiplication by a power of \(\ell\). Hence finite cones of the generators also have integral models.
\end{proof}



% object: chap-Proper-base-change (blueprint/chap-Proper-base-change/_section.md), content hidden
\section*{\texorpdfstring{4\quad Proper base change}{4 Proper base change}}
\phantomsection\label{bp285}%
\addcontentsline{toc}{section}{\texorpdfstring{4\quad Proper base change}{4 Proper base change}}
% text: prose-04-01
\phantomsection\label{bp286}%
The function–sheaf dictionary (Section \bpout{1.3}, Proposition \bpout{12.28}) turns the identity \(\#X_0(\mathbf F_q)=\sum_y\#f^{-1}(y)(\mathbf F_q)\) into a statement about sheaves, and for this the trace function of \(Rf_!K\) at a point \(y\) must be the sum of \(t_K\) over the fibre \(f^{-1}(y)\). This requires the stalks of \(Rf_!K\) to be the cohomology of the fibres, but nothing in the definition of the derived direct image says so.

% object: question-stalks-proper-direct-image (blueprint/chap-Proper-base-change/question-stalks-proper-direct-image/_section.md)
\phantomsection\label{bp287}%
\begin{bpstatement}{Question~4.1}{ (Stalks of higher direct images along a proper map)}
For \(f\) proper and \(F\) a torsion sheaf on \(X\), is the stalk \((R^qf_*F)_{\bar y}\) equal to \(H^q(X_{\bar y},F)\)?
\end{bpstatement}

% text: prose-04-02
\phantomsection\label{bp288}%
For a proper map \(f:X\to Y\) of locally compact Hausdorff spaces the answer is yes. Indeed, the sets \(f^{-1}(V)\), for \(V\) a neighbourhood of \(y\), form a basis of neighbourhoods of the compact fibre, and the cohomology of a compact subset is the colimit of the cohomology of its neighbourhoods.

In the étale topology, however, the stalk of \(R^qf_*F\) at a geometric point \(\bar y\) is the cohomology of \(X\times_YY_{(\bar y)}\), where \(Y_{(\bar y)}\) is the spectrum of the strict henselization of \(Y\) at \(\bar y\). This scheme is much larger than the fibre \(X_{\bar y}\). For instance, if \(f\) is the identity of a curve and \(\bar y\) lies over a closed point, then it is the spectrum of a strictly henselian local ring of dimension one, not a point.

Properness cannot be dropped. For example, let \(j:\mathbf G_m\hookrightarrow\mathbf A^1\) be the inclusion, \(i:\{0\}\hookrightarrow\mathbf A^1\) the origin, and \(A\) the constant sheaf on \(\mathbf G_m\) with value a nonzero finite abelian group. Then the sheaf \(i^*j_*A\) is nonzero, while the base change of \(\mathbf G_m\) to the origin is empty.

For proper \(f\), nevertheless, the answer is yes. More generally, pulling a proper family back to a new base and taking its cohomology can be done in either order. Precisely, for a cartesian square with \(f\) proper, \(g:Y'\to Y\) arbitrary and \(X'=X\times_YY'\), the base change map

\[
g^*Rf_*K\longrightarrow Rf'_*g'^*K
\]

is an isomorphism for every bounded below classical complex \(K\) with torsion cohomology sheaves (Theorem \hyperref[bp382]{4.33}). In particular, taking for \(g\) a geometric point gives the stalk formula (Corollary \hyperref[bp386]{4.34}). The theorem is due to Artin [\hyperref[bp702]{SGA4}, XII–XIII], and unlike smooth base change (Chapter \hyperref[bp637]{8}), it holds for torsion of every order, including torsion divisible by the residue characteristics.

The content of the theorem is already visible in degree zero, where it is a statement about connectedness. Indeed, sections of the constant sheaf \(\{0,1\}\) are open and closed subsets, and for a proper scheme over a henselian local ring, the open and closed subsets correspond to those of the closed fibre. When \(X\) is finite over a henselian local ring \(R\), this says that idempotents of the finite \(R\)-algebra lift from its reduction modulo the maximal ideal, which is one form of Hensel's lemma. Theorem \hyperref[bp308]{4.8} extends the statement to sections of every étale sheaf of sets. Higher degrees then reduce to degree zero and the projective line, and from \(\mathbf P^1\) the general case follows by formal steps, through \((\mathbf P^1)^n\), \(\mathbf P^n\) and Chow's lemma (Proposition \hyperref[bp378]{4.32}).

We make these calculations on the étale site, where the classical tools are available, and transport them to classical complexes on the pro-étale site by the comparison theorems of Chapter \hyperref[bp3]{2}. For \(\mathcal O_E\) we pass to the limit over the reductions modulo \(\pi_E^n\) (Proposition \hyperref[bp397]{4.38}). Since pullback need not preserve derived completeness, at this stage the statement for \(\mathcal O_E\) involves the completed inverse image. Once we know that \(Rp_*\) preserves constructibility (Proposition \hyperref[bp405]{4.40}), however, the completion is unnecessary, and we obtain the ordinary base change isomorphism on \(D^b_{\mathrm{ctf}}\) for every admissible coefficient ring.

We use proper base change three times later. First, it computes the stalks of \(Rf_!\) in the trace formula (Corollary \bpout{12.26}). Second, it makes \(Rf_!=Rp_*j_!\) independent of the compactification (Proposition \hyperref[bp492]{5.26}). Finally, it bounds the Tor amplitude of \(Rp_*\) with finite coefficients (Proposition \hyperref[bp390]{4.35}).

% text: situation-conventions
\phantomsection\label{bp289}%
\textbf{Convention.} For proper base change, schemes are arbitrary unless an additional hypothesis is stated. The theorem concerns the classical part of \(D^+(T_{\mathrm{pro}\acute{e}t},\Lambda)\), for a discrete commutative ring \(\Lambda\), but the intermediate calculations (acyclicity, curves, Chow's lemma) are made on \(T_{\acute{e}t}\), with abelian sheaves or modules over its discrete constant ring. There we write \(D^+(T_{\acute{e}t},\Lambda)\) for the bounded below derived category. Inverse image of étale sheaves along \(u:S\to T\) is written \(u^*:\Sh(T_{\acute{e}t})\to\Sh(S_{\acute{e}t})\), as for module sheaves; it is exact, and so is its extension to \(D^+(T_{\acute{e}t},\Lambda)\to D^+(S_{\acute{e}t},\Lambda)\). The derived direct image is written \(Ru_*\). Torsion always means torsion of the underlying abelian sheaves.

Proper base change holds for torsion coefficients of every characteristic. Accordingly, until Section \hyperref[bp388]{4.7}, \(\ell\) denotes an arbitrary prime, which may equal a residue characteristic. From that section on, we return to the standing finite-type hypotheses and the standing condition on \(\ell\).

% object: sec-The-base-change-map (blueprint/chap-Proper-base-change/sec-The-base-change-map/_section.md), content hidden
\subsection*{\texorpdfstring{4.1\quad The base change map}{4.1 The base change map}}
\phantomsection\label{bp290}%
\addcontentsline{toc}{subsection}{\texorpdfstring{4.1\quad The base change map}{4.1 The base change map}}
% text: prose-04-03
\phantomsection\label{bp291}%
Proper base change asserts that a specific map, from the pullback of the direct image to the direct image of the pullback, is an isomorphism. We define this map and describe it on sections, so that we can compute it on stalks.

From Definition \hyperref[bp322]{4.14} on, we fix the coefficient ring. Integral and rational coefficients are treated in Section \hyperref[bp392]{4.8}. Results for abelian sheaves apply to sheaves of modules, by the following lemma.

% object: lemma-cohomology-modules-underlying-abelian (blueprint/chap-Proper-base-change/sec-The-base-change-map/lemma-cohomology-modules-underlying-abelian/_section.md)
\phantomsection\label{bp292}%
\begin{bpstatementit}{Lemma~4.2}{ (Modules and abelian sheaves)}
For a sheaf of \(\Lambda\)-modules on a site, its cohomology and its higher direct images along a morphism of sites agree with those of its underlying abelian sheaf. The same holds for \(Ru_*\) of bounded below complexes.
\end{bpstatementit}

% proof: proof/lemma-cohomology-modules-underlying-abelian
\begin{proof}[{Proof}]\phantomsection\label{bp293}%
An injective sheaf of modules is acyclic for the section functor on every object of the site, although it need not be injective as an abelian sheaf [Stacks, \href{https://stacks.math.columbia.edu/tag/03FD}{03FD}]. Hence an injective resolution of modules computes the cohomology of the underlying abelian sheaf, on every object. Applying this to the objects over the target gives the statement for \(R^qu_*\), and bounded below resolutions give it for \(Ru_*\).
\end{proof}

% text: prose-04-04
\phantomsection\label{bp294}%
Consider a cartesian square

\[
\begin{CD}
X' @>{g'}>> X\\
@V{f'}VV @VV{f}V\\
Y' @>{g}>> Y.
\end{CD}
\tag{4.1}
\]

Since \(gf'=fg'\), the composition isomorphisms of inverse images yield an isomorphism \(f'^*g^*\simeq g'^*f^*\) on sheaves. Inverse images are exact, so the same isomorphism exists on bounded below derived categories.

% object: definition-underived-base-change-map (blueprint/chap-Proper-base-change/sec-The-base-change-map/definition-underived-base-change-map/_section.md)
\phantomsection\label{bp295}%
\begin{bpstatement}{Definition~4.3}{ (Underived base change map)}
Let \(F\) be a sheaf on \(X_{\acute{e}t}\). The \emph{underived base change map}

\[
g^*f_*F\longrightarrow f'_*g'^*F
\tag{4.2}
\]

is the adjoint, under \(f'^*\dashv f'_*\), of

\[
f'^*g^*f_*F
\simeq g'^*f^*f_*F
\xrightarrow{g'^*\varepsilon_f}g'^*F,
\]

where \(\varepsilon_f:f^*f_*\to\mathrm{id}\) is the counit.
\end{bpstatement}

% text: prose-04-05
\phantomsection\label{bp296}%
Equivalently, first apply the unit \(g^*f_*F\to f'_*f'^*g^*f_*F\), and then apply \(f'_*\) to the displayed composite. This construction is natural in \(F\), and it applies equally to sheaves of \(\Lambda\)-modules.

On sections the map has the following simple description. A section of \(g^*f_*F\) over an étale \(V\to Y'\) is represented by an étale \(W\to Y\), a map \(V\to
Y'\times_YW\), and a section over \(X\times_YW\). Given such data, the base change map restricts this section along

\[
X'\times_{Y'}V\longrightarrow X\times_YW,
\]

and it is the sheafification of this map of presheaves.

% object: definition-derived-base-change-map (blueprint/chap-Proper-base-change/sec-The-base-change-map/definition-derived-base-change-map/_section.md)
\phantomsection\label{bp297}%
\begin{bpstatement}{Definition~4.4}{ (Derived base change map)}
For \(K\in D^+(X_{\acute{e}t},\Lambda)\), the \emph{base change map}

\[
g^*Rf_*K\longrightarrow Rf'_*g'^*K
\tag{4.3}
\]

is the adjoint, under \(f'^*\dashv Rf'_*\), of

\[
f'^*g^*Rf_*K
\simeq g'^*f^*Rf_*K
\xrightarrow{g'^*\varepsilon_f}g'^*K,
\]

where \(\varepsilon_f:f^*Rf_*\to\mathrm{id}\) now denotes the derived counit. The morphism is natural in \(K\).
\end{bpstatement}

% text: prose-04-06
\phantomsection\label{bp298}%
The derived adjunction exists because inverse image is exact and its right adjoint preserves injectives. Note that the construction of the base change map uses only that the square commutes; the cartesian hypothesis is needed only when we ask whether the map is an isomorphism.

To compute (4.3), choose a bounded below injective resolution \(K\to I^\bullet\) and an injective replacement \(g'^*I^\bullet\to J^\bullet\). Then the base change map is represented by the morphism of complexes

\[
g^*f_*I^\bullet
\longrightarrow
f'_*g'^*I^\bullet\longrightarrow f'_*J^\bullet.
\tag{4.4}
\]

Here the first arrow is the underived base change map in each degree. By exactness of \(g^*\) the source is \(g^*Rf_*K\), and the last complex computes \(Rf'_*g'^*K\). The formula does not depend on the choices, since comparison maps between injective resolutions are unique up to homotopy. Moreover, it is the adjoint map of Definition \hyperref[bp297]{4.4}, because under adjunction the first arrow corresponds to the square isomorphism followed by the counit, and the replacement map is the canonical passage to the derived direct image.

Base change maps for composite squares are composites of base change maps.

% object: lemma-identity-pasting (blueprint/chap-Proper-base-change/sec-The-base-change-map/lemma-identity-pasting/_section.md)
\phantomsection\label{bp299}%
\begin{bpstatementit}{Lemma~4.5}{ (Identity and pasting)}
Let \(K\in D^+(X_{\acute{e}t},\Lambda)\). The base change map is the identity for an identity square, and is compatible with horizontal and vertical pasting of commutative squares, in particular of cartesian squares. More explicitly, if \(h:Y''\to Y'\) is a second base change, completed by maps \(h':X''\to X'\) and \(f'':X''\to Y''\) to a commutative square with \(f'h'=hf''\), the base change map for the pasted square, with bottom map \(gh\), is the composite

\[
h^*g^*Rf_*K
\longrightarrow
h^*Rf'_*g'^*K
\longrightarrow
Rf''_*h'^*g'^*K
\]

of \(h^*\) applied to the base change map of (4.1) and the base change map of the second square, under the canonical composition isomorphisms.

For composable maps \(X\xrightarrow{f}Y\xrightarrow{g}Z\) and a map \(h:Z'\to Z\), let \(g':Y'\to Z'\), \(h_Y:Y'\to Y\) and \(f':X'\to Y'\), \(h_X:X'\to X\) complete \(g\), \(h\) and \(f\), \(h_Y\) to commutative squares (for instance \(Y'=Y\times_ZZ'\) and \(X'=X\times_ZZ'\)). Then the base change map for the pasted square with left map \(g'f'\) and right map \(gf\) is the composite

\[
h^*Rg_*Rf_*K
\longrightarrow
Rg'_*h_Y^*Rf_*K
\longrightarrow
Rg'_*Rf'_*h_X^*K
\]

of the base change map for \(g\) at \(Rf_*K\) and \(Rg'_*\) applied to the base change map for \(f\) at \(K\).
\end{bpstatementit}

% proof: proof/lemma-identity-pasting
\begin{proof}[{Proof}]\phantomsection\label{bp300}%
Each base change map is a unit followed by a square isomorphism and a counit. When two squares are pasted, the square isomorphisms compose, and the triangle identities cancel the inner unit and counit. In the vertical case we also use that direct image preserves injectives, so that \(R(gf)_*\simeq Rg_*Rf_*\) is computed on the same resolution.
\end{proof}

% text: prose-04-07
\phantomsection\label{bp301}%
On the pro-étale site we use the same construction, but here \(\Lambda\) may be any coefficient ring pulled back from the point, not only a torsion ring. For such \(\Lambda\) inverse images are exact (Lemma \hyperref[bp167]{2.54}), and the adjunction \(f^*\dashv Rf_*\) exists on the unbounded derived categories.

% object: definition-pro-etale-base-change-map (blueprint/chap-Proper-base-change/sec-The-base-change-map/definition-pro-etale-base-change-map/_section.md)
\phantomsection\label{bp302}%
\begin{bpstatement}{Definition~4.6}{ (Pro-étale base change map)}
Let \(K\in D(X_{\mathrm{pro}\acute{e}t},\Lambda)\). The \emph{base change map}

\[
g^*Rf_*K\longrightarrow Rf'_*g'^*K
\tag{4.5}
\]

in \(D(Y'_{\mathrm{pro}\acute{e}t},\Lambda)\) is the adjoint transpose, for \(f'^*\dashv Rf'_*\), of the composite

\[
f'^*g^*Rf_*K\simeq g'^*f^*Rf_*K
\xrightarrow{g'^*\varepsilon_f}g'^*K,
\]

where \(\varepsilon_f:f^*Rf_*\to\mathrm{id}\) is the counit of \(f^*\dashv Rf_*\).
\end{bpstatement}

% text: prose-04-08
\phantomsection\label{bp303}%
Equivalently, with \(\eta_{f'}\) the unit of \(f'^*\dashv Rf'_*\), it is the composite

\[
g^*Rf_*K\xrightarrow{\eta_{f'}}
Rf'_*f'^*g^*Rf_*K\simeq Rf'_*g'^*f^*Rf_*K
\xrightarrow{Rf'_*g'^*\varepsilon_f}Rf'_*g'^*K.
\]

This is the pro-étale form of the base change map of Definition \hyperref[bp297]{4.4}. It is defined for every cartesian square (4.1), and the same construction applies to any commutative square. Since proper and smooth base change are proved on the étale site and then transported to the pro-étale site, we need to know that the two maps correspond.

% object: lemma-naturality-pasting-comparison (blueprint/chap-Proper-base-change/sec-The-base-change-map/lemma-naturality-pasting-comparison/_section.md)
\phantomsection\label{bp304}%
\begin{bpstatementit}{Lemma~4.7}{ (Naturality, pasting and comparison)}
The maps (4.5) form a natural transformation \(g^*Rf_*\to Rf'_*g'^*\). It is the identity for an identity square, and for a horizontal or vertical pasting of two commutative squares the base change map of the pasted square is the composite of the two base change maps, in the form displayed in Lemma \hyperref[bp299]{4.5}. For a discrete ring \(\Lambda\) and \(K\in D^+(X_{\acute{e}t},\Lambda)\), with \(f\) and \(f'\) qcqs, the comparison isomorphisms of Theorem \hyperref[bp113]{2.38} identify the base change map of \(\nu_X^*K\) with the inverse image on \(Y'_{\mathrm{pro}\acute{e}t}\) of the étale base change map of \(K\).
\end{bpstatementit}

% proof: proof/lemma-naturality-pasting-comparison
\begin{proof}[{Proof}]\phantomsection\label{bp305}%
Every assertion holds because the base change map is assembled from units, counits and composition isomorphisms for inverse images, and each of the operations in the lemma respects these ingredients. Naturality in \(K\) is immediate, since the counit \(\varepsilon_f\) is a natural transformation.

By definition, the base change map is the adjoint transpose of a counit preceded by the composition isomorphism for inverse images. Since the counit for a composite adjunction is the composite of the two counits, taking adjoint transposes proves pasting, in either order of composing the two directions of a rectangle. Similarly, the assertion about identity squares follows from the triangle identities, as in Lemma \hyperref[bp299]{4.5}. Finally, the comparisons of sites of Theorem \hyperref[bp113]{2.38} identify the composition maps for inverse images and the units and counits (Lemma \hyperref[bp116]{2.39}), and hence also the base change maps constructed from them.
\end{proof}


% object: sec-Base-change-in-degree-zero (blueprint/chap-Proper-base-change/sec-Base-change-in-degree-zero/_section.md), content hidden
\subsection*{\texorpdfstring{4.2\quad Base change in degree zero}{4.2 Base change in degree zero}}
\phantomsection\label{bp306}%
\addcontentsline{toc}{subsection}{\texorpdfstring{4.2\quad Base change in degree zero}{4.2 Base change in degree zero}}
% text: prose-04-09
\phantomsection\label{bp307}%
For a proper morphism, the underived base change map is an isomorphism for every étale sheaf of sets, with no torsion hypothesis. The reason is that, for a proper scheme over a henselian pair, restriction of sections to the closed fibre is bijective (Theorem \hyperref[bp308]{4.8}).

Recall that a pair \((R,I)\), with \(I\) an ideal of the ring \(R\), is \emph{henselian} if \(I\) is contained in the Jacobson radical of \(R\) and Hensel's lemma holds for it: every factorization modulo \(I\) of a monic polynomial into coprime monic factors lifts to a factorization in \(R\) [Stacks, \href{https://stacks.math.columbia.edu/tag/09XE}{09XE}]. A local ring with its maximal ideal is a henselian pair exactly when it is henselian. For example, a complete local ring is henselian.

% object: theorem-comparison-over-henselian-pair-degree (blueprint/chap-Proper-base-change/sec-Base-change-in-degree-zero/theorem-comparison-over-henselian-pair-degree/_section.md)
\phantomsection\label{bp308}%
\begin{bpstatementit}{Theorem~4.8}{ (Henselian comparison in degree zero)}
Let \((R,J)\) be a henselian pair, let \(Z\to\operatorname{Spec}R\) be proper, and let \(i:Z_0=Z\times_R\operatorname{Spec}(R/J)\hookrightarrow Z\). For every sheaf of sets \(F\) on \(Z_{\acute{e}t}\), restriction is a bijection

\[
\Gamma(Z,F)\xrightarrow{\sim}\Gamma(Z_0,i^*F).
\]

The comparison is natural in \(F\).
\end{bpstatementit}

% text: prose-04-10
\phantomsection\label{bp309}%
The theorem is [Stacks, \href{https://stacks.math.columbia.edu/tag/0A0C}{0A0C}]. At heart it is a statement about connectedness: every irreducible closed subset of a scheme proper over a henselian pair meets the closed fibre in a nonempty connected subset. This follows from Stein factorization and the connectedness properties of henselian pairs [Stacks, \href{https://stacks.math.columbia.edu/tag/0A0B}{0A0B}]; over a general affine base, Stein factorization only produces an algebra integral over \(R\), but that suffices. A topological argument on spectral spaces then turns connectedness into bijectivity of restriction for Zariski sheaves.

Passing to étale sheaves uses finite covers, which are harmless here: a finite cover of \(Z\) remains proper over \(R\), so the connectedness statement also holds after every finite cover. Given an étale covering of \(Z\), we choose a finite surjection \(Z'\to Z\) that factors through the covering Zariski locally on \(Z'\), and sections then descend along \(Z'\times_ZZ'\rightrightarrows Z'\) [Stacks, \href{https://stacks.math.columbia.edu/tag/09ZF}{09ZF}].

% object: example-henselian-needed (blueprint/chap-Proper-base-change/sec-Base-change-in-degree-zero/example-henselian-needed/_section.md)
\phantomsection\label{bp310}%
\begin{bpstatement}{Example~4.9}{ (Why the pair must be henselian)}
Let \(k\) be a field of characteristic \(\neq2\) and let \(R=k[t]_{(t)}\), the local ring of \(\mathbf A^1_k\) at the origin, with maximal ideal \(J=(t)\). The finite \(R\)-algebra \(B=R[x]/(x^2-1-t)\) is a domain, because \(1+t\) is not a square in \(k(t)\). Hence \(Z=\operatorname{Spec}B\) is connected and \(\Gamma(Z,\mathbf Z/2)=\mathbf Z/2\). The closed fibre, on the other hand, is \(\operatorname{Spec}k[x]/(x^2-1)\), which consists of the two points \(x=\pm1\), and so \(\Gamma(Z_0,\mathbf Z/2)=(\mathbf Z/2)^2\). Thus the conclusion of the theorem fails for \((R,J)\). The reason is that \((R,J)\) is not henselian: the factorization \(x^2-1=(x-1)(x+1)\) modulo \(t\) does not lift. In the henselization, by contrast, and in \(k[[t]]\), the element \(1+t\) has the square root \(1+t/2-t^2/8+\cdots\). There the factorization lifts, and \(Z\) splits into two copies of the base, so both sides become \((\mathbf Z/2)^2\).
\end{bpstatement}

% text: prose-04-11
\phantomsection\label{bp311}%
The stalks of \(f_*F\) are sections over geometric fibres. Indeed, a stalk is the set of sections over the strict localization, and the henselian comparison identifies these with the sections over the closed fibre.

% object: lemma-sections-geometric-fibre (blueprint/chap-Proper-base-change/sec-Base-change-in-degree-zero/lemma-sections-geometric-fibre/_section.md)
\phantomsection\label{bp312}%
\begin{bpstatementit}{Lemma~4.10}{ (Sections over a geometric fibre)}
Let \(f:X\to Y\) be proper, let \(F\) be an étale sheaf of sets on \(X\), and let \(\bar y:\operatorname{Spec}\Omega\to Y\) be a geometric point. Then restriction is a bijection

\[
(f_*F)_{\bar y}\xrightarrow{\sim}
\Gamma(X_{\bar y},F|_{X_{\bar y}}).
\tag{4.6}
\]
\end{bpstatementit}

% proof: proof/lemma-sections-geometric-fibre
\begin{proof}[{Proof}]\phantomsection\label{bp313}%
Let \(y\in Y\) be the image of \(\bar y\), let \(Y_{(\bar y)}=\operatorname{Spec}\mathcal O^{\mathrm{sh}}_{Y,\bar y}\) and \(X_{(\bar y)}=X\times_YY_{(\bar y)}\). In degree zero the stalk is computed on the strict localization, that is,

\[
(f_*F)_{\bar y}\simeq\Gamma(X_{(\bar y)},F|_{X_{(\bar y)}}).
\]

For quasi-compact quasi-separated \(f\) this holds by continuity of étale sections ([Stacks, \href{https://stacks.math.columbia.edu/tag/03Q9}{03Q9}], whose proof in degree zero applies to sheaves of sets). To recall the argument: by definition, the stalk \((f_*F)_{\bar y}\) is the filtered colimit of \(\Gamma(X\times_YV,F|)\) over pointed affine étale neighbourhoods \(V\) of \(\bar y\), where \(F|\) denotes the inverse image of \(F\). Since the limit of these neighbourhoods is \(Y_{(\bar y)}\), there is a natural map from the colimit to the sections over \(X_{(\bar y)}\), and we must see that it is bijective. A section over the limit is represented by finitely many covering objects and finitely many agreement conditions, because \(X_{(\bar y)}\) is quasi-compact and quasi-separated. Since all these data are of finite presentation, they descend to one neighbourhood. This proves surjectivity, and the same descent for two representatives proves injectivity.

The local ring of \(Y_{(\bar y)}\) is strictly henselian, with separably closed residue field \(\kappa(y)^{\mathrm{sep}}\subseteq\Omega\). Thus \(\mathcal O^{\mathrm{sh}}_{Y,\bar y}\) with its maximal ideal is a henselian pair, and \(X_{(\bar y)}\) is proper over it with closed fibre \(X_{\kappa(y)^{\mathrm{sep}}}\). Theorem \hyperref[bp308]{4.8} therefore identifies the sections on the right with those over \(X_{\kappa(y)^{\mathrm{sep}}}\). These in turn equal \(\Gamma(X_{\bar y},F|_{X_{\bar y}})\), because extending a separably closed ground field from \(\kappa(y)^{\mathrm{sep}}\) to \(\Omega\) does not change the sections of the pullback of an étale sheaf [Stacks, \href{https://stacks.math.columbia.edu/tag/0A3I}{0A3I}].
\end{proof}

% object: proposition-proper-base-change-degree-zero (blueprint/chap-Proper-base-change/sec-Base-change-in-degree-zero/proposition-proper-base-change-degree-zero/_section.md)
\phantomsection\label{bp314}%
\begin{bpstatementit}{Proposition~4.11}{ (Proper base change in degree zero)}
For (4.1) with \(f\) proper, and for every étale sheaf of sets \(F\) on \(X\), the underived base change map (4.2)

\[
g^*f_*F\longrightarrow f'_*g'^*F
\]

is an isomorphism.
\end{bpstatementit}

% proof: proof/proposition-proper-base-change-degree-zero
\begin{proof}[{Proof}]\phantomsection\label{bp315}%
We show that both stalks are the sections over the same geometric fibre. To this end, choose \(\bar y':\operatorname{Spec}\Omega\to Y'\) and use the composite \(\bar y=g\bar y'\) as the geometric point of \(Y\). Since properness is preserved by base change, Lemma \hyperref[bp312]{4.10} applies to \(f'\) as well as to \(f\). Moreover, the cartesian square identifies

\[
X'_{\bar y'}=X\times_Y\operatorname{Spec}\Omega=X_{\bar y}
\]

and identifies the two inverse image sheaves on it. Thus the stalks of both sides of the desired map are the same set of sections. It remains to see that the map is the identity. The underived base change map sends a section represented over \(X\times_YV\) to its restriction to the base change of the neighbourhood, and hence to its restriction to this same fibre. In other words, the map on stalks is the identity under the displayed identifications. Since geometric stalks detect isomorphisms, this proves the assertion.
\end{proof}


% object: sec-Torsion-sheaves-and-acyclicity (blueprint/chap-Proper-base-change/sec-Torsion-sheaves-and-acyclicity/_section.md), content hidden
\subsection*{\texorpdfstring{4.3\quad Torsion sheaves and acyclicity}{4.3 Torsion sheaves and acyclicity}}
\phantomsection\label{bp316}%
\addcontentsline{toc}{subsection}{\texorpdfstring{4.3\quad Torsion sheaves and acyclicity}{4.3 Torsion sheaves and acyclicity}}
% text: prose-04-12
\phantomsection\label{bp317}%
In degree zero we needed no hypothesis on the sheaf. In higher degrees we do: proper base change fails for the constant sheaf \(\mathbf Z\).

% object: example-pbc-fails-integral (blueprint/chap-Proper-base-change/sec-Torsion-sheaves-and-acyclicity/example-pbc-fails-integral/_section.md)
\phantomsection\label{bp318}%
\begin{bpstatement}{Example~4.12}{ (Proper base change fails for $\mathbf Z$)}
Let \(R\) be a strictly henselian discrete valuation ring with uniformizer \(t\) and algebraically closed residue field of characteristic \(\neq2\), and let \(X\subset\mathbf P^2_R\) be the closure of the affine curve \(y^2=x^3-x^2+t\). Then \(X\) is regular (at the origin of the closed fibre the equation has the linear term \(t\)), while its closed fibre \(X_s\) is the nodal cubic \(y^2=x^2(x-1)\). Let \(f:X\to\operatorname{Spec}R\) be the structure map. Since \(R\) is strictly henselian, the stalk of \(R^1f_*\mathbf Z\) at the closed point is \(H^1(X,\mathbf Z)\).

This group is zero. Indeed, \(X\) is normal and connected, so \(\mathbf Z\)-torsors on \(X\) correspond to continuous homomorphisms \(\pi_1(X)\to\mathbf Z\), and the image of such a homomorphism is a compact subgroup of \(\mathbf Z\), hence zero. On the closed fibre, however, \(H^1(X_s,\mathbf Z)=\mathbf Z\), generated by the loop through the node. This means that the base change map from the stalk of \(R^1f_*\mathbf Z\) at the closed point to \(H^1(X_s,\mathbf Z)\) is not surjective.

With \(\mathbf Z/n\) in place of \(\mathbf Z\), the class of the loop is realized by the cyclic étale cover of \(X_s\) of degree \(n\) obtained by unwinding the node, and because \(R\) is henselian, this cover lifts to \(X\). This is the phenomenon the torsion hypothesis captures.
\end{bpstatement}

% text: prose-04-13
\phantomsection\label{bp319}%
We therefore restrict to torsion sheaves.

% object: definition-torsion-sheaves (blueprint/chap-Proper-base-change/sec-Torsion-sheaves-and-acyclicity/definition-torsion-sheaves/_section.md)
\phantomsection\label{bp320}%
\begin{bpstatement}{Definition~4.13}{ (Torsion sheaves)}
A sheaf \(F\) of abelian groups, or of \(\Lambda\)-modules, is \emph{torsion} if every geometric stalk is a torsion abelian group. Equivalently,

\[
F=\varinjlim_{n\geq1}F[n],\qquad F[n]=\ker(n:F\to F),
\]

where the positive integers are ordered by divisibility.
\end{bpstatement}

% text: prose-04-14
\phantomsection\label{bp321}%
If \(\Lambda\) is torsion as an abelian group, it is killed by the additive order of \(1\), so every sheaf of \(\Lambda\)-modules is torsion. The two descriptions of torsion in Definition \hyperref[bp320]{4.13} agree. Indeed, the colimit is taken in sheaves, and geometric stalks commute with colimits and finite limits and detect isomorphisms, so the equivalence reduces to the case of abelian groups.

Torsion is a local condition on sections: every section is killed by an integer after passing to an étale covering. However, the integer may vary between members of the covering. For example, the constant sheaf \(\mathbf Q/\mathbf Z\) on a nonempty scheme is torsion, but on the disjoint union of countably many copies of the scheme its sections include \((1/2^n)_{n\geq1}\), which has infinite order.

% object: definition-torsion-pbc-property (blueprint/chap-Proper-base-change/sec-Torsion-sheaves-and-acyclicity/definition-torsion-pbc-property/_section.md)
\phantomsection\label{bp322}%
\begin{bpstatement}{Definition~4.14}{ (The PBC property)}
For a morphism \(f:X\to Y\), write \(\mathrm{PBC}_{\Lambda}(f)\) if the base change map (4.3) is an isomorphism for every base change \(g:Y'\to Y\) and every \(K\in D^+(X_{\acute{e}t},\Lambda)\) with torsion cohomology sheaves.
\end{bpstatement}

% text: prose-04-15
\phantomsection\label{bp323}%
From now until the end of Section \hyperref[bp362]{4.5}, \(\Lambda\) denotes a commutative ring that is torsion as an abelian group, unless another ring is specified. For such \(\Lambda\) and proper \(f\) we prove \(\mathrm{PBC}_\Lambda(f)\) in Propositions \hyperref[bp360]{4.26} and \hyperref[bp378]{Proposition 4.32}, and Theorem \hyperref[bp382]{4.33} then removes the condition on \(\Lambda\).

To pass from degree zero to higher degrees we apply the degree-zero result to an injective resolution \(F\to I^\bullet\) of a sheaf \(F\) on \(X\). By Proposition \hyperref[bp314]{4.11} the base change map is an isomorphism \(g^*f_*I^m\simeq f'_*g'^*I^m\) in each degree. Here the complex \(g^*f_*I^\bullet\) computes \(g^*Rf_*F\). The complex \(f'_*g'^*I^\bullet\), on the other hand, is \(f'_*\) applied to a resolution of \(g'^*F\) that need not be injective, since \(g'^*\) does not preserve injectives. It still computes \(Rf'_*g'^*F\) if every \(g'^*I^m\) is \(f'_*\)-acyclic. We call this condition injective acyclicity, and Theorem \hyperref[bp329]{4.17} shows that it is equivalent to proper base change. Its advantage is that, unlike the base change statement, injective acyclicity is a vanishing condition on single sheaves.

% object: definition-injective-acyclicity (blueprint/chap-Proper-base-change/sec-Torsion-sheaves-and-acyclicity/definition-injective-acyclicity/_section.md)
\phantomsection\label{bp324}%
\begin{bpstatement}{Definition~4.15}{ (Injective acyclicity)}
For a ring \(\Lambda\), write \(\mathrm{IA}_{\Lambda}(f)\) if for every cartesian square (4.1) and every injective sheaf \(I\) of \(\Lambda\)-modules on \(X_{\acute{e}t}\) which is torsion, one has

\[
R^qf'_*g'^*I=0\qquad(q>0).
\]
\end{bpstatement}

% text: prose-04-16
\phantomsection\label{bp325}%
In general \(g'^*I\) is not injective, and injective acyclicity only asks that it be \(f'_*\)-acyclic. By the next lemma this suffices to compute \(Rf'_*\) of the pullback of an injective resolution term by term.

% object: lemma-acyclic-resolutions (blueprint/chap-Proper-base-change/sec-Torsion-sheaves-and-acyclicity/lemma-acyclic-resolutions/_section.md)
\phantomsection\label{bp326}%
\begin{bpstatementit}{Lemma~4.16}{ (Acyclic resolutions)}
Let \(v:\mathcal A\to\mathcal B\) be a left exact additive functor between abelian categories with enough injectives. If \(F\to M^\bullet\) is a resolution and each \(M^m\) is \(v\)-acyclic, then

\[
\mathcal H^q(vM^\bullet)\simeq R^qvF.
\]

More generally, for a bounded below complex \(M^\bullet\) of \(v\)-acyclic objects, the canonical comparison \(vM^\bullet\to Rv(M^\bullet)\) is an isomorphism in \(D^+(\mathcal B)\).
\end{bpstatementit}

% proof: proof/lemma-acyclic-resolutions
\begin{proof}[{Proof}]\phantomsection\label{bp327}%
This is Leray's acyclicity lemma [Stacks, \href{https://stacks.math.columbia.edu/tag/015E}{015E}], and the first assertion is the case of a resolution of \(F\). We recall the argument, which compares \(M^\bullet\) with an injective replacement through the cone.

In a bounded below exact complex of \(v\)-acyclic objects the cycles are again \(v\)-acyclic, by induction along the short exact sequences of cycles, so the complex stays exact after applying \(v\). We apply this to the cone of a quasi-isomorphism \(M^\bullet\to J^\bullet\) into a bounded below complex of injectives. This cone is bounded below and exact, and its terms are \(v\)-acyclic, being finite direct sums of terms of \(M^\bullet\) and of \(J^\bullet\). It follows that \(vM^\bullet\to vJ^\bullet\) is a quasi-isomorphism, and the target computes \(Rv(M^\bullet)\).
\end{proof}

% text: prose-04-17
\phantomsection\label{bp328}%
To pass from sheaves killed by an integer to all torsion sheaves, we use continuity of étale cohomology: if \(u:T\to S\) is quasi-compact and quasi-separated, then

\[
\varinjlim_i R^qu_*F_i\xrightarrow{\sim}
R^qu_*\bigl(\varinjlim_i F_i\bigr)\qquad(q\geq0)
\tag{4.7}
\]

for every filtered system of abelian étale sheaves on \(T\). The same holds after any base change of \(u\) [Stacks, \href{https://stacks.math.columbia.edu/tag/09Z1}{09Z1}]. Indeed, \(R^qu_*F\) is the sheafification of \(V\mapsto H^q(T\times_SV,F|_{T\times_SV})\) on affine étale \(V\to S\), and the étale cohomology of the quasi-compact and quasi-separated scheme \(T\times_SV\) commutes with filtered colimits of abelian sheaves [Stacks, \href{https://stacks.math.columbia.edu/tag/03Q5}{03Q5}].

The following criterion reduces proper base change in higher degrees to injective acyclicity, a vanishing statement that can be checked on curves.

% object: theorem-pbc-injective-acyclicity (blueprint/chap-Proper-base-change/sec-Torsion-sheaves-and-acyclicity/theorem-pbc-injective-acyclicity/_section.md)
\phantomsection\label{bp329}%
\begin{bpstatementit}{Theorem~4.17}{ (PBC and injective acyclicity)}
Let \(f:X\to Y\) be proper. The following are equivalent:

\begin{itemize}
\item (i) \(\mathrm{PBC}_\Lambda(f)\) holds.
\item (ii) For every sheaf \(F\) of \(\Lambda\)-modules and every base change (4.1), the maps induced by (4.3) on cohomology sheaves \[
  g^*R^qf_*F\longrightarrow R^qf'_*g'^*F
  \] are isomorphisms for all \(q\geq0\).
\item (iii) \(\mathrm{IA}_\Lambda(f)\) holds.
\end{itemize}
\end{bpstatementit}

% proof: proof/theorem-pbc-injective-acyclicity
\begin{proof}[{Proof}]\phantomsection\label{bp330}%
We prove (i)\(\Rightarrow\)(ii)\(\Rightarrow\)(iii)\(\Rightarrow\)(ii)\(\Rightarrow\)(i). The step (iii)\(\Rightarrow\)(ii) is the main one: it applies the degree-zero isomorphism termwise to an injective resolution.

\emph{(i)\(\Rightarrow\)(ii)\(\Rightarrow\)(iii).} Condition (ii) is the case of (i) where \(F\) sits in degree zero. Suppose (ii) holds, and let \(I\) be an injective sheaf of \(\Lambda\)-modules. Then \(R^qf'_*g'^*I\simeq g^*R^qf_*I\) by (ii). The right hand side vanishes for \(q>0\) because \(I\) is injective, and this is (iii).

\emph{(iii)\(\Rightarrow\)(ii).} Let \(F\) be a sheaf of \(\Lambda\)-modules. We compare the two sides of the base change map through an injective resolution \(F\to I^\bullet\) by sheaves of \(\Lambda\)-modules. The complex \(f'_*g'^*I^\bullet\) computes \(Rf'_*g'^*F\). Indeed, since \(g'^*\) is exact, \(g'^*I^\bullet\) is a resolution of \(g'^*F\), and its terms are \(f'_*\)-acyclic by (iii), so Lemma \hyperref[bp326]{4.16} applies. Similarly, the complex \(g^*f_*I^\bullet\) computes \(g^*Rf_*F\), because inverse image is exact. Between the two, Proposition \hyperref[bp314]{4.11} gives an isomorphism

\[
g^*f_*I^m\xrightarrow{\sim} f'_*g'^*I^m
\]

in each degree, and these form an isomorphism of complexes by naturality. By (4.4) this isomorphism is the derived base change map for \(F\), which proves (ii).

\emph{(ii)\(\Rightarrow\)(i).} Let \(K\in D^+(X_{\acute{e}t},\Lambda)\). For bounded \(K\) we argue by induction on the cohomological length. A complex of length one is a shifted sheaf, which (ii) covers, while a longer one is, by a canonical truncation triangle, an extension of two complexes of smaller length. Since the base change map is a morphism of triangles, the five lemma on cohomology sheaves shows that it is an isomorphism for \(K\). An unbounded \(K\) reduces to the bounded case, because a given cohomology sheaf only depends on a bounded truncation. More precisely, fix a degree \(m\) and choose \(N\geq m\), so that \(\tau_{\leq N}K\) is bounded and has the desired comparison. Both functors \(g^*Rf_*\) and \(Rf'_*g'^*\) send \(D^{\geq N+1}(X_{\acute{e}t},\Lambda)\) into \(D^{\geq N+1}(Y'_{\acute{e}t},\Lambda)\), since inverse images are exact and derived direct images are left \(t\)-exact. Applying them to

\[
\tau_{\leq N}K\longrightarrow K\longrightarrow
\tau_{\geq N+1}K\xrightarrow{+1}
\]

we see that the comparison for \(K\) in degree \(m\) is the comparison for \(\tau_{\leq N}K\), hence an isomorphism. As \(m\) was arbitrary, (i) follows.
\end{proof}


% object: sec-The-projective-line (blueprint/chap-Proper-base-change/sec-The-projective-line/_section.md), content hidden
\subsection*{\texorpdfstring{4.4\quad The projective line}{4.4 The projective line}}
\phantomsection\label{bp331}%
\addcontentsline{toc}{subsection}{\texorpdfstring{4.4\quad The projective line}{4.4 The projective line}}
% text: prose-04-18
\phantomsection\label{bp332}%
Let \(\pi:\mathbf P^1_S\to S\) be the projection. A geometric stalk of \(R^q\pi_*F\) is the cohomology of \(\mathbf P^1\) over a strictly henselian local ring, and proper base change for \(\pi\) says that it is the cohomology of the geometric fibre, and that this group does not change under base change.

% object: subsec-Henselian-pairs-and-curves (blueprint/chap-Proper-base-change/sec-The-projective-line/subsec-Henselian-pairs-and-curves/_section.md), content hidden
\subsubsection*{\texorpdfstring{4.4.1\quad Henselian pairs and curves}{4.4.1 Henselian pairs and curves}}
\phantomsection\label{bp333}%
\addcontentsline{toc}{subsubsection}{\texorpdfstring{4.4.1\quad Henselian pairs and curves}{4.4.1 Henselian pairs and curves}}
% text: prose-04-19
\phantomsection\label{bp334}%
The model case is a complete discrete valuation ring \(R\). Think of \(\operatorname{Spec}R\) as a small disc around its closed point, and of \(\mathbf P^1_R\) as a family of projective lines over the disc. In topology such a family retracts onto the central fibre, and we expect that, for torsion sheaves, the cohomology of \(\mathbf P^1_R\) is that of its closed fibre. Besides the comparison in degree zero of Theorem \hyperref[bp308]{4.8}, the proof uses two theorems about henselian pairs from the Stacks Project. The first is the affine analogue of what we want to prove.

% object: theorem-gabber-affine-proper-base-change (blueprint/chap-Proper-base-change/sec-The-projective-line/subsec-Henselian-pairs-and-curves/theorem-gabber-affine-proper-base-change/_section.md)
\phantomsection\label{bp335}%
\begin{bpstatementit}{Theorem~4.18}{ (Gabber's affine analogue of proper base change)}
Let \((B,J)\) be a henselian pair and let \(G\) be a torsion abelian sheaf on \(\operatorname{Spec}(B)_{\acute{e}t}\). Then restriction induces isomorphisms

\[
H^q(\operatorname{Spec}B,G)
\xrightarrow{\sim}
H^q(\operatorname{Spec}(B/J),G|_{\operatorname{Spec}(B/J)})
\qquad(q\geq0).
\]
\end{bpstatementit}

% text: prose-04-20
\phantomsection\label{bp336}%
This is [Stacks, \href{https://stacks.math.columbia.edu/tag/09ZI}{09ZI}]; Gabber proves it by an effacement argument on affines. The second fact concerns finite étale covers of a proper scheme. In degree one it says that \(\mathbf Z/n\)-torsors, like sections, are determined by the closed fibre.

% object: theorem-finite-etale-covers-proper-henselian (blueprint/chap-Proper-base-change/sec-The-projective-line/subsec-Henselian-pairs-and-curves/theorem-finite-etale-covers-proper-henselian/_section.md)
\phantomsection\label{bp337}%
\begin{bpstatementit}{Theorem~4.19}{ (Finite étale covers over a proper henselian pair)}
Let \((B,J)\) be a henselian pair, let \(Y\) be proper over \(B\), and put \(Y_0=Y\times_{\operatorname{Spec}B}\operatorname{Spec}(B/J)\). Then restriction is an equivalence between the categories of finite étale schemes over \(Y\) and over \(Y_0\),

\[
\textit{F\'Et}_Y\xrightarrow{\sim}\textit{F\'Et}_{Y_0}.
\]

Consequently, for a finite abelian group \(M\),

\[
H^1(Y,\underline M)\xrightarrow{\sim}H^1(Y_0,\underline M).
\tag{4.8}
\]
\end{bpstatementit}

% proof: proof/theorem-finite-etale-covers-proper-henselian
\begin{proof}[{Proof}]\phantomsection\label{bp338}%
The equivalence is [Stacks, \href{https://stacks.math.columbia.edu/tag/0GS2}{0GS2}], proved by algebraization, approximation, and descent of finite presentation. For the consequence, note that the groups in (4.8) classify \(M\)-torsors, which are finite étale schemes with an action of \(M\) satisfying the equations of a simply transitive action. The equivalence transports these schemes, their actions and the equations, so it induces a bijection on isomorphism classes of \(M\)-torsors.
\end{proof}

% text: prose-04-21
\phantomsection\label{bp339}%
The last step of the proof, the change of separably closed ground field, uses the following facts about curves over an algebraically closed field, which we also need later. For torsion divisible by the characteristic, proper and affine curves behave differently, which is why some of the assertions require \(C\) to be proper.

% object: theorem-cohomology-curves (blueprint/chap-Proper-base-change/sec-The-projective-line/subsec-Henselian-pairs-and-curves/theorem-cohomology-curves/_section.md)
\phantomsection\label{bp340}%
\begin{bpstatementit}{Theorem~4.20}{ (Cohomology of curves)}
Let \(k\) be algebraically closed, let \(C\) be a separated scheme of finite type over \(k\) with \(\dim C\leq1\), and let \(F\) be a classical torsion abelian sheaf on \(C_{\mathrm{pro}\acute{e}t}\). With cohomology taken on \(C_{\mathrm{pro}\acute{e}t}\), the following assertions hold:

\begin{itemize}
\item \((C_{1})\) \(H^q(C,F)=0\) for \(q>2\).
\item \((C_{2})\) If \(C\) is affine, \(H^q(C,F)=0\) for \(q>1\).
\item \((C_{3})\) If \(\operatorname{char}k=p>0\) and \(F\) is \(p\)-primary, \(H^q(C,F)=0\) for \(q>1\).
\item \((C_{4})\) If \(F\) is constructible with finite stalks and prime to \(\operatorname{char}k\), all \(H^q(C,F)\) are finite.
\item \((C_{5})\) If \(C\) is proper and \(F\) is constructible with finite stalks, all \(H^q(C,F)\) are finite.
\item \((C_{6})\) If \(k'/k\) is an extension of algebraically closed fields and \(F\) is prime to \(\operatorname{char}k\), the pullback maps \(H^q(C,F)\to H^q(C_{k'},F|)\) are isomorphisms.
\item \((C_{7})\) If \(C\) is proper, the maps in \((C_6)\) are isomorphisms for every torsion \(F\).
\item \((C_{8})\) For any open \(U\subseteq C\), restriction \(H^2(C,F)\to H^2(U,F|_U)\) is surjective.
\end{itemize}
\end{bpstatementit}

% text: prose-04-22
\phantomsection\label{bp341}%
In characteristic zero, the condition that \(F\) be prime to the characteristic in \((C_4)\) and \((C_6)\) is automatic.

% proof: proof/theorem-cohomology-curves
\begin{proof}[{Proof of Theorem 4.20}]\phantomsection\label{bp342}%
The proof is in Appendix \bpout{D}. For a smooth proper connected curve and \(n\) invertible in \(k\), one computes \(H^q(C,\mu_n)\) with the Kummer sequence and the Jacobian (Appendix \bpout{C}), while for \(p\)-primary coefficients one uses the Artin–Schreier sequence instead. The case of general torsion sheaves is then reduced to these cases by normalization, removal of finitely many points and a finite étale cover of degree prime to the torsion order.
\end{proof}

% text: prose-04-23
\phantomsection\label{bp343}%
The final step for the projective line compares cohomology over the residue fields of two strict localizations, which are separably closed but in general not algebraically closed. For this we need the following form of \((C_7)\).

% object: lemma-extension-separably-closed-ground-field (blueprint/chap-Proper-base-change/sec-The-projective-line/subsec-Henselian-pairs-and-curves/lemma-extension-separably-closed-ground-field/_section.md)
\phantomsection\label{bp344}%
\begin{bpstatementit}{Lemma~4.21}{ (Separably closed field extension)}
Let \(k\subseteq k'\) be separably closed fields, let \(C\) be proper over \(k\) with \(\dim C\leq1\), and let \(F\) be a classical torsion abelian sheaf on \(C_{\mathrm{pro}\acute{e}t}\). Then the natural map

\[
H^q(C,F)\xrightarrow{\sim}H^q(C_{k'},F|)
\qquad(q\geq0)
\]

is an isomorphism.
\end{bpstatementit}

% proof: proof/lemma-extension-separably-closed-ground-field
\begin{proof}[{Proof}]\phantomsection\label{bp345}%
A purely inseparable extension does not change the étale topos, which allows us to replace \(k\) and \(k'\) by their algebraic closures and then apply \((C_7)\).

By the classical comparison we may compute on the étale site. Choose compatible algebraic closures \(\bar k\subseteq\bar k'\). Since \(\bar k/k\) and \(\bar k'/k'\) are purely inseparable, \(C_{\bar k}\to C\) and \(C_{\bar k'}\to C_{k'}\) are universal homeomorphisms, and hence they induce equivalences of small étale topoi [Stacks, \href{https://stacks.math.columbia.edu/tag/03SI}{03SI}]. This reduces the claim to the pullback from \(C_{\bar k}\) to \(C_{\bar k'}\), which is an isomorphism by \((C_7)\) of Theorem \hyperref[bp340]{4.20}. Since all these maps are pullback maps, the comparison square commutes [Stacks, \href{https://stacks.math.columbia.edu/tag/0A5E}{0A5E}].
\end{proof}

% object: remark-properness-needed-c7 (blueprint/chap-Proper-base-change/sec-The-projective-line/subsec-Henselian-pairs-and-curves/remark-properness-needed-c7/_section.md)
\phantomsection\label{bp346}%
\begin{bpstatement}{Remark~4.22}{ (Why properness is needed in $(C_7)$)}
Properness in \((C_7)\) is needed when the torsion order is divisible by the characteristic, and this case arises in Proposition \hyperref[bp360]{4.26}, where the field is extended for arbitrary torsion sheaves. For example, in characteristic \(p\),

\[
H^1(\mathbf A^1_k,\mathbf F_p)
=k[t]/\{h^p-h:h\in k[t]\},
\]

and extending an algebraically closed ground field adds new classes. Indeed, we have \(c^pt^{pm}\equiv ct^m\), and every element of \(k\) is of the form \(h^p-h\). It follows that every class has a unique representative \(\sum_{p\nmid m}a_mt^m\), and the same holds over an algebraically closed \(k'\supset k\). Hence for \(a\in k'\smallsetminus k\) the class of \(at\) does not come from \(k[t]\).
\end{bpstatement}


% object: subsec-Injectives-on-the-closed-fibre (blueprint/chap-Proper-base-change/sec-The-projective-line/subsec-Injectives-on-the-closed-fibre/_section.md), content hidden
\subsubsection*{\texorpdfstring{4.4.2\quad Injectives on the closed fibre}{4.4.2 Injectives on the closed fibre}}
\phantomsection\label{bp347}%
\addcontentsline{toc}{subsubsection}{\texorpdfstring{4.4.2\quad Injectives on the closed fibre}{4.4.2 Injectives on the closed fibre}}
% text: prose-04-24
\phantomsection\label{bp348}%
Let \(R\) be a henselian local ring with maximal ideal \(\mathfrak m\) and residue field \(\kappa\), and let \(i:\mathbf P^1_\kappa\hookrightarrow\mathbf P^1_R\) be the inclusion of the closed fibre. For a sheaf \(G\) on \(\mathbf P^1_{R,\acute{e}t}\) write \(G_0=i^*G\). To compare the cohomology of \(\mathbf P^1_R\) with that of its closed fibre, we apply the comparison in degree zero of Theorem \hyperref[bp308]{4.8} to an injective resolution. This works once we know that the restriction \(I_0\) of an injective sheaf \(I\) is acyclic on \(\mathbf P^1_\kappa\) (Lemma \hyperref[bp352]{4.24}). However, an injective sheaf is too large to work with directly, so we approximate it by constructible sheaves and embed these into direct images along finite morphisms, using three standard finiteness facts about small étale topoi of quasi-compact quasi-separated schemes.

% object: lemma-finiteness-small-etale-topoi (blueprint/chap-Proper-base-change/sec-The-projective-line/subsec-Injectives-on-the-closed-fibre/lemma-finiteness-small-etale-topoi/_section.md)
\phantomsection\label{bp349}%
\begin{bpstatementit}{Lemma~4.23}{ (Finiteness facts on $\mathbf P^1_R$)}
\begin{enumerate}
\item A sheaf of \(\mathbf F_\ell\)-modules on \(\mathbf P^1_{R,\acute{e}t}\) is a filtered colimit of constructible ones, and cohomology on \(\mathbf P^1_R\) and on \(\mathbf P^1_\kappa\) commutes with filtered colimits.
\item A constructible sheaf \(G\) of \(\mathbf F_\ell\)-modules on \(\mathbf P^1_{R,\acute{e}t}\) admits a monomorphism \[
   G\hookrightarrow
   \bigoplus_{a=1}^N(h_a)_*\underline{M_a},
   \qquad h_a:Y_a\to \mathbf P^1_R\ \text{finite and finitely presented},
   \tag{4.9}
   \] with \(M_a\) finite \(\mathbf F_\ell\)-modules.
\item For a finite morphism \(h:Y\to \mathbf P^1_R\) and an abelian sheaf \(G\) on \(Y_{\acute{e}t}\), we have \[
   R^qh_*G=0\quad(q>0),\qquad
   i^*h_*G\xrightarrow{\sim}(h_0)_*i_Y^*G.
   \tag{4.10}
   \] Here \(h_0:Y_0=Y\times_{\mathbf P^1_R}\mathbf P^1_\kappa\to \mathbf P^1_\kappa\) and \(i_Y:Y_0\to Y\) are the base changes of \(h\) and \(i\), and the second map is the underived base change map (4.2).
\end{enumerate}
\end{bpstatementit}

% proof: proof/lemma-finiteness-small-etale-topoi
\begin{proof}[{Proof}]\phantomsection\label{bp350}%
Part (1) is [Stacks, \href{https://stacks.math.columbia.edu/tag/03SA}{03SA}] and [Stacks, \href{https://stacks.math.columbia.edu/tag/03Q5}{03Q5}], and part (2) is [Stacks, \href{https://stacks.math.columbia.edu/tag/09Z7}{09Z7}]. Part (3) is the case \(u=h\), \(W=\mathbf P^1_R\) of Lemma \hyperref[bp364]{4.27} below, which is proved directly from the stalk description of finite direct images [Stacks, \href{https://stacks.math.columbia.edu/tag/03QP}{03QP}].
\end{proof}

% text: prose-04-25
\phantomsection\label{bp351}%
The sheaves on the right of (4.9) are direct images along finite morphisms, and by Lemma lemma-finiteness-small-etale-topoi such direct images are exact functors and commute with restriction to the closed fibre.

% object: lemma-acyclicity-closed-fibre (blueprint/chap-Proper-base-change/sec-The-projective-line/subsec-Injectives-on-the-closed-fibre/lemma-acyclicity-closed-fibre/_section.md)
\phantomsection\label{bp352}%
\begin{bpstatementit}{Lemma~4.24}{ (Acyclicity on the closed fibre)}
If \(I\) is an injective sheaf of \(\mathbf F_\ell\)-modules on \(\mathbf P^1_{R,\acute{e}t}\), then \(H^q(\mathbf P^1_\kappa,I_0)=0\) for every \(q>0\).
\end{bpstatementit}

% proof: proof/lemma-acyclicity-closed-fibre
\begin{proof}[{Proof}]\phantomsection\label{bp353}%
We prove vanishing on affine charts by Gabber's theorem, deduce the case \(q\geq2\) by Mayer–Vietoris, and treat \(q=1\) by lifting classes to finite torsors over \(\mathbf P^1_R\), where injectivity of \(I\) kills them.

\emph{Step 1: affine opens.} Let \(W=\operatorname{Spec}B\) be an affine open of \(\mathbf P^1_R\) and set \(J=\mathfrak mB\). We show that \(I_0|_{W_0}\) has no cohomology in positive degrees. Since the pair \((B,J)\) need not be henselian, we pass to its henselization \((B^h,JB^h)\), for which \(B^h/JB^h=B/J\). Moreover, \(B^h\) is the filtered colimit of the étale \(B\)-algebras \(B'\) equipped with \(B'/JB'\simeq B/J\). Gabber's theorem (Theorem \hyperref[bp335]{4.18}) and continuity of étale cohomology give

\[
H^q(W_0,I_0|_{W_0})
\simeq H^q(\operatorname{Spec}B^h,I|)
\simeq\varinjlim_{B'}H^q(\operatorname{Spec}B',I|).
\tag{4.11}
\]

Every group on the right vanishes for \(q>0\), because the restriction of an injective sheaf to an étale object is injective: the restriction functor has an exact left adjoint, namely extension by zero in the localized site.

\emph{Step 2: degrees \(q\geq2\).} The two standard affine charts \(U,V\subset \mathbf P^1_R\) have affine intersection, so (4.11) shows that the restriction of \(I\) to each of \(U_0\), \(V_0\) and \((U\cap V)_0\) has no cohomology in positive degrees. The Mayer–Vietoris sequence on \(\mathbf P^1_\kappa\) therefore gives \(H^q(\mathbf P^1_\kappa,I_0)=0\) for \(q\geq2\) (this is the case of two charts in [Stacks, \href{https://stacks.math.columbia.edu/tag/0A51}{0A51}]). In degree one the same sequence only identifies \(H^1(\mathbf P^1_\kappa,I_0)\) with the cokernel of \(\Gamma(U_0,I_0)\oplus\Gamma(V_0,I_0)\to\Gamma((U\cap V)_0,I_0)\), and nothing said so far shows that this cokernel vanishes.

\emph{Step 3: degree one.} Let \(\xi\in H^1(\mathbf P^1_\kappa,I_0)\). Following [Stacks, \href{https://stacks.math.columbia.edu/tag/0A5H}{0A5H}], we show that \(\xi=0\). First write \(I\) as a filtered colimit of constructible \(\mathbf F_\ell\)-sheaves. Since pullback and \(H^1\) on \(\mathbf P^1_\kappa\) commute with this colimit, \(\xi\) is the image of a class \(\zeta\in H^1(\mathbf P^1_\kappa,G_0)\) for some constructible \(G\) with a map \(G\to I\). Next choose (4.9) and put \(P=\bigoplus_a(h_a)_*\underline{M_a}\). Since \(I\) is injective, the map \(G\to I\) extends to \(P\to I\). Let \(\zeta'\) be the image of \(\zeta\) in \(H^1(\mathbf P^1_\kappa,P_0)\).

The class \(\zeta'\) lifts to \(H^1(\mathbf P^1_R,P)\). Indeed, (4.10) and the Leray spectral sequence give

\[
H^1(\mathbf P^1_\kappa,P_0)
\simeq\bigoplus_a H^1((Y_a)_0,\underline{M_a}),\qquad
H^1(\mathbf P^1_R,P)\simeq\bigoplus_a H^1(Y_a,\underline{M_a}).
\]

Each \(Y_a\) is finite over \(\mathbf P^1_R\), hence proper over \(R\), so the comparison (4.8) for finite torsors (Theorem \hyperref[bp337]{4.19}) lifts \(\zeta'\). This lift maps to zero in \(H^1(\mathbf P^1_R,I)\) because \(I\) is injective. Now the commutative square of restriction maps

\[
\begin{CD}
H^1(\mathbf P^1_R,P) @>>> H^1(\mathbf P^1_R,I)\\
@VVV @VVV\\
H^1(\mathbf P^1_\kappa,P_0) @>>> H^1(\mathbf P^1_\kappa,I_0)
\end{CD}
\]

shows that the image of \(\zeta'\) in \(H^1(\mathbf P^1_\kappa,I_0)\) is zero, and this image is \(\xi\), because \(G\to P\to I\) is the given map \(G\to I\). Hence \(\xi=0\).
\end{proof}

% text: prose-04-26
\phantomsection\label{bp354}%
With this acyclicity, an injective resolution on \(\mathbf P^1_R\) restricts to an acyclic resolution on the closed fibre. Applying the degree-zero comparison of Theorem \hyperref[bp308]{4.8} term by term, we obtain the comparison in all degrees. Since the lemma only concerns sheaves killed by a prime \(\ell\), we then extend the result to all torsion sheaves by dévissage.

% object: proposition-cohomology-p-1-over-henselian (blueprint/chap-Proper-base-change/sec-The-projective-line/subsec-Injectives-on-the-closed-fibre/proposition-cohomology-p-1-over-henselian/_section.md)
\phantomsection\label{bp355}%
\begin{bpstatementit}{Proposition~4.25}{ (Cohomology of $\mathbf P^1$ over henselian rings)}
Let \(R\) be a henselian local ring with residue field \(\kappa\). For every torsion abelian sheaf \(F\) on \(\mathbf P^1_{R,\acute{e}t}\), restriction induces isomorphisms, natural in \(F\), with \(F_0=i^*F\),

\[
H^q(\mathbf P^1_R,F)\xrightarrow{\sim}
H^q(\mathbf P^1_\kappa,F_0)\qquad(q\geq0).
\]
\end{bpstatementit}

% proof: proof/proposition-cohomology-p-1-over-henselian
\begin{proof}[{Proof}]\phantomsection\label{bp356}%
We treat sheaves killed by a prime first, where Lemma \hyperref[bp352]{4.24} applies, and then reach all torsion sheaves by extensions, direct sums and filtered colimits.

Suppose first that \(\ell F=0\) for a prime \(\ell\).  Take an injective resolution \(F\to I^\bullet\) in sheaves of \(\mathbf F_\ell\)-modules on \(\mathbf P^1_{R,\acute{e}t}\). Since \(i^*\) is exact, \(F_0\to I_0^\bullet\) is again a resolution, and by Lemma \hyperref[bp352]{4.24} its terms are \(\Gamma(\mathbf P^1_\kappa,-)\)-acyclic. Moreover, Theorem \hyperref[bp308]{4.8} shows that the two complexes of sections agree term by term:

\[
\Gamma(\mathbf P^1_R,I^\bullet)\xrightarrow{\sim}
\Gamma(\mathbf P^1_\kappa,I_0^\bullet).
\]

The left complex computes \(R\Gamma(\mathbf P^1_R,F)\), and the right computes \(R\Gamma(\mathbf P^1_\kappa,F_0)\). Hence the natural restriction map between them is a quasi-isomorphism.

By the long exact cohomology sequences and the five lemma, the sheaves for which the statement holds form a class closed under extensions. Since a sheaf killed by \(\ell^a\) has a finite filtration with quotients killed by \(\ell\), and a sheaf killed by an integer \(n\) is the finite direct sum of its primary parts, the class contains all sheaves of bounded torsion. An arbitrary torsion \(F\) is the colimit \(\varinjlim_n F[n]\), with positive integers ordered by divisibility, and the comparison passes to this colimit, because the cohomology of the quasi-compact quasi-separated schemes \(\mathbf P^1_R\) and \(\mathbf P^1_\kappa\) commutes with filtered colimits [Stacks, \href{https://stacks.math.columbia.edu/tag/03Q5}{03Q5}]. Finally, all maps used are restriction maps, which gives naturality in \(F\).
\end{proof}

% text: prose-04-27
\phantomsection\label{bp357}%
For example, let \(R\) be strictly henselian and \(n\) invertible in \(R\). Then the residue field \(\kappa\) is separably closed, and passing to its algebraic closure does not change the étale topos [Stacks, \href{https://stacks.math.columbia.edu/tag/03SI}{03SI}]. Hence \(H^2(\mathbf P^1_\kappa,\mu_n)\simeq\mathbf Z/n\), generated by the Kummer class of \(\mathcal O(1)\) (Proposition \bpout{D.2}, in Appendix \bpout{D}), and the proposition says that \(H^2(\mathbf P^1_R,\mu_n)\) is also \(\mathbf Z/n\), generated by the class of \(\mathcal O(1)\).


% object: subsec-Arbitrary-base-change (blueprint/chap-Proper-base-change/sec-The-projective-line/subsec-Arbitrary-base-change/_section.md), content hidden
\subsubsection*{\texorpdfstring{4.4.3\quad Arbitrary base change}{4.4.3 Arbitrary base change}}
\phantomsection\label{bp358}%
\addcontentsline{toc}{subsubsection}{\texorpdfstring{4.4.3\quad Arbitrary base change}{4.4.3 Arbitrary base change}}
% text: prose-04-28
\phantomsection\label{bp359}%
Let \(\pi:\mathbf P^1_S\to S\) be the projection and \(\bar s\) a geometric point of \(S\). The stalk of \(R^q\pi_*F\) at \(\bar s\) is the cohomology of \(\mathbf P^1\) over the strict localization at \(\bar s\), by (4.12) below. By Proposition \hyperref[bp355]{4.25} it is therefore the cohomology of \(\mathbf P^1_{k_s}\), where \(k_s\) is the separably closed residue field. Along a base change \(g:T\to S\) the residue field grows from \(k_s\) to \(k_t\), so on stalks the base change map becomes the pullback from \(\mathbf P^1_{k_s}\) to \(\mathbf P^1_{k_t}\), which is an isomorphism by Lemma \hyperref[bp344]{4.21}. This is the second place where properness enters: for \(\mathbf A^1\) and \(p\)-torsion coefficients the corresponding map is not surjective.

The stalk formula holds for every quasi-compact quasi-separated morphism \(f:Z\to S\). For a geometric point \(\bar s\) of \(S\) and an abelian sheaf \(G\) on \(Z_{\acute{e}t}\) we have

\[
(R^qf_*G)_{\bar s}
\simeq H^q(Z\times_SS_{(\bar s)},G|),\qquad
S_{(\bar s)}=\operatorname{Spec}\mathcal O^{\mathrm{sh}}_{S,\bar s}.
\tag{4.12}
\]

Indeed, the stalk is the filtered colimit of \(H^q(Z\times_SV,G|)\) over pointed étale neighbourhoods \((V,\bar v)\) of \(\bar s\), and by continuity of étale cohomology this colimit is the cohomology of the indicated limit scheme [Stacks, \href{https://stacks.math.columbia.edu/tag/03Q9}{03Q9}].

% object: proposition-proper-base-change-projective-line (blueprint/chap-Proper-base-change/sec-The-projective-line/subsec-Arbitrary-base-change/proposition-proper-base-change-projective-line/_section.md)
\phantomsection\label{bp360}%
\begin{bpstatementit}{Proposition~4.26}{ (Proper base change for $\mathbf P^1$)}
Let \(S\) be a scheme and \(\Lambda\) a commutative ring.  In the cartesian square

\[
\begin{CD}
\mathbf P^1_T @>{g'}>> \mathbf P^1_S\\
@V{\pi'}VV @VV{\pi}V\\
T @>{g}>> S,
\end{CD}
\]

the base change map

\[
g^*R\pi_*K\longrightarrow R\pi'_*g'^*K
\]

is an isomorphism for every \(K\in D^+(\mathbf P^1_{S,\acute{e}t},\Lambda)\) with torsion cohomology sheaves. In particular, \(\mathrm{IA}_\Lambda(\pi)\) holds.
\end{bpstatementit}

% proof: proof/proposition-proper-base-change-projective-line
\begin{proof}[{Proof}]\phantomsection\label{bp361}%
For a torsion sheaf we compare the two sides on geometric stalks. Both stalks are the cohomology of a projective line over a separably closed field, and the base change map becomes the pullback along a field extension. Bounded below complexes then follow by truncation.

Let \(F\) be a torsion sheaf. Choose a geometric point \(\bar t:\operatorname{Spec}\Omega\to T\) and put \(\bar s=g\bar t\). Then the strict localizations fit into a natural map \(T_{(\bar t)}\to S_{(\bar s)}\) of strictly henselian local schemes. Write \(k_s,k_t\) for their residue fields; these are separably closed, and the map induces an inclusion \(k_s\hookrightarrow k_t\). For each \(q\geq0\), using (4.12) and Proposition \hyperref[bp355]{4.25}, we obtain the commutative diagram

\[
\begin{CD}
(g^*R^q\pi_*F)_{\bar t}
@>>>
(R^q\pi'_*g'^*F)_{\bar t}\\
@V{\sim}VV @VV{\sim}V\\
H^q(\mathbf P^1_{S_{(\bar s)}},F|)
@>>> H^q(\mathbf P^1_{T_{(\bar t)}},F|)\\
@V{\sim}VV @VV{\sim}V\\
H^q(\mathbf P^1_{k_s},F|)
@>>> H^q(\mathbf P^1_{k_t},F|).
\end{CD}
\]

Here the top arrow is the stalk at \(\bar t\) of the map induced by (4.3) on \(q\)th cohomology sheaves, while the middle and bottom arrows are the natural pullback maps. The top square commutes by the definition of the base change map: a section over \(\mathbf P^1_V\), for a pointed étale neighbourhood \(V\to S\), is pulled back to \(\mathbf P^1_{V\times_ST}\), and passing to the limits over the two strict localizations gives exactly the middle arrow. The bottom square commutes by naturality of restriction to the closed fibres.

The vertical arrows are isomorphisms: the upper ones by (4.12), applied to \(\pi\) and to \(\pi'\); here \(\mathbf P^1_{S_{(\bar s)}}\) and \(\mathbf P^1_{T_{(\bar t)}}\) are the base changes of \(\mathbf P^1_S\) and \(\mathbf P^1_T\) to the strict localizations. The lower ones are isomorphisms by Proposition \hyperref[bp355]{4.25}, applied to the strictly henselian local rings of \(S_{(\bar s)}\) and \(T_{(\bar t)}\). Hence the top arrow is an isomorphism as soon as the bottom arrow is, and the bottom arrow is an isomorphism by Lemma \hyperref[bp344]{4.21}, applied to \(C=\mathbf P^1_{k_s}\) and the extension \(k_s\subseteq k_t\). Thus every geometric stalk of \(g^*R^q\pi_*F\to R^q\pi'_*g'^*F\) is an isomorphism. As geometric stalks are conservative, the base change map for \(F\) is an isomorphism in \(D^+(T_{\acute{e}t},\mathbf Z)\). Moreover, cohomology and higher direct images of sheaves of \(\Lambda\)-modules agree with those of their underlying abelian sheaves (Lemma \hyperref[bp292]{4.2}). Hence the base change map for \(F\) is also an isomorphism in \(D^+(T_{\acute{e}t},\Lambda)\).

Now let \(K\in D^+(\mathbf P^1_{S,\acute{e}t},\Lambda)\) have torsion cohomology and fix \(m\in\mathbf Z\). The two functors \(g^*R\pi_*,\ R\pi'_*g'^*:D^+(\mathbf P^1_{S,\acute{e}t},\Lambda)\to
D^+(T_{\acute{e}t},\Lambda)\) send \(D^{\geq m+1}\) into \(D^{\geq m+1}\), since inverse image is exact and derived direct image is left \(t\)-exact. Using the truncation triangle for \(\tau_{\leq m}K\), we can therefore reduce the assertion in degree \(m\) to the bounded complex \(\tau_{\leq m}K\). By induction on the length, using the truncation triangles and naturality of the base change map, we reduce further to the shifted torsion sheaves \(H^a(K)[-a]\). For these the assertion has already been proved.

For \(\mathrm{IA}_\Lambda(\pi)\), let \(I\) be a torsion injective sheaf of \(\Lambda\)-modules. Then \(R^q\pi_*I=0\) for \(q>0\), and the case of sheaves gives \(R^q\pi'_*g'^*I\simeq g^*R^q\pi_*I=0\).
\end{proof}



% object: sec-Proper-morphisms (blueprint/chap-Proper-base-change/sec-Proper-morphisms/_section.md), content hidden
\subsection*{\texorpdfstring{4.5\quad Proper morphisms}{4.5 Proper morphisms}}
\phantomsection\label{bp362}%
\addcontentsline{toc}{subsection}{\texorpdfstring{4.5\quad Proper morphisms}{4.5 Proper morphisms}}
% text: prose-04-29
\phantomsection\label{bp363}%
Proper base change for an arbitrary proper morphism follows from the case of the projective line (Proposition \hyperref[bp360]{4.26}) by formal steps. Finite morphisms satisfy proper base change directly, and Lemma \hyperref[bp369]{4.29} shows that the property is stable under composition and descends along proper surjections. Descent along a proper surjection \(u:X\to Y\) works because every injective sheaf on \(Y\) is a direct summand of \(u_*J\) for an injective sheaf \(J\) on \(X\), and the hypotheses on \(u\) and on the composite keep \(u_*J\) acyclic after base change.

The steps are then as follows. First, composition gives the projection of \((\mathbf P^1)^n\). Then descent along the finite surjection \((\mathbf P^1)^n\to\mathbf P^n\) gives the projection of \(\mathbf P^n\), and with it every projective morphism. Finally, by Chow's lemma every proper morphism is dominated by a projective one through a proper surjection, so a last descent covers all proper morphisms.

% object: lemma-finite-morphisms (blueprint/chap-Proper-base-change/sec-Proper-morphisms/lemma-finite-morphisms/_section.md)
\phantomsection\label{bp364}%
\begin{bpstatementit}{Lemma~4.27}{ (Finite morphisms)}
Let \(u:V\to W\) be finite. For an abelian sheaf \(F\) and a geometric point \(\bar w:\operatorname{Spec}(\Omega)\to W\), with \(\Omega\) algebraically closed, one has

\[
(u_*F)_{\bar w}
\simeq \prod_{\substack{\bar v:\operatorname{Spec}(\Omega)\to V\\
u\bar v=\bar w}}F_{\bar v}.
\]

The product is finite. The functor \(u_*\) from abelian sheaves on \(V_{\acute{e}t}\) to abelian sheaves on \(W_{\acute{e}t}\) is exact, and its formation commutes with arbitrary base change. Consequently, for every base change \(W'\to W\) and every constant coefficient ring \(\Lambda\), the derived base change map (4.3) is an isomorphism for every \(K\in D^+(V_{\acute{e}t},\Lambda)\).
\end{bpstatementit}

% proof: proof/lemma-finite-morphisms
\begin{proof}[{Proof}]\phantomsection\label{bp365}%
Everything follows from the stalk formula. To prove it, note that over a strict localization the source splits into strictly henselian local pieces, and on such a piece taking sections is the same as taking a stalk.

After base change to the strict henselization of \(W\) at \(\bar w\), the finite scheme \(V\) becomes a finite disjoint union of strictly henselian local schemes. On each such local scheme, the sections of an étale sheaf are its stalk at the closed geometric point. Indeed, every étale neighbourhood of that point admits a section, so a germ extends to a global section, and two sections with the same germ agree after pulling back along such a section. In particular, taking sections is exact on abelian sheaves. Combining this with the description of stalks of direct images via strict localizations [Stacks, \href{https://stacks.math.columbia.edu/tag/03QP}{03QP}], we obtain the displayed formula. Since geometric stalks are conservative and finite products of abelian groups are exact, \(u_*\) is exact.

In a cartesian square over \(W'\to W\), a geometric point of \(W'\) and its image in \(W\) have the same geometric fibre. Therefore the underived base change map is the identity on the corresponding factors of the displayed formula, hence an isomorphism. Since both direct images in the square are exact, their derived functors and the base change map are computed term by term. Finally, the same argument applies to sheaves of modules.
\end{proof}

% text: prose-04-30
\phantomsection\label{bp366}%
For instance, if \(u:\operatorname{Spec}k'\to\operatorname{Spec}k\) is a separable extension of degree \(d\) and \(\bar w\) is a geometric point, the formula says that \((u_*F)_{\bar w}\) is the product of the \(d\) stalks of \(F\) at the \(d\) embeddings of \(k'\) into \(\Omega\) over \(k\).

% object: example-finite-direct-image-double-cover (blueprint/chap-Proper-base-change/sec-Proper-morphisms/example-finite-direct-image-double-cover/_section.md)
\phantomsection\label{bp367}%
\begin{bpstatement}{Example~4.28}{ (A ramified double cover)}
Let \(k\) be algebraically closed of characteristic \(\neq2\), let \(u:\mathbf A^1_k\to\mathbf A^1_k\) be \(x\mapsto x^2\), and let \(A\) be a finite ring. Over a point \(t\neq0\) the fibre consists of the two points \(\pm\sqrt t\), whereas over \(t=0\) it is the single point \(x=0\), with a nonreduced structure that the étale topology does not see. By the lemma, therefore,

\[
(u_*A)_{\bar t}=A\oplus A\quad(t\neq0),\qquad (u_*A)_{\bar0}=A,
\]

and the rank of the stalks jumps at the origin, so \(u_*A\) is not lisse. It is, however, constructible for the stratification \(\mathbf G_m\sqcup\{0\}\). Base change to the origin gives the stalk \(A\), the cohomology of the fibre, as the lemma predicts. For the open immersion \(j:\mathbf G_m\to\mathbf A^1\), by contrast, the stalk of \(j_*A\) at \(0\) is \(A\) while the fibre is empty.
\end{bpstatement}

% text: prose-04-31
\phantomsection\label{bp368}%
Stability under composition and descent along proper surjections are best proved on the level of injective acyclicity, since by Theorem \hyperref[bp329]{4.17} it is equivalent to proper base change and it lends itself to a spectral sequence argument.

% object: lemma-composition-cancellation (blueprint/chap-Proper-base-change/sec-Proper-morphisms/lemma-composition-cancellation/_section.md)
\phantomsection\label{bp369}%
\begin{bpstatementit}{Lemma~4.29}{ (Composition and cancellation)}
Let \(X\xrightarrow{u}Y\xrightarrow{v}Z\) be proper morphisms. Then proper base change is stable under composition, and it descends along a proper surjection:

\begin{itemize}
\item (i) If \(\mathrm{PBC}_\Lambda(u)\) and \(\mathrm{PBC}_\Lambda(v)\) hold, then \(\mathrm{PBC}_\Lambda(vu)\) holds.
\item (ii) If \(u\) is surjective and \(\mathrm{PBC}_\Lambda(u)\) and \(\mathrm{PBC}_\Lambda(vu)\) hold, then \(\mathrm{PBC}_\Lambda(v)\) holds.
\end{itemize}
\end{bpstatementit}

% proof: proof/lemma-composition-cancellation
\begin{proof}[{Proof}]\phantomsection\label{bp370}%
Theorem \hyperref[bp329]{4.17} reduces both parts to the condition \(\mathrm{IA}_\Lambda\) of injective acyclicity, and we read off this vanishing from the Leray spectral sequence of the base-changed composite. For (ii) the hypotheses do not apply directly to an injective sheaf on \(Y\), so we first realize it as a direct summand of \(u_*J\) with \(J\) injective on \(X\).

To set up the spectral sequence, fix an arbitrary base change \(Z'\to Z\). Write \(g_X:X'\to X\) and \(g_Y:Y'\to Y\) for the projections, and \(u':X'\to Y'\), \(v':Y'\to Z'\) for the induced maps. The direct image \(u'_*\) preserves injectives, because its left adjoint \(u'^*\) on sheaves of \(\Lambda\)-modules is exact. Hence for every sheaf \(J\) of \(\Lambda\)-modules on \(X\) we have the Grothendieck spectral sequence for the composite direct image [Stacks, \href{https://stacks.math.columbia.edu/tag/0734}{0734}]:

\[
R^pv'_*R^qu'_*g_X^*J
\ \Longrightarrow\ R^{p+q}(v'u')_*g_X^*J.
\tag{4.13}
\]

\emph{Part (i).} Let \(J\) be an injective sheaf of \(\Lambda\)-modules on \(X\). We show \(R^m(v'u')_*g_X^*J=0\) for \(m>0\), which is \(\mathrm{IA}_\Lambda(vu)\). For this it suffices that the \(E_2\)-terms of (4.13) vanish except for \(p=q=0\), and they do. Indeed, \(R^qu'_*g_X^*J=0\) for \(q>0\) by \(\mathrm{IA}_\Lambda(u)\). For \(q=0\), proper base change in degree zero (Proposition \hyperref[bp314]{4.11}) gives

\[
u'_*g_X^*J\simeq g_Y^*u_*J,
\]

and the sheaf \(u_*J\) is injective, so \(R^pv'_*g_Y^*u_*J=0\) for \(p>0\) by \(\mathrm{IA}_\Lambda(v)\). The abutment therefore vanishes in positive degrees.

\emph{Part (ii), the splitting.} Let \(I\) be an injective sheaf of \(\Lambda\)-modules on \(Y\). Choose a monomorphism of sheaves

\[
u^*I\hookrightarrow J
\]

with \(J\) injective on \(X\), and let \(\iota:I\to u_*J\) be its adjoint. The map \(\iota\) is a monomorphism. To see this, given a geometric point \(\bar y\), choose a geometric point \(\bar x\) above it, which exists because \(u\) is surjective. Then the composite

\[
I_{\bar y}\xrightarrow{\iota_{\bar y}}(u_*J)_{\bar y}
\longrightarrow J_{\bar x}
\]

is the stalk of the chosen monomorphism, under \((u^*I)_{\bar x}=I_{\bar y}\), so it is injective. As \(I\) is an injective sheaf, there is then a map of sheaves \(s:u_*J\to I\) with \(s\iota=\mathrm{id}_I\). In other words, \(I\) is a direct summand of \(u_*J\).

\emph{Part (ii), the vanishing.} We first show \(R^pv'_*g_Y^*u_*J=0\) for \(p>0\). The hypotheses \(\mathrm{IA}_\Lambda(u)\) and \(\mathrm{IA}_\Lambda(vu)\) give

\[
R^qu'_*g_X^*J=0\quad(q>0),\qquad
R^m(v'u')_*g_X^*J=0\quad(m>0),
\]

so the spectral sequence (4.13) collapses to its row \(q=0\). This means that \(R^pv'_*u'_*g_X^*J\simeq R^p(v'u')_*g_X^*J=0\) for \(p>0\). Since \(u'_*g_X^*J\simeq g_Y^*u_*J\) by proper base change in degree zero, we obtain

\[
R^pv'_*g_Y^*u_*J
\simeq R^pv'_*u'_*g_X^*J=0\qquad(p>0).
\]

Now \(g_Y^*I\) is a direct summand of \(g_Y^*u_*J\), by applying \(g_Y^*\) to the splitting \(s\iota=\mathrm{id}_I\), and as \(R^pv'_*\) is additive, \(R^pv'_*g_Y^*I\) is a direct summand of \(R^pv'_*g_Y^*u_*J\). Hence \(R^pv'_*g_Y^*I=0\) for \(p>0\), which is \(\mathrm{IA}_\Lambda(v)\).
\end{proof}

% text: prose-04-32
\phantomsection\label{bp371}%
Part (i) already gives proper base change for \((\mathbf P^1_S)^n\to S\), a composite of projections of projective lines. To reach \(\mathbf P^n_S\) we use part (ii), and for this we need a finite surjection from \((\mathbf P^1_S)^n\) onto \(\mathbf P^n_S\) over \(S\).

% object: lemma-finite-cover-projective-space (blueprint/chap-Proper-base-change/sec-Proper-morphisms/lemma-finite-cover-projective-space/_section.md)
\phantomsection\label{bp372}%
\begin{bpstatementit}{Lemma~4.30}{ (A finite cover of projective space)}
For every scheme \(S\) and every \(n\geq1\) there is a finite surjective \(S\)-morphism

\[
\sigma_n:(\mathbf P^1_S)^n\longrightarrow\mathbf P^n_S.
\]
\end{bpstatementit}

% proof: proof/lemma-finite-cover-projective-space
\begin{proof}[{Proof}]\phantomsection\label{bp373}%
Recall that \(\mathbf P^n\) parametrizes binary forms of degree \(n\) up to scalars, and the map \(\sigma_n\) multiplies \(n\) linear forms; over a field it is the quotient of \((\mathbf P^1)^n\) by the symmetric group. Explicitly, we use the coefficients of a product of binary linear forms:

\[
\prod_{i=1}^n(x_iT+y_iU)
=\sum_{j=0}^nF_jT^{n-j}U^j,\qquad
\sigma_n\bigl(([x_i:y_i])_i\bigr)=[F_0:\cdots:F_n].
\]

Every \(F_j\) is homogeneous of degree one in each pair \((x_i,y_i)\), and the corresponding sections of \(\mathcal O(1,\ldots,1)\) have no common zero. Indeed, at a geometric point every linear factor is nonzero, and hence so is their product. Therefore the sections define the stated morphism over any \(S\), and the construction commutes with base change.

For example, on the chart \(F_0\ne0\) of the target, whose preimage is the chart where every \(x_i\) is invertible, the coordinates \(c_j=F_j/F_0\) pull back to the elementary symmetric polynomials in the \(w_i=y_i/x_i\), as in the classical description of \(\mathbf A^n\) as the quotient of \(\mathbf A^n\) by the symmetric group.

Finiteness of \(\sigma_n\) follows from properness and finiteness of the fibres. Since the source is proper over \(S\) and the target is separated over \(S\), the morphism \(\sigma_n\) is proper. For the fibres, recall that over an algebraically closed field, a nonzero binary form of degree \(n\) is a product of \(n\) linear forms \(x_iT+y_iU\), and these are unique up to order and up to multiplying each by a nonzero constant. Here a factor with \(x_i=0\) is a multiple of \(U\) and corresponds to the point \([0:1]\). In other words, the form determines the multiset of points \([x_i:y_i]\in\mathbf P^1\). Hence every geometric fibre of \(\sigma_n\) is nonempty, and \(\sigma_n\) is surjective. Moreover, each geometric fibre has at most \(n!\) geometric points, namely the orderings of this multiset. Since a scheme of finite type over a field with finitely many geometric points is zero dimensional, \(\sigma_n\) is quasi-finite, and a proper quasi-finite morphism is finite [Stacks, \href{https://stacks.math.columbia.edu/tag/02LS}{02LS}].
\end{proof}

% text: prose-04-33
\phantomsection\label{bp374}%
For \(n=1\) the map \(\sigma_1\) is the identity of \(\mathbf P^1_S\). For \(n=2\) it sends \(([x_1:y_1],[x_2:y_2])\) to \([x_1x_2:x_1y_2+y_1x_2:y_1y_2]\), and over a field it is the quotient by swapping the two factors, of degree \(2\).

% object: lemma-pbc-projective-space (blueprint/chap-Proper-base-change/sec-Proper-morphisms/lemma-pbc-projective-space/_section.md)
\phantomsection\label{bp375}%
\begin{bpstatementit}{Lemma~4.31}{ (Projective space)}
If \(\mathrm{PBC}_\Lambda(\mathbf P^1_S\to S)\) holds for every scheme \(S\), then \(\mathrm{PBC}_\Lambda(\mathbf P^n_S\to S)\) holds for every scheme \(S\) and every \(n\geq0\).
\end{bpstatementit}

% proof: proof/lemma-pbc-projective-space
\begin{proof}[{Proof}]\phantomsection\label{bp376}%
The case \(n=0\) is the identity. For \(n\geq1\), the projection \((\mathbf P^1_S)^n\to S\) is a composite of \(n\) projections of projective lines over successive bases. Hence it satisfies \(\mathrm{PBC}_\Lambda\) by the hypothesis and Lemma lemma-composition-cancellation. The finite surjection \(\sigma_n\) of Lemma \hyperref[bp372]{4.30} is a morphism over \(S\), and it satisfies \(\mathrm{PBC}_\Lambda\) by Lemma \hyperref[bp364]{4.27}. Cancelling \(\sigma_n\) by Lemma lemma-composition-cancellation, we obtain \(\mathrm{PBC}_\Lambda\) for \(\mathbf P^n_S\to S\).
\end{proof}

% text: prose-04-34
\phantomsection\label{bp377}%
We can now treat projective morphisms, since they factor as a closed immersion followed by the projection of a projective space. By Chow's lemma an arbitrary proper morphism is dominated by a projective one through a proper surjection, and the cancellation of Lemma lemma-composition-cancellation then passes proper base change to the proper morphism.

% object: proposition-proper-base-change-etale-site (blueprint/chap-Proper-base-change/sec-Proper-morphisms/proposition-proper-base-change-etale-site/_section.md)
\phantomsection\label{bp378}%
\begin{bpstatementit}{Proposition~4.32}{ (Étale proper base change)}
If \(\mathrm{PBC}_\Lambda(\mathbf P^1_S\to S)\) holds for every scheme \(S\), then \(\mathrm{PBC}_\Lambda(f)\) holds for every proper morphism of schemes \(f\).
\end{bpstatementit}

% proof: proof/proposition-proper-base-change-etale-site
\begin{proof}[{Proof}]\phantomsection\label{bp379}%
Let \(f:X\to Y\) be proper. We first reduce to the case that \(Y\) is affine, as Chow's lemma requires. Let \(Y=\bigcup U\) be an affine open cover, and assume \(\mathrm{PBC}_\Lambda\) for each restriction \(f^{-1}(U)\to U\). Given a base change \(Y'\to Y\), the inverse images \(U'\) of the sets \(U\) then cover \(Y'\). Since higher direct images commute with restriction to open subschemes, the restriction of the base change map to \(U'\) is the base change map of \(f^{-1}(U)\to U\) along \(U'\to U\), hence an isomorphism. As isomorphisms can be checked on geometric stalks, the base change map over \(Y'\) is an isomorphism.

Now \(Y\) is affine, hence quasi-compact and quasi-separated, so Theorem \bpout{A.25} with \(S=Y\) gives a proper surjection \(\pi:\widetilde X\to X\) and closed immersions \(i_Y:\widetilde X\hookrightarrow\mathbf P^n_Y\) and \(i_X=(\pi,i_Y):\widetilde X\hookrightarrow\mathbf P^n_X\) with \(f\pi\) equal to \(i_Y\) followed by the projection. Here we need Chow's lemma over a base that is affine but need not be Noetherian; we prove this form in Appendix \bpout{A}. The two maps \(\pi\) and \(f\pi\) factor as closed immersions \(i_X\) and \(i_Y\) followed by projections of projective spaces. Closed immersions are finite, so they satisfy \(\mathrm{PBC}_\Lambda\) by Lemma \hyperref[bp364]{4.27}, and projections of projective spaces satisfy \(\mathrm{PBC}_\Lambda\) by the hypothesis and Lemma \hyperref[bp375]{4.31}. By stability under composition (Lemma lemma-composition-cancellation), both \(\pi\) and \(f\pi\) therefore satisfy \(\mathrm{PBC}_\Lambda\). As \(\pi\) is surjective, we may cancel it by Lemma lemma-composition-cancellation, and this gives \(\mathrm{PBC}_\Lambda\) for \(f\).
\end{proof}


% object: sec-Pro-etale-complexes (blueprint/chap-Proper-base-change/sec-Pro-etale-complexes/_section.md), content hidden
\subsection*{\texorpdfstring{4.6\quad Pro-étale complexes}{4.6 Pro-étale complexes}}
\phantomsection\label{bp380}%
\addcontentsline{toc}{subsection}{\texorpdfstring{4.6\quad Pro-étale complexes}{4.6 Pro-étale complexes}}
% text: prose-04-35
\phantomsection\label{bp381}%
The six operations live on the pro-étale site, where we also meet torsion sheaves of modules over rings that are not torsion, such as \(\mathcal O_E\). So far \(\mathrm{PBC}_\Lambda(f)\) is known on the étale site only when \(\Lambda\) is torsion (Propositions \hyperref[bp360]{4.26} and \hyperref[bp378]{Proposition 4.32}). We therefore drop the condition on \(\Lambda\), keeping torsion cohomology sheaves, and pass to classical pro-étale complexes by the comparison theorems of Chapter \hyperref[bp3]{2}.

The reduction to torsion rings is possible because a torsion sheaf of \(\Lambda\)-modules is a filtered colimit of sheaves killed by an integer \(n\). Such a sheaf has a finite filtration with quotients killed by primes \(p\), and these quotients are modules over \(\Lambda/p\Lambda\). Moreover, the cohomology of a sheaf of modules is that of its underlying abelian sheaf, so the base change map does not depend on which ring we regard as the coefficients.

For the passage to the pro-étale site we use two facts: a classical complex in \(D^+\) is \(\nu_X^*\) of an étale one (Corollary \hyperref[bp109]{2.37}), and \(\nu^*\) commutes with \(g^*\) and with direct images along proper maps (Theorem \hyperref[bp113]{2.38}). Taking for \(g\) a geometric point then shows that the stalks of \(Rf_*K\) are the cohomology of the geometric fibres (Corollary \hyperref[bp386]{4.34}).

% object: theorem-proper-base-change-classical-pro (blueprint/chap-Proper-base-change/sec-Pro-etale-complexes/theorem-proper-base-change-classical-pro/_section.md)
\phantomsection\label{bp382}%
\begin{bpstatementit}{Theorem~4.33}{ (Pro-étale proper base change)}
Let \(\Lambda\) be a commutative ring, let \(f:X\to Y\) be a proper morphism of schemes, and let \(g:Y'\to Y\) be arbitrary. Set \(X'=X\times_YY'\), with projections \(f':X'\to Y'\) and \(g':X'\to X\). Let \(K\in D^+_{\mathrm{class}}(X_{\mathrm{pro}\acute{e}t},\Lambda)\) have torsion cohomology sheaves (as abelian sheaves). Then the base change map (4.5) of Definition \hyperref[bp302]{4.6}

\[
g^*Rf_*K\longrightarrow Rf'_*g'^*K
\]

is an isomorphism in \(D^+_{\mathrm{class}}(Y'_{\mathrm{pro}\acute{e}t},\Lambda)\).
\end{bpstatementit}

% text: prose-04-36
\phantomsection\label{bp383}%
Suppose \(\Lambda\) is killed by a positive integer, for instance that \(\Lambda\) is a finite ring. Then every complex of \(\Lambda_X\)-modules has torsion cohomology sheaves, and the theorem applies to every \(K\in D^+_{\mathrm{class}}(X_{\mathrm{pro}\acute{e}t},\Lambda)\).

% proof: proof/theorem-proper-base-change-classical-pro
\begin{proof}[{Proof of Theorem 4.33}]\phantomsection\label{bp384}%
We work on the étale site, where Propositions \hyperref[bp360]{4.26} and \hyperref[bp378]{Proposition 4.32} apply to the torsion rings \(\Lambda/p\Lambda\). By dévissage we reach all torsion sheaves of \(\Lambda\)-modules, and the comparison theorems then bring us back to the pro-étale site.

By Corollary \hyperref[bp109]{2.37}, we may write \(K=\nu_X^*K_{\acute{e}t}\) with \(K_{\acute{e}t}\in D^+(X_{\acute{e}t},\Lambda)\), and the cohomology sheaves of \(K_{\acute{e}t}\) are torsion. Indeed, \(\nu_X^*\) is exact, commutes with colimits and is fully faithful on sheaves, and therefore detects the condition \(F=\varinjlim_nF[n]\) of Definition \hyperref[bp320]{4.13}. We must show that the base change map for \(K_{\acute{e}t}\) is an isomorphism in \(D^+(Y'_{\acute{e}t},\Lambda)\), and then pass to the pro-étale site.

We reduce to the rings \(\Lambda/p\Lambda\), \(p\) prime. These are killed by \(p\), so \(\mathrm{PBC}_{\Lambda/p\Lambda}(f)\) holds by Propositions \hyperref[bp360]{4.26} and \hyperref[bp378]{Proposition 4.32}. Now a sheaf \(F\) of \(\Lambda\)-modules with \(pF=0\) is a sheaf of \(\Lambda/p\Lambda\)-modules, and restriction of scalars along \(\Lambda\to\Lambda/p\Lambda\) is exact and conservative. Moreover, it commutes with inverse images, derived direct images and base change maps, since both derived direct images are computed on underlying abelian sheaves (Lemma \hyperref[bp292]{4.2}). Hence the base change map for \(F\) is an isomorphism.

If \(nF=0\), we argue by induction on the number of prime factors of \(n\). For a prime \(p\mid n\), the outer terms of \(0\to F[p]\to F\xrightarrow{p}pF\to0\) are killed by \(p\) and by \(n/p\), so the five lemma applies to the long exact cohomology sequences. Finally, a torsion sheaf is the filtered colimit \(\varinjlim_nF[n]\). Since (4.7) applies to the qcqs morphisms \(f\) and \(f'\), the base change map is an isomorphism for every torsion sheaf of \(\Lambda\)-modules. The truncation argument at the end of the proof of Theorem \hyperref[bp329]{4.17} then passes from sheaves to bounded below complexes with torsion cohomology, in particular to \(K_{\acute{e}t}\).

To return to the pro-étale site we apply \(\nu_{Y'}^*\). Since \(f\) and \(f'\) are proper, hence qcqs, the comparison of direct images in Theorem \hyperref[bp113]{2.38} applies. Using it together with the comparison of inverse images \(g_{\mathrm{pro}\acute{e}t}^*\nu_Y^*\simeq\nu_{Y'}^*g_{\acute{e}t}^*\), we identify \(\nu_{Y'}^*g^*Rf_*K_{\acute{e}t}\) and \(\nu_{Y'}^*Rf'_*g'^*K_{\acute{e}t}\) with \(g^*Rf_*K\) and \(Rf'_*g'^*K\) in \(D^+(Y'_{\mathrm{pro}\acute{e}t},\Lambda)\). By the last assertion of Lemma \hyperref[bp304]{4.7}, under these identifications the base change map of \(K\) is the inverse image on \(Y'_{\mathrm{pro}\acute{e}t}\) of the étale base change map of \(K_{\acute{e}t}\). We have shown that the latter is an isomorphism.
\end{proof}

% text: prose-04-37
\phantomsection\label{bp385}%
Taking for \(g\) a geometric point, we answer the question at the start of the chapter. The only additional point is to make sense of a stalk on the pro-étale site, where the stalk is the value of \(\bar y^*Rf_*K\) on the point \(\operatorname{Spec}\Omega\), whose topos is that of condensed sets. We check that for a classical complex this value is the étale stalk.

% object: corollary-cohomology-geometric-fibres (blueprint/chap-Proper-base-change/sec-Pro-etale-complexes/corollary-cohomology-geometric-fibres/_section.md)
\phantomsection\label{bp386}%
\begin{bpstatementit}{Corollary~4.34}{ (Cohomology of geometric fibres)}
Under the hypotheses of Theorem \hyperref[bp382]{4.33}, for every geometric point \(\bar y:\operatorname{Spec}(\Omega)\to Y\) there is a canonical isomorphism

\[
(Rf_*K)_{\bar y}\simeq
R\Gamma\bigl(X_{\bar y,\mathrm{pro}\acute{e}t},K|_{X_{\bar y}}\bigr).
\]

For a classical torsion sheaf \(F\) of \(\Lambda\)-modules this says

\[
(R^qf_*F)_{\bar y}\simeq
H^q_{\mathrm{pro}\acute{e}t}(X_{\bar y},F|_{X_{\bar y}})\qquad(q\geq0).
\]
\end{bpstatementit}

% proof: proof/corollary-cohomology-geometric-fibres
\begin{proof}[{Proof}]\phantomsection\label{bp387}%
We apply the theorem to \(g=\bar y\) and take global sections over the point. The right side becomes the cohomology of the fibre, and we check that the left side is the étale stalk.

Take \(Y'=\operatorname{Spec}(\Omega)\) in the theorem and apply the functor of global sections on \(Y'_{\mathrm{pro}\acute{e}t}\). By Proposition \hyperref[bp153]{2.50} the topos of \(Y'_{\mathrm{pro}\acute{e}t}\) is that of condensed sets, where global sections is evaluation at the point. This functor is exact, because every surjection of profinite sets onto a point has a section. On the right we therefore get the cohomology of the fibre. On the left, \((Rf_*K)_{\bar y}\) means \(\Gamma(Y'_{\mathrm{pro}\acute{e}t},\bar y^*Rf_*K)\), and by the proof of the theorem we have \(Rf_*K\simeq\nu_Y^*Rf_*K_{\acute{e}t}\). Since \(\bar y^*\nu_Y^*\simeq\nu_{Y'}^*\bar y^*\), the complex \(\bar y^*Rf_*K\) is the constant condensed complex on the étale stalk \((Rf_*K_{\acute{e}t})_{\bar y}\), so its global sections are that stalk. Finally, by the theorem on the étale site, this stalk is the étale cohomology of the fibre.
\end{proof}


% object: sec-Finite-coefficients (blueprint/chap-Proper-base-change/sec-Finite-coefficients/_section.md), content hidden
\subsection*{\texorpdfstring{4.7\quad Finite coefficients}{4.7 Finite coefficients}}
\phantomsection\label{bp388}%
\addcontentsline{toc}{subsection}{\texorpdfstring{4.7\quad Finite coefficients}{4.7 Finite coefficients}}
% text: prose-04-38
\phantomsection\label{bp389}%
Let \(A\) be a finite admissible coefficient ring and \(p:X\to Y\) a proper morphism of separated schemes of finite type over the standing field \(k\). For \(C\in D^b_{\mathrm{ctf}}(X,A)\), proper base change identifies the geometric stalks of \(Rp_*C\) with the cohomology of the fibres. From this identification and the classical finiteness theorem we obtain that \(Rp_*C\) is constructible, with a bound on its Tor amplitude. Both facts are needed for the passage to integral and rational coefficients below, whereas direct sums in the unbounded localizing categories are treated in Section \hyperref[bp449]{5.3}.

% object: proposition-proper-direct-image-finite-coefficients (blueprint/chap-Proper-base-change/sec-Finite-coefficients/proposition-proper-direct-image-finite-coefficients/_section.md)
\phantomsection\label{bp390}%
\begin{bpstatementit}{Proposition~4.35}{ (Proper direct image for finite coefficients)}
Let \(A\) be a finite admissible ring, and let \(p:X\to Y\) be a proper morphism of separated schemes of finite type over \(k\). Then \(Rp_*\) maps \(D^b_{\mathrm{ctf}}(X,A)\) into \(D^b_{\mathrm{ctf}}(Y,A)\). More precisely, if \(C\in D^b_{\mathrm{ctf}}(X,A)\) has Tor amplitude in \([a,b]\) and \(\operatorname{cd}_\ell(X_{\bar y})\leq d\) for all geometric points \(\bar y\) of \(Y\), then \(Rp_*C\) has Tor amplitude in \([a,b+d]\).
\end{bpstatementit}

% proof: proof/proposition-proper-direct-image-finite-coefficients
\begin{proof}[{Proof}]\phantomsection\label{bp391}%
By Corollary \hyperref[bp132]{2.44} the étale and pro-étale categories \(D^b_{\mathrm{ctf}}\) agree, and by Theorem \hyperref[bp113]{2.38} so do the proper direct images of their bounded below complexes. We can therefore transport bounded constructibility and the Tor bound from the étale site.

On the étale site, let \(C\in D^b_{\mathrm{ctf}}(X_{\acute{e}t},A)\). The proper direct image of a constructible torsion sheaf is constructible [Stacks, \href{https://stacks.math.columbia.edu/tag/0GL0}{0GL0}], and \(\operatorname{cd}_\ell(p)<\infty\) [Stacks, \href{https://stacks.math.columbia.edu/tag/0F10}{0F10}]. Hence \(Rp_*C\in\Dbc(Y_{\acute{e}t},A)\).

To check finite Tor amplitude, suppose that \(C\) has Tor amplitude \([a,b]\), and let \(d\) be an integer with \(\operatorname{cd}_\ell(X_{\bar y})\leq d\) for all geometric points \(\bar y\) of \(Y\). Since derived tensor products are computed stalkwise, it suffices to show that each stalk \((Rp_*C)_{\bar y}\) has Tor amplitude in \([a,b+d]\) as a complex of \(A\)-modules. In other words, \((Rp_*C)_{\bar y}\otimes^{\mathbf L}_AN=((Rp_*C)\otimes_A^{\mathbf L}N_Y)_{\bar y}\) should have cohomology in \([a,b+d]\) for every \(A\)-module \(N\). By proper base change and the projection formula for constant modules [Stacks, \href{https://stacks.math.columbia.edu/tag/0F12}{0F12}], this stalk is

\[
R\Gamma\bigl(X_{\bar y},
C|_{X_{\bar y}}\otimes_A^{\mathbf L}N_{X_{\bar y}}\bigr).
\]

The complex \(C|_{X_{\bar y}}\otimes_A^{\mathbf L}N_{X_{\bar y}}\) has cohomology in \([a,b]\), because restriction to the fibre preserves Tor amplitude. Since \(R\Gamma(X_{\bar y},-)\) preserves \(D^{\geq a}\) and has cohomological dimension at most \(d\), it keeps the lower bound and raises the upper bound by at most \(d\). Thus \(Rp_*C\in D^b_{\mathrm{ctf}}(Y_{\acute{e}t},A)\), with Tor amplitude in \([a,b+d]\).
\end{proof}


% object: sec-ell-adic-coefficients (blueprint/chap-Proper-base-change/sec-ell-adic-coefficients/_section.md), content hidden
\subsection*{\texorpdfstring{4.8\quad $\ell$-adic coefficients}{4.8 ell-adic coefficients}}
\phantomsection\label{bp392}%
\addcontentsline{toc}{subsection}{\texorpdfstring{4.8\quad $\ell$-adic coefficients}{4.8 ell-adic coefficients}}
% text: prose-04-39
\phantomsection\label{bp393}%
Pro-étale proper base change (Theorem \hyperref[bp382]{4.33}) concerns classical complexes with torsion cohomology sheaves. For a finite ring \(A\) it therefore applies, with \(\Lambda=A\), to every \(K\in D^b_{\mathrm{ctf}}(X,A)\) for schemes of finite type over \(k\), because bounded constructible complexes on quasi-compact quasi-separated schemes are classical (Theorem \hyperref[bp129]{2.43}). For \(\mathcal O_E\), however, the theorem gives proper base change only at each finite level \(\mathcal O_E/\pi_E^n\), and we must pass to the limit. Later statements of this kind make this passage once and for all, by the reduction principle of Section \hyperref[bp632]{7.3}; here we make it directly.

The difficulty is completeness: the inverse image \(g^*:D(Y_{\mathrm{pro}\acute{e}t},\mathcal O_E)\to D(Y'_{\mathrm{pro}\acute{e}t},\mathcal O_E)\), which is exact on sheaves of modules, need not preserve derived completeness. We therefore complete after pulling back.

% object: definition-completed-inverse-image (blueprint/chap-Proper-base-change/sec-ell-adic-coefficients/definition-completed-inverse-image/_section.md)
\phantomsection\label{bp394}%
\begin{bpstatement}{Definition~4.36}{ (Completed inverse image)}
Let \(g:Y'\to Y\) be a morphism. For \(M\in D(Y_{\mathrm{pro}\acute{e}t},\mathcal O_E)\), the \emph{completed inverse image} is

\[
\hat g^*M
=R\varprojlim_n g^*(M\otimes^{\mathbf L}_{\mathcal O_{E,Y}}
(\mathcal O_E/\pi_E^n)_Y)
\in D(Y'_{\mathrm{pro}\acute{e}t},\mathcal O_E).
\]
\end{bpstatement}

% text: prose-04-40
\phantomsection\label{bp395}%
The natural map \(g^*M\to\hat g^*M\) is the completion map of \(g^*M\), since \(g^*\) commutes with the reductions. For constructible \(M\) the completion does nothing. Indeed, let \(Y\) and \(Y'\) be of finite type over \(k\) and \(M\in D^b_{\mathrm{ctf}}(Y,\mathcal O_E)\). Then \(g^*M\) lies in \(\Dbc(Y',\mathcal O_E)\) by Lemma lemma-integral-models, hence it is derived complete by Proposition \hyperref[bp226]{3.11} ((1)\(\Rightarrow\)(2)), and \(\hat g^*M=g^*M\). However, until we know that proper direct image preserves constructibility, we can only state base change with the completed inverse image.

In the proof we use the following map, which also appears in the projection formula.

% object: definition-projection-formula-map (blueprint/chap-Proper-base-change/sec-ell-adic-coefficients/definition-projection-formula-map/_section.md)
\phantomsection\label{bp396}%
\begin{bpstatement}{Definition~4.37}{ (Projection formula map)}
Let \(p:X\to Y\) be a morphism, \(K\in D(X_{\mathrm{pro}\acute{e}t},\Lambda)\) and \(M\in D(Y_{\mathrm{pro}\acute{e}t},\Lambda)\). The \emph{projection formula map}

\[
Rp_*K\otimes^{\mathbf L}M\longrightarrow Rp_*(K\otimes^{\mathbf L}p^*M)
\tag{4.14}
\]

is the adjoint, under \(p^*\dashv Rp_*\), of the composite \(p^*(Rp_*K\otimes^{\mathbf L}M)\simeq p^*Rp_*K\otimes^{\mathbf L}p^*M
\to K\otimes^{\mathbf L}p^*M\) of the monoidal isomorphism of \(p^*\) and the counit \(p^*Rp_*K\to K\) tensored with \(p^*M\).
\end{bpstatement}

% object: proposition-proper-base-change-complete-integral (blueprint/chap-Proper-base-change/sec-ell-adic-coefficients/proposition-proper-base-change-complete-integral/_section.md)
\phantomsection\label{bp397}%
\begin{bpstatementit}{Proposition~4.38}{ (Proper base change for complete complexes)}
Let \(p:X\to Y\) be proper, let \(g:Y'\to Y\) be any morphism, and let \(p':X'\to Y'\) and \(g':X'\to X\) be the projections from \(X'=X\times_YY'\). Fix a finite extension \(E/\mathbf Q_\ell\). Suppose that \(K\in D(X_{\mathrm{pro}\acute{e}t},\mathcal O_E)\) is derived \(\pi_E\)-complete and that each reduction

\[
K_n=K\otimes^{\mathbf L}_{\mathcal O_{E,X}}
(\mathcal O_E/\pi_E^n)_X
\]

belongs to \(D^+_{\mathrm{class}}(X_{\mathrm{pro}\acute{e}t},\mathcal O_E/\pi_E^n)\). Then \(Rp_*K\) is derived \(\pi_E\)-complete, and there are canonical isomorphisms

\[
\hat g^*Rp_*K\simeq R\varprojlim_n g^*Rp_*K_n,
\qquad
Rp'_*\hat g'^*K\simeq R\varprojlim_n Rp'_*g'^*K_n,
\]

and under them the derived limit of the base change maps \(g^*Rp_*K_n\to Rp'_*g'^*K_n\) is an isomorphism \(\hat g^*Rp_*K\xrightarrow{\ \sim\ }Rp'_*\hat g'^*K\).
\end{bpstatementit}

% proof: proof/proposition-proper-base-change-complete-integral
\begin{proof}[{Proof}]\phantomsection\label{bp398}%
Theorem \hyperref[bp382]{4.33} applies at each finite level, so the point is that \(Rp_*\) and the completed inverse images are compatible with the passage to the derived limit of the reductions.

The reductions of \(Rp_*K\) are controlled by the projection formula. The projection formula map (4.14) for \(M=(\mathcal O_E/\pi_E^n)_Y\) is a map

\[
(Rp_*K)\otimes^{\mathbf L}(\mathcal O_E/\pi_E^n)_Y
\longrightarrow Rp_*K_n,
\]

which we regard in \(D(Y_{\mathrm{pro}\acute{e}t},\mathcal O_E/\pi_E^n)\). We claim that it is an isomorphism. Since restriction of scalars to \(\mathcal O_E\) is conservative, we may check this after forgetting the \(\mathcal O_E/\pi_E^n\)-structure. Now the ring \(\mathcal O_E/\pi_E^n\) has the free resolution \([\mathcal O_E\xrightarrow{\pi_E^n}\mathcal O_E]\), so by Corollary \hyperref[bp161]{2.52} the source is the cone of \(\pi_E^n\) on \(Rp_*K\), and the target is \(Rp_*\) of the cone of \(\pi_E^n\) on \(K\). These agree because \(Rp_*\) is triangulated and \(\mathcal O_E\)-linear. Moreover, the resolutions for \(n+1\) and \(n\) are related by the chain map \((\pi_E,\mathrm{id})\), and applying \(Rp_*\) to it shows that these comparisons form an isomorphism of towers. Since all these constructions can be made on complexes with functorial resolutions, the comparisons define a map in the derived category of towers of Definition \hyperref[bp188]{2.61}.

Since \(K\) is derived complete, it is the derived limit of its reductions \(K_n\), and since the functor \(Rp_*\) is a right adjoint, it commutes with derived limits, and we obtain

\[
Rp_*K\simeq Rp_*R\varprojlim_nK_n
\simeq R\varprojlim_nRp_*K_n.
\]

By the reduction comparisons, the right side is the derived limit of the reductions \((Rp_*K)\otimes^{\mathbf L}(\mathcal O_E/\pi_E^n)_Y\), that is, the derived completion of \(Rp_*K\). Hence the completion map of \(Rp_*K\) is an isomorphism, and \(Rp_*K\) is derived complete. The definition of completed inverse image and the same right-adjoint property for \(Rp'_*\) now give

\[
\begin{aligned}
\hat g^*Rp_*K
&\simeq R\varprojlim_n g^*Rp_*K_n,\\
Rp'_*\hat g'^*K
&\simeq R\varprojlim_n Rp'_*g'^*K_n.
\end{aligned}
\]

The base change maps \(g^*Rp_*K_n\to Rp'_*g'^*K_n\) are natural in \(K_n\), so they form a map of towers, and each is an isomorphism by Theorem \hyperref[bp382]{4.33}. Its cone therefore has zero cohomology at every level, and applying \(R\varprojlim\) gives the asserted isomorphism.
\end{proof}

% text: prose-04-41
\phantomsection\label{bp399}%
When the pullbacks are already complete, the completed base change isomorphism is the ordinary one.

% object: lemma-completed-ordinary-base-change (blueprint/chap-Proper-base-change/sec-ell-adic-coefficients/lemma-completed-ordinary-base-change/_section.md)
\phantomsection\label{bp400}%
\begin{bpstatementit}{Lemma~4.39}{ (Completed and ordinary base change)}
In the setting of Proposition \hyperref[bp397]{4.38}, the square formed by the base change map \(g^*Rp_*K\to Rp'_*g'^*K\) of (4.5), the completed base change isomorphism \(\hat g^*Rp_*K\to Rp'_*\hat g'^*K\), and the completion maps \(g^*Rp_*K\to\hat g^*Rp_*K\) and \(Rp'_*g'^*K\to Rp'_*\hat g'^*K\) commutes. In particular, if \(g^*Rp_*K\) and \(g'^*K\) are derived complete, then the base change map \(g^*Rp_*K\to Rp'_*g'^*K\) is an isomorphism.
\end{bpstatementit}

% proof: proof/lemma-completed-ordinary-base-change
\begin{proof}[{Proof}]\phantomsection\label{bp401}%
Both maps are built from the counit of \(p^*\dashv Rp_*\), as in Definition \hyperref[bp302]{4.6}, and this counit is natural with respect to the quotient maps and the chain maps \((\pi_E,\mathrm{id})\) relating the reductions. Hence the square commutes.

Now suppose that \(g^*Rp_*K\) and \(g'^*K\) are derived complete. Then the first completion map is an isomorphism, and the second is \(Rp'_*\) applied to the completion map of \(g'^*K\), so it is an isomorphism as well. By the commutative square, the base change map is then a composite of isomorphisms.
\end{proof}

% text: prose-04-42
\phantomsection\label{bp402}%
We will show that \(Rp_*K\) is again constructible when \(K\) is constructible. Once this is known, the completed inverse images in Proposition \hyperref[bp397]{4.38} agree with \(g^*\) and \(g'^*\), and the proposition becomes the ordinary base change isomorphism

\[
g^*Rp_*K\xrightarrow{\sim}Rp'_*g'^*K
\qquad(K\in\Dbc(X,\mathcal O_E)).
\]


% object: sec-Proper-direct-image (blueprint/chap-Proper-base-change/sec-Proper-direct-image/_section.md), content hidden
\subsection*{\texorpdfstring{4.9\quad Proper direct image}{4.9 Proper direct image}}
\phantomsection\label{bp403}%
\addcontentsline{toc}{subsection}{\texorpdfstring{4.9\quad Proper direct image}{4.9 Proper direct image}}
% text: prose-04-43
\phantomsection\label{bp404}%
Let \(p:X\to Y\) be proper. We need to know, for every admissible coefficient ring, that \(Rp_*\) preserves \(D^b_{\mathrm{ctf}}\) and satisfies proper base change and the projection formula there, because the construction of \(Rf_!=Rp_*\circ j_!\) in Section \hyperref[bp459]{5.4} starts from these facts. For \(Y=\operatorname{Spec}k\) with \(k\) algebraically closed, they say that the cohomology groups \(H^i(X,\overline{\mathbf Q}_\ell)\) of a proper \(X\) are finite-dimensional and vanish outside a bounded range of degrees.

For finite coefficients, constructibility and the Tor bound are Proposition \hyperref[bp390]{4.35}, and the other admissible rings reduce to this case.

For \(\mathcal O_E\), the complex \(Rp_*K\) is derived complete, and its reduction modulo \(\pi_E\) is \(Rp_*\) of the reduction of \(K\), so the integral constructibility criterion (Proposition \hyperref[bp226]{3.11}) and derived Nakayama (Lemma \hyperref[bp190]{2.62}) reduce everything to torsion coefficients. For \(E\) and \(\overline{\mathbf Q}_\ell\) it suffices to treat the generators \((a_!L_0)[1/\ell]\). Indeed, on bounded below complexes \(Rp_*\) commutes with inverting \(\ell\) and with extension of scalars to finite extensions and to \(\overline{\mathbf Q}_\ell\), since both are filtered colimits (Corollary \hyperref[bp86]{2.29}).

% object: proposition-proper-direct-image-constructible-complexes (blueprint/chap-Proper-base-change/sec-Proper-direct-image/proposition-proper-direct-image-constructible-complexes/_section.md)
\phantomsection\label{bp405}%
\begin{bpstatementit}{Proposition~4.40}{ (Proper direct image of constructible complexes)}
Let \(p:X\to Y\) be proper and let \(\Lambda\) be admissible.

\begin{enumerate}
\item \(Rp_*:D(X_{\mathrm{pro}\acute{e}t},\Lambda)\to D(Y_{\mathrm{pro}\acute{e}t},\Lambda)\) maps \(D^b_{\mathrm{ctf}}(X,\Lambda)\) into \(D^b_{\mathrm{ctf}}(Y,\Lambda)\).
\item For a cartesian square with \(g:Y'\to Y\), \(g':X'\to X\) and \(p':X'\to Y'\), the base change map (4.5) \(g^*Rp_*\to Rp'_*g'^*\) is an isomorphism of functors \(D^b_{\mathrm{ctf}}(X,\Lambda)\to D^b_{\mathrm{ctf}}(Y',\Lambda)\).
\item The projection formula map (4.14) is an isomorphism for \(K\in D^b_{\mathrm{ctf}}(X,\Lambda)\) and \(M\in D^b_{\mathrm{ctf}}(Y,\Lambda)\).
\end{enumerate}
\end{bpstatementit}

% proof: proof/proposition-proper-direct-image-constructible-complexes
\begin{proof}[{Proof}]\phantomsection\label{bp406}%
\emph{Finite coefficients.} Part (1), including the Tor bound, is Proposition \hyperref[bp390]{4.35}, and proper base change holds by Theorem \hyperref[bp382]{4.33} with \(\Lambda=A\). The projection formula, finally, is the classical torsion formula [Stacks, \href{https://stacks.math.columbia.edu/tag/0F0G}{0F0G}] in \(D(Y_{\acute{e}t},A)\), transported to \(D(Y_{\mathrm{pro}\acute{e}t},A)\) by pullback from the étale site.

\emph{Integral coefficients.} Let \(K\in\Dbc(X,\mathcal O_E)\). We first show that \(Rp_*K\in\Dbc(Y,\mathcal O_E)\), by computing its reduction as in (4.15) below. By Proposition \hyperref[bp226]{3.11} and Lemma \hyperref[bp208]{3.5}, \(K\) is derived complete with bounded classical finite reductions. Hence Proposition \hyperref[bp397]{4.38} applies. It shows that \(Rp_*K\) is derived complete and gives a canonical comparison of reductions

\[
\rho_K:(Rp_*K)\otimes^{\mathbf L}\mathcal O_{E,Y}/\pi_E
\xrightarrow{\ \sim\ } Rp_*(K\otimes^{\mathbf L}\mathcal O_{E,X}/\pi_E).
\tag{4.15}
\]

The target of \(\rho_K\) is constructible and bounded. Indeed, it is the direct image of a torsion complex pulled back from \(X_{\acute{e}t}\) [Stacks, \href{https://stacks.math.columbia.edu/tag/09B1}{09B1}], so it is the pullback to \(Y_{\mathrm{pro}\acute{e}t}\) of the direct image in \(D(Y_{\acute{e}t},\mathbf F_E)\) [Stacks, \href{https://stacks.math.columbia.edu/tag/09A6}{09A6}]. The latter is constructible by the finiteness theorem for proper direct image with torsion coefficients [Stacks, \href{https://stacks.math.columbia.edu/tag/0GL0}{0GL0}], and bounded because \(p\) has finite cohomological dimension [Stacks, \href{https://stacks.math.columbia.edu/tag/0F10}{0F10}]. Now Proposition \hyperref[bp226]{3.11} gives \(Rp_*K\in\Dbc(Y,\mathcal O_E)\).

Proper base change holds by Lemma \hyperref[bp400]{4.39}, which applies because \(g^*Rp_*K\) and \(g'^*K\) are derived complete, as inverse image preserves integral constructibility (Lemma lemma-integral-models).

The projection formula is tested modulo \(\pi_E\) in the same way. Let \(C\) be the cone of the projection formula map for \(K\in\Dbc(X,\mathcal O_E)\) and \(M\in\Dbc(Y,\mathcal O_E)\). Its source \(Rp_*K\otimes^{\mathbf L}M\) and target \(Rp_*(K\otimes^{\mathbf L}p^*M)\) are integrally constructible, by Lemma lemma-integral-models and the fact just proved that \(Rp_*\) maps \(\Dbc(X,\mathcal O_E)\) into \(\Dbc(Y,\mathcal O_E)\). As objects of \(\Dbc\) they are derived complete, and hence so is \(C\). Derived Nakayama (Lemma \hyperref[bp190]{2.62}) therefore reduces us to showing \(C/\pi_E=0\). In other words, we must show that the reduction of the projection formula map is an isomorphism.

This reduction is the torsion projection formula map, which is an isomorphism by [Stacks, \href{https://stacks.math.columbia.edu/tag/0F0G}{0F0G}]. To identify the two maps, note that the sources and targets correspond by \(\rho_K\) and associativity of extension of scalars, and the maps agree because both are adjoint to the counit \(p^*Rp_*K\to K\), as is \(\rho_K\).

\emph{Rational coefficients.} Here \(Rp_*\) commutes with the operations that produce the generators of Lemma lemma-integral-models from integral complexes. For \(K\in\Dbc(X,\mathcal O_E)\), the complex \(K[1/\ell]\) is the colimit of \(K\xrightarrow{\ell}K\xrightarrow{\ell}\cdots\), and since this system has a common lower bound, Corollary \hyperref[bp86]{2.29} shows that

\[
Rp_*(K[1/\ell])\simeq(Rp_*K)[1/\ell].
\]

Likewise \(Rp_*\) commutes with extension of scalars to a finite \(E'/E\), which is a finite direct sum by Lemma lemma-inverse-image-coefficient-sheaves. By Corollary \hyperref[bp86]{2.29} again, it also commutes with the filtered colimit \(\QbarX\otimes_{E_X}-=\varinjlim_{E'}E'_X\otimes_{E_X}-\). By the integral case, it follows that \(Rp_*\) maps the generators into \(\Dbc(Y,E)\) and \(\Dbc(Y,\overline{\mathbf Q}_\ell)\). Since \(Rp_*\) is exact and \(\Dbc(Y,E)\) and \(\Dbc(Y,\overline{\mathbf Q}_\ell)\) are triangulated (Corollary \hyperref[bp253]{3.20}), \(Rp_*\) preserves \(\Dbc(-,E)\) and \(\Dbc(-,\overline{\mathbf Q}_\ell)\).

For proper base change and the projection formula it is enough to treat the generators. Indeed, both sides of each map are exact functors in each variable, so the objects for which the map is an isomorphism form a triangulated subcategory. On the generators, the maps are obtained from the integral isomorphisms by inverting \(\ell\) and extending scalars. These operations commute with \(Rp_*\) by the preceding paragraph, and with \(g^*\) and \(\otimes^{\mathbf L}\) because they are given by tensoring with the coefficient sheaves. Moreover, the maps so obtained are the base change and projection formula maps, since these scalar comparisons commute with the counits of \(p^*\dashv Rp_*\).
\end{proof}



% object: chap-Construction-of-f (blueprint/chap-Construction-of-f/_section.md), content hidden
\section*{\texorpdfstring{5\quad Construction of $f^!$}{5 Construction of f!}}
\phantomsection\label{bp407}%
\addcontentsline{toc}{section}{\texorpdfstring{5\quad Construction of $f^!$}{5 Construction of f!}}
% text: prose-05-01
\phantomsection\label{bp408}%
Verdier duality, Poincaré duality and the trace formula all rest on a right adjoint \(f^!\) of \(Rf_!\). Indeed, Verdier duality (Chapter \bpout{10}) is duality with respect to the dualizing complex \(a_X^!\Lambda\), where \(a_X:X\to\operatorname{Spec}k\) is the structure map, and Poincaré duality follows from it once smooth purity (Theorem \bpout{9.3}) computes \(f^!\) for smooth morphisms. Similarly, the local terms of the trace formula (Chapter \bpout{12}) are built from \(f^!\) for the two projections of a correspondence.

The functor \(Rf_!\) itself is easy to write down. A separated morphism of finite type \(f:X\to Y\) has a compactification \(f=p\circ j\), with \(j:X\hookrightarrow\overline X\) an open immersion and \(p:\overline X\to Y\) proper (Theorem \bpout{A.2}), and we put \(Rf_!=Rp_*\circ j_!\). This definition stays within constructible complexes because, for every admissible \(\Lambda\), the functor \(Rp_*\) maps \(D^b_{\mathrm{ctf}}(\overline X,\Lambda)\) into \(D^b_{\mathrm{ctf}}(Y,\Lambda)\). The resulting functor satisfies proper base change and the projection formula (Proposition \hyperref[bp405]{4.40}).

The right adjoint of \(Rf_!\) is another matter. The factor \(j_!\) has the right adjoint \(j^*\), but for \(Rp_*\) with \(p\) proper no formula for a right adjoint is available. Even where a formula exists, it is not computed pointwise: for a closed immersion \(i\) the functor \(i^!\) takes sections with support in the closed subscheme, so the value of \(i^!K\) at a point depends on \(K\) near the point, not only on its stalk there.

Grothendieck and Hartshorne built the coherent \(f^!\) from local formulas and then had to glue them. Neeman noticed, instead, that \(D_{\mathrm{qc}}(X)\) is generated by compact perfect complexes and that \(Rf_*\) preserves direct sums, so that Brown representability produces \(f^!\) with no formula. We follow his approach, with \(D^b_{\mathrm{ctf}}\) in place of the perfect complexes.

The theorem we use is the following. Let \(\mathcal T\) be a triangulated category with arbitrary direct sums, and suppose that \(\mathcal T\) is compactly generated: there is a set \(\mathcal G\) of compact objects \(G\) (that is, \(\Hom(G,-)\) commutes with direct sums) such that \(K=0\) whenever \(\Hom(G[n],K)=0\) for all \(G\in\mathcal G\) and \(n\in\mathbf Z\). Then every exact functor out of \(\mathcal T\) that preserves direct sums has an exact right adjoint [Stacks, \href{https://stacks.math.columbia.edu/tag/0A8F}{0A8F}, \href{https://stacks.math.columbia.edu/tag/0A8G}{0A8G}]. The category \(D^b_{\mathrm{ctf}}(X,\Lambda)\) is not closed under infinite direct sums. We therefore apply the theorem to the localizing category \(\Loc(D^b_{\mathrm{ctf}}(X,\Lambda))\) that it generates (Section \hyperref[bp409]{5.1}). The result is Theorem \hyperref[bp506]{5.31}: for every admissible \(\Lambda\) and every separated morphism of finite type \(f:X\to Y\) there is an exact functor

\[
f^!:\Loc(D^b_{\mathrm{ctf}}(Y,\Lambda))\to\Loc(D^b_{\mathrm{ctf}}(X,\Lambda)),
\qquad
\Hom(Rf_!M,N)\simeq\Hom(M,f^!N).
\]

For finite coefficients nothing is lost by passing to the localizing category, since \(\Loc(D^b_{\mathrm{ctf}}(X,A))\) is equivalent to the classical derived category \(D(X_{\acute{e}t},A)\) (Proposition \hyperref[bp454]{5.15}).

Everything therefore rests on the two hypotheses of the theorem: that the objects of \(D^b_{\mathrm{ctf}}(X,\Lambda)\) are compact in \(\Loc(D^b_{\mathrm{ctf}}(X,\Lambda))\), and that \(Rp_*\) preserves direct sums for proper \(p\). Neither is formal, because the summands of a direct sum may live in degrees going off to \(-\infty\). If the Ext groups out of a sheaf do not vanish in high degrees, compactness fails even for the constant sheaf. Here is an example, outside our conventions because \(\mathbf R\) is neither finite nor algebraically closed. Take \(\ell=2\) and \(X=\operatorname{Spec}\mathbf R\). Then sheaves on \(X_{\acute{e}t}\) are modules over \(\operatorname{Gal}(\mathbf C/\mathbf R)=\mathbf Z/2\), and \(\Ext^n(\mathbf F_2,\mathbf F_2)=H^n(\mathbf Z/2,\mathbf F_2)=\mathbf F_2\) for all \(n\geq0\). Since products of modules are exact and the complexes \(\mathbf F_2[n]\) live in distinct degrees, the sum \(\bigoplus_{n\geq0}\mathbf F_2[n]\) in \(D(X_{\acute{e}t},\mathbf F_2)\) is also the product. Hence

\[
\Hom\Bigl(\mathbf F_2,\bigoplus_{n\ge0}\mathbf F_2[n]\Bigr)
=\prod_{n\ge0}\mathbf F_2
\ne\bigoplus_{n\ge0}\mathbf F_2
=\bigoplus_{n\ge0}\Hom(\mathbf F_2,\mathbf F_2[n]).
\]

By our conventions (Section \bpout{1.9}) the base field is finite or algebraically closed, so its \(\ell\)-cohomological dimension is finite (Lemma \hyperref[bp139]{2.46}). For finite coefficients this rules out such examples. In that case both hypotheses reduce to the fact that classical étale cohomology commutes with direct sums of unbounded complexes (Section \hyperref[bp449]{5.3}). For \(\Lambda=E\) or \(\overline{\mathbf Q}_\ell\) a bound on the base is not enough, and we prove instead a bound on \(\Ext\) that depends only on the source (Lemma \hyperref[bp426]{5.6}).

% object: sec-Existence-and-finiteness-of-f (blueprint/chap-Construction-of-f/sec-Existence-and-finiteness-of-f/_section.md), content hidden
\subsection*{\texorpdfstring{5.1\quad Existence and finiteness of $f^!$}{5.1 Existence and finiteness of f!}}
\phantomsection\label{bp409}%
\addcontentsline{toc}{subsection}{\texorpdfstring{5.1\quad Existence and finiteness of $f^!$}{5.1 Existence and finiteness of f!}}
% text: prose-05-02
\phantomsection\label{bp410}%
We construct \(f^!\) by Brown representability, as outlined at the beginning of this chapter, without writing it down. Brown representability needs a category with arbitrary direct sums, and \(D^b_{\mathrm{ctf}}(X,\Lambda)\) is not closed under them. We therefore enlarge it to the smallest category that has them and show that \(D^b_{\mathrm{ctf}}(X,\Lambda)\) compactly generates it.

% object: definition-localizing-subcategory (blueprint/chap-Construction-of-f/sec-Existence-and-finiteness-of-f/definition-localizing-subcategory/_section.md)
\phantomsection\label{bp411}%
\begin{bpstatement}{Definition~5.1}{ (The localizing subcategory generated by $\Dctf(X,\Lambda)$)}
For admissible \(\Lambda\), let \(\Loc(D^b_{\mathrm{ctf}}(X,\Lambda))\) be the smallest full triangulated subcategory of \(D(X_{\mathrm{pro}\acute{e}t},\Lambda)\) containing \(D^b_{\mathrm{ctf}}(X,\Lambda)\) and closed under arbitrary direct sums.
\end{bpstatement}

% text: prose-05-03
\phantomsection\label{bp412}%
For integral and rational coefficients, this is also \(\Loc(\Dbc(X,\Lambda))\), by Proposition \hyperref[bp232]{3.12}.

For \(\mathcal O_E\) the objects of \(D^b_{\mathrm{ctf}}(X,\mathcal O_E)\) are derived complete, but the localizing subcategory they generate also contains complexes that are not. For this reason the proof of Theorem theorem-exceptional-inverse-image must show, for integral constructible \(K\), that \(f^!K\) is derived complete.

% object: theorem-exceptional-inverse-image (blueprint/chap-Construction-of-f/sec-Existence-and-finiteness-of-f/theorem-exceptional-inverse-image/_section.md)
\phantomsection\label{bp413}%
\begin{bpstatementit}{Theorem~5.2}{ (The exceptional inverse image)}
Let \(\Lambda\in\{A,\mathbf Z_\ell,\mathcal O_E,\mathbf Q_\ell,E,\overline{\mathbf Q}_\ell\}\) be an admissible coefficient ring and \(f:X\to Y\) a morphism of separated schemes of finite type over \(k\).

\begin{enumerate}
\item The objects of \(D^b_{\mathrm{ctf}}(X,\Lambda)\) are compact in \(\Loc(D^b_{\mathrm{ctf}}(X,\Lambda))\) and a small skeleton compactly generates it.
\item Compactifications define an exact functor \(Rf_!:\Loc(D^b_{\mathrm{ctf}}(X,\Lambda))\to\Loc(D^b_{\mathrm{ctf}}(Y,\Lambda))\), preserving direct sums and sending \(D^b_{\mathrm{ctf}}(X,\Lambda)\) into \(D^b_{\mathrm{ctf}}(Y,\Lambda)\).
\item The functor \(Rf_!\) has an exact right adjoint \(f^!:\Loc(D^b_{\mathrm{ctf}}(Y,\Lambda))\to\Loc(D^b_{\mathrm{ctf}}(X,\Lambda))\), obtained from Brown representability.
\item This right adjoint sends \(D^b_{\mathrm{ctf}}(Y,\Lambda)\) into \(D^b_{\mathrm{ctf}}(X,\Lambda)\), so its adjunction restricts to \(D^b_{\mathrm{ctf}}\).
\end{enumerate}
\end{bpstatementit}

% text: prose-05-04
\phantomsection\label{bp414}%
By Proposition \hyperref[bp232]{3.12}, the theorem holds with \(\Dbc\) in place of \(D^b_{\mathrm{ctf}}\) for \(\Lambda=\mathbf Z_\ell,\mathcal O_E,\mathbf Q_\ell,E,\overline{\mathbf Q}_\ell\).

For proper \(p\) we may use the trivial compactification, so that \(Rp_!=Rp_*\). By part (4), the functor \(Rp_*:D^b_{\mathrm{ctf}}(X,\Lambda)\to D^b_{\mathrm{ctf}}(Y,\Lambda)\) then has the right adjoint \(p^!\).

Parts (1)–(3) are proved in this chapter and part (4) in Chapter \hyperref[bp513]{6}. For smooth \(f\), smooth purity (Chapter \bpout{9}) gives a global isomorphism \(f^!\simeq f^*(d)[2d]\), compatible with trace and composition.


% object: sec-Compactness-of-constructible-complexes (blueprint/chap-Construction-of-f/sec-Compactness-of-constructible-complexes/_section.md), content hidden
\subsection*{\texorpdfstring{5.2\quad Compactness of constructible complexes}{5.2 Compactness of constructible complexes}}
\phantomsection\label{bp415}%
\addcontentsline{toc}{subsection}{\texorpdfstring{5.2\quad Compactness of constructible complexes}{5.2 Compactness of constructible complexes}}
% text: prose-05-05
\phantomsection\label{bp416}%
Brown representability produces \(f^!\) once the bounded constructible complexes are known to be compact, that is, once \(\Hom(C,-)\) commutes with direct sums for constructible \(C\). The difficulty is that the summands may have cohomology in degrees going off to \(-\infty\), so that \(\Hom(C,-)\) sees all the groups \(\Ext^n(C,-)\) at once. For instance, over \(\operatorname{Spec}\mathbf R\) these groups never stop, and we saw at the beginning of this chapter that compactness fails there. By contrast, over a finite field they stop, as the following example shows.

% object: example-compactness-finite-field (blueprint/chap-Construction-of-f/sec-Compactness-of-constructible-complexes/example-compactness-finite-field/_section.md)
\phantomsection\label{bp417}%
\begin{bpstatement}{Example~5.3}{ (Compactness over a finite field)}
Let \(X=\operatorname{Spec}\mathbf F_q\) and \(\Lambda=E\). Then \(\Ext^n(E,E)=H^n(X_{\mathrm{pro}\acute{e}t},E)\) is \(E\) for \(n=0,1\) and vanishes for \(n\geq2\). To see the vanishing, consider first \(\mathbf Z_\ell\) and the Milnor sequence (2.11) with \(F_m=(\mathbf Z/\ell^m)_X\). Since the field \(\mathbf F_q\) has \(\ell\)-cohomological dimension one, we have \(H^n(X,\mathbf Z/\ell^m)=0\) for \(n\ge2\) and all \(m\), while the finite groups \(H^1(X,\mathbf Z/\ell^m)\) have vanishing \(\varprojlim^1\). The Milnor sequence therefore gives \(H^n(X_{\mathrm{pro}\acute{e}t},\mathbf Z_{\ell,X})=0\) for \(n\geq2\). From here, inverting \(\ell\) and using Example \hyperref[bp178]{2.57} for \(n=1\) gives the claim for \(\mathbf Q_\ell\); for \(E\neq\mathbf Q_\ell\) we tensor with \(E\) over \(\mathbf Q_\ell\).

As in the example over \(\mathbf R\), products are exact (Proposition \hyperref[bp68]{2.23}) and the complexes \(E[n]\) live in distinct degrees. Hence \(\bigoplus_{n\ge0}E[n]=\prod_{n\ge0}E[n]\), and

\[
\Hom\Bigl(E,\bigoplus_{n\ge0}E[n]\Bigr)
=\prod_{n\ge0}\Ext^n(E,E)
=E\oplus E
=\bigoplus_{n\ge0}\Hom(E,E[n]).
\]

The product equals the sum because only two factors are nonzero. Over \(\mathbf R\) every factor is nonzero, and the product is larger.
\end{bpstatement}

% text: prose-05-06
\phantomsection\label{bp418}%
The general argument rests on the same mechanism: a bound \(d_C\), valid for all constructible targets at once, beyond which \(\Ext^q(C,-)\) vanishes. Until the end of Section \hyperref[bp415]{5.2}, \(\Lambda\) is \(E\) or \(\overline{\mathbf Q}_\ell\), and we show that bounded constructible complexes are compact in \(\Loc(\Dbc(X,\Lambda))\). Since \(D^b_{\mathrm{ctf}}(X,\Lambda)=\Dbc(X,\Lambda)\), these are the proposed compact generators. The advantage of working with \(\Dbc\) is that we can use truncations and the constructible heart, which is abelian by Lemma \hyperref[bp248]{3.18}. As for notation, \(\Hom\) and \(\Ext\) without decoration are taken in \(D(X_{\mathrm{pro}\acute{e}t},\Lambda)\), also when both arguments are constructible sheaves, and a subscript indicates the coefficient ring where it changes, as in the lattice arguments over \(\mathcal O_E\).

By dévissage it suffices to show, for a constructible sheaf \(C\) and a family \((K_\lambda)\) of objects of \(\Loc(\Dbc(X,\Lambda))\), that \(\bigoplus_\lambda\Hom(C,K_\lambda)\to\Hom(C,\bigoplus_\lambda K_\lambda)\) is bijective. Besides the uniform bound \(\Ext^q(C,N)=0\) for \(q>d_C\) and all constructible \(N\) (Lemma \hyperref[bp426]{5.6} below), we use that \(\Ext^q(C,-)\) commutes with filtered colimits of complexes that vanish below a fixed degree, in particular of sheaves (Lemma \hyperref[bp260]{3.22}, proved in Section \hyperref[bp255]{3.6}).

% object: subsec-A-uniform-bound-for-Ext (blueprint/chap-Construction-of-f/sec-Compactness-of-constructible-complexes/subsec-A-uniform-bound-for-Ext/_section.md), content hidden
\subsubsection*{\texorpdfstring{5.2.1\quad A uniform bound for Ext}{5.2.1 A uniform bound for Ext}}
\phantomsection\label{bp419}%
\addcontentsline{toc}{subsubsection}{\texorpdfstring{5.2.1\quad A uniform bound for Ext}{5.2.1 A uniform bound for Ext}}
% text: prose-05-07
\phantomsection\label{bp420}%
The idea of the bound is to count cohomological dimensions modulo \(\pi_E^n\), add one degree for the derived limit over \(n\), and invert \(\pi_E\) by passing to a colimit. For torsion coefficients the count is direct. By (3.1), \(R\!\operatorname{Hom}(a_!L,N)\) is the cohomology on the stratum \(S\) of a complex obtained from \(N\) by restriction and \(Rj_*\). Each of these steps raises degrees by a bounded amount that does not depend on \(N\). With coefficients \(E\) we cannot count in this way, because \(E\)-sheaves are not torsion and \(\operatorname{cd}_\ell\) only controls torsion sheaves. This is why we choose lattices on the strata. Inverting \(\pi_E\) at the end then needs the continuity of Ext from Section \hyperref[bp255]{3.6}.

The bound \(d_C\) on \(\Ext^q(C,-)\) is a sum of the cohomological dimensions \(\operatorname{cd}_\ell\) of Definition \hyperref[bp137]{2.45}. In the proof we compute \(R\!\operatorname{Hom}(a_!L_0,K)\), for a derived complete \(K\in D(X_{\mathrm{pro}\acute{e}t},\mathcal O_E)\), as a derived inverse limit over the reductions modulo \(\pi_E^n\). For this, the functors in (3.1) have to commute with derived inverse limits, and we use the next lemma with \(\Lambda=\mathcal O_E\).

% object: lemma-limits-sections-support (blueprint/chap-Construction-of-f/sec-Compactness-of-constructible-complexes/subsec-A-uniform-bound-for-Ext/lemma-limits-sections-support/_section.md)
\phantomsection\label{bp421}%
\begin{bpstatementit}{Lemma~5.4}{ (Limits and sections with support)}
Let \(\Lambda\) be admissible, let \(i:Z\hookrightarrow X\) be a closed immersion with open complement \(j:V\hookrightarrow X\), and let \(u:S\hookrightarrow Z\) be an open immersion. The functors \(R\underline\Gamma_Z:D(X_{\mathrm{pro}\acute{e}t},\Lambda)\to D(Z_{\mathrm{pro}\acute{e}t},\Lambda)\) and \(u^*R\underline\Gamma_Z:D(X_{\mathrm{pro}\acute{e}t},\Lambda)\to D(S_{\mathrm{pro}\acute{e}t},\Lambda)\) are left \(t\)-exact and commute with derived inverse limits of sequences, and so does tensoring with a finite locally free \(\Lambda_S\)-module.
\end{bpstatementit}

% proof: proof/lemma-limits-sections-support
\begin{proof}[{Proof}]\phantomsection\label{bp422}%
The functor \(R\underline\Gamma_Z\) is left \(t\)-exact because it is right adjoint to the exact functor \(i_*\). Restricting along the open immersion \(u\) is exact, so \(u^*R\underline\Gamma_Z\) is left \(t\)-exact as well.

For the limits, recall that the derived inverse limit of a sequence \((K_n)\) is the fibre of the map \(\prod_nK_n\to\prod_nK_n\) given by the identity minus the transition maps. It therefore suffices to show that the functors preserve products and fibres. Restriction to an open is exact and preserves limits, so this holds for \(j^*\) and \(u^*\). It holds for \(i^*\) by [Stacks, \href{https://stacks.math.columbia.edu/tag/09AA}{09AA}], and for \(Rj_*\) because \(Rj_*\) is a right adjoint. It follows for \(R\underline\Gamma_Z=i^*\fib(\mathrm{id}\to Rj_*j^*)\). Finally, tensoring with a finite locally free sheaf \(L\) is exact and is right adjoint to tensoring with \(L^\vee\), so it preserves products as well.
\end{proof}

% text: prose-05-08
\phantomsection\label{bp423}%
The main case is that of lattices. There we count cohomological dimensions for the reductions modulo \(\pi_E^n\), and we pass to the limit by Lemma \hyperref[bp421]{5.4}.

% object: lemma-ext-bound-lattices (blueprint/chap-Construction-of-f/sec-Compactness-of-constructible-complexes/subsec-A-uniform-bound-for-Ext/lemma-ext-bound-lattices/_section.md)
\phantomsection\label{bp424}%
\begin{bpstatementit}{Lemma~5.5}{ (Ext bound for lattices)}
Let \(i:Z\hookrightarrow X\) be a closed immersion with open complement \(j:X\setminus Z\hookrightarrow X\), let \(u:S\hookrightarrow Z\) be an open immersion, put \(a=i\circ u\), and let \(L_0\) be a lisse \(\mathcal O_{E,S}\)-module that is locally free of finite rank. Then for every locally closed immersion \(b:T\hookrightarrow X\) and every lisse \(\mathcal O_{E,T}\)-module \(M_0\) that is locally free of finite rank,

\[
\Ext^q_{\mathcal O_E}(a_!L_0,b_!M_0)=0
\quad(q>\operatorname{cd}_\ell(j)+\operatorname{cd}_\ell(S)+2).
\tag{5.1}
\]
\end{bpstatementit}

% proof: proof/lemma-ext-bound-lattices
\begin{proof}[{Proof}]\phantomsection\label{bp425}%
We first bound the complex in (3.1) for the torsion reductions of \(b_!M_0\) by cohomological dimensions, and then pass to the derived limit, which costs one degree. Both cohomological dimensions are finite by [Stacks, \href{https://stacks.math.columbia.edu/tag/0F10}{0F10}] and Lemma \hyperref[bp139]{2.46}.

By Lemma lemma-integral-models and (4), \(b_!M_0\) lies in \(\Dbc(X,\mathcal O_E)\), so it is derived complete. Since \(b_!\) is exact and \(M_0\) is \(\pi_E\)-torsion free, its reductions are

\[
b_!M_0\otimes^{\mathbf L}_{\mathcal O_{E,X}}\mathcal O_{E,X}/\pi_E^n
=b_!(M_0/\pi_E^n).
\]

These are constructible torsion sheaves. They are pulled back from \(X_{\acute{e}t}\), because extension by zero and closed direct image commute with this pullback [Stacks, \href{https://stacks.math.columbia.edu/tag/09AF}{09AF}, \href{https://stacks.math.columbia.edu/tag/09A6}{09A6}].

\emph{Torsion targets.} We bound the complex that appears in (3.1) when the target is one of the torsion sheaves \(b_!(M_0/\pi_E^n)\). By [Stacks, \href{https://stacks.math.columbia.edu/tag/09A6}{09A6}], \(Rj_*\) in \(D(X_{\mathrm{pro}\acute{e}t},\mathcal O_E/\pi_E^n)\) of a sheaf pulled back from the étale site is the pullback of its direct image in \(D(X_{\acute{e}t},\mathcal O_E/\pi_E^n)\). On the étale site the degrees are controlled by the cohomological dimension, and hence \(Rj_*j^*b_!(M_0/\pi_E^n)\) has cohomology in degrees \([0,\operatorname{cd}_\ell(j)]\). The fibre of the map from a sheaf to this complex lives in degrees \([0,\operatorname{cd}_\ell(j)+1]\), so the fibre formula of Lemma lemma-sections-support-closed-subscheme shows that

\[
R\underline\Gamma_Z\,b_!(M_0/\pi_E^n)\in D^{[0,\operatorname{cd}_\ell(j)+1]}(Z_{\mathrm{pro}\acute{e}t},\mathcal O_E/\pi_E^n),
\]

and its cohomology sheaves are torsion sheaves pulled back from the étale site, though not necessarily constructible. Restricting to \(S\) and tensoring with the locally free lisse sheaf \(L_0^\vee/\pi_E^n\) does not change this range of degrees. Since the cohomology of \(S\) with coefficients in a torsion sheaf vanishes above \(\operatorname{cd}_\ell(S)\), the hypercohomology spectral sequence on \(S\) gives

\[
H^q R\Gamma\bigl(S,(L_0^\vee/\pi_E^n)\otimes u^*R\underline\Gamma_Z\,b_!(M_0/\pi_E^n)\bigr)=0
\quad(q>\operatorname{cd}_\ell(j)+\operatorname{cd}_\ell(S)+1),
\tag{5.2}
\]

for every \(n\), \(b\) and \(M_0\).

\emph{Integral targets.} For the lattice \(b_!M_0\) itself we write the relevant complex as a derived limit of its reductions modulo \(\pi_E^n\). Consider the complex \(L_0^\vee\otimes u^*R\underline\Gamma_Z\,b_!M_0\). By Lemma \hyperref[bp421]{5.4}, the functors that form it commute with derived inverse limits of sequences. They also commute with reduction modulo \(\pi_E^n\), which is the cone of multiplication by \(\pi_E^n\). Since \(b_!M_0\) is derived complete, so is this complex, and its reductions are the complexes in (5.2). In other words, the complex is the derived limit of its reductions, and we may compute its cohomology on \(S\) from theirs. Derived global sections preserve products and cones, and hence derived inverse limits. Using (3.1) with \(\mathcal O_E\) coefficients, that is, for \(R\!\operatorname{Hom}\) in \(D(X_{\mathrm{pro}\acute{e}t},\mathcal O_E)\), we get

\[
R\!\operatorname{Hom}_X(a_!L_0,b_!M_0)
=R\Gamma\bigl(S,L_0^\vee\otimes u^*R\underline\Gamma_Z\,b_!M_0\bigr)
=R\varprojlim_n R\Gamma\bigl(S,(L_0^\vee\otimes u^*R\underline\Gamma_Z\,b_!M_0)\otimes^{\mathbf L}\mathcal O_E/\pi_E^n\bigr).
\]

A derived inverse limit of a sequence raises an upper bound on cohomology by at most one. Indeed, it is the fibre of the difference of the identity and the transition maps on the product, and products of abelian groups are exact. This is the Milnor sequence [Stacks, \href{https://stacks.math.columbia.edu/tag/08TB}{08TB}]. Applied to the bound (5.2) for the reductions, it gives (5.1).
\end{proof}

% object: lemma-uniform-ext-bound (blueprint/chap-Construction-of-f/sec-Compactness-of-constructible-complexes/subsec-A-uniform-bound-for-Ext/lemma-uniform-ext-bound/_section.md)
\phantomsection\label{bp426}%
\begin{bpstatementit}{Lemma~5.6}{ (Uniform Ext bound)}
For every \(C\in\Cons(X,\Lambda)\) there is an integer \(d_C\geq0\) such that

\[
\Ext^q(C,N)=0
\qquad\text{for all }q>d_C\text{ and all }N\in\Cons(X,\Lambda).
\]
\end{bpstatementit}

% text: prose-05-09
\phantomsection\label{bp427}%
Suppose \(C=a_!L\), with \(a\) factored as \(S\hookrightarrow Z\hookrightarrow X\) and \(j\) the open complement of \(Z\), and suppose that \(L\) contains a lisse \(\mathcal O_E\)-lattice. By Lemma \hyperref[bp269]{3.24}, this holds after refining \(S\). In this case the proof gives \(d_C=\operatorname{cd}_\ell(j)+\operatorname{cd}_\ell(S)+2\). One of the two extra degrees comes from the cone defining \(R\underline\Gamma_Z\), and the other from a derived inverse limit over the reductions modulo \(\pi_E^n\).

% proof: proof/lemma-uniform-ext-bound
\begin{proof}[{Proof of Lemma 5.6}]\phantomsection\label{bp428}%
We reduce to extensions by zero of lisse sheaves on strata, prove the bound for \(\Lambda=E\) by choosing lattices, applying Lemma \hyperref[bp424]{5.5} and inverting \(\pi_E\), and reduce \(\Lambda=\overline{\mathbf Q}_\ell\) to finite extensions of \(E\).

\emph{Reduction to strata.} Bounds pass to extensions in either variable. Indeed, if \(0\to C'\to C\to C''\to0\) is exact, then in the long exact sequence \(\Ext^q(C'',N)\to\Ext^q(C,N)\to\Ext^q(C',N)\) the outer groups vanish above the larger of the two bounds for \(C'\) and \(C''\), and so the middle one does too. Hence a bound that holds for every subquotient of a finite filtration of the source holds for the source, if we take the maximum. The same sequence in the second variable shows that a bound that holds for every subquotient of a finite filtration of the target holds for the target, with the same value. In view of Lemma lemma-constructible-sheaves, we may therefore assume \(C=a_!L\), with \(a=i\circ u:S\hookrightarrow Z\hookrightarrow X\) as in Lemma \hyperref[bp257]{3.21} and \(L\) lisse. It suffices to find an integer \(d\), depending only on \(a\) and \(L\), with \(\Ext^q(a_!L,b_!M)=0\) for \(q>d\), for all locally closed \(b:T\hookrightarrow X\) and lisse \(M\) on \(T\). Let \(j:X\setminus Z\hookrightarrow X\) be the open complement of \(Z\). When \(L\) contains a lisse \(\mathcal O_E\)-lattice, we show that \(d=\operatorname{cd}_\ell(j)+\operatorname{cd}_\ell(S)+2\) works.

\emph{Lattices.} First take \(\Lambda=E\). Since torsion arguments need integral structures, we choose lattices. After refining the stratifications (Lemma \hyperref[bp269]{3.24}), \(L\) and \(M\) contain lisse \(\mathcal O_E\)-lattices \(L_0\) and \(M_0\). Refining is harmless, because it replaces \(a_!L\) and \(b_!M\) by finite extensions of sheaves of the same form, and the first step handles finite extensions. Moreover, the refinement of the source stratum depends only on \(L\), and that of the target only on \(M\). So the bound obtained for the refined source strata is still independent of \(N\).

The lattices are locally free of finite rank over \(\mathcal O_E\). Indeed, being lisse, they are locally of the form \(\mathcal O_{E,U}\otimes_{\mathcal O_E}P\) with \(P\) finitely generated (Definition \hyperref[bp201]{3.2}). As subsheaves of \(E\)-modules they are torsion free, so \(P\) is free. Lemma \hyperref[bp424]{5.5} now gives \(\Ext^q_{\mathcal O_E}(a_!L_0,b_!M_0)=0\) for \(q>\operatorname{cd}_\ell(j)+\operatorname{cd}_\ell(S)+2\).

\emph{Inverting \(\pi_E\).} Inverting \(\pi_E\) is a filtered colimit. By continuity of Ext out of \(a_!L_0\), the Ext groups over \(E\) are therefore colimits of groups that all vanish in the same range. In detail, we have \(a_!L=a_!L_0[1/\pi_E]\) and \(b_!M=(b_!M_0)[1/\pi_E]=\varinjlim(b_!M_0\xrightarrow{\pi_E}b_!M_0\xrightarrow{\pi_E}\cdots)\). The adjunction between extension and restriction of scalars gives

\[
R\!\operatorname{Hom}_E\bigl(a_!L,b_!M\bigr)
=R\!\operatorname{Hom}_{\mathcal O_E}\bigl(a_!L_0,(b_!M_0)[1/\pi_E]\bigr).
\]

By Lemma lemma-continuity-ext-out-constructible-complex the cohomology of the right side is the colimit of the groups \(\Ext^q_{\mathcal O_E}(a_!L_0,b_!M_0)\). Hence (5.1) implies the bound for \(E\).

\emph{The case \(\Lambda=\overline{\mathbf Q}_\ell\).} The bound for \(E\) depends only on the strata and on \(\operatorname{cd}_\ell\), not on \(E\), so it passes to the colimit over finite extensions. We first fix everything that depends on \(L\) alone, before the target is given. By Lemma \hyperref[bp245]{3.17}, applied to \(L\) alone, there are a finite extension \(E_0/\mathbf Q_\ell\) and a lisse \(E_0\)-sheaf \(L'\) with \(L=L'_{\overline{\mathbf Q}_\ell}\). After refining \(S\) (Lemma \hyperref[bp269]{3.24}), \(L'\) contains a lisse \(\mathcal O_{E_0}\)-lattice \(L'_0\). As in the case \(\Lambda=E\), the refinement replaces \(a_!L\) by a finite extension of sheaves of the same form, and it depends only on \(L\). For every finite extension \(E'\supset E_0\), the sheaf \(L'_0\otimes_{\mathcal O_{E_0}}\mathcal O_{E'}\) is a lisse \(\mathcal O_{E'}\)-lattice in \((L')_{E'}\) on the same strata. The case \(\Lambda=E'\) therefore gives \(\Ext^q_{E'}\bigl((a_!L')_{E'},N'\bigr)=0\) for \(q>\operatorname{cd}_\ell(j)+\operatorname{cd}_\ell(S)+2\), for each refined stratum \(S\) and every \(N'\in\Cons(X,E')\). This bound does not depend on \(E'\) or on \(N'\).

Now let \(b\) and \(M\) be given. By Lemma \hyperref[bp245]{3.17} again, applied to \(M\), there are a finite extension \(E\supset E_0\) (enlarge the field given by the lemma to contain \(E_0\)) and a lisse \(E\)-sheaf \(M'\) with \(M=M'_{\overline{\mathbf Q}_\ell}\). The field \(E\) depends on \(M\), but neither the strata of the source nor the lattice \(L'_0\) do. Then \(a_!L=(a_!L')_{\overline{\mathbf Q}_\ell}\) and \(b_!M=(b_!M')_{\overline{\mathbf Q}_\ell}\), and by Lemma lemma-descent-finite-coefficient-field,

\[
\Ext^q_{\overline{\mathbf Q}_\ell}(a_!L,b_!M)
=\varinjlim_{E'\supset E}\Ext^q_{E'}\bigl((a_!L')_{E'},(b_!M')_{E'}\bigr).
\]

Every group on the right vanishes for \(q>\operatorname{cd}_\ell(j)+\operatorname{cd}_\ell(S)+2\) by the previous paragraph, and hence so does the colimit.
\end{proof}

% object: example-ext-bound-finite-field (blueprint/chap-Construction-of-f/sec-Compactness-of-constructible-complexes/subsec-A-uniform-bound-for-Ext/example-ext-bound-finite-field/_section.md)
\phantomsection\label{bp429}%
\begin{bpstatement}{Example~5.7}{ (The bound over a finite field)}
The proof gives an explicit bound, which is far from sharp. For instance, for \(X=\operatorname{Spec}\mathbf F_q\) and \(C=N=E\) we have \(Z=S=X\). The complement of \(Z\) is then empty, so \(R\underline\Gamma_Z\) is the identity. Moreover \(\operatorname{cd}_\ell(\mathbf F_q)=1\), and so the formula, read with \(\operatorname{cd}_\ell(j)=0\) for the empty open, gives \(d_C=0+1+2=3\), while \(\Ext^q(E,E)\) vanishes already for \(q>1\) (Example \hyperref[bp417]{5.3}). The two wasted degrees do no harm, because the arguments below only need some bound independent of \(N\).
\end{bpstatement}


% object: subsec-Ind-constructible-sheaves (blueprint/chap-Construction-of-f/sec-Compactness-of-constructible-complexes/subsec-Ind-constructible-sheaves/_section.md), content hidden
\subsubsection*{\texorpdfstring{5.2.2\quad Ind-constructible sheaves}{5.2.2 Ind-constructible sheaves}}
\phantomsection\label{bp430}%
\addcontentsline{toc}{subsubsection}{\texorpdfstring{5.2.2\quad Ind-constructible sheaves}{5.2.2 Ind-constructible sheaves}}
% text: prose-05-10
\phantomsection\label{bp431}%
The targets in the compactness statement are direct sums of objects of \(\Loc(\Dbc(X,\Lambda))\), and their cohomology sheaves need not be constructible. Already the sheaf \(\bigoplus_{n\in\mathbf N}\Lambda_X\) is not constructible, since its stalks are not finitely generated. Such sheaves are, however, filtered colimits of constructible sheaves, and both continuity and the uniform bound for Ext groups extend to these colimits.

Recall that the ind-category \(\operatorname{Ind}(\mathcal C)\) of a small category \(\mathcal C\) has as objects the filtered diagrams in \(\mathcal C\) (“formal filtered colimits”), with \(\Hom(\varinjlim_\alpha C_\alpha,\varinjlim_\beta D_\beta)
=\varprojlim_\alpha\varinjlim_\beta\Hom(C_\alpha,D_\beta)\). Proposition \hyperref[bp433]{5.9} below says that for constructible sheaves the formal and the actual colimits agree.

The dévissage arguments need a class of sheaves that is abelian and closed under extensions and direct sums, and on which \(\Ext^q(C,-)\) behaves as on constructible sheaves. By Proposition \hyperref[bp433]{5.9}, the category \(\IndCons(X,\Lambda)\) has the first three properties, and Corollary \hyperref[bp436]{5.10} gives the last.

% object: definition-ind-constructible-sheaves (blueprint/chap-Construction-of-f/sec-Compactness-of-constructible-complexes/subsec-Ind-constructible-sheaves/definition-ind-constructible-sheaves/_section.md)
\phantomsection\label{bp432}%
\begin{bpstatement}{Definition~5.8}{ (Ind-constructible sheaves)}
An \emph{ind-constructible sheaf} on \(X\) is a sheaf of \(\Lambda_X\)-modules on \(X_{\mathrm{pro}\acute{e}t}\) that is a filtered colimit, computed in \(\operatorname{Mod}(X_{\mathrm{pro}\acute{e}t},\Lambda)\), of constructible sheaves. We write \(\IndCons(X,\Lambda)\subset\operatorname{Mod}(X_{\mathrm{pro}\acute{e}t},\Lambda)\) for the full subcategory of ind-constructible sheaves.
\end{bpstatement}

% object: proposition-ind-constructible-sheaves (blueprint/chap-Construction-of-f/sec-Compactness-of-constructible-complexes/subsec-Ind-constructible-sheaves/proposition-ind-constructible-sheaves/_section.md)
\phantomsection\label{bp433}%
\begin{bpstatementit}{Proposition~5.9}{ (Ind-constructible sheaves)}
\begin{enumerate}
\item The functor \(\operatorname{Ind}(\Cons(X,\Lambda))\to\operatorname{Mod}(X_{\mathrm{pro}\acute{e}t},\Lambda)\) sending a formal filtered colimit to the actual colimit is fully faithful, with essential image \(\IndCons(X,\Lambda)\).
\item \(\IndCons(X,\Lambda)\) is closed in \(\operatorname{Mod}(X_{\mathrm{pro}\acute{e}t},\Lambda)\) under filtered colimits, direct sums, kernels, cokernels and extensions. In particular it is a Grothendieck abelian category and its inclusion is exact.
\end{enumerate}
\end{bpstatementit}

% proof: proof/proposition-ind-constructible-sheaves
\begin{proof}[{Proof}]\phantomsection\label{bp434}%
The proof uses one property of constructible sheaves: maps out of them commute with filtered colimits. Because of this, the formal colimits of the ind-category behave like actual colimits. Each closure property in (2) then reduces to the same property of \(\Cons(X,\Lambda)\), once the objects and maps involved are written as filtered colimits of constructible sheaves and maps between them.

\emph{Finite presentation.} Constructible sheaves are finitely presented in \(\operatorname{Mod}(X_{\mathrm{pro}\acute{e}t},\Lambda)\). Indeed, Lemma \hyperref[bp260]{3.22} in degree zero shows that \(\Hom(C,-)\) commutes with filtered colimits for \(C\) constructible. By Lemma lemma-constructible-sheaves,(2), the category \(\Cons(X,\Lambda)\) is essentially small, abelian, and closed under extensions. We fix a small skeleton, so that the ind-category in (1) is defined.

\emph{Full faithfulness.} For filtered diagrams \((C_\alpha)\) and \((D_\beta)\) of constructible sheaves, finite presentation implies

\[
\Hom(\varinjlim_\alpha C_\alpha,\varinjlim_\beta D_\beta)
=\varprojlim_\alpha\varinjlim_\beta\Hom_{\Cons(X,\Lambda)}(C_\alpha,D_\beta),
\]

which is the Hom set in the ind-category. This proves (1), since the essential image is \(\IndCons(X,\Lambda)\) by definition.

\emph{Canonical presentations.} The closure properties below involve several ind-constructible sheaves at once, so to compare their presentations we use a presentation that involves no choices. Namely, for \(B\in\IndCons(X,\Lambda)\), let \(J(B)\) be the category of constructible sheaves \(C\) with a map \(C\to B\); we show that \(B=\varinjlim_{J(B)}C\).

The category \(J(B)\) is filtered. Indeed, two objects \(C_1\to B\) and \(C_2\to B\) both map to \(C_1\oplus C_2\to B\), and two parallel arrows \(C\rightrightarrows C'\) over \(B\) are equalized by the cokernel of their difference, which is constructible and still maps to \(B\). Now write \(B=\varinjlim_\alpha C_\alpha\). By finite presentation, every \(C\to B\) factors through some \(C_\alpha\), and two factorizations agree at a later stage. This means that the given diagram is cofinal in \(J(B)\), and hence \(B=\varinjlim_{J(B)}C\).

\emph{Filtered colimits and direct sums.} Let \(B=\varinjlim_\lambda B_\lambda\) be a filtered colimit with \(B_\lambda\in\IndCons(X,\Lambda)\). To present \(B\), consider the category of pairs \((\lambda,C\to B_\lambda)\) with \(C\) constructible. It is filtered: we move finitely many indices to a common upper bound, and then use direct sums and cokernels in \(J(B_\lambda)\). Moreover, the colimit of the \(C\) over this category is \(\varinjlim_\lambda\varinjlim_{J(B_\lambda)}C=B\). A direct sum is the filtered colimit of its finite subsums, so it lies in \(\IndCons(X,\Lambda)\) as well.

\emph{Kernels and cokernels.} Let \(f:B\to B'\) be a map in \(\IndCons(X,\Lambda)\). It suffices to write \(f\) as a filtered colimit of maps \(v\) of constructible sheaves. Since filtered colimits are exact, \(\ker f\) and \(\coker f\) are then the filtered colimits of the constructible sheaves \(\ker v\) and \(\coker v\). Let \(P_f\) be the category of commutative squares

\[
\begin{matrix} C&\xrightarrow{v}&D\\
\downarrow&&\downarrow\\ B&\xrightarrow{f}&B',\end{matrix}
\qquad C,D\in\Cons(X,\Lambda).
\]

This category is filtered, because finite colimits of such squares are formed by finite colimits of the arrows \(v\) in the abelian category \(\Cons(X,\Lambda)\). We claim that the functor \(P_f\to J(B)\), \((v)\mapsto(C\to B)\), is cofinal. Indeed, every \(C\to B\) is the source of a square, namely \(v=\mathrm{id}_C\) with \(C\to B'\) the composite. If \(C\to C_1\xrightarrow{v_1}D_1\) and \(C\to C_2\xrightarrow{v_2}D_2\) are two squares under the same \(C\to B\), then both map to the square with source \(C_1\amalg_CC_2\) and target \(D_1\amalg_CD_2\). Here the pushouts are taken in \(\Cons(X,\Lambda)\), and they map to \(B\) and \(B'\) because the two maps from \(C\) to \(B\), and the two maps from \(C\) to \(B'\), agree.

Similarly, the functor \(P_f\to J(B')\), \((v)\mapsto(D\to B')\), is cofinal. Every \(D\to B'\) is the target of a square, namely \(v:0\to D\). If \(C_1\xrightarrow{v_1}D_1\) and \(C_2\xrightarrow{v_2}D_2\) are two squares under the same \(D\to B'\), then both map to the square with source \(C_1\oplus C_2\) and target \(D_1\amalg_DD_2\). The two cofinality statements give \(f=\varinjlim_{P_f}v\), and exactness of filtered colimits gives

\[
\ker f=\varinjlim_{P_f}\ker v,\qquad
\coker f=\varinjlim_{P_f}\coker v,
\]

both of which lie in \(\IndCons(X,\Lambda)\).

\emph{Extensions.} Let \(0\to N\to M\to B\to0\) be exact in \(\operatorname{Mod}(X_{\mathrm{pro}\acute{e}t},\Lambda)\), with \(N=\varinjlim_\beta D_\beta\) and \(B=\varinjlim_\alpha C_\alpha\) ind-constructible. We show that \(M\) is a filtered colimit of extensions of constructible sheaves. First, the pullback of the extension along \(C_\alpha\to B\) is ind-constructible. Its class lies in \(\Ext^1(C_\alpha,N)\), and Lemma \hyperref[bp260]{3.22} gives

\[
\Ext^1(C_\alpha,N)
=\varinjlim_\beta\Ext^1(C_\alpha,D_\beta).
\]

So the pulled-back extension is pushed out from an extension \(M_{\alpha\beta}\) of \(C_\alpha\) by some \(D_\beta\). Since \(\Cons(X,\Lambda)\) is closed under extensions, \(M_{\alpha\beta}\) and each of its pushouts along \(D_\beta\to D_{\beta'}\) are constructible. Their colimit over \(\beta'\) is the pullback of \(M\) to \(C_\alpha\). Second, since filtered colimits are exact, \(M\) is the colimit over \(\alpha\) of these pullbacks. By closure under filtered colimits, \(M\in\IndCons(X,\Lambda)\).

\emph{Grothendieck.} The category \(\IndCons(X,\Lambda)\) is Grothendieck abelian in the sense of [Stacks, \href{https://stacks.math.columbia.edu/tag/079A}{079A}]. Indeed, by the above its colimits exist and are computed in \(\operatorname{Mod}(X_{\mathrm{pro}\acute{e}t},\Lambda)\), so filtered colimits are exact. The direct sum of the objects of the skeleton of \(\Cons(X,\Lambda)\) is a generator.
\end{proof}

% text: prose-05-11
\phantomsection\label{bp435}%
Since every ind-constructible sheaf is a filtered colimit of constructible ones, the continuity of \(\Ext^q(C,-)\) lets us extend the uniform bound of Lemma \hyperref[bp426]{5.6} to the whole category \(\IndCons(X,\Lambda)\). Writing a direct sum as the filtered colimit of its finite subsums, we also see that \(\Ext^q(C,-)\) commutes with direct sums of ind-constructible sheaves.

% object: corollary-ind-heart-b (blueprint/chap-Construction-of-f/sec-Compactness-of-constructible-complexes/subsec-Ind-constructible-sheaves/corollary-ind-heart-b/_section.md)
\phantomsection\label{bp436}%
\begin{bpstatementit}{Corollary~5.10}{ (Ext from a constructible sheaf into ind-constructible sheaves)}
Let \(C\in\Cons(X,\Lambda)\), with \(d_C\) as in Lemma \hyperref[bp426]{5.6}. Then \(\Ext^q(C,B)=0\) for all \(q>d_C\) and all \(B\in\IndCons(X,\Lambda)\), and each functor \(\Ext^q(C,-)\) commutes with direct sums of objects of \(\IndCons(X,\Lambda)\).
\end{bpstatementit}

% proof: proof/corollary-ind-heart-b
\begin{proof}[{Proof}]\phantomsection\label{bp437}%
Write \(B=\varinjlim_\alpha B_\alpha\) as a filtered colimit of constructible sheaves, which is possible by Proposition \hyperref[bp433]{5.9}. Sheaves are complexes concentrated in degree \(0\), so they vanish below a common degree, and Lemma \hyperref[bp260]{3.22} applies. It shows that \(\Ext^q(C,B)=\varinjlim_\alpha\Ext^q(C,B_\alpha)\). Each group on the right vanishes for \(q>d_C\) by Lemma \hyperref[bp426]{5.6}, and so does the colimit.

For the second claim, let \((B_\lambda)\) be a family in \(\IndCons(X,\Lambda)\). Its direct sum is the filtered colimit of the finite subsums, and \(\Ext^q(C,-)\) commutes with finite direct sums because it is additive. Applying Lemma \hyperref[bp260]{3.22} once more, to the diagram of finite subsums, we obtain \(\Ext^q(C,\bigoplus_\lambda B_\lambda)=\bigoplus_\lambda\Ext^q(C,B_\lambda)\).
\end{proof}

% text: prose-05-12
\phantomsection\label{bp438}%
Objects of \(\Loc(\Dbc(X,\Lambda))\) can be unbounded, and we do not know a priori that their truncations stay in \(\Loc(\Dbc(X,\Lambda))\). For this reason, here and in Section \hyperref[bp459]{5.4}, we carry out truncation arguments among complexes with ind-constructible cohomology sheaves.

% object: corollary-ambient-ind-derived (blueprint/chap-Construction-of-f/sec-Compactness-of-constructible-complexes/subsec-Ind-constructible-sheaves/corollary-ambient-ind-derived/_section.md)
\phantomsection\label{bp439}%
\begin{bpstatementit}{Corollary~5.11}{ (The derived category of ind-constructible sheaves)}
The full subcategory

\[
\Dic(X,\Lambda)
=\{K\in D(X_{\mathrm{pro}\acute{e}t},\Lambda):H^n(K)\in\IndCons(X,\Lambda)\text{ for every }n\}
\]

of \(D(X_{\mathrm{pro}\acute{e}t},\Lambda)\) is triangulated, closed under direct sums, and stable under the truncations \(\tau_{\leq n}\), \(\tau_{\geq n}\). In particular \(\Loc(\Dbc(X,\Lambda))\subset\Dic(X,\Lambda)\).
\end{bpstatementit}

% proof: proof/corollary-ambient-ind-derived
\begin{proof}[{Proof}]\phantomsection\label{bp440}%
In the long exact cohomology sequence of a triangle, each cohomology sheaf of the third term is an extension of a kernel by a cokernel of maps between cohomology sheaves of the other two. Since \(\IndCons(X,\Lambda)\) is closed under kernels, cokernels and extensions (Proposition \hyperref[bp433]{5.9}), the third term lies in \(\Dic(X,\Lambda)\) whenever the other two do. Closure under shifts is clear. As direct sums of sheaves are exact, \(H^n\) commutes with direct sums. Since \(\IndCons(X,\Lambda)\) is closed under direct sums, it follows that \(\Dic(X,\Lambda)\) is closed under direct sums. Truncation only removes cohomology sheaves and keeps the others. Hence \(\Dic(X,\Lambda)\) is a localizing subcategory of \(D(X_{\mathrm{pro}\acute{e}t},\Lambda)\) containing \(\Dbc(X,\Lambda)\). It therefore contains the smallest one, \(\Loc(\Dbc(X,\Lambda))\).
\end{proof}


% object: subsec-Rational-coefficients (blueprint/chap-Construction-of-f/sec-Compactness-of-constructible-complexes/subsec-Rational-coefficients/_section.md), content hidden
\subsubsection*{\texorpdfstring{5.2.3\quad Rational coefficients}{5.2.3 Rational coefficients}}
\phantomsection\label{bp441}%
\addcontentsline{toc}{subsubsection}{\texorpdfstring{5.2.3\quad Rational coefficients}{5.2.3 Rational coefficients}}
% text: prose-05-13
\phantomsection\label{bp442}%
On ind-constructible sheaves, \(\Ext^q(C,-)\) vanishes above \(d_C\) and commutes with direct sums (Corollary \hyperref[bp436]{5.10}). Because the bound \(d_C\) is uniform, maps from a constructible sheaf \(C\) into an unbounded complex depend only on its cohomology sheaves in degrees \([-d_C,0]\). This is the step that fails over \(\mathbf R\) in the example at the beginning of this chapter, where \(\Ext^n(\mathbf F_2,\mathbf F_2)\neq0\) for all \(n\). Since the window has the same length for a whole family of targets, compactness follows from the case of sheaves by induction on the length of the window.

% object: lemma-reduction-finitely-many-degrees (blueprint/chap-Construction-of-f/sec-Compactness-of-constructible-complexes/subsec-Rational-coefficients/lemma-reduction-finitely-many-degrees/_section.md)
\phantomsection\label{bp443}%
\begin{bpstatementit}{Lemma~5.12}{ (Reduction to finitely many degrees)}
Let \(C\in\Cons(X,\Lambda)\), with \(d_C\) as in Lemma \hyperref[bp426]{5.6}, and let \(K\in\Dic(X,\Lambda)\) and \(m\in\mathbf Z\). The truncation maps induce a natural isomorphism

\[
\Hom(C,K[m])
\simeq\Hom\bigl(C,(\tau_{\geq m-d_C}\tau_{\leq m}K)[m]\bigr).
\tag{5.3}
\]
\end{bpstatementit}

% proof: proof/lemma-reduction-finitely-many-degrees
\begin{proof}[{Proof}]\phantomsection\label{bp444}%
Cutting off the cohomology above degree \(m\) is harmless for degree reasons alone. To remove the cohomology below degree \(m-d_C\) we use the uniform bound: for bounded \(K\) it kills the tail layer by layer, and for \(K\) unbounded below we write \(K\) as the derived limit of its truncations and control the resulting \(\varprojlim^1\) term.

Replacing \(K\) by \(K[m]\), we may assume \(m=0\). The truncation \(\tau_{>0}K\) lies in \(D^{\geq1}\), and \(C\) is a sheaf in degree \(0\), so there are no nonzero maps from \(C\) to an object of \(D^{\geq1}\). Hence \(\Hom(C,\tau_{>0}K[s])=0\) for \(s=0,-1\). In the long exact sequence of \(\Hom(C,-)\) for the truncation triangle \(\tau_{\leq0}K\to K\to\tau_{>0}K\to\), these are the two groups adjacent to \(\Hom(C,K)\), and so \(\Hom(C,K)=\Hom(C,\tau_{\leq0}K)\). To remove the cohomology in degrees \(\leq-d_C-1\) from \(K'=\tau_{\leq0}K\), we use the triangle \(\tau_{\leq-d_C-1}K'\to K'\to\tau_{\geq-d_C}K'\to\). It suffices to show that

\[
\Hom(C,L)=0\qquad
\text{for every }L\in\Dic(X,\Lambda)\text{ with }H^q(L)=0\text{ for }q>-d_C-1,
\tag{5.4}
\]

applied to \(L=\tau_{\leq-d_C-1}K'\) and to \(L=\tau_{\leq-d_C-1}K'[1]\).

We first prove (5.4) for bounded \(L\). Such an \(L\) is a finite extension of its shifted cohomology sheaves \(H^q(L)[-q]\), and by the long exact Hom sequences it suffices to show that each of them receives no nonzero map from \(C\). A layer \(H^q(L)[-q]\) with \(q\leq-d_C-1\) contributes \(\Hom(C,H^q(L)[-q])=\Ext^{-q}(C,H^qL)\). Since \(-q>d_C\) and \(H^q(L)\) is ind-constructible, this group vanishes by Corollary \hyperref[bp436]{5.10}. For \(L\) unbounded below we cannot argue in finitely many steps, and we use instead that the derived category is left complete.

The pro-étale site has a basis of weakly contractible objects (Theorem \hyperref[bp51]{2.17}). For this reason Proposition \hyperref[bp68]{2.23} applies: products of \(\Lambda_X\)-modules are exact, products in \(D(X_{\mathrm{pro}\acute{e}t},\Lambda)\) are computed termwise, and

\[
L=R\varprojlim_n\tau_{\geq-n}L
\]

for every \(L\). Concretely, if \(L^\bullet\) is a complex representing \(L\), the limit is the derived inverse limit of the tower of good truncations \(\tau_{\geq-n}L^\bullet\) with their canonical transition maps. Since \(\Hom(C,-)\) takes products to products, the triangle defining \(R\varprojlim\) yields the Milnor exact sequence

\[
0\longrightarrow\varprojlim_n{}^1
\Hom(C,\tau_{\geq-n}L[-1])
\longrightarrow\Hom(C,L)
\longrightarrow\varprojlim_n\Hom(C,\tau_{\geq-n}L)
\longrightarrow0.
\]

Each \(\tau_{\geq-n}L\) is bounded with cohomology in degrees \(\leq-d_C-1\), and hence the right-hand term vanishes by the bounded case. For the left-hand term, the triangle

\[
H^{-n-1}(L)[n+1]\longrightarrow \tau_{\geq-n-1}L
\longrightarrow \tau_{\geq-n}L\longrightarrow H^{-n-1}(L)[n+2]
\]

gives an exact sequence

\[
\Ext^{n}(C,H^{-n-1}L)\to\Hom(C,\tau_{\geq-n-1}L[-1])\to\Hom(C,\tau_{\geq-n}L[-1])
\to\Ext^{n+1}(C,H^{-n-1}L).
\]

The outer terms vanish for \(n>d_C\) by Corollary \hyperref[bp436]{5.10}. The transition maps are therefore isomorphisms for \(n>d_C\), so the tower is eventually constant and its \(\varprojlim^1\) vanishes. This proves (5.4).
\end{proof}

% text: prose-05-14
\phantomsection\label{bp445}%
The window \([m-d_C,m]\) in (5.3) depends on \(C\) and \(m\) only, not on \(K\), and so we can treat a whole family \((K_\lambda)\) at once in the compactness theorem.

% object: theorem-compactness-rational-coefficients (blueprint/chap-Construction-of-f/sec-Compactness-of-constructible-complexes/subsec-Rational-coefficients/theorem-compactness-rational-coefficients/_section.md)
\phantomsection\label{bp446}%
\begin{bpstatementit}{Theorem~5.13}{ (Compactness with rational coefficients)}
For \(\Lambda=E\) or \(\overline{\mathbf Q}_\ell\), every object of \(\Dbc(X,\Lambda)\) is compact in \(\Loc(\Dbc(X,\Lambda))\). More precisely, for \(C\in\Dbc(X,\Lambda)\) and every family \((K_\lambda)\) of objects of \(\Dic(X,\Lambda)\), the natural map

\[
\bigoplus_\lambda\Hom(C,K_\lambda)
\longrightarrow\Hom\Bigl(C,\bigoplus_\lambda K_\lambda\Bigr)
\]

is an isomorphism.
\end{bpstatementit}

% proof: proof/theorem-compactness-rational-coefficients
\begin{proof}[{Proof}]\phantomsection\label{bp447}%
The idea is that, by Lemma \hyperref[bp443]{5.12}, maps out of a constructible sheaf depend only on the cohomology of the target in \(d_C+1\) consecutive degrees, uniformly in the family. For targets with cohomology in a fixed bounded interval the claim follows from the case of ind-constructible sheaves by induction on the length.

The first assertion follows from the second, since \(\Loc(\Dbc(X,\Lambda))\subset\Dic(X,\Lambda)\) (Corollary \hyperref[bp439]{5.11}) and direct sums in \(\Loc(\Dbc(X,\Lambda))\) are those of \(D(X_{\mathrm{pro}\acute{e}t},\Lambda)\).

First let \(C\in\Cons(X,\Lambda)\), and suppose that all \(K_\lambda\) have cohomology in a fixed interval \([a,b]\). We prove the claim, for all shifts of the \(K_\lambda\) at once, by induction on \(b-a\). If \(a=b\), then \(K_\lambda=B_\lambda[-a]\) with \(B_\lambda\in\IndCons(X,\Lambda)\), and the claim is that \(\Ext^q(C,-)\) commutes with direct sums of objects of \(\IndCons(X,\Lambda)\), which is Corollary \hyperref[bp436]{5.10}. In general, the triangles \(\tau_{\leq b-1}K_\lambda\to K_\lambda\to H^b(K_\lambda)[-b]\to\) are compatible with direct sums, because \(H^n\) commutes with direct sums. We apply \(\Hom(C,-)\) to these triangles and take the direct sum over \(\lambda\); since direct sums of abelian groups are exact, the result is again a long exact sequence. It maps to the long exact sequence of the summed triangle. The outer terms are of the form covered by the induction hypothesis (cohomology in \([a,b-1]\) and in \([b,b]\)), and so the five lemma gives the claim for cohomology in \([a,b]\).

For arbitrary \(K_\lambda\), we apply Lemma \hyperref[bp443]{5.12} with \(m=0\), which allows us to replace each \(K_\lambda\), and \(\bigoplus_\lambda K_\lambda\), by its truncation to degrees \([-d_C,0]\). Since the interval is the same for the whole family and truncation commutes with direct sums, the bounded case applies.

From constructible sheaves \(C\) we pass to complexes by dévissage. For fixed \((K_\lambda)\), the objects \(C\) for which the map is an isomorphism are closed under shifts. They are closed under cones by the five lemma, applied to the long exact sequences of \(\Hom(-,K_\lambda)\) summed over \(\lambda\) and of \(\Hom(-,\bigoplus_\lambda K_\lambda)\). A bounded constructible complex is a finite extension of shifts of its constructible cohomology sheaves, and hence it belongs to this class.
\end{proof}

% text: prose-05-15
\phantomsection\label{bp448}%
Lemma \hyperref[bp503]{5.30} combines this theorem with the corresponding statements for the other coefficient rings and shows that a small skeleton of \(D^b_{\mathrm{ctf}}(X,\Lambda)\) is a set of compact generators of \(\Loc(D^b_{\mathrm{ctf}}(X,\Lambda))\).



% object: sec-Compactness-for-finite-coefficients (blueprint/chap-Construction-of-f/sec-Compactness-for-finite-coefficients/_section.md), content hidden
\subsection*{\texorpdfstring{5.3\quad Compactness for finite coefficients}{5.3 Compactness for finite coefficients}}
\phantomsection\label{bp449}%
\addcontentsline{toc}{subsection}{\texorpdfstring{5.3\quad Compactness for finite coefficients}{5.3 Compactness for finite coefficients}}
% text: prose-05-16
\phantomsection\label{bp450}%
In this section the coefficient ring is a finite admissible ring \(A\). Proper direct image preserves \(D^b_{\mathrm{ctf}}\) by Proposition \hyperref[bp390]{4.35}, and we now show that the constructible generators are compact and that proper direct image preserves direct sums on the localizing categories. In both cases we reduce to the classical étale theory, using bounded flat constructible models and the fact that étale cohomology commutes with direct sums of unbounded complexes.

Finite Tor amplitude is a property of the complex, not of its cohomology sheaves: the cohomology sheaves of an object of \(D^b_{\mathrm{ctf}}\) can have infinite Tor amplitude (Section \hyperref[bp230]{3.4}). The property can, however, be read off from a flat model. This is why we work with the bounded flat models of Lemma \hyperref[bp235]{3.13}.

% object: proposition-compactness-finite-coefficients (blueprint/chap-Construction-of-f/sec-Compactness-for-finite-coefficients/proposition-compactness-finite-coefficients/_section.md)
\phantomsection\label{bp451}%
\begin{bpstatementit}{Proposition~5.14}{ (Compactness for finite coefficients)}
Every object of \(D^b_{\mathrm{ctf}}(X,A)\) is compact in \(\Loc(D^b_{\mathrm{ctf}}(X,A))\), and a small skeleton of \(D^b_{\mathrm{ctf}}(X,A)\) is a set of compact generators. Moreover, \(\Loc(D^b_{\mathrm{ctf}}(X,A))^c=D^b_{\mathrm{ctf}}(X,A)\).
\end{bpstatementit}

% proof: proof/proposition-compactness-finite-coefficients
\begin{proof}[{Proof}]\phantomsection\label{bp452}%
The plan is to prove compactness on the étale site against arbitrary objects of \(D(X_{\acute{e}t},A)\), where the commutation of cohomology with direct sums that we need is classical, and then transfer the result to \(X_{\mathrm{pro}\acute{e}t}\) by \(\nu_X^*\).

Recall that compact objects are closed under shifts, finite extensions and direct summands. Hence, by Lemma \hyperref[bp235]{3.13}, it suffices to treat \(a_!L\). To do so, factor \(a=i\circ u\) with \(i:Z\hookrightarrow X\) closed and \(u:S\hookrightarrow Z\) open, and let \(j:X\setminus Z\hookrightarrow X\) be the complementary open immersion. The idea is then to write \(R\!\operatorname{Hom}(a_!L,-)\) as a composite of functors each of which commutes with direct sums. Since \(L\) has finite projective stalks, it is dualizable with dual \(L^\vee=\underline{\operatorname{Hom}}_A(L,A_S)\), and the étale form of Lemma lemma-sections-support-closed-subscheme shows that for \(K\in D(X_{\acute{e}t},A)\) we have

\[
R\!\operatorname{Hom}_X(a_!L,K)
=R\Gamma\bigl(S,L^\vee\otimes_A u^*R\underline\Gamma_ZK\bigr),
\qquad R\underline\Gamma_ZK=i^*\fib\bigl(K\to Rj_*j^*K\bigr),
\]

where \(R\underline\Gamma_Z\) denotes sections with support in \(Z\). We claim that every functor on the right, viewed as a functor between étale derived categories, preserves arbitrary direct sums. For \(Rj_*\) and \(R\Gamma(S,-)\), this is [Stacks, \href{https://stacks.math.columbia.edu/tag/0F11}{0F11}], which says that a derived direct image of finite cohomological dimension commutes with direct sums of unbounded complexes. It applies to \(A\)-modules because derived direct image and global sections commute with forgetting the coefficient ring [Stacks, \href{https://stacks.math.columbia.edu/tag/0D6J}{0D6J}]. The remaining functors are inverse images, tensor products with finite projective sheaves, and a fibre, and these also preserve direct sums: inverse images are left adjoints, tensor products commute with direct sums, and the fibre of a map between functors that preserve direct sums again preserves them. Hence \(R\!\operatorname{Hom}_X(a_!L,-)\) commutes with direct sums. Taking \(H^0\), we see that \(\Hom(a_!L,-)\) commutes with direct sums, that is, \(a_!L\) is compact. By Lemma \hyperref[bp235]{3.13} it follows that every object of \(D^b_{\mathrm{ctf}}(X_{\acute{e}t},A)\) is compact.

On the pro-étale site, Corollary \hyperref[bp132]{2.44} and Proposition \hyperref[bp146]{2.49} show that the inverse image \(\nu_X^*:D(X_{\acute{e}t},A)\to D(X_{\mathrm{pro}\acute{e}t},A)\) is exact, fully faithful and preserves direct sums. Such a functor identifies the localizing subcategories generated by \(D^b_{\mathrm{ctf}}\) on the two sites. Compactness against all étale complexes therefore implies compactness in \(\Loc(D^b_{\mathrm{ctf}}(X,A))\).

For the last assertion we have to show that there are no other compact objects. Since constructibility and finite Tor amplitude are preserved by direct summands, the category \(D^b_{\mathrm{ctf}}(X,A)\) is thick. It is essentially small, because it is a full subcategory of the essentially small bounded constructible category, and by definition a skeleton generates \(\Loc(D^b_{\mathrm{ctf}}(X,A))\). In a triangulated category with direct sums generated by a set of compact objects, the compact objects form the thick closure of that set [Stacks, \href{https://stacks.math.columbia.edu/tag/09SM}{09SM}, \href{https://stacks.math.columbia.edu/tag/09SN}{09SN}]. In our case this closure is \(D^b_{\mathrm{ctf}}(X,A)\).
\end{proof}

% text: prose-05-17
\phantomsection\label{bp453}%
For finite coefficients the localizing category is a familiar one: the next proposition shows that it is the classical unbounded derived category of the étale site. In particular, \(f^!\) with finite coefficients is defined on all of \(D(X_{\acute{e}t},A)\). The proposition also shows that \(Rf_!\), once it is defined as \(Rp_*j_!\) by a compactification in Section \hyperref[bp459]{5.4}, is computed by the classical functor \(Rp_{\acute{e}t,*}j_{\acute{e}t,!}\) on the étale site.

% object: proposition-comparison-etale-derived-category (blueprint/chap-Construction-of-f/sec-Compactness-for-finite-coefficients/proposition-comparison-etale-derived-category/_section.md)
\phantomsection\label{bp454}%
\begin{bpstatementit}{Proposition~5.15}{ (Comparison with the étale derived category)}
For a finite admissible ring \(A\), the comparison functor \(\nu_X^*:D(X_{\acute{e}t},A)\to D(X_{\mathrm{pro}\acute{e}t},A)\) identifies

\[
D(X_{\acute{e}t},A)=\Loc(D^b_{\mathrm{ctf}}(X_{\acute{e}t},A))
\simeq\Loc(D^b_{\mathrm{ctf}}(X,A)).
\]

It preserves arbitrary direct sums. For every morphism \(g:X\to Y\) of schemes of finite type over \(k\) and every \(K\in D(X_{\acute{e}t},A)\), the comparison map (2.9) \(\nu_Y^*Rg_{\acute{e}t,*}K\to Rg_{\mathrm{pro}\acute{e}t,*}\nu_X^*K\) is an isomorphism, and \(\nu^*\) commutes with inverse images and, by (2.10), with extension by zero along open immersions.
\end{bpstatementit}

% proof: proof/proposition-comparison-etale-derived-category
\begin{proof}[{Proof}]\phantomsection\label{bp455}%
There are two statements: \(\nu_X^*\) identifies the two localizing categories, and on the étale side the localizing category is everything. The first follows formally from the comparison results for the pro-étale site. For the second we exhibit compact generators of \(D(X_{\acute{e}t},A)\) that lie in \(D^b_{\mathrm{ctf}}\).

An exact fully faithful functor that preserves direct sums identifies the localizing subcategories generated by corresponding generators. In our case, Proposition \hyperref[bp146]{2.49} shows that \(\nu_X^*\) is such a functor on unbounded complexes and that it is compatible with direct images, and on \(D^b_{\mathrm{ctf}}\) it is the equivalence of Corollary \hyperref[bp132]{2.44}. Finally, inverse images commute with \(\nu^*\) because both composites are inverse images along the same morphism of topoi.

For the equality with the whole étale derived category, let \(u:U\to X\) range over a basis of quasi-compact separated étale objects. The free representable module sheaves \(u_!A_U\) are flat constructible, hence belong to \(D^b_{\mathrm{ctf}}(X_{\acute{e}t},A)\). We have \(R\!\operatorname{Hom}_X(u_!A_U,K)=R\Gamma(U,K)\). It follows that their shifts detect zero objects of \(D(X_{\acute{e}t},A)\). Indeed, if \(H^n(U,K)=0\) for all such \(U\) and all \(n\), then every cohomology sheaf of \(K\) vanishes, being the sheafification of \(U\mapsto H^n(U,K)\). These objects are also compact in all of \(D(X_{\acute{e}t},A)\). The reason is that \(R\Gamma(U,-)\) commutes with direct sums, since \(\operatorname{cd}_\ell(U)<\infty\), as in the étale part of the proof of Proposition \hyperref[bp451]{5.14}. A set of compact objects whose shifts detect zero objects generates a triangulated category with direct sums as a localizing subcategory. Hence these objects generate the entire étale derived category.
\end{proof}

% text: prose-05-18
\phantomsection\label{bp456}%
Proper direct image already preserves \(D^b_{\mathrm{ctf}}\) by Proposition \hyperref[bp390]{4.35}. Its extension to the localizing categories therefore only needs preservation of arbitrary direct sums. By the comparison just proved, this reduces to the unbounded étale statement.

% object: corollary-proper-direct-image-finite-coefficients-localizing (blueprint/chap-Construction-of-f/sec-Compactness-for-finite-coefficients/corollary-proper-direct-image-finite-coefficients-localizing/_section.md)
\phantomsection\label{bp457}%
\begin{bpstatementit}{Corollary~5.16}{ (Proper direct image on the localizing categories)}
Let \(A\) be a finite admissible ring and let \(p:X\to Y\) be proper. The functor \(Rp_*\) restricts to an exact functor

\[
Rp_*:\Loc(D^b_{\mathrm{ctf}}(X,A))\longrightarrow\Loc(D^b_{\mathrm{ctf}}(Y,A))
\]

that preserves arbitrary direct sums.
\end{bpstatementit}

% proof: proof/corollary-proper-direct-image-finite-coefficients-localizing
\begin{proof}[{Proof}]\phantomsection\label{bp458}%
On unbounded étale derived categories, proper direct image preserves arbitrary direct sums by [Stacks, \href{https://stacks.math.columbia.edu/tag/0F0D}{0F0D}]. Forgetting coefficients gives the \(A\)-linear statement, and Proposition \hyperref[bp454]{5.15} transports it to the pro-étale localizing categories.

Since \(Rp_*\) maps \(D^b_{\mathrm{ctf}}(X,A)\) into \(D^b_{\mathrm{ctf}}(Y,A)\) by Proposition \hyperref[bp390]{4.35}, the objects \(K\in\Loc(D^b_{\mathrm{ctf}}(X,A))\) whose direct image belongs to \(\Loc(D^b_{\mathrm{ctf}}(Y,A))\) form a triangulated subcategory closed under direct sums and containing \(D^b_{\mathrm{ctf}}(X,A)\), that is, all of \(\Loc(D^b_{\mathrm{ctf}}(X,A))\).
\end{proof}


% object: sec-Direct-image-with-compact-support (blueprint/chap-Construction-of-f/sec-Direct-image-with-compact-support/_section.md), content hidden
\subsection*{\texorpdfstring{5.4\quad Direct image with compact support}{5.4 Direct image with compact support}}
\phantomsection\label{bp459}%
\addcontentsline{toc}{subsection}{\texorpdfstring{5.4\quad Direct image with compact support}{5.4 Direct image with compact support}}
% text: prose-05-19
\phantomsection\label{bp460}%
The trace formula needs a relative version of cohomology with compact supports, whose stalks are the groups \(H^i_c\) of the fibres. Consider first a locally compact Hausdorff space \(X\) with an open embedding \(j:X\hookrightarrow\overline X\) into a compact space. A section of \(j_!F\) is a section of \(F\) that vanishes near \(\overline X\setminus X\), and \(H^i_c(X,F)=H^i(\overline X,j_!F)\). For example, for \(X=\mathbf R\) and \(\overline X\) the circle, we find \(H^1_c(\mathbf R,\mathbf Z)=\mathbf Z\).

The definition by sections with compact support has no direct analogue for schemes, because every scheme of finite type over \(k\) is quasi-compact, and the algebraic substitute for compactness is properness, which is a property of a morphism and not of a subset. The two-step form, however, does carry over. In it, the functor \(j_!\) imposes compact support along the open immersion, while for a proper map \(p\), as for a map of spaces with compact fibres, the direct image with compact supports is the ordinary direct image \(Rp_*\). By proper base change, the result does not depend on the compactification (Proposition \hyperref[bp492]{5.26}). Proper base change also identifies the stalks of \(Rf_!\) with the compactly supported cohomology of the fibres (Corollary \bpout{12.26}).

Let \(f:X\to Y\) be a separated morphism of finite type. For a compactification \(f=p\circ j\), with \(j\) an open immersion and \(p\) proper, we define \(Rf_!=Rp_*\circ j_!\). We must show that this functor preserves direct sums. In the following subsection we take \(\Lambda=E\) or \(\overline{\mathbf Q}_\ell\), so that \(\Loc(D^b_{\mathrm{ctf}}(X,\Lambda))=\Loc(\Dbc(X,\Lambda))\).

% object: subsec-Direct-sums (blueprint/chap-Construction-of-f/sec-Direct-image-with-compact-support/subsec-Direct-sums/_section.md), content hidden
\subsubsection*{\texorpdfstring{5.4.1\quad Direct sums}{5.4.1 Direct sums}}
\phantomsection\label{bp461}%
\addcontentsline{toc}{subsubsection}{\texorpdfstring{5.4.1\quad Direct sums}{5.4.1 Direct sums}}
% text: prose-05-20
\phantomsection\label{bp462}%
To apply Brown representability, we need \(Rf_!\) to preserve arbitrary direct sums. Since \(j_!\) is a left adjoint, it does. The question is therefore whether \(Rp_*\) preserves direct sums for proper \(p\). For complexes with a common lower and upper bound, this would follow from finiteness of \(\operatorname{cd}_\ell(p)\). However, the summands of a direct sum in \(\Loc(D^b_{\mathrm{ctf}}(X,\Lambda))\) need not have any common bound. As in the proof of compactness, we use a uniform bound instead. Let \(d\) bound the fibre dimensions of \(p\). We show that \(R^qp_*\) vanishes on \(\IndCons(X,\Lambda)\) for \(q>2d+1\). This means that \(\mathcal H^q(Rp_*K)\) depends only on the cohomology of \(K\) in degrees \([q-2d-1,q]\).

We prove the bound on sections over affine weakly contractible objects and then pass to colimits. For this we use the continuity of pro-étale cohomology, Corollary \hyperref[bp86]{2.29}, applied to sheaves, which are complexes vanishing below degree \(0\). It says that for a quasi-compact and quasi-separated morphism \(p\) the functors \(R^qp_*\) on sheaves of modules commute with filtered colimits. Since a direct sum is the filtered colimit of its finite subsums, they also commute with direct sums.

We also use the following description of \(R^qp_*\). If \(V\to Y\) is an affine weakly contractible object and \(X_V=X\times_YV\), then \(\Gamma(V,-)\) is exact (Corollary corollary-evaluation-weakly-contractible-objects), and hence

\[
(R^qp_*M)(V)=H^q((X_V)_{\mathrm{pro}\acute{e}t},M|_{X_V}).
\tag{5.5}
\]

Since such \(V\) form a basis of \(Y_{\mathrm{pro}\acute{e}t}\), a sheaf on \(Y_{\mathrm{pro}\acute{e}t}\) vanishes if its sections over all of them vanish.

% object: lemma-uniform-bound-proper-direct-image (blueprint/chap-Construction-of-f/sec-Direct-image-with-compact-support/subsec-Direct-sums/lemma-uniform-bound-proper-direct-image/_section.md)
\phantomsection\label{bp463}%
\begin{bpstatementit}{Lemma~5.17}{ (A uniform bound for proper direct image)}
Let \(p:X\to Y\) be proper, and let \(d\) be an integer bounding the dimensions of all geometric fibres of \(p\) (for instance \(d=\dim X\)). For \(M\in\IndCons(X,\Lambda)\),

\[
R^qp_*M=0\quad(q>2d+1),\qquad R^qp_*M\in\IndCons(Y,\Lambda)\quad(q\geq0),
\]

and \(R^qp_*:\IndCons(X,\Lambda)\to\IndCons(Y,\Lambda)\) preserves arbitrary direct sums.
\end{bpstatementit}

% text: prose-05-21
\phantomsection\label{bp464}%
By a stratification we reduce the lemma to the extension by zero of a lisse \(\mathcal O_E\)-lattice on a stratum. Its reductions modulo \(\pi_E^n\) are torsion sheaves, for which the classical bound \(2d\) holds. In the passage to the limit over \(n\), the \(\varprojlim^1\) term raises the bound by one, to \(2d+1\). The remaining steps, inverting \(\ell\), extending scalars to \(\overline{\mathbf Q}_\ell\) and passing to ind-objects, are filtered colimits. By Corollary \hyperref[bp86]{2.29}, the bound survives each of them.

% proof: proof/lemma-uniform-bound-proper-direct-image
\begin{proof}[{Proof of Lemma 5.17}]\phantomsection\label{bp465}%
\emph{Reduction to a lattice.} The vanishing statement is stable under extensions: if \(R^qp_*\) vanishes for \(q>2d+1\) on two sheaves, it vanishes on every extension of one by the other. By Lemma lemma-constructible-sheaves, every \(M\in\Cons(X,\Lambda)\) is a finite extension of sheaves \(a_!L\), with \(a:S\hookrightarrow X\) a locally closed immersion and \(L\) lisse. Moreover, for \(\Lambda=\overline{\mathbf Q}_\ell\), \(L\) is the scalar extension \(\QbarX\otimes_{E}L_E\) of a lisse \(E_S\)-module \(L_E\) for some finite \(E/\mathbf Q_\ell\) (Lemma \hyperref[bp245]{3.17}). For \(\Lambda=E\), we simply put \(L_E=L\). In both cases we may refine \(S\) by a finite stratification, since this replaces \(a_!L\) by a finite extension of sheaves of the same form. By Lemma \hyperref[bp269]{3.24}, after such a refinement \(L_E\) contains a lattice \(L_0\), locally free of finite rank over \(\mathcal O_{E,S}\), with \(L_0[1/\ell]=L_E\). Hence it is enough to prove the bound for \(a_!L_0\), and then to pass to \(a_!L_E=(a_!L_0)[1/\ell]\) and, if \(\Lambda=\overline{\mathbf Q}_\ell\), to its scalar extension \(a_!L\).

In the passage to the limit below we use that the sheaf \(a_!L_0\) is derived \(\pi_E\)-complete. To see this, note that \(a\) is a closed immersion composed with a quasi-compact open immersion. Closed pushforward is exact and preserves products [Stacks, \href{https://stacks.math.columbia.edu/tag/09BM}{09BM}], and extension by zero is exact and preserves limits [Stacks, \href{https://stacks.math.columbia.edu/tag/09BP}{09BP}]. Therefore both functors preserve the derived inverse limit that expresses the completeness of \(L_0\).

\emph{The torsion bound.} Over the objects \(V\) that compute \(R^qp_*\), the reductions \(a_!(L_0/\pi_E^n)\) have cohomology only up to degree \(2d\). Fix an affine weakly contractible \(V\to Y\) and let \(p_V:X_V\to V\) be the base change of \(p\), which is proper with the same bound \(d\). Since \(L_0\) is flat, \(a_!L_0\otimes^{\mathbf L}_{\mathcal O_E}
\mathcal O_E/\pi_E^n\) is the sheaf \(a_!(L_0/\pi_E^n)\), which is the pullback to the pro-étale site of a classical constructible étale sheaf of \(\mathcal O_E/\pi_E^n\)-modules. Since the scheme \(V\) is not Noetherian, we use the proper vanishing theorem in the generality of [Stacks, \href{https://stacks.math.columbia.edu/tag/095U}{095U}]: if a proper morphism has geometric fibres of dimension at most \(d\), then its higher direct images of torsion sheaves vanish in degrees \(q>2d\). With the comparison [Stacks, \href{https://stacks.math.columbia.edu/tag/09A6}{09A6}] of étale and pro-étale direct images, this theorem for \(p_V\) says that \(R^qp_{V,*}\) of this sheaf vanishes for \(q>2d\). As \(V\) is weakly contractible, \(\Gamma(V,-)\) is exact on all pro-étale sheaves [Stacks, \href{https://stacks.math.columbia.edu/tag/0991}{0991}]. Hence \(H^q(X_V,-)=\Gamma(V,R^qp_{V,*}(-))\), and

\[
H^q\bigl(X_V,a_!(L_0/\pi_E^n)\bigr)=0\quad(q>2d).
\]

\emph{The limit.} Derived completeness lets us pass from the reductions to \(a_!L_0\). Since restriction to \(X_V\) is exact and commutes with products, hence with derived inverse limits, completeness gives \((a_!L_0)|_{X_V}\simeq R\varprojlim_n a_!(L_0/\pi_E^n)|_{X_V}\). Derived global sections, too, commute with derived inverse limits, since a product of K-injective complexes is K-injective and sections preserve products. The Milnor exact sequence of Proposition \hyperref[bp71]{2.24}

\[
0\to{\textstyle\varprojlim^1_n}H^{q-1}\bigl(X_V,a_!(L_0/\pi_E^n)\bigr)
\to H^q(X_V,a_!L_0)
\to\varprojlim_nH^q\bigl(X_V,a_!(L_0/\pi_E^n)\bigr)\to0
\]

then shows that \(H^q(X_V,a_!L_0)=0\) for \(q>2d+1\), since for such \(q\) both outer terms vanish by the torsion bound.

\emph{Back to \(\Lambda\).} The remaining passages are filtered colimits. By Corollary \hyperref[bp86]{2.29}, cohomology on the quasi-compact and quasi-separated scheme \(X_V\) commutes with them. For instance, the sheaf \((a_!L_0)[1/\ell]\) is the filtered colimit of \(a_!L_0\) along multiplication by \(\ell\), and so it has the same vanishing. For \(\Lambda=\overline{\mathbf Q}_\ell\), scalar extension from \(E\) to \(\overline{\mathbf Q}_\ell\) is the filtered colimit of the extensions to finite \(E'\supset E\). Since \(E'\) is finite free over \(E\), each of these is, as a sheaf of abelian groups, a finite direct sum of copies of the original sheaf. Hence the continuity lemma again gives the bound. As \(V\) was an arbitrary affine weakly contractible object, formula (5.5) shows that \(R^qp_*(a_!L)=0\) for \(q>2d+1\). By the first step, \(R^qp_*M=0\) for \(q>2d+1\) and every \(M\in\Cons(X,\Lambda)\).

\emph{Passage to \(\IndCons(X,\Lambda)\).} For constructible \(M\) we have \(R^qp_*M\in\Cons(Y,\Lambda)\), because \(Rp_*\) maps \(\Dbc(X,\Lambda)\) into \(\Dbc(Y,\Lambda)\) by Proposition \hyperref[bp405]{4.40}. Moreover, every object of \(\IndCons(X,\Lambda)\) is a filtered colimit of constructible sheaves, and \(R^qp_*\) commutes with filtered colimits and direct sums by Corollary \hyperref[bp86]{2.29}. All three assertions for \(\IndCons(X,\Lambda)\) therefore follow.
\end{proof}

% text: prose-05-22
\phantomsection\label{bp466}%
The bound \(2d+1\) is not sharp. For example, for the structure map of a smooth projective curve over \(\bar k\), where \(d=1\), the lemma only says that cohomology vanishes above degree \(3\), while in fact it vanishes above degree \(2\). The extra degree does no harm, because the arguments below only need a bound that does not depend on the sheaf. As in the compactness argument, such a bound passes from sheaves to unbounded complexes.

% object: lemma-proper-direct-image-unbounded-complexes (blueprint/chap-Construction-of-f/sec-Direct-image-with-compact-support/subsec-Direct-sums/lemma-proper-direct-image-unbounded-complexes/_section.md)
\phantomsection\label{bp467}%
\begin{bpstatementit}{Lemma~5.18}{ (Proper direct image of unbounded complexes)}
Let \(p:X\to Y\) be proper, and let \(d\) bound the dimensions of the geometric fibres of \(p\). For integers \(a\) and \(b\), write \(\Dic^{\leq b}(X,\Lambda)\) and \(\Dic^{\geq a}(X,\Lambda)\) for the objects of \(\Dic(X,\Lambda)\) (Corollary \hyperref[bp439]{5.11}) with cohomology in degrees \(\leq b\), respectively \(\geq a\).

\begin{enumerate}
\item \(Rp_*\) maps \(\Dic(X,\Lambda)\) into \(\Dic(Y,\Lambda)\), with \(Rp_*\Dic^{\leq b}(X,\Lambda)\subset\Dic^{\leq b+2d+1}(Y,\Lambda)\) and \(Rp_*\Dic^{\geq a}(X,\Lambda)\subset\Dic^{\geq a}(Y,\Lambda)\).
\item For \(K\in\Dic(X,\Lambda)\) and every integer \(q\), the truncation maps \(K\leftarrow\tau_{\leq q}K\to\tau_{\geq q-2d-1}\tau_{\leq q}K\) induce an isomorphism in \(\IndCons(Y,\Lambda)\) \[
   \mathcal H^q(Rp_*K)\simeq
   \mathcal H^q\bigl(Rp_*\tau_{\geq q-2d-1}\tau_{\leq q}K\bigr).
   \tag{5.6}
   \]
\item \(Rp_*:\Dic(X,\Lambda)\to\Dic(Y,\Lambda)\) preserves arbitrary direct sums.
\end{enumerate}
\end{bpstatementit}

% proof: proof/lemma-proper-direct-image-unbounded-complexes
\begin{proof}[{Proof}]\phantomsection\label{bp468}%
For bounded below \(K\) everything is read off from the hypercohomology spectral sequence, which has only \(2d+2\) nonzero columns. A general \(K\), in turn, is the derived limit of its truncations \(\tau_{\geq-n}K\), and the column bound shows that in each fixed degree the tower of direct images is eventually constant. The window (5.6) and the preservation of direct sums follow from this.

\emph{Bounded below \(K\).} The hypercohomology spectral sequence

\[
E_2^{u,v}=R^up_*\mathcal H^v(K)\Longrightarrow\mathcal H^{u+v}(Rp_*K)
\]

converges. By Lemma \hyperref[bp463]{5.17}, its only nonzero columns are \(0\leq u\leq2d+1\), and their terms lie in \(\IndCons(Y,\Lambda)\). So each \(\mathcal H^n(Rp_*K)\) has a finite filtration whose graded pieces are subquotients of objects of \(\IndCons(Y,\Lambda)\). Since this category is closed under kernels, cokernels and extensions by Proposition \hyperref[bp433]{5.9}, we get \(\mathcal H^n(Rp_*K)\in\IndCons(Y,\Lambda)\). Both amplitude bounds can be read off from the columns.

\emph{Arbitrary \(K\).} By Proposition \hyperref[bp77]{2.26}, \(K\simeq R\varprojlim_n
\tau_{\geq-n}K\). Since \(Rp_*\) is a right adjoint, it preserves products and hence derived inverse limits. Thus \(Rp_*K\simeq R\varprojlim_nRp_*\tau_{\geq-n}K\), and it suffices to see that in each fixed degree the tower \(Rp_*\tau_{\geq-n}K\) stabilizes. The fibre of \(\tau_{\geq-n-1}K\to\tau_{\geq-n}K\) is \(\mathcal H^{-n-1}(K)[n+1]\), a single ind-constructible sheaf in degree \(-n-1\). By the bounded below case its direct image has cohomology only in degrees \([-n-1,-n-1+2d+1]=[-n-1,2d-n]\). In particular it has none in degrees \(q-1\), \(q\), \(q+1\) once \(n>2d+1-q\). For fixed \(q\), the long exact sequences then show that the transition maps of the towers \(\mathcal H^q(Rp_*\tau_{\geq-n}K)\) and \(\mathcal H^{q-1}(Rp_*\tau_{\geq-n}K)\) are isomorphisms for \(n\) large. Since an eventually constant tower has vanishing \(\varprojlim^1\), the Milnor sequence of Proposition \hyperref[bp71]{2.24} gives

\[
\mathcal H^q(Rp_*K)\simeq
\mathcal H^q(Rp_*\tau_{\geq-n}K)\quad\text{for all sufficiently large }n.
\]

The bounded below case, applied to the right-hand side, shows that \(\mathcal H^q(Rp_*K)\in\IndCons(Y,\Lambda)\), with the bound \(b+2d+1\) if \(K\in
\Dic^{\leq b}(X,\Lambda)\). The lower bound needs no new argument, because an object of \(\Dic^{\geq a}\) is bounded below.

\emph{The isomorphism (5.6).} It comes from the bounds, applied to the truncation triangles

\[
\tau_{\leq q}K\to K\to\tau_{\geq q+1}K,
\qquad
\tau_{\leq q-2d-2}K\to\tau_{\leq q}K\to\tau_{\geq q-2d-1}\tau_{\leq q}K.
\]

The direct image of \(\tau_{\geq q+1}K\) lives in degrees \(\geq q+1\), and that of \(\tau_{\leq q-2d-2}K\) in degrees \(\leq q-1\). The long exact cohomology sequences of the two triangles therefore yield isomorphisms \(\mathcal H^q(Rp_*\tau_{\leq q}K)\simeq\mathcal H^q(Rp_*K)\) and \(\mathcal H^q(Rp_*\tau_{\leq q}K)\simeq
\mathcal H^q(Rp_*\tau_{\geq q-2d-1}\tau_{\leq q}K)\).

\emph{Direct sums, assertion (3).} The window (5.6) reduces us to families with cohomology in a common finite interval, and for such a family \((M_\alpha)\) we show by induction on the length of the interval that the map \(\bigoplus_\alpha Rp_*M_\alpha\to Rp_*(\bigoplus_\alpha M_\alpha)\) is an isomorphism. When the length is zero, each \(M_\alpha\) is a shifted object of \(\IndCons(X,\Lambda)\), and the claim is the assertion about direct sums in Lemma \hyperref[bp463]{5.17}, in each degree. Then, for the inductive step, we take the direct sum of the truncation triangles that split off the top cohomology sheaf and apply the five lemma to the long exact cohomology sequences.

Now let \((K_\alpha)\) be arbitrary and fix \(q\). Since direct sums of sheaves are exact, cohomology and truncations commute with direct sums. The objects \(\tau_{\geq q-2d-1}\tau_{\leq q}K_\alpha\) have cohomology in the common interval \([q-2d-1,q]\). Applying (5.6) to each \(K_\alpha\) and to \(\bigoplus_\alpha K_\alpha\), and the first paragraph to the truncated family, we see that \(\bigoplus_\alpha Rp_*K_\alpha\to Rp_*(\bigoplus_\alpha K_\alpha)\) is an isomorphism on \(\mathcal H^q\). Since \(q\) is arbitrary, the map is an isomorphism.
\end{proof}


% object: subsec-Integral-coefficients (blueprint/chap-Construction-of-f/sec-Direct-image-with-compact-support/subsec-Integral-coefficients/_section.md), content hidden
\subsubsection*{\texorpdfstring{5.4.2\quad Integral coefficients}{5.4.2 Integral coefficients}}
\phantomsection\label{bp469}%
\addcontentsline{toc}{subsubsection}{\texorpdfstring{5.4.2\quad Integral coefficients}{5.4.2 Integral coefficients}}
% text: prose-05-23
\phantomsection\label{bp470}%
For \(\mathcal O_E\) the rational arguments for compactness (Theorem \hyperref[bp446]{5.13}) and for direct sums (Lemma \hyperref[bp467]{5.18}) go through with one change. Integral constructible sheaves are no longer locally free on strata; for example, \(\mathcal O_{E,X}/\pi_E\) is constructible. But \(\mathcal O_E\) is a discrete valuation ring, so they have free resolutions of length one. The estimates only use the resulting Tor amplitude \([-1,0]\).

% object: lemma-structure-integral-constructible-complexes (blueprint/chap-Construction-of-f/sec-Direct-image-with-compact-support/subsec-Integral-coefficients/lemma-structure-integral-constructible-complexes/_section.md)
\phantomsection\label{bp471}%
\begin{bpstatementit}{Lemma~5.19}{ (Structure of integral constructible complexes)}
\begin{enumerate}
\item \(\Dbc(X,\mathcal O_E)\) is a triangulated subcategory of \(D(X_{\mathrm{pro}\acute{e}t},\mathcal O_E)\) of bounded objects, stable under truncation; its objects in degree \(0\) form the category \(\Cons(X,\mathcal O_E)\), an essentially small abelian subcategory of \(\operatorname{Mod}(X_{\mathrm{pro}\acute{e}t},\mathcal O_E)\) closed under extensions.
\item For \(C\in\Cons(X,\mathcal O_E)\) there is a finite stratification on whose strata \(C\) is pro-étale locally \(\mathcal O_{E,S}\otimes_{\mathcal O_E}P\), with \(P\) finitely generated.
\item On each such stratum \(C\) is dualizable of Tor amplitude \([-1,0]\), and its dual has Tor amplitude \([0,1]\).
\end{enumerate}
\end{bpstatementit}

% proof: proof/lemma-structure-integral-constructible-complexes
\begin{proof}[{Proof}]\phantomsection\label{bp472}%
The conditions of Proposition \hyperref[bp226]{3.11}, derived completeness and constructible reduction, are stable under shifts and cones, so \(\Dbc(X,\mathcal O_E)\) is triangulated. By Proposition \hyperref[bp223]{3.10}, on the strata of a finite stratification an object is pro-étale locally of the form \(\mathcal O_{E,S}\otimes P^\bullet\) with \(P^\bullet\) a bounded complex of finite free \(\mathcal O_E\)-modules. Since \(\mathcal O_{E,S}\) is flat over \(\mathcal O_E\), we get \(H^i=\mathcal O_{E,S}\otimes H^i(P^\bullet)\) on such a stratum. This local form of the cohomology sheaves shows that the truncations stay in \(\Dbc\). It also gives (2), and then the abelian part of (1). In (1), essential smallness holds because a constructible sheaf is determined up to isomorphism by a stratification, local presentations, descent data and finitely many extension classes. Finally, for (3) we use that a finitely generated module over the discrete valuation ring \(\mathcal O_E\) has a free resolution \(0\to\mathcal O_E^r\to\mathcal O_E^s\to P\to0\) of length one.
\end{proof}

% text: prose-05-24
\phantomsection\label{bp473}%
For instance, \(\mathcal O_{E,X}/\pi_E\) is the cokernel of \(\pi_E:\mathcal O_{E,X}\to\mathcal O_{E,X}\), so it has Tor amplitude \([-1,0]\). Its dual \(\cHom(\mathcal O_{E,X}/\pi_E,\mathcal O_{E,X})=
(\mathcal O_{E,X}/\pi_E)[-1]\) has Tor amplitude \([0,1]\), as in (3).

The dévissage in the proof of compactness also needs that the operations it uses preserve \(\Dbc(-,\mathcal O_E)\).

% object: lemma-inverse-image-tensor-product (blueprint/chap-Construction-of-f/sec-Direct-image-with-compact-support/subsec-Integral-coefficients/lemma-inverse-image-tensor-product/_section.md)
\phantomsection\label{bp474}%
\begin{bpstatementit}{Lemma~5.20}{ (Inverse image and tensor product)}
Inverse image, \(\otimes^{\mathbf L}_{\mathcal O_{E,X}}\) without completion, extension by zero along a quasi-compact open immersion and direct image along a closed immersion preserve \(\Dbc(-,\mathcal O_E)\); on \(\Dbc\) inverse image agrees with the completed inverse image of the Stacks Project.
\end{bpstatementit}

% proof: proof/lemma-inverse-image-tensor-product
\begin{proof}[{Proof}]\phantomsection\label{bp475}%
The preservation statements follow from Lemma lemma-integral-models, since closed direct image is exact and preserves limits [Stacks, \href{https://stacks.math.columbia.edu/tag/09AA}{09AA}, \href{https://stacks.math.columbia.edu/tag/09BM}{09BM}]. The two inverse images agree on objects that are pro-étale locally of the form \(\mathcal O_{E,S}\otimes P\) with \(P\) finitely generated, and by the local structure of Lemma lemma-structure-integral-constructible-complexes and (2) they therefore agree on \(\Dbc\).
\end{proof}

% text: prose-05-25
\phantomsection\label{bp476}%
With these two lemmas, the integral case of the two hypotheses of Brown representability can be checked by following the rational proofs.

% object: proposition-compactness-proper-direct-image (blueprint/chap-Construction-of-f/sec-Direct-image-with-compact-support/subsec-Integral-coefficients/proposition-compactness-proper-direct-image/_section.md)
\phantomsection\label{bp477}%
\begin{bpstatementit}{Proposition~5.21}{ (Compactness and proper direct image)}
Every object of \(D^b_{\mathrm{ctf}}(X,\mathcal O_E)=\Dbc(X,\mathcal O_E)\) is compact in \(\Loc(D^b_{\mathrm{ctf}}(X,\mathcal O_E))\), and for proper \(p:X\to Y\) the functor \(Rp_*\) preserves direct sums of objects of \(\Loc(D^b_{\mathrm{ctf}}(X,\mathcal O_E))\).
\end{bpstatementit}

% proof: proof/proposition-compactness-proper-direct-image
\begin{proof}[{Proof}]\phantomsection\label{bp478}%
The rational proofs (Theorem \hyperref[bp446]{5.13}, Lemmas \hyperref[bp443]{5.12} and \hyperref[bp467]{Lemma 5.18}) go through once we supply for \(\mathcal O_E\) the four inputs they use: the abelian structure of Lemma lemma-structure-integral-constructible-complexes, continuity of \(\Ext^i(C,-)\) for bounded below systems, a bound \(d_C\) with \(\Ext^i(C,N)=0\) for \(i>d_C\) and \(N\) constructible, and a bound on the amplitude of \(Rp_*\).

For continuity, take \(C=a_!L\) on a stratum. By Lemma lemma-sections-support-closed-subscheme, \(R\!\operatorname{Hom}(a_!L,-)\) is \(R\Gamma(S,-)\) of the dual of \(L\) tensored with \(u^*R\underline\Gamma_Z\). The proof of that formula uses only that \(L\) is dualizable, which holds by Lemma lemma-structure-integral-constructible-complexes, and so continuity follows from Corollary \hyperref[bp86]{2.29}.

The bound \(d_C\) comes from counting degrees as in Lemma \hyperref[bp426]{5.6}. Namely, the reductions of \(N\) lie in \([-1,0]\), and the dual of \(L\) adds one degree. Next, the functors \(R\Gamma\) and \(R\underline\Gamma_Z\) add the cohomological dimensions of Definition \hyperref[bp137]{2.45}. Finally, the derived limit over the reductions adds one more [Stacks, \href{https://stacks.math.columbia.edu/tag/08TB}{08TB}].

The same count bounds \(Rp_*\). For fibres of dimension \(\le d\), the torsion bound [Stacks, \href{https://stacks.math.columbia.edu/tag/095U}{095U}] gives \(2d\), and the derived limit adds one. Hence \(Rp_*\) of a constructible sheaf lies in \([0,2d+1]\).
\end{proof}

% text: prose-05-26
\phantomsection\label{bp479}%
These are the facts about \(\mathcal O_E\) that we use in the construction of \(Rf_!\) below.


% object: subsec-Construction-of-Rf (blueprint/chap-Construction-of-f/sec-Direct-image-with-compact-support/subsec-Construction-of-Rf/_section.md), content hidden
\subsubsection*{\texorpdfstring{5.4.3\quad Construction of $Rf_!$}{5.4.3 Construction of Rf!}}
\phantomsection\label{bp480}%
\addcontentsline{toc}{subsubsection}{\texorpdfstring{5.4.3\quad Construction of $Rf_!$}{5.4.3 Construction of Rf!}}
% text: prose-05-27
\phantomsection\label{bp481}%
From now on, \(\Lambda\) is any admissible coefficient ring. Lemmas \hyperref[bp463]{5.17} and \hyperref[bp467]{Lemma 5.18} only treat \(E\) and \(\overline{\mathbf Q}_\ell\), but we use them only through the fact that \(Rp_*\) preserves direct sums for proper \(p\). For finite and integral coefficients this fact was proved separately, and the next proof collects the references.

Proper base change shows, by one and the same argument, that \(Rp_*\circ j_!\) does not depend on the compactification and is compatible with composition and base change. The comparison maps are maps out of an extension by zero \(j_!K\). They are therefore determined by their restriction to \(X\), where they are the identity. On the closed complement of \(X\), proper base change shows that both sides vanish, and so the map is an isomorphism. For instance, if \(h:Z'\to Z\) is proper and \(h^{-1}(X)=X\), then \(Rh_*j'_!K\) restricted to \(Z\setminus X\) is the direct image of the restriction of \(j'_!K\) to \(Z'\setminus X\), which is zero. Since proper base change is only known on \(D^b_{\mathrm{ctf}}\), we first show that all the functors involved preserve direct sums, so that it suffices to check isomorphisms on \(D^b_{\mathrm{ctf}}\) (Lemma \hyperref[bp485]{5.23}).

That both factors of \(Rp_*\circ j_!\) preserve the localizing categories comes from the definition of \(\Loc(D^b_{\mathrm{ctf}}(X,\Lambda))\) in the following form: if the objects of \(\Loc(D^b_{\mathrm{ctf}}(X,\Lambda))\) with some property form a triangulated subcategory that is closed under direct sums and contains \(D^b_{\mathrm{ctf}}(X,\Lambda)\), then every object of \(\Loc(D^b_{\mathrm{ctf}}(X,\Lambda))\) has the property.

% object: proposition-proper-direct-image-extension-zero (blueprint/chap-Construction-of-f/sec-Direct-image-with-compact-support/subsec-Construction-of-Rf/proposition-proper-direct-image-extension-zero/_section.md)
\phantomsection\label{bp482}%
\begin{bpstatementit}{Proposition~5.22}{ (Proper direct image and extension by zero)}
For every admissible \(\Lambda\), if \(p:X\to Y\) is proper, \(Rp_*:D(X_{\mathrm{pro}\acute{e}t},\Lambda)\to D(Y_{\mathrm{pro}\acute{e}t},\Lambda)\) restricts to an exact functor \(\Loc(D^b_{\mathrm{ctf}}(X,\Lambda))\to\Loc(D^b_{\mathrm{ctf}}(Y,\Lambda))\) that preserves direct sums. If \(j:U\to X\) is an open immersion, \(j_!:D(U_{\mathrm{pro}\acute{e}t},\Lambda)\to D(X_{\mathrm{pro}\acute{e}t},\Lambda)\) restricts to an exact functor \(\Loc(D^b_{\mathrm{ctf}}(U,\Lambda))\to\Loc(D^b_{\mathrm{ctf}}(X,\Lambda))\) that preserves direct sums and maps \(D^b_{\mathrm{ctf}}(U,\Lambda)\) into \(D^b_{\mathrm{ctf}}(X,\Lambda)\).
\end{bpstatementit}

% proof: proof/proposition-proper-direct-image-extension-zero
\begin{proof}[{Proof}]\phantomsection\label{bp483}%
For each functor we need two facts: it preserves direct sums, and it maps \(D^b_{\mathrm{ctf}}\) into \(D^b_{\mathrm{ctf}}\). The rest then follows from the definition of the localizing categories, in the form recalled before the proposition.

For \(Rp_*\) with \(E\) and \(\overline{\mathbf Q}_\ell\), Lemma lemma-proper-direct-image-unbounded-complexes shows that \(Rp_*\) preserves direct sums on \(\Dic(X,\Lambda)\supset\Loc(D^b_{\mathrm{ctf}}(X,\Lambda))\), and Proposition \hyperref[bp405]{4.40} shows that \(Rp_*\) preserves \(D^b_{\mathrm{ctf}}=\Dbc\). For finite \(A\), these two facts are Proposition \hyperref[bp390]{4.35} and Corollary \hyperref[bp457]{5.16}; for \(\mathcal O_E\), they are Proposition \hyperref[bp477]{5.21} and Proposition \hyperref[bp405]{4.40}. In each case, the objects \(K\in\Loc(D^b_{\mathrm{ctf}}(X,\Lambda))\) with \(Rp_*K\in\Loc(D^b_{\mathrm{ctf}}(Y,\Lambda))\) form a triangulated subcategory closed under direct sums and containing \(D^b_{\mathrm{ctf}}(X,\Lambda)\), and hence they form all of \(\Loc(D^b_{\mathrm{ctf}}(X,\Lambda))\).

Extension by zero is exact and left adjoint to \(j^*\), so it preserves direct sums. It preserves constructibility, since we may add the closed complement as a stratum on which the sheaf is zero. For finite \(A\) it also preserves finite Tor amplitude, as one checks on stalks, while for \(\mathcal O_E\), \(E\) and \(\overline{\mathbf Q}_\ell\) we have \(D^b_{\mathrm{ctf}}=\Dbc\), so constructibility suffices; for \(\mathcal O_E\) this is Lemma \hyperref[bp474]{5.20}. With this, the localizing-subcategory argument shows that \(j_!\) restricts as stated.
\end{proof}

% text: prose-05-28
\phantomsection\label{bp484}%
The proper base change map and the comparison maps between compactifications are natural transformations between exact functors on \(\Loc(D^b_{\mathrm{ctf}}(-,\Lambda))\) that preserve direct sums. To show that such a transformation is an isomorphism, it suffices to check this on \(D^b_{\mathrm{ctf}}(-,\Lambda)\).

% object: lemma-checking-isomorphisms-constructible-complexes (blueprint/chap-Construction-of-f/sec-Direct-image-with-compact-support/subsec-Construction-of-Rf/lemma-checking-isomorphisms-constructible-complexes/_section.md)
\phantomsection\label{bp485}%
\begin{bpstatementit}{Lemma~5.23}{ (Checking isomorphisms on constructible complexes)}
Let \(\Lambda\) be admissible, let \(F_1,F_2:\Loc(D^b_{\mathrm{ctf}}(X,\Lambda))\to\Loc(D^b_{\mathrm{ctf}}(Y,\Lambda))\) be exact functors preserving direct sums, and let \(\eta:F_1\to F_2\) be a natural transformation of exact functors. If \(\eta_K\) is an isomorphism for every \(K\in D^b_{\mathrm{ctf}}(X,\Lambda)\), then it is an isomorphism for every \(K\in\Loc(D^b_{\mathrm{ctf}}(X,\Lambda))\).
\end{bpstatementit}

% proof: proof/lemma-checking-isomorphisms-constructible-complexes
\begin{proof}[{Proof}]\phantomsection\label{bp486}%
The objects \(K\) for which \(\eta_K\) is an isomorphism are closed under shifts and cones by the five lemma, and under direct sums because \(F_1\) and \(F_2\) preserve them. Since they contain \(D^b_{\mathrm{ctf}}(X,\Lambda)\), they form all of \(\Loc(D^b_{\mathrm{ctf}}(X,\Lambda))\).
\end{proof}

% text: prose-05-29
\phantomsection\label{bp487}%
For a cartesian square with \(p\) proper, the base change map \(g^*Rp_*\to Rp'_*g'^*\) of (4.5) is a map of functors \(D(X_{\mathrm{pro}\acute{e}t},\Lambda)\to D(Y'_{\mathrm{pro}\acute{e}t},\Lambda)\). The lemma lets us pass from \(D^b_{\mathrm{ctf}}(X,\Lambda)\), where it is an isomorphism, to \(\Loc(D^b_{\mathrm{ctf}}(X,\Lambda))\).

% object: proposition-proper-base-change-localizing-categories (blueprint/chap-Construction-of-f/sec-Direct-image-with-compact-support/subsec-Construction-of-Rf/proposition-proper-base-change-localizing-categories/_section.md)
\phantomsection\label{bp488}%
\begin{bpstatementit}{Proposition~5.24}{ (Proper base change on the localizing categories)}
Let \(\Lambda\) be admissible. Consider a cartesian square

\[
\begin{array}{ccc}
X'&\xrightarrow{g'}&X\\
\downarrow p'&&\downarrow p\\
Y'&\xrightarrow{g}&Y
\end{array}
\]

with \(p\) proper. The base change map \(g^*Rp_*K\to Rp'_*g'^*K\) is an isomorphism for every \(K\in\Loc(D^b_{\mathrm{ctf}}(X,\Lambda))\).
\end{bpstatementit}

% proof: proof/proposition-proper-base-change-localizing-categories
\begin{proof}[{Proof}]\phantomsection\label{bp489}%
On \(D^b_{\mathrm{ctf}}(X,\Lambda)\) the base change map is an isomorphism by Proposition \hyperref[bp405]{4.40}, so by Lemma \hyperref[bp485]{5.23} it suffices to show that \(g^*Rp_*\) and \(Rp'_*g'^*\) are exact functors \(\Loc(D^b_{\mathrm{ctf}}(X,\Lambda))\to\Loc(D^b_{\mathrm{ctf}}(Y',\Lambda))\) preserving direct sums. For \(Rp_*\) and \(Rp'_*\) this is Proposition \hyperref[bp482]{5.22}, so it remains to treat inverse image. It is exact and preserves direct sums, and it sends \(D^b_{\mathrm{ctf}}(-,\Lambda)\) into \(D^b_{\mathrm{ctf}}(-,\Lambda)\). Indeed, for finite \(A\) this holds because of bounded flat constructible representatives, and for \(\mathcal O_E\), \(E\) and \(\overline{\mathbf Q}_\ell\) it holds because pullback preserves the local description of constructible sheaves on strata (Lemma \hyperref[bp474]{5.20} in the integral case). Given this, the localizing-subcategory argument before Proposition \hyperref[bp482]{5.22} shows that inverse image also preserves \(\Loc(D^b_{\mathrm{ctf}}(-,\Lambda))\).
\end{proof}

% object: remark-compactifications (blueprint/chap-Construction-of-f/sec-Direct-image-with-compact-support/subsec-Construction-of-Rf/remark-compactifications/_section.md)
\phantomsection\label{bp490}%
\begin{bpstatement}{Remark~5.25}{ (Compactifications)}
\begin{enumerate}
\item By Nagata's theorem (Theorem \bpout{A.2}), every separated morphism of finite type \(f:X\to Y\) with \(Y\) qcqs has a compactification, that is, a factorization \(f=p\circ j\) with \(j:X\to\overline X\) a quasi-compact open immersion and \(p:\overline X\to Y\) proper.
\item In comparisons we may replace \(\overline X\) by the scheme-theoretic closure of \(X\), so that \(X\) is dense, and by Lemma \bpout{A.5} this does not change \(Rp_*j_!\).
\end{enumerate}
\end{bpstatement}

% text: prose-05-30
\phantomsection\label{bp491}%
The comparisons between compactifications below are the proper direct images of the unique maps out of \(j_!K\) that restrict to the identity on \(X\). They are unique because, for an open immersion \(j:X\to Z\), the functor \(j_!\) is left adjoint to \(j^*\) on the pro-étale derived categories, so that a map out of \(j_!K\) is determined by its restriction to \(X\).

% object: proposition-construction-independence-rf-shriek (blueprint/chap-Construction-of-f/sec-Direct-image-with-compact-support/subsec-Construction-of-Rf/proposition-construction-independence-rf-shriek/_section.md)
\phantomsection\label{bp492}%
\begin{bpstatementit}{Proposition~5.26}{ (Construction and independence of $Rf_!$)}
Let \(\Lambda\) be admissible and let \(f:X\to Y\) be separated of finite type, with compactification \(f=p\circ j\) as above. The functor

\[
Rf_!=Rp_*\circ j_!:\Loc(D^b_{\mathrm{ctf}}(X,\Lambda))\to\Loc(D^b_{\mathrm{ctf}}(Y,\Lambda))
\]

is exact, preserves direct sums, and maps \(D^b_{\mathrm{ctf}}(X,\Lambda)\) into \(D^b_{\mathrm{ctf}}(Y,\Lambda)\). Two compactifications give canonically isomorphic functors, and these isomorphisms satisfy the cocycle condition.
\end{bpstatementit}

% proof: proof/proposition-construction-independence-rf-shriek
\begin{proof}[{Proof}]\phantomsection\label{bp493}%
For a single compactification there is nothing new to prove, and two compactifications are compared through a third one that refines both, by the argument described at the start of this subsection.

The properties of \(Rp_*\circ j_!\) for a single compactification follow from Proposition \hyperref[bp482]{5.22} and the preservation of \(D^b_{\mathrm{ctf}}(-,\Lambda)\) by proper direct image (Proposition \hyperref[bp405]{4.40}).

For independence it suffices to compare a compactification with a refinement of it. Indeed, the schematic closure of \(X\) in a compactification is a refinement of it, since its inverse image over \(X\) is \(X\). Moreover, any two dense compactifications have a common refinement (Proposition \bpout{A.7}).

\emph{Refinements.} Let \(j:X\to Z\) and \(j':X\to Z'\) be compactifications over \(Y\), and let \(h:Z'\to Z\) be a refinement (Definition \bpout{A.1}), that is, a proper \(Y\)-morphism with \(h\circ j'=j\) and \(h^{-1}(X)=X\). To construct the comparison map, note that restriction to the open subscheme \(X\) commutes with \(Rh_*\). This holds because \(h^{-1}(X)=X\) and restriction is exact and preserves K-injectives, its left adjoint \(j_!\) being exact. As \(j'^*j'_!K=K\), we get \(j^*Rh_*j'_!K\simeq K\). By adjunction this isomorphism corresponds to a map

\[
e_{h,K}:j_!K\longrightarrow Rh_*j'_!K
\tag{5.7}
\]

whose restriction to \(X\) is the identity. We claim that it is an isomorphism for \(K\in\Loc(D^b_{\mathrm{ctf}}(X,\Lambda))\). By Proposition \hyperref[bp482]{5.22}, both sides are exact functors of \(K\) preserving direct sums, so by Lemma \hyperref[bp485]{5.23} we may assume that \(K\in D^b_{\mathrm{ctf}}(X,\Lambda)\), where proper base change is available.

By the exact sequence for an open and its closed complement, applied in each degree [Stacks, \href{https://stacks.math.columbia.edu/tag/09AH}{09AH}], a morphism in \(D(Z_{\mathrm{pro}\acute{e}t},\Lambda)\) that is an isomorphism on \(X\) and on \(Z\setminus X\) is an isomorphism. On \(X\) our map is the identity, so it remains to see that on \(Z\setminus X\) both sides vanish. For \(j_!K\) this holds by definition of extension by zero, while for \(Rh_*j'_!K\) we use proper base change on \(D^b_{\mathrm{ctf}}\) (Proposition \hyperref[bp488]{5.24}) for the inclusion of \(Z\setminus X\) into \(Z\). It shows that the restriction of \(Rh_*j'_!K\) is the direct image of the restriction of \(j'_!K\) to \(h^{-1}(Z\setminus X)=Z'\setminus X\), which is zero. This proves the claim. Applying \(Rp_*\) and using \(R(p\circ h)_*=Rp_*Rh_*\), we obtain an isomorphism between the two candidate functors \(Rf_!\), namely \(Rp_*j_!\) and \(R(ph)_*j'_!\).

If \(Z''\xrightarrow{k}Z'\xrightarrow{h}Z\) are two successive refinements, with \(j'':X\to Z''\) the open immersion, then \(Rh_*(e_{k,K})\circ e_{h,K}\) and \(e_{hk,K}\) are maps from \(j_!K\) to \(R(hk)_*j''_!K\) that both restrict to the identity on \(X\). By the observation before the proposition, they are equal.

\emph{Common refinements.} Two arbitrary compactifications need not be related by a refinement, but they have a common one. Indeed, given dense compactifications \(Z_1\) and \(Z_2\), Proposition \bpout{A.7} provides a dense compactification \(Z\) with proper maps \(h_i:Z\to Z_i\) satisfying \(h_i^{-1}(X)=X\), and similarly for finitely many compactifications. The refinement isomorphisms for \(h_1\) and \(h_2\) then compose to an isomorphism between the functors defined by \(Z_1\) and \(Z_2\), and transitivity of the refinement maps on a common refinement shows that this comparison does not depend on the chosen refinement and satisfies the cocycle condition.
\end{proof}

% object: example-cuspidal-compactification (blueprint/chap-Construction-of-f/sec-Direct-image-with-compact-support/subsec-Construction-of-Rf/example-cuspidal-compactification/_section.md)
\phantomsection\label{bp494}%
\begin{bpstatement}{Example~5.27}{ (A singular compactification of the line)}
Over \(k\), let \(C\subset\mathbf P^2\) be the cuspidal cubic \(zy^2=x^3\), with cusp \(c=(0:0:1)\). Its normalization is the map \(h:\mathbf P^1\to C\), \((s:u)\mapsto(su^2:u^3:s^3)\). It sends \(u=0\) to the cusp and the affine line \(u\neq0\) isomorphically onto \(C\setminus\{c\}\). So \(X=\mathbf A^1\) has the two compactifications \(\mathbf P^1\) and \(C\). Since \(h^{-1}(C\setminus\{c\})=\mathbf A^1\), the map \(h\) is a refinement. The refinement map \(j_!\Lambda\to Rh_*j'_!\Lambda\) is the identity on \(X\), and at the cusp both stalks vanish: for the right side, proper base change shows that the stalk of \(Rh_*j'_!\Lambda\) at \(c\) is the cohomology of \(j'_!\Lambda\) on the single point \(h^{-1}(c)=\{\infty\}\), where \(j'_!\Lambda\) is zero. So the singularity of \(C\) does not affect \(R\Gamma_c(\mathbf A^1,\Lambda)\).
\end{bpstatement}

% text: prose-05-31
\phantomsection\label{bp495}%
To make \(Rf_!\) functorial in \(f\), and for the base change theorems later on, we need composition and base change isomorphisms. The base change map is defined as follows. Let \(f:X\to Y\) be separated of finite type with compactification \(f=p\circ j\), let \(v:Y'\to Y\) be a morphism, and let \(f':X'\to Y'\), \(u:X'\to X\), \(p':\overline X'\to Y'\), \(j':X'\to\overline X'\) and \(v':\overline X'\to\overline X\) be the base changes, so that \(f'=p'\circ j'\) is a compactification. For \(K\in D(X_{\mathrm{pro}\acute{e}t},\Lambda)\), the isomorphism \(j'^*v'^*j_!K\simeq u^*j^*j_!K=u^*K\) corresponds under \(j'_!\dashv j'^*\) to a map \(j'_!u^*K\to v'^*j_!K\). This map is an isomorphism: it is the identity on the open \(X'\), and both sides vanish on the closed complement \(v'^{-1}(\overline X\setminus X)=\overline X'\setminus X'\) [Stacks, \href{https://stacks.math.columbia.edu/tag/09AH}{09AH}]. The base change map for \(Rf_!\) is the composite

\[
v^*Rp_*j_!K\longrightarrow Rp'_*v'^*j_!K\simeq Rp'_*j'_!u^*K
\tag{5.8}
\]

of the base change map (4.5) for \(p\) and \(Rp'_*\) of the inverse of this isomorphism.

% object: proposition-composition-base-change-rf-shriek (blueprint/chap-Construction-of-f/sec-Direct-image-with-compact-support/subsec-Construction-of-Rf/proposition-composition-base-change-rf-shriek/_section.md)
\phantomsection\label{bp496}%
\begin{bpstatementit}{Proposition~5.28}{ (Composition and base change for $Rf_!$)}
For every admissible \(\Lambda\), the following statements hold on \(\Loc(D^b_{\mathrm{ctf}}(-,\Lambda))\) and restrict to \(D^b_{\mathrm{ctf}}(-,\Lambda)\).

\begin{enumerate}
\item For separated morphisms of finite type \(X\xrightarrow fY
  \xrightarrow gZ\) there is a canonical isomorphism \(Rg_!\circ Rf_!\simeq R(gf)_!\). These isomorphisms are associative and compatible with the identity \(R(\mathrm{id})_!=\mathrm{id}\).
\item Let \(v:Y'\to Y\) be a morphism, \(f':X'\to Y'\) the base change of \(f\) along \(v\), and \(u:X'\to X\) the projection. The base change map (5.8) is an isomorphism \(v^*Rf_!\simeq Rf'_!u^*\). It does not depend on the compactification, and it is compatible with composing cartesian squares and with the composition isomorphisms of (1).
\end{enumerate}
\end{bpstatementit}

% proof: proof/proposition-composition-base-change-rf-shriek
\begin{proof}[{Proof}]\phantomsection\label{bp497}%
For composition, the only new ingredient is a map that moves an extension by zero past a proper direct image. It is an isomorphism by the same argument as for refinements: it restricts to the identity on an open, and both sides vanish on the closed complement by proper base change. For base change of \(Rf_!\) we combine proper base change with the elementary fact that extension by zero commutes with pullback.

\emph{The composition isomorphism.} Choose compactifications \(f=p\circ j\) and \(g=q\circ j_Y\), with \(j:X\to\overline X\) and \(j_Y:Y\to\overline Y\) open and \(p,q\) proper. In the composite \(Rq_*j_{Y!}Rp_*j_!\), the middle factor \(j_{Y!}Rp_*\) applies a proper direct image before an extension by zero, and we want to interchange the two. To do this, compactify the morphism \(j_Y\circ p:\overline X\to\overline Y\) as \(\overline X\xrightarrow{j'}W\xrightarrow{r}\overline Y\), with \(j'\) open, \(r\) proper and \(\overline X\) dense in \(W\). Then \(r^{-1}(Y)=\overline X\), because the open immersion \(\overline X\to r^{-1}(Y)\) is a morphism from a proper \(Y\)-scheme to a separated \(Y\)-scheme, hence proper, and its image is closed, open and dense.

The interchange map is defined as follows. For \(M\) in \(D(\overline X_{\mathrm{pro}\acute{e}t},\Lambda)\), restriction to the open \(Y\subset\overline Y\) identifies \(j_Y^*Rr_*j'_!M\) with \(Rp_*M\), because \(r^{-1}(Y)=\overline X\), the restriction of \(r\) over \(Y\) is \(p\), and \(j'_!M\) restricts to \(M\) on \(\overline X\). By the adjunction \(j_{Y!}\dashv j_Y^*\), we obtain a map

\[
\iota_M:j_{Y!}Rp_*M\longrightarrow Rr_*j'_!M,
\tag{5.9}
\]

which is the unique map restricting to the identity of \(Rp_*M\) over \(Y\).

For \(M\in D^b_{\mathrm{ctf}}(\overline X,\Lambda)\), this map is an isomorphism. Indeed, over \(Y\) it is the identity, and over \(\overline Y\setminus Y\) both sides vanish. For the right-hand side, the vanishing follows from proper base change on \(D^b_{\mathrm{ctf}}\) (Proposition \hyperref[bp488]{5.24}) for the inclusion of \(\overline Y\setminus Y\) into \(\overline Y\), since \(j'_!M\) vanishes on \(r^{-1}(\overline Y\setminus Y)=W\setminus\overline X\). Both sides preserve direct sums on \(\Loc(D^b_{\mathrm{ctf}}(\overline X,\Lambda))\) (Proposition \hyperref[bp482]{5.22}). Hence Lemma \hyperref[bp485]{5.23} shows that \(\iota_M\) is an isomorphism for all \(M\in\Loc(D^b_{\mathrm{ctf}}(\overline X,\Lambda))\). For \(M=j_!K\) it gives

\[
Rq_*j_{Y!}Rp_*j_!K\simeq Rq_*Rr_*j'_!j_!K=R(qr)_*(j'j)_!K,
\]

and the right side is \(R(gf)_!K\) computed with the compactification \(gf=(qr)\circ(j'j)\).

\emph{Verification of canonicity and associativity.} Every map used is characterized by its restriction to an open subscheme, since the refinement maps (5.7) and the maps \(\iota_M\) restrict to identities. For example, replacing any of the chosen compactifications by a refinement commutes with \(\iota_M\), because both composites restrict to the same identity and are then equal by adjunction. For associativity, given three composable morphisms, we choose common refinements (by graph closure) of the finitely many compactifications that occur in the composites of two of them. On these, each of the two bracketings is built from composites of proper direct images, composites of extensions by zero, and maps of the form \(\iota_M\). The first two kinds of pieces are associative: the isomorphisms \(Rq_*Rr_*\simeq R(qr)_*\) are associative because they are the isomorphisms of derived functors of composites of left exact functors (equivalently, they are adjoint to the composition isomorphisms of the inverse images), and composition of extensions by zero is associative by adjunction. Moreover, interchanging two adjacent open and proper steps in either order gives the same map by the uniqueness of \(\iota_M\). So the two bracketings agree. Similarly, the unit conditions are the special case where one compactification is the identity. Finally, transitivity of the refinement maps carries these equalities over to the originally chosen compactifications.

\emph{Base change.} The base change map for \(Rf_!=Rp_*j_!\) is built from one for each factor, so we treat \(j_!\) and \(Rp_*\) separately. Since open immersions and proper morphisms are stable under base change, we may pull back a compactification \(f=p\circ j\) along \(v\) to a compactification \(f'=p'\circ j'\), where now \(j'\) and \(p'\) denote the base changes of \(j\) and \(p\). For \(j_!\), extension by zero commutes with inverse image, as shown before the proposition. Combining this with proper base change on \(\Loc(D^b_{\mathrm{ctf}}(\overline X,\Lambda))\) (Proposition \hyperref[bp488]{5.24}), we obtain the isomorphism \(v^*Rp_*j_!\simeq Rp'_*v'^*j_!\simeq Rp'_*j'_!u^*\) on \(\Loc(D^b_{\mathrm{ctf}}(X,\Lambda))\). This isomorphism is compatible with change of compactification because the refinement maps are characterized by their restriction to the open subscheme. It is also compatible with composing two cartesian squares, because both factors are. For the first factor, a map between extensions by zero from an open is determined by its restriction to that open, since \(\Hom(j_!B,j_!B')=\Hom(B,B')\). For the second, a base change map is the adjoint transpose of a counit, and the counit of a composite adjunction is the composite of the counits. Finally, compatibility with the composition isomorphisms of (1) follows in the same way as canonicity: both composites are built from refinement maps, maps \(\iota_M\) and proper base change maps, and they agree because their restrictions to the open part are the same identity.
\end{proof}

% text: prose-05-32
\phantomsection\label{bp498}%
On a space, a section with compact support is in particular a section, and so compactly supported cohomology maps to ordinary cohomology. The algebraic counterpart is a map \(Rf_!\to Rf_*\), which is an isomorphism for proper \(f\), and we need it to be compatible with the composition maps just constructed: in Section \hyperref[bp515]{6.1} this is one of the compatibilities that the comparison with the étale site has to respect.

% object: proposition-forgetting-supports-composition (blueprint/chap-Construction-of-f/sec-Direct-image-with-compact-support/subsec-Construction-of-Rf/proposition-forgetting-supports-composition/_section.md)
\phantomsection\label{bp499}%
\begin{bpstatementit}{Proposition~5.29}{ (Forgetting supports and composition)}
There is a natural map \(\eta_f:Rf_!\to Rf_*\) on \(\Loc(D^b_{\mathrm{ctf}}(X,\Lambda))\). For a compactification \(f=pj\) it is

\[
Rp_*j_!\longrightarrow Rp_*Rj_*=Rf_*,
\]

where \(j_!\to Rj_*\) is adjoint to the inverse of the counit \(j^*Rj_*\xrightarrow{\sim}\mathrm{id}\). It is independent of the compactification, and the square

\[
\begin{CD}
Rg_!Rf_! @>{c^!_{g,f}}>> R(gf)_!\\
@V{Rg_*(\eta_f)\,\circ\,(\eta_g Rf_!)}VV @VV{\eta_{gf}}V\\
Rg_*Rf_* @>{c^*_{g,f}}>> R(gf)_*
\end{CD}
\]

commutes. Here \(c^!_{g,f}\) is the composition map of Proposition \hyperref[bp496]{5.28}, and \(c^*_{g,f}\) is the composition map for ordinary direct images. If \(f\) is proper, \(\eta_f\) is an isomorphism; under these isomorphisms the two composition maps agree whenever \(f\) and \(g\) are proper.
\end{bpstatementit}

% proof: proof/proposition-forgetting-supports-composition
\begin{proof}[{Proof}]\phantomsection\label{bp500}%
All maps involved are maps out of an extension by zero, so each of them is determined by its restriction to the open subscheme, where it is an identity or an inverse counit. We use this repeatedly.

Independence of the compactification reduces to a refinement \(h:Z'\to Z\). The two maps \(j_!K\to Rh_*Rj'_*K=Rj_*K\) obtained from (5.7) and from \(j_!K\to Rj_*K\) have the same adjoint under \(j_!\dashv j^*\), namely the inverse of \(j^*Rj_*K\to K\). Hence they are equal. Applying \(Rp_*\) gives independence.

The square is checked on the compactifications that define \(c^!_{g,f}\) and the interchange map (5.9). Apart from composition and naturality squares for ordinary direct images, it comes down to comparing the two maps

\[
j_{Y!}Rp_*M\longrightarrow Rr_*j'_!M
\longrightarrow Rr_*Rj'_*M=Rj_{Y*}Rp_*M,
\qquad
j_{Y!}Rp_*M\longrightarrow Rj_{Y*}Rp_*M.
\]

Their adjoints under \(j_{Y!}\dashv j_Y^*\) are both the inverse counit for \(j_Y^*\dashv Rj_{Y*}\), because the interchange map restricts to the identity on \(Y\). So the two maps agree. Applying \(Rq_*\) with \(M=j_!K\), and using naturality for \(j_!K\to Rj_*K\), we get the displayed square for the chosen compactifications. The refinement calculation extends it to any choices. If \(f\) is proper, the compactification with \(j=\mathrm{id}\) has \(\eta_f\) equal to the identity, and the last assertion follows from the square.
\end{proof}



% object: sec-Brown-representability (blueprint/chap-Construction-of-f/sec-Brown-representability/_section.md), content hidden
\subsection*{\texorpdfstring{5.5\quad Brown representability}{5.5 Brown representability}}
\phantomsection\label{bp501}%
\addcontentsline{toc}{subsection}{\texorpdfstring{5.5\quad Brown representability}{5.5 Brown representability}}
% text: prose-05-33
\phantomsection\label{bp502}%
Let \(\Lambda\) be an admissible coefficient ring and fix \(N\) in \(\Loc(D^b_{\mathrm{ctf}}(Y,\Lambda))\). Since the functor \(Rf_!\) preserves direct sums (Proposition \hyperref[bp492]{5.26}), the functor \(M\mapsto\Hom(Rf_!M,N)\) turns direct sums into products. If the objects of \(D^b_{\mathrm{ctf}}(X,\Lambda)\) are compact generators of the localizing category, Brown representability produces an object \(f^!N\) representing it. The next lemma checks this hypothesis for every admissible \(\Lambda\); for \(\mathcal O_E\) it uses the compactness proved in Subsection \hyperref[bp469]{5.4.2} by reduction to torsion coefficients.

Compactness enters a second time. The functor \(Rf_!\) also maps \(D^b_{\mathrm{ctf}}(X,\Lambda)\) into \(D^b_{\mathrm{ctf}}(Y,\Lambda)\), so it sends the compact generators to compact objects. For this reason the adjoint \(f^!\) preserves direct sums (Theorem \hyperref[bp506]{5.31}).

% object: lemma-hypotheses-brown-representability (blueprint/chap-Construction-of-f/sec-Brown-representability/lemma-hypotheses-brown-representability/_section.md)
\phantomsection\label{bp503}%
\begin{bpstatementit}{Lemma~5.30}{ (Hypotheses of Brown representability)}
For every admissible \(\Lambda\), the category \(\Loc(D^b_{\mathrm{ctf}}(X,\Lambda))\) is locally small and triangulated, has arbitrary direct sums, and has split idempotents. The thick category \(D^b_{\mathrm{ctf}}(X,\Lambda)\) is essentially small, and a small skeleton is a set of compact generators of \(\Loc(D^b_{\mathrm{ctf}}(X,\Lambda))\).
\end{bpstatementit}

% proof: proof/lemma-hypotheses-brown-representability
\begin{proof}[{Proof}]\phantomsection\label{bp504}%
The set-theoretic conditions are general facts about derived categories of Grothendieck abelian categories. Compactness, on the other hand, was proved separately for each kind of coefficients, and generation follows from the definition of the localizing category.

As the derived category of a Grothendieck abelian category of module sheaves, the ambient category is locally small, and so is its full subcategory \(\Loc(D^b_{\mathrm{ctf}}(-,\Lambda))\). Moreover, the latter has direct sums by definition, so its idempotents split [Stacks, \href{https://stacks.math.columbia.edu/tag/05QW}{05QW}]. Finally, the subcategory \(D^b_{\mathrm{ctf}}(X,\Lambda)\) is thick, since constructibility and finite Tor amplitude are preserved by shifts, cones and direct summands.

To see that \(D^b_{\mathrm{ctf}}(X,\Lambda)\) is essentially small, note first that the constructible sheaves in the heart have a set of representatives for each admissible \(\Lambda\). A bounded complex with such cohomology is built from its cohomology sheaves by finitely many extensions, formed with truncations in the ambient category. Each extension step is classified by an element of a Hom set of the ambient category, and these Hom sets are sets because the ambient category is locally small. Hence these complexes form an essentially small category, of which \(D^b_{\mathrm{ctf}}(X,\Lambda)\) is a full subcategory.

Compactness for finite \(A\) is Proposition \hyperref[bp451]{5.14}. For \(E\) and \(\overline{\mathbf Q}_\ell\) it is Theorem \hyperref[bp446]{5.13}, and for \(\mathcal O_E\) it is Proposition \hyperref[bp477]{5.21}; in the integral and rational cases, \(D^b_{\mathrm{ctf}}=\Dbc\).

For generation, suppose that \(K\in\Loc(D^b_{\mathrm{ctf}}(X,\Lambda))\) receives no nonzero maps from shifts of objects of a small skeleton. The objects \(M\) with \(\Hom(M[n],K)=0\) for all \(n\) form a localizing subcategory containing \(D^b_{\mathrm{ctf}}\). Hence this subcategory contains \(K\). Then \(\mathrm{id}_K=0\), and so \(K=0\).
\end{proof}

% text: prose-05-34
\phantomsection\label{bp505}%
With the hypotheses checked, the theorem below follows from Brown representability in the form recalled at the start of the chapter.

% object: theorem-existence-exceptional-inverse-image (blueprint/chap-Construction-of-f/sec-Brown-representability/theorem-existence-exceptional-inverse-image/_section.md)
\phantomsection\label{bp506}%
\begin{bpstatementit}{Theorem~5.31}{ (Existence of the exceptional inverse image)}
Let \(\Lambda\in\{A,\mathbf Z_\ell,\mathcal O_E,\mathbf Q_\ell,E,\overline{\mathbf Q}_\ell\}\) be any admissible coefficient ring. For every separated morphism of finite type \(f:X\to Y\), the coproduct-preserving functor \(Rf_!:\Loc(D^b_{\mathrm{ctf}}(X,\Lambda))\to\Loc(D^b_{\mathrm{ctf}}(Y,\Lambda))\) has an exact right adjoint

\[
f^!:\Loc(D^b_{\mathrm{ctf}}(Y,\Lambda))\longrightarrow\Loc(D^b_{\mathrm{ctf}}(X,\Lambda)),
\qquad
\Hom(Rf_!M,N)\simeq\Hom(M,f^!N),
\]

for \(M\in\Loc(D^b_{\mathrm{ctf}}(X,\Lambda))\) and \(N\in\Loc(D^b_{\mathrm{ctf}}(Y,\Lambda))\). It is unique up to unique isomorphism compatible with the adjunction, and it preserves arbitrary direct sums.
\end{bpstatementit}

% proof: proof/theorem-existence-exceptional-inverse-image
\begin{proof}[{Proof}]\phantomsection\label{bp507}%
We construct \(f^!N\) one object at a time as the object representing maps out of \(Rf_!(-)\) into \(N\), and then check direct sums on the compact generators.

For \(N\in\Loc(D^b_{\mathrm{ctf}}(Y,\Lambda))\), the functor

\[
\Loc(D^b_{\mathrm{ctf}}(X,\Lambda))^{\mathrm{op}}\longrightarrow\operatorname{Ab},
\qquad M\longmapsto\Hom(Rf_!M,N)
\]

is cohomological and sends direct sums to products. By Lemma \hyperref[bp503]{5.30}, \(\Loc(D^b_{\mathrm{ctf}}(X,\Lambda))\) is compactly generated, and Brown representability [Stacks, \href{https://stacks.math.columbia.edu/tag/0A8F}{0A8F}] says that on such a category every cohomological functor sending direct sums to products is representable, say by \(f^!N\). By the Yoneda lemma, these objects form a right adjoint, and a right adjoint of an exact functor of triangulated categories is exact [Stacks, \href{https://stacks.math.columbia.edu/tag/0A8G}{0A8G}].

To see that \(f^!\) preserves direct sums, we test the canonical map \(\bigoplus_\alpha f^!N_\alpha\to f^!(\bigoplus_\alpha N_\alpha)\) on a compact generator \(G\in D^b_{\mathrm{ctf}}(X,\Lambda)\) and its shifts. On the source, compactness of \(G\) and adjunction give

\[
\Hom\Bigl(G[n],\bigoplus_\alpha f^!N_\alpha\Bigr)
=\bigoplus_\alpha\Hom(G[n],f^!N_\alpha)
=\bigoplus_\alpha\Hom(Rf_!G[n],N_\alpha).
\]

On the target, adjunction and compactness of \(Rf_!G\in D^b_{\mathrm{ctf}}(Y,\Lambda)\) give the same group. The map therefore induces bijections on \(\Hom(G[n],-)\) for all compact generators \(G\) and all \(n\). Applying the detection property to its cone, we see that it is an isomorphism.
\end{proof}


% object: sec-Composition-and-restriction (blueprint/chap-Construction-of-f/sec-Composition-and-restriction/_section.md), content hidden
\subsection*{\texorpdfstring{5.6\quad Composition and restriction}{5.6 Composition and restriction}}
\phantomsection\label{bp508}%
\addcontentsline{toc}{subsection}{\texorpdfstring{5.6\quad Composition and restriction}{5.6 Composition and restriction}}
% text: prose-05-35
\phantomsection\label{bp509}%
Brown representability gives \(f^!\) one morphism at a time. To compute it in Chapter \hyperref[bp513]{6}, however, we factor \(f\) locally. The two are linked by uniqueness of adjoints: an isomorphism between composites of left adjoints, such as \(Rg_!\circ Rf_!\simeq R(gf)_!\) or \(Rf_!=Rp_*\circ j_!\), gives an isomorphism between the right adjoints, composed in the opposite order. Since \(j^!=j^*\) for an open immersion \(j\), it follows that \(f^!\) is local on \(X\) and on \(Y\). For example, a smooth \(f\) is locally the composite of an étale map and a projection \(\mathbf A^n_Y\to Y\). By composition and locality, \(f^!\) is then determined by its values on these two kinds of maps.

Base change for \(f^!\) does not follow in this way. Passing to right adjoints in the base change isomorphism \(v^*Rf_!\simeq Rf'_!u^*\) relates \(f^!Rv_*\) and \(Ru_*f'^!\). These involve direct images, not the inverse images \(u^*f^!\) and \(f'^!v^*\) that one wants to compare. We prove base change for smooth morphisms in Proposition \hyperref[bp542]{6.9}. That \(f^!\) preserves \(D^b_{\mathrm{ctf}}\) is proved in Theorems \hyperref[bp573]{6.20} and \hyperref[bp596]{Theorem 6.27}.

% object: proposition-composition-open-restriction-compactifications-f (blueprint/chap-Construction-of-f/sec-Composition-and-restriction/proposition-composition-open-restriction-compactifications-f/_section.md)
\phantomsection\label{bp510}%
\begin{bpstatementit}{Proposition~5.32}{ (Composition, open restriction and compactifications for $f^!$)}
For every admissible \(\Lambda\), the adjoints \(f^!\) on \(\Loc(D^b_{\mathrm{ctf}}(-,\Lambda))\) have the following canonical compatibilities.

\begin{enumerate}
\item For composable \(X\xrightarrow fY\xrightarrow gZ\) there are canonical isomorphisms \((gf)^!\simeq f^!g^!\), associative and compatible with the identity.
\item For an open immersion \(j:U\to X\), \(j^!\simeq j^*\).
\item Let \(v:V\to Y\) be an open immersion, \(X_V=f^{-1}(V)\) with inclusion \(u:X_V\to X\), and \(f_V:X_V\to V\) the restriction of \(f\). Then \(u^*f^!\simeq f_V^!v^*\).
\item If \(f=p\circ j\) is a compactification, then \(f^!\simeq j^*\circ p^!\) on \(\Loc(D^b_{\mathrm{ctf}}(Y,\Lambda))\).
\end{enumerate}

In particular, \(f^!\) can be computed locally on \(X\) and on \(Y\) for open coverings.
\end{bpstatementit}

% proof: proof/proposition-composition-open-restriction-compactifications-f
\begin{proof}[{Proof}]\phantomsection\label{bp511}%
Each statement is the right adjoint of a statement about lower shriek functors already proved, and we conclude by uniqueness of adjoints.

(1) For \(M\in\Loc(D^b_{\mathrm{ctf}}(X,\Lambda))\) and \(N\in\Loc(D^b_{\mathrm{ctf}}(Z,\Lambda))\),

\[
\Hom(M,f^!g^!N)
\simeq\Hom(Rf_!M,g^!N)
\simeq\Hom(Rg_!Rf_!M,N)
\simeq\Hom(R(gf)_!M,N),
\]

where the last step holds by Proposition \hyperref[bp496]{5.28}. By the Yoneda lemma, \(f^!g^!\simeq(gf)^!\). Passing to right adjoints transforms the associativity and identity constraints of Proposition proposition-composition-base-change-rf-shriek into those for \(f^!\).

(2) For an open immersion, \(Rj_!=j_!\) (take the trivial compactification \(p=\mathrm{id}\)), and \(j_!\dashv j^*\) on the localizing categories; here we use that \(j^*\) preserves the localizing categories, as inverse images do (see the proof of Proposition \hyperref[bp488]{5.24}). By uniqueness of right adjoints, \(j^!\simeq j^*\).

(3) Since \(f\circ u=v\circ f_V\), parts (1) and (2) show that \(u^*f^!\simeq(fu)^!=(vf_V)^!\simeq f_V^!v^*\). For an open \(j:U\to X\), parts (1) and (2) also give \(j^*f^!\simeq(fj)^!\). These two isomorphisms show that \(f^!\) is local on \(X\) and on \(Y\).

(4) On \(\Loc(D^b_{\mathrm{ctf}}(-,\Lambda))\), \(j^*\) is right adjoint to \(j_!\) and \(p^!\) to \(Rp_*=Rp_!\), so \(j^*\circ p^!\) and \(f^!\) are both right adjoint to \(Rp_*\circ j_!=Rf_!\), and they are isomorphic by uniqueness of right adjoints.
\end{proof}

% text: prose-05-36
\phantomsection\label{bp512}%
Once we know that \(f^!\) preserves \(D^b_{\mathrm{ctf}}(-,\Lambda)\) (Chapter \hyperref[bp513]{6}), these isomorphisms restrict to the full subcategories \(D^b_{\mathrm{ctf}}(-,\Lambda)\), and the same statements hold for the adjunction between \(D^b_{\mathrm{ctf}}(X,\Lambda)\) and \(D^b_{\mathrm{ctf}}(Y,\Lambda)\).



% object: chap-Constructibility-of-f (blueprint/chap-Constructibility-of-f/_section.md), content hidden
\section*{\texorpdfstring{6\quad Constructibility of $f^!$}{6 Constructibility of f!}}
\phantomsection\label{bp513}%
\addcontentsline{toc}{section}{\texorpdfstring{6\quad Constructibility of $f^!$}{6 Constructibility of f!}}
% text: prose-06-01
\phantomsection\label{bp514}%
Brown representability gave us \(f^!\) (Theorem \hyperref[bp506]{5.31}), but only as a functor

\[
f^!:\Loc(D^b_{\mathrm{ctf}}(Y,\Lambda))\longrightarrow\Loc(D^b_{\mathrm{ctf}}(X,\Lambda))
\]

for \(f:X\to Y\) separated of finite type and \(\Lambda\) an admissible coefficient ring. It tells us nothing about the values of \(f^!\). A priori, \(f^!\) of a constructible complex is just some object of the large category on the right, possibly unbounded and with cohomology sheaves that are not constructible.

The rest of the notes needs more than this. For example, Verdier duality (Chapter \bpout{10}) says that \(\mathbb D_X=\cHom(-,a_X^!\Lambda)\) is an involution of \(D^b_{\mathrm{ctf}}(X,\Lambda)\), and this makes sense only if the dualizing complex \(a_X^!\Lambda\) lies in \(D^b_{\mathrm{ctf}}(X,\Lambda)\). Similarly, the local terms of the trace formula are defined through the isomorphism \(\cHom(c_2^*K,c_1^!K)\simeq c^!(K\boxtimes\mathbb D_XK)\) for a correspondence \(c=(c_1,c_2)\) (Lemma \bpout{11.7}), which also needs \(f^!\) to preserve constructible complexes. In this chapter we prove what both need (part (4) of Theorem \hyperref[bp413]{5.2}): \(f^!\) maps \(D^b_{\mathrm{ctf}}(Y,\Lambda)\) into \(D^b_{\mathrm{ctf}}(X,\Lambda)\). A universal property cannot give this, so we have to compute \(f^!\), and the computation is local.

By Proposition \hyperref[bp510]{5.32}, the functor \(f^!\) is compatible with composition, equals \(j^*\) for an open immersion \(j\), and is local on \(X\) and on \(Y\). We may therefore assume that \(Y\) is affine and that \(f\) is a closed immersion into some \(\mathbf A^n_Y\) followed by the projection to \(Y\). Since \(\mathbf A^n_Y\) is open in \((\mathbf P^1_Y)^n\), and \((\mathbf P^1_Y)^n\to Y\) is a composite of \(n\) projections \(\mathbf P^1_T\to T\), two cases remain.

For the projection \(p:\mathbf P^1_T\to T\) we can write the adjunction down, and we find \(p^!=p^*(1)[2]\) (Proposition \hyperref[bp534]{6.7}). Here the counit is the trace \(Rp_*\Lambda(1)[2]\to\Lambda\), and the unit is the cycle class of the diagonal in \(\mathbf P^1_T\times_T\mathbf P^1_T\). The triangle identities then hold because the diagonal meets each fibre of either projection in a single point. For a closed immersion \(i:Z\hookrightarrow X\) with open complement \(j:U\hookrightarrow X\), on the other hand, the functor \(i^!\) takes sections with support in \(Z\), and it sits in the localization triangle

\[
i_*i^!K\longrightarrow K\longrightarrow Rj_*j^*K\longrightarrow .
\]

Let \(K\) be constructible. Then \(j^*K\) is constructible, and since \(i^*i_*\simeq\mathrm{id}\) and \(i_*\) preserves constructibility, \(i^!K\) is constructible exactly when the third term \(Rj_*j^*K\) is. So the constructibility of \(f^!\) comes down to one question: does direct image along an open immersion preserve bounded constructible complexes and, with finite coefficients, finite Tor amplitude? Proper base change does not help here, since \(j\) is not proper.

The answer is not formal, because \(Rj_*\), and with it \(i^!\), sees more than the stalks of \(K\) on \(Z\): it sees the cohomology of \(K\) on a punctured neighbourhood of \(Z\). For the origin \(i:0\hookrightarrow\mathbf A^1\) over an algebraically closed field, the stalk of \(Rj_*\Lambda\) at \(0\) is the cohomology of a punctured disc, which the Kummer covers of \(\mathbf G_m\) compute to be \(\Lambda\oplus\Lambda(-1)[-1]\), and the triangle gives \(i^!\Lambda\simeq\Lambda(-1)[-2]\), whereas \(i^*\Lambda=\Lambda\) (Example \hyperref[bp591]{6.25} below).

Over a field the answer is Deligne's finiteness theorem in SGA 4\(\tfrac12\) [\hyperref[bp703]{SGA4½}], and we follow his proof. Its induction on dimension leads to fibrations over bases of positive dimension, so Deligne proves a statement in families over a base \(S\). In families, however, the formation of \(Rj_*\) need not commute with base change. Let \(S=\mathbf A^1\) with coordinate \(t\), let \(Y=\mathbf A^2\) with coordinates \((t,x)\), viewed over \(S\), and let \(j\) be the inclusion of \(X=Y\setminus\{(0,0)\}\). By purity (Theorem \bpout{9.3}, proved in Chapter \bpout{9} below), the punctured strict henselization of \(Y\) at the origin has cohomology \(\Lambda\) in degree \(0\) and \(\Lambda(-2)\) in degree \(3\) only. Hence \(R^1j_*\Lambda=0\). On the fibre over \(t=0\), however, \(j\) becomes \(\mathbf G_m\hookrightarrow\mathbf A^1\), whose \(R^1j_*\Lambda\) is \(\Lambda(-1)\) at the origin. Since over \(t\neq0\) the map \(j\) is an isomorphism, base change fails only at \(t=0\).

Deligne's generic base change theorem says that this is the general picture. Let \(j:X\hookrightarrow Y\) be an open immersion of schemes of finite type over an integral base \(S\) of finite type over \(k\), and let \(F\) be a constructible \(\Fell\)-sheaf on \(X\). Then there is a nonempty open \(S_0\subset S\) over which \(Rj_*F\) is bounded constructible and its formation commutes with base change along every \(T\to S_0\) of finite type (Theorem \hyperref[bp563]{6.17}; Deligne proves more). Over \(S=\operatorname{Spec}k\) the only nonempty open is \(S\) itself, so we obtain the finiteness and base change of \(Rj_*\) over a field (Corollary \hyperref[bp566]{6.18}).

The projective line returns in the proof of the generic theorem. There we must control \(Rf_*L\) for smooth \(f\) and locally constant \(L\), and we do so through the formula for \(f^!\) of a smooth morphism (Section \hyperref[bp536]{6.3}).

For finite coefficients (Theorem \hyperref[bp573]{6.20}) we also bound the Tor amplitude of \(Rj_*K\): it grows by at most \(\operatorname{cd}_\ell(j)\) (Proposition \hyperref[bp570]{6.19}). For integral and rational coefficients we have \(D^b_{\mathrm{ctf}}=\Dbc\), and we reduce modulo a uniformizer. This works because an integral complex is constructible if it is derived complete with constructible reduction (Proposition \hyperref[bp226]{3.11}). Only the adjunction for the projective line resists this reduction. Its unit and counit must be constructed as single morphisms with \(\mathcal O_E\)-coefficients (Proposition \hyperref[bp581]{6.22}). The result is Theorem \hyperref[bp596]{6.27}.

% object: sec-Conventions-2 (blueprint/chap-Constructibility-of-f/sec-Conventions-2/_section.md), content hidden
\subsection*{\texorpdfstring{6.1\quad Conventions}{6.1 Conventions}}
\phantomsection\label{bp515}%
\addcontentsline{toc}{subsection}{\texorpdfstring{6.1\quad Conventions}{6.1 Conventions}}
% text: prose-06-02
\phantomsection\label{bp516}%
Until Section \hyperref[bp568]{6.6} the coefficient ring is a finite commutative \(\ell\)-primary ring, which we denote by \(A\). With such coefficients nothing is lost by computing on the étale site. For \(X\) of finite type over \(k\), the comparison functor

\[
\nu_X^*:D(X_{\acute{e}t},A)\longrightarrow D(X_{\mathrm{pro}\acute{e}t},A)
\]

is fully faithful (Proposition \hyperref[bp146]{2.49}), and it identifies \(D(X_{\acute{e}t},A)\) with the localizing category \(\Loc(D^b_{\mathrm{ctf}}(X,A))\) on which \(f^!\) was constructed (Proposition \hyperref[bp454]{5.15}). It also identifies the constructible complexes on the two sites (Theorem \hyperref[bp129]{2.43}). In the following sections we therefore work in \(D(X_{\acute{e}t},A)\), and every result is transported to the pro-étale site by \(\nu_X^*\).

This transport is harmless because \(\nu^*\) identifies all the operations we use with their étale counterparts, compatibly with units and counits (Lemma \hyperref[bp116]{2.39}). These operations are the cohomology (Theorem \hyperref[bp103]{2.35}); the direct images and their base change maps, also for unbounded complexes (Theorem \hyperref[bp113]{2.38} and Proposition \hyperref[bp146]{2.49}); and the inverse images, tensor products, extension by zero and the localization triangles (Section \hyperref[bp111]{2.9}). Moreover, by Corollary \hyperref[bp132]{2.44}, finite Tor amplitude can be checked on the étale site.

We write \(\Dbc(X,A)\) also for the bounded constructible objects of \(D(X_{\acute{e}t},A)\), which correspond to those of the pro-étale category. For a separated morphism of finite type \(f:X\to S\), the functor \(Rf_!\) and its right adjoint \(f^!\) of Theorem \hyperref[bp506]{5.31} correspond to the functors \(Rf_!\) and \(f^!\) between \(D(X_{\acute{e}t},A)\) and \(D(S_{\acute{e}t},A)\), and the two adjunctions correspond. Accordingly, from now on \(g^*\), \(Rf_*\), \(Rf_!\), \(f^!\) and \(\otimes^{\mathbf L}\) denote the functors on the étale derived categories.

Below we apply Proposition \hyperref[bp496]{5.28}, stated on \(\Loc(D^b_{\mathrm{ctf}}(-,A))\), to the étale functors. This is legitimate because the comparison identifies the composition maps of Proposition \hyperref[bp496]{5.28} on the two sites: the refinement maps (5.7) and the open–proper interchange (5.9) are characterized by their restriction to the open. The comparison also preserves the counits defining base change and the inverse counits defining forgetting supports (Proposition \hyperref[bp499]{5.29}). Taking right adjoints, we see that it identifies the composition maps for \(f^!\) as well.

Except in Lemmas \hyperref[bp523]{6.3}, \hyperref[bp526]{Lemma 6.4} and \hyperref[bp553]{Lemma 6.13}, which allow more general bases as stated there, all schemes and base changes are of finite type over \(k\).

We use three classical theorems about the étale site.

% object: theorem-results-from-etale-theory (blueprint/chap-Constructibility-of-f/sec-Conventions-2/theorem-results-from-etale-theory/_section.md)
\phantomsection\label{bp517}%
\begin{bpstatementit}{Theorem~6.1}{ (Results from the étale theory)}
Let \(f:X\to S\) be a morphism of schemes of finite type over \(k\).

\begin{enumerate}
\item (Proper base change) If \(f\) is proper, then for every morphism of schemes \(g:S'\to S\), not necessarily of finite type over \(k\), and every \(K\in D(X_{\acute{e}t},A)\), the base change map \(g^*Rf_*K\to Rf'_*g'^*K\) is an isomorphism, where \(f':X'\to S'\) and \(g':X'\to X\) are the projections of \(X'=X\times_SS'\). In particular, the geometric stalks of \(Rf_*K\) are the cohomology of the geometric fibres.
\item (Finiteness) If \(f\) is separated, then \(Rf_!\) maps \(\Dbc(X,A)\) into \(\Dbc(S,A)\).
\item (Cohomological dimension) The functors \(Rf_*\) and, if \(f\) is separated, \(Rf_!\) have finite cohomological amplitude on torsion sheaves. In particular \(\operatorname{cd}_\ell(f)<\infty\) in the sense of Definition \hyperref[bp137]{2.45}.
\end{enumerate}
\end{bpstatementit}

% proof: proof/theorem-results-from-etale-theory
\begin{proof}[{Proof}]\phantomsection\label{bp518}%
For bounded below complexes, (1) is Theorem \hyperref[bp382]{4.33} together with the comparison of direct images (Theorem \hyperref[bp113]{2.38}); for unbounded complexes it is [Stacks, \href{https://stacks.math.columbia.edu/tag/0F0C}{0F0C}], whose finiteness hypothesis on \(f\) is part (3). Part (2) is [Stacks, \href{https://stacks.math.columbia.edu/tag/0GL0}{0GL0}]. As for part (3), for \(Rf_*\) it is [Stacks, \href{https://stacks.math.columbia.edu/tag/0F10}{0F10}], and for \(Rf_!\) we write \(f=\bar f\circ j\) with \(\bar f\) proper and \(j\) an open immersion; then \(Rf_!=R\bar f_*j_!\) with \(j_!\) exact, and [Stacks, \href{https://stacks.math.columbia.edu/tag/0F10}{0F10}] applies to \(\bar f\). Finally, the last assertion follows by the comparison of direct images, since classical sheaves are of the form \(\nu_X^*F\).
\end{proof}

% object: definition-exceptional-base-change-finite (blueprint/chap-Constructibility-of-f/sec-Conventions-2/definition-exceptional-base-change-finite/_section.md)
\phantomsection\label{bp519}%
\begin{bpstatement}{Definition~6.2}{ (The base change map for $f^!$, finite coefficients)}
Let \(f:X\to S\) be separated and of finite type, let \(g:S'\to S\) be a morphism, let \(X'=X\times_SS'\), and let \(f':X'\to S'\), \(g':X'\to X\) be the projections. For \(N\in D(S_{\acute{e}t},A)\), let \(\varepsilon_N:Rf_!f^!N\to N\) be the counit of \(Rf_!\dashv f^!\). The \emph{base change map}

\[
\beta_N:g'^*f^!N\longrightarrow f'^!g^*N
\tag{6.1}
\]

is the transpose, for \(f'\), of the morphism

\[
Rf'_!g'^*f^!N
\xrightarrow{\ \sim\ }g^*Rf_!f^!N
\xrightarrow{g^*\varepsilon_N}g^*N,
\tag{6.2}
\]

where the first arrow is the inverse of the base change isomorphism for \(Rf_!\) (Proposition proposition-composition-base-change-rf-shriek).
\end{bpstatement}

% text: prose-06-03
\phantomsection\label{bp520}%
It is an isomorphism for smooth \(f\) (Proposition \hyperref[bp542]{6.9}). The global isomorphism \(f^!\simeq f^*(d)[2d]\) for smooth \(f\) of relative dimension \(d\), and its compatibility with composition, are treated with smooth purity in Chapter \bpout{9}. Here we need only its restriction to charts \(U\to\mathbf A^d_S\).


% object: sec-Duality-for-the-projective-line (blueprint/chap-Constructibility-of-f/sec-Duality-for-the-projective-line/_section.md), content hidden
\subsection*{\texorpdfstring{6.2\quad Duality for the projective line}{6.2 Duality for the projective line}}
\phantomsection\label{bp521}%
\addcontentsline{toc}{subsection}{\texorpdfstring{6.2\quad Duality for the projective line}{6.2 Duality for the projective line}}
% text: prose-06-04
\phantomsection\label{bp522}%
Let \(p:\mathbf P^1_S\to S\) be the projection. We construct the adjunction \(Rp_*\dashv p^*(1)[2]\) from two maps: the counit is a trace \(Rp_*A(1)[2]\to A\), obtained from the cohomology of \(\mathbf P^1\), and the unit is the cycle class of the diagonal in \(\mathbf P^1_S\times_S\mathbf P^1_S\), a special case of the cycle class of an effective Cartier divisor. Both triangle identities say that pushing the class of the diagonal forward along either projection \(\mathbf P^1\times\mathbf P^1\to\mathbf P^1\) and applying the trace gives \(1\). This holds because the diagonal meets each fibre of either projection in a single point.

In topology the same two maps appear. For instance, for a bundle \(E\to B\) with fibre \(S^2\), integration over the fibre \(H^2(S^2)\to\mathbf Z\) is the counit of Poincaré duality on the fibres, and the unit is the Thom class of the diagonal of \(E\times_BE\), which is Poincaré dual to the identity correspondence. In étale cohomology, by contrast, the fundamental class of a point on \(\mathbf P^1\) comes from a line bundle through the Kummer sequence. So both maps are built from line bundles: the trace comes from \(c_1(\mathcal O(1))\), and the class of the diagonal comes from the divisor \(x_0y_1-x_1y_0=0\).

For every other morphism, we shall compute \(f^!\) from this adjunction, from \(j_!\dashv j^*\) for open immersions and from localization.

% object: lemma-cycle-class-cartier-divisor (blueprint/chap-Constructibility-of-f/sec-Duality-for-the-projective-line/lemma-cycle-class-cartier-divisor/_section.md)
\phantomsection\label{bp523}%
\begin{bpstatementit}{Lemma~6.3}{ (The cycle class of a Cartier divisor)}
Let \(\Lambda\) be a finite commutative \(\ell\)-primary ring, and let \(i:D\hookrightarrow Z\) be an effective Cartier divisor on a scheme \(Z\) on which \(\ell\) is invertible. There is a morphism

\[
c_D:i_*\Lambda_D\longrightarrow\Lambda_Z(1)[2]\quad\text{in }D(Z_{\acute{e}t},\Lambda),
\]

and hence, via \(\nu_Z^*\), in \(D(Z_{\mathrm{pro}\acute{e}t},\Lambda)\), whose composite with \(\Lambda_Z\to i_*\Lambda_D\) is the Kummer first Chern class \(c_1(\mathcal O_Z(D))\). It is compatible with change of coefficients and with every base change along which the pullback of the divisor is still an effective Cartier divisor.
\end{bpstatementit}

% proof: proof/lemma-cycle-class-cartier-divisor
\begin{proof}[{Proof}]\phantomsection\label{bp524}%
The map \(c_D\) will be the boundary map of the extension of \(i_*\mathbf Z_D\) by \(\mathbf G_m\) formed by meromorphic functions with poles along \(D\), reduced modulo \(\ell^r\); this extension is the étale version of the divisor sequence \(0\to\mathcal O_Z\to\mathcal O_Z(D)\to
\mathcal O_Z(D)|_D\to0\).

We construct \(c_D\) in \(D(Z_{\acute{e}t},\Lambda)\). The functor \(\nu_Z^*\), which commutes with inverse images and with change of coefficients, then transports it and its compatibilities to \(Z_{\mathrm{pro}\acute{e}t}\).

Let \(j:Z\setminus D\hookrightarrow Z\) be the open complement. Let \(\mathcal M_D\subset j_*\mathbf G_m\) be the subsheaf which, on an open where \(D\) has a regular equation \(t\), consists of the products \(ut^{-m}\) with \(u\) a unit and \(m\) a locally constant integer on \(D\); equivalently, it is generated by \(\mathbf G_m\) and \(t^{-1}\). Since replacing \(t\) by a unit multiple gives the same subsheaf, \(\mathcal M_D\) is globally defined. Sending \(t^{-1}\) to \(1\), we obtain a short exact sequence of abelian sheaves on \(Z_{\acute{e}t}\)

\[
\Sigma_D:\ 0\longrightarrow\mathbf G_m\longrightarrow \mathcal M_D
\longrightarrow i_*\mathbf Z_D\longrightarrow0,
\tag{6.3}
\]

whose exactness can be checked on geometric stalks. Away from \(D\), the function \(t\) is a unit, so there \(\mathcal M_D=\mathbf G_m\), and both the quotient and the stalk of \(i_*\mathbf Z_D\) are zero. At a point of \(D\), on the other hand, no nonzero power of a regular nonunit \(t\) is a unit, and therefore the exponent \(m\) in \(ut^{-m}\) is well defined. The kernel of the map \(ut^{-m}\mapsto m\) is then \(\mathbf G_m\), and the map is surjective because exponents lift locally to powers of \(t^{-1}\). For example, for \(D=\{0\}\subset Z=\mathbf A^1=\operatorname{Spec}k[t]\), the sections of \(\mathcal M_D\) near the origin are the functions \(ut^{-m}\), that is, the rational functions whose only zero or pole is at \(0\), and the map records the order of the pole, \(m=-\operatorname{ord}_0\).

The boundary of \(\Sigma_D\) is a morphism \(i_*\mathbf Z_D\to\mathbf G_m[1]\) in \(D(Z_{\acute{e}t},\mathbf Z)\), to which we apply \(-\otimes^{\mathbf L}_{\mathbf Z}\mathbf Z/\ell^r\mathbf Z\). The source becomes \(i_*\mathbf Z/\ell^r\mathbf Z\). For the target, the Kummer sequence [Stacks, \href{https://stacks.math.columbia.edu/tag/03PL}{03PL}] shows that the cone of multiplication by \(\ell^r\) on \(\mathbf G_m\) has the single cohomology sheaf \(\mu_{\ell^r}\) in degree \(-1\), and hence the target becomes \(\mu_{\ell^r}[2]\). This is \(c_D\) for \(\mathbf Z/\ell^r\mathbf Z\). For a finite \(\mathbf Z/\ell^r\mathbf Z\)-algebra \(\Lambda\), we apply \(-\otimes^{\mathbf L}_{\mathbf Z/\ell^r}\Lambda\). Since \(\mu_{\ell^r}\) is locally free of rank one over \(\mathbf Z/\ell^r\mathbf Z\), the target becomes \(\Lambda_Z(1)[2]\).

To identify the composite with \(c_1\), we cover \(Z\) by opens with local equations \(t_\alpha\). The local lifts of the section \(1\) of \(i_*\mathbf Z_D\) through \(\Sigma_D\) are the \(t_\alpha^{-1}\). Their ratios are the transition functions of \(\mathcal O_Z(D)\); this fixes the sign convention. Hence the boundary of \(\Sigma_D\) followed by the Kummer boundary is the class of \(\mathcal O_Z(D)\). Changing a local equation multiplies a lift by a unit, and this changes the cocycle by a coboundary.

Naturality holds because the local equations, the units and the sequence \(\Sigma_D\) pull back, compatibly with the boundary maps. On the source we use the identity \(g^*i_*=i'_*g_D^*\) for a base change \(g:Z'\to Z\) with \(i':D'\hookrightarrow Z'\) and \(g_D:D'\to D\); it can be checked on geometric stalks.
\end{proof}

% text: prose-06-05
\phantomsection\label{bp525}%
The lemma will be applied twice: to the section at infinity of \(\mathbf P^1_S\), which gives the trace, and to the diagonal of \(\mathbf P^1_S\times_S\mathbf P^1_S\), which gives the unit. The first application enters the computation of \(Rp_*A\).

% object: lemma-cohomology-projective-line-trace (blueprint/chap-Constructibility-of-f/sec-Duality-for-the-projective-line/lemma-cohomology-projective-line-trace/_section.md)
\phantomsection\label{bp526}%
\begin{bpstatementit}{Lemma~6.4}{ (Cohomology of the projective line and the trace)}
Let \(\Lambda\) be a finite commutative \(\ell\)-primary ring, let \(S\) be a Noetherian scheme on which \(\ell\) is invertible, and let \(p:\mathbf P^1_S\to S\) be the projection. The unit \(\Lambda_S\to Rp_*\Lambda_{\mathbf P^1_S}\) and the class \(c_1(\mathcal O_{\mathbf P^1_S}(1))\) define an isomorphism

\[
\sigma_S:\Lambda_S\oplus\Lambda_S(-1)[-2]
\ \xrightarrow{\ \sim\ }\ Rp_*\Lambda_{\mathbf P^1_S}
\quad\text{in }D^+(S_{\acute{e}t},\Lambda),
\tag{6.4}
\]

compatible with base change in the sense that for \(g:S'\to S\) with \(S'\) Noetherian, the base change map \(g^*Rp_*\Lambda\to Rp'_*\Lambda\) carries \(g^*\sigma_S\) to \(\sigma_{S'}\). Let

\[
\operatorname{Tr}_p:Rp_*\Lambda_{\mathbf P^1_S}(1)[2]\longrightarrow\Lambda_S
\tag{6.5}
\]

be the morphism given by the projection onto the summand \(\Lambda_S\) of the twist of \(\sigma_S\). Then \(\operatorname{Tr}_p\) is compatible with every Noetherian base change.
\end{bpstatementit}

% proof: proof/lemma-cohomology-projective-line-trace
\begin{proof}[{Proof}]\phantomsection\label{bp527}%
By the comparison theorems for the cohomology and the proper direct image of bounded below complexes (Theorems \hyperref[bp103]{2.35} and \hyperref[bp113]{Theorem 2.38}), it suffices to prove the assertions on \(\sigma_S\) and \(\operatorname{Tr}_p\) in \(D^+(S_{\acute{e}t},\Lambda)\) for Noetherian \(S\).

Both components of \(\sigma_S\) are globally defined. The second comes from Lemma \hyperref[bp523]{6.3} applied to the section at infinity, whose line bundle is \(\mathcal O(1)\). Geometric stalks detect isomorphisms, and by proper base change the stalk of \(Rp_*\Lambda\) at a geometric point \(\bar s\) is \(R\Gamma(\mathbf P^1_{\bar s},\Lambda)\). It therefore suffices to check that \(\sigma_S\) is an isomorphism on a geometric fibre, first for \(\Lambda=\mathbf Z/\ell^r\mathbf Z\). By the genus zero case of Proposition \bpout{D.2}, the cohomology lives only in degrees \(0\) and \(2\), with \(H^0=\mathbf Z/\ell^r\mathbf Z\) and \(H^2=(\mathbf Z/\ell^r\mathbf Z)(-1)\). The Kummer description sends a line bundle to its degree modulo \(\ell^r\). As \(\mathcal O(1)\) has degree one, its class generates \(H^2\), and \(\sigma_S\) is an isomorphism. The unit and the class of \(\mathcal O(1)\) commute with base change. Therefore so do \(\sigma_S\), its inverse, and \(\operatorname{Tr}_p\).

For a finite \(\mathbf Z/\ell^r\mathbf Z\)-algebra \(\Lambda\), the projection formula [Stacks, \href{https://stacks.math.columbia.edu/tag/0GL5}{0GL5}] gives \(Rp_*\Lambda=Rp_*(\mathbf Z/\ell^r\mathbf Z)
\otimes^{\mathbf L}_{\mathbf Z/\ell^r}\Lambda\). Under this identification, \(\sigma_S\) for \(\Lambda\) is \(\sigma_S\) for \(\mathbf Z/\ell^r\mathbf Z\) tensored with \(\Lambda\). It is therefore again an isomorphism compatible with base change.
\end{proof}

% text: prose-06-06
\phantomsection\label{bp528}%
Over \(S=\operatorname{Spec}\bar k\), (6.4) reads \(R\Gamma(\mathbf P^1_{\bar k},A)\simeq A\oplus A(-1)[-2]\). Here the shift \([-2]\) places the second summand in degree \(2\), so that it is \(H^2\), and the twist \((-1)\) records how Galois acts on \(H^2\): the Kummer map identifies \(H^2\) with \(\operatorname{Pic}(\mathbf P^1)\otimes\mu_{\ell^r}^{\otimes-1}\), that is, with \(A(-1)\). Twisting by \((1)\) and shifting by \([2]\) therefore moves this summand to degree \(0\) with trivial twist. The map \(\operatorname{Tr}_p\) is the projection onto this summand, and it sends the class of a point to \(1\). It is the counit of \(Rp_*\dashv p^*(1)[2]\) in Proposition \hyperref[bp534]{6.7}. After its lift to \(\mathcal O_E\) in Subsection \hyperref[bp579]{6.7.1}, it is also the trace map \(\operatorname{Tr}_p\) of \(\mathbf P^1\) in Section \bpout{9.6}.

As a check of the twist, let us work over \(\mathbf F_q\). The geometric Frobenius acts by \(1\) on \(H^0\) and by \(q\) on \(H^2=A(-1)\). The alternating sum of the traces is \(1+q=\#\mathbf P^1(\mathbf F_q)\), as in Corollary \bpout{12.31}.

% object: example-compact-supports-running-example (blueprint/chap-Constructibility-of-f/sec-Duality-for-the-projective-line/example-compact-supports-running-example/_section.md)
\phantomsection\label{bp529}%
\begin{bpstatement}{Example~6.5}{ (Compact supports on the running example)}
We use the open immersion \(j_\infty:\mathbf A^1\to\mathbf P^1\) of Example \hyperref[bp251]{3.19}, base changed to \(\bar k\), and we let \(i_\infty:\{\infty\}\to\mathbf P^1_{\bar k}\) be the inclusion of the complementary point.

By definition, \(R\Gamma_c(\mathbf A^1,A)=R\Gamma(\mathbf P^1,j_{\infty!}A)\). The localization triangle \(j_{\infty!}A\to A_{\mathbf P^1}\to i_{\infty*}A\) and (6.4) show that

\[
R\Gamma_c(\mathbf A^1_{\bar k},A)
\simeq\fib\bigl(A\oplus A(-1)[-2]\to A\bigr)
\simeq A(-1)[-2],
\]

since restriction to \(\infty\) is an isomorphism on \(H^0\). In other words, \(H^2_c(\mathbf A^1)=A(-1)\) is the only nonzero compactly supported cohomology group. The trace on it is the same \(\operatorname{Tr}_p\).
\end{bpstatement}

% text: prose-06-07
\phantomsection\label{bp530}%
\textbf{Exercise.} Using the localization sequence of Example \hyperref[bp251]{3.19}, show that \(H^1_c(\mathbf G_{m,\bar k},A)\simeq A\) and \(H^2_c(\mathbf G_{m,\bar k},A)\simeq A(-1)\), and that these are the only nonzero groups.

The adjunction is built from the trace \(\operatorname{Tr}_p\) of (6.5) and the cycle class of the diagonal provided by Lemma \hyperref[bp523]{6.3}. As in all of this part of the section, we compute on the étale sites (Section \hyperref[bp515]{6.1}). Let \(p:\mathbf P^1_S\to S\) be the projection, let \(p_1,p_2:\mathbf P^1_S\times_S\mathbf P^1_S\to\mathbf P^1_S\) be the two projections and let \(\Delta:\mathbf P^1_S\hookrightarrow\mathbf P^1_S\times_S\mathbf P^1_S\) be the diagonal. For \(i=1,2\), the projection \(p_i\) is a projective line over the other factor, with trace \(\operatorname{Tr}_{p_i}\) given by (6.5). Both triangle identities rest on the following computation.

% object: lemma-class-diagonal-projective-line (blueprint/chap-Constructibility-of-f/sec-Duality-for-the-projective-line/lemma-class-diagonal-projective-line/_section.md)
\phantomsection\label{bp531}%
\begin{bpstatementit}{Lemma~6.6}{ (The class of the diagonal of $\mathbf P^1_S$)}
Let \(S\) be of finite type over \(k\), and let \(p_i\) and \(\operatorname{Tr}_{p_i}\) be as above. The diagonal \(\Delta\) is a relative effective Cartier divisor in \(\mathbf P^1_S\times_S\mathbf P^1_S\) over \(S\), and it remains one after every base change \(S'\to S\). Let

\[
\delta:\Delta_*A_{\mathbf P^1_S}\longrightarrow
A_{\mathbf P^1_S\times_S\mathbf P^1_S}(1)[2]
\]

be its cycle class (Lemma \hyperref[bp523]{6.3}). Then for \(i=1,2\) the composite

\[
s_i:A_{\mathbf P^1_S}=Rp_{i*}\Delta_*A_{\mathbf P^1_S}
\xrightarrow{Rp_{i*}\delta}
Rp_{i*}A_{\mathbf P^1_S\times_S\mathbf P^1_S}(1)[2]
\xrightarrow{\operatorname{Tr}_{p_i}}A_{\mathbf P^1_S}
\tag{6.6}
\]

is the identity.
\end{bpstatementit}

% proof: proof/lemma-class-diagonal-projective-line
\begin{proof}[{Proof}]\phantomsection\label{bp532}%
Let \(x_0,x_1\) and \(y_0,y_1\) be the pullbacks along \(p_1\) and \(p_2\) of the homogeneous coordinates of \(\mathbf P^1_S\). Then \(\Delta\) is the zero locus of the section \(x_0y_1-x_1y_0\) of \(p_1^*\mathcal O(1)\otimes p_2^*\mathcal O(1)\). On the open where \(x_0\) and \(y_0\) are invertible, in the affine coordinates \(x=x_1/x_0\) and \(y=y_1/y_0\), this equation becomes \(y-x\). It is a nonzerodivisor, and the quotient \(\mathcal O_S[x,y]/(y-x)\simeq\mathcal O_S[x]\) is flat over \(S\). The open where \(x_1\) and \(y_1\) are invertible is treated in the same way, and these two opens cover \(\Delta\). Hence \(\Delta\) is a relative effective Cartier divisor, and it remains one after every base change.

To show \(s_i=1\), note that both sides are endomorphisms of a sheaf placed in degree zero. We have

\[
\Hom_{D((\mathbf P^1_S)_{\acute{e}t},A)}
(A_{\mathbf P^1_S},A_{\mathbf P^1_S})
=\Hom_{\operatorname{Mod}((\mathbf P^1_S)_{\acute{e}t},A)}
(A_{\mathbf P^1_S},A_{\mathbf P^1_S}),
\]

so \(s_i\) is a morphism of sheaves and is determined by its geometric stalks. By proper base change and the compatibility of \(\delta\) and \(\operatorname{Tr}_{p_i}\) with base change, the stalk of (6.6) at a geometric point \(\bar x\) of \(\mathbf P^1_S\) is computed on the fibre \(p_i^{-1}(\bar x)\simeq\mathbf P^1_{\bar x}\). On this fibre the equation \(x_0y_1-x_1y_0\) restricts to a nonzero linear form, so the diagonal cuts out a single reduced point, whose line bundle is \(\mathcal O(1)\). By Lemma \hyperref[bp523]{6.3}, the class of this point composed with the unit is \(c_1(\mathcal O(1))\). In turn, the trace \(\operatorname{Tr}_p\) sends \(c_1(\mathcal O(1))\) to \(1\), since by definition (6.5) it is the projection onto the summand of (6.4) spanned by this class. Thus the stalk of \(s_i\) at every geometric point \(\bar x\) sends \(1\) to \(1\), and so \(s_i=1\).
\end{proof}

% text: prose-06-08
\phantomsection\label{bp533}%
We build the unit and counit of the adjunction from \(\delta\) and from the trace. For \(M\in D((\mathbf P^1_S)_{\acute{e}t},A)\), let

\[
\begin{split}
\eta_M:M
&=Rp_{2*}\Delta_*M\\
&\longrightarrow Rp_{2*}\bigl(p_1^*M(1)[2]\bigr)
\xrightarrow{\ \sim\ }p^*Rp_*M(1)[2] .
\end{split}
\tag{6.7}
\]

Here the first arrow is \(Rp_{2*}\) applied to \(\delta\otimes p_1^*M\), where we use the projection formula \(\Delta_*M=\Delta_*A\otimes^{\mathbf L}p_1^*M\), and the second arrow is the proper base change isomorphism \(p^*Rp_*\simeq Rp_{2*}p_1^*\). For \(N\in D(S_{\acute{e}t},A)\), let

\[
\varepsilon_N:Rp_*\bigl(p^*N(1)[2]\bigr)
\xrightarrow{\ \sim\ }N\otimes_A^{\mathbf L}
Rp_*A_{\mathbf P^1_S}(1)[2]
\xrightarrow{1\otimes \operatorname{Tr}_p}N,
\tag{6.8}
\]

where the first arrow is the projection formula. These maps exist for unbounded complexes by [Stacks, \href{https://stacks.math.columbia.edu/tag/0F0C}{0F0C}, \href{https://stacks.math.columbia.edu/tag/0GL5}{0GL5}].

% object: proposition-adjunction-projective-line (blueprint/chap-Constructibility-of-f/sec-Duality-for-the-projective-line/proposition-adjunction-projective-line/_section.md)
\phantomsection\label{bp534}%
\begin{bpstatementit}{Proposition~6.7}{ (The adjunction for the projective line)}
Let \(S\) be of finite type over \(k\), and let \(p:\mathbf P^1_S\to S\) be the projection. The functors \(Rp_*:D((\mathbf P^1_S)_{\acute{e}t},A)\to
D(S_{\acute{e}t},A)\) and \(p^*(1)[2]:D(S_{\acute{e}t},A)\to
D((\mathbf P^1_S)_{\acute{e}t},A)\) form an adjunction

\[
Rp_* \ \dashv\ p^*(1)[2],
\tag{6.9}
\]

with unit (6.7) and counit (6.8). For every morphism \(g:S'\to S\) with \(S'\) of finite type over \(k\), the unit and counit are compatible with base change along \(g\). They are also compatible with change of finite coefficients.
\end{bpstatementit}

% proof: proof/proposition-adjunction-projective-line
\begin{proof}[{Proof}]\phantomsection\label{bp535}%
The projection formula and proper base change reduce each of the two triangle identities to the identity \(s_i=1\) of Lemma \hyperref[bp531]{6.6}, tensored with the identity of a coefficient complex.

\emph{First triangle identity.} This is the equation \(\varepsilon_{Rp_*M}\circ Rp_*(\eta_M)=1_{Rp_*M}\). Using composition of direct images, \(pp_1=pp_2\), proper base change and the projection formula, this composite equals

\[
\begin{split}
Rp_*M
&=R(pp_1)_*\Delta_*M\\
&\longrightarrow R(pp_1)_*p_1^*M(1)[2]\\
&\xrightarrow{\ \sim\ }
Rp_*\bigl(M\otimes_A^{\mathbf L}
Rp_{1*}A_{\mathbf P^1_S\times_S\mathbf P^1_S}(1)[2]\bigr)\\
&\xrightarrow{Rp_*(1_M\otimes \operatorname{Tr}_{p_1})}Rp_*M .
\end{split}
\tag{6.10}
\]

Here the first arrow is \(Rp_*\) of the first arrow of (6.7), with \(R(pp_2)_*\) rewritten as \(R(pp_1)_*\). The middle isomorphism is the projection formula for \(p_1\). It replaces the proper base change isomorphism of (6.7), after \(Rp_*\), followed by the projection formula isomorphism of (6.8). The last arrow replaces \(1\otimes\operatorname{Tr}_p\) in (6.8) by \(\operatorname{Tr}_{p_1}\), the base change of \(\operatorname{Tr}_p\) along \(p\) (Lemma \hyperref[bp526]{6.4}). In this form, the trace is applied to \(Rp_{1*}\delta\), that is, to the factor \(\Delta_*A\), and \(M\) only appears as a coefficient.

This rewriting is an equality of maps because the projection formula isomorphism is compatible with proper base change. To see this, recall that the proper base change map is the adjoint transpose of the counit of \(g^*\dashv Rg_*\) for the relevant pullback \(g\). The projection map is the adjoint transpose of the same counit tensored with the identity, and the counit of a composite adjunction is the composite of the counits. Therefore the two routes from \(R(pp_2)_*p_1^*M=R(pp_1)_*p_1^*M\) to \(Rp_*M\otimes Rp_*A(1)[2]\) have the same adjoint transpose, as in the pasting Lemma \hyperref[bp299]{4.5}.

On the route through (6.10) we identify \(Rp_{1*}A=p^*Rp_*A\) by proper base change and then use the projection formula for \(p\). By naturality of the projection formula in its second variable, the first arrow of (6.10) is then \(1_M\otimes Rp_{1*}\delta\) under the projection identification. The composite is therefore \(Rp_*(1_M\otimes s_1)\), which is \(1_{Rp_*M}\) by Lemma \hyperref[bp531]{6.6}.

\emph{Second triangle identity.} Here the roles of the two projections are exchanged: we must show that \(p^*(\varepsilon_N)(1)[2]\circ\eta_{p^*N(1)[2]}\) is the identity. Take \(M=p^*N(1)[2]\) in (6.7). Since \(N\) and the Tate twist are pulled back from \(S\), both \(p_1^*M\) and \(p_2^*M\) are \((pp_1)^*N(1)[2]=(pp_2)^*N(1)[2]\). The projection formula moves this common factor outside \(Rp_{2*}\), and \(\eta_{p^*N(1)[2]}\) becomes \(1\otimes Rp_{2*}\delta\). By the base change compatibility of (6.5), \(p^*(\varepsilon_N)(1)[2]\) becomes \(1\otimes\operatorname{Tr}_{p_2}\), applied to the target of \(Rp_{2*}\delta\). We obtain

\[
p^*(\varepsilon_N)(1)[2]\circ\eta_{p^*N(1)[2]}
=1_{p^*N(1)[2]}\otimes
(\operatorname{Tr}_{p_2}\circ Rp_{2*}\delta)
=1_{p^*N(1)[2]}\otimes s_2=1_{p^*N(1)[2]}.
\tag{6.11}
\]

All shifts moved in (6.10) and (6.11) are even, so no Koszul signs appear.

The unit is constructed from \(\delta\) and proper base change, and the counit from (6.5) and the projection formula. All of these ingredients are compatible with base change and change of coefficients, and hence so are \(\eta\) and \(\varepsilon\). Transported by \(\nu^*\), the same then holds for the corresponding adjunction on the pro-étale categories.
\end{proof}


% object: sec-Smooth-morphisms (blueprint/chap-Constructibility-of-f/sec-Smooth-morphisms/_section.md), content hidden
\subsection*{\texorpdfstring{6.3\quad Smooth morphisms}{6.3 Smooth morphisms}}
\phantomsection\label{bp536}%
\addcontentsline{toc}{subsection}{\texorpdfstring{6.3\quad Smooth morphisms}{6.3 Smooth morphisms}}
% text: prose-06-09
\phantomsection\label{bp537}%
Let \(f:X\to S\) be smooth of relative dimension \(d\). Extending the formula \(p^!=p^*(1)[2]\) for \(p:\mathbf P^1_S\to S\), we show in Proposition \hyperref[bp542]{6.9} that \(f^!N\simeq f^*N\otimes^{\mathbf L}\omega_f\) with \(\omega_f\) locally isomorphic to \(A(d)[2d]\). We need this formula in the proof of generic constructibility below, because it must control \(Rf_*L\) for a locally constant sheaf \(L\). Indeed, we know nothing about \(Rf_*L\) directly, whereas the finiteness theorem and proper base change control \(Rf_!\). With the formula, however, the adjunction \(Rf_!\dashv f^!\) becomes a duality between \(Rf_*L\) and \(Rf_!\) of the twisted dual of \(L\) (Lemma \hyperref[bp545]{6.10}). On a manifold, Poincaré duality exchanges ordinary and compactly supported cohomology in the same way.

For example, let \(f:\mathbf G_m\to\operatorname{Spec}\bar k\) and \(L=A=\Fell\). The localization triangle for the origin of \(\mathbf A^1\) and \(R\Gamma_c(\mathbf A^1_{\bar k},A)\simeq A(-1)[-2]\) give \(R\Gamma_c(\mathbf G_m,A)\simeq\fib\bigl(A(-1)[-2]\to A\bigr)\). Hence \(H^1_c=A\) and \(H^2_c=A(-1)\) are the only nonzero groups, and \(R\Gamma_c(\mathbf G_m,A(1)[2])\) has cohomology \(A(1)\) in degree \(-1\) and \(A\) in degree \(0\). Its dual has \(A\) in degree \(0\) and \(A(-1)\) in degree \(1\). This is indeed \(R\Gamma(\mathbf G_m,A)\), with \(H^1=A(-1)\) given by the Kummer torsors.

The duality must also commute with base change. This requires the base change isomorphism \(g'^*f^!\simeq f'^!g^*\) for smooth \(f\), and for it the local form \(\omega_f|_U\simeq A_U(d)[2d]\) on charts suffices. A global identification of \(\omega_f\) with \(A_X(d)[2d]\), by contrast, needs a trace map for \(f\); it is part of smooth purity (Chapter \bpout{9}). For the local form we use that, Zariski locally on \(X\) and \(S\), a smooth morphism \(f:X\to S\) factors as an étale morphism to some \(\mathbf A^d_S\) followed by the projection to \(S\), and the next lemma treats these two kinds of morphisms.

% object: lemma-affine-space-etale-localization (blueprint/chap-Constructibility-of-f/sec-Smooth-morphisms/lemma-affine-space-etale-localization/_section.md)
\phantomsection\label{bp538}%
\begin{bpstatementit}{Lemma~6.8}{ (Affine space and étale localization)}
Let \(S\) be of finite type over the base field \(k\).

\begin{enumerate}
\item For the projection \(p_d:\mathbf A^d_S\to S\), the right adjoint \(p_d^!\) of \(Rp_{d!}:D((\mathbf A^d_S)_{\acute{e}t},A)
\to D(S_{\acute{e}t},A)\) is \(p_d^*(d)[2d]\). Its counit \(Rp_{d!}p_d^*N(d)[2d]\to N\) is compatible with base change: for a morphism \(g:S'\to S\) of finite type over \(k\), with base change \(p'_d:\mathbf A^d_{S'}\to S'\), the pullback \(g^*\) of the counit for \(N\) is, under the base change isomorphism for \(Rp_{d!}\), the counit for \(g^*N\). The same holds for the unit, which is determined by the counit through the triangle identities.
\item Let \(f:X\to S\) be separated of finite type, let \(N\in D(S_{\acute{e}t},A)\), and let \(u:U\to X\) be separated, étale and of finite type. There is a canonical isomorphism \[
   u^*f^!N\simeq(fu)^!N
   \tag{6.12}
   \] under which the counit of \(R(fu)_!\dashv(fu)^!\) is the composite \[
   R(fu)_!u^*f^!N\simeq Rf_!u_!u^*f^!N\longrightarrow Rf_!f^!N
   \longrightarrow N
   \] of the composition isomorphism, the counit of \(u_!\dashv u^*\) and the counit of \(Rf_!\dashv f^!\). For a second separated étale morphism \(u':U'\to U\) of finite type, the isomorphism (6.12) for \(uu'\) is the composite of \(u'^*\) of (6.12) for \(u\) and (6.12) for \(u'\) and \(fu\).
\end{enumerate}
\end{bpstatementit}

% text: prose-06-10
\phantomsection\label{bp539}%
Part (2) generalizes \(j^!=j^*\) from open immersions to étale maps. For \(d=1\), \(k\) algebraically closed and \(S=\operatorname{Spec}k\), part (1) says that \(\Hom(R\Gamma_c(\mathbf A^1,F),M)=\Hom(F,M_{\mathbf A^1}(1)[2])\). For \(F=A\) both sides are \(\Hom(A,M(1)[2])\), consistent with the example above: the left side because \(R\Gamma_c(\mathbf A^1_{\bar k},A)\simeq A(-1)[-2]\), and the right side because \(R\Gamma(\mathbf A^1_{\bar k},A)=A\).

% proof: proof/lemma-affine-space-etale-localization
\begin{proof}[{Proof of Lemma 6.8}]\phantomsection\label{bp540}%
Both parts follow formally from uniqueness of right adjoints. Part (1) comes from (6.9) and \(j_!\dashv j^*\) for \(\mathbf A^1\subset\mathbf P^1\), and part (2) from \(u_!\dashv u^*\) for étale \(u\).

(1) Let \(j:\mathbf A^1_S\hookrightarrow\mathbf P^1_S\) be the open immersion, so that \(Rp_{1!}=Rp_*j_!\) for \(d=1\). Since \(j_!\dashv j^*\) and \(Rp_*\dashv p^*(1)[2]\) by (6.9), composing right adjoints gives \(p_1^!=j^*p^*(1)[2]\). This equals \(p_1^*(1)[2]\) because \(p\circ j=p_1\). For general \(d\), write \(p_d\) as a composite of \(d\) projections \(r_1,\dots,r_d\) of affine lines. By the composition isomorphism for \(Rf_!\) (Proposition proposition-composition-base-change-rf-shriek) we have \(Rp_{d!}\simeq Rr_{1!}\circ\dots\circ Rr_{d!}\). Hence \(p_d^!\) is the composite of the right adjoints of the \(Rr_{i!}\), which is \(p_d^*(d)[2d]\) by the case \(d=1\). Under this isomorphism, the counit of \(Rp_{d!}\dashv p_d^*(d)[2d]\) is the composite of the \(d\) counits for the affine lines. Each of these is a trace (6.5) composed with the counit of \(j_!\dashv j^*\). By the associativity in Proposition proposition-composition-base-change-rf-shriek, the composite does not depend on how the \(d\) factors are grouped. Each factor and the composition isomorphism commute with base change (Proposition proposition-composition-base-change-rf-shriek), and therefore so does the counit. The unit is determined by the counit through the triangle identities, and it commutes with base change as well.

(2) For an étale \(u\), the functor \(Ru_!\) is extension by zero. By its elementary construction [Stacks, \href{https://stacks.math.columbia.edu/tag/03S2}{03S2}], it is an exact functor \(u_!:\operatorname{Mod}(U_{\acute{e}t},A)\to\operatorname{Mod}(X_{\acute{e}t},A)\) left adjoint to \(u^*\). As both functors are exact, they need not be derived, and the adjunction extends to derived categories. For every \(M\in D(U_{\acute{e}t},A)\),

\[
\Hom(M,u^*f^!N)
=\Hom(u_!M,f^!N)
=\Hom(Rf_!u_!M,N)
=\Hom(R(fu)_!M,N),
\]

where the last step is the composition isomorphism of Proposition proposition-composition-base-change-rf-shriek. The Yoneda lemma therefore gives (6.12). Taking \(M=u^*f^!N\) and the identity on the left, we see that the counit of \(R(fu)_!\dashv(fu)^!\) is the composite of the counits as stated. For the cocycle condition, the isomorphism for \(uu'\) and the composite of those for \(u\) and \(u'\) both correspond under the Yoneda lemma to the bijections obtained from the composition isomorphism for \(f\circ u\circ u'\). These bijections agree by associativity of the composition isomorphisms (Proposition proposition-composition-base-change-rf-shriek).
\end{proof}

% text: prose-06-11
\phantomsection\label{bp541}%
Applied along local charts \(U\xrightarrow{e}\mathbf A^d_S\to S\), the two parts of the lemma compute \(f^!\) locally on \(X\) for every smooth \(f\). The proposition below turns this into two global statements, a tensor formula and base change.

% object: proposition-exceptional-inverse-image-smooth-morphism (blueprint/chap-Constructibility-of-f/sec-Smooth-morphisms/proposition-exceptional-inverse-image-smooth-morphism/_section.md)
\phantomsection\label{bp542}%
\begin{bpstatementit}{Proposition~6.9}{ (The exceptional inverse image of a smooth morphism)}
Let \(f:X\to S\) be separated, smooth and of finite type, of pure relative dimension \(d\), and let \(\omega_f=f^!A_S\in D(X_{\acute{e}t},A)\) be its dualizing complex. Let \(N\in D(S_{\acute{e}t},A)\).

\begin{enumerate}
\item On an open \(U\subset X\) with \(f|_U=p_d\circ e\), where \(e:U\to\mathbf A^d_S\) is separated and étale, we have \((f^!N)|_U\simeq e^*p_d^*N(d)[2d]\), with counit the sum over the fibres of \(e\) followed by the counit of Lemma lemma-affine-space-etale-localization. In particular \(\omega_f|_U\simeq A_U(d)[2d]\), and \(\omega_f[-2d]\) is a locally constant sheaf of free \(A\)-modules of rank one.
\item The canonical map \[
   \mu_N:f^*N\otimes_A^{\mathbf L}\omega_f\longrightarrow f^!N
   \tag{6.13}
   \] is an isomorphism.
\item Let \(g:S'\to S\) be a morphism, let \(X'=X\times_SS'\), and let \(f':X'\to S'\), \(g':X'\to X\) be the projections. Then the base change map \(\beta_N:g'^*f^!N\to f'^!g^*N\) of Definition \hyperref[bp519]{6.2} is an isomorphism. It identifies \(g'^*\omega_f\) with \(\omega_{f'}\) and is compatible with \(\mu_N\).
\end{enumerate}
\end{bpstatementit}

% proof: proof/proposition-exceptional-inverse-image-smooth-morphism
\begin{proof}[{Proof}]\phantomsection\label{bp543}%
The maps \(\mu_N\) and \(\beta_N\) are defined globally as adjoint transposes. Since being an isomorphism is local on \(X\), respectively on \(X'\), we show that on a chart both maps become identifications of inverse images under the explicit identification (6.14). For this we follow the counits through the factorization \(U\to\mathbf A^d_V\to V\).

We write \(\varepsilon_N:Rf_!f^!N\to N\) for the counit of \(Rf_!\dashv f^!\) (for \(f=p\) this is (6.8)). The map \(\mu_N\) is the transpose of the morphism

\[
Rf_!(f^*N\otimes^{\mathbf L}\omega_f)
\xrightarrow{\ \sim\ }
N\otimes^{\mathbf L}Rf_!\omega_f
\xrightarrow{1\otimes\varepsilon_{A_S}}N ,
\]

where the first arrow is the projection formula. The map \(\beta_N\) is the transpose, for \(f'\), of the morphism (6.2) (Definition \hyperref[bp519]{6.2}).

\emph{Local charts.} The computation of \(f^!\) on a chart also proves (1). Zariski locally on \(X\) and \(S\), a smooth morphism is an étale map to affine space followed by the projection [Stacks, \href{https://stacks.math.columbia.edu/tag/054L}{054L}]. Thus every point of \(X\) has an affine open neighbourhood \(u:U\hookrightarrow X\) mapping into an affine open \(v:V\hookrightarrow S\), together with a separated étale map \(e:U\to\mathbf A^d_V\) such that \(fu=vp_de\), where \(p_d:\mathbf A^d_V\to V\) is the projection. By Lemma \hyperref[bp538]{6.8} and uniqueness of right adjoints, we have

\[
u^*f^!N=e^*p_d^*v^*N(d)[2d],\qquad
u^*\omega_f=A_U(d)[2d].
\tag{6.14}
\]

In particular, \(\omega_f\) has a single nonzero cohomology sheaf, in degree \(-2d\), which is locally free of rank one. This proves (1), because the counit on \(U\) is the composite of the étale counit for \(e\) and the counit for \(p_d\). The first of these is the sum over the fibres.

\emph{\(\mu_N\) is an isomorphism.} On each such \(U\), under (6.14), the source and the target of \(\mu_N\) are both \(e^*p_d^*v^*N(d)[2d]\). We claim that the map is the identification obtained by reassociating tensor product and twist. It suffices to compare the transposes. Both are the projection formula followed by the counit for \(v\), the \(d\) counits for the affine lines, and the étale counit for \(e\). They agree by the formula for the counit of a composite in Lemma \hyperref[bp538]{6.8}. Hence \(\mu_N\) is an isomorphism on \(X\).

\emph{\(\beta_N\) is an isomorphism.} For each type of chart, under the explicit description of \(f^!\) as a twisted inverse image, \(\beta_N\) is the evident identification of inverse images. For the projective line, the transpose of (6.2) is the usual identification \(g'^*p^*N(1)[2]=p'^*g^*N(1)[2]\) of pullbacks. This holds because the unit (6.7) and counit (6.8) commute with base change. By restriction to the affine line and composition, the same holds for affine space. For a separated étale map \(u\), the transpose is again the usual identification of inverse images. Indeed, the stalk of \(u_!M\) at a geometric point is the direct sum of the stalks of \(M\) at the geometric points over it. This description is compatible with base change, and under it the counit \(u_!u^*N\to N\) is the sum of the identity maps of the stalks of \(N\).

The cases combine, because base change for \(Rf_!\) is compatible with composition (Proposition proposition-composition-base-change-rf-shriek) and the counit of a composite is the composite of the counits. The open immersions \(u\) and \(v\) are étale, so the étale case covers them as well. Substituting these facts into (6.2), we see that the restriction of \(\beta_N\) to \(U\times_XX'\) is the composite of the étale and affine-space cases just treated. Hence \(\beta_N\) is an isomorphism on a cover of \(X'\), and therefore it is an isomorphism.

\emph{Compatibilities.} Taking \(N=A_S\), we see that \(\beta_N\) identifies \(g'^*\omega_f\) with \(\omega_{f'}\), because \(g^*A_S=A_{S'}\). The compatibility of \(\mu_N\) with \(\beta_N\) says that two maps \(g'^*(f^*N\otimes^{\mathbf L}\omega_f)\to f'^!g^*N\) agree. Here the first is \(\mu_N\) followed by \(\beta_N\), and the second is the identification \(g'^*(f^*N\otimes^{\mathbf L}\omega_f)=f'^*g^*N\otimes^{\mathbf L}\omega_{f'}\), coming from \(g'^*f^*=f'^*g^*\) and the case \(N=A_S\), followed by \(\mu_N\) for \(f'\). Their transposes are maps \(Rf'_!g'^*(f^*N\otimes^{\mathbf L}\omega_f)\to g^*N\). Both are given by the projection formula followed by the counit, once before and once after the base change isomorphism for \(Rf_!\). They agree for two reasons. First, the projection formula isomorphism is compatible with the base change isomorphism for \(Rf_!\). Second, by the definition (6.2) of \(\beta_N\), the counit for \(f'\) composed with \(Rf'_!\) of \(\beta_N\) is \(g^*\) of the counit for \(f\) under that isomorphism.
\end{proof}

% text: prose-06-12
\phantomsection\label{bp544}%
In Step 2 of the proof of Theorem \hyperref[bp563]{6.17}, the next lemma is applied to \(f=bj\), with \(b:Y\to S\) projective and \(j:X\hookrightarrow Y\) open.

% object: lemma-smooth-morphisms-locally-constant-sheaves (blueprint/chap-Constructibility-of-f/sec-Smooth-morphisms/lemma-smooth-morphisms-locally-constant-sheaves/_section.md)
\phantomsection\label{bp545}%
\begin{bpstatementit}{Lemma~6.10}{ (Smooth morphisms and locally constant sheaves)}
Let \(S\) be integral and let \(f:X\to S\) be separated, smooth and of finite type. For a finite locally constant \(\Fell\)-sheaf \(L\) on \(X_{\acute{e}t}\), there is a nonempty open \(S_0\subset S\) with the following property. For the restriction \(f_{S_0}:X_{S_0}\to S_0\), the complex \(Rf_{S_0,*}L\) is bounded with finite locally constant cohomology sheaves, and its formation commutes with base change. Namely, for every \(k\)-scheme \(v:T\to S_0\) of finite type, with base change \(f_T:X_T\to T\) and \(v_X:X_T\to X_{S_0}\), the canonical base change map \(v^*Rf_{S_0,*}L\to Rf_{T,*}v_X^*L\) is an isomorphism.
\end{bpstatementit}

% proof: proof/lemma-smooth-morphisms-locally-constant-sheaves
\begin{proof}[{Proof}]\phantomsection\label{bp546}%
We write \(Rf_*L\) as the dual of \(Rf_!N\) for a suitable constructible \(N\). The point is that constructibility and base change are known for \(Rf_!N\), and over a dense open of \(S\) the cohomology sheaves of \(Rf_!N\) are locally constant, and there dualizing preserves both properties.

By Proposition \hyperref[bp126]{2.42}, the locally constant sheaf \(L\) comes from a locally constant sheaf on the étale site, and we work with the latter.

Since \(X\) is the disjoint union of finitely many open pieces of pure relative dimension, we may assume that \(f\) has pure relative dimension \(d\). Let \(\omega_f=f^!\mathbf F_{\ell,S}\) as in Proposition \hyperref[bp542]{6.9}, and let

\[
N=\cHom_X(L,\mathbf F_{\ell,X})\otimes^{\mathbf L}\omega_f .
\]

By Proposition \hyperref[bp542]{6.9}, \(\omega_f\) is locally \(\mathbf F_{\ell,X}(d)[2d]\). Hence \(N\) is bounded constructible and \(\cHom_X(N,\omega_f)=L\).

\emph{Duality.} Consider the pairing

\[
Rf_!N\otimes^{\mathbf L}Rf_*L
\longrightarrow Rf_!(N\otimes^{\mathbf L}L)
\longrightarrow Rf_!\omega_f
\xrightarrow{\varepsilon_{\mathbf F_{\ell,S}}}\mathbf F_{\ell,S},
\tag{6.15}
\]

where the first arrow is the projection formula followed by the counit \(f^*Rf_*L\to L\), and the second is evaluation. Its adjoint is a map

\[
\delta_L:Rf_*L
\longrightarrow \cHom_S(Rf_!N,\mathbf F_{\ell,S}),
\tag{6.16}
\]

and we claim that it is an isomorphism. Since the map is an adjoint transpose, we test it on Hom groups out of a family of generators, where adjunctions turn one side into the other. To this end we write \(L=\cHom_X(N,f^!\mathbf F_{\ell,S})\), and test \(\delta_L\) on the objects \(P[n]\), where \(P=u_!\Fell\) and \(u:V\to S\) is quasi-compact, separated and étale. The objects \(u_!\Fell\) are constructible of finite Tor amplitude, and their shifts detect isomorphisms in \(D(S_{\acute{e}t},\Fell)\). Since \(f^*P\otimes^{\mathbf L}N\) lies in \(D^b_{\mathrm{ctf}}(X,\Fell)\), we may use tensor-Hom adjunction, the adjunction \(Rf_!\dashv f^!\) of Theorem \hyperref[bp506]{5.31} between \(\Loc(D^b_{\mathrm{ctf}}(X,\Fell))\) and \(\Loc(D^b_{\mathrm{ctf}}(S,\Fell))\), and the projection formula. With Hom groups in \(D(S_{\acute{e}t},\Fell)\), we then obtain

\[
\Hom(P,Rf_*\cHom(N,f^!\mathbf F_{\ell,S}))
=\Hom(Rf_!(f^*P\otimes^{\mathbf L}N),\mathbf F_{\ell,S})
=\Hom(P\otimes^{\mathbf L}Rf_!N,\mathbf F_{\ell,S}).
\]

By tensor-Hom adjunction, the last group is \(\Hom(P,\cHom_S(Rf_!N,\mathbf F_{\ell,S}))\). Unwinding the adjunctions, one sees that the composite bijection is the map induced by \(\delta_L\), for every shift of \(P\). Hence \(\delta_L\) is an isomorphism.

\emph{Shrinking the base.} Since \(\delta_L\) is an isomorphism, it now suffices to understand the dual of \(Rf_!N\). By the finiteness theorem for \(Rf_!\) [Stacks, \href{https://stacks.math.columbia.edu/tag/0GL0}{0GL0}, \href{https://stacks.math.columbia.edu/tag/0G2A}{0G2A}], the complex \(Rf_!N\) is bounded constructible. A constructible sheaf on the integral scheme \(S\) is locally constant on a nonempty open subset. Thus each of the finitely many nonzero cohomology sheaves of \(Rf_!N\) is locally constant on a nonempty open of \(S\). Let \(S_0\) be the intersection of these opens; it is nonempty because \(S\) is irreducible.

Over the field \(\Fell\), a finite locally constant sheaf is locally free of finite rank. Its dual is exact, and evaluation identifies it with its double dual. We now argue by induction on the number of nonzero cohomology sheaves, using that the relevant canonical maps form morphisms of the truncation triangles. In this way we see that a bounded complex with finite locally constant cohomology is dualizable, that its dual is again of this kind, and that forming the dual commutes with every pullback. Through \(\delta_L\), this proves the first assertion over \(S_0\).

\emph{Base change.} The duality \(\delta_L\) exchanges \(Rf_*L\) with the dual of \(Rf_!N\), and \(Rf_!N\) satisfies base change. We show that the duality is compatible with base change; then \(Rf_*L\) satisfies base change too. Let \(v:T\to S_0\) be a \(k\)-scheme of finite type over \(S_0\), and let \(f_T:X_T\to T\) and \(v_X:X_T\to X\) be the base changes. First, base change for \(Rf_!\) and (6.1) identify \(v_X^*N\), \(v^*Rf_!N\) and \(v_X^*\omega_f\) with the corresponding objects for \(f_T\). Next, the base change map \(v^*Rf_*L\to Rf_{T,*}v_X^*L\) is characterized by the fact that it carries the pullback of the counit \(f^*Rf_*L\to L\) to the counit for \(f_T\). Hence the pullback of (6.15) is the pairing for \(f_T\), composed with this base change map. Finally, the complex \(v^*Rf_!N\) is dualizable with dual \(v^*\cHom(Rf_!N,\mathbf F_{\ell,S})\). Therefore the base change map for \(Rf_*L\) is identified with the dual of the base change isomorphism for \(Rf_!N\), and it is an isomorphism.
\end{proof}

% object: example-shrinking-base-smooth (blueprint/chap-Constructibility-of-f/sec-Smooth-morphisms/example-shrinking-base-smooth/_section.md)
\phantomsection\label{bp547}%
\begin{bpstatement}{Example~6.11}{ (Why the base must be shrunk)}
Let \(k\) be algebraically closed, \(S=\mathbf A^1\) with coordinate \(t\), and let \(X\subset\mathbf A^1_S\) be the complement of the two sections \(x=0\) and \(x=t\), with \(f:X\to S\) the projection and \(L=\Fell\). The morphism \(f\) is smooth. For \(t\neq0\) its fibre is \(\mathbf A^1\) minus two points, that is, a projective line minus three points. Its \(H^1\) is \(\Fell(-1)^{\oplus2}\). The fibre over \(t=0\), however, is \(\mathbf G_m\), with \(H^1=\Fell(-1)\). On \(S_0=\mathbf G_m\) the lemma applies, but no open containing the origin works. Suppose \(R^1f_*L\) were locally constant near \(0\) and commuted with base change to the points of \(S\). Then its stalks at \(0\) and at nearby \(t\) would be the \(H^1\) of the fibres, which have different ranks. What happens is that the two missing points collide at \(t=0\), and \(Rf_*\) cannot see this from the generic fibre.
\end{bpstatement}


% object: sec-Coordinate-projections (blueprint/chap-Constructibility-of-f/sec-Coordinate-projections/_section.md), content hidden
\subsection*{\texorpdfstring{6.4\quad Coordinate projections}{6.4 Coordinate projections}}
\phantomsection\label{bp548}%
\addcontentsline{toc}{subsection}{\texorpdfstring{6.4\quad Coordinate projections}{6.4 Coordinate projections}}
% text: prose-06-13
\phantomsection\label{bp549}%
To control \(Rj_*F\) for an open immersion \(j:X\hookrightarrow Y\) over \(S\), with \(Y\) projective over \(S\), we cover most of \(Y\) by opens where the induction hypothesis applies and treat the rest separately. On a chart \(Y_\alpha\), a coordinate function \(t:Y_\alpha\to\mathbf A^1_S\) lets us view \(X\cap Y_\alpha\) over the new base \(\mathbf A^1_S\), over which its generic fibre has one dimension less. Hence the induction hypothesis in Theorem \hyperref[bp563]{6.17} gives constructibility and base change over a dense open of \(\mathbf A^1_S\). The preimages of these opens, however, do not cover \(Y\). Take for instance \(Y=\mathbf P^1_S\) with the two standard charts: there a point is missed only if each of the coordinates \(x\) and \(1/x\) defined there takes an excluded value. After shrinking \(S\), the leftover set is therefore a closed subset finite over \(S\). Lemma \hyperref[bp553]{6.13} shows that the same holds in general.

On the leftover finite part we cannot argue by dimension again. We do not have to, since a finite direct image is exact and conservative. This means that on a subset finite over \(S\), the constructibility of a sheaf can be read off from its direct image to \(S\). Lemma \hyperref[bp550]{6.12} turns this into a criterion. If a complex is already constructible away from such a subset, a proper direct image to \(S\) detects whether it is constructible everywhere. In the smooth case this direct image is \(R(bj)_*F\), which Lemma \hyperref[bp545]{6.10} controls. For this reason the induction passes through the smooth case first.

% object: lemma-detection-proper-direct-image (blueprint/chap-Constructibility-of-f/sec-Coordinate-projections/lemma-detection-proper-direct-image/_section.md)
\phantomsection\label{bp550}%
\begin{bpstatementit}{Lemma~6.12}{ (Detection by a proper direct image)}
Let \(b:Y\to S\) be proper, let \(j:U\hookrightarrow Y\) be an open immersion with reduced closed complement \(i:Z\hookrightarrow Y\), and suppose that \(bi:Z\to S\) is finite.

\begin{enumerate}
\item If \(K\in D^+(Y_{\acute{e}t},\Fell)\), \(j^*K\in\Dbc(U,\Fell)\), and \(Rb_*K\in\Dbc(S,\Fell)\), then \(K\in\Dbc(Y,\Fell)\).
\item If \(\varphi:K\to K'\) is a morphism in \(D^+(Y_{\acute{e}t},\Fell)\) such that \(j^*\varphi\) and \(Rb_*\varphi\) are isomorphisms, then \(\varphi\) is an isomorphism.
\end{enumerate}
\end{bpstatementit}

% proof: proof/lemma-detection-proper-direct-image
\begin{proof}[{Proof}]\phantomsection\label{bp551}%
Using the localization triangle we split \(K\) into a part on \(U\), which is known, and a part on \(Z\). Since \(Z\) is finite over \(S\), the part on \(Z\) is detected by its direct image to \(S\). We compute this direct image from \(Rb_*K\) and the known part.

We work on the étale sites, in \(D^+(-_{\acute{e}t},\Fell)\) and \(\operatorname{Mod}(-_{\acute{e}t},\Fell)\).

\emph{Finite direct images.} Let \(h:Z\to S\) be finite. Then \(h_*\) is exact and conservative, and \(Rh_*M\in\Dbc\) implies \(M\in\Dbc\). The first two properties hold because, for a sheaf \(F\) of \(\Fell\)-modules on \(Z\), \(h_*F\) has stalks

\[
(h_*F)_{\bar s}
=\prod_{\bar z\in |Z_{\bar s}|}F_{\bar z}
\tag{6.17}
\]

by [Stacks, \href{https://stacks.math.columbia.edu/tag/03QP}{03QP}]. Suppose \(h_*F\) is constructible. The counit \(h^*h_*F\to F\) is surjective, since on stalks it is a projection from the finite product (6.17). Its source is constructible, and a quotient of a constructible sheaf is constructible [Stacks, \href{https://stacks.math.columbia.edu/tag/09BH}{09BH}]. Hence \(F\) is constructible. For a complex \(M\) we have \(H^n(Rh_*M)=h_*H^n(M)\) because \(h_*\) is exact, so the statement for complexes follows from conservativity and the sheaf case in each degree.

\emph{The localization triangle.} The exact sequence for an open and its closed complement [Stacks, \href{https://stacks.math.columbia.edu/tag/095L}{095L}], applied termwise to \(K\), gives a triangle

\[
j_!j^*K\longrightarrow K\longrightarrow i_*i^*K
\longrightarrow j_!j^*K[1].
\tag{6.18}
\]

\emph{Part (1).} The first term \(j_!j^*K\) is bounded constructible, because \(j^*K\) is and extension by zero preserves constructibility. Its direct image \(Rb_*j_!j^*K\) is bounded constructible by the finiteness theorem [Stacks, \href{https://stacks.math.columbia.edu/tag/0GL0}{0GL0}, \href{https://stacks.math.columbia.edu/tag/0G2A}{0G2A}], which applies since \(Rb_*=Rb_!\) for proper \(b\). Applying \(Rb_*\) to (6.18), we obtain a triangle whose first two terms, \(Rb_*j_!j^*K\) and \(Rb_*K\), are bounded constructible. Hence the third term \(Rb_*i_*i^*K=R(bi)_*i^*K\) is bounded constructible too. The first step, applied to \(h=bi\), now gives \(i^*K\in\Dbc(Z,\Fell)\). As \(i_*\) is exact and preserves constructibility, (6.18) shows that \(K\in\Dbc(Y,\Fell)\).

\emph{Part (2).} Let \(M\) be a cone of \(\varphi\); we must show \(M=0\). We have \(j^*M=0\) because \(j^*\varphi\) is an isomorphism. By (6.18), this means \(M=i_*i^*M\). Because \(Rb_*\varphi\) is an isomorphism, we also have \(0=Rb_*M=R(bi)_*i^*M\). Since \((bi)_*\) is exact, the cohomology sheaves of \(R(bi)_*i^*M\) are the sheaves \((bi)_*H^n(i^*M)\). Conservativity of \((bi)_*\) then shows that every \(H^n(i^*M)\) vanishes, and so \(M=i_*i^*M=0\).
\end{proof}

% text: prose-06-14
\phantomsection\label{bp552}%
The next lemma produces such a finite complement. Projections to coordinate lines lower the dimension of the generic fibre, and the subset that they do not cover becomes finite over a smaller base.

% object: lemma-coordinate-projections (blueprint/chap-Constructibility-of-f/sec-Coordinate-projections/lemma-coordinate-projections/_section.md)
\phantomsection\label{bp553}%
\begin{bpstatementit}{Lemma~6.13}{ (Coordinate projections)}
Let \(S\) be an integral Noetherian affine scheme with function field \(k(S)\), and let \(Y\to S\) be projective. Choose a finite affine open cover \(Y=\bigcup_\alpha Y_\alpha\) and closed immersions \(Y_\alpha\hookrightarrow\mathbf A^{m_\alpha}_S\) with \(m_\alpha\geq1\), and write \(t_{\alpha r}:Y_\alpha\to\mathbf A^1_S\) (\(1\le r\le m_\alpha\)) for the coordinate functions.

\begin{enumerate}
\item If \(X\subset Y\) is open and \(\dim X_{k(S)}\leq n\), then for every \(\alpha,r\) the generic fibre of \(t_{\alpha r}:X\cap Y_\alpha\to\mathbf A^1_S\) has dimension at most \(n-1\).
\item Suppose given a dense open \(V_{\alpha r}\subset\mathbf A^1_S\) for each \(\alpha,r\), and let \[
   G=\bigcup_{\alpha,r}
   t_{\alpha r}^{-1}(V_{\alpha r})\subset Y.
   \] After shrinking \(S\), the reduced closed complement \(Z=Y\setminus G\) is finite over \(S\).
\end{enumerate}
\end{bpstatementit}

% proof: proof/lemma-coordinate-projections
\begin{proof}[{Proof}]\phantomsection\label{bp554}%
Part (1) is a transcendence degree count. For part (2), note that a point of \(Z\) has all its coordinates in the finite sets of excluded values. After shrinking \(S\), these sets are cut out by monic polynomials, and monic equations in every coordinate force finiteness.

(1) We use that the dimension of an integral scheme of finite type over a field is the transcendence degree of its function field. Consider an irreducible component of \(X_{k(S)}\cap Y_\alpha\) with function field \(L\). If \(t_{\alpha r}\) is algebraic over \(k(S)\) on it, the component does not meet the generic fibre of \(t_{\alpha r}\), as it does not dominate \(\mathbf A^1_{k(S)}\). Otherwise \(t_{\alpha r}\) is transcendental, and

\[
\operatorname{trdeg}_{k(S)(t_{\alpha r})}L
=\operatorname{trdeg}_{k(S)} L-1\leq n-1.
\]

The left-hand side is the dimension of the corresponding component of the generic fibre of \(t_{\alpha r}\), and every component of that generic fibre arises in this way.

(2) Let \(t\) be the coordinate of \(\mathbf A^1_S\). The ideal of the complement of \(V_{\alpha r}\) contains a nonzero polynomial \(P_{\alpha r}\in\mathcal O(S)[t]\). Indeed, since \(V_{\alpha r}\) is dense, its complement is a proper closed subset of the integral scheme \(\mathbf A^1_S\), and its ideal is nonzero. After inverting the leading coefficients of these finitely many polynomials, we may assume that all of them are monic. The function \(P_{\alpha r}(t_{\alpha r})\) is zero in the coordinate ring of the reduced affine scheme \(Z\cap Y_\alpha\), because it vanishes at every point, \(t_{\alpha r}\) mapping \(Z\cap Y_\alpha\) into the complement of \(V_{\alpha r}\). This ring is therefore a quotient of

\[
\mathcal O(S)[t_{\alpha1},\ldots,t_{\alpha m_\alpha}]/
(P_{\alpha1}(t_{\alpha1}),\ldots,
P_{\alpha m_\alpha}(t_{\alpha m_\alpha})),
\]

which is generated as an \(\mathcal O(S)\)-module by the monomials whose exponent in each variable is less than the degree of the corresponding polynomial. Each \(Z\cap Y_\alpha\) is thus finite over \(S\), and \(Z\to S\) has finite fibres. As \(Z\) is closed in the projective scheme \(Y\), the morphism \(Z\to S\) is also proper, and a proper morphism with finite fibres is finite [Stacks, \href{https://stacks.math.columbia.edu/tag/02LS}{02LS}].
\end{proof}


% object: sec-Generic-base-change (blueprint/chap-Constructibility-of-f/sec-Generic-base-change/_section.md), content hidden
\subsection*{\texorpdfstring{6.5\quad Generic base change}{6.5 Generic base change}}
\phantomsection\label{bp555}%
\addcontentsline{toc}{subsection}{\texorpdfstring{6.5\quad Generic base change}{6.5 Generic base change}}
% text: prose-06-15
\phantomsection\label{bp556}%
Let \(j:X\hookrightarrow Y\) be an open immersion of schemes over an integral base \(S\), and let \(F\) be a constructible \(\Fell\)-sheaf on \(X\). For such data, Deligne's generic base change theorem says that over a dense open of \(S\) the complex \(Rj_*F\) is bounded constructible and its formation commutes with base change. In particular, over a field it gives the finiteness of \(Rj_*\) (Corollary \hyperref[bp566]{6.18}), and this is all that the computation of \(f^!\) requires. On the other hand, the punctured plane over the \(t\)-line, discussed at the beginning of this chapter, shows that we cannot ask for more than a dense open.

In characteristic \(p\) the generic fibre need not contain a smooth dense open. For example, let \(X=\mathbf A^1_{\mathbf F_p}\) map to the \(t\)-line by \(x\mapsto x^p\). The generic fibre is \(\operatorname{Spec}\mathbf F_p(t^{1/p})\), which is not smooth over \(\mathbf F_p(t)\). For this reason Step 3 of the proof of Theorem \hyperref[bp563]{6.17} below allows a finite universal homeomorphism of the base. Such a map changes neither the étale topoi nor the conclusion.

We fix notation for base change. With \(j\) as above and a morphism \(T\to S\), we write \(j_T:X_T\hookrightarrow Y_T\) for the base change and \(g_X:X_T\to X\), \(g_Y:Y_T\to Y\) for the projections. The base change map \(g_Y^*Rj_*K\to Rj_{T,*}g_X^*K\) is then the map (4.5).

The proof of Theorem \hyperref[bp563]{6.17} is an induction on the dimension of the generic fibre, in three steps. The next three lemmas are the parts of this proof that do not use the induction hypothesis, and each serves one step. First, Lemma \hyperref[bp557]{6.14} is Step 2: if \(X\) is smooth and \(F\) locally constant, it extends the conclusion from an open \(G\) with complement finite over \(S\), provided by Step 1, across that complement. Next, Lemma \hyperref[bp559]{6.15} allows Step 3 to replace \(S\) by a finite universal homeomorphism, as the example above requires. Lemma \hyperref[bp561]{6.16} then provides, after this replacement, the dense open \(V\) of Step 3, smooth with \(F\) locally constant, whose complement has a smaller generic fibre.

% object: lemma-extension-across-finite-subset (blueprint/chap-Constructibility-of-f/sec-Generic-base-change/lemma-extension-across-finite-subset/_section.md)
\phantomsection\label{bp557}%
\begin{bpstatementit}{Lemma~6.14}{ (Extension across a finite subset)}
Let \(S\) be integral and affine, let \(b:Y\to S\) be projective, and let \(j:X\hookrightarrow Y\) be an open immersion with \(X\) smooth over \(S\). Let \(F\) be a finite locally constant \(\Fell\)-sheaf on \(X_{\acute{e}t}\). Suppose that there is an open \(G\subset Y\) whose reduced complement \(i:Z\hookrightarrow Y\) is finite over \(S\), such that \((Rj_*F)|_G\in\Dbc(G,\Fell)\) and such that for every \(k\)-scheme \(T\to S\) of finite type the base change map \(g_Y^*Rj_*F\to Rj_{T,*}g_X^*F\) is an isomorphism on \(G_T\). Then there is a nonempty open \(S_0\subset S\) such that \(Rj_{S_0,*}(F|_{X_{S_0}})\in\Dbc(Y_{S_0},\Fell)\) and such that for every \(k\)-scheme \(T\to S_0\) of finite type the base change map \(g_Y^*Rj_{S_0,*}(F|_{X_{S_0}})\to Rj_{T,*}g_X^*F\) is an isomorphism.
\end{bpstatementit}

% proof: proof/lemma-extension-across-finite-subset
\begin{proof}[{Proof}]\phantomsection\label{bp558}%
Lemma \hyperref[bp545]{6.10} controls the direct image \(R(bj)_*F\) of \(Rj_*F\) to \(S\). Then Lemma \hyperref[bp550]{6.12} turns this, together with the hypothesis on \(G\), into control of \(Rj_*F\) along \(Z\). Throughout, we work on the étale sites.

For constructibility we check the hypotheses of part (1) of Lemma \hyperref[bp550]{6.12} for \(K=Rj_*F\) and \(U=G\). The complex \(Rj_*F\) lies in \(D^+(Y_{\acute{e}t},\Fell)\) because \(F\) is a sheaf. Its restriction to \(G\) is in \(\Dbc\) by assumption, and \(Z=Y\setminus G\) is finite over \(S\). It remains to treat the direct image. Since \(bj:X\to S\) is smooth and \(F\) is locally constant, we may apply Lemma \hyperref[bp545]{6.10}. After shrinking \(S\) to a nonempty open \(S_0\), the complex \(Rb_*Rj_*F=R(bj)_*F\) lies in \(\Dbc(S_0,\Fell)\), and its formation commutes with base change along every \(T\to S_0\) of finite type over \(k\). Part (1) of Lemma \hyperref[bp550]{6.12} then gives \(Rj_*F\in\Dbc(Y_{S_0},\Fell)\).

For base change we use part (2) of the lemma instead. Let \(T\to S_0\) be a \(k\)-scheme of finite type, and let \(\varphi\) be the base change map \(g_Y^*Rj_*F\to Rj_{T,*}g_X^*F\). By assumption it is an isomorphism on \(G_T\). Proper base change for \(b\) and the compatibility of base change maps with composition [Stacks, \href{https://stacks.math.columbia.edu/tag/0E46}{0E46}] identify \(Rb_{T,*}\varphi\) with the base change map for \(R(bj)_*F\), which is an isomorphism. Since \(b_T\) is proper and \(b_Ti_T\) is finite, part (2) of Lemma \hyperref[bp550]{6.12} applies over \(T\). It shows that \(\varphi\) is an isomorphism.
\end{proof}

% object: lemma-universal-homeomorphism-base (blueprint/chap-Constructibility-of-f/sec-Generic-base-change/lemma-universal-homeomorphism-base/_section.md)
\phantomsection\label{bp559}%
\begin{bpstatementit}{Lemma~6.15}{ (Universal homeomorphisms of the base)}
Let \(S\) be integral, let \(j:X\hookrightarrow Y\) be an open immersion of \(S\)-schemes, and let \(F\) be a constructible \(\Fell\)-sheaf on \(X_{\acute{e}t}\). Let \(S'\to S\) be a finite universal homeomorphism with \(S'\) integral, let \(X'=(X_{S'})_{\mathrm{red}}\) and \(Y'=(Y_{S'})_{\mathrm{red}}\), let \(j':X'\hookrightarrow Y'\) be the open immersion, and let \(F'\) be the pullback of \(F\) to \(X'\). Suppose that there is a nonempty open \(U'\subset S'\) such that \(Rj'_{U',*}(F'|_{X'_{U'}})\in\Dbc(Y'_{U'},\Fell)\) and such that for every \(k\)-scheme \(T'\to U'\) of finite type the base change map for \(j'\) along \(T'\) is an isomorphism. Then there is a nonempty open \(U\subset S\) with the same two properties for \(j\) and \(F\).
\end{bpstatementit}

% proof: proof/lemma-universal-homeomorphism-base
\begin{proof}[{Proof}]\phantomsection\label{bp560}%
A finite universal homeomorphism is a homeomorphism of underlying spaces, so we may remove from \(S\) the closed image of \(S'\setminus U'\). The remaining open \(U\) is nonempty, and its inverse image \(S'_U\) lies in \(U'\). Over \(U\) the morphisms \(X'_U\to X_U\) and \(Y'_U\to Y_U\) are finite universal homeomorphisms, and pullback along them is therefore an equivalence of étale topoi [Stacks, \href{https://stacks.math.columbia.edu/tag/03SI}{03SI}]. Such equivalences preserve constructibility in both directions and commute with \(Rj_*\). Hence pullback along these morphisms detects boundedness and constructibility, and we get \(Rj_{U,*}(F|_{X_U})\in\Dbc(Y_U,\Fell)\).

For base change, let \(T\to U\) be a \(k\)-scheme of finite type, with projections \(g_X:X_T\to X\) and \(g_Y:Y_T\to Y\), and let

\[
\beta_T:g_Y^*Rj_*F\longrightarrow Rj_{T,*}g_X^*F
\]

be the base change map for \(j\) along \(T\). Put \(T'=T\times_US'_U\). Then \(v:T'\to T\) is again a universal homeomorphism, and so are the induced maps \(v_X:X'_{T'}\to X_T\) and \(v_Y:Y'_{T'}\to Y_T\), each a thickening followed by a base change of \(v\). Let \(\beta'_{T'}:g_{Y'}^*Rj'_*F'\to Rj'_{T',*}g_{X'}^*F'\) be the base change map for \(j'\) along \(T'\). Base change maps compose [Stacks, \href{https://stacks.math.columbia.edu/tag/0EZR}{0EZR}, \href{https://stacks.math.columbia.edu/tag/0EZS}{0EZS}], and \(v_Y^*Rj_{T,*}\simeq Rj'_{T',*}v_X^*\) because \(v_X^*\) and \(v_Y^*\) are equivalences of étale topoi [Stacks, \href{https://stacks.math.columbia.edu/tag/03SI}{03SI}]. Hence \(v_Y^*\beta_T\) is identified with \(\beta'_{T'}\). The map \(\beta'_{T'}\) is an isomorphism by assumption, so \(v_Y^*\beta_T\) is one. Since \(v_Y^*\) is an equivalence, it is conservative, and therefore \(\beta_T\) is an isomorphism.
\end{proof}

% object: lemma-generic-smooth-open (blueprint/chap-Constructibility-of-f/sec-Generic-base-change/lemma-generic-smooth-open/_section.md)
\phantomsection\label{bp561}%
\begin{bpstatementit}{Lemma~6.16}{ (A smooth dense open after a universal homeomorphism)}
Let \(S\) be integral, let \(X\) be an \(S\)-scheme all of whose irreducible components dominate \(S\), let \(n=\dim X_{k(S)}\), and let \(F\) be a constructible \(\Fell\)-sheaf on \(X_{\acute{e}t}\). After shrinking \(S\) to a nonempty open, there are a finite universal homeomorphism \(S'\to S\) with \(S'\) integral and a dense open \(V\subset(X_{S'})_{\mathrm{red}}\) such that \(V\) is smooth over \(S'\), the pullback of \(F\) to \(V\) is locally constant, and the generic fibre of \((X_{S'})_{\mathrm{red}}\setminus V\) over \(S'\) has dimension less than \(n\).
\end{bpstatementit}

% proof: proof/lemma-generic-smooth-open
\begin{proof}[{Proof}]\phantomsection\label{bp562}%
By [Stacks, \href{https://stacks.math.columbia.edu/tag/0550}{0550}], after shrinking \(S\), there is a finite universal homeomorphism \(S'\to S\) with \(S'\) integral, such that \((X_{S'})_{\mathrm{red}}\to S'\) is flat of finite presentation with geometrically reduced generic fibre. This generic fibre has a dense smooth open [Stacks, \href{https://stacks.math.columbia.edu/tag/056V}{056V}]. By the fibrewise criterion [Stacks, \href{https://stacks.math.columbia.edu/tag/01V9}{01V9}], its generic points therefore lie in the smooth locus of \((X_{S'})_{\mathrm{red}}\to S'\). After shrinking the base once more to remove non-dominant components, this smooth locus contains a dense open of \((X_{S'})_{\mathrm{red}}\). Since \(F\) is constructible, there is a smaller dense open \(V\) on which the pullback of \(F\) is locally constant. Its complement contains no generic point of a component of the generic fibre. Hence the generic fibre of the complement has dimension less than \(n\).
\end{proof}

% object: theorem-generic-constructibility-base-change-rj (blueprint/chap-Constructibility-of-f/sec-Generic-base-change/theorem-generic-constructibility-base-change-rj/_section.md)
\phantomsection\label{bp563}%
\begin{bpstatementit}{Theorem~6.17}{ (Generic constructibility and base change for $Rj_*$)}
Let \(S\) be integral, and let \(j:X\hookrightarrow Y\) be an open immersion of \(S\)-schemes. For a constructible \(\Fell\)-sheaf \(F\) on \(X_{\acute{e}t}\), there is a nonempty open \(S_0\subset S\) such that

\[
Rj_{S_0,*}(F|_{X_{S_0}})\in\Dbc(Y_{S_0},\Fell),
\tag{6.19}
\]

and such that for every \(k\)-scheme \(T\to S_0\) of finite type, with base change \(j_T:X_T\hookrightarrow Y_T\) and projections \(g_X:X_T\to X_{S_0}\) and \(g_Y:Y_T\to Y_{S_0}\), the base change map

\[
\beta_T:g_Y^*Rj_{S_0,*}(F|_{X_{S_0}})
\longrightarrow Rj_{T,*}g_X^*(F|_{X_{S_0}})
\tag{6.20}
\]

is an isomorphism.
\end{bpstatementit}

% proof: proof/theorem-generic-constructibility-base-change-rj
\begin{proof}[{Proof}]\phantomsection\label{bp564}%
We work with the étale sheaf corresponding to \(F\), in \(D^+(-_{\acute{e}t},\Fell)\), and prove (6.19) and \(\beta_T\) together by induction on \(n=\dim X_{k(S)}\), allowing at each stage all integral bases \(S\) of finite type over \(k\). We first handle the complement of a finite set by passing to the base \(\mathbf A^1_S\), then the smooth, locally constant case, and finally the general case by dévissage.

\emph{Preliminary reductions.} Several modifications of the data are harmless, and we will make them freely in the induction. First, the assertions extend from open to locally closed immersions with the same source. Indeed, such an immersion factors as an open immersion into its scheme-theoretic closure followed by a closed immersion. Since closed direct image is exact, preserves constructibility, and commutes with arbitrary pullback [Stacks, \href{https://stacks.math.columbia.edu/tag/0959}{0959}], the assertions pass along this factorization. We will use the induction hypothesis in this form. Next, replacing \(Y\) by the scheme-theoretic closure of \(X\), we may assume that \(X\) is dense in \(Y\). We may also shrink \(S\) to be affine, and we may remove the closures of the images of the irreducible components of \(Y\) that do not dominate \(S\), since there are finitely many of them and none contains the generic point of \(S\). Finally, the assertions are local on \(Y\). Since we only use finite affine covers, the finitely many resulting opens of \(S\) can be intersected.

\emph{Base case.} If the generic fibre is empty or zero-dimensional, then after shrinking \(S\) the immersion \(j\) becomes an isomorphism or \(Y\) becomes empty. Suppose first that \(X_{k(S)}\) is empty. By density \(Y_{k(S)}\) is empty. On a finite affine cover, a coordinate algebra that vanishes over \(k(S)\) has its unit killed by a nonzero element of \(\mathcal O(S)\). Hence \(Y\) becomes empty after shrinking \(S\). If \(n=0\), the dense open \(X_{k(S)}\) of the zero-dimensional Noetherian space \(Y_{k(S)}\) is all of \(Y_{k(S)}\). The same argument applied to \(Y\setminus X\) shows that \(j\) becomes an isomorphism after shrinking \(S\).

Now let \(n\geq1\), and assume the theorem, including \(\beta_T\), for generic fibres of dimension less than \(n\).

\emph{Step 1: for \(Y\) projective over \(S\), both assertions hold on an open \(G\subset Y\) whose complement is finite over \(S\).} We replace the base \(S\) by the larger base \(\mathbf A^1_S\), over which the generic fibre is smaller, so that the induction hypothesis applies. To do this, choose charts and coordinates as in Lemma \hyperref[bp553]{6.13}. Via \(t_{\alpha r}\), we regard \(X\cap Y_\alpha\hookrightarrow Y_\alpha\) as an open immersion of schemes over the integral base \(\mathbf A^1_S\), which is of finite type over \(k\). By part (1) of that lemma, its generic fibre has dimension less than \(n\). Therefore the induction hypothesis, applied over the base \(\mathbf A^1_S\) to the restriction of \(F\), gives a dense open \(V_{\alpha r}\subset\mathbf A^1_S\) over which both assertions hold.

Since \(Rj_*\) commutes with restriction to the open \(Y_\alpha\), the assertions hold for \(Rj_*F\) on \(t_{\alpha r}^{-1}(V_{\alpha r})\). They continue to hold after a base change \(T\to S\) with \(T\) of finite type over \(k\), as we see by applying the hypothesis to the inverse image of \(V_{\alpha r}\) in \(\mathbf A^1_T\). Let \(G\) be the finite union of the opens \(t_{\alpha r}^{-1}(V_{\alpha r})\). Then on \(G\) the map \(\beta_T\) is an isomorphism, and the complex \(Rj_*F\) is bounded constructible. For the boundedness, note that there are finitely many pieces, and therefore the cohomology lies in a common finite range of degrees. Finally, part (2) of the lemma shows that, after shrinking \(S\), the reduced complement \(Z=Y\setminus G\) is finite over \(S\).

\emph{Step 2: the theorem holds in dimension \(n\) when \(X\) is smooth over \(S\) and \(F\) is locally constant.} We may assume that \(Y\) is projective over \(S\). Indeed, let \(Y_\beta\) be a member of a finite affine cover of \(Y\), embed it in some \(\mathbf A^m_S\) over the affine base, and let \(\overline Y\) be the closure of \(X\cap Y_\beta\) in \(\mathbf P^m_S\). Then \(X\cap Y_\beta\) is open and dense in \(\overline Y\), which is projective over \(S\). Moreover \(X\cap Y_\beta\) is still smooth over \(S\), with generic fibre of dimension at most \(n\). The result for \(X\cap Y_\beta\subset\overline Y\) gives it on \(Y_\beta\), because \(Rj_*\) and its base change maps commute with restriction to opens of the target.

For \(Y\) projective, Step 1 gives an open \(G\subset Y\) with complement finite over \(S\) on which both assertions hold. Since \(X\) is smooth over \(S\) and \(F\) is locally constant, Lemma \hyperref[bp557]{6.14} applies. It shows that the assertions hold on all of \(Y\) after shrinking \(S\).

\emph{Step 3: the general case in dimension \(n\).} We find a dense open \(V\subset X\) to which Step 2 applies, and treat the difference between \(F\) and \(Rv_*v^*F\), which lives on the smaller set \(X\setminus V\), by the induction hypothesis. Let \(S'\to S\) and \(V\) be as in Lemma \hyperref[bp561]{6.16}. Lemma \hyperref[bp559]{6.15} shows that it suffices to prove the theorem for \((X_{S'})_{\mathrm{red}}\hookrightarrow(Y_{S'})_{\mathrm{red}}\) and the pullback of \(F\) over \(S'\). Its generic fibre still has dimension \(n\) because \(S'\to S\) is a universal homeomorphism. We therefore replace \(S\) by \(S'\) and all schemes by these reductions. Then \(V\subset X\) is a dense open, smooth over \(S\), on which \(F\) is locally constant. The generic fibre of \(X\setminus V\) has dimension less than \(n\).

Let \(v:V\hookrightarrow X\) be the open immersion and \(i:X\setminus V\hookrightarrow X\) its reduced closed complement. By Step 2, after shrinking the base, both \(v\) and \(jv\) have bounded constructible direct images of \(v^*F\). Their base change maps are isomorphisms after every base change \(T\to S\) with \(T\) of finite type over \(k\). Complete the unit of \(v^*\dashv Rv_*\) to a triangle

\[
\Delta_F:\ F\longrightarrow Rv_*v^*F\longrightarrow D
\longrightarrow F[1].
\tag{6.21}
\]

The first term is the constructible sheaf \(F\), and the second is bounded constructible by Step 2. Hence \(D\) is bounded constructible too. Indeed, by the long exact cohomology sequence, its cohomology sheaves are extensions of kernels by cokernels of maps between constructible sheaves, and these are constructible [Stacks, \href{https://stacks.math.columbia.edu/tag/09BH}{09BH}]. Moreover \(v^*D=0\), as the unit is an isomorphism on \(V\). The localization triangle (6.18) for \(v\) and \(i\) now shows that

\[
D=i_*i^*D.
\]

The locally closed immersion \(ji:X\setminus V\to Y\) has generic fibre of dimension less than \(n\). We apply the induction hypothesis to each of the finitely many constructible cohomology sheaves of \(i^*D\), together with their truncation triangles. This gives

\[
R(ji)_*i^*D=Rj_*D\in\Dbc(Y,\Fell),
\]

with base change maps that are isomorphisms. Now apply \(Rj_*\) to the triangle \(\Delta_F\). In the resulting triangle, the middle term \(R(jv)_*v^*F\) is bounded constructible by Step 2, and the third term \(Rj_*D\) was treated just now. Hence \(Rj_*F\) is bounded constructible.

To see that the base change map \(\beta_T\) is an isomorphism, we compare the triangle \(\Delta_F\) with its analogue over \(T\). Let \(g_X:X_T\to X\), \(g_Y:Y_T\to Y\) and \(g_V:V_T\to V\) be the projections, and \(v_T:V_T\to X_T\) the base change of \(v\). The base change map

\[
\beta_v:g_X^*Rv_*v^*F\longrightarrow Rv_{T,*}g_V^*v^*F
\]

is an isomorphism by Step 2, and it carries \(g_X^*\) of the first arrow of the triangle \(\Delta_F\) to the unit \(g_X^*F\to Rv_{T,*}v_T^*g_X^*F\) of \(v_T^*\dashv Rv_{T,*}\). Since cones are not functorial, we have to choose the triangles compatibly, and we take as triangle for this unit the pullback of the triangle \(\Delta_F\), transported through \(\beta_v\). Now the base change map

\[
\beta_j:g_Y^*Rj_*\longrightarrow Rj_{T,*}g_X^*
\]

is a natural transformation of exact functors. It therefore induces a morphism from \(g_Y^*\) of \(Rj_*\) of the triangle \(\Delta_F\) to \(Rj_{T,*}\) of this triangle. On the second term it is \(\beta_j\) at \(Rv_*v^*F\) followed by \(Rj_{T,*}\beta_v\). Since base change maps compose [Stacks, \href{https://stacks.math.columbia.edu/tag/0EZR}{0EZR}, \href{https://stacks.math.columbia.edu/tag/0EZS}{0EZS}], this composite is the base change map for \(jv\), an isomorphism by Step 2. On the third term it is \(\beta_j\) at \(D\). It is an isomorphism because, composed with \(Rj_{T,*}\) of the base change isomorphism for the closed immersion \(i\) [Stacks, \href{https://stacks.math.columbia.edu/tag/0959}{0959}], it is the base change map for \(ji\) at \(i^*D\) treated above. Hence, by the five lemma for triangles, it is an isomorphism on the first term, which is \(\beta_T\).

This completes the induction.
\end{proof}

% text: prose-06-16
\phantomsection\label{bp565}%
The only nonempty open of \(\operatorname{Spec}k\) is \(\operatorname{Spec}k\) itself, and hence over a field the theorem holds without any shrinking.

% object: corollary-open-direct-image-over-field (blueprint/chap-Constructibility-of-f/sec-Generic-base-change/corollary-open-direct-image-over-field/_section.md)
\phantomsection\label{bp566}%
\begin{bpstatementit}{Corollary~6.18}{ (Open direct image over a field)}
Let \(j:X\hookrightarrow Y\) be an open immersion. Then

\[
Rj_*:\Dbc(X,A)\longrightarrow\Dbc(Y,A),
\tag{6.22}
\]

and for every \(k\)-scheme \(T\) of finite type, with projections \(g_X:X_T\to X\) and \(g_Y:Y_T\to Y\), and every \(K\in\Dbc(X,A)\), the base change map \(g_Y^*Rj_*K\to Rj_{T,*}g_X^*K\) is an isomorphism.
\end{bpstatementit}

% proof: proof/corollary-open-direct-image-over-field
\begin{proof}[{Proof}]\phantomsection\label{bp567}%
We argue on the étale site. The theorem gives the case of \(\Fell\)-sheaves, and we pass to \(\mathbf Z/\ell^r\)-sheaves by filtrations, to \(A\)-sheaves by forgetting the \(A\)-structure, and to complexes by truncation.

For a constructible \(\Fell\)-sheaf this is Theorem \hyperref[bp563]{6.17} with \(S=\operatorname{Spec}k\). A constructible sheaf \(F\) killed by \(\ell^r\) has the finite filtration \(F\supset\ell F\supset\cdots\supset\ell^rF=0\). Its successive quotients \(\ell^iF/\ell^{i+1}F\) are constructible and killed by \(\ell\), that is, constructible \(\Fell\)-sheaves. Each short exact sequence in the filtration is a triangle in the derived category, and \(Rj_*\) carries it to a triangle. By induction on \(r\), the long exact cohomology sequences of these triangles give (6.22). The base change maps form morphisms between the resulting triangles, and they are isomorphisms by the five lemma.

For a finite \(\ell\)-primary ring \(A\), we apply this to the underlying \(\mathbf Z/\ell^r\mathbf Z\)-sheaf and carry the \(A\)-action along all maps. This is possible for two reasons. First, \(Rj_*:D^+(X_{\acute{e}t},A)\to D^+(Y_{\acute{e}t},A)\) commutes with the forgetful functors to \(D^+(-_{\acute{e}t},\mathbf Z/\ell^r\mathbf Z)\) [Stacks, \href{https://stacks.math.columbia.edu/tag/0D6J}{0D6J}]. Second, an \(A\)-sheaf is constructible if and only if its underlying sheaf is. To see this, consider a stratum where the underlying sheaf is finite locally constant. There the action of each of the finitely many elements of \(A\) is a section of the finite locally constant sheaf of endomorphisms, hence constant étale locally. Therefore the \(A\)-module is locally constant on that stratum. Truncation triangles extend both conclusions from sheaves to bounded constructible complexes.
\end{proof}


% object: sec-Finite-coefficients-2 (blueprint/chap-Constructibility-of-f/sec-Finite-coefficients-2/_section.md), content hidden
\subsection*{\texorpdfstring{6.6\quad Finite coefficients}{6.6 Finite coefficients}}
\phantomsection\label{bp568}%
\addcontentsline{toc}{subsection}{\texorpdfstring{6.6\quad Finite coefficients}{6.6 Finite coefficients}}
% text: prose-06-17
\phantomsection\label{bp569}%
Let \(A\) be a finite commutative \(\ell\)-primary ring and \(f\) a separated morphism of finite type. We show that the right adjoint \(f^!\) of Theorem \hyperref[bp506]{5.31} preserves \(D^b_{\mathrm{ctf}}(-,A)\) (Theorem \hyperref[bp573]{6.20}). By the reduction at the start of this chapter and by the previous section, we only have to show that direct image along an open immersion \(j\) preserves finite Tor amplitude. The upper end of the Tor amplitude grows by at most the \(\ell\)-cohomological dimension \(\operatorname{cd}_\ell(j)\) of Definition \hyperref[bp137]{2.45}. This dimension is finite for morphisms of finite type over \(k\).

Finite Tor amplitude cannot be checked on cohomology sheaves, since these may have infinite Tor dimension over \(A\): for \(A=\mathbf Z/\ell^2\), already the constant sheaf \(\mathbf F_\ell\) has infinite Tor dimension. We therefore obtain the Tor bounds below from flat models. By [Stacks, \href{https://stacks.math.columbia.edu/tag/0F4M}{0F4M}], every object of \(D^b_{\mathrm{ctf}}(X,A)\) is represented by a bounded complex of constructible flat \(A\)-sheaves. On the strata of a finite stratification, such a model is locally a bounded complex of finite projective \(A\)-modules.

% object: proposition-open-direct-image-preserves-finite (blueprint/chap-Constructibility-of-f/sec-Finite-coefficients-2/proposition-open-direct-image-preserves-finite/_section.md)
\phantomsection\label{bp570}%
\begin{bpstatementit}{Proposition~6.19}{ (Open direct image preserves finite Tor amplitude)}
Let \(j:U\hookrightarrow X\) be an open immersion. Then \(Rj_*:D(U_{\mathrm{pro}\acute{e}t},A)\to D(X_{\mathrm{pro}\acute{e}t},A)\) restricts to

\[
Rj_*:D^b_{\mathrm{ctf}}(U,A)\longrightarrow D^b_{\mathrm{ctf}}(X,A).
\]

More precisely, if \(K\) has Tor amplitude in \([a,b]\) and \(d\ge
\operatorname{cd}_\ell(j)\), then \(Rj_*K\) has Tor amplitude in \([a,b+d]\).
\end{bpstatementit}

% proof: proof/proposition-open-direct-image-preserves-finite
\begin{proof}[{Proof}]\phantomsection\label{bp571}%
We test Tor amplitude by tensoring with modules. By a projection formula for constant modules we may move the tensor product inside \(Rj_*\). There the complex \(K\otimes^{\mathbf L}\underline M\) has amplitude \([a,b]\), and \(Rj_*\) adds at most \(d\) degrees.

The complex \(Rj_*K\) is bounded constructible by Corollary \hyperref[bp566]{6.18}, and \(\operatorname{cd}_\ell(j)\) is finite by Theorem theorem-results-from-etale-theory. To bound the Tor amplitude, we work on the étale site, where geometric points are conservative. By Corollary \hyperref[bp132]{2.44}, we have \(K\simeq\nu_U^*K_{\acute{e}t}\) for some \(K_{\acute{e}t}\in D^b_{\mathrm{ctf}}(U_{\acute{e}t},A)\). By Proposition \hyperref[bp146]{2.49}, the pro-étale pullback of \(Rj_{\acute{e}t,*}K_{\acute{e}t}\) is \(Rj_*K\).

We first prove the projection formula for constant coefficients. Namely, for every \(A\)-module \(M\), viewed as a constant sheaf, the projection morphism

\[
(Rj_*K)\otimes_A^{\mathbf L}\underline M_X
\xrightarrow{\ \sim\ }
Rj_*\bigl(K\otimes_A^{\mathbf L}\underline M_U\bigr)
\tag{6.23}
\]

is an isomorphism. Recall that for a scheme \(V\) with \(\operatorname{cd}_\ell(V)<\infty\), the projection formula of [Stacks, \href{https://stacks.math.columbia.edu/tag/0F12}{0F12}] identifies \(R\Gamma(V,K)\otimes^{\mathbf L}_AM\) with \(R\Gamma(V,K\otimes^{\mathbf L}_A\underline M)\). Moreover, a map of complexes of sheaves is an isomorphism if it induces isomorphisms on derived sections over a basis of the site. Since the sections of \(Rj_*\) over \(V\) are the sections over \(V\times_XU\), we apply the projection formula to every quasi-compact étale \(V\to X\) and to \(V\times_XU\). This gives isomorphisms between the derived sections of the two sides of the étale version of (6.23), compatibly with the projection morphism. Hence the étale version holds. Finally, inverse image is exact and symmetric monoidal, and it commutes with unbounded direct image, so the pullback of the étale isomorphism is (6.23).

On the étale site it suffices to bound Tor against constant \(A\)-modules. Indeed, let \(G\) be an étale sheaf of \(A\)-modules and \(\bar x\) a geometric point. The stalk of \(Rj_{\acute{e}t,*}K_{\acute{e}t}\otimes^{\mathbf L}_AG\) at \(\bar x\) is that of \(Rj_{\acute{e}t,*}K_{\acute{e}t}\otimes^{\mathbf L}_A\underline M\) with \(M=G_{\bar x}\), and geometric stalks are conservative. Now \(K\otimes_A^{\mathbf L}\underline M_U\) has cohomology in \([a,b]\). Hence the étale version of the right-hand side of (6.23) lies in \([a,b+d]\). By the projection formula, so does \(Rj_{\acute{e}t,*}K_{\acute{e}t}\otimes^{\mathbf L}_A\underline M\). In other words, \(Rj_{\acute{e}t,*}K_{\acute{e}t}\) has Tor amplitude in \([a,b+d]\).

To pass to \(X_{\mathrm{pro}\acute{e}t}\), we pull back a bounded flat constructible model of \(Rj_{\acute{e}t,*}K_{\acute{e}t}\) in degrees \([a,b+d]\) [Stacks, \href{https://stacks.math.columbia.edu/tag/0F4M}{0F4M}]. Since pullback preserves flatness, its pullback to \(X_{\mathrm{pro}\acute{e}t}\) is again a complex of flat sheaves in degrees \([a,b+d]\). Hence its tensor product with any sheaf of \(A\)-modules on \(X_{\mathrm{pro}\acute{e}t}\) is computed termwise and lies in degrees \([a,b+d]\). This argument also covers the sheaves on \(X_{\mathrm{pro}\acute{e}t}\) that are not pulled back from \(X_{\acute{e}t}\), for which geometric stalks would not detect the Tor amplitude.
\end{proof}

% text: prose-06-18
\phantomsection\label{bp572}%
For \(j:\mathbf G_m\hookrightarrow\mathbf A^1\) over \(\bar k\) and \(K=A\), the stalk of \(Rj_*A\) at the origin is the cohomology of the punctured strict henselization. It has cohomology \(A\) in degree \(0\) and \(A(-1)\) in degree \(1\), both free. It is therefore isomorphic to \(A\oplus A(-1)[-1]\). The Tor amplitude grows from \([0,0]\) to \([0,1]\), and the bound of the proposition, with \(d=\operatorname{cd}_\ell(j)=1\), is attained.

Combining the proposition with the computation of \(p^!\) for affine space, we obtain the main theorem for finite coefficients.

% object: theorem-f-shriek-finite-coefficients (blueprint/chap-Constructibility-of-f/sec-Finite-coefficients-2/theorem-f-shriek-finite-coefficients/_section.md)
\phantomsection\label{bp573}%
\begin{bpstatementit}{Theorem~6.20}{ ($f^!$ with finite coefficients)}
For every separated morphism \(f:X\to Y\) of finite type, the right adjoint \(f^!:\Loc(D^b_{\mathrm{ctf}}(Y,A))\to\Loc(D^b_{\mathrm{ctf}}(X,A))\) of Theorem \hyperref[bp506]{5.31} maps \(D^b_{\mathrm{ctf}}(Y,A)\) into \(D^b_{\mathrm{ctf}}(X,A)\). Consequently, its adjunction restricts to

\[
Rf_!:D^b_{\mathrm{ctf}}(X,A)\ \rightleftarrows\ D^b_{\mathrm{ctf}}(Y,A):f^!.
\]
\end{bpstatementit}

% proof: proof/theorem-f-shriek-finite-coefficients
\begin{proof}[{Proof}]\phantomsection\label{bp574}%
Locally on \(X\) and \(Y\), the morphism \(f\) is a closed immersion followed by the projection from an affine space. We therefore treat these two cases and then glue.

\emph{Closed immersions.} Let \(i:Z\hookrightarrow X\) be a closed immersion with open complement \(j:U\hookrightarrow X\), and let \(R\underline\Gamma_Z\) be the local cohomology functor of Section \hyperref[bp179]{2.14}. We show that \(R\underline\Gamma_Z\) preserves \(D^b_{\mathrm{ctf}}\) and that it is \(i^!\). For \(K\in D^b_{\mathrm{ctf}}(X,A)\), there is a triangle

\[
i_*R\underline\Gamma_ZK\longrightarrow K\longrightarrow Rj_*j^*K
\longrightarrow i_*R\underline\Gamma_ZK[1].
\]

By Proposition \hyperref[bp570]{6.19}, the two terms on the right lie in \(D^b_{\mathrm{ctf}}\), and hence so does \(i_*R\underline\Gamma_ZK\). Since \(i^*i_*=\mathrm{id}\) and \(i^*\) preserves \(D^b_{\mathrm{ctf}}\), we get \(R\underline\Gamma_ZK\in D^b_{\mathrm{ctf}}(Z,A)\). Explicitly, if \(K\) has Tor amplitude in \([a,b]\), then \(R\underline\Gamma_ZK\) has Tor amplitude in \([a,b+\operatorname{cd}_\ell(j)+1]\). The functor \(R\underline\Gamma_Z:D(X_{\mathrm{pro}\acute{e}t},A)\to D(Z_{\mathrm{pro}\acute{e}t},A)\) is right adjoint to \(i_*\) (Lemma lemma-sections-support-closed-subscheme), and \(\Loc(D^b_{\mathrm{ctf}}(-,A))\) is a full subcategory of \(D((-)_{\mathrm{pro}\acute{e}t},A)\). By uniqueness of right adjoints, \(R\underline\Gamma_ZK\) is therefore \(i^!K\) in \(\Loc(D^b_{\mathrm{ctf}}(Z,A))\).

\emph{Projections.} Let \(p:\mathbf A^d_Y\to Y\) be the projection. From the adjunction \(Rq_*\dashv q^*(1)[2]\) for the projective line \(q:\mathbf P^1_Y\to Y\) (Proposition \hyperref[bp534]{6.7}) and from Lemma \hyperref[bp538]{6.8}, we deduce that \(p^!\simeq p^*(-)(d)[2d]\) on \(\Loc(D^b_{\mathrm{ctf}}(Y,A))\). This functor preserves \(D^b_{\mathrm{ctf}}\), since pullback preserves a bounded constructible flat model and the Tate twist is locally free of rank one.

For example, for the origin \(i:0\hookrightarrow\mathbf A^1\) over \(\bar k\) and \(K=A\), the triangle gives \(i^!A\simeq A(-1)[-2]\) (Example \hyperref[bp591]{6.25} below). Its Tor amplitude \([2,2]\) lies in the predicted interval \([0,0+1+1]\).

\emph{Gluing.} After restricting to affine opens in \(X\) and \(Y\), a morphism of finite type factors as a closed immersion into an affine space over \(Y\) followed by its projection. Since \(f^!\) is compatible with composition and restriction to opens (Proposition \hyperref[bp510]{5.32}), the restriction of \(f^!K\) is obtained by applying the two functors just considered. Thus \(f^!K\) is bounded constructible of finite Tor amplitude locally on \(X\). A finite affine cover of the quasi-compact scheme \(X\) gives a single cohomological interval and a single interval of Tor amplitude valid everywhere. Finally, \(Rf_!\) preserves \(D^b_{\mathrm{ctf}}\) by Proposition \hyperref[bp390]{4.35}. Both sides of the adjunction therefore lie in \(D^b_{\mathrm{ctf}}\), and its Hom bijections restrict.
\end{proof}


% object: sec-ell-adic-coefficients-2 (blueprint/chap-Constructibility-of-f/sec-ell-adic-coefficients-2/_section.md), content hidden
\subsection*{\texorpdfstring{6.7\quad $\ell$-adic coefficients}{6.7 ell-adic coefficients}}
\phantomsection\label{bp575}%
\addcontentsline{toc}{subsection}{\texorpdfstring{6.7\quad $\ell$-adic coefficients}{6.7 ell-adic coefficients}}
% text: prose-06-19
\phantomsection\label{bp576}%
For \(\Lambda=\mathcal O_E\), \(E\) and \(\overline{\mathbf Q}_\ell\) we have \(D^b_{\mathrm{ctf}}=\Dbc\), and we show that \(f^!\) preserves \(\Dbc\). The argument of Section \hyperref[bp568]{6.6} uses two geometric facts, namely the finiteness of \(Rj_*\) for an open immersion and the adjunction for the projective line, and we deduce both from the case of torsion coefficients. This is possible because an integral complex lies in \(\Dbc\) if it is derived complete with constructible reduction modulo \(\pi_E\) (Proposition \hyperref[bp226]{3.11}), and rational complexes are built from integral ones (Lemma lemma-integral-models).

For \(Rj_*\) this reduction is immediate, since \(Rj_*\) commutes with \(R\varprojlim\) and with cones (Proposition \hyperref[bp577]{6.21}). For the projective line it is not, because the unit and counit are morphisms, and a compatible family of morphisms modulo \(\ell^n\) need not come from a morphism of the limits. We therefore construct the class of the diagonal and the trace once, at the level of inverse systems (Subsection \hyperref[bp579]{6.7.1}). Besides this, the integral case uses the structure of integral constructible complexes and the fact that proper direct image preserves direct sums (Subsection \hyperref[bp469]{5.4.2}).

% object: proposition-direct-image-along-open-immersion (blueprint/chap-Constructibility-of-f/sec-ell-adic-coefficients-2/proposition-direct-image-along-open-immersion/_section.md)
\phantomsection\label{bp577}%
\begin{bpstatementit}{Proposition~6.21}{ (Direct image along an open immersion)}
Let \(j:U\hookrightarrow X\) be a quasi-compact open immersion. Then \(Rj_*:D(U_{\mathrm{pro}\acute{e}t},\Lambda)\to D(X_{\mathrm{pro}\acute{e}t},\Lambda)\) maps \(\Dbc(U,\Lambda)\) into \(\Dbc(X,\Lambda)\) for \(\Lambda=\mathcal O_E\), \(E\) and \(\overline{\mathbf Q}_\ell\).
\end{bpstatementit}

% proof: proof/proposition-direct-image-along-open-immersion
\begin{proof}[{Proof}]\phantomsection\label{bp578}%
For \(\mathcal O_E\) we use the criterion of Proposition \hyperref[bp226]{3.11}, which asks for derived completeness and constructible reduction. Let \(K\in\Dbc(U,\mathcal O_E)\). Since \(Rj_*\) commutes with \(R\varprojlim\) [Stacks, \href{https://stacks.math.columbia.edu/tag/099L}{099L}], \(Rj_*K\) is derived complete. Since \(Rj_*\) also commutes with cones, the reduction of \(Rj_*K\) modulo \(\pi_E\) is \(Rj_*\) of the reduction of \(K\). That reduction comes from \(U_{\acute{e}t}\). Its direct image is therefore the pullback of the étale one [Stacks, \href{https://stacks.math.columbia.edu/tag/09A6}{09A6}], which is bounded constructible by Corollary \hyperref[bp566]{6.18}. Both conditions of the criterion hold, and so \(Rj_*K\in\Dbc(X,\mathcal O_E)\).

For \(E\) and \(\overline{\mathbf Q}_\ell\) we reduce to the integral case. Since \(Rj_*\) is exact, the objects \(K\) with \(Rj_*K\in\Dbc\) form a triangulated subcategory. It therefore suffices to treat the generators of Lemma lemma-integral-models. These come from integral objects by inverting \(\ell\) and by a colimit over finite extensions of \(E\). Both operations are filtered colimits of systems with a common lower bound, and \(Rj_*\) commutes with such colimits (Corollary \hyperref[bp86]{2.29}). Hence \(Rj_*\) of a generator is obtained from \(Rj_*\) of an integral object by the same operations, and it lies in \(\Dbc\).
\end{proof}

% object: subsec-The-projective-line (blueprint/chap-Constructibility-of-f/sec-ell-adic-coefficients-2/subsec-The-projective-line/_section.md), content hidden
\subsubsection*{\texorpdfstring{6.7.1\quad The projective line}{6.7.1 The projective line}}
\phantomsection\label{bp579}%
\addcontentsline{toc}{subsubsection}{\texorpdfstring{6.7.1\quad The projective line}{6.7.1 The projective line}}
% text: prose-06-20
\phantomsection\label{bp580}%
With torsion coefficients, \(Rq_*\dashv q^*(1)[2]\) is built from the cycle class \(\delta:\Delta_*\Lambda\to\Lambda(1)[2]\) of the diagonal of \(\mathbf P^1_Y\times_Y\mathbf P^1_Y\), the trace \(\operatorname{Tr}_q:Rq_*\Lambda(1)[2]\to\Lambda\) and two identities between them (Proposition \hyperref[bp534]{6.7}). The classes modulo \(\ell^n\), however, are morphisms in different derived categories, compatible under reduction. Since morphisms in a derived category are not computed levelwise, such a family need not give a morphism after \(R\varprojlim\). We therefore construct the class once, in the derived category of inverse systems of sheaves, and then pass to the limit.

% object: proposition-ell-adic-cycle-class-diagonal (blueprint/chap-Constructibility-of-f/sec-ell-adic-coefficients-2/subsec-The-projective-line/proposition-ell-adic-cycle-class-diagonal/_section.md)
\phantomsection\label{bp581}%
\begin{bpstatementit}{Proposition~6.22}{ (The $\ell$-adic cycle class of the diagonal)}
Let \(q:\mathbf P^1_Y\to Y\) be the projection, \(\Delta\) the diagonal of \(\mathbf P^1_Y\times_Y\mathbf P^1_Y\) and \(p_1,p_2\) its projections. There are a morphism \(\delta_{\mathcal O_E}:\Delta_*\mathcal O_E\to
\mathcal O_E(1)[2]\) reducing to the cycle classes of Lemma \hyperref[bp523]{6.3} and a trace \(\operatorname{Tr}_q:Rq_*\mathcal O_E(1)[2]\to\mathcal O_E\) lifting those of Lemma \hyperref[bp526]{6.4}, with \(\operatorname{Tr}_{p_i}\circ Rp_{i*}(\delta_{\mathcal O_E})=1\) for \(i=1,2\). Extending scalars gives \(\delta_\Lambda\) and \(\operatorname{Tr}_q\) with the same identities for \(\Lambda=E,\overline{\mathbf Q}_\ell\). All are compatible with base change along separated morphisms of finite type \(Y'\to Y\).
\end{bpstatementit}

% proof: proof/proposition-ell-adic-cycle-class-diagonal
\begin{proof}[{Proof}]\phantomsection\label{bp582}%
The cycle class comes from a single map defined with \(\mathbf Z\)-coefficients, the boundary of the divisor sequence. Hence all its reductions are reductions of one map of systems.

The diagonal is an effective Cartier divisor, which gives a boundary morphism \(\Delta_*\mathbf Z\to\mathbf G_m[1]\) (Lemma \hyperref[bp523]{6.3}). Tensoring with the system \((\mathbf Z/\ell^n)_n\) and using the Kummer sequence, we get one morphism of systems \(\Delta_*\mathbf Z/\ell^\bullet\to\mathbf Z/\ell^\bullet(1)[2]\). Its level \(n\) is the class of Lemma \hyperref[bp523]{6.3} for \(\mathbf Z/\ell^n\). The transition maps between levels are part of the data, and this is what a family of separately constructed classes would lack. Next we take \(R\varprojlim\). The site is replete with exact products (Propositions \hyperref[bp65]{2.22} and \hyperref[bp68]{Proposition 2.23}). Hence \(R\varprojlim\circ(-\otimes\mathbf Z/\ell^\bullet)\) is derived \(\ell\)-completion [Stacks, \href{https://stacks.math.columbia.edu/tag/099M}{099M}], and the limit \(\Delta_*\mathbf Z_\ell\to\mathbf Z_\ell(1)[2]\) is the completion of the boundary map of Lemma \hyperref[bp523]{6.3}. It is linear over the completed coefficient sheaf by [Stacks, \href{https://stacks.math.columbia.edu/tag/099I}{099I}]. We obtain the trace in the same way from the decomposition of \(Rq_*\mathbf Z/\ell^\bullet\).

Each \(\operatorname{Tr}_{p_i}\circ Rp_{i*}(\delta)\) is an endomorphism of \(\mathbf Z/\ell^\bullet\) in degree \(0\), that is, a family of endomorphisms of the levels. These are the identity by Proposition \hyperref[bp534]{6.7}. We then extend scalars to \(\mathcal O_E\), invert \(\ell\) and take the colimit over finite extensions (Corollary \hyperref[bp86]{2.29}). Base change holds because the Cartier and Kummer maps are natural and pullback commutes with \(R\varprojlim\) here (Lemma \hyperref[bp167]{2.54}).
\end{proof}

% text: prose-06-21
\phantomsection\label{bp583}%
To turn \(\delta_\Lambda\) and \(\operatorname{Tr}_q\) into an adjunction we repeat the torsion calculation. It used the projection formula and proper base change on the categories where \(Rq_*\) and \(q^*(1)[2]\) act. Proper base change on \(\Loc(D^b_{\mathrm{ctf}})\) is Proposition \hyperref[bp488]{5.24}, and we extend the projection formula in the same way.

% object: proposition-projection-formula-rational-coefficients (blueprint/chap-Constructibility-of-f/sec-ell-adic-coefficients-2/subsec-The-projective-line/proposition-projection-formula-rational-coefficients/_section.md)
\phantomsection\label{bp584}%
\begin{bpstatementit}{Proposition~6.23}{ (Projection formula with rational coefficients)}
For \(\Lambda=E\) or \(\overline{\mathbf Q}_\ell\), \(\Loc(D^b_{\mathrm{ctf}}(X,\Lambda))\) is closed under \(\otimes^{\mathbf L}\), and for proper \(p:X\to Y\) the map \((Rp_*M)\otimes^{\mathbf L}N\to Rp_*(M\otimes^{\mathbf L}p^*N)\) is an isomorphism for \(M\in\Loc(D^b_{\mathrm{ctf}}(X,\Lambda))\) and \(N\in\Loc(D^b_{\mathrm{ctf}}(Y,\Lambda))\).
\end{bpstatementit}

% proof: proof/proposition-projection-formula-rational-coefficients
\begin{proof}[{Proof}]\phantomsection\label{bp585}%
Both claims are first checked on \(\Dbc\) and then extended. Tensor products and both sides of the projection map commute with direct sums and cones in each variable. Hence, fixing all variables but one, the objects satisfying a claim form a localizing subcategory. It suffices that this subcategory contains \(\Dbc\) (Lemma \hyperref[bp485]{5.23}). By the dévissage of Corollary \hyperref[bp253]{3.20}, \(\Dbc\) is closed under \(\otimes^{\mathbf L}\), since \(a_!L\otimes^{\mathbf L}b_!M\) is the extension by zero of a lisse sheaf on the intersection of the two strata. The projection formula on \(\Dbc\) is Proposition \hyperref[bp405]{4.40}.
\end{proof}

% text: prose-06-22
\phantomsection\label{bp586}%
Now the torsion proof of the adjunction for \(\mathbf P^1\) can be run with \(\delta_\Lambda\) and \(\operatorname{Tr}_q\), once \(s_1=s_2=1\) is known.

% object: proposition-projective-line-rational-coefficients (blueprint/chap-Constructibility-of-f/sec-ell-adic-coefficients-2/subsec-The-projective-line/proposition-projective-line-rational-coefficients/_section.md)
\phantomsection\label{bp587}%
\begin{bpstatementit}{Proposition~6.24}{ (The projective line, rational coefficients)}
For \(\Lambda=E\) or \(\overline{\mathbf Q}_\ell\), the functor \(q^*(1)[2]:\Loc(D^b_{\mathrm{ctf}}(Y,\Lambda))\to\Loc(D^b_{\mathrm{ctf}}(\mathbf P^1_Y,\Lambda))\) is right adjoint to \(Rq_*\). Hence \(q^!\simeq q^*(1)[2]\) for the right adjoint \(q^!\) of Theorem \hyperref[bp506]{5.31}.
\end{bpstatementit}

% proof: proof/proposition-projective-line-rational-coefficients
\begin{proof}[{Proof}]\phantomsection\label{bp588}%
Define the unit and counit by the formulas (6.7) and (6.8) of Proposition \hyperref[bp534]{6.7}, using \(\Delta_*M\simeq p_1^*M\otimes^{\mathbf L}\Delta_*\Lambda\) [Stacks, \href{https://stacks.math.columbia.edu/tag/09AD}{09AD}] and the maps \(\delta_\Lambda\) and \(\operatorname{Tr}_q\) in place of their torsion versions. The proof of the identities \(s_1=s_2=1\) computed stalks of torsion sheaves, but apart from these identities the torsion calculation of the triangle identities used only the projection formula and proper base change. These hold here by Propositions \hyperref[bp584]{6.23} and \hyperref[bp488]{Proposition 5.24}. Thus the calculation reduces the triangle identities to the identities \(\operatorname{Tr}_{p_i}\circ Rp_{i*}(\delta_\Lambda)=1\) of Proposition \hyperref[bp581]{6.22}. Finally, uniqueness of adjoints gives the last claim.
\end{proof}


% object: subsec-Finiteness-of-f (blueprint/chap-Constructibility-of-f/sec-ell-adic-coefficients-2/subsec-Finiteness-of-f/_section.md), content hidden
\subsubsection*{\texorpdfstring{6.7.2\quad Finiteness of $f^!$}{6.7.2 Finiteness of f!}}
\phantomsection\label{bp589}%
\addcontentsline{toc}{subsubsection}{\texorpdfstring{6.7.2\quad Finiteness of $f^!$}{6.7.2 Finiteness of f!}}
% text: prose-06-23
\phantomsection\label{bp590}%
For admissible \(\Lambda\) and a closed immersion \(i:Z\hookrightarrow X\) with complement \(j\), the local cohomology functor \(R\underline\Gamma_Z\) is right adjoint to \(i_*=Ri_!\) (Lemma lemma-sections-support-closed-subscheme). Hence \(i^!\simeq R\underline\Gamma_Z\), and \(i^!\) is computed by the localization triangle

\[
i_*i^!K\longrightarrow K\longrightarrow Rj_*j^*K
\longrightarrow i_*i^!K[1].
\]

% object: example-closed-support (blueprint/chap-Constructibility-of-f/sec-ell-adic-coefficients-2/subsec-Finiteness-of-f/example-closed-support/_section.md)
\phantomsection\label{bp591}%
\begin{bpstatement}{Example~6.25}{ (Local cohomology at the origin of the affine line)}
Let \(X=\mathbf A^1_k\), \(k\) algebraically closed, \(Z\) the origin and \(U=\mathbf G_{m,k}\). Then \(Rj_*\Lambda_U\) has \(\mathcal H^0=\Lambda_X\) and \(\mathcal H^1=i_*\Lambda(-1)\), from the Kummer covers \(t\mapsto t^{\ell^n}\). The triangle gives \(i^!\Lambda_X\simeq\Lambda(-1)[-2]\), as Theorem \bpout{9.3}, applied to the structure maps of \(Z\) and \(X\), predicts.
\end{bpstatement}

% text: prose-06-24
\phantomsection\label{bp592}%
For \(\mathcal O_E\) we also need the projection formula on the localizing categories, which follows from the case of \(\Dbc\) and the results of Subsection \hyperref[bp469]{5.4.2}.

% object: lemma-projection-formula-integral-coefficients (blueprint/chap-Constructibility-of-f/sec-ell-adic-coefficients-2/subsec-Finiteness-of-f/lemma-projection-formula-integral-coefficients/_section.md)
\phantomsection\label{bp593}%
\begin{bpstatementit}{Lemma~6.26}{ (Projection formula with integral coefficients)}
For proper \(p:X\to Y\) the map \((Rp_*M)\otimes^{\mathbf L}N\to Rp_*(M\otimes^{\mathbf L}p^*N)\) is an isomorphism for \(M\in\Loc(D^b_{\mathrm{ctf}}(X,\mathcal O_E))\) and \(N\in\Loc(D^b_{\mathrm{ctf}}(Y,\mathcal O_E))\).
\end{bpstatementit}

% proof: proof/lemma-projection-formula-integral-coefficients
\begin{proof}[{Proof}]\phantomsection\label{bp594}%
On \(\Dbc\) this is Proposition \hyperref[bp405]{4.40}. Both sides make sense on the localizing categories, because \(M\otimes^{\mathbf L}p^*N\) lies in \(\Loc(D^b_{\mathrm{ctf}}(X,\mathcal O_E))\). For \(M\) and \(N\) in \(\Dbc\) this is Lemma \hyperref[bp474]{5.20}, and since the tensor product commutes with direct sums and cones in each variable, the argument of Proposition \hyperref[bp584]{6.23} extends it to the localizing categories. Indeed, fix one variable. Both sides are then exact and preserve direct sums in the other; for the variable \(M\) this holds because \(Rp_*\) preserves direct sums (Proposition \hyperref[bp477]{5.21}). Therefore Lemma \hyperref[bp485]{5.23}, applied first in \(M\) and then in \(N\), extends the isomorphism from \(\Dbc\) to the localizing categories.
\end{proof}

% text: prose-06-25
\phantomsection\label{bp595}%
We can now repeat the argument for finite coefficients with \(\mathcal O_E\), \(E\) and \(\overline{\mathbf Q}_\ell\).

% object: theorem-f-shriek-ell-adic-coefficients (blueprint/chap-Constructibility-of-f/sec-ell-adic-coefficients-2/subsec-Finiteness-of-f/theorem-f-shriek-ell-adic-coefficients/_section.md)
\phantomsection\label{bp596}%
\begin{bpstatementit}{Theorem~6.27}{ ($f^!$ with $\ell$-adic coefficients)}
Let \(\Lambda=\mathcal O_E\), \(E\) or \(\overline{\mathbf Q}_\ell\), and let \(f:X\to Y\) be separated of finite type. Then \(f^!\) maps \(D^b_{\mathrm{ctf}}(Y,\Lambda)\) into \(D^b_{\mathrm{ctf}}(X,\Lambda)\), and this restriction is right adjoint to \(Rf_!:D^b_{\mathrm{ctf}}(X,\Lambda)\to D^b_{\mathrm{ctf}}(Y,\Lambda)\).
\end{bpstatementit}

% proof: proof/theorem-f-shriek-ell-adic-coefficients
\begin{proof}[{Proof}]\phantomsection\label{bp597}%
As for finite coefficients, we factor \(f\) locally into a closed immersion and a projection from affine space, and treat the two factors separately.

Choose finite affine covers of \(X\) and \(Y\) and closed immersions \(U\hookrightarrow\mathbf A^d_V\). Here \(u:U\to X\) and \(v:V\to Y\) are the inclusions of the affine opens, and \(i:U\hookrightarrow\mathbf A^d_V\) and \(p:\mathbf A^d_V\to V\) are the two factors. Then \(u^*f^!\simeq i^!p^!v^*\) (Proposition \hyperref[bp510]{5.32}), and being in \(\Dbc\) is local on \(X\), so it suffices to show that \(p^!\) and \(i^!\) preserve \(\Dbc\).

For \(E\) and \(\overline{\mathbf Q}_\ell\) we argue as in Section \hyperref[bp568]{6.6}. For \(p\) we use Proposition \hyperref[bp587]{6.24}, the adjunction \(j_!\dashv j^*\) for \(\mathbf A^1_V\subset\mathbf P^1_V\), and Proposition \hyperref[bp510]{5.32} applied \(d\) times. They give \(p^!\simeq p^*(d)[2d]\), which preserves \(\Dbc\). For \(i\), let \(j'\) be the inclusion of the open complement of \(U\) in \(\mathbf A^d_V\). Then the localization triangle before Example \hyperref[bp591]{6.25}, with \(j'\) in place of \(j\), shows that \(i^!K\in\Dbc(U,\Lambda)\) for \(K\in\Dbc(\mathbf A^d_V,\Lambda)\). Indeed, its third term \(Rj'_*j'^*K\) is in \(\Dbc\) by Proposition \hyperref[bp577]{6.21}, and \(i^*i_*=\mathrm{id}\).

For \(\mathcal O_E\) that argument applies once its ingredients are available with integral coefficients. They are as follows. First, the adjunction for the projective line comes from Proposition \hyperref[bp581]{6.22}, with the triangle identities verified as in Proposition \hyperref[bp587]{6.24}. Next, the projection formula is Lemma \hyperref[bp593]{6.26}, and inverse image and tensor product preserve \(\Dbc(-,\mathcal O_E)\) by Lemma \hyperref[bp474]{5.20}. Moreover, the finiteness of \(Rj_*\) and proper base change are Propositions \hyperref[bp577]{6.21} and \hyperref[bp488]{Proposition 5.24}. We also use that derived completeness is local on \(X\); it is part of the criterion for \(\Dbc\) with \(\mathcal O_E\)-coefficients. Finally, \(Rf_!\) preserves \(\Dbc\) (Proposition \hyperref[bp492]{5.26}), and so the adjunction restricts to \(\Dbc\).
\end{proof}



% object: sec-Proof-of-the-theorem (blueprint/chap-Constructibility-of-f/sec-Proof-of-the-theorem/_section.md), content hidden
\subsection*{\texorpdfstring{6.8\quad Proof of the theorem}{6.8 Proof of the theorem}}
\phantomsection\label{bp598}%
\addcontentsline{toc}{subsection}{\texorpdfstring{6.8\quad Proof of the theorem}{6.8 Proof of the theorem}}
% text: prose-06-26
\phantomsection\label{bp599}%
The parts of Theorem \hyperref[bp413]{5.2} were proved separately for each kind of coefficient ring, and we collect the references.

% proof: proof/theorem-exceptional-inverse-image
\begin{proof}[{Proof of Theorem 5.2}]\phantomsection\label{bp600}%
Part (1) is Lemma \hyperref[bp503]{5.30}. Its compactness statement was proved separately for each kind of coefficient ring: for finite \(A\) in Proposition \hyperref[bp451]{5.14}, for \(E\) and \(\overline{\mathbf Q}_\ell\) in Theorem \hyperref[bp446]{5.13}, and for \(\mathcal O_E\) in Proposition \hyperref[bp477]{5.21}.

Part (2) is Proposition \hyperref[bp492]{5.26}. It rests on the fact that proper direct image preserves direct sums, which was proved for finite \(A\) in Corollary \hyperref[bp457]{5.16}, for \(E\) and \(\overline{\mathbf Q}_\ell\) in Lemma \hyperref[bp467]{5.18}, and for \(\mathcal O_E\) in Proposition \hyperref[bp477]{5.21} (see Proposition \hyperref[bp482]{5.22}).

Part (3) is Theorem \hyperref[bp506]{5.31}.

Part (4) is Theorem \hyperref[bp573]{6.20} for finite \(A\) and Theorem \hyperref[bp596]{6.27} for \(\mathcal O_E\), \(E\) and \(\overline{\mathbf Q}_\ell\). The second theorem includes \(\mathbf Z_\ell\) and \(\mathbf Q_\ell\) as the case \(E=\mathbf Q_\ell\); recall that \(D^b_{\mathrm{ctf}}=\Dbc\) in the integral and rational cases. For the last assertion, note that the categories \(D^b_{\mathrm{ctf}}(X,\Lambda)\) and \(D^b_{\mathrm{ctf}}(Y,\Lambda)\) are full subcategories of the localizing categories, and \(Rf_!\) and \(f^!\) preserve them. Hence the Hom bijections of the adjunction restrict to them.
\end{proof}



% object: chap-The-six-operations-on-Dbctf (blueprint/chap-The-six-operations-on-Dbctf/_section.md), content hidden
\section*{\texorpdfstring{7\quad The six operations on $\Dctf$}{7 The six operations on Dctf}}
\phantomsection\label{bp601}%
\addcontentsline{toc}{section}{\texorpdfstring{7\quad The six operations on $\Dctf$}{7 The six operations on Dctf}}
% text: prose-07-01
\phantomsection\label{bp602}%
Smooth base change and Verdier biduality will be proved with torsion coefficients, as were proper base change and Deligne's theorem on \(Rj_*\). The Weil conjectures, however, need them over \(\overline{\mathbf Q}_\ell\), where Frobenius eigenvalues live. This chapter provides the principle that carries every such statement across. A canonical map built from the six operations is an isomorphism for all admissible \(\Lambda\) as soon as it is one for every finite ring (Proposition \hyperref[bp634]{7.12}).

For \(\mathcal O_E\) both sides are derived complete, and by derived Nakayama it suffices to check the map modulo \(\pi_E\). For instance, let \(K\) be in \(D^b_{\mathrm{ctf}}(X,\mathcal O_E)\). Both sides of the biduality map \(K\to\mathbb D_X\mathbb D_XK\) are derived complete. Modulo \(\pi_E\), moreover, the map is biduality with \(\mathbf F_E\)-coefficients, which is an isomorphism. Hence the biduality map is an isomorphism. From there we pass to \(E\) and \(\overline{\mathbf Q}_\ell\): every object over \(E\) is built by cones from integral models with \(\ell\) inverted, and every object over \(\overline{\mathbf Q}_\ell\) is defined over a finite extension \(E\) (Proposition \hyperref[bp629]{7.11}). In this way smooth base change (Theorem \hyperref[bp681]{8.11}) and Verdier biduality (Theorem \bpout{10.8}) pass from torsion to \(\ell\)-adic coefficients.

The principle needs two things, and neither is formal. First, both sides must lie in \(D^b_{\mathrm{ctf}}(-,\Lambda)\). Thus all six operations must preserve these categories, and not only \(Rf_!\) and \(f^!\), which we now control (Theorem \hyperref[bp413]{5.2}, Section \hyperref[bp598]{6.8}). Second, each operation must commute with change of coefficients. For \(f^*\) and \(\otimes^{\mathbf L}\) this is clear, since a change of coefficients is itself a tensor product. For a right adjoint, however, it is a real statement. For example, consider a map of finite rings \(A\to B\); then the isomorphism \(R\Gamma(X,K)\otimes^{\mathbf L}_AB\simeq
R\Gamma(X,K\otimes^{\mathbf L}_AB)\) is a projection formula, and it uses finite cohomological dimension. Next consider \(\mathcal O_E\to E\). Since \(E\) is the colimit of the modules \(\pi_E^{-n}\mathcal O_E\), the comparison says that \(Rf_*\) commutes with a filtered colimit. This holds because the terms are uniformly bounded below.

Internal Hom causes a further difficulty, because it hides the exceptional inverse image. For a closed immersion \(i\) we have

\[
\cHom(i_*\Lambda,K)\simeq i_*i^!K .
\]

Before we can say anything about \(\cHom(K,L)\otimes^{\mathbf L}\Lambda'\), we must therefore know how \(i^!\) behaves under change of coefficients. The comparison map for \(f^!\) is defined by adjunction, and to show that it is an isomorphism we compute \(f^!\) once more, for a closed immersion and for \(\mathbf P^1\). For \(\mathbf P^1\), for instance, it comes down to the compatibility of the cycle class of the diagonal and of the trace with change of coefficients (Proposition \hyperref[bp626]{7.10}).

The first requirement needs no construction: on the pro-étale site the six operations are the ordinary derived functors of sheaves of modules, and we only check that they preserve \(D^b_{\mathrm{ctf}}(-,\Lambda)\) (Proposition \hyperref[bp605]{7.1}). For torsion coefficients on the étale site, these finiteness theorems are in SGA 4 and SGA 4\(\tfrac12\) [\hyperref[bp702]{SGA4}], [\hyperref[bp703]{SGA4½}]. In general, the work is in \(Rf_*\) and in \(\cHom\) (Section \hyperref[bp603]{7.1}). There, a compactification reduces \(Rf_*\) to proper direct image and \(Rj_*\), while the formula above and its analogue for open immersions express \(\cHom\) through \(Ra_*\), \(a^!\) and tensor products. Tensor products, finally, are handled by the local description of a complex of finite Tor amplitude: on the strata of a finite stratification it is locally \(\Lambda_S\otimes P\) with \(P\) perfect (Section \hyperref[bp230]{3.4}).

From now on we work in \(D^b_{\mathrm{ctf}}\) without further comment and pass between \(\mathcal O_E\), \(E\), \(\overline{\mathbf Q}_\ell\) and finite rings whenever a statement is easier for one of them, as in the trace formula and in the proof of the Weil conjectures.

% object: sec-Stability-of-Dbctf (blueprint/chap-The-six-operations-on-Dbctf/sec-Stability-of-Dbctf/_section.md), content hidden
\subsection*{\texorpdfstring{7.1\quad Stability of $\Dctf$}{7.1 Stability of Dctf}}
\phantomsection\label{bp603}%
\addcontentsline{toc}{subsection}{\texorpdfstring{7.1\quad Stability of $\Dctf$}{7.1 Stability of Dctf}}
% text: prose-07-02
\phantomsection\label{bp604}%
Pullback and tensor product are local, and they preserve the description of a complex of \(D^b_{\mathrm{ctf}}\) on the strata of a stratification as \(\Lambda_S\otimes P\) with \(P\) perfect. The direct images and internal Hom are not local. For the direct images a compactification \(f=pj\) helps. Since \(Rf_*=Rp_*Rj_*\), it suffices to know proper direct image and direct image along an open immersion. Internal Hom hides both of these operations and also \(f^!\). For an open immersion \(j\) we have \(\cHom(j_!\Lambda,K)=Rj_*j^*K\), and for a closed immersion \(i\) we have the formula \(\cHom(i_*\Lambda,K)\simeq i_*i^!K\) recalled at the beginning of this chapter. If \(a\) is the locally closed immersion of a stratum and \(P\) is a perfect complex on it, the formula

\[
\cHom(a_!P,K)\simeq Ra_*\bigl(P^\vee\otimes^{\mathbf L}a^!K\bigr)
\]

combines the two. Every constructible complex is built from objects \(a_!P\) by finitely many cones. By this formula, internal Hom is therefore expressed through \(Ra_*\), \(a^!\) and tensor products. This is one reason we needed \(f^!\) on \(D^b_{\mathrm{ctf}}\) first. The result is Proposition \hyperref[bp605]{7.1}.

% object: proposition-six-operations-dctf (blueprint/chap-The-six-operations-on-Dbctf/sec-Stability-of-Dbctf/proposition-six-operations-dctf/_section.md)
\phantomsection\label{bp605}%
\begin{bpstatementit}{Proposition~7.1}{ (The six operations on $\Dctf$)}
For every admissible \(\Lambda\in\{A,\mathbf Z_\ell,\mathcal O_E,\mathbf Q_\ell,E,\overline{\mathbf Q}_\ell\}\), the operations \(f^*\), \(Rf_*\), \(Rf_!\), \(f^!\), \(\otimes^{\mathbf L}\) and \(\cHom\) preserve \(D^b_{\mathrm{ctf}}(-,\Lambda)\) for morphisms of schemes of finite type over \(k\). Their adjunctions are the restrictions of the adjunctions \(f^*\dashv Rf_*\) and \(-\otimes^{\mathbf L}K\dashv\cHom(K,-)\) on \(D(-_{\mathrm{pro}\acute{e}t},\Lambda)\) or of the adjunction \(Rf_!\dashv f^!\) on \(\Loc(D^b_{\mathrm{ctf}}(-,\Lambda))\).
\end{bpstatementit}

% proof: proof/proposition-six-operations-dctf
\begin{proof}[{Proof}]\phantomsection\label{bp606}%
We treat the operations in the order in which they depend on each other: first the local operations \(f^*\) and \(\otimes^{\mathbf L}\), then the two direct images, and finally internal Hom, which we express through the others.

Pullback and tensor product preserve the property of being, on the strata of a finite stratification, locally \(\Lambda_S\otimes P\) with \(P\) a perfect complex. For pullback this is clear, since the strata pull back to a stratification, while for tensor products we take a common refinement of the two stratifications, and on each of its strata we use that \(P\otimes^{\mathbf L}Q\) is perfect when \(P\) and \(Q\) are. Similarly, extension by zero preserves finite Tor amplitude, since it is exact and satisfies the projection formula \(a_!K\otimes^{\mathbf L}N\simeq a_!(K\otimes^{\mathbf L}a^*N)\). Concretely, for finite \(A\) these statements can be checked on the bounded flat constructible models of Lemma \hyperref[bp235]{3.13}. For \(\mathcal O_E\), \(E\) and \(\overline{\mathbf Q}_\ell\), on the other hand, they follow from Lemma \hyperref[bp474]{5.20} and from Definition \hyperref[bp199]{3.1} of constructible sheaves: on each stratum, such a sheaf is locally \(\Lambda_X\otimes_{\Lambda(*)}P\) with \(P\) finitely generated over the discrete valuation ring or field \(\Lambda(*)\). Over such a ring a finitely generated module has a finite free resolution of length at most one. Hence \(P\) is perfect.

For the direct images we use a compactification to separate the proper part from the open part. Proper direct image preserves \(D^b_{\mathrm{ctf}}(-,\Lambda)\) by Proposition \hyperref[bp405]{4.40}. Direct image along an open immersion preserves it by Proposition \hyperref[bp570]{6.19} for finite \(A\) and Proposition \hyperref[bp577]{6.21} for \(\mathcal O_E\), \(E\) and \(\overline{\mathbf Q}_\ell\). If \(f=pj\) is a compactification, then \(Rf_*=Rp_*Rj_*\) and \(Rf_!=Rp_*j_!\). Each of these composites consists of functors already known to preserve \(D^b_{\mathrm{ctf}}\), which proves both assertions on direct images. That \(f^!\) preserves \(D^b_{\mathrm{ctf}}(-,\Lambda)\) is part of Theorem \hyperref[bp413]{5.2}.

For internal Hom we reduce to a first argument supported on a single stratum, where \(\cHom\) can be written in terms of \(Ra_*\), \(a^!\) and \(\otimes^{\mathbf L}\). This is possible because the functor \(\cHom(-,K)\) sends triangles in the first argument to triangles, and \(D^b_{\mathrm{ctf}}\) is closed under cones and shifts. Hence, for a triangle \(M'\to M\to M''\), the complex \(\cHom(M,K)\) lies in \(D^b_{\mathrm{ctf}}\) whenever \(\cHom(M',K)\) and \(\cHom(M'',K)\) do. Using a finite stratification and finite cones, we may therefore reduce to the case where the first argument is \(a_!P\). Here \(a:S\hookrightarrow X\) is locally closed and \(P\) is a locally constant perfect complex on \(S\). For finite \(A\) this reduction follows from the bounded flat model and its stupid filtration, and for \(\mathcal O_E\), \(E\) and \(\overline{\mathbf Q}_\ell\) it follows from the same description of constructible sheaves on strata. By tensor-Hom adjunction and the projection formula for \(a_!\) we have

\[
\cHom_X(a_!P,K)
\simeq Ra_*\cHom_S(P,a^!K)
\simeq Ra_*(P^\vee\otimes^{\mathbf L}a^!K).
\]

Here \(P^\vee=\cHom_S(P,\Lambda_S)\), and the second identity holds because \(P\) is dualizable. The last term belongs to \(D^b_{\mathrm{ctf}}(-,\Lambda)\) by the cases already established, namely those of \(a^!\), of tensor product and of \(Ra_*\).

Finally, since the inclusions are full, the Hom bijections of the ambient categories restrict to the stated adjunctions.
\end{proof}

% object: example-internal-hom-affine-line (blueprint/chap-The-six-operations-on-Dbctf/sec-Stability-of-Dbctf/example-internal-hom-affine-line/_section.md)
\phantomsection\label{bp607}%
\begin{bpstatement}{Example~7.2}{ (Internal Hom on the affine line)}
Over an algebraically closed \(k\), let \(j:U=\mathbf G_{m,k}\hookrightarrow
X=\mathbf A^1_k\) and let \(i\) be the inclusion of the origin. The localization sequence \(0\to j_!\Lambda_U\to\Lambda_X\to i_*\Lambda\to0\) shows that \(\Lambda_X\) is built from \(j_!\Lambda_U\) and \(i_*\Lambda\), and the formula above computes \(\cHom\) out of each piece. By Example \hyperref[bp591]{6.25},

\[
\cHom(j_!\Lambda_U,\Lambda_X)=Rj_*\Lambda_U,
\qquad
\cHom(i_*\Lambda,\Lambda_X)=i_*i^!\Lambda_X\simeq i_*\Lambda(-1)[-2].
\]

The first has \(\mathcal H^0=\Lambda_X\) and \(\mathcal H^1=i_*\Lambda(-1)\). In the second, the local Ext sheaves from the skyscraper sheaf \(i_*\Lambda\) to the constant sheaf vanish except in degree two. Both results are constructible but not lisse. The first is computed by a nonproper direct image, and the second by \(i^!\).
\end{bpstatement}


% object: sec-Change-of-coefficients (blueprint/chap-The-six-operations-on-Dbctf/sec-Change-of-coefficients/_section.md), content hidden
\subsection*{\texorpdfstring{7.2\quad Change of coefficients}{7.2 Change of coefficients}}
\phantomsection\label{bp608}%
\addcontentsline{toc}{subsection}{\texorpdfstring{7.2\quad Change of coefficients}{7.2 Change of coefficients}}
% text: prose-07-03
\phantomsection\label{bp609}%
Smooth base change for \(E\) and \(\overline{\mathbf Q}_\ell\) is deduced from the integral case by extending scalars (Chapter \hyperref[bp637]{8}), and later proofs often argue in the same way, through the reduction principle of Proposition \hyperref[bp634]{7.12}. For this the six operations must commute with extension of scalars. For pullback, extension by zero and tensor product this is formal, since extension of scalars is itself a tensor product. For \(Rf_*\), by contrast, it rests on finite cohomological dimension and on continuity, and for \(f^!\) it rests on a second computation of \(f^!\), for closed immersions and for \(\mathbf P^1\). These results are collected in Proposition \hyperref[bp626]{7.10}.

A change of coefficients is derived extension of scalars along a map of condensed rings, pulled back to \(X\). For a map \(\Lambda\to\Lambda'\) write

\[
K_{\Lambda'}=\Lambda'_X\otimes^{\mathbf L}_{\Lambda_X}K.
\]

This functor \(D(X_{\mathrm{pro}\acute{e}t},\Lambda)\to D(X_{\mathrm{pro}\acute{e}t},\Lambda')\) is left adjoint to the exact restriction of scalars \(D(X_{\mathrm{pro}\acute{e}t},\Lambda')\to D(X_{\mathrm{pro}\acute{e}t},\Lambda)\). We consider the following maps of coefficient rings:

\begin{enumerate}
\item any homomorphism \(A\to B\) of finite admissible rings;
\item the maps \(\mathcal O_E\to\mathcal O_F\) and \(E\to F\) for a finite extension \(F/E\) inside \(\overline{\mathbf Q}_\ell\);
\item localization \(\mathcal O_E\to E\), inclusions \(E\to\overline{\mathbf Q}_\ell\), and continuous maps \(\mathcal O_E\to A\) to finite admissible rings, such as the reductions \(\mathcal O_E\to\mathcal O_E/\pi_E^n\) (continuity is needed for the map to define a map of condensed rings);
\item composites of these maps.
\end{enumerate}

For \(E=\mathbf Q_\ell\), the list contains \(\mathbf Z_\ell\to\mathcal O_F\), \(\mathbf Q_\ell\to F\), \(\mathbf Z_\ell\to\mathbf Q_\ell\), \(\mathbf Z_\ell\to A\) and \(\mathbf Q_\ell\to\overline{\mathbf Q}_\ell\). Thus these maps connect all the admissible coefficient rings. Note that the map \(A\to B\) in (1) need not be flat, and for this reason extension of scalars is derived. As a consequence, the comparison isomorphisms below hold for complexes but not degree by degree, as the following example shows.

% object: example-derived-reduction-point (blueprint/chap-The-six-operations-on-Dbctf/sec-Change-of-coefficients/example-derived-reduction-point/_section.md)
\phantomsection\label{bp610}%
\begin{bpstatement}{Example~7.3}{ (Reducing cohomology modulo $\ell$)}
Let \(X=\operatorname{Spec}\mathbf F_q\) and consider \(\mathbf Z/\ell^2\to\mathbf F_\ell\). Let \(K\) be the lisse sheaf \(\mathbf Z/\ell^2\) on which Frobenius, the topological generator of \(\operatorname{Gal}(\overline{\mathbf F}_q/\mathbf F_q)
\simeq\widehat{\mathbf Z}\), acts by multiplication by the unit \(1+\ell\). Recall that for a finite module \(M\) with a continuous action of \(\widehat{\mathbf Z}\), cohomology is computed by the complex \(M\xrightarrow{F-1}M\) in degrees \(0\) and \(1\), where \(F\) is the action of the generator. For \(K\) this differential is multiplication by \(\ell\). Hence \(H^0(X,K)=\ell\mathbf Z/\ell^2\) and \(H^1(X,K)=(\mathbf Z/\ell^2)/\ell\), both of order \(\ell\). On the other hand, \(K_{\mathbf F_\ell}=K/\ell K\) is \(\mathbf F_\ell\) with trivial action, and \(H^0(X,K/\ell K)=\mathbf F_\ell\). The map \(H^0(X,K)\otimes\mathbf F_\ell\to H^0(X,K/\ell K)\) is zero, since the generator \(\ell\) of \(H^0(X,K)\) reduces to \(0\).

The derived comparison is nevertheless an isomorphism. The terms of the complex are free over \(\mathbf Z/\ell^2\). Therefore \(R\Gamma(X,K)\otimes^{\mathbf L}\mathbf F_\ell\) is \(\mathbf F_\ell\xrightarrow{0}\mathbf F_\ell\), which is \(R\Gamma(X,K/\ell K)\). The class \(1\in H^0(X,K/\ell K)\) comes from the element \(1\) of the degree-zero term, which is not a cocycle before reduction.
\end{bpstatement}

% text: prose-07-04
\phantomsection\label{bp611}%
Before stating the result we fix the comparison maps. Items (1)–(3) below, for \(f^*\), \(j_!\) and \(\otimes^{\mathbf L}\), are identifications of two composites of left adjoints. The remaining items are derived from these: items (4) and (6), for \(Rf_*\) and \(\cHom\), are adjoint transposes of the inverse of (1), respectively (3), composed with the extension of scalars of the counit, and item (5) combines (4) for \(p\) with (2).

% object: definition-change-of-coefficients-maps (blueprint/chap-The-six-operations-on-Dbctf/sec-Change-of-coefficients/definition-change-of-coefficients-maps/_section.md)
\phantomsection\label{bp612}%
\begin{bpstatement}{Definition~7.4}{ (Change-of-coefficients maps)}
Let \(\Lambda\to\Lambda'\) be one of the maps listed above, and let \(f:X\to Y\) be a morphism of schemes of finite type over \(k\).

\begin{enumerate}
\item For \(K\in D(Y_{\mathrm{pro}\acute{e}t},\Lambda)\), the isomorphism \((f^*K)_{\Lambda'}\simeq f^*(K_{\Lambda'})\) is the one given by associativity of derived pullback for ringed topoi [Stacks, \href{https://stacks.math.columbia.edu/tag/06YV}{06YV}]: both sides are the derived pullback along the morphism of ringed topoi \((X_{\mathrm{pro}\acute{e}t},\Lambda'_X)\to(Y_{\mathrm{pro}\acute{e}t},\Lambda_Y)\), factored in the two possible ways.
\item For an open immersion \(j:U\to X\) and \(K\in D(U_{\mathrm{pro}\acute{e}t},\Lambda)\), the isomorphism \((j_!K)_{\Lambda'}\simeq j_!(K_{\Lambda'})\) is the isomorphism of left adjoints corresponding to the equality of their right adjoints, restriction of scalars followed by \(j^*\) and \(j^*\) followed by restriction of scalars.
\item For \(K,L\in D(X_{\mathrm{pro}\acute{e}t},\Lambda)\), the isomorphism \((K\otimes^{\mathbf L}_\Lambda L)_{\Lambda'}\simeq
  K_{\Lambda'}\otimes^{\mathbf L}_{\Lambda'}L_{\Lambda'}\) is the associativity isomorphism: both sides are \(K\otimes^{\mathbf L}_\Lambda L\otimes^{\mathbf L}_\Lambda\Lambda'\).
\item For \(K\in D(X_{\mathrm{pro}\acute{e}t},\Lambda)\), the map \((Rf_*K)_{\Lambda'}\to Rf_*(K_{\Lambda'})\) is the adjoint transpose, under \(f^*\dashv Rf_*\) on \(D(-_{\mathrm{pro}\acute{e}t},\Lambda')\), of \[
   f^*\bigl((Rf_*K)_{\Lambda'}\bigr)\simeq(f^*Rf_*K)_{\Lambda'}
   \xrightarrow{(\varepsilon_K)_{\Lambda'}}K_{\Lambda'},
   \] where the first arrow is the inverse of (1) and \(\varepsilon\) is the counit of \(f^*\dashv Rf_*\) on \(D(-_{\mathrm{pro}\acute{e}t},\Lambda)\).
\item If \(f\) is separated, \(f=p\circ j\) is a compactification and \(K\in D^b_{\mathrm{ctf}}(X,\Lambda)\), the map \((Rf_!K)_{\Lambda'}\to Rf_!(K_{\Lambda'})\) is the composite \[
   (Rp_*j_!K)_{\Lambda'}\longrightarrow Rp_*\bigl((j_!K)_{\Lambda'}\bigr)
   \simeq Rp_*j_!(K_{\Lambda'})
   \] of (4) for \(p\) and \(Rp_*\) of (2).
\item For \(K,L\in D(X_{\mathrm{pro}\acute{e}t},\Lambda)\), the map \(\cHom(K,L)_{\Lambda'}\to\cHom(K_{\Lambda'},L_{\Lambda'})\) is the adjoint transpose, under \(-\otimes^{\mathbf L}K_{\Lambda'}\dashv
  \cHom(K_{\Lambda'},-)\), of \[
   \cHom(K,L)_{\Lambda'}\otimes^{\mathbf L}_{\Lambda'}K_{\Lambda'}
   \simeq\bigl(\cHom(K,L)\otimes^{\mathbf L}_\Lambda K\bigr)_{\Lambda'}
   \xrightarrow{(\mathrm{ev})_{\Lambda'}}L_{\Lambda'},
   \] where the first arrow is the inverse of (3) and \(\mathrm{ev}\) is evaluation.
\end{enumerate}
\end{bpstatement}

% text: prose-07-05
\phantomsection\label{bp613}%
The map for \(f^!\) is defined in the same way, as the transpose of a counit, once we know that (5) is an isomorphism; see (7.1) below. We now show, one operation at a time, that these maps are isomorphisms on \(D^b_{\mathrm{ctf}}\).

In the lemmas below it suffices to treat the maps in items (1)–(3) of the list of coefficient maps, that is, all maps except composites. Indeed, extension of scalars along a composite of two maps of the list is the composite of the two extensions. Each map in Definition \hyperref[bp612]{7.4} for the composite is the composite of the maps for the two factors. For (1)–(3) this holds by the coherence of the associativity isomorphisms, and for (4)–(6) it holds because the adjoint transpose of a composite of such maps is the composite of the transposes.

% object: lemma-extension-scalars-preserves-dctf (blueprint/chap-The-six-operations-on-Dbctf/sec-Change-of-coefficients/lemma-extension-scalars-preserves-dctf/_section.md)
\phantomsection\label{bp614}%
\begin{bpstatementit}{Lemma~7.5}{ (Extension of scalars preserves $\Dctf$)}
For each map \(\Lambda\to\Lambda'\) listed above, extension of scalars sends \(D^b_{\mathrm{ctf}}(X,\Lambda)\) into \(D^b_{\mathrm{ctf}}(X,\Lambda')\).
\end{bpstatementit}

% proof: proof/lemma-extension-scalars-preserves-dctf
\begin{proof}[{Proof}]\phantomsection\label{bp615}%
An object that is locally \(\Lambda_S\otimes P\) on a stratum \(S\), with \(P\) perfect, remains of this form after derived extension of scalars, since \(P\otimes^{\mathbf L}_\Lambda\Lambda'\) is a perfect \(\Lambda'\)-complex. Moreover, extension of scalars commutes with extension by zero and with finite cones. Every object of \(D^b_{\mathrm{ctf}}(X,\Lambda)\) is built from such objects by extension by zero from strata and finite cones, and so its extension of scalars lies in \(D^b_{\mathrm{ctf}}(X,\Lambda')\).
\end{proof}

% object: lemma-change-coefficients-direct-image (blueprint/chap-The-six-operations-on-Dbctf/sec-Change-of-coefficients/lemma-change-coefficients-direct-image/_section.md)
\phantomsection\label{bp616}%
\begin{bpstatementit}{Lemma~7.6}{ (Change of coefficients for direct image)}
For each map \(\Lambda\to\Lambda'\) listed above, every morphism \(f:X\to Y\) of schemes of finite type over \(k\) and every \(K\in D^b_{\mathrm{ctf}}(X,\Lambda)\), the map \((Rf_*K)_{\Lambda'}\to Rf_*(K_{\Lambda'})\) of Definition definition-change-of-coefficients-maps is an isomorphism.
\end{bpstatementit}

% proof: proof/lemma-change-coefficients-direct-image
\begin{proof}[{Proof}]\phantomsection\label{bp617}%
The argument depends on the kind of map \(\Lambda\to\Lambda'\). For maps of finite rings we use a projection formula, whereas in the other cases extension of scalars is either a finite construction (finite sums and a cone), with which every exact functor commutes, or a filtered colimit, with which \(Rf_*\) commutes on uniformly bounded below systems.

\emph{Maps \(A\to B\) of finite rings.} We compute on the étale site using Proposition \hyperref[bp454]{5.15}, which identifies the pro-étale and étale categories in this case, and we test the comparison on derived sections over quasi-compact étale objects. For a quasi-compact étale \(V\to Y\), we apply the projection formula for constant coefficients [Stacks, \href{https://stacks.math.columbia.edu/tag/0F12}{0F12}] on both \(V\) and \(X_V=X\times_YV\) and obtain

\[
R\Gamma(V,(Rf_*K)_B)
\simeq R\Gamma(X_V,K)\otimes_A^{\mathbf L}B
\simeq R\Gamma(X_V,K_B).
\]

Here the first identification uses the projection formula on \(V\) and then \(R\Gamma(V,Rf_*K)=R\Gamma(X_V,K)\), and the second is the projection formula on \(X_V\). This projection formula needs only cohomological bounds for \(\ell\)-primary coefficients, which hold for schemes of finite type over \(k\), and so it applies to every finite \(A\)-algebra \(B\). Since these isomorphisms commute with restriction in \(V\), the comparison map is an isomorphism in \(D(Y_{\acute{e}t},B)\), and hence, by Proposition \hyperref[bp454]{5.15}, in \(D(Y_{\mathrm{pro}\acute{e}t},B)\).

\emph{Maps \(\mathcal O_E\to A\) with \(A\) finite.} Since \(\mathcal O_E\) is a discrete valuation ring, \(A\) has a finite free \(\mathcal O_E\)-resolution of length at most one. Its condensed version is again exact. Indeed, a continuous map from a profinite set to the finite quotient \(A\) is constant on the pieces of a finite clopen partition, and hence lifts to a continuous map to the free module. The pullback of this resolution to \(X\) is a finite free resolution of \(A_X\) over \(\mathcal O_{E,X}\). Consequently, tensoring with \(A_X\) is computed by finite sums and a cone, and these commute with the exact functor \(Rf_*\).

\emph{Finite extensions.} For \(\mathcal O_E\to\mathcal O_F\) and \(E\to F\), the ring \(\mathcal O_F\) is free of rank \([F:E]\) over \(\mathcal O_E\). Hence \(K_{\mathcal O_F}\) is, as an \(\mathcal O_E\)-complex, a finite direct sum of copies of \(K\). The assertion follows because \(Rf_*\) commutes with finite sums.

\emph{The maps \(\mathcal O_E\to E\) and \(E\to\overline{\mathbf Q}_\ell\).} These are not finite, and we use continuity. For \(\mathcal O_E\to E\), the complex \(K_E\) is the colimit of \(K\xrightarrow{\ell}K\xrightarrow{\ell}\cdots\), and for \(E\to\overline{\mathbf Q}_\ell\), the complex \(K_{\overline{\mathbf Q}_\ell}\) is the filtered colimit of the \(K_F\) over finite extensions \(F/E\) in \(\overline{\mathbf Q}_\ell\). Since \(K\) is bounded, all terms of these systems are bounded below by a common integer. Hence the colimits commute with \(Rf_*\) by Corollary \hyperref[bp86]{2.29}.
\end{proof}

% object: lemma-change-coefficients-compact-support (blueprint/chap-The-six-operations-on-Dbctf/sec-Change-of-coefficients/lemma-change-coefficients-compact-support/_section.md)
\phantomsection\label{bp618}%
\begin{bpstatementit}{Lemma~7.7}{ (Change of coefficients for compactly supported direct image)}
For each map \(\Lambda\to\Lambda'\) listed above, every separated morphism \(f:X\to Y\) of finite type and every \(K\in D^b_{\mathrm{ctf}}(X,\Lambda)\), the map \((Rf_!K)_{\Lambda'}\to Rf_!(K_{\Lambda'})\) of Definition definition-change-of-coefficients-maps is an isomorphism, and it does not depend on the compactification.
\end{bpstatementit}

% proof: proof/lemma-change-coefficients-compact-support
\begin{proof}[{Proof}]\phantomsection\label{bp619}%
The map is the composite of the isomorphism of Lemma \hyperref[bp616]{7.6} for \(p\), applied to \(j_!K\in D^b_{\mathrm{ctf}}\) (Proposition \hyperref[bp482]{5.22}), and of \(Rp_*\) of the isomorphism of Definition definition-change-of-coefficients-maps. For independence, let \(h:Z'\to Z\) be a refinement of compactifications. Both comparison maps (4) and (2) commute with the refinement map (5.7). Indeed, they are formed from identifications of pullbacks and adjunction maps. A map out of an extension by zero is determined by its restriction to the open subscheme, and there all maps involved are identities. We also use that (4) is compatible with composition of morphisms. Under \(R(ph)_*\simeq Rp_*Rh_*\), the map (4) for \(ph\) is \(Rp_*\) of (4) for \(h\) followed by (4) for \(p\), since transposes of composites are composites. As in the proof of Proposition \hyperref[bp492]{5.26}, this gives independence for any two compactifications.
\end{proof}

% text: prose-07-06
\phantomsection\label{bp620}%
We can now define the comparison map for \(f^!\). Let \(f:X\to Y\) be separated of finite type and \(K\in D^b_{\mathrm{ctf}}(Y,\Lambda)\). Then \(f^!K\in D^b_{\mathrm{ctf}}(X,\Lambda)\) by Theorem \hyperref[bp413]{5.2}, and we let

\[
\beta_{f,K}:(f^!K)_{\Lambda'}\longrightarrow f^!(K_{\Lambda'})
\tag{7.1}
\]

be the transpose, under the adjunction \(Rf_!\dashv f^!\) on \(\Loc(D^b_{\mathrm{ctf}}(-,\Lambda'))\), of the morphism

\[
Rf_!((f^!K)_{\Lambda'})
\xrightarrow{\sim}(Rf_!f^!K)_{\Lambda'}
\xrightarrow{(\varepsilon_K)_{\Lambda'}}K_{\Lambda'},
\]

whose first arrow is the inverse of the isomorphism of Lemma \hyperref[bp618]{7.7}.

The proof below uses three properties of \(\beta\). First, the maps \(\beta\) are compatible with composition: under \((gf)^!\simeq f^!g^!\) we have \(\beta_{gf,K}=f^!(\beta_{g,K})\circ\beta_{f,g^!K}\). This holds because the counit of a composite adjunction is the composite of the counits, and because the comparison for \(Rf_!\) is compatible with the composition isomorphisms of Proposition \hyperref[bp496]{5.28}; both are determined by their restrictions to the open subscheme. Second, let \(\Lambda\to\Lambda'\to\Lambda''\) be a composite of two maps of the list. Then \(\beta\) for the composite is \(\beta\) for \(\Lambda'\to\Lambda''\) composed with the extension of scalars to \(\Lambda''\) of \(\beta\) for \(\Lambda\to\Lambda'\). Indeed, both are transposes of the composite of the extended counits and the comparisons for \(Rf_!\). Third, for an open immersion \(u\), under \(u^!\simeq u^*\), the map \(\beta_{u,K}\) is the isomorphism of Definition definition-change-of-coefficients-maps. Both are transposes of the same map under \(u_!\dashv u^*\).

% object: lemma-change-coefficients-exceptional-inverse-image (blueprint/chap-The-six-operations-on-Dbctf/sec-Change-of-coefficients/lemma-change-coefficients-exceptional-inverse-image/_section.md)
\phantomsection\label{bp621}%
\begin{bpstatementit}{Lemma~7.8}{ (Change of coefficients for the exceptional inverse image)}
For each map \(\Lambda\to\Lambda'\) listed above, every separated morphism \(f:X\to Y\) of finite type and every \(K\in D^b_{\mathrm{ctf}}(Y,\Lambda)\), the map \(\beta_{f,K}\) of (7.1) is an isomorphism.
\end{bpstatementit}

% proof: proof/lemma-change-coefficients-exceptional-inverse-image
\begin{proof}[{Proof}]\phantomsection\label{bp622}%
Since \(f^!\) is only known through its adjunction, there is no formula into which we could substitute \(K_{\Lambda'}\) directly. Instead we show that \(\beta_{f,K}\) is an isomorphism in two cases where \(f^!\) has been computed: closed immersions and the projection \(\mathbf P^1_Y\to Y\). We then factor a general \(f\) locally into such maps, using that \(\beta\) is compatible with composition and with restriction to opens.

\emph{Closed immersions.} For a closed immersion \(i\) with open complement \(j\), the functor \(i^!\) is determined by the localization triangle

\[
i_*i^!K\longrightarrow K\longrightarrow Rj_*j^*K.
\]

It expresses \(i_*i^!K\) as the fibre of a map between objects whose behaviour under extension of scalars we already know. Extension of scalars commutes with \(i_*\), \(j^*\) and, by Lemma \hyperref[bp616]{7.6}, \(Rj_*\). Hence the extension of scalars of this triangle is the triangle for \(K_{\Lambda'}\) over \(\Lambda'\). Since the first map is the counit, the comparison on fibres is \(i_*\beta_{i,K}\). The comparisons on the other two terms are isomorphisms, and therefore \(i_*\beta_{i,K}\) is one as well. Since \(i_*\) is fully faithful, \(\beta_{i,K}\) is an isomorphism.

\emph{The projective line.} Let \(q:\mathbf P^1_Y\to Y\). Recall that the adjunction \(Rq_*\dashv q^*(1)[2]\) between \(D^b_{\mathrm{ctf}}(\mathbf P^1_Y,\Lambda)\) and \(D^b_{\mathrm{ctf}}(Y,\Lambda)\) is defined by the diagonal class and the trace. Through these two maps, \(q^!\) is identified with \(q^*(1)[2]\). It therefore suffices to know that the two maps commute with all the maps of coefficient rings listed above. This holds by construction (Propositions \hyperref[bp534]{6.7} and \hyperref[bp581]{Proposition 6.22}). Indeed, the classes with finite coefficients are constructed from the Kummer sequence and derived extension of scalars, while the integral classes are their compatible limits, and the classes over fields are obtained from these by extension of scalars. Only one map in the list needs a further remark. For \(\mathcal O_E\to A\), the proof of Lemma \hyperref[bp616]{7.6} shows that \(A_X\) is perfect. Hence the derived tensor product with \(A_X\) commutes with the inverse limits defining the integral classes, and so the resulting classes are exactly the classes with finite coefficients. Consequently the identification \(q^!\simeq q^*(1)[2]\) is compatible with extension of scalars. Under this identification, finally, \(\beta_{q,K}\) becomes the isomorphism of Definition definition-change-of-coefficients-maps for \(q^*\), twisted by \((1)[2]\).

\emph{General \(f\).} The open immersion \(\mathbf A^1_Y\hookrightarrow\mathbf P^1_Y\) is a case of restriction to opens, and so \(\beta\) is an isomorphism for \(\mathbf A^1_Y\to Y\). Since the projection \(\mathbf A^d_Y\to Y\) is a composite of \(d\) projections of this kind, the compatibility of \(\beta\) with composition gives the case of affine space. In general, whether \(\beta_{f,K}\) is an isomorphism may be checked locally on \(X\). Using a finite affine cover, we may therefore assume that \(f\) factors as a closed immersion into an affine space over \(Y\) followed by the projection. Then the two cases already treated and compatibility with composition give the claim.
\end{proof}

% object: lemma-change-coefficients-internal-hom (blueprint/chap-The-six-operations-on-Dbctf/sec-Change-of-coefficients/lemma-change-coefficients-internal-hom/_section.md)
\phantomsection\label{bp623}%
\begin{bpstatementit}{Lemma~7.9}{ (Change of coefficients for internal Hom)}
For each map \(\Lambda\to\Lambda'\) listed above and all \(K,L\in D^b_{\mathrm{ctf}}(X,\Lambda)\), the map \(\cHom(K,L)_{\Lambda'}\to\cHom(K_{\Lambda'},L_{\Lambda'})\) of Definition definition-change-of-coefficients-maps is an isomorphism.
\end{bpstatementit}

% proof: proof/lemma-change-coefficients-internal-hom
\begin{proof}[{Proof}]\phantomsection\label{bp624}%
Internal Hom is a right adjoint. As for \(f^!\), there is no formal reason for it to commute with extension of scalars, and we argue as in the proof of Proposition \hyperref[bp605]{7.1}. The comparison map for \(\cHom(-,L)\) is compatible with triangles in the first argument, so we choose a finite stratification on whose strata \(K\) is locally \(\Lambda_S\otimes P\) with \(P\) perfect. By the five lemma and finite cones we may then reduce to the case \(K=a_!P\), with \(a:S\hookrightarrow X\) locally closed and \(P\) a locally constant perfect complex on \(S\). The identity

\[
\cHom_X(a_!P,L)
\simeq Ra_*(P^\vee\otimes^{\mathbf L}a^!L)
\]

commutes with extension of scalars by Lemmas \hyperref[bp616]{7.6} and \hyperref[bp621]{Lemma 7.8}, the isomorphism of Definition definition-change-of-coefficients-maps, and the comparison for the dual of a perfect complex. For the last of these, recall that \(P\) is locally a finite complex of finite projective modules, and for such modules \(\cHom(P,\Lambda)_{\Lambda'}=\cHom(P_{\Lambda'},\Lambda')\) term by term. Moreover, all of these identifications are defined by the same evaluation pairing. Hence their composite is the map of Definition definition-change-of-coefficients-maps.
\end{proof}

% text: prose-07-07
\phantomsection\label{bp625}%
The lemmas combine into the following statement, which is the form in which change of coefficients is used later.

% object: proposition-six-operations-change-coefficients (blueprint/chap-The-six-operations-on-Dbctf/sec-Change-of-coefficients/proposition-six-operations-change-coefficients/_section.md)
\phantomsection\label{bp626}%
\begin{bpstatementit}{Proposition~7.10}{ (The six operations and change of coefficients)}
Let \(\Lambda\to\Lambda'\) be one of the maps of coefficient rings listed above. Extension of scalars sends \(D^b_{\mathrm{ctf}}(X,\Lambda)\) into \(D^b_{\mathrm{ctf}}(X,\Lambda')\). For morphisms \(f:X\to Y\) of schemes of finite type over \(k\) (separated in the case of \(Rf_!\) and \(f^!\)), the change-of-coefficients maps of Definition \hyperref[bp612]{7.4} and (7.1) for \(f^*\), \(Rf_*\), \(Rf_!\), \(f^!\), \(\otimes^{\mathbf L}\) and \(\cHom\) are isomorphisms on objects of \(D^b_{\mathrm{ctf}}(-,\Lambda)\). For a composite \(\Lambda\to\Lambda'\to\Lambda''\) of two such maps, each of these isomorphisms is the composite of the isomorphisms for \(\Lambda\to\Lambda'\) and \(\Lambda'\to\Lambda''\). These isomorphisms are compatible with the composition isomorphisms \((gf)^*\simeq f^*g^*\), \(R(gf)_*\simeq Rg_*Rf_*\), \(R(gf)_!\simeq Rg_!Rf_!\) and \((gf)^!\simeq f^!g^!\) for composable morphisms \(f\) and \(g\). Under these isomorphisms, the extensions of scalars of the units and counits of the adjunctions \(f^*\dashv Rf_*\), \(Rf_!\dashv f^!\) and \(-\otimes^{\mathbf L}K\dashv\cHom(K,-)\) (the last counit being evaluation) over \(\Lambda\) are the corresponding units and counits over \(\Lambda'\).
\end{bpstatementit}

% proof: proof/proposition-six-operations-change-coefficients
\begin{proof}[{Proof}]\phantomsection\label{bp627}%
The first assertion is Lemma \hyperref[bp614]{7.5}. Next, the maps for \(f^*\) and \(\otimes^{\mathbf L}\) are isomorphisms by definition, and those for \(Rf_*\), \(Rf_!\), \(f^!\) and \(\cHom\) by Lemmas \hyperref[bp616]{7.6}, \hyperref[bp618]{Lemma 7.7}, \hyperref[bp621]{Lemma 7.8} and \hyperref[bp623]{Lemma 7.9}. As for the compatibilities, the one with composites of coefficient maps was explained before Lemma \hyperref[bp614]{7.5}, and for \(f^!\) after (7.1). The compatibility with composition of morphisms, on the other hand, holds for \(f^*\) by the coherence of the associativity isomorphisms for pullbacks. It holds for \(Rf_*\) because adjoint transposes of composites are composites. For \(Rf_!\) the composition isomorphism is built from \(R(qr)_*\simeq Rq_*Rr_*\) and the interchange of an open and a proper direct image. Since the interchange is a map out of an extension by zero, the claim follows from the case of \(Rf_*\) and restriction to the open subscheme. Finally, for \(f^!\) it was shown after (7.1).

For the units and counits, consider one of the three adjunctions \(F\dashv G\) over \(\Lambda\). We write \(F'\dashv G'\) for the same adjunction over \(\Lambda'\), with unit \(\eta'\) and counit \(\varepsilon'\), and write \(M'\) for \(M_{\Lambda'}\). Then the comparison for the left adjoint is an isomorphism \(c_F:F'(M')\simeq(FM)'\). In each case the comparison \(c_G:(GN)'\to G'(N')\) for the right adjoint was defined as the transpose of \((\varepsilon_N)'\circ c_F\). In other words \(\varepsilon'_{N'}\circ F'(c_G)=(\varepsilon_N)'\circ c_F\), which is the statement for the counit. For the unit we must show \(c_G\circ(\eta_M)'=G'(c_F)\circ\eta'_{M'}\) as maps \(M'\to G'((FM)')\). We compare transposes. The transpose of the right side is \(c_F\), whereas the transpose of the left side is \((\varepsilon_{FM})'\circ c_F\circ F'((\eta_M)')\). By naturality of \(c_F\) this equals \((\varepsilon_{FM}\circ F\eta_M)'\circ c_F\), which is \(c_F\) by a triangle identity.
\end{proof}

% text: prose-07-08
\phantomsection\label{bp628}%
Without the condition on Tor amplitude, change of coefficients would not preserve the categories. For instance, let \(X\) be a geometric point. The sheaf \(\mathbf F_\ell\) belongs to \(\Dbc(X,\mathbf Z/\ell^2)\) but not to \(D^b_{\mathrm{ctf}}(X,\mathbf Z/\ell^2)\). Its extension of scalars \(\mathbf F_\ell\otimes^{\mathbf L}_{\mathbf Z/\ell^2}\mathbf F_\ell\) is unbounded.

For coefficients in \(\overline{\mathbf Q}_\ell\), Lemma \hyperref[bp265]{3.23} expresses \(D^b_{\mathrm{ctf}}(X,\overline{\mathbf Q}_\ell)\) in terms of the categories \(D^b_{\mathrm{ctf}}(X,E)\) for finite \(E\). As a statement about ordinary categories, it reads as follows.

% object: proposition-dctf-x-qbar-2-colimit (blueprint/chap-The-six-operations-on-Dbctf/sec-Change-of-coefficients/proposition-dctf-x-qbar-2-colimit/_section.md)
\phantomsection\label{bp629}%
\begin{bpstatementit}{Proposition~7.11}{ ($\Dctf(X,\Qbar)$ as a $2$-colimit)}
Extension of scalars identifies

\[
D^b_{\mathrm{ctf}}(X,\overline{\mathbf Q}_\ell)\simeq\mathop{\text{2-colim}}_{E\subset\overline{\mathbf Q}_\ell}D^b_{\mathrm{ctf}}(X,E)
\]

as ordinary categories. In other words, every object of \(D^b_{\mathrm{ctf}}(X,\overline{\mathbf Q}_\ell)\) is isomorphic to \(K_{\overline{\mathbf Q}_\ell}\) for some finite \(E\) and some \(K\in D^b_{\mathrm{ctf}}(X,E)\), and for \(K,L\in D^b_{\mathrm{ctf}}(X,E)\) extension of scalars induces a bijection

\[
\varinjlim_{F\supset E}\Hom_{D^b_{\mathrm{ctf}}(X,F)}(K_F,L_F)
\xrightarrow{\sim}\Hom_{D^b_{\mathrm{ctf}}(X,\overline{\mathbf Q}_\ell)}(K_{\overline{\mathbf Q}_\ell},L_{\overline{\mathbf Q}_\ell}),
\]

the colimit running over finite extensions \(F\) of \(E\) in \(\overline{\mathbf Q}_\ell\).
\end{bpstatementit}

% proof: proof/proposition-dctf-x-qbar-2-colimit
\begin{proof}[{Proof}]\phantomsection\label{bp630}%
An equivalence with a \(2\)-colimit amounts to two statements, essential surjectivity on objects and a bijection on morphisms. Both are parts of Lemma \hyperref[bp265]{3.23}, and we recall how each is proved before putting them together.

The statement about objects is Lemma lemma-descent-finite-coefficient-field. Its proof first descends the lisse sheaves on the strata, together with their transition matrices. It then descends the connecting morphisms of a finite Postnikov tower, and in this way the complexes themselves descend. The bijection on morphisms is Lemma lemma-descent-finite-coefficient-field with \(r=0\), since \(D^b_{\mathrm{ctf}}\) is a full subcategory. Its proof uses the continuity of Ext out of a constructible complex (Lemma \hyperref[bp260]{3.22}), applied to the filtered colimit \(\QbarX=\varinjlim_F F_X\). Surjectivity says that every morphism descends to a finite stage, and injectivity says that two morphisms with the same image agree at a later finite stage. Now consider the category whose objects are the pairs \((E,K)\) with \(K\in D^b_{\mathrm{ctf}}(X,E)\) and whose Hom sets are these colimits. Together with descent of objects, the statements above say exactly that the functor from this category to \(D^b_{\mathrm{ctf}}(X,\overline{\mathbf Q}_\ell)\) is fully faithful and essentially surjective.
\end{proof}

% text: prose-07-09
\phantomsection\label{bp631}%
In practice the proposition is used as follows. To prove a statement about finitely many objects and morphisms of \(D^b_{\mathrm{ctf}}(X,\overline{\mathbf Q}_\ell)\), we may assume that they are extensions of scalars of objects and morphisms of \(D^b_{\mathrm{ctf}}(X,E)\) for a finite extension \(E\) of \(\mathbf Q_\ell\). Moreover, two morphisms over \(E\) that agree over \(\overline{\mathbf Q}_\ell\) already agree over a finite extension of \(E\).


% object: sec-Reduction-to-finite-coefficients (blueprint/chap-The-six-operations-on-Dbctf/sec-Reduction-to-finite-coefficients/_section.md), content hidden
\subsection*{\texorpdfstring{7.3\quad Reduction to finite coefficients}{7.3 Reduction to finite coefficients}}
\phantomsection\label{bp632}%
\addcontentsline{toc}{subsection}{\texorpdfstring{7.3\quad Reduction to finite coefficients}{7.3 Reduction to finite coefficients}}
% text: prose-07-10
\phantomsection\label{bp633}%
Many later theorems say that a canonical map is an isomorphism. We prove them with finite coefficients and then extend them to all admissible rings in three steps. First, for \(\mathcal O_E\) we reduce modulo \(\pi_E\) and use derived Nakayama. This step needs derived completeness, since \(E\) itself is a nonzero \(\mathcal O_E\)-module with \(E/\pi_E E=0\). Then for \(E\) we invert \(\ell\) on integral models, and for \(\overline{\mathbf Q}_\ell\) we descend to a finite extension. The following proposition makes this argument once and for all. It only asks that the functors \(F\), \(G\) and the map \(\theta\) of the statement commute with change of coefficients, and by Proposition \hyperref[bp626]{7.10} this holds for everything built from the six operations.

% object: proposition-isomorphisms-reduce-finite-coefficients (blueprint/chap-The-six-operations-on-Dbctf/sec-Reduction-to-finite-coefficients/proposition-isomorphisms-reduce-finite-coefficients/_section.md)
\phantomsection\label{bp634}%
\begin{bpstatementit}{Proposition~7.12}{ (Isomorphisms reduce to finite coefficients)}
For each admissible \(\Lambda\) let \(F_\Lambda,G_\Lambda:D^b_{\mathrm{ctf}}(X,\Lambda)\to D^b_{\mathrm{ctf}}(Y,\Lambda)\) be exact functors and \(\theta_\Lambda:F_\Lambda\to G_\Lambda\) a natural transformation of exact functors. Assume that for each coefficient map \(\Lambda\to\Lambda'\) of Proposition \hyperref[bp626]{7.10} there are isomorphisms \(F_\Lambda(K)\otimes\Lambda'\simeq F_{\Lambda'}(K\otimes\Lambda')\) and \(G_\Lambda(K)\otimes\Lambda'\simeq G_{\Lambda'}(K\otimes\Lambda')\) carrying \(\theta_\Lambda(K)\otimes\Lambda'\) to \(\theta_{\Lambda'}(K\otimes\Lambda')\). This holds, by Proposition \hyperref[bp626]{7.10}, whenever \(F\), \(G\) and \(\theta\) are built from the six operations, their units and counits and base change maps, covariantly in \(K\). If \(\theta_A(K)\) is an isomorphism for every finite ring \(A\) and every \(K\in D^b_{\mathrm{ctf}}(X,A)\), then \(\theta_\Lambda(K)\) is an isomorphism for every admissible \(\Lambda\) and every \(K\in D^b_{\mathrm{ctf}}(X,\Lambda)\).
\end{bpstatementit}

% proof: proof/proposition-isomorphisms-reduce-finite-coefficients
\begin{proof}[{Proof}]\phantomsection\label{bp635}%
The three kinds of coefficients are reached from finite ones in three different ways: \(\mathcal O_E\) by completion, \(E\) by inverting \(\ell\), and \(\overline{\mathbf Q}_\ell\) as a colimit over finite extensions. We treat them in this order, each case using the previous one.

Let \(\Lambda=\mathcal O_E\), and put \(C=\Cone(\theta(K))\); we must show \(C=0\). Since the complexes \(F(K)\) and \(G(K)\) lie in \(\Dbc(Y,\mathcal O_E)\), they are derived \(\pi_E\)-complete by Proposition \hyperref[bp226]{3.11}. Moreover, derived complete complexes are closed under cones, and so \(C\) is derived complete too. Derived Nakayama (Lemma \hyperref[bp190]{2.62}) therefore reduces \(C=0\) to \(C/\pi_E=0\). Now reduction modulo \(\pi_E\) is a change of coefficients along \(\mathcal O_E\to\mathcal O_E/\pi_E\). By the compatibility assumed in the statement, it identifies \(\theta(K)/\pi_E\) with \(\theta_{\mathcal O_E/\pi_E}(K/\pi_E)\). This is an isomorphism by the hypothesis for finite rings, applied to \(A=\mathcal O_E/\pi_E\). Thus \(C/\pi_E=0\), and so \(C=0\). In particular, the case \(E=\mathbf Q_\ell\) gives \(\Lambda=\mathbf Z_\ell\).

Let \(\Lambda=E\). Here completeness is of no help, since a complex over \(E\) is killed by reduction modulo \(\pi_E\). Instead we use integral models. Both \(F_E\) and \(G_E\) are exact and \(\theta_E\) is a transformation of exact functors. Hence the objects \(K\) for which \(\theta_E(K)\) is an isomorphism form a triangulated subcategory. By Lemma lemma-integral-models, every object of \(\Dbc(X,E)\) is built by finitely many cones and shifts from objects \(K_0[1/\ell]\) with \(K_0\in\Dbc(X,\mathcal O_E)\). It therefore suffices to treat these. Inverting \(\ell\) is the change of coefficients along \(\mathcal O_E\to E\). By hypothesis, \(\theta_E(K_0[1/\ell])\) is therefore identified with \(\theta_{\mathcal O_E}(K_0)[1/\ell]\). This is an isomorphism by the integral case.

Let \(\Lambda=\overline{\mathbf Q}_\ell\). By Proposition \hyperref[bp629]{7.11}, every \(K\) is isomorphic to \(K_E\otimes_E\overline{\mathbf Q}_\ell\) for some finite \(E\) and some \(K_E\in D^b_{\mathrm{ctf}}(X,E)\). By hypothesis, applied to \(E\to\overline{\mathbf Q}_\ell\), the map \(\theta_{\overline{\mathbf Q}_\ell}(K)\) is identified with \(\theta_E(K_E)\otimes_E\overline{\mathbf Q}_\ell\), which is an isomorphism by the case of \(E\). Again \(E=\mathbf Q_\ell\) is included.
\end{proof}

% text: prose-07-11
\phantomsection\label{bp636}%
After a rewriting, the proposition also applies to some contravariant constructions. For instance, the double-dual map \(K\to\mathbb D\mathbb DK\) is a natural transformation of covariant exact functors. What does not reduce in this way is the construction of an adjoint or of a map. For this reason the cycle class of the diagonal and the adjunction \(Rf_!\dashv f^*(d)[2d]\) for smooth \(f\) must be built as single morphisms with \(\mathcal O_E\)-coefficients (Sections \hyperref[bp575]{6.7} and \bpout{Section 9.5}). Likewise, a trace with values in \(\Lambda\) is not a map of complexes.



% object: chap-Smooth-base-change (blueprint/chap-Smooth-base-change/_section.md), content hidden
\section*{\texorpdfstring{8\quad Smooth base change}{8 Smooth base change}}
\phantomsection\label{bp637}%
\addcontentsline{toc}{section}{\texorpdfstring{8\quad Smooth base change}{8 Smooth base change}}
% text: prose-08-01
\phantomsection\label{bp638}%
To construct the trace classes in the trace formula we need the Künneth formula for ordinary direct images, \(Rf_{1*}M_1\boxtimes Rf_{2*}M_2\simeq R(f_1\times f_2)_*(M_1\boxtimes M_2)\) (Proposition \hyperref[bp700]{8.18}, Section \hyperref[bp685]{8.9}). For \(Rf_!\) this formula follows from proper base change and the projection formula, but for \(Rf_*\) this argument does not work, because \(Rf_*\) does not commute with base change in general. For example, let \(j:U=\mathbf A^1\setminus\{0\}\hookrightarrow\mathbf A^1\) over an algebraically closed field, and let \(n\) be invertible in it. Then the stalk of \(R^1j_*(\mathbf Z/n)\) at \(0\) is \(H^1\) of the punctured strictly henselian disc, which is \(\mathbf Z/n(-1)\neq0\), whereas the fibre of \(j\) over \(0\) is empty, so base change along the closed immersion of \(0\) fails.

The Künneth formula only needs base change along a product projection \(Y\times T\to Y\) (Lemma \hyperref[bp687]{8.13}). The question is whether base change along such a map holds for every \(f\).

In topology it does, at least when \(T\) is a manifold. If \(g\) is a local homeomorphism, base change along \(g\) is restriction to open subsets. If \(g\) is the projection \(Y\times\mathbf R^m\to Y\), then the stalk of \(R^qf'_*\) at \((y,t)\) is the colimit of the cohomology groups of \(f^{-1}(V)\times B\) over small neighbourhoods \(V\times B\) of \((y,t)\), and the contractible ball \(B\) does not change them. For schemes the analogue is less clear: a smooth morphism of schemes is, Zariski locally on the source, étale over some \(\mathbf A^d_Y\), but the affine line has no small contractible neighbourhoods. Nevertheless, the theorem of this chapter, due to Artin (SGA 4, Exposé XVI [\hyperref[bp702]{SGA4}]), says that for coefficients prime to the characteristic the affine line behaves as if it had them. More precisely, let \(g:Y'\to Y\) be smooth, let \(f\) be quasi-compact and quasi-separated, and let \(F\) be a torsion sheaf whose torsion orders are invertible on \(Y\). Then the base change map \(g^*R^qf_*F\to R^qf'_*g'^*F\) is an isomorphism (Theorem \hyperref[bp643]{8.1}). Moreover, on \(D^b_{\mathrm{ctf}}(-,\Lambda)\) the same holds for every admissible \(\Lambda\) (Theorem \hyperref[bp681]{8.11}).

A product projection \(Y\times T\to Y\) is not smooth when \(T\) is singular. What the Künneth formula uses is a consequence of smooth base change, namely base change under extension of the ground field [Stacks, \href{https://stacks.math.columbia.edu/tag/0F1C}{0F1C}].

The essential case is a change of the algebraically closed ground field, and there the key input is a fact about curves. Let \(\Omega/L\) be an extension of algebraically closed fields, \(C\) a smooth affine curve over \(L\), and \(n\) invertible in \(L\). Then pullback \(H^1(C,\mathbf Z/n)\to H^1(C_\Omega,\mathbf Z/n)\) is an isomorphism (Theorem \hyperref[bp340]{4.20}). In other words, a cyclic cover of \(C_\Omega\) of degree \(n\) does not depend on the extra transcendental parameters of \(\Omega\).

For \(p\)-torsion coefficients this invariance fails, and so does smooth base change. Let \(k\) be algebraically closed of characteristic \(p>0\), and let \(f:\mathbf A^1_k\to\operatorname{Spec}k\) and \(g:\mathbf A^1_k\to\operatorname{Spec}k\) be the structure maps, with coordinates \(x\) on the source of \(f\) and \(t\) on the source of \(g\). By the Artin–Schreier sequence, \(H^1(\mathbf A^1_K,\mathbf F_p)\) is the cokernel of \(h\mapsto h^p-h\) on \(K[x]\), for every field \(K\supset k\). At the geometric generic point of the \(t\)-line, the stalk of \(g^*R^1f_*\mathbf F_p\) is \(H^1(\mathbf A^1_k,\mathbf F_p)\), the stalk of \(R^1f'_*\mathbf F_p\) is \(H^1(\mathbf A^1_K,\mathbf F_p)\) with \(K=k(t)^{\mathrm{sep}}\), and the base change map is pullback. We claim that the class of \(tx\) is not in its image. Suppose \(tx=a+h^p-h\) with \(a\in k[x]\) and \(h\in K[x]\), and let \(b_e\) be the coefficient of \(x^{p^e}\) in \(h\). Comparing the coefficients of \(x\) gives \(b_0\in -t+k\), and comparing those of \(x^{p^{e+1}}\) gives \(b_{e+1}\in b_e^p+k\). By induction \(b_e\in -t^{p^e}+k\), so \(b_e\neq0\) for every \(e\geq0\). This is absurd for a polynomial. In geometric terms, the Artin–Schreier cover \(y^p-y=tx\) depends on the parameter \(t\). Hence the hypothesis on the torsion orders cannot be dropped.

% object: sec-The-base-change-map-2 (blueprint/chap-Smooth-base-change/sec-The-base-change-map-2/_section.md), content hidden
\subsection*{\texorpdfstring{8.1\quad The base change map}{8.1 The base change map}}
\phantomsection\label{bp639}%
\addcontentsline{toc}{subsection}{\texorpdfstring{8.1\quad The base change map}{8.1 The base change map}}
% text: prose-08-02
\phantomsection\label{bp640}%
Consider a cartesian square, with \(X'=Y'\times_YX\),

\[
\begin{CD}
X' @>{g'}>> X\\
@V{f'}VV @VV{f}V\\
Y' @>{g}>> Y.
\end{CD}
\tag{8.1}
\]

The base change map \(g^*Rf_*K\to Rf'_*g'^*K\), for \(K\in D(X_{\mathrm{pro}\acute{e}t},\Lambda)\), is the map (4.5) of Section \hyperref[bp290]{4.1}. It comes by adjunction from the counit of \(f^*\dashv Rf_*\), and it is defined for every such square and every complex, with \(g\) smooth or not. If \(f\) is quasi-compact and quasi-separated (qcqs), its stalk at a geometric point \(\bar y'\) of \(Y'\) with image \(\bar y\) in \(Y\) is the pullback map from the cohomology of \(X\times_Y\operatorname{Spec}\mathcal O^{\mathrm{sh}}_{Y,\bar y}\) to that of \(X'\times_{Y'}\operatorname{Spec}\mathcal O^{\mathrm{sh}}_{Y',\bar y'}\).

In smooth base change for torsion sheaves (Theorem \hyperref[bp643]{8.1}) we assume that \(f\) is qcqs and that \(g\) is smooth. In the theorem on \(D^b_{\mathrm{ctf}}\) all schemes are of finite type over \(k\). Every inverse image below is the exact inverse image for the pullback of the coefficient ring, so that \(g^*=Lg^*\) and \(g'^*=Lg'^*\).

The proofs of these theorems transport the base change map in three ways: from the pro-étale to the étale site, where the classical arguments are made; along composites of squares, when a smooth morphism is factored into curves; and along a change of the coefficient ring, to reach \(\ell\)-adic coefficients. Since the map is built from units and counits, it is compatible with all three, by Lemma \hyperref[bp304]{4.7} for the first two and by Lemma \hyperref[bp679]{8.10} in Section \hyperref[bp677]{8.8} below for the third.


% object: sec-Torsion-sheaves (blueprint/chap-Smooth-base-change/sec-Torsion-sheaves/_section.md), content hidden
\subsection*{\texorpdfstring{8.2\quad Torsion sheaves}{8.2 Torsion sheaves}}
\phantomsection\label{bp641}%
\addcontentsline{toc}{subsection}{\texorpdfstring{8.2\quad Torsion sheaves}{8.2 Torsion sheaves}}
% text: prose-08-03
\phantomsection\label{bp642}%
Smooth base change for complexes (Section \hyperref[bp672]{8.7}) and for \(\ell\)-adic coefficients (Section \hyperref[bp677]{8.8}) follows from the case of a single torsion sheaf, which is the next theorem. Its proof leaves schemes of finite type over \(k\): it passes through base schemes such as spectra of fields and of integral closures. The comparison results of Chapter \hyperref[bp3]{2} hold over arbitrary schemes for bounded below complexes, so they apply to these bases as well.

% object: theorem-smooth-base-change-classical-torsion (blueprint/chap-Smooth-base-change/sec-Torsion-sheaves/theorem-smooth-base-change-classical-torsion/_section.md)
\phantomsection\label{bp643}%
\begin{bpstatementit}{Theorem~8.1}{ (Smooth base change for classical torsion sheaves)}
In (8.1), suppose \(g\) is smooth and \(f\) is quasi-compact and quasi-separated. Let \(F\) be a classical torsion abelian sheaf on \(X_{\mathrm{pro}\acute{e}t}\) which is the filtered union of its subsheaves \(F[n]\) for the positive integers \(n\) invertible on \(Y\). Then for every \(q\geq0\) the degree-\(q\) base change map induced by (4.5)

\[
g^*R^qf_*F\longrightarrow R^qf'_*g'^*F
\]

is an isomorphism.
\end{bpstatementit}

% text: prose-08-04
\phantomsection\label{bp644}%
As for proper base change, we transport the problem to the étale site by the comparison theorems of Chapter \hyperref[bp3]{2}. Since \(F\) is classical, \(F=\nu_X^*F_{\acute{e}t}\) for a torsion abelian sheaf \(F_{\acute{e}t}\) on \(X_{\acute{e}t}\). As \(f\) and \(f'\) are qcqs, Theorem \hyperref[bp113]{2.38} and Lemma \hyperref[bp304]{4.7} identify both sides and the base change map with the inverse images of the corresponding étale objects. We may therefore work on the étale sites in the proof below, where \(g^{-1}\) denotes inverse image of étale sheaves.

We follow the second proof of [Stacks, \href{https://stacks.math.columbia.edu/tag/0EYU}{0EYU}], an induction on the cohomological degree \(q\). The idea is to concentrate a first failure of base change in a single sheaf \(R^qh_*\mathbf Z/m\), and then to show that this sheaf vanishes. In degree zero, by contrast, the base change map can be computed directly on sections. Moreover, base change may be checked étale locally and is stable under composition, so that étale coordinates reduce us to the case where \(g\) is a smooth affine curve.

Now suppose that base change holds in degrees below \(q\) but fails in degree \(q\). The failure persists when we shrink the source of the direct image by limit arguments, and also when we make integral changes of the source and of the base. In this way it is concentrated in a single situation, in which the coefficients are constant, equal to \(\mathbf Z/m\), and we have a smooth affine curve \(C\) over the spectrum \(S\) of a normal domain whose fraction field \(L\) is algebraically closed, together with the morphism \(h:C_\Omega\to C\) for an algebraically closed extension \(\Omega\) of \(L\). But in this situation \(R^qh_*\mathbf Z/m=0\). For \(q\geq2\) the reason is that affine curves over an algebraically closed field have no cohomology in degree at least two (Appendix \bpout{D}). For \(q=1\) it follows from the invariance of \(H^1\) of curves and an extension theorem for covers (Section \hyperref[bp666]{8.6}).


% object: sec-Degree-zero (blueprint/chap-Smooth-base-change/sec-Degree-zero/_section.md), content hidden
\subsection*{\texorpdfstring{8.3\quad Degree zero}{8.3 Degree zero}}
\phantomsection\label{bp645}%
\addcontentsline{toc}{subsection}{\texorpdfstring{8.3\quad Degree zero}{8.3 Degree zero}}
% text: prose-08-05
\phantomsection\label{bp646}%
In degree zero the base change map can be computed on sections, even for sheaves of sets. Sections of \(f_*F\) over \(V\) are sections of \(F\) over \(V\times_YX\), and the question is whether pulling back to a \(Y'\)-scheme \(U\) over \(V\) creates new sections. It does not when \(U\to V\) is flat, quasi-compact and quasi-separated with geometrically connected fibres, because pullback along such a map does not change sections of the pullback of an étale sheaf. Moreover, a flat \(g\) locally of finite presentation with geometrically reduced fibres factors locally into such a map followed by an étale one, and in this way Lemma \hyperref[bp647]{8.2} proves the isomorphism for every such \(g\).

% object: lemma-degree-zero-calculation (blueprint/chap-Smooth-base-change/sec-Degree-zero/lemma-degree-zero-calculation/_section.md)
\phantomsection\label{bp647}%
\begin{bpstatementit}{Lemma~8.2}{ (The degree-zero calculation)}
Suppose that \(g\) in (8.1) is flat and locally of finite presentation with geometrically reduced fibres, and that \(f\) is qcqs. Then the base change map \(g^{-1}f_*F_{\acute{e}t}\to f'_*g'^{-1}F_{\acute{e}t}\) is an isomorphism for every sheaf of sets \(F_{\acute{e}t}\) on \(X_{\acute{e}t}\).
\end{bpstatementit}

% proof: proof/lemma-degree-zero-calculation
\begin{proof}[{Proof}]\phantomsection\label{bp648}%
Both sides are sheaves on \(Y'_{\acute{e}t}\), so it is enough to compare their sections on a family of étale objects \(U\to Y'\) that cover \(Y'\). We choose these objects so that both sections can be computed from sections of \(F_{\acute{e}t}\) over a single scheme \(V\times_YX\).

Such objects exist locally on every étale \(U\to Y'\): by [Stacks, \href{https://stacks.math.columbia.edu/tag/0EY5}{0EY5}], locally on \(U\) there is a factorization \(U\to V\to Y\) in which \(V\to Y\) is étale and \(U\to V\) is flat and qcqs with geometrically connected fibres. Pullback along such a map does not change the sections of the pullback of an étale sheaf [Stacks, \href{https://stacks.math.columbia.edu/tag/0A3H}{0A3H}, \href{https://stacks.math.columbia.edu/tag/0EZX}{0EZX}], and we use this twice. Applying it to \(U\to V\), we compute the left side from \((f_*F_{\acute{e}t})(V)\); applying it to the base change \(U\times_YX\to V\times_YX\), which again has geometrically connected fibres, we compute the right side. On these objects we therefore have

\[
(g^{-1}f_*F_{\acute{e}t})(U)
=F_{\acute{e}t}(V\times_Y X)
=(f'_*g'^{-1}F_{\acute{e}t})(U),
\]

and the base change map is bijective on sections over every such \(U\), because these equalities are induced by it. As these \(U\) cover \(Y'\), the lemma follows by descent, as in the argument in degree zero of [Stacks, \href{https://stacks.math.columbia.edu/tag/0EYS}{0EYS}].
\end{proof}


% object: sec-The-induction-hypothesis (blueprint/chap-Smooth-base-change/sec-The-induction-hypothesis/_section.md), content hidden
\subsection*{\texorpdfstring{8.4\quad The induction hypothesis}{8.4 The induction hypothesis}}
\phantomsection\label{bp649}%
\addcontentsline{toc}{subsection}{\texorpdfstring{8.4\quad The induction hypothesis}{8.4 The induction hypothesis}}
% text: prose-08-06
\phantomsection\label{bp650}%
In positive degrees we argue by induction on the degree. This induction hypothesis has to survive the reductions outlined above: composition of base changes, étale localization, and replacement of the source and of the base of the direct image by other schemes. But a statement about the single square (8.1) does not survive them, so instead we fix an integer \(n\) invertible on \(Y\) and use a property of a morphism \(g:Y'\to Y\) which concerns all direct images over all \(Y\)-schemes at once, as in [Stacks, \href{https://stacks.math.columbia.edu/tag/0F02}{0F02}]. For \(r\geq0\), we say that \(\mathrm{BC}(g,n,r)\) holds if, for every scheme \(T\) over \(Y\), every qcqs morphism \(Z\to T\), every sheaf of \(\mathbf Z/n\)-modules \(M\) on \(Z_{\acute{e}t}\), and every \(0\leq i\leq r\), the base change map in degree \(i\) for the square with base change \(Y'\times_YT\to T\) and direct image along \(Z\to T\) is an isomorphism. We allow arbitrary schemes \(T\) because the reductions below lead to base schemes, such as the spectrum of a normal domain with algebraically closed fraction field, that are in general not Noetherian.

The property \(\mathrm{BC}(g,n,r)\) may be checked étale locally on \(Y\) and on \(Y'\). Moreover, for composable morphisms \(g_1\) and \(g_2\), \(\mathrm{BC}(g_1,n,r)\) and \(\mathrm{BC}(g_2,n,r)\) imply \(\mathrm{BC}(g_2\circ g_1,n,r)\). Both statements rest on two properties of the base change map in a fixed degree. First, the map is local for étale coverings of \(Y'\) and of \(Y\), since restriction to an étale object is localization of the site. Second, base change maps compose when two base changes are composed [Stacks, \href{https://stacks.math.columbia.edu/tag/0EZR}{0EZR}, \href{https://stacks.math.columbia.edu/tag/0EZS}{0EZS}].

For smooth \(g\), étale coordinates [Stacks, \href{https://stacks.math.columbia.edu/tag/054L}{054L}] give, locally on \(Y'\) and \(Y\) and with all schemes affine, an étale \(Y\)-morphism \(Y'\to\mathbf A^d_Y\), where \(d\) is the relative dimension of \(g\). Composing it with the projections \(\mathbf A^{j}_Y\to\mathbf A^{j-1}_Y\) that forget the last coordinate, we obtain a factorization

\[
Y'\longrightarrow\mathbf A^{d-1}_Y
\longrightarrow\cdots\longrightarrow\mathbf A^1_Y
\longrightarrow Y
\]

of \(g\) in which every displayed arrow is smooth of relative dimension one. Indeed, the first is an étale map followed by a projection, and the others are projections. The case \(d=0\) is étale localization. Hence, in the induction on the degree, it suffices to treat affine \(Y\) and \(Y'\) with \(g\) smooth of relative dimension one.

The coefficients reduce to \(\mathbf Z/n\)-modules. Indeed, a sheaf \(F\) as in Theorem \hyperref[bp643]{8.1} is the filtered union of the subsheaves \(F[n]\), ordered by divisibility of integers invertible on \(Y\), and inverse image commutes with this colimit, and so do higher direct images along qcqs morphisms [Stacks, \href{https://stacks.math.columbia.edu/tag/09Z1}{09Z1}]. It therefore suffices to treat \(\mathbf Z/n\)-modules for each such \(n\).


% object: sec-Reduction-to-curves (blueprint/chap-Smooth-base-change/sec-Reduction-to-curves/_section.md), content hidden
\subsection*{\texorpdfstring{8.5\quad Reduction to curves}{8.5 Reduction to curves}}
\phantomsection\label{bp651}%
\addcontentsline{toc}{subsection}{\texorpdfstring{8.5\quad Reduction to curves}{8.5 Reduction to curves}}
% text: prose-08-07
\phantomsection\label{bp652}%
Recall that \(\mathrm{BC}(g,n,r)\) asserts base change in degrees \(\leq r\), for \(\mathbf Z/n\)-coefficients, along every base change \(Y'\times_YT\to T\) of \(g\). Assume now that \(q\geq1\) and that \(\mathrm{BC}(g,n,q-1)\) holds. Under this assumption we show that if \(\mathrm{BC}(g,n,q)\) fails, it also fails in a special form: with constant coefficients and with direct image from the spectrum of an algebraically closed field (Lemma \hyperref[bp662]{8.6}). The next four lemmas, all on the étale sites, carry out this reduction step by step.

We argue by contradiction rather than by dévissage, because a nonzero cokernel of the base change map can be transported along an embedding of coefficient sheaves \(M\hookrightarrow N\), whereas an isomorphism cannot.

% object: lemma-injectivity-base-change-map-degree (blueprint/chap-Smooth-base-change/sec-Reduction-to-curves/lemma-injectivity-base-change-map-degree/_section.md)
\phantomsection\label{bp653}%
\begin{bpstatementit}{Lemma~8.3}{ (Injectivity of the base change map in degree $q$)}
Suppose \(q\geq1\) and \(\mathrm{BC}(g,n,q-1)\) holds. For every square obtained by base change as in the definition of \(\mathrm{BC}\) and every sheaf of \(\mathbf Z/n\)-modules \(M\) on \(Z_{\acute{e}t}\), the base change map \(b_M^q:g^{-1}R^qf_*M\to R^qf'_*g'^{-1}M\) in degree \(q\) is injective. An inclusion \(M\hookrightarrow N\) induces

\[
\operatorname{coker}(b_M^q)
\hookrightarrow\operatorname{coker}(b_N^q).
\]

For an embedding \(M\hookrightarrow I\) with \(I\) injective in the category of étale sheaves of \(\mathbf Z/n\)-modules, the cokernel of \(b_M^q\) embeds into \(R^qf'_*g'^{-1}I\).
\end{bpstatementit}

% proof: proof/lemma-injectivity-base-change-map-degree
\begin{proof}[{Proof}]\phantomsection\label{bp654}%
We embed \(M\) into an injective sheaf and compare the long exact sequences of the two sides. The induction hypothesis identifies everything in degree \(q-1\), and since the injective sheaf has no higher direct images upstairs, a diagram chase shows that \(b^q_M\) is injective with cokernel inside \(R^qf'_*g'^{-1}I\).

Choose \(0\to M\to I\to Q\to0\) with \(I\) as in the statement. By the long exact sequences, the rows of the following diagram are exact. Here the upper row ends in \(0\) because \(R^qf_*I=0\) for the injective sheaf \(I\), whereas the lower row does not, because \(g'^{-1}I\) need not be injective.

\[
\begin{CD}
g^{-1}R^{q-1}f_*I @>>> g^{-1}R^{q-1}f_*Q
@>>> g^{-1}R^qf_*M @>>>0\\
@V{\sim}VV @V{\sim}VV @VVV @VVV\\
R^{q-1}f'_*g'^{-1}I @>>> R^{q-1}f'_*g'^{-1}Q
@>>>R^qf'_*g'^{-1}M @>>>R^qf'_*g'^{-1}I.
\end{CD}
\]

The first two vertical arrows are isomorphisms by \(\mathrm{BC}(g,n,q-1)\), applied to \(I\) and to \(Q\). Exactness of the upper row says that \(g^{-1}R^qf_*M\) is the cokernel of the map \(g^{-1}R^{q-1}f_*I\to g^{-1}R^{q-1}f_*Q\), and exactness of the lower row says that the kernel of \(R^qf'_*g'^{-1}M\to R^qf'_*g'^{-1}I\) is the cokernel of the corresponding map in the lower row. Since the first two vertical isomorphisms identify these two maps, \(b_M^q\) is an isomorphism onto this kernel. In particular \(b^q_M\) is injective, and its cokernel embeds into \(R^qf'_*g'^{-1}I\), which is the last assertion. For \(M\subset N\), choose an injective \(I\) containing \(N\). What we have just shown, applied to \(M\subset I\) and to \(N\subset I\), embeds the cokernels of \(b_M^q\) and \(b_N^q\) compatibly into \(R^qf'_*g'^{-1}I\). As in [Stacks, \href{https://stacks.math.columbia.edu/tag/0F03}{0F03}], the map between the cokernels induced by \(M\subset N\) is compatible with these two embeddings and is therefore injective.
\end{proof}

% text: prose-08-08
\phantomsection\label{bp655}%
We use the lemma in the form of its second assertion: enlarging the coefficient sheaf does not lose a failure. In the next lemma this is how we pass from a scheme to an open subscheme or a closed complement.

% object: lemma-reduction-field (blueprint/chap-Smooth-base-change/sec-Reduction-to-curves/lemma-reduction-field/_section.md)
\phantomsection\label{bp656}%
\begin{bpstatementit}{Lemma~8.4}{ (Reduction to a field)}
Let \(q\geq1\). If \(\mathrm{BC}(g,n,q-1)\) holds but \(\mathrm{BC}(g,n,q)\) fails, then some failure occurs for an affine base change \(T\to Y\), direct image along a morphism \(\operatorname{Spec}\Omega\to T\) with \(\Omega\) a field, and a sheaf of \(\mathbf Z/n\)-modules on \(\operatorname{Spec}\Omega\).
\end{bpstatementit}

% proof: proof/lemma-reduction-field
\begin{proof}[{Proof}]\phantomsection\label{bp657}%
We shrink the source of the direct image, keeping a fixed nonzero element of the cokernel, until it becomes the spectrum of a field. A minimal closed subscheme carrying the failure exists by a limit argument, and its generic point still carries it. The one nonformal step is the third property below, which says that a failure on a scheme is seen either on an open subscheme or on its closed complement, and whose proof uses the injectivity of the base change map.

We first fix notation for the cokernels. Let \(T\) be a \(Y\)-scheme, put \(T'=Y'\times_YT\), and let \(g_T:T'\to T\) be the projection. For a qcqs morphism \(h:Z'\to T\) we write \(h':Z'_{T'}=Z'\times_TT'\to T'\) and \(g_{Z'}:Z'_{T'}\to Z'\) for its base changes, and for a sheaf \(N\) of \(\mathbf Z/n\)-modules on \(Z'_{\acute{e}t}\) we put

\[
b^q_{h,N}:g_T^{-1}R^qh_*N\longrightarrow R^qh'_*g_{Z'}^{-1}N,
\qquad
C(h,N)=\operatorname{coker}\bigl(b^q_{h,N}\bigr),
\]

a sheaf on \(T'_{\acute{e}t}\). These cokernels are functorial in two ways. A map of sheaves \(N\to N_1\) on \(Z'_{\acute{e}t}\) induces \(C(h,N)\to C(h,N_1)\), and a morphism \(u:Z''\to Z'\) induces \(C(h,N)\to C(hu,u^{-1}N)\), by pulling back along \(u\): upstairs through \(N\to u_*u^{-1}N\) and the edge map \(R^qh_*u_*\to R^q(hu)_*\), and downstairs in the same way. These pullback maps are compatible with the base change maps, and with composition in \(u\), because base change maps compose [Stacks, \href{https://stacks.math.columbia.edu/tag/0EZR}{0EZR}, \href{https://stacks.math.columbia.edu/tag/0EZS}{0EZS}].

Since \(\mathrm{BC}(g,n,q)\) fails, there are a \(Y\)-scheme \(T\), a qcqs morphism \(f:Z\to T\) and a sheaf \(M\) of \(\mathbf Z/n\)-modules on \(Z_{\acute{e}t}\) such that \(b^q_{f,M}\) is not an isomorphism. By Lemma \hyperref[bp653]{8.3} it is injective, so \(C(f,M)\neq0\), and we fix a geometric point \(\bar t'\) of \(T'\) and a nonzero element \(\xi\in C(f,M)_{\bar t'}\). For a qcqs morphism \(u:Z'\to Z\) let \(P(Z')\) mean that the image \(\xi_{Z'}\) of \(\xi\) in \(C(fu,u^{-1}M)_{\bar t'}\) is nonzero. The verification in [Stacks, \href{https://stacks.math.columbia.edu/tag/0F07}{0F07}] shows that \(P\) has three properties.

\begin{enumerate}
\item thickenings preserve \(P\), by topological invariance;
\item limits with affine transition maps preserve \(P\), by continuity;
\item if \(P(Z')\) holds and \(V\subset Z'\) is a quasi-compact open with closed complement \(W\), then \(P(V)\) or \(P(W)\) holds.
\end{enumerate}

To prove the third property, write \(h=fu:Z'\to T\) and \(M'=u^{-1}M\), and let \(j:V\to Z'\) and \(i:W\to Z'\) be the inclusions. We split \(M'\) into an open part and a closed part. Choose an injective sheaf \(J\) of \(\mathbf Z/n\)-modules on \(V_{\acute{e}t}\) containing \(j^{-1}M'\). Then the map

\[
M'\longrightarrow j_*J\oplus i_*i^{-1}M',
\]

whose components are adjoint to \(j^{-1}M'\subset J\) and to the identity of \(i^{-1}M'\), is injective. Indeed, at a geometric point of \(V\) its first component is the inclusion \(M'_{\bar v}\subset J_{\bar v}\), and at a geometric point of \(W\) its second component is the identity of \(M'_{\bar w}\). By the second assertion of Lemma \hyperref[bp653]{8.3}, the induced map

\[
C(h,M')\longrightarrow C(h,j_*J)\oplus C(h,i_*i^{-1}M')
\]

is injective as well. Since \(P(Z')\) holds, \(\xi_{Z'}\neq0\), and therefore its image in \(C(h,j_*J)_{\bar t'}\) or its image in \(C(h,i_*i^{-1}M')_{\bar t'}\) is nonzero. We show that the second case gives \(P(W)\) and the first gives \(P(V)\).

The closed summand is formal. The base change \(i':W_{T'}\to Z'_{T'}\) of \(i\) is again a closed immersion, and closed direct images are exact and commute with every base change (Lemma \hyperref[bp364]{4.27}). Hence, for a sheaf \(N\) on \(W_{\acute{e}t}\),

\[
R^qh_*(i_*N)=R^q(hi)_*N,
\qquad
R^qh'_*\bigl(g_{Z'}^{-1}i_*N\bigr)=R^qh'_*\bigl(i'_*g_W^{-1}N\bigr)
=R^q(hi)'_*g_W^{-1}N,
\]

and these identifications are compatible with the base change maps [Stacks, \href{https://stacks.math.columbia.edu/tag/0EZR}{0EZR}, \href{https://stacks.math.columbia.edu/tag/0EZS}{0EZS}]. Thus \(C(h,i_*N)=C(hi,N)\). For \(N=i^{-1}M'\) the image of \(\xi_{Z'}\) in \(C(h,i_*i^{-1}M')_{\bar t'}\) becomes \(\xi_W\), and so it is nonzero exactly when \(P(W)\) holds.

The open summand needs the induction hypothesis, because \(j\) is not finite. Let \(j':V_{T'}\to Z'_{T'}\) and \(g_V:V_{T'}\to V\) be the base changes of \(j\), and compare the Leray spectral sequences for \(hj=h\circ j\) and \(h'j'=h'\circ j'\). Upstairs, \(R^bj_*J=0\) for \(b>0\) because \(J\) is injective, so the Leray spectral sequence collapses and gives

\[
R^qh_*(j_*J)\xrightarrow{\ \sim\ }R^q(hj)_*J .
\]

Downstairs we use \(\mathrm{BC}(g,n,q-1)\). The scheme \(Z'\) is a \(Y\)-scheme through \(T\), and \(Z'_{T'}=Y'\times_YZ'\), so the square formed by \(j\), \(j'\), \(g_{Z'}\) and \(g_V\) is one of the squares in the definition of \(\mathrm{BC}\). Hence

\[
g_{Z'}^{-1}R^bj_*J\xrightarrow{\ \sim\ }R^bj'_*g_V^{-1}J
\qquad(0\le b\le q-1).
\]

For \(b=0\) this identifies \(g_{Z'}^{-1}j_*J\) with \(j'_*g_V^{-1}J\), and for \(0<b<q\) it shows that \(R^bj'_*g_V^{-1}J=0\). Consequently, in the Leray spectral sequence

\[
E_2^{a,b}=R^ah'_*R^bj'_*g_V^{-1}J\ \Longrightarrow\
R^{a+b}(h'j')_*g_V^{-1}J
\]

the rows \(0<b<q\) vanish. No differential can reach \(E^{q,0}\): the differential \(d_r\) into \(E_r^{q,0}\) starts at \(E_r^{q-r,r-1}\), which is zero for \(2\le r\le q\) because then \(0<r-1<q\), and for \(r>q\) because then \(q-r<0\). So \(E_2^{q,0}=E_\infty^{q,0}\), which is the bottom step of the filtration of the abutment, and the edge map

\[
e:R^qh'_*\bigl(j'_*g_V^{-1}J\bigr)\longrightarrow R^q(h'j')_*g_V^{-1}J
\]

is injective. Since base change maps compose [Stacks, \href{https://stacks.math.columbia.edu/tag/0EZR}{0EZR}, \href{https://stacks.math.columbia.edu/tag/0EZS}{0EZS}] and edge maps are natural, the square

\[
\begin{CD}
g_T^{-1}R^qh_*(j_*J) @>{b^q_{h,j_*J}}>> R^qh'_*\bigl(j'_*g_V^{-1}J\bigr)\\
@V{\wr}VV @VV{e}V\\
g_T^{-1}R^q(hj)_*J @>{b^q_{hj,J}}>> R^q(h'j')_*g_V^{-1}J
\end{CD}
\]

commutes. Its left arrow is an isomorphism and its right arrow is injective, so it induces an injective map \(C(h,j_*J)\to C(hj,J)\). (In the top row we have identified \(g_{Z'}^{-1}j_*J\) with \(j'_*g_V^{-1}J\) by the case \(b=0\); this does not change the cokernel.) The vertical arrows are the pullback maps along \(j\), so this map is pullback along \(j\) followed by \(j^{-1}j_*J=J\).

It remains to pass from \(J\) back to \(j^{-1}M'\). By naturality, the composite \(C(h,M')\to C(h,j_*J)\to C(hj,J)\) equals the composite \(C(h,M')\to C(hj,j^{-1}M')\to C(hj,J)\) of pullback along \(j\) and the map induced by \(j^{-1}M'\subset J\). In the first case the image of \(\xi_{Z'}\) in \(C(h,j_*J)_{\bar t'}\) is nonzero, hence so is its image in \(C(hj,J)_{\bar t'}\). This image is the image of \(\xi_V\) under the second composite, and so \(\xi_V\neq0\), which is \(P(V)\).

The field comes out of the three properties as follows. The intersection of a decreasing chain of closed subschemes is the limit of the chain, with affine transition maps. By the second property and Zorn's lemma there is therefore a minimal closed subscheme with \(P\). It is nonempty, since \(P\) fails for the empty scheme, and by the first property it is reduced. Let \(\xi\) be the generic point of an irreducible component. We claim that every affine open neighbourhood of \(\xi\) has \(P\). Indeed, by the third property, either the neighbourhood or its closed complement has \(P\), and the complement is a closed subscheme not containing \(\xi\), hence strictly smaller, so by minimality it does not have \(P\). Since the minimal closed subscheme is reduced, the limit of these neighbourhoods is \(\operatorname{Spec}\kappa(\xi)\). It has \(P\) by the second property (compare [Stacks, \href{https://stacks.math.columbia.edu/tag/0F06}{0F06}]), and so \(\kappa(\xi)\) is the field asserted in the lemma.
\end{proof}

% text: prose-08-09
\phantomsection\label{bp658}%
The field \(\Omega\) produced by this lemma is arbitrary, and the base \(T\) is still an arbitrary affine scheme. To bring them into the special situation of Lemma \hyperref[bp662]{8.6} we want to replace \(\Omega\) by its algebraic closure and \(T\) by a normal scheme. Both replacements are integral morphisms, and the next lemma shows that a failure survives them.

% object: lemma-integral-changes-source-base (blueprint/chap-Smooth-base-change/sec-Reduction-to-curves/lemma-integral-changes-source-base/_section.md)
\phantomsection\label{bp659}%
\begin{bpstatementit}{Lemma~8.5}{ (Integral changes of source and base)}
Assume \(\mathrm{BC}(g,n,q-1)\). A failure of base change in degree \(q\) persists after an integral surjective change of the source of the morphism along which we take direct image. It also persists after an integral change of the base through which this morphism factors.
\end{bpstatementit}

% proof: proof/lemma-integral-changes-source-base
\begin{proof}[{Proof}]\phantomsection\label{bp660}%
Both assertions are [Stacks, \href{https://stacks.math.columbia.edu/tag/0F04}{0F04}, \href{https://stacks.math.columbia.edu/tag/0F05}{0F05}]. They use three properties of direct image along an integral morphism \(\pi\): it is exact and conservative [Stacks, \href{https://stacks.math.columbia.edu/tag/04C2}{04C2}], even if \(\pi\) is not surjective, and it commutes with arbitrary base change [Stacks, \href{https://stacks.math.columbia.edu/tag/0EYP}{0EYP}]. By exactness, the Leray spectral sequence for a composite with \(\pi\) degenerates, which means that direct images along the composite are computed by applying \(\pi_*\) first.

We use the notation \(b^q_{h,N}\) and \(C(h,N)=\operatorname{coker}b^q_{h,N}\) of the proof of Lemma \hyperref[bp656]{8.4}, for a \(Y\)-scheme \(T\), its base change \(g_T:T'=Y'\times_YT\to T\), a qcqs morphism \(h:Z\to T\) and a sheaf \(N\) of \(\mathbf Z/n\)-modules on \(Z_{\acute{e}t}\). By Lemma \hyperref[bp653]{8.3}, \(b^q_{h,N}\) is injective, so a failure of base change in degree \(q\) means \(C(h,N)\neq0\).

For the first assertion, let \(f:Z\to T\) and \(M\) give a failure, so \(C(f,M)\neq0\), and let \(\pi:\widetilde Z\to Z\) be integral and surjective. We show that \(C(f\pi,\pi^{-1}M)\neq0\). Since \(\pi\) is surjective, the unit \(M\to\pi_*\pi^{-1}M\) is injective, as one sees on stalks. By the second assertion of Lemma \hyperref[bp653]{8.3} it induces an injection \(C(f,M)\hookrightarrow C(f,\pi_*\pi^{-1}M)\), and so \(C(f,\pi_*\pi^{-1}M)\neq0\). It remains to identify \(C(f,\pi_*\pi^{-1}M)\) with \(C(f\pi,\pi^{-1}M)\). Let \(\pi':\widetilde Z_{T'}\to Z_{T'}\) be the base change of \(\pi\), again integral. Upstairs, \(\pi_*\) is exact, so the Leray spectral sequence for \(f\circ\pi\) degenerates and gives \(R^qf_*(\pi_*\pi^{-1}M)=R^q(f\pi)_*\pi^{-1}M\). Downstairs, \(\pi_*\) commutes with base change, \(g_Z^{-1}\pi_*\pi^{-1}M\simeq\pi'_*g_{\widetilde Z}^{-1}\pi^{-1}M\), and \(\pi'_*\) is exact, so in the same way

\[
R^qf'_*\bigl(g_Z^{-1}\pi_*\pi^{-1}M\bigr)
\simeq R^qf'_*\pi'_*\bigl(g_{\widetilde Z}^{-1}\pi^{-1}M\bigr)
=R^q(f\pi)'_*\bigl(g_{\widetilde Z}^{-1}\pi^{-1}M\bigr).
\]

These identifications are compatible with the base change maps [Stacks, \href{https://stacks.math.columbia.edu/tag/0EZR}{0EZR}, \href{https://stacks.math.columbia.edu/tag/0EZS}{0EZS}], so they carry \(b^q_{f,\pi_*\pi^{-1}M}\) to \(b^q_{f\pi,\pi^{-1}M}\). Hence \(C(f\pi,\pi^{-1}M)\simeq C(f,\pi_*\pi^{-1}M)\neq0\), and base change fails for the new source.

For the second assertion, let \(f:Z\to T\) factor as \(f=\pi\circ f_1\) with \(\pi:T_1\to T\) integral and \(f_1:Z\to T_1\), and suppose base change fails for \(f\) and \(M\), that is, \(C(f,M)\neq0\). We show that it fails for \(f_1\) and \(M\) over the new base \(T_1\). Put \(T'_1=Y'\times_YT_1=T'\times_TT_1\), and let \(\pi':T'_1\to T'\) be the base change of \(\pi\), again integral. The base change map for \(f_1\) is

\[
b^q_{f_1,M}:g_{T_1}^{-1}R^qf_{1*}M\longrightarrow R^qf'_{1*}g_Z^{-1}M
\qquad\text{on }T'_{1,\acute{e}t}.
\]

We push it down to \(T'\) by \(\pi'_*\). Since \(\pi'_*\) is exact and \(\pi_*\) commutes with base change, \(\pi'_*g_{T_1}^{-1}R^qf_{1*}M\simeq g_T^{-1}\pi_*R^qf_{1*}M
=g_T^{-1}R^qf_*M\), where the last equality is the degenerate Leray spectral sequence for \(f=\pi f_1\). In the same way \(\pi'_*R^qf'_{1*}g_Z^{-1}M=R^qf'_*g_Z^{-1}M\). Under these identifications \(\pi'_*\bigl(b^q_{f_1,M}\bigr)=b^q_{f,M}\), again because base change maps compose. Since \(\pi'_*\) is exact, it follows that \(\pi'_*C(f_1,M)\simeq C(f,M)\neq0\), so \(C(f_1,M)\neq0\), which is the assertion.
\end{proof}

% text: prose-08-10
\phantomsection\label{bp661}%
With the two previous lemmas we may take the source of the direct image to be the spectrum of an algebraically closed field. Its cohomology vanishes in positive degrees, so the source of the base change map is zero and the failure sits in the target, a direct image along an extension of the ground field.

% object: lemma-reduction-curve-over-normal-base (blueprint/chap-Smooth-base-change/sec-Reduction-to-curves/lemma-reduction-curve-over-normal-base/_section.md)
\phantomsection\label{bp662}%
\begin{bpstatementit}{Lemma~8.6}{ (Reduction to a curve over a normal base)}
Suppose \(g:Y'\to Y\) is affine and smooth of relative dimension one, \(n\) is invertible on \(Y\), and \(\mathrm{BC}(g,n,q-1)\) holds. If \(\mathrm{BC}(g,n,q)\) fails, there are an affine scheme \(S=\operatorname{Spec}R\) over \(Y\) with \(R\) a normal domain whose fraction field \(L\) is algebraically closed, an algebraically closed extension field \(\Omega\) of \(L\), an integer \(m\mid n\), and, for the smooth affine curve \(C=Y'\times_YS\), the cartesian square

\[
\begin{CD}
C_\Omega @>{h}>> C\\
@VVV @VVV\\
\operatorname{Spec}\Omega @>>> S
\end{CD}
\]

such that

\[
R^qh_*\underline{\mathbf Z/m}\ne0.
\]
\end{bpstatementit}

% text: prose-08-11
\phantomsection\label{bp663}%
Here \(\underline{\mathbf Z/m}\) is the classical constant sheaf.

% proof: proof/lemma-reduction-curve-over-normal-base
\begin{proof}[{Proof of Lemma 8.6}]\phantomsection\label{bp664}%
We first simplify the source of the direct image and the coefficients, then show that the source of the base change map vanishes, and finally replace the base by a normal one.

By Lemma \hyperref[bp656]{8.4} there is a failure with direct image from the spectrum of a field \(\Omega\). Since the map from the spectrum of an algebraic closure of \(\Omega\) to \(\operatorname{Spec}\Omega\) is integral and surjective, Lemma \hyperref[bp659]{8.5} allows us to replace \(\Omega\) by an algebraic closure. On this geometric point, a sheaf of \(\mathbf Z/n\)-modules is a \(\mathbf Z/n\)-module, hence a direct sum of cyclic groups \(\mathbf Z/m\) with \(m\mid n\). As the base change map commutes with this direct sum by continuity of qcqs direct images, base change still fails for some summand \(\mathbf Z/m\).

The source of the base change map is the pullback of \(R^qf_*\underline{\mathbf Z/m}\) for \(f:\operatorname{Spec}\Omega\to T\), and we show that it is zero. The stalk of this sheaf at a geometric point \(\bar t\) of \(T\) is the cohomology of \(\operatorname{Spec}\Omega\times_T\operatorname{Spec}\mathcal O^{\mathrm{sh}}_{T,\bar t}\). This scheme is a limit of étale \(\Omega\)-schemes, which, since \(\Omega\) is algebraically closed, are disjoint unions of copies of \(\operatorname{Spec}\Omega\). By continuity of cohomology the stalk therefore vanishes for \(q\geq1\). As the source of the base change map is zero, the failure means that the target \(R^qh_*\underline{\mathbf Z/m}\) is nonzero.

The extension theorem for covers will require a normal base with algebraically closed fraction field. If the base is \(\operatorname{Spec}B\), we replace it by \(\operatorname{Spec}R\), where \(R\) is the integral closure in \(\Omega\) of the image of \(B\). This is allowed, because the morphism \(\operatorname{Spec}\Omega\to\operatorname{Spec}B\) factors through \(\operatorname{Spec}R\), and \(\operatorname{Spec}R\to\operatorname{Spec}B\) is integral, so the failure persists by Lemma \hyperref[bp659]{8.5}. Moreover, the fraction field \(L\) of \(R\) is the algebraic closure of \(\operatorname{Frac}(\operatorname{im}(B\to \Omega))\) in \(\Omega\), hence algebraically closed, although in general \(\Omega/L\) is transcendental. We are now in the situation of [Stacks, \href{https://stacks.math.columbia.edu/tag/0F08}{0F08}].
\end{proof}

% object: example-reduction-end-affine-line (blueprint/chap-Smooth-base-change/sec-Reduction-to-curves/example-reduction-end-affine-line/_section.md)
\phantomsection\label{bp665}%
\begin{bpstatement}{Example~8.7}{ (Where the reduction ends)}
The square of the lemma is itself an instance of smooth base change, and the simplest case shows what remains to be proved. Let \(Y=S=\operatorname{Spec}L\) with \(L\) algebraically closed, let \(g:Y'=C=\mathbf A^1_L\to Y\) be the structure map, and take direct image along \(f:X=\operatorname{Spec}\Omega\to Y\) for an algebraically closed extension \(\Omega\) of \(L\), so that \(X'=\mathbf A^1_\Omega\) and \(f'=h\). The sheaf \(R^qf_*\mathbf Z/m\) on \(\operatorname{Spec}L\) is \(H^q(\operatorname{Spec}\Omega,\mathbf Z/m)\), which vanishes for \(q\geq1\). So smooth base change for this square says exactly that \(R^qh_*\mathbf Z/m=0\) for \(q\geq1\). For \(m\) invertible in \(L\) this is the content of Proposition \hyperref[bp668]{8.8} below.

Let \(q=1\) and \(m\) invertible in \(L\). The sections over \(C\) itself are no problem: choosing \(\mu_m\simeq\mathbf Z/m\), the Kummer sequence gives \(H^1(\mathbf A^1_\Omega,\mathbf Z/m)=\Omega^*/(\Omega^*)^m=0\), since \(\Gamma(\mathbf A^1_\Omega,\mathcal O)^*=\Omega^*\) is divisible and \(\operatorname{Pic}(\mathbf A^1_\Omega)=0\). However, the presheaf whose sheafification is \(R^1h_*\mathbf Z/m\) has other sections. For example, on the open \(U=\mathbf G_{m,L}\) of \(C\) the group \(H^1(\mathbf G_{m,\Omega},\mathbf Z/m)\) is cyclic of order \(m\), and such a class has to be killed by an étale covering of \(U\), as we shall see after the statement of the proposition.

For \(m=p\) in characteristic \(p\) the argument breaks down already over \(C\) itself. Indeed, the group \(H^1(\mathbf A^1_\Omega,\mathbf F_p)\) is the cokernel of \(u\mapsto u^p-u\) on \(\Omega[x]\), and for \(\omega\in\Omega\) transcendental over \(L\) it contains the class of \(\omega x\), which does not come from \(C\) (compare the class of \(tx\) at the beginning of the chapter).
\end{bpstatement}


% object: sec-Change-of-base-field (blueprint/chap-Smooth-base-change/sec-Change-of-base-field/_section.md), content hidden
\subsection*{\texorpdfstring{8.6\quad Change of base field}{8.6 Change of base field}}
\phantomsection\label{bp666}%
\addcontentsline{toc}{subsection}{\texorpdfstring{8.6\quad Change of base field}{8.6 Change of base field}}
% text: prose-08-12
\phantomsection\label{bp667}%
Lemma \hyperref[bp662]{8.6} turned a failure of \(\mathrm{BC}(g,n,q)\) into a single nonvanishing sheaf \(R^qh_*\mathbf Z/m\ne0\), where \(h:C_\Omega\to C\) is the base change of a smooth affine curve \(C\) over a normal base \(S\) to a larger algebraically closed field \(\Omega\). Proposition \hyperref[bp668]{8.8} below shows that this sheaf vanishes, which contradicts the lemma and finishes the induction. As the sheaf \(R^qh_*\mathbf Z/m\) is the sheafification of \(U\mapsto H^q(U_\Omega,\mathbf Z/m)\) on \(C_{\acute{e}t}\), we must show that classes on \(U_\Omega\) die after étale localization on \(U\). In degrees \(q\geq2\) there is nothing to kill, since affine curves over an algebraically closed field have cohomological dimension one.

In degree one the problem is about covers. Recall that \(S=\operatorname{Spec}R\) with \(R\) a normal domain whose fraction field \(L\) is algebraically closed, and \(L\subset\Omega\) (Lemma \hyperref[bp662]{8.6}). Over the generic point of \(S\), a class on \(C_\Omega\) already comes from \(C_L\), by the invariance of \(H^1\) of a curve under extension of algebraically closed fields (Theorem \hyperref[bp340]{4.20}). In other words, it is given by a \(\mathbf Z/m\)-torsor over \(C_L\), so it suffices to find an étale cover of \(C\) on which this torsor becomes trivial. To find one, we take a connected component of the torsor and extend it over \(C\) by the extension theorem for Galois covers of order invertible on \(S\). This theorem needs \(S\) normal and a section of \(C\to S\). The section is available: since \(L\) is algebraically closed, the generic fibre \(C_L\) has a rational point, and after étale localization on \(S\) this point extends to a section.

We use that \(m\) is invertible in two places, and both fail for \(p\)-torsion. First, the invariance of \(H^1\) fails: the Artin–Schreier class of \(tx\) on \(\mathbf A^1_{k(t)^{\mathrm{sep}}}\) does not come from \(\mathbf A^1_k\). Second, the extension theorem can fail for covers of degree divisible by a residue characteristic, because such covers may be wildly ramified along the special fibres.

% object: proposition-base-change-larger-algebraically-closed (blueprint/chap-Smooth-base-change/sec-Change-of-base-field/proposition-base-change-larger-algebraically-closed/_section.md)
\phantomsection\label{bp668}%
\begin{bpstatementit}{Proposition~8.8}{ (Base change to a larger algebraically closed field)}
Let \(S=\operatorname{Spec}R\) for a normal domain \(R\) with algebraically closed fraction field \(L\), and let \(a:C\to S\) be affine and smooth of relative dimension one. Let \(\Omega/L\) be an extension of algebraically closed fields, let \(m\) be a positive integer invertible on \(S\), and let \(h:C_\Omega\to C\) be the projection. Then

\[
R^qh_*\underline{\mathbf Z/m}=0\qquad(q\geq1)
\]

on \(C_{\mathrm{pro}\acute{e}t}\).
\end{bpstatementit}

% text: prose-08-13
\phantomsection\label{bp669}%
The simplest case already shows how the proof goes. Take \(C=\mathbf G_{m,S}\) with coordinate \(x\). Then \(H^1(\mathbf G_{m,\Omega},\mathbf Z/m)\) is cyclic of order \(m\), generated by the class of the Kummer cover \(y^m=x\). Here we view this cover as a \(\mathbf Z/m\)-torsor through a primitive \(m\)-th root of unity, which lies in \(R\) because \(R\) is normal and \(L\) is algebraically closed. The point is that the cover is already defined over \(\mathbf G_{m,S}\), and it is finite étale there because \(m\) is invertible on \(S\), so the class becomes trivial after pullback to it. The proof below produces such a cover for an arbitrary class.

% proof: proof/proposition-base-change-larger-algebraically-closed
\begin{proof}[{Proof of Proposition 8.8}]\phantomsection\label{bp670}%
We describe \(R^qh_*\) as a sheafification and show that each local class dies after a suitable étale covering. In degrees at least two this is the cohomological dimension of affine curves. In degree one we first descend the class to \(C_L\), then trivialize it on a finite étale Galois cover of \(C_L\), and finally extend this cover over \(C\).

\emph{Step 1: reduction to local classes.} We may work on the étale site, because \(h\) is qcqs and direct images can be compared (Theorem \hyperref[bp113]{2.38}). There, [Stacks, \href{https://stacks.math.columbia.edu/tag/03Q8}{03Q8}] identifies \(R^qh_*\underline{\mathbf Z/m}\) with the sheafification of the presheaf

\[
U\longmapsto H^q(U\times_S\operatorname{Spec}\Omega,
\mathbf Z/m),
\qquad U\to C\text{ \'etale},
\]

so we must show that every such class vanishes on an étale covering of \(U\).

\emph{Step 2: degrees \(q\geq2\).} In these degrees the groups themselves vanish for affine \(U\), which we may assume because the question is local on \(U\). Indeed, \(U_\Omega\) is then affine and étale over \(C_\Omega\). Hence it is empty or a finite disjoint union of smooth affine curves over the algebraically closed field \(\Omega\), and \((C_2)\) of Theorem \hyperref[bp340]{4.20} kills its cohomology in degrees \(q\geq2\).

\emph{Step 3: a section and connected fibres.} From now on \(q=1\). The extension theorem we want to use requires a section and geometrically connected fibres. The map \(U\to S\) need not have them, but we may arrange them locally. By [Stacks, \href{https://stacks.math.columbia.edu/tag/0EY4}{0EY4}], after an étale covering of \(U\) the map \(U\to S\) factors as \(U\to V\to S\) with \(U\) and \(V\) affine, \(V\to S\) étale, and \(U\to V\) smooth with geometrically connected fibres and a section. Moreover, the separated étale map \(V\to S\) is a finite disjoint union of open immersions [Stacks, \href{https://stacks.math.columbia.edu/tag/09Z9}{09Z9}], because \(S\) is normal and integral with separably closed function field. As we only have to kill classes on \(U_\Omega\), we may therefore replace \(S\) by one of these opens and \(C\) by the part of \(U\) over it, which does not change the fraction field \(L\). Afterwards there is a section \(s:S\to C\) and the fibres are geometrically connected.

\emph{Step 4: descent to \(C_L\) and a trivializing cover.} Let \(\xi\in H^1(C_\Omega,\mathbf Z/m)\). It is the pullback of a unique class \(\xi_L\in H^1(C_L,\mathbf Z/m)\), because the pullback \(H^1(C_L,\mathbf Z/m)\to H^1(C_\Omega,\mathbf Z/m)\) is an isomorphism by \((C_6)\) of Theorem \hyperref[bp340]{4.20}, applied to the affine curve \(C_L\) as in [Stacks, \href{https://stacks.math.columbia.edu/tag/03SC}{03SC}]; this step uses that \(m\) is prime to the characteristic. Next, the classification of \(H^1\) by torsors [Stacks, \href{https://stacks.math.columbia.edu/tag/03AJ}{03AJ}] and Lemma \hyperref[bp123]{2.41} show that \(\xi_L\) is the class of a finite étale \(\mathbf Z/m\)-torsor \(P\to C_L\). Choose a connected component \(Z\) of \(P\) and let \(G\subseteq\mathbf Z/m\) be its stabilizer. As the curve \(C_L\) is connected, \(Z\to C_L\) is a finite étale Galois cover with group \(G\). The class \(\xi_L\) vanishes on \(Z\), because the pullback of \(P\) to \(Z\) has the diagonal section.

\emph{Step 5: extending the cover over \(C\).} We use the extension theorem for generic Galois covers [Stacks, \href{https://stacks.math.columbia.edu/tag/0EZJ}{0EZJ}]: let \(S\) be integral, normal and qcqs with generic point \(\eta\), and let \(C\to S\) be qcqs and smooth with geometrically connected fibres and a section \(s\). If \(Z\to C_\eta\) is a finite étale Galois cover with group \(G\) of order invertible on \(S\), and \(Z\) has a \(\kappa(\eta)\)-rational point over \(s(\eta)\), then \(Z\) extends to a finite étale Galois cover \(\widetilde Z\to C\) with group \(G\). The idea is that a cover of order prime to the residue characteristics ramifies only tamely, and splitting over the section then rules out ramification along the irreducible special fibres, and at points of codimension at least two the cover extends by purity of the branch locus.

The hypotheses hold with \(C_\eta=C_L\). Indeed, the order \(|G|\) divides \(m\) and is therefore invertible on \(S\). Moreover, the fibre of \(Z\) over \(s(\eta)\) is nonempty, as \(Z\to C_L\) is finite étale and surjective, and it has an \(L\)-rational point because \(L\) is algebraically closed. So \(\widetilde Z\) exists.

We check that \(\widetilde Z\to C\) is an étale covering, that is, that it is surjective. The scheme \(C\) is connected, since \(C\to S\) is open and surjective with connected fibres over a connected base. Hence the image of \(\widetilde Z\) in \(C\), which is open and closed and contains the generic fibre, is all of \(C\). Finally, \(\xi\) dies on this covering. Indeed, \(\widetilde Z\times_C C_\Omega=Z\times_{C_L}C_\Omega\), so the pullback of \(\xi\) to \(\widetilde Z\times_CC_\Omega\) is the pullback of \(\xi_L|_Z=0\).
\end{proof}

% proof: proof/theorem-smooth-base-change-classical-torsion
\begin{proof}[{Proof of Theorem 8.1}]\phantomsection\label{bp671}%
We argue by induction on \(q\). The degree-zero calculation starts the induction. By the reductions above, a hypothetical failure in degree \(q\) gives a nonzero sheaf \(R^qh_*\underline{\mathbf Z/m}\), and Proposition \hyperref[bp668]{8.8} says that this sheaf is zero.

A smooth morphism is flat, locally of finite presentation, and has geometrically reduced fibres, and so the case of degree zero is Lemma \hyperref[bp647]{8.2}. Let \(q\geq1\) and assume the theorem in all degrees less than \(q\). Fix \(n\) invertible on \(Y\). Since smoothness is stable under base change, we may apply the induction hypothesis after every base change \(T\to Y\), which shows that \(\mathrm{BC}(g,n,q-1)\) holds. By étale coordinates and composition of base change maps, it suffices to treat the case where \(g\) is affine and smooth of relative dimension one.

If \(\mathrm{BC}(g,n,q)\) failed, Lemma \hyperref[bp662]{8.6} would give \(S\), \(C\), \(h\) and \(m\) with \(R^qh_*\underline{\mathbf Z/m}\neq0\). These satisfy the hypotheses of Proposition \hyperref[bp668]{8.8}: by construction \(S\) is the spectrum of a normal domain with algebraically closed fraction field, \(m\) is invertible on \(S\) because it divides \(n\), and \(C\) is a smooth affine curve over \(S\). Therefore the conclusion \(R^qh_*\underline{\mathbf Z/m}=0\) of the proposition is a contradiction. Hence \(\mathrm{BC}(g,n,q)\) holds for affine \(g\) of relative dimension one, and by the reduction above it holds for every smooth \(g\).

Since \(f\) and \(f'\) are qcqs, both sides commute with the filtered union \(F=\bigcup_nF[n]\), which gives the theorem for \(F\) on the étale sites. Finally, Theorem \hyperref[bp113]{2.38} transfers the isomorphism and the base change map from the étale sites to \(Y'_{\mathrm{pro}\acute{e}t}\).
\end{proof}


% object: sec-Bounded-below-complexes (blueprint/chap-Smooth-base-change/sec-Bounded-below-complexes/_section.md), content hidden
\subsection*{\texorpdfstring{8.7\quad Bounded below complexes}{8.7 Bounded below complexes}}
\phantomsection\label{bp672}%
\addcontentsline{toc}{subsection}{\texorpdfstring{8.7\quad Bounded below complexes}{8.7 Bounded below complexes}}
% text: prose-08-14
\phantomsection\label{bp673}%
Theorem \hyperref[bp643]{8.1} concerns a single torsion sheaf, but its applications concern complexes. Smooth base change on \(D^b_{\mathrm{ctf}}\) (Theorem \hyperref[bp681]{8.11}) is reduced, for integral coefficients, to complexes with finite coefficients, and the Künneth formula (Section \hyperref[bp685]{8.9}) is a statement about complexes. So we need smooth base change for complexes of torsion sheaves, and bounded below complexes suffice, since objects of \(D^b_{\mathrm{ctf}}\) are bounded.

% object: theorem-smooth-base-change-classical-torsion-2 (blueprint/chap-Smooth-base-change/sec-Bounded-below-complexes/theorem-smooth-base-change-classical-torsion-2/_section.md)
\phantomsection\label{bp674}%
\begin{bpstatementit}{Theorem~8.9}{ (Smooth base change for classical torsion complexes)}
In (8.1), let \(g\) be smooth and \(f\) quasi-compact and quasi-separated. Let \(\Lambda\) be a discrete commutative ring annihilated by an integer \(n\) invertible on \(Y\). For every

\[
K\in D^+_{\mathrm{class}}(X_{\mathrm{pro}\acute{e}t},\Lambda)
\]

the canonical map (4.5) is an isomorphism.
\end{bpstatementit}

% proof: proof/theorem-smooth-base-change-classical-torsion-2
\begin{proof}[{Proof}]\phantomsection\label{bp675}%
We first treat a single sheaf, where the only point is \(\Lambda\)-linearity, and then pass to complexes by comparing hypercohomology spectral sequences.

For a classical sheaf the assertion is Theorem \hyperref[bp643]{8.1}, applied to the underlying abelian sheaf, which is killed by \(n\). Since the base change map is constructed from units and counits, it is \(\Lambda\)-linear, and so it is an isomorphism of \(\Lambda\)-module sheaves.

For a complex, we first note that by Theorem \hyperref[bp113]{2.38} and Corollary \hyperref[bp109]{2.37} the direct and inverse images involved are classical, and that Lemma \hyperref[bp304]{4.7} identifies the comparison with (4.5). So we may argue on the étale sites. Choose an integer \(m\) with \(K\in D^{\geq m}\). Since both \(g^*\) and \(g'^*\) are exact, they commute with taking cohomology sheaves, and the two sides of (4.5) are computed by the hypercohomology spectral sequences for \(Rf_*\) and \(Rf'_*\), applied to \(K\) and to \(g'^*K\). These spectral sequences have \(E_2\) terms

\[
g^*R^pf_*H^q(K),\qquad R^pf'_*g'^*H^q(K),
\qquad p\geq0,\quad q\geq m,
\]

and the base change map induces on them the comparison for sheaves, which is an isomorphism by the preceding paragraph. Since in each total degree only finitely many pairs \((p,q)\) occur, these shifted first quadrant spectral sequences converge; this is where we use that \(K\) is bounded below. A map of convergent spectral sequences that is an isomorphism on \(E_2\) is an isomorphism on abutments, so the base change map induces isomorphisms \(g^*R^if_*K\to R^if'_*g'^*K\) for all \(i\), and so it is an isomorphism in the derived category.
\end{proof}

% text: prose-08-15
\phantomsection\label{bp676}%
For the theorem on \(D^b_{\mathrm{ctf}}\) we only need the case of a finite admissible ring \(A\), where the hypothesis holds with \(n=\ell^r\) for \(\ell^rA=0\).


% object: sec-ell-adic-coefficients-3 (blueprint/chap-Smooth-base-change/sec-ell-adic-coefficients-3/_section.md), content hidden
\subsection*{\texorpdfstring{8.8\quad $\ell$-adic coefficients}{8.8 ell-adic coefficients}}
\phantomsection\label{bp677}%
\addcontentsline{toc}{subsection}{\texorpdfstring{8.8\quad $\ell$-adic coefficients}{8.8 ell-adic coefficients}}
% text: prose-08-16
\phantomsection\label{bp678}%
For every admissible coefficient ring, smooth base change takes the following form. For finite rings it is a case of smooth base change for torsion complexes (Theorem \hyperref[bp674]{8.9}), and we deduce the other rings by Proposition \hyperref[bp634]{7.12}.

On the pro-étale site the base change map for \(\mathcal O_E\)-coefficients is a single map of complexes, defined by the same adjunction as for finite rings. Lemma \hyperref[bp679]{8.10} identifies its reduction modulo \(\pi_E\) with the base change map for \(\mathcal O_E/\pi_E\), so derived Nakayama gives the integral case from the torsion theorem. The rational cases then follow by inverting \(\ell\) and by descent to a finite extension.

% object: lemma-change-coefficients (blueprint/chap-Smooth-base-change/sec-ell-adic-coefficients-3/lemma-change-coefficients/_section.md)
\phantomsection\label{bp679}%
\begin{bpstatementit}{Lemma~8.10}{ (Change of coefficients)}
Suppose that the schemes in (8.1) are of finite type over \(k\). Let \(\Lambda\to\Lambda'\) be a coefficient map of Proposition \hyperref[bp626]{7.10} and let \(K\in D^b_{\mathrm{ctf}}(X,\Lambda)\). Then the square

\[
\begin{CD}
(g^*Rf_*K)_{\Lambda'} @>>> (Rf'_*g'^*K)_{\Lambda'}\\
@VVV @VVV\\
g^*Rf_*(K_{\Lambda'}) @>>> Rf'_*g'^*(K_{\Lambda'})
\end{CD}
\]

commutes. Here the upper arrow is the extension of scalars of the base change map (4.5) for \(K\), the lower arrow is the base change map for \(K_{\Lambda'}\), and the vertical arrows are the composites of the change-of-coefficients isomorphisms of Proposition \hyperref[bp626]{7.10} for \(g^*\) and \(Rf_*\), respectively for \(Rf'_*\) and \(g'^*\).
\end{bpstatementit}

% proof: proof/lemma-change-coefficients
\begin{proof}[{Proof}]\phantomsection\label{bp680}%
The base change map is built from units, counits and composition isomorphisms for inverse images. By Proposition \hyperref[bp626]{7.10} the change-of-coefficients maps are compatible with all of them, and so the square commutes.
\end{proof}

% object: theorem-smooth-base-change-all-admissible (blueprint/chap-Smooth-base-change/sec-ell-adic-coefficients-3/theorem-smooth-base-change-all-admissible/_section.md)
\phantomsection\label{bp681}%
\begin{bpstatementit}{Theorem~8.11}{ (Smooth base change for all admissible coefficients)}
In (8.1), assume that all schemes are of finite type over \(k\) and that \(g\) is smooth. For every admissible coefficient ring

\[
\Lambda\in\{A,\mathbf Z_\ell,\mathcal O_E,\mathbf Q_\ell,E,\overline{\mathbf Q}_\ell\}
\]

and every \(K\in D^b_{\mathrm{ctf}}(X,\Lambda)\), the base change map (4.5) is an isomorphism

\[
g^*Rf_*K\xrightarrow{\sim}Rf'_*g'^*K.
\tag{8.2}
\]
\end{bpstatementit}

% text: prose-08-17
\phantomsection\label{bp682}%
For integral and rational coefficients we have \(D^b_{\mathrm{ctf}}=\Dbc\), so in those cases the theorem concerns all bounded constructible complexes. The isomorphisms (8.2) respect composition of cartesian squares and change of coefficients, because the base change map itself does (Lemmas \hyperref[bp304]{4.7} and \hyperref[bp679]{Lemma 8.10}).

% proof: proof/theorem-smooth-base-change-all-admissible
\begin{proof}[{Proof of Theorem 8.11}]\phantomsection\label{bp683}%
For finite rings the theorem is the torsion theorem once we know that \(K\) is classical, and the other rings follow by the reduction principle. By Proposition \hyperref[bp605]{7.1}, both sides lie in \(D^b_{\mathrm{ctf}}(Y',\Lambda)\), so we only have to show that the canonical map (4.5) is invertible.

\emph{Finite rings \(A\).} The torsion theorem is stated for classical complexes, so we first check that \(K\) is one. By Theorem \hyperref[bp129]{2.43}, constructible pro-étale sheaves with discrete Noetherian coefficients are classical, and since the cohomology sheaves of \(K\) are constructible, \(K\) is a bounded classical complex. Since \(A\) is killed by a power of \(\ell\), which is invertible on \(Y\), smooth base change for torsion complexes (Theorem \hyperref[bp674]{8.9}) applies to \(K\).

\emph{Other rings.} We apply Proposition \hyperref[bp634]{7.12} with \(F_\Lambda=g^*Rf_*\), \(G_\Lambda=Rf'_*g'^*\) and \(\theta_\Lambda\) the base change map. Both functors are exact, as the proposition requires, and the base change map commutes with change of coefficients, which is Lemma \hyperref[bp679]{8.10}. Since \(\theta_A\) is an isomorphism for every finite \(A\) by the first case, the proposition shows that \(\theta_\Lambda\) is an isomorphism for every admissible \(\Lambda\).
\end{proof}

% object: remark-smooth-base-change-proper-base (blueprint/chap-Smooth-base-change/sec-ell-adic-coefficients-3/remark-smooth-base-change-proper-base/_section.md)
\phantomsection\label{bp684}%
\begin{bpstatement}{Remark~8.12}{ (Smooth base change, proper base change and base change for $f^!$)}
In (8.1), smooth base change assumes that \(g\) is smooth and allows any \(f\) of finite type, whereas proper base change assumes that \(f\) is proper and allows any \(g\). Smooth purity (Chapter \bpout{9}) gives a third base change isomorphism \(g'^*f^!\simeq f'^!g^*\), for the exceptional inverse image, when \(f\) is smooth and \(g\) is arbitrary. By contrast, neither smooth nor proper base change applies to the open immersion \(j:\mathbf G_m\hookrightarrow\mathbf A^1\) of the beginning of this chapter, where base change along the immersion of the origin fails: Theorem \hyperref[bp681]{8.11} does not apply because the immersion of the origin is not smooth, and proper base change does not apply because \(j\) is not proper.
\end{bpstatement}


% object: sec-The-Kunneth-formula (blueprint/chap-Smooth-base-change/sec-The-Kunneth-formula/_section.md), content hidden
\subsection*{\texorpdfstring{8.9\quad The Künneth formula}{8.9 The Künneth formula}}
\phantomsection\label{bp685}%
\addcontentsline{toc}{subsection}{\texorpdfstring{8.9\quad The Künneth formula}{8.9 The Künneth formula}}
% text: prose-08-18
\phantomsection\label{bp686}%
The trace formula needs the Künneth formula for ordinary direct images (Proposition \hyperref[bp700]{8.18}), which is a base change statement for \(Rf_*\). Base change for \(Rf_*\) fails in general, as the open immersion \(j:\mathbf G_m\hookrightarrow\mathbf A^1\) of the beginning of the chapter shows. Along a product projection \(Y\times T\to Y\), however, base change holds (Lemma \hyperref[bp687]{8.13}). To prove this we compactify, after which only an open immersion remains, and for an open immersion the statement comes down to two results: base change along extensions of the ground field, a consequence of smooth base change, and the Künneth formula with one proper factor, a consequence of proper base change. In this section \(A\) is a finite admissible ring, and we write \(M\boxtimes N=p_1^*M\otimes^{\mathbf L}p_2^*N\) for the external tensor product.

We start with the base change statement, which is the case of the Künneth formula in which the second coefficient is constant.

% object: lemma-base-change-along-product-projection (blueprint/chap-Smooth-base-change/sec-The-Kunneth-formula/lemma-base-change-along-product-projection/_section.md)
\phantomsection\label{bp687}%
\begin{bpstatementit}{Lemma~8.13}{ (Base change along a product projection)}
Let \(f:X\to Y\) be a morphism of separated \(k\)-schemes of finite type, let \(T\) be a separated \(k\)-scheme of finite type, and let \(M\in\Dbc(X,A)\).  Let \(p_X:X\times T\to X\) and \(p_Y:Y\times T\to Y\) be the projections.  Then the base change map

\[
p_Y^*Rf_*M\longrightarrow R(f\times\operatorname{id}_T)_*p_X^*M
\]

is an isomorphism.
\end{bpstatementit}

% proof: proof/lemma-base-change-along-product-projection
\begin{proof}[{Proof}]\phantomsection\label{bp688}%
The idea is to compactify \(f\), so that only an open immersion remains, and for an open immersion base change along a product projection holds even when \(T\) is singular. We argue on the étale site, where \(M\) is a bounded complex with torsion coefficients, and transport the result to the pro-étale site by Theorem \hyperref[bp113]{2.38}.

Choose a compactification \(f=\bar f j\), with \(j:X\hookrightarrow\bar X\) open and \(\bar f:\bar X\to Y\) proper (Theorem \bpout{A.2}). Since base change maps compose, it suffices to prove the lemma for \(j\) and for \(\bar f\) separately. For the open immersion \(j\), the base change map \(p_{\bar X}^*Rj_*M\to R(j\times\operatorname{id}_T)_*p_X^*M\) is an isomorphism by [Stacks, \href{https://stacks.math.columbia.edu/tag/0F1H}{0F1H}], applied to the cohomology sheaves of \(M\) and combined with the bounded-below spectral sequence. For \(\bar f\), proper base change, applied to \(Rj_*M\), shows that

\[
p_Y^*R\bar f_*Rj_*M\longrightarrow
R(\bar f\times\operatorname{id}_T)_*p_{\bar X}^*Rj_*M
\]

is an isomorphism. By Lemma \hyperref[bp299]{4.5} the composite of the two base change maps is the base change map for \(f\).

The step for \(j\) does not depend on Proposition \hyperref[bp700]{8.18} below, so there is no circularity. Indeed, the proof of [Stacks, \href{https://stacks.math.columbia.edu/tag/0F1H}{0F1H}] has two inputs. The first is base change under extension of the ground field [Stacks, \href{https://stacks.math.columbia.edu/tag/0F1C}{0F1C}], which by [Stacks, \href{https://stacks.math.columbia.edu/tag/0F09}{0F09}] follows from smooth base change (Theorem \hyperref[bp674]{8.9}), because after passing to perfect closures the larger field is a filtered union of smooth algebras over the smaller one. The second input, [Stacks, \href{https://stacks.math.columbia.edu/tag/0F1G}{0F1G}], is a Künneth formula with one proper factor, which comes from proper base change and the projection formula.
\end{proof}

% text: prose-08-19
\phantomsection\label{bp689}%
For \(f=j:\mathbf G_m\hookrightarrow\mathbf A^1\) and \(M=A\), the lemma says that the stalk of \(R(j\times\operatorname{id}_T)_*A\) at a point \((0,t)\) has the same cohomology as the stalk of \(Rj_*A\) at the origin, namely \(A\) in degree \(0\) and \(A(-1)\) in degree \(1\), whatever \(T\) is. The Künneth formula allows an arbitrary constructible second factor in place of the constant one, and we first define the map it concerns.

% object: definition-kunneth-map (blueprint/chap-Smooth-base-change/sec-The-Kunneth-formula/definition-kunneth-map/_section.md)
\phantomsection\label{bp690}%
\begin{bpstatement}{Definition~8.14}{ (Künneth map)}
Let \(f_i:X_i\to Y_i\), \(i=1,2\), be morphisms of separated \(k\)-schemes of finite type, put \(h=f_1\times f_2:X_1\times X_2\to Y_1\times Y_2\), and let \(M_i\in D^b_{\mathrm{ctf}}(X_i,A)\). The \emph{Künneth map}

\[
Rf_{1*}M_1\boxtimes Rf_{2*}M_2
\longrightarrow Rh_*(M_1\boxtimes M_2)
\tag{8.3}
\]

is the adjoint transpose, for \(h^*\dashv Rh_*\), of the morphism

\[
h^*(Rf_{1*}M_1\boxtimes Rf_{2*}M_2)
\simeq f_1^*Rf_{1*}M_1\boxtimes f_2^*Rf_{2*}M_2
\xrightarrow{\ \varepsilon_{f_1}\boxtimes\varepsilon_{f_2}\ }
M_1\boxtimes M_2,
\]

where \(\varepsilon_{f_i}\) is the counit of \(f_i^*\dashv Rf_{i*}\) and the first isomorphism comes from \(h^*p_i^*\simeq p_i^*f_i^*\).
\end{bpstatement}

% text: prose-08-20
\phantomsection\label{bp691}%
The case we prove first is the one in which one factor is an identity. For \(f:X\to Y\), \(T\) as in Lemma \hyperref[bp687]{8.13}, \(M\in D^b_{\mathrm{ctf}}(X,A)\) and \(Q\in D^b_{\mathrm{ctf}}(T,A)\), with \(f_1=f\) and \(f_2=\operatorname{id}_T\), put \(h=f\times\operatorname{id}_T\). The Künneth map is then

\[
\kappa_{M,Q}:Rf_*M\boxtimes Q
\longrightarrow Rh_*(M\boxtimes Q),
\tag{8.4}
\]

the adjoint transpose of \(h^*(Rf_*M\boxtimes Q)=p_X^*f^*Rf_*M\otimes^{\mathbf L} p_T^*Q
\xrightarrow{\varepsilon_f\otimes^{\mathbf L}1}M\boxtimes Q\). For constant \(Q\) it is the base change map of Lemma \hyperref[bp687]{8.13}, and every constructible \(Q\) is built from finite direct images of constant sheaves, which the next lemma treats.

% object: lemma-kunneth-finite-direct-image (blueprint/chap-Smooth-base-change/sec-The-Kunneth-formula/lemma-kunneth-finite-direct-image/_section.md)
\phantomsection\label{bp692}%
\begin{bpstatementit}{Lemma~8.15}{ (Finite direct images of constant sheaves)}
In the situation of (8.4), let \(Q\) be a finite direct sum of sheaves \(\pi_*A_{T'}\) with \(\pi:T'\to T\) finite. Then the Künneth map \(\kappa_{M,Q}\) is an isomorphism.
\end{bpstatementit}

% proof: proof/lemma-kunneth-finite-direct-image
\begin{proof}[{Proof}]\phantomsection\label{bp693}%
Both sides are additive in \(Q\), so we may assume \(Q=\pi_*A_{T'}\). Set

\[
v=\operatorname{id}_Y\times\pi,\qquad
w=\operatorname{id}_X\times\pi,\qquad
h'=f\times\operatorname{id}_{T'}.
\]

Direct image along a finite morphism is exact and commutes with pullback, and by the projection formula we have

\[
Rf_*M\boxtimes\pi_*A_{T'}
\simeq v_*p_{Y,T'}^*Rf_*M,
\qquad
M\boxtimes\pi_*A_{T'}\simeq w_*p_{X,T'}^*M.
\]

Since \(hw=vh'\), the comparison between these two expressions is \(v_*\) of the base change map of Lemma \hyperref[bp687]{8.13} for \(h'\), hence it is an isomorphism. Moreover, this comparison is \(\kappa_{M,Q}\), because under the adjunction \(h^*\dashv Rh_*\) and the adjunctions for the finite maps, both correspond to the counit \(\varepsilon_f\otimes^{\mathbf L}1\) that defines \(\kappa_{M,Q}\).
\end{proof}

% text: prose-08-21
\phantomsection\label{bp694}%
Over a field every constructible sheaf embeds into a sheaf of this form, and this allows an induction on the cohomological degree.

% object: lemma-kunneth-field-coefficients (blueprint/chap-Smooth-base-change/sec-The-Kunneth-formula/lemma-kunneth-field-coefficients/_section.md)
\phantomsection\label{bp695}%
\begin{bpstatementit}{Lemma~8.16}{ (Field coefficients)}
In the situation of (8.4), suppose that \(A=L\) is a finite field of characteristic \(\ell\). Then the Künneth map \(\kappa_{M,Q}\) is an isomorphism.
\end{bpstatementit}

% proof: proof/lemma-kunneth-field-coefficients
\begin{proof}[{Proof}]\phantomsection\label{bp696}%
We first treat constructible sheaves \(M\) and \(Q\), by the usual argument for comparing a \(\delta\)-functor with an exact functor, and then pass to complexes by truncation. We may compute with the small étale models of \(M\) and \(Q\), since Theorems \hyperref[bp129]{2.43} and \hyperref[bp113]{Theorem 2.38} identify the operations and the adjunction maps defining \(\kappa_{M,Q}\) on the two sites.

\emph{Sheaves.} Let \(M\) and \(Q\) be constructible sheaves of \(L\)-modules. By Stacks, [09Z7](https://stacks.math.columbia.edu/tag/09Z7), applied on \(T\), there is a short exact sequence of constructible \(L\)-sheaves

\[
0\longrightarrow Q\longrightarrow P\longrightarrow Q'\longrightarrow0,
\qquad
P=\bigoplus_a\pi_{a*}L_{T_a}^{\oplus n_a},
\quad \pi_a\text{ finite}.
\]

Indeed, every constructible sheaf on a qcqs scheme embeds into a finite sum of finite direct images of constant sheaves with finite coefficients, and over \(L\) these coefficients are finite sums of \(L\). For \(r\geq0\) put

\[
F^r(Q)=p_Y^*R^rf_*M\otimes_Lp_T^*Q,
\qquad
G^r(Q)=R^rh_*(p_X^*M\otimes_Lp_T^*Q).
\]

Since tensoring over \(L\) is exact, each \(F^r\) is exact in \(Q\), while the \(G^r\) form a long exact sequence. On cohomology sheaves the Künneth map induces natural maps \(a_Q^r:F^r(Q)\to G^r(Q)\), and \(a_P^r\) is an isomorphism for every \(r\) by Lemma \hyperref[bp692]{8.15}. It follows that each \(a_Q^r\) is injective. Indeed, in its naturality square for \(Q\hookrightarrow P\) the map \(F^r(Q)\to F^r(P)\) is injective and \(a_P^r\) is an isomorphism. Surjectivity is proved by induction on \(r\), simultaneously for all constructible \(Q\). For \(r=0\), the beginning of the long exact sequence of the \(G^r\) identifies \(G^0(Q)\) with the kernel of \(G^0(P)\to G^0(Q')\), and since \(F^0\) is exact, \(F^0(Q)\) is the kernel of \(F^0(P)\to F^0(Q')\). As \(a_P^0\) is an isomorphism and \(a_{Q'}^0\) is injective, \(a^0_Q\) maps the second kernel onto the first. For \(r>0\), the induction hypothesis for \(Q'\) and the surjection \(F^{r-1}(P)\to F^{r-1}(Q')\) imply that \(G^{r-1}(P)\to G^{r-1}(Q')\) is surjective. By the long exact sequence,

\[
G^r(Q)=\ker\bigl(G^r(P)\to G^r(Q')\bigr).
\]

Since \(F^r(Q)=\ker(F^r(P)\to F^r(Q'))\) by exactness of \(F^r\), the isomorphism \(a_P^r\) and the injectivity of \(a_{Q'}^r\) identify this kernel with \(F^r(Q)\). This completes the induction, so \(\kappa_{M,Q}\) is an isomorphism for sheaves.

\emph{Complexes.} The source and target of \(\kappa_{M,Q}\) are exact functors of \(M\) and of \(Q\), and \(\kappa_{M,Q}\) is natural in both. Using the finitely many truncation triangles of \(M\) and of \(Q\) and two-out-of-three, we can therefore extend the case of sheaves to \(M\in D^b_{\mathrm{ctf}}(X,L)\) and \(Q\in D^b_{\mathrm{ctf}}(T,L)\). This is possible because over the field \(L\) the truncations stay in \(D^b_{\mathrm{ctf}}\), since constructible cohomology sheaves have Tor amplitude zero, and all operations are bounded on the schemes under consideration.
\end{proof}

% text: prose-08-22
\phantomsection\label{bp697}%
Over a general finite ring the truncations of a complex in \(D^b_{\mathrm{ctf}}\) need not have finite Tor amplitude, so the argument by truncation does not apply directly. Instead we reduce to the residue fields and use derived Nakayama.

% object: lemma-kunneth-one-identity-factor (blueprint/chap-Smooth-base-change/sec-The-Kunneth-formula/lemma-kunneth-one-identity-factor/_section.md)
\phantomsection\label{bp698}%
\begin{bpstatementit}{Lemma~8.17}{ (Künneth formula with one identity factor)}
In the situation of (8.4), the Künneth map \(\kappa_{M,Q}\) is an isomorphism.
\end{bpstatementit}

% proof: proof/lemma-kunneth-one-identity-factor
\begin{proof}[{Proof}]\phantomsection\label{bp699}%
We show that the cone \(N\) of \(\kappa_{M,Q}\) vanishes modulo every maximal ideal of \(A\), and conclude by derived Nakayama on stalks. The source and target of \(\kappa_{M,Q}\) are in \(D^b_{\mathrm{ctf}}(Y\times T,A)\) by Proposition \hyperref[bp605]{7.1}, so every geometric stalk of \(N\) is a perfect \(A\)-complex. Now let \(\mathfrak m\subset A\) be a maximal ideal and put \(k=A/\mathfrak m\). The change-of-coefficients maps of Proposition \hyperref[bp626]{7.10} for pullback, tensor product and ordinary direct image respect the counit defining the Künneth map. Hence they identify the reduction \(\kappa_{M,Q}\otimes^{\mathbf L}_A k\) with \(\kappa_{M\otimes^{\mathbf L}_Ak,\,Q\otimes^{\mathbf L}_Ak}\), which is an isomorphism by Lemma \hyperref[bp695]{8.16}. This means that \(N\otimes^{\mathbf L}_A k=0\) for every maximal ideal. Since the Artinian ring \(A\) is a product of local rings, derived Nakayama for the perfect complex \(N_{\bar z}\) at a geometric point \(\bar z\) shows that \(N_{\bar z}=0\). Finally, since \(N\) lies in \(D^b_{\mathrm{ctf}}(Y\times T,A)\), it comes from a complex on the small étale site by Theorem \hyperref[bp129]{2.43}, where vanishing is detected on geometric stalks, so \(N=0\).
\end{proof}

% object: proposition-ordinary-external-kunneth (blueprint/chap-Smooth-base-change/sec-The-Kunneth-formula/proposition-ordinary-external-kunneth/_section.md)
\phantomsection\label{bp700}%
\begin{bpstatementit}{Proposition~8.18}{ (Ordinary external Künneth)}
Let \(f_i:X_i\to Y_i\), \(i=1,2\), be morphisms of separated \(k\)-schemes of finite type, and let \(M_i\in D^b_{\mathrm{ctf}}(X_i,A)\). Then the Künneth map (8.3)

\[
Rf_{1*}M_1\boxtimes Rf_{2*}M_2
\longrightarrow R(f_1\times f_2)_*(M_1\boxtimes M_2)
\]

is an isomorphism.
\end{bpstatementit}

% proof: proof/proposition-ordinary-external-kunneth
\begin{proof}[{Proof}]\phantomsection\label{bp701}%
We apply Lemma \hyperref[bp698]{8.17} once in each factor of the factorization

\[
X_1\times X_2\xrightarrow{\operatorname{id}_{X_1}\times f_2}
X_1\times Y_2\xrightarrow{f_1\times\operatorname{id}_{Y_2}}
Y_1\times Y_2
\]

of \(f_1\times f_2\). We apply it to \(f_1\) with second coefficient \(Rf_{2*}M_2\), which lies in \(D^b_{\mathrm{ctf}}(Y_2,A)\) by Proposition \hyperref[bp605]{7.1}. After interchanging the factors, we apply it to \(f_2\) with coefficient \(M_1\). Then the composite of \(R(f_1\times\operatorname{id}_{Y_2})_*\) of the second isomorphism with the first is an isomorphism from \(Rf_{1*}M_1\boxtimes Rf_{2*}M_2\) to \(R(f_1\times f_2)_*(M_1\boxtimes M_2)\). Since its adjoint transpose is \(\varepsilon_{f_1}\boxtimes\varepsilon_{f_2}\), it is the Künneth map (8.3).
\end{proof}




\clearpage\phantomsection\hypertarget{bp-outside}{}%
\section*{Outside this export}
\addcontentsline{toc}{section}{Outside this export}
References marked \ensuremath{{}^\ast} are to parts of the blueprint this document leaves out:
\begin{description}
\item[Section 1.2] Weil cohomology, in 1 Introduction (\bpslug{sec-{}\allowbreak{}Weil-{}\allowbreak{}cohomology})
\item[Section 1.3] The function–sheaf dictionary, in 1 Introduction (\bpslug{sec-{}\allowbreak{}The-{}\allowbreak{}function-{}\allowbreak{}sheaf-{}\allowbreak{}dictionary})
\item[Section 1.4] The six operations, in 1 Introduction (\bpslug{sec-{}\allowbreak{}The-{}\allowbreak{}six-{}\allowbreak{}operations})
\item[Section 1.9] Conventions, in 1 Introduction (\bpslug{sec-{}\allowbreak{}Conventions})
\item[Chapter 9] Smooth purity, in Ⅱ Duality and the trace formula (\bpslug{chap-{}\allowbreak{}Smooth-{}\allowbreak{}purity})
\item[Theorem 9.3] Smooth purity for all admissible coefficients, in Ⅱ Duality and the trace formula (\bpslug{theorem-{}\allowbreak{}smooth-{}\allowbreak{}purity-{}\allowbreak{}all-{}\allowbreak{}admissible-{}\allowbreak{}coefficients})
\item[Section 9.5] $\ell$-adic coefficients, in Ⅱ Duality and the trace formula (\bpslug{sec-{}\allowbreak{}ell-{}\allowbreak{}adic-{}\allowbreak{}coefficients-{}\allowbreak{}4})
\item[Section 9.6] Compatibilities, in Ⅱ Duality and the trace formula (\bpslug{sec-{}\allowbreak{}Compatibilities})
\item[Chapter 10] Verdier and Poincaré duality, in Ⅱ Duality and the trace formula (\bpslug{chap-{}\allowbreak{}Verdier-{}\allowbreak{}and-{}\allowbreak{}Poincare-{}\allowbreak{}duality})
\item[Theorem 10.8] Verdier biduality, in Ⅱ Duality and the trace formula (\bpslug{theorem-{}\allowbreak{}verdier-{}\allowbreak{}biduality})
\item[Proposition 10.20] Purity and trace commute with Frobenius, in Ⅱ Duality and the trace formula (\bpslug{proposition-{}\allowbreak{}purity-{}\allowbreak{}trace-{}\allowbreak{}commute-{}\allowbreak{}frobenius})
\item[Lemma 11.7] External duality, in Ⅱ Duality and the trace formula (\bpslug{lemma-{}\allowbreak{}external-{}\allowbreak{}duality-{}\allowbreak{}trace-{}\allowbreak{}morphism})
\item[Chapter 12] The trace formula, in Ⅱ Duality and the trace formula (\bpslug{chap-{}\allowbreak{}The-{}\allowbreak{}trace-{}\allowbreak{}formula})
\item[Corollary 12.26] Trace functions and direct image with compact support, in Ⅱ Duality and the trace formula (\bpslug{corollary-{}\allowbreak{}trace-{}\allowbreak{}functions-{}\allowbreak{}direct-{}\allowbreak{}image-{}\allowbreak{}compact})
\item[Proposition 12.28] The function–sheaf dictionary, in Ⅱ Duality and the trace formula (\bpslug{proposition-{}\allowbreak{}function-{}\allowbreak{}sheaf-{}\allowbreak{}dictionary})
\item[Corollary 12.31] Point counting and the zeta function, in Ⅱ Duality and the trace formula (\bpslug{corollary-{}\allowbreak{}point-{}\allowbreak{}counting-{}\allowbreak{}zeta-{}\allowbreak{}function})
\item[Proposition 13.10] Lattices and representations on a normal scheme, in Ⅲ Weights and the Riemann hypothesis (\bpslug{proposition-{}\allowbreak{}lattices-{}\allowbreak{}representations-{}\allowbreak{}normal-{}\allowbreak{}scheme})
\item[Appendix A] Nagata compactification and Chow's lemma, in Appendices (\bpslug{chap-{}\allowbreak{}Nagata-{}\allowbreak{}compactification-{}\allowbreak{}and-{}\allowbreak{}Chows-{}\allowbreak{}lemma})
\item[Definition A.1] Compactification and refinement, in Appendices (\bpslug{definition-{}\allowbreak{}compactification-{}\allowbreak{}refinement})
\item[Theorem A.2] Nagata's compactification theorem, in Appendices (\bpslug{theorem-{}\allowbreak{}nagata-{}\allowbreak{}s-{}\allowbreak{}compactification-{}\allowbreak{}theorem})
\item[Lemma A.5] Passing to the schematic closure, in Appendices (\bpslug{lemma-{}\allowbreak{}passing-{}\allowbreak{}schematic-{}\allowbreak{}closure})
\item[Proposition A.7] Dense compactifications form a cofiltered category, in Appendices (\bpslug{proposition-{}\allowbreak{}dense-{}\allowbreak{}compactifications-{}\allowbreak{}form-{}\allowbreak{}cofiltered-{}\allowbreak{}category})
\item[Theorem A.25] Chow's lemma, projective form, in Appendices (\bpslug{theorem-{}\allowbreak{}chow-{}\allowbreak{}s-{}\allowbreak{}lemma-{}\allowbreak{}projective-{}\allowbreak{}form})
\item[Appendix C] The Picard scheme of a curve, in Appendices (\bpslug{chap-{}\allowbreak{}The-{}\allowbreak{}Picard-{}\allowbreak{}scheme-{}\allowbreak{}of-{}\allowbreak{}a-{}\allowbreak{}curve})
\item[Appendix D] Cohomology of torsion sheaves on curves, in Appendices (\bpslug{chap-{}\allowbreak{}Cohomology-{}\allowbreak{}of-{}\allowbreak{}torsion-{}\allowbreak{}sheaves-{}\allowbreak{}on-{}\allowbreak{}curves})
\item[Proposition D.2] Smooth proper curves, torsion prime to the characteristic, in Appendices (\bpslug{proposition-{}\allowbreak{}smooth-{}\allowbreak{}proper-{}\allowbreak{}curves-{}\allowbreak{}torsion-{}\allowbreak{}prime})
\end{description}
\clearpage\phantomsection%
\section*{References}
\addcontentsline{toc}{section}{References}
\bpbibitem{BBD82}{bp706}{A. Beilinson, J. Bernstein, and P. Deligne. \emph{Faisceaux pervers}, volume 100 of \emph{Astérisque}. 1982.}

\bpbibitem{BS15}{bp705}{B. Bhatt and P. Scholze. The pro-étale topology for schemes. \emph{Astérisque}, 369:99–201, 2015. \href{https://people.mpim-bonn.mpg.de/scholze/proetale.pdf}{Author's manuscript}.}

\bpbibitem{Con07}{bp707}{B. Conrad. Deligne's notes on Nagata compactifications. \emph{Journal of the Ramanujan Mathematical Society}, 22:205–257, 2007. \href{https://math.stanford.edu/\~conrad/papers/nagatafinal.pdf}{Author's manuscript}.}

\bpbibitem{Del74}{bp708}{P. Deligne. La conjecture de Weil. I. \emph{Publ. Math. IHÉS}, 43:273–307, 1974.}

\bpbibitem{Del80}{bp709}{P. Deligne. La conjecture de Weil. II. \emph{Publ. Math. IHÉS}, 52:137–252, 1980.}

\bpbibitem{Den22}{bp710}{C. Deninger. There is no “Weil-”cohomology theory with real coefficients for arithmetic curves. \emph{Ann. Sc. Norm. Super. Pisa Cl. Sci.}, 2022. To appear. arXiv:\href{https://arxiv.org/abs/2204.02714}{2204.02714}.}

\bpbibitem{Deu41}{bp711}{M. Deuring. Die Typen der Multiplikatorenringe elliptischer Funktionenkörper. \emph{Abh. Math. Sem. Univ. Hamburg}, 14:197–272, 1941.}

\bpbibitem{Eke90}{bp704}{T. Ekedahl. On the adic formalism. In \emph{The Grothendieck Festschrift, Vol. II}, volume 87 of \emph{Progress in Mathematics}, pages 197–218. Birkhäuser, 1990.}

\bpbibitem{Has36}{bp712}{H. Hasse. Zur Theorie der abstrakten elliptischen Funktionenkörper. I, II, III. \emph{J. Reine Angew. Math.}, 175:55–62, 69–88, 193–208, 1936.}

\bpbibitem{Hum72}{bp713}{J. E. Humphreys. \emph{Introduction to Lie Algebras and Representation Theory}, volume 9 of \emph{Graduate Texts in Mathematics}. Springer, 1972.}

\bpbibitem{Ill94}{bp714}{L. Illusie. Autour du théorème de monodromie locale. \emph{Astérisque}, 223:9–57, 1994. §4. \href{https://www.numdam.org/item/AST_1994__223__9_0/}{Original article}.}

\bpbibitem{Kat88}{bp715}{N. M. Katz. Travaux de Laumon. \emph{Astérisque}, 161–162:105–132, 1988. Séminaire Bourbaki, exp. 691.}

\bpbibitem{KW01}{bp716}{R. Kiehl and R. Weissauer. \emph{Weil conjectures, perverse sheaves and \(\ell\)-adic Fourier transform}, volume 42 of \emph{Ergebnisse der Mathematik}. Springer, 2001.}

\bpbibitem{Lau87}{bp717}{G. Laumon. Transformation de Fourier, constantes d'équations fonctionnelles et conjecture de Weil. \emph{Publ. Math. IHÉS}, 65:131–210, 1987.}

\bpbibitem{Mil86a}{bp718}{J. S. Milne. \href{https://www.jmilne.org/math/xnotes/AVs.pdf}{Abelian Varieties}. In \emph{Arithmetic Geometry}. 1986. Corrected online version.}

\bpbibitem{Mil86b}{bp719}{J. S. Milne. \href{https://www.jmilne.org/math/xnotes/JVs.pdf}{Jacobian Varieties}. In \emph{Arithmetic Geometry}. 1986. Corrected online version.}

\bpbibitem{Ran39}{bp720}{R. A. Rankin. Contributions to the theory of Ramanujan's function \(\tau(n)\) and similar arithmetical functions. II. \emph{Proc. Cambridge Philos. Soc.}, 35:357–372, 1939.}

\bpbibitem{SGA4}{bp702}{M. Artin, A. Grothendieck, and J.-L. Verdier, editors. \emph{Théorie des topos et cohomologie étale des schémas}, volumes 269, 270, 305 of \emph{Lecture Notes in Mathematics}. Springer, 1972–1973. SGA 4.}

\bpbibitem{SGA4½}{bp703}{P. Deligne. \emph{Cohomologie étale}, volume 569 of \emph{Lecture Notes in Mathematics}. Springer, 1977. SGA 4½; \emph{Théorèmes de finitude}, §3, and \emph{Rapport sur la formule des traces}, §3, and \emph{Sommes trigonométriques}, §1.}

\bpbibitem{SGA5}{bp721}{A. Grothendieck, editor. \emph{Cohomologie \(\ell\)-adique et fonctions \(L\)}, volume 589 of \emph{Lecture Notes in Mathematics}. Springer, 1977. SGA 5.}

\bpbibitem{SGA7II}{bp722}{P. Deligne and N. Katz, editors. \emph{Groupes de monodromie en géométrie algébrique}, volume 340 of \emph{Lecture Notes in Mathematics}. Springer, 1973. SGA 7 II; Exposé XIII, (2.1.7).}

\bpbibitem{Stacks}{bp723}{The Stacks Project Authors. \href{https://stacks.math.columbia.edu}{The Stacks Project}. Individual results are cited by their stable tags in the text.}

\bpbibitem{Var05}{bp724}{Y. Varshavsky. Lefschetz–Verdier trace formula and a generalization of a theorem of Fujiwara. 2005. Version 2. arXiv:\href{https://arxiv.org/abs/math/0505564}{math/0505564}. \href{https://arxiv.org/abs/math/0505564v2}{Author's preprint}.}

\bpbibitem{Wei48a}{bp725}{A. Weil. \emph{Sur les courbes algébriques et les variétés qui s'en déduisent}. Hermann, Paris, 1948.}

\bpbibitem{Wei48b}{bp726}{A. Weil. \emph{Variétés abéliennes et courbes algébriques}. Hermann, Paris, 1948.}

\bpbibitem{Wei49}{bp727}{A. Weil. Numbers of solutions of equations in finite fields. \emph{Bull. Amer. Math. Soc.}, 55:497–508, 1949.}

\end{document}
